% Complete standalone editable source corresponding to the EGA III-1 zh-Hans-CN reader.
% Generated deterministically from the sealed producer source; no authority scan is included.
\documentclass[UTF8,fontset=none,11pt,oneside,openany]{ctexbook}

\usepackage{iftex}
\usepackage{fontspec}
\ifXeTeX
  \usepackage{xeCJK}
\fi
\defaultfontfeatures{Ligatures=TeX}
\setmainfont{texgyretermes-regular.otf}[
  BoldFont=texgyretermes-bold.otf,
  ItalicFont=texgyretermes-italic.otf,
  BoldItalicFont=texgyretermes-bolditalic.otf
]
\setsansfont{texgyreheros-regular.otf}
\setmonofont{lmmono10-regular.otf}[ItalicFont=lmmono10-italic.otf]
\setCJKmainfont{FandolSong-Regular.otf}[
  BoldFont=FandolSong-Bold.otf,
  ItalicFont=FandolKai-Regular.otf,
  BoldItalicFont=FandolSong-Bold.otf
]
\setCJKsansfont{FandolHei-Regular.otf}
\setCJKmonofont{FandolFang-Regular.otf}
\ifXeTeX
  \xeCJKsetup{PunctStyle=kaiming,CheckSingle=true}
  \XeTeXlinebreaklocale "zh"
  \XeTeXlinebreakskip=0pt plus 1pt
\fi

\usepackage{amsmath}
\ifLuaTeX
  \usepackage{unicode-math}
  \setmathfont{texgyretermes-math.otf}
  \providecommand{\square}{\mdlgwhtsquare}
\else
  \usepackage{amssymb,mathrsfs}
\fi
\ifLuaTeX
  \renewcommand{\leadsto}{\mathrel{\rightsquigarrow}}
\fi
\usepackage{graphicx}
\usepackage{tikz-cd}
\ifLuaTeX
  \usepackage{luatex85}
\fi
\usepackage[all]{xy}
\ifLuaTeX
  % Preserve the established EGA Xy-pic stroke geometry under OpenType math.
  \makeatletter
  \AtBeginDocument{%
    \edef\xP@preclw{0.43799pt}%
    \edef\xP@lw{\xP@dim\xP@preclw}%
  }
  \makeatother
\fi
\usepackage[shortlabels]{enumitem}
\usepackage{marginnote}
\newcommand{\SourceCorrectionNote}[1]{\footnote{#1}}
\ifLuaTeX\else\usepackage{microtype}\fi
\usepackage{fancyhdr}
\usepackage[
  a4paper,
  inner=26mm,
  outer=24mm,
  top=24mm,
  bottom=24mm,
  headheight=14pt,
  headsep=8mm,
  footskip=12mm
]{geometry}
\usepackage[
  unicode=true,
  colorlinks=true,
  linkcolor=black,
  citecolor=black,
  urlcolor=blue,
  hypertexnames=false
]{hyperref}

\hypersetup{
  pdftitle={代数几何原理 III：上同调研究（大陆简体中文工作版，非典范检查点）},
  pdfauthor={Alexander Grothendieck；Jean Dieudonné（协作编写）},
  pdfsubject={依据 NUMDAM 法文印本与法文外交式 TeX 制作的大陆简体中文版本},
  pdfkeywords={代数几何；上同调；EGA；简体中文}
}

\AtBeginDocument{\fontsize{10.95pt}{16.5pt}\selectfont}
\setlength{\parindent}{2em}
\setlength{\parskip}{0pt}
\setlength{\emergencystretch}{2em}
\flushbottom

\pagestyle{fancy}
\fancyhf{}
\fancyhead[L]{\small 代数几何原理 III}
\fancyhead[R]{\small 上同调研究}
\fancyfoot[C]{\thepage}
\renewcommand{\headrulewidth}{0.4pt}
\renewcommand{\footrulewidth}{0pt}
\fancypagestyle{plain}{%
  \fancyhf{}%
  \fancyhead[L]{\small 代数几何原理 III}%
  \fancyhead[R]{\small 上同调研究}%
  \fancyfoot[C]{\thepage}%
  \renewcommand{\headrulewidth}{0.4pt}%
  \renewcommand{\footrulewidth}{0pt}%
}

\renewcommand{\contentsname}{目录}
\renewcommand{\bibname}{参考文献}
\renewcommand{\figurename}{图}
\renewcommand{\tablename}{表}
\renewcommand{\thesection}{\arabic{section}}
\renewcommand{\thesubsection}{\thesection.\arabic{subsection}}

\newenvironment{env}[1][]{%
  \par\medskip\noindent\textbf{(#1)}\ %
}{\par}
\newenvironment{definition}[1][]{%
  \par\medskip\noindent\textbf{定义 (#1).}\ %
}{\par}
\newenvironment{lemma}[1][]{%
  \par\medskip\noindent\textbf{引理 (#1).}\ %
}{\par}
\newenvironment{theorem}[1][]{%
  \par\medskip\noindent\textbf{定理 (#1).}\ %
}{\par}
\newenvironment{proposition}[1][]{%
  \par\medskip\noindent\textbf{命题 (#1).}\ %
}{\par}
\newenvironment{corollary}[1][]{%
  \par\medskip\noindent\textbf{推论 (#1).}\ %
}{\par}
\newenvironment{remark}[1][]{%
  \par\medskip\noindent\textbf{注 (#1).}\ %
}{\par}
\newenvironment{remarks}[1][]{%
  \par\medskip\noindent\textbf{注 (#1).}\ %
}{\par}
\newenvironment{notation}[1][]{%
  \par\medskip\noindent\textbf{记号 (#1).}\ %
}{\par}
\newenvironment{proof}{%
  \par\smallskip\noindent\textit{证明.}\ %
}{\hfill$\square$\par}

\newcommand{\oldpage}[2][]{%
  \hypertarget{ega-source-page-#1-#2}{}%
  \marginnote{\scriptsize\textbf{#1}~|~#2}\ignorespaces}

\let\EGAOriginalLabel\label
\renewcommand{\label}[1]{\phantomsection\EGAOriginalLabel{#1}}

\begin{document}
\setcounter{page}{349}
% EGA III-1 mainland Simplified-Chinese producer source.
% Complete through the sealed terminal boundary; producer output remains uncertified pending independent canon replay.
% The underlying EGA text is not assigned a new licence, and the NUMDAM authority scans are not redistributed here.
\clearpage
\thispagestyle{plain}
\markboth{预备知识}{预备知识}
\oldpage[0\textsubscript{III}]{5}
\phantomsection
\label{chapter:0III-fr}

\begin{center}
{\large 第 0 章（续）\footnotemark\par}
\vspace{1.5em}
{\Large\bfseries 预备知识\par}
\end{center}

\footnotetext{为便于查找参考文献，从现在起，凡引用随第 I 章发表的第 0 章各节，均以符号 $0_{\mathrm I}$ 标示。}

\setcounter{section}{7}
\section{可表函子}
\label{section:0.8-fr}

\setcounter{subsection}{0}
\subsection{可表函子}
\label{subsection:0.8.1-fr}

\begin{env}[8.1.1]
\label{0.8.1.1-fr}
我们用 $\mathrm{Ens}$ 表示集合范畴。设 $C$ 是一个范畴；对 $C$ 的两个对象 $X$、$Y$，置
$h_X(Y)=\operatorname{Hom}(Y,X)$；对 $C$ 中任一态射 $u:Y\to Y'$，以 $h_X(u)$ 表示从
$\operatorname{Hom}(Y',X)$ 到 $\operatorname{Hom}(Y,X)$ 的映射 $v\mapsto vu$。由这些定义立即可见，
$h_X:C\to\mathrm{Ens}$ 是一个\emph{反变函子}，也就是说，它是范畴
$\operatorname{Hom}(C^0,\mathrm{Ens})$ 的一个对象；该范畴由范畴 $C$ 的对偶范畴 $C^0$ 到集合范畴
$\mathrm{Ens}$ 的协变函子组成（T, 1.7, d) 与 [29]）。
\end{env}
\begin{env}[8.1.2]
\label{0.8.1.2-fr}
现在设 $w:X\to X'$ 是 $C$ 中的一个态射；对每个 $Y\in C$ 以及每个
$v\in\operatorname{Hom}(Y,X)=h_X(Y)$，有
$wv\in\operatorname{Hom}(Y,X')=h_{X'}(Y)$；以 $h_w(Y)$ 表示从 $h_X(Y)$ 到
$h_{X'}(Y)$ 的映射 $v\mapsto wv$。立即可见，对 $C$ 中任一态射 $u:Y\to Y'$，图表
\[
  \xymatrix{
    h_X(Y')\ar[r]^{h_X(u)}\ar[d]_{h_w(Y')} &
    h_X(Y)\ar[d]^{h_w(Y)}\\
    h_{X'}(Y')\ar[r]_{h_{X'}(u)} &
    h_{X'}(Y)
  }
\]
是交换的；换言之，$h_w$ 是一个\emph{函子性态射（自然变换）}
$h_X\to h_{X'}$（T, 1.2），也就是范畴
$\operatorname{Hom}(C^0,\mathrm{Ens})$ 中的一个态射（T, 1.7, d)）。因此，$h_X$ 与
$h_w$ 的定义共同给出了一个\emph{典范协变函子}
\[
  h:C\to\operatorname{Hom}(C^0,\mathrm{Ens}),\qquad X\mapsto h_X.
  \tag{8.1.2.1}
\]
\end{env}
\begin{env}[8.1.3]
\label{0.8.1.3-fr}
设 $X$ 是 $C$ 的一个对象，$F$ 是从 $C$ 到 $\mathrm{Ens}$ 的一个反变函子（即
$\operatorname{Hom}(C^0,\mathrm{Ens})$ 的一个对象）。设 $g:h_X\to F$ 是一个
\emph{函子性态射}：于是对每个 $Y\in C$，
\oldpage[0\textsubscript{III}]{6}
$g(Y)$ 是一个映射 $h_X(Y)\to F(Y)$，并且对 $C$ 中每个态射 $u:Y\to Y'$，图表
\[
  \xymatrix{
    h_X(Y')\ar[r]^{h_X(u)}\ar[d]_{g(Y')} &
    h_X(Y)\ar[d]^{g(Y)}\\
    F(Y')\ar[r]_{F(u)} & F(Y)
  }
  \tag{8.1.3.1}\label{0.8.1.3.1-fr}
\]
是交换的。特别地，有映射
$g(X):h_X(X)=\operatorname{Hom}(X,X)\to F(X)$，从而得到元素
\[
  \alpha(g)=(g(X))(1_X)\in F(X)
  \tag{8.1.3.2}
\]
以及典范映射
\[
  \alpha:\operatorname{Hom}(h_X,F)\to F(X).
  \tag{8.1.3.3}
\]
反过来，考虑一个元素 $\xi\in F(X)$；对 $C$ 中每个态射 $v:Y\to X$，$F(v)$ 是映射
$F(X)\to F(Y)$；考虑映射
\[
  v\mapsto(F(v))(\xi)
  \tag{8.1.3.4}
\]
从 $h_X(Y)$ 到 $F(Y)$。若以 $(\beta(\xi))(Y)$ 表示这一映射，则
\[
  \beta(\xi):h_X\to F
  \tag{8.1.3.5}
\]
是一个\emph{函子性态射}，因为对 $C$ 中任一态射 $u:Y\to Y'$，有
% Canon current rebase 2026-08-31: EGA-III1-8.1.3-P350-AUTHORIAL-CONTRAVARIANT-COMPOSITION-001.
$(F(vu))(\xi)=(F(u)\circ F(v))(\xi)$；这就验证了当 $g=\beta(\xi)$ 时
(\hyperref[0.8.1.3.1-fr]{8.1.3.1}) 的交换性。这样便定义了典范映射
\[
  \beta:F(X)\to\operatorname{Hom}(h_X,F).
  \tag{8.1.3.6}
\]
\end{env}
\begin{proposition}[8.1.4]
\label{0.8.1.4-fr}
映射 $\alpha$ 与 $\beta$ 互为逆双射。
\end{proposition}
对 $\xi\in F(X)$ 计算 $\alpha(\beta(\xi))$；对每个 $Y\in C$，$(\beta(\xi))(Y)$ 是从
$h_X(Y)$ 到 $F(Y)$ 的映射 $g_1(Y):v\mapsto(F(v))(\xi)$。因此
\[
  \alpha(\beta(\xi))=(g_1(X))(1_X)=(F(1_X))(\xi)
  =1_{F(X)}(\xi)=\xi.
\]
取 $g\in\operatorname{Hom}(h_X,F)$。现在计算 $\beta(\alpha(g))$；\par\noindent
对每个 $Y\in C$，
$(\beta(\alpha(g)))(Y)$ 是映射
$v\mapsto(F(v))((g(X))(1_X))$；由
(\hyperref[0.8.1.3.1-fr]{8.1.3.1}) 的交换性，这一映射无非是
$v\mapsto(g(Y))((h_X(v))(1_X))=(g(Y))(v)$，其中最后一个等式来自 $h_X(v)$ 的定义。
换言之，它等于 $g(Y)$，命题得证。
\begin{env}[8.1.5]
\label{0.8.1.5-fr}
回忆：范畴 $C$ 的一个\emph{子范畴} $C'$，由以下条件定义：它的对象是 $C$ 的对象；若
$X'$、$Y'$ 是 $C'$ 的两个对象，则 $C'$ \emph{中}态射 $X'\to Y'$ 的集合
$\operatorname{Hom}_{C'}(X',Y')$ 是 $C$ \emph{中}态射 $X'\to Y'$ 的集合
$\operatorname{Hom}_C(X',Y')$ 的一个子集，并且典范的“态射复合”映射
\[
  \operatorname{Hom}_{C'}(X',Y')\times\operatorname{Hom}_{C'}(Y',Z')
  \to\operatorname{Hom}_{C'}(X',Z')
\]
\oldpage[0\textsubscript{III}]{7}
是典范映射
\[
  \operatorname{Hom}_C(X',Y')\times\operatorname{Hom}_C(Y',Z')
  \to\operatorname{Hom}_C(X',Z').
\]
的限制。
若对 $C'$ 的每一对对象都有
$\operatorname{Hom}_{C'}(X',Y')=\operatorname{Hom}_C(X',Y')$，就称 $C'$ 是 $C$ 的一个
\emph{全子范畴}。此时，由 $C$ 中与 $C'$ 的对象同构的全部对象组成的 $C$ 的子范畴 $C''$，仍是
$C$ 的一个全子范畴，而且不难验证，它作为范畴与 $C'$ \emph{等价}（T, 1.2）。
若对 $C_1$ 的每一对对象 $X_1$、$Y_1$，从
$\operatorname{Hom}(X_1,Y_1)$ 到
$\operatorname{Hom}(F(X_1),F(Y_1))$ 的映射 $u\mapsto F(u)$ 都是\emph{双射}，就称协变函子
$F:C_1\to C_2$ \emph{全忠实}；这蕴含 $C_2$ 的子范畴 $F(C_1)$ 是\emph{全的}。此外，若两个对象
$X_1$、$X'_1$ 有同一个像 $X_2$，则存在唯一同构 $u:X_1\to X'_1$，使得
$F(u)=1_{X_2}$。对 $F(C_1)$ 的每个对象 $X_2$，取 $C_1$ 中满足 $F(X_1)=X_2$ 的某一个对象
$X_1$，并把它记作 $G(X_2)$（这里借助选择公理定义 $G$）；对 $F(C_1)$ 中每个态射
$v:X_2\to Y_2$，令 $G(v)$ 为满足 $F(u)=v$ 的唯一态射
$u:G(X_2)\to G(Y_2)$。于是 $G$ 是从 $F(C_1)$ 到 $C_1$ 的一个\emph{函子}；$FG$ 是
$F(C_1)$ 上的恒等函子，而上面的论证表明存在函子同构
$\varphi:1_{C_1}\to GF$，使得 $F$、$G$、$\varphi$ 以及恒等变换
$1_{F(C_1)}\to FG$ 给出范畴 $C_1$ 与 $C_2$ 的全子范畴 $F(C_1)$ 的一个\emph{等价}
（T, 1.2）。
\end{env}
\begin{env}[8.1.6]
\label{0.8.1.6-fr}
把命题 (\hyperref[0.8.1.4-fr]{8.1.4}) 应用于函子 $F$ 为 $h_{X'}$ 的情形，其中 $X'$ 是 $C$ 的任一对象；
此时，映射
$\beta:\operatorname{Hom}(X,X')\to\operatorname{Hom}(h_X,h_{X'})$ 正是
(\hyperref[0.8.1.2-fr]{8.1.2}) 中定义的映射 $w\mapsto h_w$。由于这一映射是\emph{双射}，用
(\hyperref[0.8.1.5-fr]{8.1.5}) 的术语便得到：
\end{env}
\begin{proposition}[8.1.7]
\label{0.8.1.7-fr}
典范函子
$h:C\to\operatorname{Hom}(C^0,\mathrm{Ens})$ 是全忠实的。
\end{proposition}
\begin{env}[8.1.8]
\label{0.8.1.8-fr}
设 $F$ 是从 $C$ 到 $\mathrm{Ens}$ 的一个反变函子；若存在对象 $X\in C$，使 $F$ 与
$h_X$ \emph{同构}，就称 $F$ \emph{可表}。由
(\hyperref[0.8.1.7-fr]{8.1.7}) 可知，给定 $X\in C$ 以及函子同构
$g:h_X\to F$，便在唯一同构意义下确定 $X$。命题
(\hyperref[0.8.1.7-fr]{8.1.7}) 也意味着：$h$ 给出 $C$ 与
$\operatorname{Hom}(C^0,\mathrm{Ens})$ 中由\emph{反变可表函子}组成的全子范畴之间的一个
\emph{等价}。另一方面，由命题 (\hyperref[0.8.1.4-fr]{8.1.4}) 可知，给定函子性态射
$g:h_X\to F$，等价于给定元素 $\xi\in F(X)$；$g$ 是\emph{同构}，又等价于 $\xi$ 满足如下条件：
\emph{对 $C$ 的每个对象 $Y$，从 $\operatorname{Hom}(Y,X)$ 到 $F(Y)$ 的映射
$v\mapsto(F(v))(\xi)$ 都是双射}。当 $\xi$ 满足这一条件时，称二元组 $(X,\xi)$
\emph{表示}可表函子 $F$。滥用语言时，若存在 $\xi\in F(X)$ 使 $(X,\xi)$ 表示 $F$，也就是说，
若 $h_X$ 与 $F$ 同构，也称对象 $X\in C$ 表示 $F$。
设 $F$、$F'$ 是从 $C$ 到 $\mathrm{Ens}$ 的两个反变可表函子，且
$h_X\to F$、$h_{X'}\to F'$ 是两个函子同构。于是由
(\hyperref[0.8.1.6-fr]{8.1.6}) 可知，$\operatorname{Hom}(X,X')$ 与函子性态射 $F\to F'$ 的集合
$\operatorname{Hom}(F,F')$ 之间存在典范双射对应。
\end{env}
\begin{env}[8.1.9]
\label{0.8.1.9-fr}
\textit{例 I：逆极限。} 反变可表函子的概念尤其包含通常所谓“泛问题的解”这一概念的
“对偶”概念。
\oldpage[0\textsubscript{III}]{8}
更一般地，我们将看到，\emph{逆极限}的概念是可表函子概念的一种特殊情形。回忆：在范畴
$C$ 中，一个\emph{逆系统}由以下数据定义：一个预序集 $I$，$C$ 的一族对象
$(A_\alpha)_{\alpha\in I}$，以及对满足 $\alpha\leq\beta$ 的每一对指标
$(\alpha,\beta)$，一个态射 $u_{\alpha\beta}:A_\beta\to A_\alpha$。该系统在 $C$ 中的一个
\emph{逆极限}，由 $C$ 的一个对象 $B$（记作 $\varprojlim A_\alpha$）以及对每个
$\alpha\in I$ 的一个态射 $u_\alpha:B\to A_\alpha$ 组成，并满足：
1\textsuperscript{o} 当 $\alpha\leq\beta$ 时，$u_\alpha=u_{\alpha\beta}u_\beta$；
2\textsuperscript{o} 对 $C$ 的任一对象 $X$ 以及任一族态射
$(v_\alpha)_{\alpha\in I}$，其中 $v_\alpha:X\to A_\alpha$ 且当
$\alpha\leq\beta$ 时 $v_\alpha=u_{\alpha\beta}v_\beta$，存在唯一态射 $v:X\to B$
（记作 $\varprojlim v_\alpha$），使得对每个 $\alpha\in I$ 都有 $v_\alpha=u_\alpha v$
（T, 1.8）。这可解释如下：$u_{\alpha\beta}$ 典范地定义映射
\[
  \overline{u}_{\alpha\beta}:\operatorname{Hom}(X,A_\beta)
  \to\operatorname{Hom}(X,A_\alpha)
\]
它们给出一个集合的\emph{逆系统}
$(\operatorname{Hom}(X,A_\alpha),\overline{u}_{\alpha\beta})$，而
$(v_\alpha)$ 按定义是集合
$\varprojlim_\alpha\operatorname{Hom}(X,A_\alpha)$ 的一个元素。显然，
$X\mapsto\varprojlim_\alpha\operatorname{Hom}(X,A_\alpha)$ 是从 $C$ 到
$\mathrm{Ens}$ 的一个\emph{反变函子}，而逆极限 $B$ 的存在性等价于说，
$(v_\alpha)\mapsto\varprojlim v_\alpha$ 是关于 $X$ 的\emph{函子同构}
\[
  \varprojlim_\alpha\operatorname{Hom}(X,A_\alpha)
  \xrightarrow{\sim}\operatorname{Hom}(X,B)
  \tag{8.1.9.1}
\]
换言之，函子
$X\mapsto\varprojlim_\alpha\operatorname{Hom}(X,A_\alpha)$ 是\emph{可表的}。
\end{env}
\begin{env}[8.1.10]
\label{0.8.1.10-fr}
\textit{例 II：终对象。} 设 $C$ 是一个范畴，$\{a\}$ 是一个单元素集。考虑反变函子
$F:C\to\mathrm{Ens}$：它把 $C$ 的每个对象 $X$ 映到集合 $\{a\}$，并把 $C$ 中每个态射
$X\to X'$ 映到唯一的映射 $\{a\}\to\{a\}$。说这个函子\emph{可表}，就是指存在对象
$e\in C$，使得对每个 $Y\in C$，$\operatorname{Hom}(Y,e)=h_e(Y)$ 都是\emph{单元素集}；
称 $e$ 为 $C$ 的一个\emph{终对象}。显然，$C$ 的任意两个终对象都同构（因此一般可借助选择公理选定
$C$ 的一个终对象，并把它记作 $e_C$）。例如，在范畴 $\mathrm{Ens}$ 中，终对象就是单元素集；在域
$K$ 上的\emph{增广代数}范畴中（其态射是与增广相容的代数同态），$K$ 是终对象；在
\emph{$S$-预概形}范畴中（I, 2.5.1），$S$ 是终对象。
\end{env}
\begin{env}[8.1.11]
\label{0.8.1.11-fr}
对范畴 $C$ 的两个对象 $X$、$Y$，置
$h'_X(Y)=\operatorname{Hom}(X,Y)$；对任一态射 $u:Y\to Y'$，令 $h'_X(u)$ 为从
$\operatorname{Hom}(X,Y)$ 到 $\operatorname{Hom}(X,Y')$ 的映射 $v\mapsto uv$。于是
$h'_X$ 是一个\emph{协变函子} $C\to\mathrm{Ens}$，从而可像
(\hyperref[0.8.1.2-fr]{8.1.2}) 那样定义典范协变函子
$h':C^0\to\operatorname{Hom}(C,\mathrm{Ens})$。此时，从 $C$ 到 $\mathrm{Ens}$ 的一个
\emph{协变函子} $F$，换言之，$\operatorname{Hom}(C,\mathrm{Ens})$ 的一个对象，称为
\emph{可表的}，是指存在对象 $X\in C$（它在唯一同构意义下必然唯一），使 $F$
\emph{与 $h'_X$ 同构}。我们把上述论述对此概念的“对偶”形式留给读者自行展开；它此时包含
\emph{归纳极限}的概念，尤其包含通常所谓“泛问题的解”这一概念。
\end{env}
\subsection{范畴中的代数结构}
\label{subsection:0.8.2-fr}

\begin{env}[8.2.1]
\label{0.8.2.1-fr}
\oldpage[0\textsubscript{III}]{9}
给定从 $C$ 到 $\mathrm{Ens}$ 的两个反变函子 $F$、$F'$，回忆：对每个对象 $Y\in C$，置
$(F\times F')(Y)=F(Y)\times F'(Y)$；对 $C$ 中每个态射 $u:Y\to Y'$，置
$(F\times F')(u)=F(u)\times F'(u)$，即从 $F(Y')\times F'(Y')$ 到
$F(Y)\times F'(Y)$ 的映射
$(t,t')\mapsto((F(u))(t),(F'(u))(t'))$。因此，$F\times F':C\to\mathrm{Ens}$ 是一个
\emph{反变函子}（事实上，它正是范畴
$\operatorname{Hom}(C^0,\mathrm{Ens})$ 中对象 $F$、$F'$ 的\emph{积}）。给定对象 $X\in C$，我们把一个函子性态射
\[
  \gamma_X:h_X\times h_X\to h_X.
  \tag{8.2.1.1}
\]
称为 $X$ 上的一个\emph{内合成律}。

换言之（T, 1.2），对每个对象 $Y\in C$，$\gamma_X(Y)$ 是映射
$h_X(Y)\times h_X(Y)\to h_X(Y)$（因而按定义是集合 $h_X(Y)$ 上的一个
\emph{内合成律}），并满足：对 $C$ 中每个态射 $u:Y\to Y'$，图表
\[
  \xymatrix{
    h_X(Y')\times h_X(Y')
      \ar[rr]^{h_X(u)\times h_X(u)}\ar[dd]_{\gamma_X(Y')} &&
    h_X(Y)\times h_X(Y)\ar[dd]^{\gamma_X(Y)}\\\\
    h_X(Y')\ar[rr]_{h_X(u)} && h_X(Y)
  }
\]
是交换的；这意味着，对合成律 $\gamma_X(Y)$ 和 $\gamma_X(Y')$ 而言，$h_X(u)$ 是从
$h_X(Y')$ 到 $h_X(Y)$ 的一个\emph{同态}。

同样，给定 $C$ 的两个对象 $Z$、$X$，我们把函子性态射
\[
  \omega_{X,Z}:h_Z\times h_X\to h_X.
  \tag{8.2.1.2}
\]
称为\emph{$X$ 上以 $Z$ 为算子域的外合成律}。

如上可见，对每个 $Y\in C$，$\omega_{X,Z}(Y)$ 是 $h_X(Y)$ 上以 $h_Z(Y)$ 为算子域的一个外合成律，
并且对每个态射 $u:Y\to Y'$，$h_X(u)$ 与 $h_Z(u)$ 构成从
$(h_Z(Y'),h_X(Y'))$ 到 $(h_Z(Y),h_X(Y))$ 的一个\emph{双同态}。
\end{env}

\begin{env}[8.2.2]
\label{0.8.2.2-fr}
设 $X'$ 是 $C$ 的另一个对象，并设在 $X'$ 上给定了内合成律 $\gamma_{X'}$。若对每个
$Y\in C$，$h_w(Y):h_X(Y)\to h_{X'}(Y)$ 都是关于合成律 $\gamma_X(Y)$ 与
$\gamma_{X'}(Y)$ 的同态，就称 $C$ 中的态射 $w:X\to X'$ 是关于这些合成律的一个
\emph{同态}。若 $X''$ 是第三个
\oldpage[0\textsubscript{III}]{10}
$C$ 的对象，带有内合成律 $\gamma_{X''}$，且 $w':X'\to X''$ 是 $C$ 中关于
$\gamma_{X'}$ 与 $\gamma_{X''}$ 的同态，则显然，态射
$w'w:X\to X''$ 是关于合成律 $\gamma_X$ 与 $\gamma_{X''}$ 的同态。若同构
$w:X\xrightarrow{\sim}X'$ 在 $C$ 中满足：$w$ 关于这些合成律是同态，且其逆态射 $w^{-1}$ 关于合成律
$\gamma_{X'}$ 与 $\gamma_X$ 也是同态，就称它为
\emph{关于合成律 $\gamma_X$ 与 $\gamma_{X'}$ 的同构}。

对于带有外合成律的 $C$ 中对象偶，也以同样方式定义\emph{双同态}。
\end{env}

\begin{env}[8.2.3]
\label{0.8.2.3-fr}
若对象 $X\in C$ 上的内合成律 $\gamma_X$ 使得对\emph{每个} $Y\in C$，$\gamma_X(Y)$ 都是
$h_X(Y)$ 上的\emph{群律}，就称带有此律的 $X$ 为一个\emph{$C$-群}，或一个
\emph{$C$ 中的群对象}。同样可定义\emph{$C$-环}、\emph{$C$-模}等。
\end{env}

\begin{env}[8.2.4]
\label{0.8.2.4-fr}
设对象 $X\in C$ 与自身的\emph{积} $X\times X$ 在 $C$ 中存在。于是按定义，在典范同构意义下有
$h_{X\times X}=h_X\times h_X$，因为这是逆极限的一个特殊情形
（\hyperref[0.8.1.9-fr]{8.1.9}）。因此，$X$ 上的一个内合成律可视为函子性态射
$\gamma_X:h_{X\times X}\to h_X$，从而典范地确定
（\hyperref[0.8.1.6-fr]{8.1.6}）一个元素
$c_X\in\operatorname{Hom}(X\times X,X)$，使得 $h_{c_X}=\gamma_X$。所以在这种情形下，给定
$X$ 上的一个内合成律，等价于给定一个态射 $X\times X\to X$。当 $C$ 是范畴
$\mathrm{Ens}$ 时，这就回到集合上内合成律的经典概念。当积 $Z\times X$ 在 $C$ 中存在时，
外合成律也有类似结果。
\end{env}

\begin{env}[8.2.5]
\label{0.8.2.5-fr}
沿用上述记号，并进一步假设 $X\times X\times X$ 在 $C$ 中存在。把积刻画为表示某函子的对象
（\hyperref[0.8.1.9-fr]{8.1.9}），便得到典范同构
\[
  (X\times X)\times X\xrightarrow{\sim}X\times X\times X
  \xrightarrow{\sim}X\times(X\times X)~;
\]
若典范地把 $X\times X\times X$ 与 $(X\times X)\times X$ 认同，则对每个 $Y\in C$，映射
$\gamma_X(Y)\times 1_{h_X(Y)}$ 与 $h_{c_X\times 1_X}(Y)$ 相认同。因此，下列说法等价：对每个
$Y\in C$，内合成律 $\gamma_X(Y)$ 是\emph{结合的}；映射图表
\[
  \xymatrix{
    h_X(Y)\times h_X(Y)\times h_X(Y)
      \ar[rr]^{\gamma_X(Y)\times 1}\ar[dd]_{1\times\gamma_X(Y)} &&
    h_X(Y)\times h_X(Y)\ar[dd]^{\gamma_X(Y)}\\\\
    h_X(Y)\times h_X(Y)\ar[rr]_{\gamma_X(Y)} && h_X(Y)
  }
\]
\oldpage[0\textsubscript{III}]{11}
是交换的；态射图表
\[
  \xymatrix{
    X\times X\times X
      \ar[rr]^{c_X\times 1_X}\ar[dd]_{1_X\times c_X} &&
    X\times X\ar[dd]^{c_X}\\\\
    X\times X\ar[rr]_{c_X} && X
  }
\]
是交换的。
\end{env}

\begin{env}[8.2.6]
\label{0.8.2.6-fr}
在（\hyperref[0.8.2.5-fr]{8.2.5}）的假设下，若要表达对每个 $Y\in C$，内合成律
$\gamma_X(Y)$ 都是一个\emph{群律}，一方面必须表达它是结合的；另一方面必须表达存在映射
$\alpha_X(Y):h_X(Y)\to h_X(Y)$，具有群中\emph{取逆}的性质。由于已知，对 $C$ 中每个态射
$u:Y\to Y'$，$h_X(u)$ 必须是群同态 $h_X(Y')\to h_X(Y)$，首先可见
$\alpha_X:h_X\to h_X$ 必须是一个\emph{函子性态射}。另一方面，无须引入单位元，也可表达群
$G$ 中逆元映射 $s\mapsto s^{-1}$ 的特征性质：只须写出两个复合映射
\[
  (s,t)\mapsto(s,s^{-1},t)\mapsto(s,s^{-1}t)\mapsto s(s^{-1}t),
\]
\[
  (s,t)\mapsto(s,s^{-1},t)\mapsto(s,ts^{-1})\mapsto(ts^{-1})s
\]
都等于从 $G\times G$ 到 $G$ 的第二投影 $(s,t)\mapsto t$。由
（\hyperref[0.8.1.3-fr]{8.1.3}），有 $\alpha_X=h_{a_X}$，其中
$a_X\in\operatorname{Hom}(X,X)$；于是上述第一个条件表达复合态射
\[
  X\times X\xrightarrow{(1_X,a_X)\times 1_X}X\times X\times X
  \xrightarrow{1_X\times c_X}X\times X\xrightarrow{c_X}X
\]
就是 $C$ 中的第二投影 $X\times X\to X$；第二个条件可同样转述。
\end{env}

\begin{env}[8.2.7]
\label{0.8.2.7-fr}
现在假设 $C$ 中存在一个\emph{终对象} $e$
（\hyperref[0.8.1.10-fr]{8.1.10}）。仍假设对每个 $Y\in C$，$\gamma_X(Y)$ 都是
$h_X(Y)$ 上的一个群律，并以 $\eta_X(Y)$ 表示 $\gamma_X(Y)$ 的单位元。由于对 $C$ 中每个态射
$u:Y\to Y'$，$h_X(u)$ 都是群同态，故有
$\eta_X(Y)=(h_X(u))(\eta_X(Y'))$。特别取 $Y'=e$；此时 $u$ 是
$\operatorname{Hom}(Y,e)$ 的唯一元素 $\varepsilon$，可见元素 $\eta_X(e)$ 完全确定每个
$Y\in C$ 的 $\eta_X(Y)$。置 $e_X=\eta_X(X)$，它是群
$h_X(X)=\operatorname{Hom}(X,X)$ 的单位元；图表
\[
  \xymatrix{
    h_X(e)\ar[r]^{h_X(\varepsilon)}\ar[d]_{h_{e_X}(e)} &
    h_X(Y)\ar[d]^{h_{e_X}(Y)}\\
    h_X(e)\ar[r]_{h_X(\varepsilon)} & h_X(Y)
  }
\]
的交换性（参见 \hyperref[0.8.1.2-fr]{8.1.2}）表明，在集合 $h_X(Y)$ 上，映射
$h_{e_X}(Y)$ 正是
\oldpage[0\textsubscript{III}]{12}
把每个元素都变为单位元的映射 $s\mapsto\eta_X(Y)$。由此可验证，对每个 $Y\in C$，
$\eta_X(Y)$ 都是 $\gamma_X(Y)$ 的单位元，等价于说复合态射
\[
  X\xrightarrow{(1_X,1_X)}X\times X
  \xrightarrow{1_X\times e_X}X\times X\xrightarrow{c_X}X,
\]
以及交换 $1_X$ 与 $e_X$ 所得的类似复合态射，都\emph{等于} $1_X$。
\end{env}

\begin{env}[8.2.8]
\label{0.8.2.8-fr}
当然，不难列举更多范畴中的代数结构。这里已相当详细地处理了群的例子；不过以后遇到代数结构的例子时，
我们通常会把展开类似论述的工作留给读者。
\end{env}

\section{可构造集}
\label{section:0.9-fr}

\subsection{可构造集}
\label{subsection:0.9.1-fr}

\begin{definition}[9.1.1]
\label{0.9.1.1-fr}
若连续映射 $f:X\to Y$ 使得对 $Y$ 的每个拟紧开集 $U$，$f^{-1}(U)$ 都拟紧，就称它是
\emph{拟紧的}。若拓扑空间 $X$ 的子集 $Z$ 的典范包含映射 $Z\to X$ 是拟紧的，就称它
\emph{在 $X$ 中逆紧}；换言之，对 $X$ 的每个拟紧开集 $U$，$U\cap Z$ 都拟紧。
\end{definition}

$X$ 的\emph{闭}子集在 $X$ 中逆紧，但 $X$ 的拟紧子集未必在 $X$ 中逆紧。若
$X$ 拟紧，则 $X$ 中每个逆紧开集都拟紧。显然，$X$ 中逆紧集的每个\emph{有限}并仍在
$X$ 中逆紧，因为拟紧集的有限并是拟紧的。$X$ 中逆紧开集的每个有限交仍是 $X$ 中的逆紧开集。
在\emph{局部 Noether} 空间 $X$ 中，每个拟紧集都是 Noether 子空间；因此，$X$ 的
\emph{每个子集}都在 $X$ 中逆紧。

\begin{definition}[9.1.2]
\label{0.9.1.2-fr}
给定拓扑空间 $X$。若 $X$ 的一个子集属于 $X$ 的如下最小子集族
$\mathfrak{F}$：它包含 $\mathfrak{F}$ 的所有逆紧开子集，并且对有限交与取补封闭
（这蕴含 $\mathfrak{F}$ 也对有限并封闭），就称该子集是\emph{可构造的}。
% EGA3-ZH-DEF-0002: printed self-reference to F retained in this noncanonical candidate.
\end{definition}

\begin{proposition}[9.1.3]
\label{0.9.1.3-fr}
$X$ 的一个子集可构造，当且仅当它是有限多个形如 $U\cap\complement{V}$ 的集合之并，其中
$U$、$V$ 是 $X$ 中的逆紧开集。
\end{proposition}

条件的充分性是显然的。为证明必要性，考虑由形如 $U\cap\complement{V}$ 的集合的有限并所组成的集合
$\mathfrak{G}$，其中 $U$、$V$ 是 $X$ 中的逆紧开集；只须证明
$\mathfrak{G}$ 中每个集合的补集仍属于 $\mathfrak{G}$。于是设
$Z=\bigcup_{i\in I}(U_i\cap\complement{V_i})$，其中 $I$ 有限，$U_i$、$V_i$ 是
$X$ 中的逆紧开集。我们有
$\complement{Z}=\bigcap_{i\in I}(V_i\cup\complement{U_i})$；因此，
$\complement{Z}$ 是有限多个集合之并，而每个这样的集合是若干个 $V_i$ 与若干个
$\complement{U_i}$ 的交，因而形如
\oldpage[0\textsubscript{III}]{13}
$V\cap\complement{U}$，其中 $U$ 是若干个 $U_i$ 的并，$V$ 是若干个 $V_i$ 的交。上面已经指出，
$X$ 中逆紧开集的有限并与有限交仍是 $X$ 中的逆紧开集，结论由此得到。

\begin{corollary}[9.1.4]
\label{0.9.1.4-fr}
$X$ 的每个可构造子集都在 $X$ 中逆紧。
\end{corollary}


只须证明：若 $U$、$V$ 是 $X$ 中的逆紧开集，则 $U\cap\complement{V}$ 在 $X$ 中逆紧。事实上，
若 $W$ 是 $X$ 中的拟紧开集，则 $W\cap U\cap\complement{V}$ 是拟紧空间
$W\cap U$ 的闭子集，因而拟紧。

特别地：

\begin{corollary}[9.1.5]
\label{0.9.1.5-fr}
$X$ 的开子集 $U$ 可构造，当且仅当它在 $X$ 中逆紧。$X$ 的闭子集 $F$ 可构造，当且仅当开集
$\complement{F}$ 逆紧。
\end{corollary}

\begin{env}[9.1.6]
\label{0.9.1.6-fr}
一种重要情形是：$X$ 的每个拟紧开子集都逆紧；换言之，$X$ 的任意两个拟紧开子集之交都拟紧
（参见 \hyperref[I.5.5.6-fr]{I, 5.5.6}）。若 $X$ 本身拟紧，这就是说，$X$ 中的逆紧开子集恰好是
$X$ 的拟紧开子集，而 $X$ 的可构造子集恰好是有限多个形如
$U\cap\complement{V}$ 的集合之并，其中 $U$、$V$ 是拟紧开集。
\end{env}

\begin{corollary}[9.1.7]
\label{0.9.1.7-fr}
Noether 空间 $X$ 的一个子集可构造，当且仅当它是 $X$ 的有限多个局部闭子集之并。
\end{corollary}

\begin{proposition}[9.1.8]
\label{0.9.1.8-fr}
设 $X$ 是拓扑空间，$U$ 是 $X$ 的一个开子集。
\begin{enumerate}
  \item[(i)] 若 $T$ 是 $X$ 的可构造子集，则 $T\cap U$ 是 $U$ 的可构造子集。
  \item[(ii)] 进一步假设 $U$ 在 $X$ 中逆紧。$U$ 的子集 $Z$ 在 $X$ 中可构造，当且仅当它在
    $U$ 中可构造。
\end{enumerate}
\end{proposition}

\begin{enumerate}
  \item[(i)] 利用（\hyperref[0.9.1.3-fr]{9.1.3}），问题归结为证明：若 $T$ 是 $X$ 中的逆紧开集，
    则 $T\cap U$ 是 $U$ 中的逆紧开集。换言之，对每个拟紧开集
    $W\subset U$，$T\cap U\cap W=T\cap W$ 拟紧；这立即由假设得出。
  \item[(ii)] 由 (i)，条件是必要的，只须证明它也是充分的。由
    （\hyperref[0.9.1.3-fr]{9.1.3}），只须考虑 $Z$ \emph{在 $U$ 中}为逆紧开集的情形，
    因为这将蕴含 $U\setminus Z$ 在 $X$ 中可构造，而且若 $Z$、$Z'$ 是两个
    \emph{在 $U$ 中}的逆紧开集，则 $Z\cap(U\setminus Z')$ 确实在 $X$ 中可构造。现在，若
    $W$ 是 $X$ 中的拟紧开集，$Z$ 是 $U$ 中的逆紧开集，则
    $Z\cap W=Z\cap(W\cap U)$；按假设，$W\cap U$ 是 $U$ 中的拟紧开集。因此
    $W\cap Z$ 的确拟紧，从而 $Z$ 是 $X$ 中的逆紧开集，并且\emph{更不待言}在 $X$ 中可构造。
\end{enumerate}

\begin{corollary}[9.1.9]
\label{0.9.1.9-fr}
设 $X$ 是拓扑空间，$(U_i)_{i\in I}$ 是由 $X$ 中逆紧开集组成的 $X$ 的一个有限覆盖。$X$ 的子集
$Z$ 在 $X$ 中可构造，当且仅当对每个 $i\in I$，$Z\cap U_i$ 在 $U_i$ 中可构造。
\end{corollary}

\begin{env}[9.1.10]
\label{0.9.1.10-fr}
特别地，假设 $X$ 拟紧，并且 $X$ 的每个点
\oldpage[0\textsubscript{III}]{14}
都有一个由 $X$ 中逆紧开邻域组成的邻域基本系（这些邻域\emph{更不待言}也是拟紧的）。那么，
$X$ 的子集 $Z$ 在 $X$ 中可构造这一条件具有\emph{局部}性质；换言之，当且仅当对每个
$x\in X$，存在 $x$ 的一个开邻域 $V$，使 $V\cap Z$ 在 $V$ 中可构造。事实上，若满足这一条件，
则由关于 $X$ 的假设及（\hyperref[0.9.1.8-fr]{9.1.8, (i)}），对每个 $x\in X$，都存在 $x$ 的一个
\emph{在 $X$ 中逆紧}的开邻域 $V$，使 $V\cap Z$ 在 $V$ 中可构造。此时只须用有限多个这样的邻域覆盖
$X$，再应用（\hyperref[0.9.1.9-fr]{9.1.9}）。
\end{env}

\begin{definition}[9.1.11]
\label{0.9.1.11-fr}
设 $X$ 是拓扑空间。若对每个 $x\in X$，都存在 $x$ 的一个开邻域 $V$，使 $T\cap V$ 在 $V$ 中可构造，
就称 $X$ 的子集 $T$ \emph{在 $X$ 中局部可构造}。
\end{definition}

由（\hyperref[0.9.1.8-fr]{9.1.8, (i)}）立即可知：若 $V$ 满足 $V\cap T$ 在 $V$ 中可构造，则对每个开集
$W\subset V$，$W\cap T$ 在 $W$ 中可构造。若 $T$ 在 $X$ 中局部可构造，则对 $X$ 的每个开集
$U$，$T\cap U$ 在 $U$ 中局部可构造，这也由上面的说明得出。同一说明表明，$X$ 中局部可构造子集的集合
对有限并与有限交封闭；另一方面，它显然也对取补封闭。

\begin{proposition}[9.1.12]
\label{0.9.1.12-fr}
设 $X$ 是拓扑空间。$X$ 中的每个可构造集都在 $X$ 中局部可构造。若 $X$ 拟紧，且其拓扑有一个由
$X$ 中逆紧集组成的基，则逆命题也成立。
\end{proposition}

第一项断言由定义（\hyperref[0.9.1.11-fr]{9.1.11}）得出，第二项由
（\hyperref[0.9.1.10-fr]{9.1.10}）得出。

\begin{corollary}[9.1.13]
\label{0.9.1.13-fr}
设拓扑空间 $X$ 的拓扑有一个由 $X$ 中逆紧集组成的基。那么，$X$ 中每个局部可构造子集
$T$ 都在 $X$ 中逆紧。
\end{corollary}

事实上，设 $U$ 是 $X$ 中的拟紧开集。$T\cap U$ 在 $U$ 中局部可构造，故由
（\hyperref[0.9.1.12-fr]{9.1.12}）在 $U$ 中可构造，再由
（\hyperref[0.9.1.4-fr]{9.1.4}）可知它拟紧。

\subsection{Noether 空间中的可构造集}
\label{subsection:0.9.2-fr}

\begin{env}[9.2.1]
\label{0.9.2.1-fr}
我们已经看到（\hyperref[0.9.1.7-fr]{9.1.7}），在 Noether 空间 $X$ 中，$X$ 的可构造子集就是
\emph{$X$ 的有限多个局部闭子集之并}。

从 Noether 空间 $X'$ 到 $X$ 的连续映射对 $X$ 中可构造集的逆像，在 $X'$ 中可构造。若
$Y$ 是 Noether 空间 $X$ 的可构造子集，则 $Y$ 的那些作为 $Y$ 的子空间可构造的子集，
恰好就是那些作为 $X$ 的子空间可构造的子集。
\end{env}

\begin{proposition}[9.2.2]
\label{0.9.2.2-fr}
设 $X$ 是不可约 Noether 空间，$E$ 是 $X$ 的一个可构造子集。$E$ 在 $X$ 中处处稠密，当且仅当
$E$ 包含 $X$ 的一个非空开子集。
\end{proposition}

\oldpage[0\textsubscript{III}]{15}
条件的充分性是显然的，因为每个非空开集都在 $X$ 中稠密。反过来，设
$E=\bigcup_{i=1}^n(U_i\cap F_i)$ 是 $X$ 的一个可构造子集，其中 $U_i$ 是非空开集，$F_i$ 是
$X$ 中的闭集；于是
$\overline{E}\subset\bigcup_i F_i$。因此，若 $\overline{E}=X$，则 $X$ 等于某个 $F_i$，从而
$E\supset U_i$，证明完毕。

若 $X$ 有一个泛点 $x$
（\hyperref[0.2.1.2-fr]{$0_{\mathrm I}$, 2.1.2}），则
（\hyperref[0.9.2.2-fr]{9.2.2}）的条件等价于关系 $x\in E$。

\begin{proposition}[9.2.3]
\label{0.9.2.3-fr}
设 $X$ 是 Noether 空间。$X$ 的子集 $E$ 可构造，当且仅当对 $X$ 的每个不可约闭子集
$Y$，$E\cap Y$ 在 $Y$ 中无处稠密，或者包含 $Y$ 的一个非空开子集。
\end{proposition}

条件的必要性来自：$E\cap Y$ 必须是 $Y$ 的可构造子集，以及
（\hyperref[0.9.2.2-fr]{9.2.2}）；因为 $Y$ 的非稠密子集在不可约空间
$Y$ 中必然无处稠密（\hyperref[0.2.1.1-fr]{$0_{\mathrm I}$, 2.1.1}）。为证明条件充分，对由满足
$Y\cap E$ 可构造的 $X$ 的闭子集 $Y$ 所组成的集合
$\mathfrak{F}$，应用 Noether 归纳原理
（\hyperref[0.2.2.2-fr]{$0_{\mathrm I}$, 2.2.2}）（这里相对于 $Y$ 或相对于 $X$ 可构造是同一回事）。
因此可假设，对 $X$ 的每个闭子集 $Y\ne X$，$E\cap Y$ 都可构造。先假设 $X$ 不是不可约的，并设
$X_i$（$1\leq i\leq m$）为其不可约分支；这些分支必然只有有限多个
（\hyperref[0.2.2.5-fr]{$0_{\mathrm I}$, 2.2.5}）。按假设，$E\cap X_i$ 都可构造，故它们的并
$E$ 也可构造。再假设 $X$ 不可约；按假设，或者 $E$ 无处稠密，于是
$\overline{E}\ne X$，并且 $E=E\cap\overline{E}$ 可构造；或者 $E$ 包含 $E$ 的一个非空开集
$U$，于是它是 $U$ 与 $E\cap(X\setminus U)$ 的并；但 $X\setminus U$ 是异于 $X$ 的闭集，
故 $E\cap(X\setminus U)$ 可构造。因此 $E$ 本身可构造，证明完毕。
% EGA3-ZH-DEF-0003: printed 'open U of E' retained in this noncanonical candidate.

\begin{corollary}[9.2.4]
\label{0.9.2.4-fr}
设 $X$ 是 Noether 空间，$(E_\alpha)$ 是 $X$ 的可构造子集的一个递增有向族，满足：
\begin{enumerate}
  \item[$1^\circ$] $X$ 是族 $(E_\alpha)$ 的并。
  \item[$2^\circ$] $X$ 的每个不可约闭子集都包含于某个 $E_\alpha$ 的闭包中。
\end{enumerate}
那么，存在指标 $\alpha$，使得 $X=E_\alpha$。

若 $X$ 的每个不可约闭子集都有泛点，则可删去假设 $2^\circ$。
\end{corollary}

对至少包含于某个 $E_\alpha$ 中的 $X$ 的闭子集所组成的集合
$\mathfrak{M}$，应用 Noether 归纳原理
（\hyperref[0.2.2.2-fr]{$0_{\mathrm I}$, 2.2.2}）。于是可假设，$X$ 的每个闭子集
$Y\ne X$ 都包含于某个 $E_\alpha$。若 $X$ 不是不可约的，则命题显然成立，因为 $X$ 的每个不可约分支
$X_i$（$1\leq i\leq m$）都包含于某个 $E_{\alpha_i}$，而存在一个包含所有
$E_{\alpha_i}$ 的 $E_\alpha$。现假设 $X$ 不可约。按假设，存在 $\beta$，使
$X=\overline{E_\beta}$；所以由（\hyperref[0.9.2.2-fr]{9.2.2}），$E_\beta$ 包含 $X$ 的一个非空开集
$U$。于是闭集 $X\setminus U$ 包含于某个 $E_\gamma$；只须取一个同时包含
$E_\beta$ 与 $E_\gamma$ 的 $E_\alpha$。若 $X$ 的每个不可约闭子集 $Y$
\oldpage[0\textsubscript{III}]{16}
都有泛点 $y$，则存在 $\alpha$，使 $y\in E_\alpha$，从而
$Y=\overline{\{y\}}\subset\overline{E_\alpha}$；条件 $2^\circ$ 因而是条件 $1^\circ$ 的推论。

\begin{proposition}[9.2.5]
\label{0.9.2.5-fr}
设 $X$ 是 Noether 空间，$x$ 是 $X$ 的一点，$E$ 是 $X$ 的一个可构造子集。$E$ 是 $x$ 的邻域，
当且仅当对 $X$ 的每个包含 $x$ 的不可约闭子集 $Y$，$E\cap Y$ 都在 $Y$ 中稠密
（若 $Y$ 有泛点 $y$，由（\hyperref[0.9.2.2-fr]{9.2.2}），这也就是说 $y\in E$）。
\end{proposition}

条件的必要性是显然的；我们来证明充分性。对 $X$ 的如下闭子集所组成的集合
$\mathfrak{M}$ 应用 Noether 归纳原理：$Y$ 包含 $x$，并且 $E\cap Y$ 是 $Y$ 中 $x$ 的一个邻域。
于是可假设，$X$ 的每个包含 $x$ 的闭子集 $Y\ne X$ 都属于 $\mathfrak{M}$。若 $X$ 不是不可约的，则
$X$ 的每个包含 $x$ 的不可约分支 $X_i$ 都异于 $X$，故 $E\cap X_i$ 是相对于 $X_i$ 的
$x$ 的一个邻域。因此，$E$ 是 $X$ 的所有包含 $x$ 的不可约分支之并中的 $x$ 的一个邻域；
又因为这个并是 $X$ 中 $x$ 的一个邻域，所以 $E$ 也是。若 $X$ 不可约，则按假设，$E$ 在
$X$ 中稠密，故包含 $X$ 的一个非空开子集 $U$（\hyperref[0.9.2.2-fr]{9.2.2}）。若 $x\in U$，
命题显然；否则，按假设，$x$ 是 $E\cap(X\setminus U)$ 相对于 $X\setminus U$ 的内点，
所以 $X\setminus E$ 在 $X$ 中的闭包不含 $x$，而该闭包的补集是包含于 $E$ 的 $x$ 的一个邻域，
证明完毕。

\begin{corollary}[9.2.6]
\label{0.9.2.6-fr}
设 $X$ 是 Noether 空间，$E$ 是 $X$ 的一个子集。$E$ 是 $X$ 中的开集，当且仅当对 $X$ 的每个与
$E$ 相交的不可约闭子集 $Y$，$E\cap Y$ 都包含 $Y$ 的一个非空开子集。
\end{corollary}

条件的必要性是显然的。反过来，若条件成立，则由
（\hyperref[0.9.2.3-fr]{9.2.3}），它蕴含 $E$ 可构造。此外，
（\hyperref[0.9.2.5-fr]{9.2.5}）表明，此时 $E$ 是其每一点的邻域，结论由此得到。

\subsection{可构造函数}
\label{subsection:0.9.3-fr}

\begin{definition}[9.3.1]
\label{0.9.3.1-fr}
设 $h$ 是从拓扑空间 $X$ 到集合 $T$ 的一个映射。若对每个 $t\in T$，$h^{-1}(t)$ 都可构造，
并且除 $t$ 的有限多个取值外均为空，就称 $h$ 是\emph{可构造的}；此时对 $T$ 的每个子集
$S$，$h^{-1}(S)$ 都可构造。若每个 $x\in X$ 都有一个开邻域 $V$，使 $h|V$ 可构造，
就称 $h$ 是\emph{局部可构造的}。
\end{definition}

每个可构造函数都局部可构造；若 $X$ 拟紧，并且有一个由 $X$ 中逆紧开集组成的基，则逆命题成立
（特别地，当 $X$ 是 Noether 空间时成立）。

\begin{proposition}[9.3.2]
\label{0.9.3.2-fr}
设 $h$ 是从 Noether 空间 $X$ 到集合 $T$ 的一个映射。$h$ 可构造，当且仅当对 $X$ 的每个不可约闭子集
$Y$，存在 $Y$ 的一个非空子集 $U$，它在 $Y$ 中开，并且 $h$ 在其中为常值。
\end{proposition}

条件是必要的：事实上，按假设，$h$ 在 $Y$ 上只取有限多个值 $t_i$，而每个集合
$h^{-1}(t_i)\cap Y$ 都在 $Y$ 中可构造
（\hyperref[0.9.2.1-fr]{9.2.1}）。这些集合不可能全都是空间 $Y$ 的无处稠密子集，所以至少有一个包含非空开集
（\hyperref[0.9.2.3-fr]{9.2.3}）。

\oldpage[0\textsubscript{III}]{17}
为证明条件充分，对由如下 $X$ 的闭子集 $Y$ 所组成的集合 $\mathfrak{M}$ 应用 Noether 归纳原理：
限制 $h|Y$ 可构造。于是可假设，对 $X$ 的每个闭子集 $Y\ne X$，$h|Y$ 都可构造。若 $X$ 不是不可约的，
则 $h$ 在 $X$ 的每个不可约分支 $X_i$（只有有限多个）上的限制都可构造；此时由定义
（\hyperref[0.9.3.1-fr]{9.3.1}）立即可知 $h$ 可构造。若 $X$ 不可约，则按假设存在 $X$ 的一个非空开子集
$U$，使 $h$ 在其中为常值；另一方面，按假设，$h$ 在 $X\setminus U$ 上的限制可构造，
从而立即可知 $h$ 可构造。

\begin{corollary}[9.3.3]
\label{0.9.3.3-fr}
设 $X$ 是每个不可约闭子集都有泛点的 Noether 空间。若 $h$ 是从 $X$ 到集合 $T$ 的一个映射，
并且对每个 $t\in T$，$h^{-1}(t)$ 都可构造，则 $h$ 可构造。
\end{corollary}

事实上，若 $Y$ 是 $X$ 的一个不可约闭子集，$y$ 是它的泛点，则
$Y\cap h^{-1}(h(y))$ 可构造且包含 $y$，故由
（\hyperref[0.9.2.2-fr]{9.2.2}），这个集合包含 $Y$ 的一个非空开子集；只须应用
（\hyperref[0.9.3.2-fr]{9.3.2}）。

\begin{proposition}[9.3.4]
\label{0.9.3.4-fr}
设 $X$ 是每个不可约闭子集都有泛点的 Noether 空间，$h$ 是从 $X$ 到一个有序集的可构造映射。
$h$ 在 $X$ 上上半连续，当且仅当对每个 $x\in X$ 以及 $x$ 的每个一般化
（\hyperref[0.2.1.2-fr]{$0_{\mathrm I}$, 2.1.2}）$x'$，都有
$h(x')\leq h(x)$。
\end{proposition}

函数 $h$ 只取有限多个值。因此，说它上半连续，就是说对每个 $x\in X$，由满足
$h(y)\leq h(x)$ 的 $y\in X$ 所组成的集合 $E$ 是 $x$ 的一个邻域。按假设，$E$ 是 $X$ 的一个可构造子集；
另一方面，说 $X$ 的一个不可约闭子集 $Y$ 包含 $x$，就是说明其泛点 $y$ 是 $x$ 的一个一般化。
结论于是由（\hyperref[0.9.2.5-fr]{9.2.5}）得出。

\section{关于平坦模的补充}
\label{section:0.10-fr}

对于（\hyperref[subsection:0.10.1-fr]{10.1}）和
（\hyperref[subsection:0.10.2-fr]{10.2}）中未给出证明的性质，我们请读者参阅 Bourbaki，
\emph{Alg. comm.}，第 II、III 章。

\subsection{平坦模与自由模之间的关系}
\label{subsection:0.10.1-fr}

\begin{env}[10.1.1]
\label{0.10.1.1-fr}
设 $A$ 是一个环，$\mathfrak{J}$ 是 $A$ 的一个理想，$M$ 是一个 $A$-模；对每个整数
$p\geq0$，都有一个典范的 $(A/\mathfrak{J})$-模同态
\[
\label{0.10.1.1.1-fr}
  \varphi_p:(M/\mathfrak{J}M)\otimes_{A/\mathfrak{J}}
  (\mathfrak{J}^p/\mathfrak{J}^{p+1})
  \longrightarrow \mathfrak{J}^pM/\mathfrak{J}^{p+1}M
  \tag{10.1.1.1}
\]
它显然是\emph{满射}。以
$\operatorname{gr}(A)=\bigoplus_{p\geq0}\mathfrak{J}^p/\mathfrak{J}^{p+1}$ 表示由
$\mathfrak{J}^p$ 所给滤过的 $A$ 的伴随分次环，以
$\operatorname{gr}(M)=\bigoplus_{p\geq0}\mathfrak{J}^pM/\mathfrak{J}^{p+1}M$ 表示由
$\mathfrak{J}^pM$ 所给滤过的 $M$ 的伴随分次 $\operatorname{gr}(A)$-模；于是
$\operatorname{gr}_p(A)=\mathfrak{J}^p/\mathfrak{J}^{p+1}$，
$\operatorname{gr}_p(M)=\mathfrak{J}^pM/\mathfrak{J}^{p+1}M$；各 $\varphi_p$ 定义一个分次
$\operatorname{gr}(A)$-模的\emph{满射}同态
\[
\label{0.10.1.1.2-fr}
  \varphi:\operatorname{gr}_0(M)\otimes_{\operatorname{gr}_0(A)}
  \operatorname{gr}(A)\longrightarrow\operatorname{gr}(M).
  \tag{10.1.1.2}
\]
\end{env}

\oldpage[0\textsubscript{III}]{18}
\begin{env}[10.1.2]
\label{0.10.1.2-fr}
假设下列假设之一成立：
\begin{enumerate}
  \item[(i)] $\mathfrak{J}$ 是幂零的；
  \item[(ii)] $A$ 是 Noether 环，$\mathfrak{J}$ 包含于 $A$ 的 Jacobson 根中，并且 $M$ 是有限型的。
\end{enumerate}
那么下列性质等价：
\begin{enumerate}
  \item[$a)$] $M$ 是自由 $A$-模。
  \item[$b)$] $M/\mathfrak{J}M=M\otimes_A(A/\mathfrak{J})$ 是自由
    $(A/\mathfrak{J})$-模，并且 $\operatorname{Tor}_1^A(M,A/\mathfrak{J})=0$。
  \item[$c)$] $M/\mathfrak{J}M$ 是自由 $(A/\mathfrak{J})$-模，并且典范同态
    （\hyperref[0.10.1.1.2-fr]{10.1.1.2}）是单射（因而是双射）。
\end{enumerate}
\end{env}

\begin{env}[10.1.3]
\label{0.10.1.3-fr}
假设 $A/\mathfrak{J}$ 是一个\emph{域}（换言之，$\mathfrak{J}$ 是极大理想），并且
（\hyperref[0.10.1.2-fr]{10.1.2}）的假设 (i)、(ii) 之一成立。那么下列性质等价：
\begin{enumerate}
  \item[$a)$] $M$ 是自由 $A$-模。
  \item[$b)$] $M$ 是射影 $A$-模。
  \item[$c)$] $M$ 是平坦 $A$-模。
  \item[$d)$] $\operatorname{Tor}_1^A(M,A/\mathfrak{J})=0$。
  \item[$e)$] 典范同态（\hyperref[0.10.1.1.2-fr]{10.1.1.2}）是双射。
\end{enumerate}
这个结果特别适用于下列两种情形：
\begin{enumerate}
  \item[(i)] $M$ 是一个\emph{任意}模，它所在的局部环 $A$ 的极大理想
    $\mathfrak{J}$ 是\emph{幂零的}（例如，某个 Artin 局部环）。
  \item[(ii)] $M$ 是\emph{局部 Noether} 环上的\emph{有限型}模。
\end{enumerate}
\end{env}

\subsection{平坦性的局部判据}
\label{subsection:0.10.2-fr}

\begin{env}[10.2.1]
\label{0.10.2.1-fr}
沿用（\hyperref[0.10.1.1-fr]{10.1.1}）的假设和记号，考虑下列条件：
\begin{enumerate}
  \item[$a)$] $M$ 是平坦 $A$-模。
  \item[$b)$] $M/\mathfrak{J}M$ 是平坦 $(A/\mathfrak{J})$-模，并且
    $\operatorname{Tor}_1^A(M,A/\mathfrak{J})=0$。
  \item[$c)$] $M/\mathfrak{J}M$ 是平坦 $(A/\mathfrak{J})$-模，并且典范同态
    （\hyperref[0.10.1.1.2-fr]{10.1.1.2}）是双射。
  \item[$d)$] 对每个 $n\geq1$，$M/\mathfrak{J}^nM$ 是平坦
    $(A/\mathfrak{J}^n)$-模。
\end{enumerate}
此时有蕴涵关系
\[
  a\Rightarrow b\Rightarrow c\Rightarrow d
\]
而且若 $\mathfrak{J}$ 是\emph{幂零的}，则四个条件 $a)$、$b)$、$c)$、$d)$
\emph{等价}。若 $A$ 是 Noether 环，并且 $M$ 还\emph{理想分离}，结论也相同；所谓理想分离，
是指对 $A$ 的每个理想 $\mathfrak{a}$，$A$-模 $\mathfrak{a}\otimes_A M$ 关于
$\mathfrak{J}$-预进制拓扑都是\emph{分离的}。
\end{env}

\begin{env}[10.2.2]
\label{0.10.2.2-fr}
设 $A$ 是 Noether 环，$B$ 是交换 Noether $A$-代数，$\mathfrak{J}$ 是 $A$ 的一个理想，
使得 $\mathfrak{J}B$ 包含于 $B$ 的 Jacobson 根中，并且 $M$ 是有限型 $B$-模。那么，把 $M$
看作 $A$-模时，（\hyperref[0.10.2.1-fr]{10.2.1}）的条件 $a)$、$b)$、$c)$、$d)$ 等价。

这个结果主要适用于 $A$ 和 $B$ 都是 Noether \emph{局部}环、同态 $A\to B$ 是
\emph{局部}同态的情形。更特别地，若这时 $\mathfrak{J}$ 是 $A$ 的\emph{极大}理想，则在条件
$b)$ 和 $c)$ 中，可以删去 $M/\mathfrak{J}M$ 平坦这一自动成立的假设；条件 $d)$ 则表示各模
$M/\mathfrak{J}^nM$ 在 $A/\mathfrak{J}^n$ 上都是\emph{自由的}。
\end{env}

\oldpage[0\textsubscript{III}]{19}
\begin{env}[10.2.3]
\label{0.10.2.3-fr}
设关于 $A$、$B$、$\mathfrak{J}$、$M$ 的假设与
（\hyperref[0.10.2.2-fr]{10.2.2}）开头所述相同；令 $\widehat A$ 是 $A$ 关于
$\mathfrak{J}$-预进制拓扑的分离完备化，$\widehat M$ 是 $M$ 关于
$\mathfrak{J}B$-预进制拓扑的分离完备化。那么，$M$ 是平坦 $A$-模，当且仅当
$\widehat M$ 是平坦 $\widehat A$-模。
\end{env}

\begin{env}[10.2.4]
\label{0.10.2.4-fr}
设 $\rho:A\to B$ 是 Noether 局部环之间的局部同态，$k$ 是 $A$ 的剩余域，$M$、$N$
是两个有限型 $B$-模，并假设 $N$ 是 $A$-\emph{平坦的}。设 $u:M\to N$ 是一个 $B$-同态。
那么下列条件等价：
\begin{enumerate}
  \item[$a)$] $u$ 是单射，并且 $\operatorname{Coker}(u)$ 是平坦 $A$-模。
  \item[$b)$] $u\otimes1:M\otimes_A k\to N\otimes_A k$ 是单射。
\end{enumerate}
\end{env}

\begin{env}[10.2.5]
\label{0.10.2.5-fr}
设 $\rho:A\to B$、$\sigma:B\to C$ 是 Noether 局部环之间的局部同态，$k$ 是 $A$
的剩余域，$M$ 是有限型 $C$-模。假设 $B$ 是\emph{平坦} $A$-模。那么下列条件等价：
\begin{enumerate}
  \item[$a)$] $M$ 是平坦 $B$-模。
  \item[$b)$] $M$ 是平坦 $A$-模，并且 $M\otimes_A k$ 是平坦
    $(B\otimes_A k)$-模。
\end{enumerate}
\end{env}

\begin{proposition}[10.2.6]
\label{0.10.2.6-fr}
设 $A$、$B$ 是两个 Noether 局部环，$\rho:A\to B$ 是局部同态，$\mathfrak{J}$ 是 $B$
的一个包含于极大理想中的理想，$M$ 是有限型 $B$-模。假设对每个 $n\geq0$，
$M_n=M/\mathfrak{J}^{n+1}M$ 都是平坦 $A$-模。那么 $M$ 是平坦 $A$-模。
\end{proposition}

需要证明：对有限型 $A$-模的每个单同态 $u:N'\to N$，同态
$v=1\otimes u:M\otimes_A N'\to M\otimes_A N$ 都是单射。可是，$M\otimes_A N'$ 和
$M\otimes_A N$ 都是有限型 $B$-模，因而关于 $\mathfrak{J}$-预进制拓扑分离
（\hyperref[0.7.3.5-fr]{$0_{\mathrm I}$, 7.3.5}）；所以只须证明分离完备化之间的同态
$\widehat v:\widehat{(M\otimes_A N')}\to\widehat{(M\otimes_A N)}$ 是单射。又有
$\widehat v=\varprojlim v_n$，其中 $v_n$ 是同态
$1\otimes u:M_n\otimes_A N'\to M_n\otimes_A N$；按假设 $M_n$ 是 $A$-平坦的，故对每个
$n$，$v_n$ 都是单射；所以 $\widehat v$ 也同样是单射\SourceCorrectionNote{法文印本此处写作
$v$；但前文已把所需结论归结为分离完备化之间的映射 $\widehat v$ 为单射，并且紧接着写出
$\widehat v=\varprojlim v_n$。本译依稳定勘误
EGA-III1-10.2.6-P363-AUTHORIAL-V-HAT-001 及 D17 交叉核验采用 $\widehat v$；
法文外交式来源保持不变。}，因为函子 $\varprojlim$ 左正合。

\begin{corollary}[10.2.7]
\label{0.10.2.7-fr}
设 $A$ 是 Noether 环，$B$ 是 Noether 局部环，$\rho:A\to B$ 是同态，$f$ 是 $B$ 的
极大理想中的一个元素，$M$ 是有限型 $B$-模。假设 $M$ 上的位似
$f_M:x\mapsto fx$ 是单射，并且 $M/fM$ 是平坦 $A$-模。那么 $M$ 是平坦 $A$-模。
\end{corollary}

对 $i\geq0$，令 $M_i=f^iM$；由于 $f_M$ 是单射，$M_i/M_{i+1}$ 同构于 $M/fM$，
所以对每个 $i\geq0$ 都是 $A$-平坦的；由正合列
\[
  0\longrightarrow M_i/M_{i+1}\longrightarrow M/M_{i+1}
  \longrightarrow M/M_i\longrightarrow0
\]
对 $i$ 归纳可知，对每个 $i\geq0$，$M/M_i$ 都是 $A$-\emph{平坦的}
（\hyperref[0.6.1.2-fr]{$0_{\mathrm I}$, 6.1.2}）；所以可以应用
（\hyperref[0.10.2.6-fr]{10.2.6}）。也可以直接作如下论证：对每个有限型 $A$-模 $N$，
$M\otimes_A N$ 都是有限型 $B$-模；因为 $f$ 属于 $B$ 的根 $\mathfrak n$，
$M\otimes_A N$ 上的 $(f)$-进制拓扑比
\oldpage[0\textsubscript{III}]{20}
$\mathfrak n$-进制拓扑更细，而后者是\emph{分离的}
（\hyperref[0.7.3.5-fr]{$0_{\mathrm I}$, 7.3.5}）。此外，由于 $M/M_i$ 是 $A$-平坦的，有
\[
  f^i(M\otimes_A N)=\operatorname{Im}(M_i\otimes_A N\to M\otimes_A N)
  =\operatorname{Ker}(M\otimes_A N\to(M/M_i)\otimes_A N)
\]
（\hyperref[0.6.1.2-fr]{$0_{\mathrm I}$, 6.1.2}）。现设 $N$ 是有限型 $A$-模，$N'$ 是 $N$
的一个子模，$j:N'\to N$ 是典范嵌入；在交换图
\[
\xymatrix{
  M\otimes_A N'\ar[r]\ar[d]_{1_M\otimes j} &
  (M/M_i)\otimes_A N'\ar[d]^{1_{M/M_i}\otimes j}\\
  M\otimes_A N\ar[r] & (M/M_i)\otimes_A N
}
\]
中，由于 $M/M_i$ 是 $A$-平坦的，$1_{M/M_i}\otimes j$ 是单射；因而对每个 $i$，都有
\[
  \operatorname{Ker}(M\otimes_A N'\to M\otimes_A N)
  \subset\operatorname{Ker}(M\otimes_A N'\to(M/M_i)\otimes_A N')
\]
如上所见，右端各项的交约化为 $0$，所以左端也如此，从而 $M$ 是 $A$-平坦的。

\begin{proposition}[10.2.8]
\label{0.10.2.8-fr}
设 $A$ 是既约 Noether 环，$M$ 是有限型 $A$-模。假设对每个作为离散赋值环的 $A$-代数
$B$，$M\otimes_A B$ 都是平坦 $B$-模（因而是自由的
（\hyperref[0.10.1.3-fr]{10.1.3}））。那么 $M$ 是平坦 $A$-模。
\end{proposition}

我们知道，$M$ 平坦，当且仅当对 $A$ 的每个极大理想 $\mathfrak m$，$M_{\mathfrak m}$
都是平坦 $A_{\mathfrak m}$-模
（\hyperref[0.6.3.3-fr]{$0_{\mathrm I}$, 6.3.3}）；所以可以归结到 $A$ 为\emph{局部}环的情形
（\hyperref[0.1.2.8-fr]{$0_{\mathrm I}$, 1.2.8}）。于是，设 $\mathfrak m$ 是 $A$ 的极大理想，
$\mathfrak p_i$ $(1\leq i\leq r)$ 是它的各极小素理想，$k$ 是剩余域 $A/\mathfrak m$。
我们知道（\hyperref[II.7.1.7-fr]{II, 7.1.7}），对每个 $i$，都存在一个离散赋值环 $B_i$，
它与整环 $A/\mathfrak p_i$ 有相同的分式域 $K_i$，并且支配后者。令
$M_i=M\otimes_A B_i$。按假设，$M_i$ 是自由 $B_i$-模；因此，以 $k_i$ 表示 $B_i$ 的剩余域，有
\[
\label{0.10.2.8.1-fr}
  \operatorname{rg}_{k_i}(M_i\otimes_{B_i}k_i)
  =\operatorname{rg}_{K_i}(M_i\otimes_{B_i}K_i).
  \tag{10.2.8.1}
\]
可是，复合同态 $A\to A/\mathfrak p_i\to B_i$ 显然是局部同态，故 $k_i$ 是 $k$ 的扩张
\SourceCorrectionNote{法文印本把扩张方向写作“$k$ 是 $k_i$ 的扩张”；但该局部同态诱导
$k\to k_i$，而下一式 $(M\otimes_A k)\otimes_k k_i$ 也要求 $k_i$ 是 $k$-代数。
本译依稳定勘误 EGA-III1-10.2.8-P365-AUTHORIAL-RESIDUE-EXTENSION-DIRECTION-001
及 D17 交叉核验采用校正方向；该稳定编号中的 P365 为历史定位，实际可见印页为第 364 页；
法文外交式来源保持不变。}，
并且有 $M_i\otimes_{B_i}k_i=M\otimes_A k_i=(M\otimes_A k)\otimes_k k_i$；另一方面，
$M_i\otimes_{B_i}K_i=M\otimes_A K_i$。等式
（\hyperref[0.10.2.8.1-fr]{10.2.8.1}）因而给出
\[
  \operatorname{rg}_k(M\otimes_A k)=\operatorname{rg}_{K_i}(M\otimes_A K_i)
  \qquad\text{其中 }1\leq i\leq r
\]
又因为 $A$ 既约，我们知道该条件蕴含 $M$ 是\emph{自由} $A$-模
（Bourbaki，\emph{Alg. comm.}，第 II 章，\S~3，
n\textsuperscript{o}~2，命题 7）。

\subsection{局部环的平坦扩张的存在性}
\label{subsection:0.10.3-fr}

\begin{proposition}[10.3.1]
\label{0.10.3.1-fr}
设 $A$ 是 Noether 局部环，$\mathfrak J$ 是它的极大理想，$k=A/\mathfrak J$ 是它的剩余域。
设 $K$ 是域 $k$ 的一个扩张。存在从 $A$ 到某个 Noether 局部环 $B$ 的局部同态，使得
$B/\mathfrak J B$ 与 $K$ 在 $k$ 上同构，并且 $B$ 是平坦 $A$-模。
\end{proposition}

我们将分几步证明这个命题。

\begin{env}[10.3.1.1]
\label{0.10.3.1.1-fr}
先假设 $K=k(T)$，其中 $T$ 是一个不定元。在多项式环 $A'=A[T]$ 中，考虑由系数全都属于理想
$\mathfrak J$ 的多项式所组成的素理想 $\mathfrak J'=\mathfrak J A'$；显然
\oldpage[0\textsubscript{III}]{21}
$A'/\mathfrak J'$ 典范同构于 $k[T]$。我们证明分式环 $B=A'_{\mathfrak J'}$ 满足所求；
它显然是极大理想为 $\mathfrak L=\mathfrak J B$ 的 Noether 局部环。此外，
$B/\mathfrak L=(A'/\mathfrak J')_{\mathfrak J'}=(k[T])_{\mathfrak J'}$ 正是 $k[T]$ 的分式域
$K$。最后，$B$ 是平坦 $A'$-模，而 $A'$ 是自由 $A$-模，所以 $B$ 是平坦 $A$-模
（\hyperref[0.6.2.1-fr]{$0_{\mathrm I}$, 6.2.1}）。
\end{env}

\begin{env}[10.3.1.2]
\label{0.10.3.1.2-fr}
再假设 $K=k(t)=k[t]$，其中 $t$ 在 $k$ 上代数；设 $f\in k[T]$ 是 $t$ 的极小多项式；
存在一个首一多项式 $F\in A[T]$，它在 $k[T]$ 中的典范像是 $f$。仍令
$A'=A[T]$，并令 $\mathfrak J'$ 为 $A'$ 中的理想 $\mathfrak J A'+(F)$。我们将看到，
这时商环 $B=A'/(F)$ 满足所求。首先，$B$ 显然是\emph{自由} $A$-模，因而平坦。环
$A'/\mathfrak J'$ 同构于
\[
  (A'/\mathfrak J A')/((\mathfrak J A'+(F))/\mathfrak J A')
  =k[T]/(f)=K~;
\]
所以 $\mathfrak J'$ 在 $B$ 中的像 $\mathfrak L$ 是极大理想，并且显然有
$\mathfrak L=\mathfrak J B$。最后，$B$ 是半局部环，因为它是有限型 $A$-模
（Bourbaki，\emph{Alg. comm.}，第 IV 章，\S~2，
n\textsuperscript{o}~5，命题 9 的推论 3），并且它的极大理想与 $B/\mathfrak J B$ 的极大理想
一一对应（[13]，第 I 卷，第 259 页）；因此，上述论证证明了 $B$ 是局部环。
\end{env}

\begin{lemma}[10.3.1.3]
\label{0.10.3.1.3-fr}
设 $(A_\lambda,f_{\mu\lambda})$ 是局部环的一个滤过归纳系统，其中各
$f_{\mu\lambda}$ 都是局部同态；设 $\mathfrak m_\lambda$ 是 $A_\lambda$ 的极大理想，并令
$K_\lambda=A_\lambda/\mathfrak m_\lambda$。那么 $A'=\varinjlim A_\lambda$ 是局部环，
$\mathfrak m'=\varinjlim\mathfrak m_\lambda$ 是它的极大理想，而
$K=\varinjlim K_\lambda$ 是它的剩余域。此外，若对 $\lambda<\mu$ 有
$\mathfrak m_\mu=\mathfrak m_\lambda A_\mu$，则对每个 $\lambda$ 都有
$\mathfrak m'=\mathfrak m_\lambda A'$。若进一步对 $\lambda<\mu$，$A_\mu$ 是平坦
$A_\lambda$-模，并且所有 $A_\lambda$ 都是 Noether 环，则 $A'$ 是 Noether 环，且对每个
$\lambda$ 都是平坦 $A_\lambda$-模。
\end{lemma}

按假设，对 $\lambda<\mu$ 有
$f_{\mu\lambda}(\mathfrak m_\lambda)\subset\mathfrak m_\mu$，所以各
$\mathfrak m_\lambda$ 构成一个归纳系统，其极限 $\mathfrak m'$ 显然是 $A'$ 的一个理想。
此外，若 $x'\notin\mathfrak m'$，则存在 $\lambda$ 和某个 $x_\lambda\in A_\lambda$，使得
$x'=f_\lambda(x_\lambda)$（这里 $f_\lambda:A_\lambda\to A'$ 表示典范同态）；因为
$x'\notin\mathfrak m'$，必有 $x_\lambda\notin\mathfrak m_\lambda$，故 $x_\lambda$ 在
$A_\lambda$ 中有逆元 $y_\lambda$，而 $y'=f_\lambda(y_\lambda)$ 是 $x'$ 在 $A'$ 中的逆元。
这证明 $A'$ 是极大理想为 $\mathfrak m'$ 的局部环；关于 $K$ 的断言立即来自函子
$\varinjlim$ 的正合性。假设 $\mathfrak m_\mu=\mathfrak m_\lambda A_\mu$ 表示典范映射
$\mathfrak m_\lambda\otimes_{A_\lambda}A_\mu\to\mathfrak m_\mu$ 是满射；关系
$\mathfrak m'=\mathfrak m_\lambda A'$ 于是也由函子 $\varinjlim$ 的正合性以及它与张量积交换这一事实得到。

现假设对 $\lambda<\mu$ 有 $\mathfrak m_\mu=\mathfrak m_\lambda A_\mu$，并且 $A_\mu$
是平坦 $A_\lambda$-模。由（\hyperref[0.6.2.3-fr]{$0_{\mathrm I}$, 6.2.3}），对每个
$\lambda$，$A'$ 都是平坦 $A_\lambda$-模；又因 $A'$ 与 $A_\lambda$ 都是局部环，且
$\mathfrak m'=\mathfrak m_\lambda A'$，$A'$ 甚至是\emph{忠实平坦} $A_\lambda$-模
（\hyperref[0.6.6.2-fr]{$0_{\mathrm I}$, 6.6.2}）。最后，再假设各 $A_\lambda$ 是
\emph{Noether} 环；于是各 $\mathfrak m_\lambda$-预进制拓扑都是分离的
（\hyperref[0.7.3.5-fr]{$0_{\mathrm I}$, 7.3.5}）。我们先证明，由此可知 $A'$ 上的
$\mathfrak m'$-进制拓扑是\emph{分离的}。事实上，若 $x'\in A'$ 属于所有
$\mathfrak m'^{n}$ $(n>0)$，则对某个指标 $\mu$，它是某个 $x_\mu\in A_\mu$ 的像；而
$\mathfrak m'^{n}=\mathfrak m_\mu^n A'$ 在 $A_\mu$ 中的逆像是 $\mathfrak m_\mu^n$
（\hyperref[0.6.6.1-fr]{$0_{\mathrm I}$, 6.6.1}），所以 $x_\mu$ 属于所有
$\mathfrak m_\mu^n$。按假设，$x_\mu=0$，从而 $x'=0$。以 $\widehat A'$ 表示 $A'$ 关于
$\mathfrak m'$-进制拓扑的完备化；上述结果表明 $A'\subset\widehat A'$。我们将证明，对每个
$\lambda$，$\widehat A'$ 都是\emph{Noether} 环且是 $A_\lambda$-\emph{平坦的}；由此
\oldpage[0\textsubscript{III}]{22}
可知 $\widehat A'$ 是 $A'$-平坦的
（\hyperref[0.6.2.3-fr]{$0_{\mathrm I}$, 6.2.3}）；又因为
$\mathfrak m'\widehat A'\ne\widehat A'$，$\widehat A'$ 是忠实平坦 $A'$-模
（\hyperref[0.6.6.2-fr]{$0_{\mathrm I}$, 6.6.2}）。最终可由此断定 $A'$ 是\emph{Noether} 环
（\hyperref[0.6.5.2-fr]{$0_{\mathrm I}$, 6.5.2}），这就完成了引理的证明。

有 $\widehat A'=\varprojlim_n A'/\mathfrak m'^{n}$；由于 $A'$ 是 $A_\lambda$-平坦的，有
\[
  \mathfrak m'^{n}/\mathfrak m'^{n+1}
  =(\mathfrak m_\lambda^n/\mathfrak m_\lambda^{n+1})
    \otimes_{A_\lambda}A'
  =(\mathfrak m_\lambda^n/\mathfrak m_\lambda^{n+1})
    \otimes_{K_\lambda}(K_\lambda\otimes_{A_\lambda}A')
  =(\mathfrak m_\lambda^n/\mathfrak m_\lambda^{n+1})
    \otimes_{K_\lambda}K~;
\]
因为 $\mathfrak m_\lambda^n/\mathfrak m_\lambda^{n+1}$ 是有限维 $K_\lambda$-向量空间，
所以对每个 $n\geq0$，$\mathfrak m'^{n}/\mathfrak m'^{n+1}$ 都是有限维 $K$-向量空间。
由（\hyperref[0.7.2.12-fr]{$0_{\mathrm I}$, 7.2.12}）和
（\hyperref[0.7.2.8-fr]{$0_{\mathrm I}$, 7.2.8}）可知，$\widehat A'$ 是 Noether 环。
此外，我们知道 $\widehat A'$ 的极大理想是 $\mathfrak m'\widehat A'$，并且
$\widehat A'/\mathfrak m'^{n}\widehat A'$ 同构于 $A'/\mathfrak m'^{n}$；由于
$A'/\mathfrak m'^{n}=(A_\lambda/\mathfrak m_\lambda^n) \otimes_{A_\lambda}A'$，
$A'/\mathfrak m'^{n}$ 是平坦 $(A_\lambda/\mathfrak m_\lambda^n)$-模
（\hyperref[0.6.2.1-fr]{$0_{\mathrm I}$, 6.2.1}）。因此，判据
（\hyperref[0.10.2.2-fr]{10.2.2}）可用于 Noether $A_\lambda$-代数 $\widehat A'$，并表明
$\widehat A'$ 是 $A_\lambda$-平坦的。证毕。

\begin{env}[10.3.1.4]
\label{0.10.3.1.4-fr}
现在处理一般情形。存在一个序数 $\gamma$，并且对每个序数 $\lambda\leq\gamma$，存在 $K$
的一个包含 $k$ 的子域 $k_\lambda$，使得：$1^\circ$ 对每个 $\lambda<\gamma$，
$k_{\lambda+1}$ 是由单个元素生成的 $k_\lambda$ 的扩张；$2^\circ$ 对每个无前驱的序数
$\mu$，有 $k_\mu=\bigcup_{\lambda<\mu}k_\lambda$；$3^\circ$ $K=k_\gamma$。
事实上，只须考虑从各序数 $\xi\leq\beta$（其中 $\beta$ 适当）所成的集合到 $K$ 的一个双射
$\xi\mapsto t_\xi$，并用超限归纳定义 $k_\lambda$（$\lambda\leq\beta$）：若 $\lambda$ 无前驱，
则取为各 $\mu<\lambda$ 的 $k_\mu$ 之并；若 $\lambda=\nu+1$，则取为
$k_\nu(t_\xi)$，其中 $\xi$ 是满足 $t_\xi\notin k_\nu$ 的最小序数。于是，按定义，
$\gamma$ 是满足 $k_\gamma=K$ 的最小序数 $\leq\beta$。

在此基础上，我们将用超限递归定义一族 Noether 局部环 $A_\lambda$（$\lambda\leq\gamma$）以及
局部同态 $f_{\mu\lambda}:A_\lambda\to A_\mu$（$\lambda\leq\mu$），使它们满足下列条件：
\begin{enumerate}
  \item[(i)] $(A_\lambda,f_{\mu\lambda})$ 是归纳系统，且 $A_0=A$。
  \item[(ii)] 对每个 $\lambda$，有一个 $k$-同构
    $A_\lambda/\mathfrak J A_\lambda\xrightarrow{\sim}k_\lambda$。
  \item[(iii)] 对 $\lambda\leq\mu$，$A_\mu$ 是平坦 $A_\lambda$-模。
\end{enumerate}

现假设对 $\lambda<\mu<\xi$，$A_\lambda$ 与 $f_{\mu\lambda}$ 已经定义。先假设
$\xi=\zeta+1$，从而 $k_\xi=k_\zeta(t)$。若 $t$ 在 $k_\zeta$ 上超越，则按照
（\hyperref[0.10.3.1.1-fr]{10.3.1.1}）的方法，把 $A_\xi$ 定义为
$(A_\zeta[t])_{\mathfrak J A_\zeta[t]}$；$f_{\xi\zeta}$ 是典范映射，而对
$\lambda<\zeta$，取 $f_{\xi\lambda}=f_{\xi\zeta}\circ f_{\zeta\lambda}$。
鉴于（\hyperref[0.10.3.1.1-fr]{10.3.1.1}）所证，条件 (i) 到 (iii) 的验证是立即的。
再假设 $t$ 是代数的；设 $h$ 是它在 $k_\zeta[T]$ 中的极小多项式，$H$ 是
$A_\zeta[T]$ 中像为 $k_\zeta[T]$ 中的 $h$ 的一个首一多项式；此时取 $A_\xi$
等于 $A_\zeta[T]/(H)$，各 $f_{\xi\lambda}$ 如前定义；条件 (i) 到 (iii) 的验证来自
（\hyperref[0.10.3.1.2-fr]{10.3.1.2}）所证。

现假设 $\xi$ 无前驱；取 $A_\xi$ 为由各 $\lambda<\xi$ 的局部环组成的归纳系统
$(A_\lambda,f_{\mu\lambda})$ 的归纳极限；对 $\lambda<\xi$，把 $f_{\xi\lambda}$ 定义为典范映射。
$A_\xi$ 是 Noether 局部环、各 $f_{\xi\lambda}$ 是局部同态，以及条件 (i) 到 (iii) 对
$\lambda\leq\xi$ 成立，
\oldpage[0\textsubscript{III}]{23}
均由归纳假设和引理（\hyperref[0.10.3.1.3-fr]{10.3.1.3}）得到。完成这一构造以后，显然环
$B=A_\gamma$ 满足（\hyperref[0.10.3.1-fr]{10.3.1}）的陈述。
\end{env}

注意，由（\hyperref[0.10.2.1-fr]{10.2.1, $c)$}），有一个典范同构
\[
\label{0.10.3.1.5-fr}
  \operatorname{gr}(A)\otimes_k K\xrightarrow{\sim}\operatorname{gr}(B).
  \tag{10.3.1.5}
\]

另一方面，可以用其 $\mathfrak J B$-进制完备化 $\widehat B$ 代替 $B$，而不改变
（\hyperref[0.10.3.1-fr]{10.3.1}）的结论，因为 $\widehat B$ 是平坦 $B$-模
（\hyperref[0.7.3.3-fr]{$0_{\mathrm I}$, 7.3.3}），因而是平坦 $A$-模
（\hyperref[0.6.2.1-fr]{$0_{\mathrm I}$, 6.2.1}）。

此外，我们还证明了下列结果。

\begin{corollary}[10.3.2]
\label{0.10.3.2-fr}
若 $K$ 是有限次扩张，则可以假设 $B$ 是有限 $A$-代数。
\end{corollary}

\section{同调代数补充}
\label{section:0.11-fr}

\subsection{谱序列回顾}
\label{subsection:0.11.1-fr}

\begin{env}[11.1.1]
\label{0.11.1.1-fr}
下文将使用一种比 (T, 2.4) 中所定义者更一般的谱序列概念。沿用 (T, 2.4)
的记号，我们称阿贝尔范畴 $\mathcal C$ 中的一个\emph{谱序列}为由下列资料构成的系统 E：
\begin{enumerate}
  \item[(a)] 一族 $\mathcal C$ 的对象 $(E_r^{pq})$，其中
    $p\in\mathbb Z$、$q\in\mathbb Z$ 且 $r\geq2$。
  \item[(b)] 一族态射
    $d_r^{pq}:E_r^{pq}\to E_r^{p+r,q-r+1}$，满足
    $d_r^{p+r,q-r+1}d_r^{pq}=0$。置
    $Z_{r+1}(E_r^{pq})=\operatorname{Ker}(d_r^{pq})$、
    $B_{r+1}(E_r^{pq})=\operatorname{Im}(d_r^{p-r,q+r-1})$，于是
    \[
      B_{r+1}(E_r^{pq})\subset Z_{r+1}(E_r^{pq})\subset E_r^{pq}.
    \]
  \item[(c)] 一族同构
    $\alpha_r^{pq}:Z_{r+1}(E_r^{pq})/B_{r+1}(E_r^{pq})
    \xrightarrow{\sim}E_{r+1}^{pq}$。
\end{enumerate}

于是，对 $k\geq r+1$，关于 $k$ 递归地定义子对象 $B_k(E_r^{pq})$ 和
$Z_k(E_r^{pq})$：它们分别是商对象中那些经 $\alpha_r^{pq}$ 与子对象
$B_k(E_{r+1}^{pq})$ 和 $Z_k(E_{r+1}^{pq})$ 相等同者，在典范态射
$E_r^{pq}\to E_r^{pq}/B_{r+1}(E_r^{pq})$ 下的逆像。显然，在同构意义下有
\[
\label{0.11.1.1.1-fr}
  Z_k(E_r^{pq})/B_k(E_r^{pq})=E_k^{pq}
  \qquad\text{当 }k\geq r+1,
  \tag{11.1.1.1}
\]
并且，若再置 $B_r(E_r^{pq})=0$ 及 $Z_r(E_r^{pq})=E_r^{pq}$，则有包含关系
\[
\label{0.11.1.1.2-fr}
  0=B_r(E_r^{pq})\subset B_{r+1}(E_r^{pq})\subset B_{r+2}(E_r^{pq})
  \subset\cdots\subset Z_{r+2}(E_r^{pq})\subset Z_{r+1}(E_r^{pq})
  \subset Z_r(E_r^{pq})=E_r^{pq}.
  \tag{11.1.1.2}
\]

E 的其余资料如下：
\begin{enumerate}
  \item[(d)] $E_2^{pq}$ 的两个子对象 $B_\infty(E_2^{pq})$ 和
    $Z_\infty(E_2^{pq})$，满足
    $B_\infty(E_2^{pq})\subset Z_\infty(E_2^{pq})$，并且对每个 $k\geq2$，
    \[
      B_k(E_2^{pq})\subset B_\infty(E_2^{pq})
      \qquad\text{且}\qquad
      Z_\infty(E_2^{pq})\subset Z_k(E_2^{pq}).
    \]
    置
    \[
    \label{0.11.1.1.3-fr}
      E_\infty^{pq}=Z_\infty(E_2^{pq})/B_\infty(E_2^{pq}).
      \tag{11.1.1.3}
    \]
  \item[(e)]
    \oldpage[0\textsubscript{III}]{24}
    一族 $\mathcal C$ 的对象 $(E^n)$，每个对象都配备一个\emph{递减滤过}
    $(F^p(E^n))_{p\in\mathbb Z}$。照常以 $\operatorname{gr}(E^n)$ 表示滤过对象
    $E^n$ 所关联的分次对象，即各
    $\operatorname{gr}_p(E^n)=F^p(E^n)/F^{p+1}(E^n)$ 的直和。
  \item[(f)] 对每一对 $(p,q)\in\mathbb Z\times\mathbb Z$，给定同构
    $\beta^{pq}:E_\infty^{pq}\xrightarrow{\sim}
    \operatorname{gr}_p(E^{p+q})$。
\end{enumerate}

族 $(E^n)$（不计其滤过）称为谱序列 E 的\emph{收敛项}。

假设范畴 $\mathcal C$ 容许无限直和，或者对每个 $r\geq2$ 和每个
$n\in\mathbb Z$，满足 $p+q=n$ 且 $E_r^{pq}\ne0$ 的数对 $(p,q)$ 只有有限个
（只要求 $r=2$ 时如此便已足够）。于是
\[
  E_r^{(n)}=\sum_{p+q=n}E_r^{pq}
\]
有定义。以 $d_r^{(n)}$ 表示态射 $E_r^{(n)}\to E_r^{(n+1)}$，其在每个满足
$p+q=n$ 的数对 $(p,q)$ 所对应的 $E_r^{pq}$ 上的限制为 $d_r^{pq}$，则
$d_r^{(n+1)}\circ d_r^{(n)}=0$。换言之，$(E_r^{(n)})_{n\in\mathbb Z}$
是 $\mathcal C$ 中一个微分次数为 $+1$ 的\emph{复形 $E_r^{(\bullet)}$}；
由 c) 可知，对每个 $r\geq2$，$\mathrm H^n(E_r^{(\bullet)})$
同构于 $E_{r+1}^{(n)}$。
\end{env}

\begin{env}[11.1.2]
\label{0.11.1.2-fr}
从谱序列 E 到谱序列 $E'=(E_r^{\prime pq},E^{\prime n})$ 的一个
\emph{态射 $u:E\to E'$}，由态射族
$u_r^{pq}:E_r^{pq}\to E_r^{\prime pq}$ 和 $u^n:E^n\to E^{\prime n}$
构成，其中 $u^n$ 与 $E^n$ 和 $E^{\prime n}$ 的滤过相容，并且诸图
\[
  \xymatrix{
    E_r^{pq}\ar[r]^-{d_r^{pq}}\ar[d]_{u_r^{pq}} &
    E_r^{p+r,q-r+1}\ar[d]^{u_r^{p+r,q-r+1}}\\
    E_r^{\prime pq}\ar[r]_{d_r^{\prime pq}} &
    E_r^{\prime p+r,q-r+1}
  }
\]
交换。此外，取商后 $u_r^{pq}$ 给出态射
\[
  \overline u_r^{pq}:
  Z_{r+1}(E_r^{pq})/B_{r+1}(E_r^{pq})\longrightarrow
  Z_{r+1}(E_r^{\prime pq})/B_{r+1}(E_r^{\prime pq})
\]
并且必须有
$\alpha_r^{\prime pq}\circ\overline u_r^{pq}
=u_{r+1}^{pq}\circ\alpha_r^{pq}$。最后，还必须有
$u_2^{pq}(B_\infty(E_2^{pq}))\subset B_\infty(E_2^{\prime pq})$、
$u_2^{pq}(Z_\infty(E_2^{pq}))\subset Z_\infty(E_2^{\prime pq})$。取商后，
$u_2^{pq}$ 因而给出态射
$u_\infty^{pq}:E_\infty^{pq}\to E_\infty^{\prime pq}$，并且图
\[
  \xymatrix{
    E_\infty^{pq}\ar[r]^-{u_\infty^{pq}}\ar[d]_{\beta^{pq}} &
    E_\infty^{\prime pq}\ar[d]^{\beta^{\prime pq}}\\
    \operatorname{gr}_p(E^{p+q})
      \ar[r]_{\operatorname{gr}_p(u^{p+q})} &
    \operatorname{gr}_p(E^{\prime p+q})
  }
\]
必须交换。

由以上定义，对 $r$ 归纳可知：如果诸 $u_2^{pq}$ 都是\emph{同构}，那么对
$r\geq2$，诸 $u_r^{pq}$ 也都是同构；\emph{如果还知道}
$u_2^{pq}(B_\infty(E_2^{pq}))=B_\infty(E_2^{\prime pq})$、
$u_2^{pq}(Z_\infty(E_2^{pq}))=Z_\infty(E_2^{\prime pq})$，
\emph{并且诸 $u^n$ 都是同构}，那么可以断定 $u$ 是\emph{同构}。
\end{env}

\begin{env}[11.1.3]
\label{0.11.1.3-fr}
\oldpage[0\textsubscript{III}]{25}
回顾：若 $(F^p(X))_{p\in\mathbb Z}$ 是对象 $X\in\mathcal C$ 的一个
（递减）滤过，则称该滤过为\emph{分离的}，若
$\inf(F^p(X))=0$；称为\emph{离散的}，若存在 $p$ 使
$F^p(X)=0$；称为\emph{穷尽的}（或\emph{余分离的}），若
$\sup(F^p(X))=X$；称为\emph{余离散的}，若存在 $p$ 使 $F^p(X)=X$。

若谱序列 $E=(E_r^{pq},E^n)$ 满足
\[
  B_\infty(E_2^{pq})=\sup_k(B_k(E_2^{pq})),\qquad
  Z_\infty(E_2^{pq})=\inf_k(Z_k(E_2^{pq}))
\]
则称它\emph{弱收敛}（换言之，对象 $B_\infty(E_2^{pq})$ 和
$Z_\infty(E_2^{pq})$ 由谱序列 E 的资料 a) 到 c) 所确定）。
若谱序列 E 弱收敛，并且还满足下列条件，则称它\emph{正则}：
\begin{enumerate}
  \item[$1^\circ$] 对每一对 $(p,q)$，递减序列
    $(Z_k(E_2^{pq}))_{k\geq2}$ \emph{最终稳定}；E 弱收敛的假设于是蕴含：
    当 $k$ 充分大时（其界依赖于 $p$ 和 $q$），有
    $Z_\infty(E_2^{pq})=Z_k(E_2^{pq})$。
  \item[$2^\circ$] 对每个 $n$，$E^n$ 的滤过
    $(F^p(E^n))_{p\in\mathbb Z}$ 既\emph{离散}又\emph{穷尽}。
\end{enumerate}

若谱序列 E 弱收敛，并且还满足下列条件，则称它\emph{余正则}：
\begin{enumerate}
  \item[$3^\circ$] 对每一对 $(p,q)$，递增序列
    $(B_k(E_2^{pq}))_{k\geq2}$ \emph{最终稳定}；这蕴含
    $B_\infty(E_2^{pq})=B_k(E_2^{pq})$，从而
    $E_\infty^{pq}=\inf E_k^{pq}$。
  \item[$4^\circ$] 对每个 $n$，$E^n$ 的滤过是\emph{余离散的}。
\end{enumerate}

最后，如果 E 既正则又余正则，则称它\emph{双正则}；换言之，它满足下列条件：
\begin{enumerate}
  \item[(a)] 对每一对 $(p,q)$，序列
    $(B_k(E_2^{pq}))_{k\geq2}$ 和 $(Z_k(E_2^{pq}))_{k\geq2}$
    都\emph{最终稳定}，并且当 $k$ 充分大时，有
    $B_\infty(E_2^{pq})=B_k(E_2^{pq})$ 及
    $Z_\infty(E_2^{pq})=Z_k(E_2^{pq})$
    （这蕴含 $E_\infty^{pq}=E_k^{pq}$）。
  \item[(b)] 对每个 $n$，滤过 $(F^p(E^n))_{p\in\mathbb Z}$ 既
    \emph{离散}又\emph{余离散}（这也表述为该滤过是\emph{有限的}）。
\end{enumerate}

因此，(T, 2.4) 中定义的谱序列就是双正则谱序列。
\end{env}

\begin{env}[11.1.4]
\label{0.11.1.4-fr}
假设范畴 $\mathcal C$ 中存在滤过归纳极限，并且函子 $\varinjlim$
\emph{正合}（这等价于说 (T, 1.5) 的公理 AB 5) 成立；参见 T, 1.8）。
对象 $X\in\mathcal C$ 的滤过 $(F^p(X))_{p\in\mathbb Z}$ 穷尽这一条件，
此时也可写作 $\varinjlim_{p\to-\infty}F^p(X)=X$。若谱序列 E 弱收敛，则有
$B_\infty(E_2^{pq})=\varinjlim_{k\to\infty}B_k(E_2^{pq})$；若此外
$u:E\to E'$ 是从 E 到 $\mathcal C$ 中另一个弱收敛谱序列 $E'$ 的态射，
则由 $\varinjlim$ 的正合性，有
$u_2^{pq}(B_\infty(E_2^{pq}))=B_\infty(E_2^{\prime pq})$。而且：
\end{env}

\begin{proposition}[11.1.5]
\label{0.11.1.5-fr}
设 $\mathcal C$ 是滤过归纳极限正合的阿贝尔范畴，E、$E'$ 是
$\mathcal C$ 中两个正则谱序列，$u:E\to E'$ 是谱序列的态射。若诸
$u_2^{pq}$ 都是同构，则 $u$ 也是同构。
\end{proposition}

\begin{proof}
由（\hyperref[0.11.1.2-fr]{11.1.2}）已知诸 $u_r^{pq}$ 都是同构，并且
\[
  u_2^{pq}(B_\infty(E_2^{pq}))=B_\infty(E_2^{\prime pq})~;
\]
\oldpage[0\textsubscript{III}]{26}
E 与 $E'$ 正则这一假设还蕴含
$u_2^{pq}(Z_\infty(E_2^{pq}))=Z_\infty(E_2^{\prime pq})$。由于
$u_2^{pq}$ 是同构，$u_\infty^{pq}$ 也是同构；因而
$\operatorname{gr}_p(u^{p+q})$ 也是同构。可是，$E^n$ 和
$E^{\prime n}$ 的滤过都是离散且穷尽的，所以诸 $u^n$ 也是同构
（Bourbaki，\emph{Alg. comm.}，第 III 章，\S~2，
n\textsuperscript{o}~8，定理 1）。
\end{proof}





\begin{env}[11.1.6]
\label{0.11.1.6-fr}
由（\hyperref[0.11.1.1.2-fr]{11.1.1.2}）及定义
（\hyperref[0.11.1.1.3-fr]{11.1.1.3}）可知：若谱序列 E 中
$E_r^{pq}=0$，则对 $k\geq r$ 有 $E_k^{pq}=0$，并且
$E_\infty^{pq}=0$。若存在整数 $r\geq2$，并且对每个整数
$n\in\mathbb Z$ 都存在整数 $q(n)$，使得对每个 $q\ne q(n)$ 有
$E_r^{n-q,q}=0$，则称谱序列\emph{退化}。首先由上面的说明可知，对
$k\geq r$（包括 $k=\infty$）和 $q\ne q(n)$，也有
$E_k^{n-q,q}=0$。此外，由 $E_{r+1}^{pq}$ 的定义可知
$E_{r+1}^{n-q(n),q(n)}=E_r^{n-q(n),q(n)}$；若 E 弱收敛，则还有
$E_\infty^{n-q(n),q(n)}=E_r^{n-q(n),q(n)}$。换言之，对每个
% Canon current rebase 2026-08-31: fixed-total-degree filtration index.
$n\in\mathbb Z$，当 $p\ne n-q(n)$ 时
$\operatorname{gr}_p(E^n)=0$，并且
$\operatorname{gr}_{n-q(n)}(E^n)=E_r^{n-q(n),q(n)}$。若此外 $E^n$
的滤过离散且穷尽，则 E 是正则谱序列，并且在同构意义下
$E^n=E_r^{n-q(n),q(n)}$。
\end{env}

\begin{env}[11.1.7]
\label{0.11.1.7-fr}
假设范畴 $\mathcal C$ 中存在滤过归纳极限且它们正合，并设
$(E_\lambda,u_{\mu\lambda})$ 是 $\mathcal C$ 中谱序列的一个归纳系统
（其指标集是滤过集）。那么，这一归纳系统在 $\mathcal C$ 的对象的谱序列所成的
加性范畴中存在\emph{归纳极限}。事实上，只须把
$E_r^{pq}$、$d_r^{pq}$、$\alpha_r^{pq}$、
$B_\infty(E_2^{pq})$、$Z_\infty(E_2^{pq})$、$E^n$、
$F^p(E^n)$ 和 $\beta^{pq}$ 分别定义为
$E_{r,\lambda}^{pq}$、$d_{r,\lambda}^{pq}$、
$\alpha_{r,\lambda}^{pq}$、$B_\infty(E_{2,\lambda}^{pq})$、
$Z_\infty(E_{2,\lambda}^{pq})$、$E_\lambda^n$、
$F^p(E_\lambda^n)$ 和 $\beta_\lambda^{pq}$ 的归纳极限；
（\hyperref[0.11.1.1-fr]{11.1.1}）诸条件的验证由
$\mathcal C$ 中函子 $\varinjlim$ 的正合性得出。
\end{env}

\begin{remark}[11.1.8]
\label{0.11.1.8-fr}
假设范畴 $\mathcal C$ 是一个\emph{Noether} 环 $A$ 上的 $A$-模范畴
（相应地，是一个环 $A$ 上的模范畴）。由
（\hyperref[0.11.1.1-fr]{11.1.1}）的定义可知：如果对某个给定的 $r$，
诸 $E_r^{pq}$ 都是\emph{有限型} $A$-模（相应地，具有\emph{有限长度}），
那么对所有 $s\geq r$，诸 $E_s^{pq}$ 也都如此，
$E_\infty^{pq}$ 亦然。若此外收敛项 $(E^n)$ 的滤过对每个 $n$
都离散且余离散，则可断定每个 $E^n$ 也是\emph{有限型} $A$-模
（相应地，具有\emph{有限长度}）。
\end{remark}


\begin{env}[11.1.9]
\label{0.11.1.9-fr}
我们将要考察一些条件，以保证谱序列 E 关于 $p+q=n$ 以“统一”的方式双正则。
为此将使用下列引理：
\end{env}

\begin{lemma}[11.1.10]
\label{0.11.1.10-fr}
设 $(E_r^{pq})$ 是 $\mathcal C$ 的一族对象，它们由
（\hyperref[0.11.1.1-fr]{11.1.1}）的资料 a)、b)、c) 联系起来。
对固定的整数 $n$，下列性质等价：
\begin{enumerate}
  \item[(a)] 存在整数 $r(n)$，使得当 $r\geq r(n)$，且
    $p+q=n$ 或 $p+q=n-1$ 时，所有态射 $d_r^{pq}$ 都为零。
  \item[(b)] 存在整数 $r(n)$，使得当 $p+q=n$ 或 $p+q=n+1$ 时，
    对 $s\geq r\geq r(n)$ 有
    $B_r(E_2^{pq})=B_s(E_2^{pq})$。
  \item[(c)] 存在整数 $r(n)$，使得当 $p+q=n$ 或 $p+q=n-1$ 时，
    对 $s\geq r\geq r(n)$ 有
    $Z_r(E_2^{pq})=Z_s(E_2^{pq})$。
  \item[(d)] 存在整数 $r(n)$，使得当 $p+q=n$ 时，对
    $s\geq r\geq r(n)$ 有
    $B_r(E_2^{pq})=B_s(E_2^{pq})$ 且
    $Z_r(E_2^{pq})=Z_s(E_2^{pq})$。
\end{enumerate}
\end{lemma}

\begin{proof}
\oldpage[0\textsubscript{III}]{27}
事实上，由（\hyperref[0.11.1.1-fr]{11.1.1}）的条件 a)、b)、c)，
$Z_{r+1}(E_2^{pq})=Z_r(E_2^{pq})$ 等价于 $d_r^{pq}=0$；而
$B_r(E_2^{p+r,q-r+1})=B_{r+1}(E_2^{p+r,q-r+1})$
也等价于 $d_r^{pq}=0$。引理立即由此得出。
\end{proof}

\subsection{滤过复形的谱序列}
\label{subsection:0.11.2-fr}

\begin{env}[11.2.1]
\label{0.11.2.1-fr}
给定一个阿贝尔范畴 $\mathcal C$，我们约定用 $K^\bullet$ 之类的记号
表示 $\mathcal C$ 的对象组成的\emph{上链复形} $(K^i)_{i\in\mathbb Z}$，其中微分算子
的次数为 $+1$；用 $K_\bullet$ 之类的记号表示 $\mathcal C$ 的对象
组成的链复形 $(K_i)_{i\in\mathbb Z}$，其中微分算子的次数为 $-1$。对任意
微分算子 $d$ 为 $+1$ 次的上链复形 $K^\bullet=(K^i)$，可以通过置
$K'_i=K^{-i}$ 联系一个链复形 $K'_\bullet=(K'_i)$，其微分算子
$K'_i\to K'_{i-1}$ 就是算子 $d:K^{-i}\to K^{-i+1}$；\emph{反之亦然}。因此，
可按具体情形考虑这两类复形中的任意一类，并把其中一类的任意结果
转译为另一类的结果。同样，我们用 $K^{\bullet\bullet}=(K^{ij})$（分别地，
$K_{\bullet\bullet}=(K_{ij})$）之类的记号表示 $\mathcal C$ 的对象组成的
\emph{双复形}，其中\emph{两个}微分算子的次数都是 $+1$（分别地，$-1$）；
变更指标的符号仍可从一类转到另一类，而任意多重复形也有类似的记号
和注意事项。记号 $K^\bullet$ 或 $K_\bullet$ 也用于 $\mathcal C$ 中 $\mathbb Z$ 型的
\emph{分次对象}，它们不必是复形（或者可在微分算子为零时将其视为复形）；
例如，我们写
$\mathrm H^\bullet(K^\bullet)=(\mathrm H^i(K^\bullet))_{i\in\mathbb Z}$ 表示微分算子
为 $+1$ 次的上链复形 $K^\bullet$ 的\emph{上同调}，写
$\mathrm H_\bullet(K_\bullet)=(\mathrm H_i(K_\bullet))_{i\in\mathbb Z}$ 表示微分算子为
$-1$ 次的链复形 $K_\bullet$ 的\emph{同调}；当按上述操作从 $K^\bullet$ 转到
$K'_\bullet$ 时，有 $\mathrm H_i(K'_\bullet)=\mathrm H^{-i}(K^\bullet)$。

在此回顾：对一个上链复形 $K^\bullet$（分别地，链复形 $K_\bullet$），我们通常写
$Z^i(K^\bullet)=\operatorname{Ker}(K^i\to K^{i+1})$（“上闭链对象”）和
$B^i(K^\bullet)=\operatorname{Im}(K^{i-1}\to K^i)$（“上边缘对象”）
（分别地，$Z_i(K_\bullet)=\operatorname{Ker}(K_i\to K_{i-1})$（“闭链对象”）和
$B_i(K_\bullet)=\operatorname{Im}(K_{i+1}\to K_i)$（“边缘对象”）），从而
$\mathrm H^i(K^\bullet)=Z^i(K^\bullet)/B^i(K^\bullet)$（分别地，
$\mathrm H_i(K_\bullet)=Z_i(K_\bullet)/B_i(K_\bullet)$）。

若 $K^\bullet=(K^i)$（分别地，$K_\bullet=(K_i)$）是 $\mathcal C$ 中的一个上链复形（分别地，链复形），
而 $T:\mathcal C\to\mathcal C'$ 是从 $\mathcal C$ 到另一个阿贝尔范畴 $\mathcal C'$ 的函子，
则用 $T(K^\bullet)$（分别地，$T(K_\bullet)$）表示 $\mathcal C'$ 中的上链复形（分别地，链复形）
$(T(K^i))$（分别地，$(T(K_i))$）。

我们不再回顾 $\partial$-函子（T, 2.1）的定义，只指出：当态射 $\partial$
使次数降低一个单位时，我们也用 $\partial$-函子代替 $\partial^*$-函子；
如果可能产生疑义，每次由上下文说明这一点。

最后，若 $\mathcal C$ 的一个\emph{分次对象} $(A_i)_{i\in\mathbb Z}$ 存在 $i_0$，
使得对 $i<i_0$（分别地，$i>i_0$）有 $A_i=0$，则称它\emph{下有界}
（分别地，\emph{上有界}）。
\end{env}

\begin{env}[11.2.2]
\label{0.11.2.2-fr}
设 $K^\bullet$ 是 $\mathcal C$ 中的一个上链复形，其微分算子 $d$ 为 $+1$ 次，并假设
它配备了一个\emph{滤过} $F(K^\bullet)=(F^p(K^\bullet))_{p\in\mathbb Z}$，该滤过由
$K^\bullet$ 的分次子对象构成，也就是说
\oldpage[0\textsubscript{III}]{28}
$F^p(K^\bullet)=(K^i\cap F^p(K^\bullet))_{i\in\mathbb Z}$；此外，还假设对每个
$p\in\mathbb Z$ 都有 $d(F^p(K^\bullet))\subset F^p(K^\bullet)$。现快速回顾如何由
$K^\bullet$ \emph{函子地}定义一个谱序列 $E(K^\bullet)$（M, XV, 4 和 G, I, 4.3）。
对 $r\geq2$，典范态射
$F^p(K^\bullet)/F^{p+r}(K^\bullet)\to
F^p(K^\bullet)/F^{p+1}(K^\bullet)$ 在上同调上定义态射
\[
  \mathrm H^{p+q}(F^p(K^\bullet)/F^{p+r}(K^\bullet))
  \longrightarrow
  \mathrm H^{p+q}(F^p(K^\bullet)/F^{p+1}(K^\bullet)).
\]
用 $Z_r^{pq}(K^\bullet)$ 表示这个态射的像。同样，由正合列
\[
  0\to F^p(K^\bullet)/F^{p+1}(K^\bullet)
  \to F^{p-r+1}(K^\bullet)/F^{p+1}(K^\bullet)
  \to F^{p-r+1}(K^\bullet)/F^p(K^\bullet)\to0
\]
可通过上同调正合列得到一个态射
\[
  \mathrm H^{p+q-1}(F^{p-r+1}(K^\bullet)/F^p(K^\bullet))
  \longrightarrow
  \mathrm H^{p+q}(F^p(K^\bullet)/F^{p+1}(K^\bullet))
\]
用 $B_r^{pq}(K^\bullet)$ 表示这个态射的像；可证明
$B_r^{pq}(K^\bullet)\subset Z_r^{pq}(K^\bullet)$，并置
$E_r^{pq}(K^\bullet)=Z_r^{pq}(K^\bullet)/B_r^{pq}(K^\bullet)$；我们不详述
$d_r^{pq}$ 和 $\alpha_r^{pq}$ 的定义。

注意，对固定的 $p$、$q$，所有 $Z_r^{pq}(K^\bullet)$ 和 $B_r^{pq}(K^\bullet)$
都是同一对象
$\mathrm H^{p+q}(F^p(K^\bullet)/F^{p+1}(K^\bullet))$ 的子对象，该对象记为
$Z_1^{pq}(K^\bullet)$；置 $B_1^{pq}(K^\bullet)=0$，于是上述 $Z_r^{pq}(K^\bullet)$
和 $B_r^{pq}(K^\bullet)$ 的定义也适用于 $r=1$；再置
$E_1^{pq}(K^\bullet)=Z_1^{pq}(K^\bullet)$。定义 $d_1^{pq}$ 和 $\alpha_1^{pq}$，
使（\hyperref[0.11.1.1-fr]{11.1.1}）的条件在 $r=1$ 时成立。另一方面，
定义子对象 $Z_\infty^{pq}(K^\bullet)$ 为态射
\[
  \mathrm H^{p+q}(F^p(K^\bullet))\longrightarrow
  \mathrm H^{p+q}(F^p(K^\bullet)/F^{p+1}(K^\bullet))
  =E_1^{pq}(K^\bullet)
\]
的像；定义 $B_\infty^{pq}(K^\bullet)$ 为态射
\[
  \mathrm H^{p+q-1}(K^\bullet/F^p(K^\bullet))\longrightarrow
  \mathrm H^{p+q}(F^p(K^\bullet)/F^{p+1}(K^\bullet))
  =E_1^{pq}(K^\bullet)
\]
的像，这个态射如上由某个上同调正合列得出。取
$Z_\infty(E_2^{pq}(K^\bullet))$ 和 $B_\infty(E_2^{pq}(K^\bullet))$ 分别为
$Z_\infty^{pq}(K^\bullet)$ 和 $B_\infty^{pq}(K^\bullet)$ 在 $E_2^{pq}(K^\bullet)$ 中的典范像。

最后，用 $F^p(\mathrm H^n(K^\bullet))$ 表示由典范单射
$F^p(K^\bullet)\to K^\bullet$ 所给出的态射
$\mathrm H^n(F^p(K^\bullet))\to\mathrm H^n(K^\bullet)$ 在 $\mathrm H^n(K^\bullet)$ 中的像；
由上同调正合列，这也是态射
$\mathrm H^n(K^\bullet)\to\mathrm H^n(K^\bullet/F^p(K^\bullet))$ 的核。这样就在
$E^n(K^\bullet)=\mathrm H^n(K^\bullet)$ 上定义了一个滤过；我们在此也不给出同构
$\beta^{pq}$ 的定义。
\end{env}

\begin{env}[11.2.3]
\label{0.11.2.3-fr}
$E(K^\bullet)$ 的\emph{函子性}应作如下理解：给定 $\mathcal C$ 中两个
\emph{滤过}上链复形 $K^\bullet$、$K^{\prime\bullet}$，以及一个与滤过\emph{相容}的上链复形态射
$u:K^\bullet\to K^{\prime\bullet}$，则以显然的方式得到态射 $u_r^{pq}$（对 $r\geq1$）
和 $u^n$，并可证明这些态射按（\hyperref[0.11.1.2-fr]{11.1.2}）的意义与
$d_r^{pq}$、$\alpha_r^{pq}$ 和 $\beta^{pq}$ 相容，因而确实定义了谱序列的一个态射
$E(u):E(K^\bullet)\to E(K^{\prime\bullet})$。此外，可证明：若 $u$ 和 $v$ 是上述类型的
$K^\bullet\to K^{\prime\bullet}$ 态射，并且\emph{以阶数 $\leq k$ 同伦}，那么对 $r>k$
有 $u_r^{pq}=v_r^{pq}$，而对所有 $n$ 有 $u^n=v^n$（M, XV, 3.1）。
\end{env}

\begin{env}[11.2.4]
\label{0.11.2.4-fr}
\oldpage[0\textsubscript{III}]{29}
假设 $\mathcal C$ 中的滤过余极限正合。那么，若 $K^\bullet$ 的滤过
$(F^p(K^\bullet))$ \emph{穷尽}，则对每个 $n$，滤过
$(F^p(\mathrm H^n(K^\bullet)))$ 也穷尽；因为依假设有
$K^\bullet=\varinjlim_{p\to-\infty}F^p(K^\bullet)$，而关于 $\mathcal C$ 的假设导致上同调
与余极限可交换。此外，由同样的理由，有
$B_\infty(E_2^{pq}(K^\bullet))=\sup_k B_k(E_2^{pq}(K^\bullet))$。若对每个 $n$ 都存在
一个整数 $u(n)$，使得对 $p>u(n)$ 有 $\mathrm H^n(F^p(K^\bullet))=0$，则称
$K^\bullet$ 的滤过 $(F^p(K^\bullet))$ \emph{正则}。特别地，当 $K^\bullet$ 的滤过
\emph{离散}时就是如此。当 $K^\bullet$ 的滤过正则且穷尽，而且 $\mathcal C$ 中的
滤过余极限正合时，可证明（M, XV, 4）谱序列 $E(K^\bullet)$ \emph{正则}。
\end{env}

\subsection{双复形的谱序列}
\label{subsection:0.11.3-fr}

\begin{env}[11.3.1]
\label{0.11.3.1-fr}
关于双复形的约定，我们采用（T, 2.4）的约定，而不采用（M）的约定；
因此，双复形 $K^{\bullet\bullet}=(K^{ij})$ 的两个（次数为 $+1$ 的）微分算子
$d'$、$d''$ 假设为\emph{可交换}的。假设下列两个条件之一成立：
\begin{enumerate}
  \item[$1^\circ$] $\mathcal C$ 中存在无限直和。
  \item[$2^\circ$] 对每个 $n\in\mathbb Z$，满足 $p+q=n$ 且 $K^{pq}\ne0$
    的偶对 $(p,q)$ 只有\emph{有限}多个。
\end{enumerate}
那么，双复形 $K^{\bullet\bullet}$ 定义一个（单）复形
$(K^{\prime n})_{n\in\mathbb Z}$，其中
$K^{\prime n}=\sum_{i+j=n}K^{ij}$；这个复形的（次数为 $+1$ 的）微分算子 $d$
由下式给出：对 $x\in K^{ij}$，有
$dx=d'x+(-1)^i d''x$。此后，当我们谈到\emph{由双复形
$K^{\bullet\bullet}$ 定义的（单）复形}时，总是默认上述条件之一成立。
对于多重复形，采用类似的约定。

用 $K^{i,\bullet}$（分别地，$K^{\bullet,j}$）表示单复形
$(K^{ij})_{j\in\mathbb Z}$（分别地，$(K^{ij})_{i\in\mathbb Z}$），并用
$Z_{\mathrm{II}}^p(K^{i,\bullet})$、$B_{\mathrm{II}}^p(K^{i,\bullet})$、
$\mathrm H_{\mathrm{II}}^p(K^{i,\bullet})$（分别地，
$Z_{\mathrm I}^p(K^{\bullet,j})$、$B_{\mathrm I}^p(K^{\bullet,j})$、
$\mathrm H_{\mathrm I}^p(K^{\bullet,j})$）分别表示它的第 $p$ 个上闭链对象、
上边缘对象和上同调对象；微分算子
$d':K^{i,\bullet}\to K^{i+1,\bullet}$ 是复形态射，因而在上闭链、上边缘和
上同调上给出算子
\[
\begin{aligned}
  d'&:Z_{\mathrm{II}}^p(K^{i,\bullet})
      \longrightarrow Z_{\mathrm{II}}^p(K^{i+1,\bullet}),\\
  d'&:B_{\mathrm{II}}^p(K^{i,\bullet})
      \longrightarrow B_{\mathrm{II}}^p(K^{i+1,\bullet}),\\
  d'&:\mathrm H_{\mathrm{II}}^p(K^{i,\bullet})
      \longrightarrow \mathrm H_{\mathrm{II}}^p(K^{i+1,\bullet}),
\end{aligned}
\]
并且显然，对于这些算子，
$(Z_{\mathrm{II}}^p(K^{i,\bullet}))_{i\in\mathbb Z}$、
$(B_{\mathrm{II}}^p(K^{i,\bullet}))_{i\in\mathbb Z}$ 和
$(\mathrm H_{\mathrm{II}}^p(K^{i,\bullet}))_{i\in\mathbb Z}$ 都是复形；我们用
$\mathrm H_{\mathrm{II}}^p(K^{\bullet\bullet})$ 表示复形
$(\mathrm H_{\mathrm{II}}^p(K^{i,\bullet}))_{i\in\mathbb Z}$，并用
$Z_{\mathrm I}^q(\mathrm H_{\mathrm{II}}^p(K^{\bullet\bullet}))$、
$B_{\mathrm I}^q(\mathrm H_{\mathrm{II}}^p(K^{\bullet\bullet}))$ 和
$\mathrm H_{\mathrm I}^q(\mathrm H_{\mathrm{II}}^p(K^{\bullet\bullet}))$
分别表示它的第 $q$ 个上闭链对象、上边缘对象和上同调对象。同样定义复形
$\mathrm H_{\mathrm I}^p(K^{\bullet\bullet})$ 及其上同调对象
$\mathrm H_{\mathrm{II}}^q(\mathrm H_{\mathrm I}^p(K^{\bullet\bullet}))$。
另一方面，回顾 $\mathrm H^n(K^{\bullet\bullet})$ 表示由
$K^{\bullet\bullet}$ 定义的（单）复形的第 $n$ 个上同调对象。
\end{env}

\begin{env}[11.3.2]
\label{0.11.3.2-fr}
在由双复形 $K^{\bullet\bullet}$ 定义的复形上，可以考虑两个典范滤过
$(F_{\mathrm I}^p(K^{\bullet\bullet}))$ 和
$(F_{\mathrm{II}}^p(K^{\bullet\bullet}))$，它们由下式给出：
\[
\label{0.11.3.2.1-fr}
  F_{\mathrm I}^p(K^{\bullet\bullet})
  =\left(\sum_{\substack{i+j=n\\i\geq p}}K^{ij}\right)_{n\in\mathbb Z},
  \qquad
  F_{\mathrm{II}}^p(K^{\bullet\bullet})
  =\left(\sum_{\substack{i+j=n\\j\geq p}}K^{ij}\right)_{n\in\mathbb Z}.
  \tag{11.3.2.1}
\]
\oldpage[0\textsubscript{III}]{30}
按照定义，它们确是由 $K^{\bullet\bullet}$ 定义的（单）复形的分次子对象，
因而使这个复形成为滤过复形；此外，显然这些滤过是穷尽且分离的。

每个滤过都对应一个谱序列
（\hyperref[0.11.2.2-fr]{11.2.2}）；用
$'E(K^{\bullet\bullet})$ 和 $''E(K^{\bullet\bullet})$ 分别表示与
$(F_{\mathrm I}^p(K^{\bullet\bullet}))$ 和
$(F_{\mathrm{II}}^p(K^{\bullet\bullet}))$ 对应的谱序列，称为
\emph{双复形 $K^{\bullet\bullet}$ 的谱序列}；二者都以上同调
$(\mathrm H^n(K^{\bullet\bullet}))$ 为收敛项。此外，可证明（M, XV, 6）
\[
\label{0.11.3.2.2-fr}
  'E_2^{pq}(K^{\bullet\bullet})
  =\mathrm H_{\mathrm I}^p(\mathrm H_{\mathrm{II}}^q(K^{\bullet\bullet})),
  \qquad
  ''E_2^{pq}(K^{\bullet\bullet})
  =\mathrm H_{\mathrm{II}}^q(\mathrm H_{\mathrm I}^p(K^{\bullet\bullet})).
  \tag{11.3.2.2}
\]

双复形的任意态射 $u:K^{\bullet\bullet}\to K^{\prime\bullet\bullet}$
自然与 $K^{\bullet\bullet}$ 和 $K^{\prime\bullet\bullet}$ 的同类型滤过相容，
因而对两个谱序列分别定义一个态射；此外，两个\emph{同伦}态射定义相应滤过
（单）复形之间的一个\emph{阶数 $\leq1$ 的同伦}，因此对两个谱序列分别定义
\emph{同一个}态射（M, XV, 6.1）。
\end{env}

\begin{proposition}[11.3.3]
\label{0.11.3.3-fr}
设 $K^{\bullet\bullet}=(K^{ij})$ 是阿贝尔范畴 $\mathcal C$ 中的一个双复形。
\begin{enumerate}
  \item[(i)] 若存在 $i_0$ 和 $j_0$，使得当 $i<i_0$ 或 $j<j_0$
    （分别地，当 $i>i_0$ 或 $j>j_0$）时有 $K^{ij}=0$，则两个谱序列
    $'E(K^{\bullet\bullet})$ 和 $''E(K^{\bullet\bullet})$ 都是双正则的。
  \item[(ii)] 若存在 $i_0$ 和 $i_1$，使得当 $i<i_0$ 或 $i>i_1$ 时有
    $K^{ij}=0$（分别地，若存在 $j_0$ 和 $j_1$，使得当 $j<j_0$ 或
    $j>j_1$ 时有 $K^{ij}=0$），则两个谱序列
    $'E(K^{\bullet\bullet})$ 和 $''E(K^{\bullet\bullet})$ 都是双正则的。
\end{enumerate}
再假设 $\mathcal C$ 中存在滤过余极限且它们正合。那么：
\begin{enumerate}
  \item[(iii)] 若存在 $i_0$，使得当 $i>i_0$ 时有 $K^{ij}=0$
    （分别地，若存在 $j_0$，使得当 $j<j_0$ 时有 $K^{ij}=0$），则谱序列
    $'E(K^{\bullet\bullet})$ 正则。
  \item[(iv)] 若存在 $i_0$，使得当 $i<i_0$ 时有 $K^{ij}=0$
    （分别地，若存在 $j_0$，使得当 $j>j_0$ 时有 $K^{ij}=0$），则谱序列
    $''E(K^{\bullet\bullet})$ 正则。
\end{enumerate}
\end{proposition}

\begin{proof}
该命题立即由定义（\hyperref[0.11.1.3-fr]{11.1.3}）、
（\hyperref[0.11.2.4-fr]{11.2.4}）以及下列关于滤过 $F_{\mathrm I}$ 的观察得出
（交换 $K^{\bullet\bullet}$ 的两个指标的作用，即得到关于 $F_{\mathrm{II}}$ 的类似观察）：
\begin{enumerate}
  \item[$1^\circ$] 若存在 $i_0$，使得当 $i>i_0$ 时有 $K^{ij}=0$，则滤过
    $F_{\mathrm I}(K^{\bullet\bullet})$ 离散。
  \item[$2^\circ$] 若存在 $i_0$，使得当 $i<i_0$ 时有 $K^{ij}=0$，则滤过
    $F_{\mathrm I}(K^{\bullet\bullet})$ 余离散。立即可推出，对每个 $n$，相应的滤过
    $F_{\mathrm I}(\mathrm H^n(K^{\bullet\bullet}))$ 也余离散；此外，对应于滤过
    $F_{\mathrm I}(K^{\bullet\bullet})$ 的 $B_r^{pq}$ 的定义
    （\hyperref[0.11.2.2-fr]{11.2.2}）表明，对每个偶对 $(p,q)$，序列
    $(B_r^{pq})_{r\geq2}$ 稳定。
  \item[$3^\circ$] 若存在 $j_0$，使得当 $j<j_0$ 时有 $K^{ij}=0$，则
    \[
      F_{\mathrm I}^{p+r}(K^{\bullet\bullet})
      \cap\left(\sum_{i+j=n}K^{ij}\right)=0
    \]
    只要 $p+r+j_0>n$；因此，当 $r>q-j_0+1$ 时有
    $Z_r^{pq}=Z_\infty(E_2^{pq})$；另一方面，当 $p>n-j_0+1$ 时有
    $\mathrm H^n(F_{\mathrm I}^p(K^{\bullet\bullet}))=0$。
  \item[$4^\circ$] 若存在 $j_0$，使得当 $j>j_0$ 时有 $K^{ij}=0$，则
    \[
      F_{\mathrm I}^{p-r+1}(K^{\bullet\bullet})
      \cap\left(\sum_{i+j=n}K^{ij}\right)=\sum_{i+j=n}K^{ij}.
    \]
    \oldpage[0\textsubscript{III}]{31}
    只要 $p-r+1+j_0<n$；因此，当 $r>j_0-q+1$ 时有
    $B_r^{pq}=B_\infty(E_2^{pq})$
    \SourceCorrectionNote{法文印本在结论中写作 $r<j_0-q+1$；但因
    $n=p+q$，前述条件 $p-r+1+j_0<n$ 等价于 $r>j_0-q+1$。
    本译依稳定勘误 EGA-III1-11.3.3-P375-AUTHORIAL-R-INEQUALITY-001
    及 D27 交叉核验采用校正后的不等号；法文外交式来源保持不变。}；另一方面，当 $p+j_0<n-1$ 时有
    $\mathrm H^n(F_{\mathrm I}^p(K^{\bullet\bullet}))
    =\mathrm H^n(K^{\bullet\bullet})$。
\end{enumerate}
\end{proof}

\begin{env}[11.3.4]
\label{0.11.3.4-fr}
假设双复形 $K^{\bullet\bullet}=(K^{ij})$ 满足：当 $i<0$ 或 $j<0$ 时有
$K^{ij}=0$。已知这时对每个 $p\in\mathbb Z$ 都可以定义一个典范“边缘同态”
\[
\label{0.11.3.4.1-fr}
  'E_2^{p0}(K^{\bullet\bullet})\longrightarrow
  \mathrm H^p(K^{\bullet\bullet})
  \tag{11.3.4.1}
\]
（M, XV, 6）。快速回顾其原因：一方面，在谱序列
$'E(K^{\bullet\bullet})$ 中，对 $2\leq r\leq+\infty$ 有
$Z_r^{p0}=Z_{\mathrm I}^p(Z_{\mathrm{II}}^0(K^{\bullet\bullet}))$；另一方面，
$\mathrm H^p(F_{\mathrm I}^{p+1}(K^{\bullet\bullet}))=0$，所以同构
$\beta^{p0}:'E_\infty^{p0}\xrightarrow{\sim}
\mathrm H^p(F_{\mathrm I}^p)/\mathrm H^p(F_{\mathrm I}^{p+1})$ 给出同态
$'E_\infty^{p0}\to\mathrm H^p(F_{\mathrm I}^p(K^{\bullet\bullet}))
\to\mathrm H^p(K^{\bullet\bullet})$。所有 $Z_r^{p0}$ 相等，于是可以定义
$r\leq s$ 时的典范同态 $'E_r^{p0}\to'E_s^{p0}$，特别是同态
$'E_2^{p0}\to'E_\infty^{p0}$；与前述同态复合，就得到边缘同态
$'E_2^{p0}\to\mathrm H^p(K^{\bullet\bullet})$。此外，立即可验证：对元素
$z\in Z_{\mathrm{II}}^0(K^{\bullet\bullet})\subset K^{p0}$ 且 $d'z=0$ 的模
$B_2^{p0}$ 类，这样定义的边缘同态先使它对应于 $'E_\infty^{p0}$ 中 $z$ 的模
$B_\infty^{p0}$ 类，再使后者对应于 $\mathrm H^p(K^{\bullet\bullet})$ 中 $z$ 的
上同调类。因此，最终可见，\emph{边缘同态
（\hyperref[0.11.3.4.1-fr]{11.3.4.1}）是由典范单射}
$Z_{\mathrm{II}}^0(K^{\bullet\bullet})\to K^{\bullet\bullet}$
\emph{取上同调而得的}（这里把 $K^{\bullet\bullet}$ 视为单复形）。同样，边缘同态
\[
\label{0.11.3.4.2-fr}
  ''E_2^{p0}(K^{\bullet\bullet})\longrightarrow
  \mathrm H^p(K^{\bullet\bullet})
  \tag{11.3.4.2}
\]
自然可解释为由典范单射
$Z_{\mathrm I}^0(K^{\bullet\bullet})\to K^{\bullet\bullet}$ 得出。
\end{env}

\begin{env}[11.3.5]
\label{0.11.3.5-fr}
现设 $K_{\bullet\bullet}=(K_{ij})$ 是 $\mathcal C$ 的一个双复形，其两个微分算子
的次数均为 $-1$。用 $K_{i,\bullet}$（分别地，$K_{\bullet,j}$）表示单复形
$(K_{ij})_{j\in\mathbb Z}$（分别地，$(K_{ij})_{i\in\mathbb Z}$），用
$\mathrm H_p^{\mathrm{II}}(K_{i,\bullet})$（分别地，
$\mathrm H_p^{\mathrm I}(K_{\bullet,j})$）表示其第 $p$ 个同调对象，用
$\mathrm H_p^{\mathrm{II}}(K_{\bullet\bullet})$（分别地，
$\mathrm H_p^{\mathrm I}(K_{\bullet\bullet})$）表示复形
$(\mathrm H_p^{\mathrm{II}}(K_{i,\bullet}))_{i\in\mathbb Z}$（分别地，
$(\mathrm H_p^{\mathrm I}(K_{\bullet,j}))_{j\in\mathbb Z}$），用
$\mathrm H_q^{\mathrm I}(\mathrm H_p^{\mathrm{II}}(K_{\bullet\bullet}))$
（分别地，
$\mathrm H_q^{\mathrm{II}}(\mathrm H_p^{\mathrm I}(K_{\bullet\bullet}))$）
表示它的第 $q$ 个同调对象；对闭链对象和边缘对象使用类似的记号；最后，用
$\mathrm H_n(K_{\bullet\bullet})$ 表示（当它存在时）由 $K_{\bullet\bullet}$ 定义的
（微分算子次数为 $-1$ 的）单复形的第 $n$ 个同调对象。

设 $K^{\prime\bullet\bullet}=(K^{\prime ij})$，其中
$K^{\prime ij}=K_{-i,-j}$，这是与 $K_{\bullet\bullet}$ 相联系的微分算子次数为
$+1$ 的双复形；按照定义，$K_{\bullet\bullet}$ 的\emph{谱序列}是
$K^{\prime\bullet\bullet}$ 的谱序列，记为 $'E(K_{\bullet\bullet})$ 和
$''E(K_{\bullet\bullet})$，但我们把记号改写为
\[
  'E_{pq}^r(K_{\bullet\bullet})
  ='E_r^{-p,-q}(K^{\prime\bullet\bullet}),\qquad
  ''E_{pq}^r(K_{\bullet\bullet})
  =''E_r^{-p,-q}(K^{\prime\bullet\bullet})
\]
（$2\leq r\leq\infty$）。用这些记号，有
\[
  'E_{pq}^2(K_{\bullet\bullet})
  =\mathrm H_p^{\mathrm I}(\mathrm H_q^{\mathrm{II}}(K_{\bullet\bullet})),
  \qquad
  ''E_{pq}^2(K_{\bullet\bullet})
  =\mathrm H_q^{\mathrm{II}}(\mathrm H_p^{\mathrm I}(K_{\bullet\bullet})).
\]
为避免符号错误，在讨论这些谱序列与其收敛项之间的关系时，一般最好回到复形
$K^{\prime\bullet\bullet}$。不过，记下与
（\hyperref[0.11.3.3-fr]{11.3.3}）相对应的判别准则：
\end{env}

\begin{env}[11.3.6]
\label{0.11.3.6-fr}
在下列情形中，谱序列 $'E(K_{\bullet\bullet})$ 和
$''E(K_{\bullet\bullet})$ \emph{双正则}：
\begin{enumerate}
  \item[(a)] 存在 $i_0$ 和 $j_0$，使得当 $i>i_0$ 或 $j>j_0$
    （分别地，当 $i<i_0$
    \oldpage[0\textsubscript{III}]{32}
    或 $j<j_0$）时有 $K_{ij}=0$。
  \item[(b)] 存在 $i_0$ 和 $i_1$，使得对 $i<i_0$ 以及 $i>i_1$ 均有
    $K_{ij}=0$。
  \item[(c)] 存在 $j_0$ 和 $j_1$，使得对 $j<j_0$ 以及 $j>j_1$ 均有
    $K_{ij}=0$。
\end{enumerate}

若存在 $i_0$，使得当 $i<i_0$ 时有 $K_{ij}=0$，或者存在 $j_0$，使得当
$j>j_0$ 时有 $K_{ij}=0$，则谱序列 $'E(K_{\bullet\bullet})$ 正则。

若存在 $i_0$，使得当 $i>i_0$ 时有 $K_{ij}=0$，或者存在 $j_0$，使得当
$j<j_0$ 时有 $K_{ij}=0$，则谱序列 $''E(K_{\bullet\bullet})$ 正则。
\end{env}

\subsection{函子关于复形 $K^\bullet$ 的超上同调}
\label{subsection:0.11.4-fr}

\begin{env}[11.4.1]
\label{0.11.4.1-fr}
设 $\mathcal C$ 是一个阿贝尔范畴；回顾，$\mathcal C$ 的对象 $A$ 的一个
\emph{右消解}（或\emph{上同调消解}）是 $\mathcal C$ 的对象组成的复形，其微分
算子的次数为 $+1$，
\[
  0\longrightarrow L^0\longrightarrow L^1\longrightarrow L^2
  \longrightarrow\cdots
\]
并配备一个态射 $\varepsilon:A\to L^0$，称为消解的\emph{增广}（也可把它视为
复形态射
\[
\xymatrix{
  0\ar[r] & A\ar[r]\ar[d]_{\varepsilon} & 0\ar[r]\ar[d]
    & 0\ar[r]\ar[d] & \cdots\\
  0\ar[r] & L^0\ar[r] & L^1\ar[r] & L^2\ar[r] & \cdots
})
\]
使得序列
\[
  0\longrightarrow A\xrightarrow{\varepsilon}L^0\longrightarrow L^1
  \longrightarrow\cdots
\]
正合；同样，$A$ 的一个\emph{左消解}（或\emph{同调消解}）是
$\mathcal C$ 的对象组成的复形
$0\leftarrow L_0\leftarrow L_1\leftarrow\cdots$，其微分算子的次数为 $-1$，
并配备一个增广 $\varepsilon:L_0\to A$，使得序列
\[
  0\longleftarrow A\xleftarrow{\varepsilon}L_0\longleftarrow L_1
  \longleftarrow\cdots
\]
正合。

当对象 $A$ 的一个右消解 $(L^i)_{i\geq0}$ 满足对 $i\geq n+1$ 有
$L^i=0$ 时，称这个消解的\emph{长度 $\leq n$}
\SourceCorrectionNote{法文印本及外交式转录在这里把右上同调消解写成
$(L_i)$、$L_i=0$；同一处前文的右消解为 $L^0\to L^1\to\cdots$，
当前校勘本亦改作 $(L^i)$、$L^i=0$。本译采用类型一致的读法，并保留
可逆来源说明；该候选已送交 Canon keeper 去重裁定。}。同样定义长度 $\leq n$ 的
左消解。若一个消解对某个整数 $n$ 的长度 $\leq n$，则称它是\emph{有限}的。

若 $A$ 的一个消解中除 $A$ 外的各个 $\mathcal C$ 对象都是\emph{投射的}
（分别地，\emph{内射的}），则称该消解\emph{投射}（分别地，\emph{内射}）。
当 $\mathcal C$ 是某个环上的（例如左）模范畴时，若消解中除 $A$ 外的各个模
都是\emph{平坦的}（分别地，\emph{自由的}），则同样称 $A$ 的消解\emph{平坦}
（分别地，\emph{自由}）。
\end{env}

\begin{env}[11.4.2]
\label{0.11.4.2-fr}
设 $K^\bullet=(K^i)_{i\in\mathbb Z}$ 是 $\mathcal C$ 的对象组成的复形，其微分
算子的次数为 $+1$。

$K^\bullet$ 的一个\emph{右 Cartan--Eilenberg 消解}，是由一个微分算子次数为
$+1$ 且当 $j<0$ 时有 $L^{ij}=0$ 的双复形
$L^{\bullet\bullet}=(L^{ij})$，以及一个单复形态射
$\varepsilon:K^\bullet\to L^{\bullet,0}$ 组成的偶对，并且满足下列条件：
\oldpage[0\textsubscript{III}]{33}
\begin{enumerate}
  \item[(i)] 对每个指标 $i$，序列
    \[
    \begin{aligned}
      0\longrightarrow{}&K^i\xrightarrow{\varepsilon}L^{i0}
        \longrightarrow L^{i1}\longrightarrow\cdots\\
      0\longrightarrow{}&\mathrm B^i(K^\bullet)\xrightarrow{\varepsilon}
        \mathrm B_{\mathrm I}^i(L^{\bullet,0})
        \longrightarrow\mathrm B_{\mathrm I}^i(L^{\bullet,1})
        \longrightarrow\cdots\\
      0\longrightarrow{}&\mathrm Z^i(K^\bullet)\xrightarrow{\varepsilon}
        \mathrm Z_{\mathrm I}^i(L^{\bullet,0})
        \longrightarrow\mathrm Z_{\mathrm I}^i(L^{\bullet,1})
        \longrightarrow\cdots\\
      0\longrightarrow{}&\mathrm H^i(K^\bullet)\xrightarrow{\varepsilon}
        \mathrm H_{\mathrm I}^i(L^{\bullet,0})
        \longrightarrow\mathrm H_{\mathrm I}^i(L^{\bullet,1})
        \longrightarrow\cdots
    \end{aligned}
    \]
    都正合；换言之，$(L^{i,\bullet})$、
    $(\mathrm B_{\mathrm I}^i(L^{\bullet,\bullet}))$、
    $(\mathrm Z_{\mathrm I}^i(L^{\bullet,\bullet}))$ 和
    $(\mathrm H_{\mathrm I}^i(L^{\bullet,\bullet}))$ 分别是
    $K^i$、$\mathrm B^i(K^\bullet)$、$\mathrm Z^i(K^\bullet)$ 和
    $\mathrm H^i(K^\bullet)$ 的\emph{消解}。
  \item[(ii)] 对每个 $j$，单复形 $L^{\bullet,j}$ 都\emph{分裂}；换言之，正合列
    \[
      \label{0.11.4.2.1-fr}
      0\longrightarrow\mathrm B_{\mathrm I}^i(L^{\bullet,j})
      \longrightarrow\mathrm Z_{\mathrm I}^i(L^{\bullet,j})
      \longrightarrow\mathrm H_{\mathrm I}^i(L^{\bullet,j})
      \longrightarrow0
      \tag{11.4.2.1}
    \]
    \[
      \label{0.11.4.2.2-fr}
      0\longrightarrow\mathrm Z_{\mathrm I}^i(L^{\bullet,j})
      \longrightarrow L^{ij}
      \longrightarrow\mathrm B_{\mathrm I}^{i+1}(L^{\bullet,j})
      \longrightarrow0
      \tag{11.4.2.2}
    \]
    都分裂。
\end{enumerate}

可证明（M, XVII, 1.2）：若 $\mathcal C$ 的每个对象都是某个内射对象的子对象，
则 $\mathcal C$ 的每个复形 $K^\bullet$ 都有一个\emph{内射} Cartan--Eilenberg
消解，也就是说，该消解由内射对象 $L^{ij}$ 组成（这时，上面的条件 (ii) 蕴含
$\mathrm B_{\mathrm I}^i(L^{\bullet,j})$、
$\mathrm Z_{\mathrm I}^i(L^{\bullet,j})$ 和
$\mathrm H_{\mathrm I}^i(L^{\bullet,j})$ 也是内射对象）。此外，对
$\mathcal C$ 的任意复形态射 $f:K^\bullet\to K'^\bullet$、$K^\bullet$ 的任意
Cartan--Eilenberg 消解 $L^{\bullet\bullet}$ 以及 $K'^\bullet$ 的任意内射
Cartan--Eilenberg 消解 $L'^{\bullet\bullet}$，存在与 $f$ 及增广相容的双复形态射
$F:L^{\bullet\bullet}\to L'^{\bullet\bullet}$；而且，若 $f$ 和 $g$ 是
$K^\bullet$ 到 $K'^\bullet$ 的两个同伦态射，则相应的
$L^{\bullet\bullet}$ 到 $L'^{\bullet\bullet}$ 的态射也同伦（\emph{同上引文}）。

当 $K^\bullet$ \emph{下有界}（分别地，\emph{上有界}）时，如果对 $i<i_0$
（分别地，对 $i>i_0$）有 $K^i=0$，则可以取 $L^{\bullet\bullet}$，使得对
$i<i_0$（分别地，对 $i>i_0$）有 $L^{ij}=0$（M, XVII, 1.3）。

另一方面，假设存在整数 $n$，使 $\mathcal C$ 的每个对象都有一个
\emph{长度 $\leq n$ 的内射消解}；那么可以假设对 $j>n$ 有
$L^{ij}=0$（M, XVII, 1.4）。
\end{env}

\begin{env}[11.4.3]
\label{0.11.4.3-fr}
现设 $T$ 是从 $\mathcal C$ 到阿贝尔范畴 $\mathcal C'$ 的一个
\emph{加性协变函子}。给定 $\mathcal C$ 的一个复形 $K^\bullet$ 及
$K^\bullet$ 的一个内射 Cartan--Eilenberg 消解 $L^{\bullet\bullet}$，假设由双复形
$T(L^{\bullet\bullet})$ 定义的（单）复形存在
（参见 \hyperref[0.11.3.1-fr]{11.3.1}）；那么，这个双复形的两个谱序列
$'E(T(L^{\bullet\bullet}))$ 和 $''E(T(L^{\bullet\bullet}))$ 称为
\emph{$T$ 关于复形 $K^\bullet$ 的超上同调谱序列}；由
（\hyperref[0.11.4.2-fr]{11.4.2}）和
（\hyperref[0.11.3.2-fr]{11.3.2}），它们实际上只依赖于 $K^\bullet$，
而不依赖于所选的内射 Cartan--Eilenberg 消解 $L^{\bullet\bullet}$；此外，
它们\emph{函子地}依赖于 $K^\bullet$。它们有相同的收敛项
$\mathrm H^\bullet(T(L^{\bullet\bullet}))$，也称为
\emph{$T$ 关于 $K^\bullet$ 的超上同调}，记作
$\mathrm R^\bullet T(K^\bullet)$。可证明，上述两个谱序列的 $E_2$ 项由下式给出：
\[
  \label{0.11.4.3.1-fr}
  'E_2^{pq}=\mathrm H^p(\mathrm R^qT(K^\bullet))
  \tag{11.4.3.1}
\]
\[
  \label{0.11.4.3.2-fr}
  ''E_2^{pq}=\mathrm R^pT(\mathrm H^q(K^\bullet))
  \tag{11.4.3.2}
\]
\oldpage[0\textsubscript{III}]{34}
其中，对 $p\in\mathbb Z$，$\mathrm R^pT$ 如通常一样表示 $T$ 的第 $p$ 个
\emph{导出函子}；$\mathrm R^qT(K^\bullet)$ 表示复形
$(\mathrm R^qT(K^i))_{i\in\mathbb Z}$。

\phantomsection\label{0.11.4.4-fr}
除非另有明确说明，此后都假设 $\mathcal C$ 的每个对象都是 $\mathcal C$ 的某个
内射对象的子对象，从而 $\mathcal C$ 的任意复形都有内射 Cartan--Eilenberg
消解。由于当 $j<0$ 时有 $L^{ij}=0$，
（\hyperref[0.11.3.3-fr]{11.3.3}）的判别准则表明，在下列两种情形中的每一种，
$T$ 关于 $K^\bullet$ 的\emph{两个}超上同调谱序列都存在并且\emph{双正则}：
\begin{enumerate}
  \item[$1^\circ$] $K^\bullet$ 下有界；
  \item[$2^\circ$] $\mathcal C$ 的每个对象都有一个长度不超过某个整数 $n$
    （该整数与所考虑的对象无关）的内射消解。
\end{enumerate}
事实上，在第一种情形中，可以假设
（\hyperref[0.11.4.2-fr]{11.4.2}）存在 $i_0$，使得当 $i<i_0$ 时有
$L^{ij}=0$；在第二种情形中，可以假设存在 $j_1$，使得当 $j>j_1$ 时有
$L^{ij}=0$。此外，在这两种情形中的每一种，对给定的 $n$，满足
$L^{ij}\neq0$ 且 $i+j=n$ 的偶对 $(i,j)$ 显然只有有限多个，这就证明了上述断言。

若假设 $\mathcal C'$ 中存在滤过余极限且它们正合（特别地，这蕴含
$\mathcal C'$ 中存在无限直和），则由双复形 $T(L^{\bullet\bullet})$ 定义的复形
存在，并且（\hyperref[0.11.3.3-fr]{11.3.3}）的判别准则表明谱序列
$'E(T(L^{\bullet\bullet}))$ 总是\emph{正则}。

若 $K^\bullet$ 是一个除单独一项 $K^{i_0}$ 外所有项 $K^i$ 都为零的复形，
则 $\mathrm R^nT(K^\bullet)$ 同构于 $\mathrm R^{n-i_0}T(K^{i_0})$；这立即由
定义得出，只须取一个满足当 $i\neq i_0$ 时有 $L^{ij}=0$ 的
Cartan--Eilenberg 消解 $L^{\bullet\bullet}$。

若 $K^\bullet$ 和 $K'^\bullet$ 是 $\mathcal C$ 的两个复形，而 $f$、$g$ 是
$K^\bullet$ 到 $K'^\bullet$ 的两个\emph{同伦}态射，则由 $f$ 和 $g$ 导出的态射
$\mathrm R^\bullet T(K^\bullet)\to
\mathrm R^\bullet T(K'^\bullet)$ 相同；上同调谱序列的相应态射也相同。
\end{env}

\begin{proposition}[11.4.5]
\label{0.11.4.5-fr}
假设 $\mathcal C'$ 中存在滤过余极限且它们正合。若对每个 $n>0$ 和每个
$i\in\mathbb Z$ 都有 $\mathrm R^nT(K^i)=0$，则有函子性同构
\[
  \label{0.11.4.5.1-fr}
  \mathrm R^iT(K^\bullet)\xrightarrow{\sim}
  \mathrm H^i(T(K^\bullet))\qquad(i\in\mathbb Z).
  \tag{11.4.5.1}
\]
\end{proposition}

\begin{proof}
事实上，这时第一个谱序列
（\hyperref[0.11.4.3.1-fr]{11.4.3.1}）中仅有的非零 $E_2$ 项是
$'E_2^{p0}=\mathrm H^p(T(K^\bullet))$；换言之，该谱序列\emph{退化}；由于它
正则（\hyperref[0.11.4.4-fr]{11.4.4}），结论由
（\hyperref[0.11.1.6-fr]{11.1.6}）得出。
\end{proof}

\begin{env}[11.4.6]
\label{0.11.4.6-fr}
例如，现在考虑从 $\mathcal C\times\mathcal C'$ 到 $\mathcal C''$ 的一个
\emph{协变双函子} $(M,N)\mapsto T(M,N)$，其中 $\mathcal C$、$\mathcal C'$、
$\mathcal C''$ 是三个阿贝尔范畴；为简便起见，假设 $T$ 对每个变元都可加，
并且 $\mathcal C$ 的每个对象和 $\mathcal C'$ 的每个对象都是某个内射对象的
子对象，而 $\mathcal C''$ 中存在滤过余极限且它们正合。于是，对分别属于
$\mathcal C$ 和 $\mathcal C'$、微分算子次数为 $+1$ 的两个复形
$K^\bullet$、$K'^\bullet$，我们定义\emph{$T$ 关于这两个复形的超上同调}：
用内射 Cartan--Eilenberg 消解 $L^{\bullet\bullet}$（相应地，
$L'^{\bullet\bullet}$）替代 $K^\bullet$（相应地，$K'^\bullet$）。这时
$T(L^{\bullet\bullet},L'^{\bullet\bullet})$ 是 $\mathcal C''$ 中的四重复形；
令 $T(L^{ij},L'^{hk})$ 的两个次数分别为整数 $i+h$ 和 $j+k$，便把它视为
$\mathcal C''$ 中的一个\emph{双复形}。

\emph{$T$ 关于 $K^\bullet$ 和 $K'^\bullet$ 的超上同调}按定义是这个双复形
（换言之，其相伴单复形）的上同调
$\mathrm H^\bullet(T(L^{\bullet\bullet},L'^{\bullet\bullet}))$，记为
$\mathrm R^\bullet T(K^\bullet,K'^\bullet)$；它是两个谱序列的收敛项，其
$E_2$ 项由下式给出：
\oldpage[0\textsubscript{III}]{35}
\[
  \label{0.11.4.6.1-fr}
  'E_2^{pq}=\mathrm H^p(\mathrm R^qT(K^\bullet,K'^\bullet))
  \tag{11.4.6.1}
\]
\[
  \label{0.11.4.6.2-fr}
  ''E_2^{pq}=\sum_{q'+q''=q}
  \mathrm R^pT(\mathrm H^{q'}(K^\bullet),\mathrm H^{q''}(K'^\bullet))
  \tag{11.4.6.2}
\]
（参见 M, XVII, 2）。

这里，$\mathrm R^qT(K^\bullet,K'^\bullet)$ 是双复形
$(\mathrm R^qT(K^i,K'^j))_{(i,j)\in\mathbb Z\times\mathbb Z}$；
（\hyperref[0.11.4.6.1-fr]{11.4.6.1}）右端是把它视为单复形时的上同调。

此外，第一个谱序列总是正则；若存在整数 $n$，使得 $\mathcal C$ 的每个对象和
$\mathcal C'$ 的每个对象都有长度 $\leq n$ 的内射消解，或者若 $K^\bullet$
和 $K'^\bullet$ 下有界，则两个谱序列双正则；在后一种情形中，还可省去
$\mathcal C''$ 中存在余极限这一假设
\SourceCorrectionNote{法文印本此处写作可省去 $\mathcal C$ 与 $\mathcal C'$
中的余极限存在假设，但（11.4.6）唯一作出的此类假设位于目标范畴
$\mathcal C''$；下有界性恰使每个总次数只涉及有限项。本译依 Canon 处置
EGA3-1-CANON-DISP-0001 采用 $\mathcal C''$；法文外交式来源保持不变。}。

若 $K_1^\bullet$、$K_1'^\bullet$ 分别是 $\mathcal C$、$\mathcal C'$ 中的另外
两个复形，$f$、$g$ 是从 $K^\bullet$ 到 $K_1^\bullet$ 的两个同伦态射，
$f'$、$g'$ 是从 $K'^\bullet$ 到 $K_1'^\bullet$ 的两个同伦态射，则一方面
由 $f$ 和 $f'$、另一方面由 $g$ 和 $g'$ 导出的态射
$\mathrm R^\bullet T(K^\bullet,K'^\bullet)\to
\mathrm R^\bullet T(K_1^\bullet,K_1'^\bullet)$ 相同；超上同调谱序列的态射也
相同。

上述构造容易推广到任意可加协变多函子。
\end{env}

\begin{proposition}[11.4.7]
\label{0.11.4.7-fr}
假设对 $\mathcal C$ 的每个内射对象 $I$（相应地，对 $\mathcal C'$ 的每个内射
对象 $I'$），$A'\mapsto T(I,A')$（相应地，$A\mapsto T(A,I')$）是正合函子。
那么，采用（\hyperref[0.11.4.6-fr]{11.4.6}）的记号，有典范同构
\[
  \label{0.11.4.7.1-fr}
  \mathrm R^\bullet T(K^\bullet,K'^\bullet)
  \xrightarrow{\sim}\mathrm H^\bullet(T(L^{\bullet\bullet},K'^\bullet))
  \xrightarrow{\sim}\mathrm H^\bullet(T(K^\bullet,L'^{\bullet\bullet}))
  \tag{11.4.7.1}
\]
其中，后两项分别是由三重复形 $T(L^{\bullet\bullet},K'^\bullet)$ 和
$T(K^\bullet,L'^{\bullet\bullet})$ 定义的单复形的上同调。
\end{proposition}

\begin{proof}
例如来定义其中第一个同构。可把四重复形
$T(L^{\bullet\bullet},L'^{\bullet\bullet})$ 视为一个\emph{双复形}，其中
$T(L^{ij},L'^{hk})$ 的两个次数取为 $i+j+h$ 和 $k$。由于对每个 $h$，
$L'^{h,\bullet}$ 都是 $K'^h$ 的一个\emph{消解}，故根据关于 $T$ 的假设，对这个
双复形有
$\mathrm H_{\mathrm{II}}^q(T(L^{\bullet\bullet},L'^{\bullet\bullet}))=0$
（$q\neq0$），以及
$\mathrm H_{\mathrm{II}}^0(T(L^{\bullet\bullet},L'^{\bullet\bullet}))
=T(L^{\bullet\bullet},K'^\bullet)$。因此，这个双复形的第一个谱序列
\emph{退化}；又因当 $k<0$ 时有 $L'^{hk}=0$，该谱序列还正则
（\hyperref[0.11.3.3-fr]{11.3.3}），故结论由
（\hyperref[0.11.1.6-fr]{11.1.6}）得出。
\end{proof}

对于具有任意个数 $n$ 个变元的协变多函子，也有类似结果：计算超上同调时，
不必用 Cartan--Eilenberg 消解替代\emph{所有}复形，只须替代其中 $n-1$ 个，
条件是：任取 $n-1$ 个变元并固定为\emph{内射}对象后，关于剩余变元的协变函子
是\emph{正合}的。

\subsection{超上同调中取余极限}
\label{subsection:0.11.5-fr}

\begin{lemma}[11.5.1]
\label{0.11.5.1-fr}
设 $K^\bullet=(K^i)_{i\in\mathbb Z}$ 是 $\mathcal C$ 的一个复形；对每个整数
$r\in\mathbb Z$，令 $K_{(r)}^\bullet$ 为满足当 $i<r$ 时
$K_{(r)}^i=0$、当 $i\geq r$ 时 $K_{(r)}^i=K^i$ 的复形。设 $T$ 是从
$\mathcal C$ 到 $\mathcal C'$ 的一个可加协变函子，
\oldpage[0\textsubscript{III}]{36}
它与余极限交换（假设 $\mathcal C$ 和 $\mathcal C'$ 中存在滤过余极限且它们
正合）。那么，当 $r$ 趋于 $-\infty$ 时，$\mathrm R^\bullet T(K^\bullet)$
同构于余极限
\[
  \varinjlim_{r\to-\infty}\mathrm R^\bullet T(K_{(r)}^\bullet)
\]
\end{lemma}

构造 $K^\bullet$ 的一个内射 Cartan--Eilenberg 消解时，对每个 $i$ 任意选取
$\mathrm B^i(K^\bullet)$ 的一个内射消解
$(X_{\mathrm B}^{ij})_{j\geq0}$，以及 $\mathrm H^i(K^\bullet)$ 的一个内射消解
$(X_{\mathrm H}^{ij})_{j\geq0}$。选定这些以后，构造方法表明，$K^i$ 的内射消解
$(L^{ij})_{j\geq0}$ 以及微分算子 $L^{i,j}\to L^{i+1,j}$
\emph{仅依赖于消解} $(X_{\mathrm B}^{i,\bullet})$、
$(X_{\mathrm H}^{i,\bullet})$ 和 $(X_{\mathrm B}^{i+1,\bullet})$
（M, XVII, 1.2）。

显然，对 $i\geq r+1$，有
$\mathrm B^i(K_{(r)}^\bullet)=\mathrm B^i(K^\bullet)$ 和
$\mathrm H^i(K_{(r)}^\bullet)=\mathrm H^i(K^\bullet)$。另一方面，对每个 $r$
都有一个典范单射
$\varphi_{r-1,r}^\bullet:K_{(r)}^\bullet\to K_{(r-1)}^\bullet$，其中当
$i\neq r-1$ 时，$\varphi_{r-1,r}^i$ 是恒等态射。前面的观察表明，若
$L^{\bullet\bullet}=(L^{ij})$ 是 $K^\bullet$ 的一个内射
Cartan--Eilenberg 消解，则对每个 $r$，可以定义 $K_{(r)}^\bullet$ 的一个内射
Cartan--Eilenberg 消解
$L_{(r)}^{\bullet\bullet}=(L_{(r)}^{ij})$，使得当 $i<r$ 时
$L_{(r)}^{ij}=0$，而当 $i\geq r+1$ 时 $L_{(r)}^{ij}=L^{ij}$。此外，可以定义
对应于 $\varphi_{r-1,r}^\bullet$ 的双复形态射
$\Phi_{r-1,r}^{\bullet\bullet}:L_{(r)}^{\bullet\bullet}\to
L_{(r-1)}^{\bullet\bullet}$；该态射的定义方法（\emph{同上引文}）还表明，
可以把它构造成：当 $i\neq r$ 且 $i\neq r-1$ 时，
$\Phi_{r-1,r}^{ij}$ 是恒等态射。

这样就定义了 $\mathcal C$ 中双复形的一个归纳系
$(L_{(r)}^{\bullet\bullet})$；当 $r$ 趋于 $-\infty$ 时，其余极限显然是
$L^{\bullet\bullet}$。由于直和与余极限可交换，$L^{\bullet\bullet}$ 的相伴
单复形也是诸 $L_{(r)}^{\bullet\bullet}$ 的相伴单复形的余极限。根据假设，
$T$ 与余极限交换；上同调也如此（因为函子 $\varinjlim$ 正合），故在同构意义下
确有
\[
  \mathrm H^\bullet(T(L^{\bullet\bullet}))=
  \varinjlim\mathrm H^\bullet(T(L_{(r)}^{\bullet\bullet}))
\]

引理（\hyperref[0.11.5.1-fr]{11.5.1}）使我们能够通过取余极限，把对
\emph{下有界}复形成立的性质推广到任意复形 $K^\bullet$。作为第一个例子，
我们将证明：

\begin{proposition}[11.5.2]
\label{0.11.5.2-fr}
在（\hyperref[0.11.5.1-fr]{11.5.1}）关于 $\mathcal C$、$\mathcal C'$ 和 $T$
的假设下，$\mathrm R^\bullet T(K^\bullet)$ 是 $\mathcal C$ 的复形所成阿贝尔
范畴上的一个上同调函子。
\end{proposition}

\begin{proof}
我们先证明可以归结到下有界复形的情形。若有复形的正合列
\[
  0\longrightarrow K'^\bullet\longrightarrow K^\bullet
  \longrightarrow K''^\bullet\longrightarrow0
\]
则对每个 $r$，显然导出正合列
\[
  0\longrightarrow K_{(r)}'^\bullet\longrightarrow K_{(r)}^\bullet
  \longrightarrow K_{(r)}''^\bullet\longrightarrow0,
\]
从而由假设得到正合列
\[
  \cdots\longrightarrow\mathrm R^nT(K_{(r)}'^\bullet)
  \longrightarrow\mathrm R^nT(K_{(r)}^\bullet)
  \longrightarrow\mathrm R^nT(K_{(r)}''^\bullet)
  \xrightarrow{\partial}\mathrm R^{n+1}T(K_{(r)}'^\bullet)
  \longrightarrow\cdots
\]
这些正合列组成一个归纳系；引理（\hyperref[0.11.5.1-fr]{11.5.1}）和函子
$\varinjlim$ 的正合性证明我们有正合列
\[
  \cdots\longrightarrow\mathrm R^nT(K'^\bullet)
  \longrightarrow\mathrm R^nT(K^\bullet)
  \longrightarrow\mathrm R^nT(K''^\bullet)
  \xrightarrow{\partial}\mathrm R^{n+1}T(K'^\bullet)
  \longrightarrow\cdots
\]
\end{proof}

为处理下有界复形的情形，只须考虑满足当 $i<0$ 时 $K^i=0$ 的那些复形；
它们显然组成一个阿贝尔范畴 $\mathcal K$。

\begin{lemma}[11.5.2.1]
\label{0.11.5.2.1-fr}
在 $\mathcal K$ 中，令 $\mathfrak J$ 为具有下列
\oldpage[0\textsubscript{III}]{37}
性质的复形 $Q^\bullet=(Q^i)_{i\geq0}$ 所成的集合：
\begin{enumerate}
  \item[$1^\circ$] 每个 $Q^i$ 都是 $\mathcal C$ 的内射对象；
  \item[$2^\circ$] 对每个 $i\geq0$，有
    $\mathrm Z^i(Q^\bullet)=\mathrm B^i(Q^\bullet)$，并且
    $\mathrm Z^i(Q^\bullet)$ 是 $Q^i$ 的直和因子。
\end{enumerate}
那么：
\begin{enumerate}
  \item[(i)] 每个 $Q^\bullet\in\mathfrak J$ 都是 $\mathcal K$ 的内射对象。
  \item[(ii)] $\mathcal K$ 的每个对象都同构于某个属于 $\mathfrak J$ 的复形的
    子复形。
\end{enumerate}
\end{lemma}

\begin{proof}
\textup{(i)} 设 $A^\bullet=(A^i)$ 是 $\mathcal K$ 的一个对象，
$A'^\bullet=(A'^i)$ 是 $A^\bullet$ 的一个子对象，
$Q^\bullet=(Q^i)$ 是 $\mathfrak J$ 的一个对象，并且给定了态射
$f=(f^i):A'^\bullet\to Q^\bullet$；需要把它延拓成态射
$g=(g^i):A^\bullet\to Q^\bullet$。为简便起见，我们将使用模范畴的语言
（参见 [27]）。

把 $Q^i$ 与 $\mathrm B^i(Q^\bullet)\oplus\mathrm B^{i+1}(Q^\bullet)$ 认同；
对 $i$ 作归纳，并假设当 $j<i$ 时 $g^j$ 已经定义；当 $j<i-1$ 时，它们与
微分算子 $d^j:A^j\to A^{j+1}$ 和 $d'^j:Q^j\to Q^{j+1}$ 相容，并且还满足：
\begin{enumerate}
  \item[$1^\circ$] $g^{i-1}(\mathrm Z^{i-1}(A^\bullet))\subset
    \mathrm Z^{i-1}(Q^\bullet)$；
  \item[$2^\circ$] 若对每个 $j$ 令 $C^j=(d^j)^{-1}(A'^{j+1})$，则在
    $C^{i-1}$ 上，$d'^{i-1}\circ g^{i-1}$ 与
    $f^i\circ d^{i-1}$ 相同。
\end{enumerate}
把态射 $f^i:A'^i\to Q^i$ 与两个投影复合，得到两个态射
$f'^i:A'^i\to\mathrm B^i(Q^\bullet)$ 和
$f''^i:A'^i\to\mathrm B^{i+1}(Q^\bullet)$。由于
$d'^{i-1}\circ g^{i-1}$ 把 $A^{i-1}$ 映入
$\mathrm B^i(Q^\bullet)$，并且在 $\mathrm Z^{i-1}(A^\bullet)$ 上为零，
它定义了一个态射
$h^i:\mathrm B^i(A^\bullet)\to\mathrm B^i(Q^\bullet)$；又由于在
$C^{i-1}$ 上 $d'^{i-1}\circ g^{i-1}$ 与
$f^i\circ d^{i-1}$ 相同，故 $h^i$ 在
$\mathrm B^i(A^\bullet)\cap A'^i$ 上与 $f'^i$ 相同。因为
$\mathrm B^i(Q^\bullet)$ 是 $Q^i$ 的直和因子，所以它是内射的；于是存在态射
$g'^i:A^i\to\mathrm B^i(Q^\bullet)$，它在
$\mathrm B^i(A^\bullet)$ 上与 $h^i$ 相同，在 $A'^i$ 上与 $f'^i$ 相同。
另一方面，考虑态射
$f'^{i+1}\circ d^i:C^i\to\mathrm B^{i+1}(Q^\bullet)$
\SourceCorrectionNote{法文印本在本句及下文各写一次 $f''^{i+1}\circ d^i$；
但 $f''^{i+1}$ 的值落在 $\mathrm B^{i+2}(Q^\bullet)$，只有
$f'^{i+1}$ 的值落在所需的 $\mathrm B^{i+1}(Q^\bullet)$。本译依 Canon 处置
EGA3-1-CANON-DISP-0003 在两处均采用单撇映射；法文外交式来源保持不变。}，它在
$\mathrm Z^i(A^\bullet)$ 上为零；由于 $\mathrm B^{i+1}(Q^\bullet)$ 是内射的，
存在态射 $g''^i:A^i\to\mathrm B^{i+1}(Q^\bullet)$，它在 $C^i$ 上与
$f'^{i+1}\circ d^i$ 相同，而在 $\mathrm Z^i(A^\bullet)$ 上与 $0$ 相同。
这时取 $g^i=g'^i+g''^i$，即可继续归纳。

\textup{(ii)} 为把 $A^\bullet=(A^i)$ 嵌入一个属于 $\mathfrak J$ 的复形，对每个
$i\geq1$，取 $\mathcal C$ 的一个内射对象 $Q'^i$，使得存在单射
$f'^i:A^i\to Q'^i$。于是令
\[
  Q^i=0\quad(i<0),\qquad Q^0=Q'^1,\qquad
  Q^i=Q'^i\oplus Q'^{i+1}\quad(i\geq1),
\]
并配以显然的微分算子。对每个 $i$，令
$f^i=f'^i+(f'^{i+1}\circ d^i)$（当 $i\leq0$ 时令 $f'^i=0$）；立即可见，
对每个 $i$，$f^i$ 都是单射，并且诸 $f^i$ 与微分算子相容。
\end{proof}

\begin{corollary}[11.5.2.2]
\label{0.11.5.2.2-fr}
$\mathcal K$ 的每个对象 $K^\bullet$ 都有一个由 $\mathfrak J$ 的对象组成的右消解。
若 $L^{\bullet\bullet}$ 是这样的消解，则对 $K^\bullet$ 的任意一个由
$\mathcal K$ 的对象组成的消解 $L'^{\bullet\bullet}$，存在与增广相容的双复形
态射 $F:L'^{\bullet\bullet}\to L^{\bullet\bullet}$，而任意两个这样的态射
$F$、$F'$ 都同伦。
\end{corollary}

\begin{proof}
这正是把（M, V, 1.1 a)）应用于阿贝尔范畴 $\mathcal K$。
\end{proof}

\begin{env}[11.5.2.3]
\label{0.11.5.2.3-fr}
有了这些预备结果，考虑 $K^\bullet$ 的一个内射 Cartan--Eilenberg 消解
$L'^{\bullet\bullet}$，以及 $K^\bullet$ 的一个由 $\mathfrak J$ 的对象组成的消解
$L^{\bullet\bullet}$；我们证明有同构
\[
  \mathrm H^\bullet(T(L'^{\bullet\bullet}))\xrightarrow{\sim}
  \mathrm H^\bullet(T(L^{\bullet\bullet})).
\]
事实上，由（\hyperref[0.11.5.2.2-fr]{11.5.2.2}）导出双复形态射
$T(L'^{\bullet\bullet})\to T(L^{\bullet\bullet})$，因而导出这些双复形的第一个
谱序列之间的态射
$'E(T(L'^{\bullet\bullet}))\to{}'E(T(L^{\bullet\bullet}))$。由
（\hyperref[0.11.3.3-fr]{11.3.3}），这些谱序列正则，故根据
（\hyperref[0.11.1.5-fr]{11.1.5}），只须证明上述态射在 $E_2$ 项上是同构；
换言之，只须证明
$\mathrm H_{\mathrm{II}}^q(T(L^{i,\bullet}))$ 等于 $\mathrm R^qT(K^i)$。
由于 $L^{i,\bullet}$ 是 $K^i$ 的右消解，问题归结为证明下列引理。
\end{env}

\oldpage[0\textsubscript{III}]{38}
\begin{lemma}[11.5.2.4]
\label{0.11.5.2.4-fr}
若 $L^\bullet=(L^i)_{i\in\mathbb Z}$ 是 $\mathcal C$ 的一个对象 $A$ 的右消解，
并且对每个 $i\in\mathbb Z$ 和每个 $n>0$ 都有
$\mathrm R^nT(L^i)=0$，那么
$\mathrm H^\bullet T(L^\bullet)=\mathrm R^\bullet T(A)$。
\end{lemma}

\begin{proof}
这是（T, 2.5.2）的一个特殊情形。
\end{proof}

\begin{env}[11.5.2.5]
\label{0.11.5.2.5-fr}
现在，（\hyperref[0.11.5.2-fr]{11.5.2}）的证明立即完成，因为
$L'^{\bullet\bullet}$ 是阿贝尔范畴 $\mathcal K$ 中 $K^\bullet$ 的一个内射消解；
换言之，$K^\bullet\mapsto\mathrm H^\bullet(T(L'^{\bullet\bullet}))$ 正是范畴
$\mathcal K$ 中 $T$ 的右导出上同调函子（T, 2.3）。
\end{env}

\begin{proposition}[11.5.3]
\label{0.11.5.3-fr}
在（\hyperref[0.11.5.1-fr]{11.5.1}）关于 $\mathcal C$、$\mathcal C'$ 和 $T$
的假设下，设 $L^{\bullet\bullet}=(L^{ij})$ 是满足当 $j<0$ 时
$L^{ij}=0$ 的一个双复形，并且对每个 $i$，$L^{i,\bullet}$ 是 $K^i$ 的一个消解；
最后，假设对每个偶对 $(i,j)$ 和每个 $n>0$，都有
$\mathrm R^nT(L^{ij})=0$。那么存在函子性同构
\[
  \label{0.11.5.3.1-fr}
  \mathrm R^\bullet T(K^\bullet)\xrightarrow{\sim}
  \mathrm H^\bullet(T(L^{\bullet\bullet})).
  \tag{11.5.3.1}
\]
\end{proposition}

\begin{proof}
令 $L_{(r)}^{\bullet\bullet}=(L_{(r)}^{ij})$ 为满足当 $i<r$ 时
$L_{(r)}^{ij}=0$、当 $i\geq r$ 时 $L_{(r)}^{ij}=L^{ij}$ 的双复形；立即可见，
当 $r$ 趋于 $-\infty$ 时，$L^{\bullet\bullet}$ 是诸
$L_{(r)}^{\bullet\bullet}$ 的余极限。因此，根据关于 $T$ 的假设和
（\hyperref[0.11.5.1-fr]{11.5.1}），只须在 $K^\bullet$ 下有界的情形下证明
命题，例如，假设当 $i<0$ 时 $K^i=0$，并且当 $i<0$ 时 $L^{ij}=0$。这时，
设 $L'^{\bullet\bullet}=(L'^{ij})$ 是由 $\mathfrak J$ 的对象组成的
$K^\bullet$ 的一个右消解（\hyperref[0.11.5.2.2-fr]{11.5.2.2}）；存在与增广
相容的双复形态射 $L^{\bullet\bullet}\to L'^{\bullet\bullet}$，从而导出第一个
谱序列之间的态射
$'E(T(L^{\bullet\bullet}))\to{}'E(T(L'^{\bullet\bullet}))$。引理
（\hyperref[0.11.5.2.4-fr]{11.5.2.4}）像在
（\hyperref[0.11.5.2.3-fr]{11.5.2.3}）中一样表明，这个态射是
\emph{同构}，结论随之得出。
\end{proof}

\begin{remark}[11.5.4]
\label{0.11.5.4-fr}
前面的论证证明，在下有界复形 $K^\bullet$ 所成的范畴中，
（\hyperref[0.11.5.2-fr]{11.5.2}）和
（\hyperref[0.11.5.3-fr]{11.5.3}）的结论成立，\emph{无须假设 $T$ 与滤过
余极限交换}。此外，当只考虑满足当 $i<0$ 时 $K^i=0$ 的复形 $K^\bullet$
所成的范畴 $\mathcal K$ 时，把 $\mathrm R^\bullet T(K^\bullet)$ 刻画为
$\mathcal K$ 中 $T$ 的\emph{右导出函子}系，表明这个上同调函子是
\emph{泛的}（T, 2.3）。

（\hyperref[0.11.5.2-fr]{11.5.2}）在不对 $T$ 作附加假设时仍成立的另一种
情形是：存在整数 $m>0$，使得 $\mathcal C$ 的每个对象都有长度 $\leq m$ 的
内射消解。事实上，在（\hyperref[0.11.5.1-fr]{11.5.1}）的证明中，所出现的
$\mathcal C$ 的对象的所有内射消解都可以取为长度 $\leq m$；由此立即得到：
一旦 $-r$ 足够大，双复形 $T(L_{(r)}^{\bullet\bullet})$ 的次数总和为 $n$ 的诸项
便与 $T(L^{\bullet\bullet})$ 的相应项相等，而且数目有限；因而对每个 $n$，
一旦 $-r$ 足够大，就有
$\mathrm H^n(T(L^{\bullet\bullet}))=
\mathrm H^n(T(L_{(r)}^{\bullet\bullet}))$。采用
（\hyperref[0.11.5.2-fr]{11.5.2}）的记号，对足够大的 $-r$（依赖于 $n$）也有
$\mathrm R^nT(K_{(r)}^\bullet)=\mathrm R^nT(K^\bullet)$；对于
$K'^\bullet$ 和 $K''^\bullet$ 也一样，结论随之得出。同理，当 $\mathcal C$
满足前述假设，并假设消解 $L^{i,\bullet}$ 的长度 $\leq m$ 时，
（\hyperref[0.11.5.3-fr]{11.5.3}）无需关于 $T$ 的附加条件也成立
\SourceCorrectionNote{法文印本在本段三处写作“$r$ 足够大”；但
$L_{(r)}^{i\bullet}$ 在 $i<r$ 时为零，并且（11.5.1）取的是 $r\to-\infty$
的余极限，故稳定条件应为 $-r$ 足够大。本译依 Canon 处置
EGA3-1-CANON-DISP-0004 校正三处方向；法文外交式来源保持不变。}。
\end{remark}

\begin{env}[11.5.5]
\label{0.11.5.5-fr}
（\hyperref[0.11.5.2-fr]{11.5.2}）的结果可推广到协变多函子。例如，考虑
（\hyperref[0.11.4.6-fr]{11.4.6}）中的情形；在那里假设 $\mathcal C$、
$\mathcal C'$ 和 $\mathcal C''$ 中
\oldpage[0\textsubscript{III}]{39}
存在滤过余极限且它们正合，并且 $T$ 与余极限交换。那么
$\mathrm R^\bullet T(K^\bullet,K'^\bullet)$ 是分别关于复形 $K^\bullet$、
$K'^\bullet$ 都为上同调函子的双函子。为说明这一点，像在
（\hyperref[0.11.5.2-fr]{11.5.2}）中一样，归结到 $K^\bullet$ 和
$K'^\bullet$ 都下有界的情形；然后为 $K^\bullet$ 和 $K'^\bullet$ 取
（\hyperref[0.11.5.2.2-fr]{11.5.2.2}）所述类型的内射消解，便归结到
（M, V, 4.1）中证明的一般性质。
\end{env}

\begin{env}[11.5.6]
\label{0.11.5.6-fr}
同样，（\hyperref[0.11.4.7-fr]{11.4.7}）和
（\hyperref[0.11.5.3-fr]{11.5.3}）的结果可推广如下。假设（在
（\hyperref[0.11.5.5-fr]{11.5.5}）的假设下）有两个双复形
$L^{\bullet\bullet}=(L^{ij})$、$L'^{\bullet\bullet}=(L'^{ij})$，当 $j<0$ 时
$L^{ij}=0$ 且 $L'^{ij}=0$；对每个 $i$，$L^{i,\bullet}$ 是 $K^i$ 的消解，
$L'^{i,\bullet}$ 是 $K'^i$ 的消解；最后，对 $n>0$ 及任意指标系
$(i,j,h,k)$，有 $\mathrm R^nT(L^{ij},L'^{hk})=0$。那么有一个关于
$K^\bullet$ 和 $K'^\bullet$ 函子性的同构
\[
  \label{0.11.5.6.1-fr}
  \mathrm R^\bullet T(K^\bullet,K'^\bullet)\xrightarrow{\sim}
  \mathrm H^\bullet(T(L^{\bullet\bullet},L'^{\bullet\bullet})).
  \tag{11.5.6.1}
\]
像在（\hyperref[0.11.5.3-fr]{11.5.3}）中一样，归结到 $K^\bullet$ 和
$K'^\bullet$ 都下有界的情形，即可证明此式。

进一步假设，对每个偶对 $(i,j)$ 和每个偶对 $(h,k)$，函子
$A\mapsto T(A,L'^{hk})$ 和 $A'\mapsto T(L^{ij},A')$ 分别在
$\mathcal C$ 和 $\mathcal C'$ 中正合。那么还有函子性同构
\[
  \label{0.11.5.6.2-fr}
  \mathrm R^\bullet T(K^\bullet,K'^\bullet)
  \xrightarrow{\sim}\mathrm H^\bullet(T(L^{\bullet\bullet},K'^\bullet))
  \xrightarrow{\sim}\mathrm H^\bullet(T(K^\bullet,L'^{\bullet\bullet})).
  \tag{11.5.6.2}
\]
证明与（\hyperref[0.11.4.7-fr]{11.4.7}）的证明类似。
\end{env}

\begin{env}[11.5.7]
\label{0.11.5.7-fr}
还应注意，（\hyperref[0.11.5.5-fr]{11.5.5}）和
（\hyperref[0.11.5.6-fr]{11.5.6}）的结果无需假设 $T$ 与余极限交换也成立，
只要限于下有界复形 $K^\bullet$、$K'^\bullet$，或者假设
$\mathcal C$（相应地，$\mathcal C'$）的每个对象都有长度有界的内射消解，
并且在（\hyperref[0.11.5.6-fr]{11.5.6}）中双复形
$L^{\bullet\bullet}$ 和 $L'^{\bullet\bullet}$ 的第二个次数上有界。
\end{env}

\subsection{函子关于复形 $K_\bullet$ 的超同调}
\label{subsection:0.11.6-fr}

\begin{env}[11.6.1]
\label{0.11.6.1-fr}
设 $\mathcal C$ 是阿贝尔范畴，
$K_\bullet=(K_i)_{i\in\mathbb Z}$ 是由 $\mathcal C$ 的对象组成的复形，其微分算子
的次数为 $-1$。$K_\bullet$ 的一个\emph{左 Cartan--Eilenberg 消解}由一个微分
算子次数为 $-1$ 的双复形 $L_{\bullet\bullet}=(L_{ij})$ 和一个单复形态射
$\varepsilon:L_{\bullet,0}\to K_\bullet$ 组成，其中当 $j<0$ 时 $L_{ij}=0$，
并且满足由（\hyperref[0.11.4.2-fr]{11.4.2}）中的条件“反转箭头”所得的条件。
若 $\mathcal C$ 的每个对象都是某个投射对象的商，则 $\mathcal C$ 的每个复形
$K_\bullet$ 都有一个\emph{投射的} Cartan--Eilenberg 消解，即由投射对象
$L_{ij}$ 组成的消解；它具有与（\hyperref[0.11.4.2-fr]{11.4.2}）中类似的
函子性。此外，若 $K_\bullet$ 下有界（相应地，上有界），并且当 $i<i_0$
（相应地，$i>i_0$）时 $K_i=0$，则可假设当 $i<i_0$（相应地，$i>i_0$）时
$L_{ij}=0$。若 $\mathcal C$ 的每个对象都有长度 $\leq n$ 的投射消解，则可假设
当 $j>n$ 时 $L_{ij}=0$。
\end{env}

\begin{env}[11.6.2]
\label{0.11.6.2-fr}
\oldpage[0\textsubscript{III}]{40}
假设 $T$ 是从 $\mathcal C$ 到一个阿贝尔范畴 $\mathcal C'$ 的协变加性函子。
$T$ 关于 $\mathcal C$ 的复形 $K_\bullet$ 的\emph{超同调}
$L_\bullet T(K_\bullet)$ 和\emph{超同调谱序列}（当它们存在时）的定义，仍由
（\hyperref[0.11.4.3-fr]{11.4.3}）“反转箭头”得到；由此所得两个谱序列的
$E_2$ 项为
\[
\label{0.11.6.2.1-fr}
\tag{11.6.2.1}
{}'E^2_{pq}=H_p(L_qT(K_\bullet))
\]
\[
\label{0.11.6.2.2-fr}
\tag{11.6.2.2}
{}''E^2_{pq}=L_pT(H_q(K_\bullet)),
\]
其中，对 $p\geq 0$，$L_pT$ 表示 $T$ 的第 $p$ 个左导出函子；对 $p<0$，
则为 $0$；而 $L_qT(K_\bullet)$ 表示复形
$(L_qT(K_i))_{i\in\mathbb Z}$。

超同调的性质并非全都能由超上同调的性质简单地“反转箭头”得到（除非对范畴
$\mathcal C'$ 附加（T, 1.5）中（AB~5*）型的假设），其原因在于前述两个谱序列
的正则性条件；这一次必须对它们应用（\hyperref[0.11.3.4-fr]{11.3.4}）的判别
准则。这些准则表明，若假设 $\mathcal C'$ 中存在滤过余极限且它们正合，则由
双复形 $T(L_{\bullet\bullet})$ 确定的复形存在，而第二个谱序列
$\,{}''E(T(L_{\bullet\bullet}))$ 此时正则。若假设 $K_\bullet$ 下有界，或者存在
整数 $n$，使得 $\mathcal C$ 的每个对象都有长度 $\leq n$ 的投射消解，则两个
超同调谱序列都存在（无需对 $\mathcal C'$ 作假设），并且双正则。
\end{env}

\begin{proposition}[11.6.3]
\label{0.11.6.3-fr}
设 $\mathcal C$、$\mathcal C'$ 是两个阿贝尔范畴，$T$ 是从 $\mathcal C$ 到
$\mathcal C'$ 的协变加性函子。那么：
\begin{enumerate}
\item[{\rm(i)}] \emph{超同调 $L_\bullet T(K_\bullet)$ 是 $\mathcal C$ 的下有界
复形所成阿贝尔范畴中的同调函子。}
\item[{\rm(ii)}] \emph{设 $K_\bullet$ 是 $\mathcal C$ 的下有界复形。若对每个
$n>0$ 和每个 $i\in\mathbb Z$ 都有 $L_nT(K_i)=0$，则存在函子性同构}
\[
\label{0.11.6.3.1-fr}
\tag{11.6.3.1}
L_iT(K_\bullet)\simeq H_i(T(K_\bullet))\qquad(i\in\mathbb Z).
\]
\item[{\rm(iii)}] \emph{设 $K_\bullet$ 是 $\mathcal C$ 的下有界复形。设
$L_{\bullet\bullet}=(L_{ij})$ 是一个双复形，当 $j<0$ 时 $L_{ij}=0$，并且对每个
$i$，$L_{i,\bullet}$ 是 $K_i$ 的一个消解；最后假设对每个偶对 $(i,j)$ 和每个
$n>0$ 都有 $L_nT(L_{ij})=0$。那么存在函子性同构}
\[
\label{0.11.6.3.2-fr}
\tag{11.6.3.2}
L_\bullet T(K_\bullet)\simeq H_\bullet(T(L_{\bullet\bullet})).
\]
\end{enumerate}
这些证明分别与下有界复形情形下（\hyperref[0.11.5.2-fr]{11.5.2}）、
（\hyperref[0.11.4.5-fr]{11.4.5}）和
（\hyperref[0.11.5.3-fr]{11.5.3}）的证明相同。我们把这些论证的细节留给读者。
\end{proposition}

\begin{env}[11.6.4]
\label{0.11.6.4-fr}
对于关于每个变元都为加性的协变多函子，有完全类似的结果。例如，对一个双函子
$T$，有两个超同调谱序列，其 $E_2$ 项为
\oldpage[0\textsubscript{III}]{41}
\[
\label{0.11.6.4.1-fr}
\tag{11.6.4.1}
{}'E^2_{pq}=H_p(L_qT(K_\bullet,K'_\bullet))
\]
\[
\label{0.11.6.4.2-fr}
\tag{11.6.4.2}
{}''E^2_{pq}=
\sum_{q'+q''=q}L_pT(H_{q'}(K_\bullet),H_{q''}(K'_\bullet)).
\]
在这里仍是第二个谱序列正则；当 $K_\bullet$、$K'_\bullet$ 是下有界复形，
或者所考虑的阿贝尔范畴的对象都有某个固定长度的投射消解时，两个谱序列双正则。

此外：
\end{env}

\begin{proposition}[11.6.5]
\label{0.11.6.5-fr}
设 $\mathcal C$、$\mathcal C'$、$\mathcal C''$ 是三个阿贝尔范畴，$T$ 是从
$\mathcal C\times\mathcal C'$ 到 $\mathcal C''$ 的协变双加性双函子。
\begin{enumerate}
\item[{\rm(i)}] \emph{$L_\bullet T(K_\bullet,K'_\bullet)$ 分别关于下有界复形
$K_\bullet$、$K'_\bullet$ 是同调双函子（它们分别由 $\mathcal C$ 的对象和
$\mathcal C'$ 的对象组成）。}
\item[{\rm(ii)}] \emph{假设 $K_\bullet$ 和 $K'_\bullet$ 下有界。设
$L_{\bullet\bullet}=(L_{ij})$、$L'_{\bullet\bullet}=(L'_{ij})$ 是两个双复形，
当 $j<0$ 时 $L_{ij}=0$ 且 $L'_{ij}=0$；对每个 $i$，$L_{i,\bullet}$ 是 $K_i$
的一个消解，$L'_{i,\bullet}$ 是 $K'_i$ 的一个消解；最后，对 $n>0$ 和每个指标
系 $(i,j,h,k)$，都有 $L_nT(L_{ij},L'_{hk})=0$。那么存在函子性同构}
\[
\label{0.11.6.5.1-fr}
\tag{11.6.5.1}
L_\bullet T(K_\bullet,K'_\bullet)
\simeq H_\bullet(T(L_{\bullet\bullet},L'_{\bullet\bullet})).
\]
\item[{\rm(iii)}] \emph{进一步假设，对每个偶对 $(i,j)$ 和每个偶对 $(h,k)$，
函子
\[
A\longmapsto T(A,L'_{hk}),\qquad
A'\longmapsto T(L_{ij},A')
\]
分别在 $\mathcal C$ 和 $\mathcal C'$ 中正合。那么存在函子性同构}
\[
\label{0.11.6.5.2-fr}
\tag{11.6.5.2}
L_\bullet T(K_\bullet,K'_\bullet)
\simeq H_\bullet(T(L_{\bullet\bullet},K'_\bullet))
\simeq H_\bullet(T(K_\bullet,L'_{\bullet\bullet})).
\]
\end{enumerate}
这些证明与（\hyperref[0.11.5.5-fr]{11.5.5}）和
（\hyperref[0.11.5.6-fr]{11.5.6}）的证明类似。
\end{proposition}

\subsection{函子关于双复形 $K_{\bullet\bullet}$ 的超同调}
\label{subsection:0.11.7-fr}

\begin{env}[11.7.1]
\label{0.11.7.1-fr}
设 $\mathcal C$ 是一个阿贝尔范畴，其中每个对象都是某个投射对象的商。考虑一个
由 $\mathcal C$ 的对象组成且两个次数都下有界的双复形
$K_{\bullet\bullet}=(K_{ij})$；总可以归结到当 $i<0$ 或 $j<0$ 时
$K_{ij}=0$ 的情形，以下即作此约定。可以把 $K_{\bullet\bullet}$ 看成一个
（单）复形，它由阿贝尔范畴 $\mathcal K$ 的对象
$K_{i,\bullet}=(K_{ij})_{j\geq 0}$ 组成；该范畴是由 $\mathcal C$
的对象组成的非负次数复形所成的范畴。由引理
（\hyperref[0.11.5.2.1-fr]{11.5.2.1}）（更确切地说，由“反转箭头”所得的
“对偶”引理）和（M, V, 2.2）可知，$K_{\bullet\bullet}$ 在范畴 $\mathcal K$
中有一个\emph{投射 Cartan--Eilenberg 消解}；这样的消解是 $\mathcal C$ 中
各次数都 $\geq 0$、且由投射对象组成的一个三复形
$M_{\bullet\bullet\bullet}=(M_{ijk})$，使得对每个 $i$，
$M_{i,\bullet\bullet}$、$\mathrm B_i^{\mathrm I}(M_{\bullet\bullet\bullet})$、
$\mathrm Z_i^{\mathrm I}(M_{\bullet\bullet\bullet})$、
$H_i^{\mathrm I}(M_{\bullet\bullet\bullet})$ 分别是 $K_{i,\bullet}$、
$\mathrm B_i^{\mathrm I}(K_{\bullet\bullet})$、
$\mathrm Z_i^{\mathrm I}(K_{\bullet\bullet})$、
$H_i^{\mathrm I}(K_{\bullet\bullet})$ 在范畴 $\mathcal K$ 中的投射消解；
特别地，对每个偶对 $(i,j)$，$M_{ij,\bullet}$ 是 $K_{ij}$ 在 $\mathcal C$
中的投射消解。
\end{env}

\begin{proposition}[11.7.2]
\label{0.11.7.2-fr}
设 $T$ 是从 $\mathcal C$ 到一个阿贝尔范畴 $\mathcal C'$ 的协变加性函子。
沿用（\hyperref[0.11.7.1-fr]{11.7.1}）的记号，\emph{与三复形
$T(M_{\bullet\bullet\bullet})$ 相关联的单复形的同调
$H_\bullet(T(M_{\bullet\bullet\bullet}))$，与
$K_{\bullet\bullet}$ 所关联单复形的超同调
$L_\bullet T(K_{\bullet\bullet})$
\oldpage[0\textsubscript{III}]{42}
（\hyperref[0.11.6.2-fr]{11.6.2}）典范同构；它还是六个双正则谱序列的共同
收敛项。这些谱序列记为 $\mathop{}^{(t)}E$（其中 $t=a,b,a',b',c$ 或 $d$），
其 $E_2$ 项由下列公式给出：}
\[
\begin{aligned}
{}^{(a)}E^2_{pq}&=L_pT(H_q(K_{\bullet\bullet})),&
{}^{(b)}E^2_{pq}&=H_p(L_q^{\mathrm{II}}T(K_{\bullet\bullet})),\\
{}^{(a')}E^2_{pq}&=L_pT(H_q^{\mathrm I}(K_{\bullet\bullet})),&
{}^{(b')}E^2_{pq}&=H_p(L_qT(K_{\bullet\bullet})),\\
{}^{(c)}E^2_{pq}&=L_pT(H_q^{\mathrm{II}}(K_{\bullet\bullet})),&
{}^{(d)}E^2_{pq}&=H_p(L_q^{\mathrm I}T(K_{\bullet\bullet})).
\end{aligned}
\]
\end{proposition}

回顾一下，对任意复形 $A_\bullet=(A_i)$，我们用记号 $F(A_\bullet)$ 表示由对象
$F(A_i)$ 组成的复形；例如，$L_q^{\mathrm{II}}T(K_{\bullet\bullet})$ 表示复形
$(L_q^{\mathrm{II}}T(K_{i,\bullet}))_{i\geq0}$，其中
$L_q^{\mathrm{II}}T(K_{i,\bullet})$ 是 $T$ 关于单复形 $K_{i,\bullet}$ 的指标为
$q$ 的超同调。

\begin{proof}
以 $L_\bullet$ 表示与 $K_{\bullet\bullet}$ 相关联的单复形，于是
\[
L_i=\bigoplus_{r+s=i}K_{rs},
\]
并令
\[
N_{ij}=\bigoplus_{r+s=i}M_{rsj};
\]
显然，对每个 $i$，$N_{i,\bullet}$ 是 $L_i$ 在 $\mathcal C$ 中的一个投射消解；
因而由（\hyperref[0.11.6.3-fr]{11.6.3}）和
（\hyperref[0.11.6.4-fr]{11.6.4}）得到函子性同构
\[
L_\bullet T(L_\bullet)\simeq H_\bullet(T(N_{\bullet\bullet}));
\]
由于与双复形 $T(N_{\bullet\bullet})$ 相关联的单复形也就是与三复形
$T(M_{\bullet\bullet\bullet})$ 相关联的单复形，这就证明了命题的第一个断言。

此外，$L_\bullet T(L_\bullet)$ 是 $T$ 关于单复形 $L_\bullet$ 的两个超同调
谱序列（\hyperref[0.11.6.2-fr]{11.6.2}）的收敛项；这两个谱序列分别就是
$\mathop{}^{(b')}E$ 和 $\mathop{}^{(a)}E$。

现在把 $M_{\bullet\bullet\bullet}$ 看成双复形 $U_{\bullet\bullet}$，其中
\[
U_{ij}=\bigoplus_{r+s=j}M_{irs};
\]
$H_\bullet(T(M_{\bullet\bullet\bullet}))$ 仍是双复形
$T(U_{\bullet\bullet})$ 的两个谱序列的收敛项。对每个 $i$，
$M_{i,\bullet\bullet}$ 都是一个相对于单复形 $K_{i,\bullet}$ 满足
（\hyperref[0.11.6.3-fr]{11.6.3}）条件的双复形；因此
$H_q^{\mathrm{II}}(T(U_{i,\bullet}))=
L_q^{\mathrm{II}}T(K_{i,\bullet})$，而 $T(U_{\bullet\bullet})$ 的第一个谱序列
正是 $\mathop{}^{(b)}E$。另一方面，对每个 $r$，
$M_{\bullet,r,\bullet}$ 是单复形 $K_{\bullet,r}$ 的一个 Cartan--Eilenberg
消解；（M, XV, 2）中的计算表明
\[
H_q^{\mathrm I}(T(M_{\bullet,rs}))
=T(H_q^{\mathrm I}(M_{\bullet,rs})),
\]
从而
\[
H_q^{\mathrm I}(T(U_{\bullet,j}))
=\bigoplus_{r+s=j}T(H_q^{\mathrm I}(M_{\bullet,rs}));
\]
换言之，单复形 $H_q^{\mathrm I}(T(U_{\bullet\bullet}))$ 就是与双复形
$T(H_q^{\mathrm I}(M_{\bullet\bullet\bullet}))$ 相关联的单复形。对每个 $q$，
$H_q^{\mathrm I}(M_{\bullet\bullet\bullet})$ 是单复形
$H_q^{\mathrm I}(K_{\bullet\bullet})$ 在范畴 $\mathcal K$ 中的一个投射消解；
应用（\hyperref[0.11.6.3-fr]{11.6.3}），可见
\[
H_p^{\mathrm{II}}
\bigl(T(H_q^{\mathrm I}(M_{\bullet\bullet\bullet}))\bigr)
=L_pT(H_q^{\mathrm I}(K_{\bullet\bullet})),
\]
这就得到谱序列 $\mathop{}^{(a')}E$。最后，在三复形
$M_{\bullet\bullet\bullet}$ 的定义中交换指标 $i$ 与 $j$ 的角色，并对这个新
三复形应用以上论证，便得到谱序列 $\mathop{}^{(c)}E$ 和
$\mathop{}^{(d)}E$。
\end{proof}

称 $L_\bullet T(K_{\bullet\bullet})$ 为 $T$ 关于双复形
$K_{\bullet\bullet}$ 的\emph{超同调}。

\begin{env}[11.7.3]
\label{0.11.7.3-fr}
\emph{注记}。
\begin{enumerate}
\item[{\rm(i)}] 由（\hyperref[0.11.6.3-fr]{11.6.3}）可知，
$L_\bullet T(K_{\bullet\bullet})$ 是 $\mathcal C$ 的两个次数都下有界的双复形
$K_{\bullet\bullet}$ 所成范畴中的同调函子。
\oldpage[0\textsubscript{III}]{43}
\item[{\rm(ii)}] 设 $M_{\bullet\bullet\bullet}$ 是 $\mathcal C$ 中的一个三复形，
使得对每个偶对 $(r,s)$，$M_{rs,\bullet}$ 是 $K_{rs}$ 的一个消解，并且对每个
三元组 $(i,j,k)$ 和每个 $n>0$，都有 $L_nT(M_{ijk})=0$。那么有同构
\[
L_\bullet T(K_{\bullet\bullet})
\simeq H_\bullet(T(M_{\bullet\bullet\bullet}));
\]
事实上，沿用（\hyperref[0.11.7.2-fr]{11.7.2}）证明中的记号，
$N_{i,\bullet}$ 是 $L_i$ 的一个消解，并且对 $n>0$ 和每个偶对 $(i,j)$，
$L_nT(N_{ij})=0$；只须应用（\hyperref[0.11.6.3-fr]{11.6.3}, (iii)）。
\item[{\rm(iii)}] （\hyperref[0.11.7.2-fr]{11.7.2}）的结果立即推广到协变
多函子。例如，设 $\mathcal C'$ 是另一个每个对象都是某个投射对象之商的
阿贝尔范畴，$K'_{\bullet\bullet}$ 是 $\mathcal C'$ 中两个次数都下有界的双复形，
而 $T$ 是从 $\mathcal C\times\mathcal C'$ 到一个阿贝尔范畴 $\mathcal C''$
的协变加性双函子。若 $L_\bullet$ 和 $L'_\bullet$ 分别是与
$K_{\bullet\bullet}$ 和 $K'_{\bullet\bullet}$ 相关联的单复形，则把 $T$ 关于
两个双复形 $K_{\bullet\bullet}$、$K'_{\bullet\bullet}$ 的超同调定义为超同调
$L_\bullet T(L_\bullet,L'_\bullet)$；应用
（\hyperref[0.11.6.4-fr]{11.6.4}）和
（\hyperref[0.11.6.5-fr]{11.6.5}），与
（\hyperref[0.11.7.2-fr]{11.7.2}）相同，可得六个收敛到这个超同调的谱序列；
我们把写出它们的工作留给读者。
\end{enumerate}
\end{env}

\subsection{单纯复形上同调的补充}
\label{subsection:0.11.8-fr}

\begin{env}[11.8.1]
\label{0.11.8.1-fr}
设 $A$ 是有限集，$\Sigma(A)$ 是由 $A$ 的元素组成的有限序列
$\sigma=(\alpha_0,\ldots,\alpha_h)$（$A$ 的“单形”）之集；令
$|\sigma|=\{\alpha_0,\ldots,\alpha_h\}$。回顾链复形 $\mathrm C_\bullet(A)$：
它是由 $\Sigma(A)$ 的元素生成的自由分次阿贝尔群，其中
$(\alpha_0,\ldots,\alpha_h)$ 的次数为 $h$，而微分定义为
\[
d(\alpha_0,\ldots,\alpha_h)=
\sum_{j=0}^{h}(-1)^j
(\alpha_0,\ldots,\widehat{\alpha_j},\ldots,\alpha_h).
\]
$\mathrm C_\bullet(A)$ 的子群 $\mathrm D_\bullet(A)$ 由如下链生成：两个
$\alpha_i$ 相等的链 $\sigma=(\alpha_0,\ldots,\alpha_h)$，以及链
$\pi(\sigma)-\varepsilon_\pi\sigma$；这里对每个置换 $\pi$，
\[
\pi(\sigma)=(\alpha_{\pi(0)},\ldots,\alpha_{\pi(h)})
\]
且 $\varepsilon_\pi$ 是 $\pi$ 的符号。它是 $\mathrm C_\bullet(A)$ 的一个
子复形，其元素称为\emph{退化链}。令
$\mathrm L_\bullet(A)=\mathrm C_\bullet(A)/\mathrm D_\bullet(A)$，便有复形的
自然同态
\[
p:\mathrm C_\bullet(A)\longrightarrow\mathrm L_\bullet(A).
\]
另一方面，按如下方式定义复形同态
\[
j:\mathrm L_\bullet(A)\longrightarrow\mathrm C_\bullet(A)
\]
在 $A$ 上选定一个全序。对单形 $\sigma=(\alpha_0,\ldots,\alpha_h)$ 模
$\mathrm D_\bullet(A)$ 的类，若两个 $\alpha_i$ 相等，就令它对应于 $0$；否则，
令它对应于把诸 $\alpha_i$ 按递增次序排列所得的序列
$(\beta_0,\ldots,\beta_h)$。显然，$p\circ j$ 是 $\mathrm L_\bullet(A)$ 的恒等映射。
\end{env}

\begin{env}[11.8.2]
\label{0.11.8.2-fr}
设 $B$ 是另一个有限集。若 $d'$、$d''$ 分别是 $\mathrm C_\bullet(A)$ 和
$\mathrm C_\bullet(B)$ 的微分，则张量积复形
$\mathrm C_\bullet(A)\otimes\mathrm C_\bullet(B)$ 可以看成由
$\Sigma(A)\times\Sigma(B)$ 的元素生成的自由阿贝尔群，其微分为
\[
d(\sigma,\tau)=(d'\sigma,\tau)+(-1)^{h+1}(\sigma,d''\tau)
\quad\text{若 }\operatorname{Card}|\sigma|=h+1.
\]
自然同态 $\mathrm C_\bullet(A)\to\mathrm L_\bullet(A)$、
$\mathrm C_\bullet(B)\to\mathrm L_\bullet(B)$ 定义了同态
\[
p:\mathrm C_\bullet(A)\otimes\mathrm C_\bullet(B)
\longrightarrow\mathrm L_\bullet(A)\otimes\mathrm L_\bullet(B),
\]
后一个张量积同构于
\[
\frac{\mathrm C_\bullet(A)\otimes\mathrm C_\bullet(B)}
{\mathrm C_\bullet(A)\otimes\mathrm D_\bullet(B)
+\mathrm D_\bullet(A)\otimes\mathrm C_\bullet(B)}.
\]
同样，利用（\hyperref[0.11.8.1-fr]{11.8.1}）中（借助 $A$ 和 $B$ 上的全序）
定义的同态 $\mathrm L_\bullet(A)\to\mathrm C_\bullet(A)$ 和
$\mathrm L_\bullet(B)\to\mathrm C_\bullet(B)$，定义同态
\[
j:\mathrm L_\bullet(A)\otimes\mathrm L_\bullet(B)
\longrightarrow\mathrm C_\bullet(A)\otimes\mathrm C_\bullet(B)
\]
使得 $p\circ j$ 是恒等映射。
\end{env}

\oldpage[0\textsubscript{III}]{44}
沿用这些记号：
\begin{proposition}[11.8.3]
\label{0.11.8.3-fr}
\emph{存在同伦
\[
h:\mathrm C_\bullet(A)\otimes\mathrm C_\bullet(B)
\longrightarrow\mathrm C_\bullet(A)\otimes\mathrm C_\bullet(B),
\]
使得 $h(\sigma,\tau)$ 是诸单形偶对 $(\sigma_i,\tau_i)$ 的线性组合，其中
$|\sigma_i|\subset|\sigma|$、$|\tau_i|\subset|\tau|$；并且对
$f=j\circ p$ 有}
\[
\label{0.11.8.3.1-fr}
\tag{11.8.3.1}
f-1=h\circ d+d\circ h.
\]
\end{proposition}

\begin{proof}
只须用 $\sigma$ 与 $\tau$ 的次数之和作归纳，对每个单形偶对
$(\sigma,\tau)$ 定义 $h$；当该和为 $0$ 时，可以取 $h=0$。令
\[
\omega=f(\sigma,\tau)-(\sigma,\tau)-h(d(\sigma,\tau));
\]
由归纳假设和 $d$ 的定义，有
\[
\omega\in\mathrm C_\bullet(|\sigma|)\otimes\mathrm C_\bullet(|\tau|).
\]
又有
\[
d\omega=f(d(\sigma,\tau))-d(\sigma,\tau)-d(h(d(\sigma,\tau)))
=h(d(d(\sigma,\tau)))=0
\]
，这是由（\hyperref[0.11.8.3.1-fr]{11.8.3.1}）和归纳假设得到的。另一方面，
对 $q>0$，有 $H_q(\mathrm C_\bullet(A))=0$（G, I, 3.7.4），因而由 Künneth
公式（G, I, 5.5.2）也有
\[
H_q(\mathrm C_\bullet(A)\otimes\mathrm C_\bullet(B))=0\quad(q>0),
\]
在这条注记中以 $|\sigma|$ 代替 $A$、以 $|\tau|$ 代替 $B$，可见存在元素
\[
\omega'\in\mathrm C_\bullet(|\sigma|)\otimes\mathrm C_\bullet(|\tau|)
\]
使得 $\omega=d\omega'$；取 $h(\sigma,\tau)=\omega'$，便对偶对
$(\sigma,\tau)$ 验证了（\hyperref[0.11.8.3.1-fr]{11.8.3.1}），归纳可以继续。
\end{proof}

\begin{env}[11.8.4]
\label{0.11.8.4-fr}
沿用（\hyperref[0.11.8.2-fr]{11.8.2}）的记号；若
$|\sigma|\subset|\sigma'|$ 且 $|\tau|\subset|\tau'|$，就记
$(\sigma,\tau)\leq(\sigma',\tau')$。$\Sigma(A)\times\Sigma(B)$ 上的一个
\emph{系数系} $\Gamma$ 由一族阿贝尔群 $(\Gamma_{\sigma,\tau})$ 和一族同态组成；其中
$\Gamma_{\sigma,\tau}$ 只依赖于集合 $|\sigma|$ 和 $|\tau|$，而这族同态为
\[
\Gamma_{\sigma,\tau}\longrightarrow\Gamma_{\sigma',\tau'}
\quad\text{当 }(\sigma,\tau)\leq(\sigma',\tau'),
\]
，并且对这个预序关系构成一个归纳系。于是定义上链复形
$\mathrm C^\bullet(A,B;\Gamma)$：其元素是族
$\lambda=(\lambda(\sigma,\tau))$，其中 $(\sigma,\tau)$ 遍历
$\Sigma(A)\times\Sigma(B)$，并且对每个偶对 $(\sigma,\tau)$，都有
$\lambda(\sigma,\tau)\in\Gamma_{\sigma,\tau}$。微分定义如下：若
\[
d(\sigma,\tau)=\sum_i\pm(\sigma_i,\tau_i),
\]
则对每个 $i$，有 $|\sigma_i|\subset|\sigma|$、
$|\tau_i|\subset|\tau|$，并令
\[
d\lambda(\sigma,\tau)=\sum_i\pm\lambda_i(\sigma_i,\tau_i),
\]
其中 $\lambda_i(\sigma_i,\tau_i)$ 表示 $\lambda(\sigma_i,\tau_i)$ 在
$\Gamma_{\sigma,\tau}$ 中的典范像。

若上链 $\lambda\in\mathrm C^\bullet(A,B;\Gamma)$ 满足如下条件，就称它为
\emph{双交错的}：当单形 $\sigma$、$\tau$ 中有一个含两个相等项时，
$\lambda(\sigma,\tau)=0$；并且对指标的任意置换 $\pi$、$\pi'$，有
\[
\lambda(\pi(\sigma),\tau)=\varepsilon_\pi\lambda(\sigma,\tau),
\qquad
\lambda(\sigma,\pi'(\tau))=\varepsilon_{\pi'}\lambda(\sigma,\tau)
\]
显然，这些上链生成 $\mathrm C^\bullet(A,B;\Gamma)$ 的一个子复形
$\mathrm L^\bullet(A,B;\Gamma)$。
\end{env}

\begin{proposition}[11.8.5]
\label{0.11.8.5-fr}
\emph{典范单射
\[
\mathrm L^\bullet(A,B;\Gamma)\longrightarrow
\mathrm C^\bullet(A,B;\Gamma)
\]
在这两个复形的上同调之间诱导一个同构。}
\end{proposition}

\begin{proof}
注意，若 $p$ 和 $j$ 具有（\hyperref[0.11.8.2-fr]{11.8.2}）中所定义的意义，
则映射
\[
{}'p:\lambda\longmapsto\lambda\circ p,\qquad
{}'j:\lambda\longmapsto\lambda\circ j
\]
分别定义在 $\mathrm L^\bullet(A,B;\Gamma)$ 和
$\mathrm C^\bullet(A,B;\Gamma)$ 中；第一个映射正是典范单射。由于
$\mathop{}'j\circ{}'p$ 是恒等映射，只须证明 $\mathop{}'p\circ{}'j$ 与恒等映射
同伦。由（\hyperref[0.11.8.3-fr]{11.8.3}），映射
$\mathop{}'h:\lambda\mapsto\lambda\circ h$ 定义在
$\mathrm C^\bullet(A,B;\Gamma)$ 中，因而可以转置恒等式
（\hyperref[0.11.8.3.1-fr]{11.8.3.1}），这就给出所求结果。
\end{proof}

\begin{env}[11.8.6]
\label{0.11.8.6-fr}
\oldpage[0\textsubscript{III}]{45}
命题（\hyperref[0.11.8.5-fr]{11.8.5}）把
$\mathrm L^\bullet(A,B;\Gamma)$ 的上同调计算归结为
$\mathrm C^\bullet(A,B;\Gamma)$ 的上同调计算。另一方面，回顾由
Eilenberg--Zilber 定理（G, I, 3.10.2），后者典范同构于如下上链复形的上同调。
构造链复形 $\mathrm P_\bullet(A,B)$，它由 $\sigma$ 和 $\tau$ 次数相同的
$(\sigma,\tau)\in\Sigma(A)\times\Sigma(B)$ 的线性组合组成；这个复形的微分为
\[
d(\sigma,\tau)=\sum_{j,k}(-1)^{j+k}(\sigma_j,\tau_k)
\]
，如果
\[
d\sigma=\sum_j(-1)^j\sigma_j,\qquad
d\tau=\sum_k(-1)^k\tau_k;
\]
于是有两个复形的典范同态
\[
f:\mathrm P_\bullet(A,B)\longrightarrow
\mathrm C_\bullet(A)\otimes\mathrm C_\bullet(B),\qquad
g:\mathrm C_\bullet(A)\otimes\mathrm C_\bullet(B)
\longrightarrow\mathrm P_\bullet(A,B),
\]
并且可以证明（\emph{同上}），存在同伦 $h$、$h'$，使得
\[
f\circ g-1=d\circ h+h\circ d,\qquad
g\circ f-1=d\circ h'+h'\circ d.
\]
此外，
\[
f(\sigma,\tau)\in\mathrm C_\bullet(|\sigma|)
\otimes\mathrm C_\bullet(|\tau|),\qquad
g(\sigma,\tau)\in\mathrm P_\bullet(|\sigma|,|\tau|)
\]
而同伦 $h$、$h'$ 可以取成满足
\[
h(\sigma,\tau)\in\mathrm C_\bullet(|\sigma|)
\otimes\mathrm C_\bullet(|\tau|),\qquad
h'(\sigma,\tau)\in\mathrm P_\bullet(|\sigma|,|\tau|).
\]
这一点来自如下事实：$h(\sigma,\tau)$ 和 $h'(\sigma,\tau)$ 的定义可以用
$\sigma$ 与 $\tau$ 的次数之和作归纳；而
\[
% Canon current rebase 2026-08-31: these are homological chain-complex groups.
H_q(\mathrm C_\bullet(|\sigma|)\otimes\mathrm C_\bullet(|\tau|))
\quad\text{和}\quad
H_q(\mathrm P_\bullet(|\sigma|,|\tau|))
\]
对 $q>0$ 为零（\emph{同上}）。于是像在
（\hyperref[0.11.8.3-fr]{11.8.3}）中一样论证，即得结论。

现在把 $\mathrm P^\bullet(A,B;\Gamma)$ 定义为所有族
$\lambda=(\lambda(\sigma,\tau))$ 之集，其中 $(\sigma,\tau)$ 遍历两项次数相同的
偶对，且 $\lambda(\sigma,\tau)\in\Gamma_{\sigma,\tau}$。由于
\[
d\sigma=\sum_j(-1)^j\sigma_j,\qquad
d\tau=\sum_k(-1)^k\tau_k,
\]
表达式
\[
d\lambda(\sigma,\tau)=
\sum_{j,k}(-1)^{j+k}\lambda(\sigma_j,\tau_k)
\]
有定义，并给出复形 $\mathrm P^\bullet(A,B;\Gamma)$ 的微分。于是，由以上注记，
映射
\[
{}'f:\lambda\mapsto\lambda\circ f,\quad
{}'g:\lambda\mapsto\lambda\circ g,\quad
{}'h:\lambda\mapsto\lambda\circ h,\quad
{}'h':\lambda\mapsto\lambda\circ h'
\]
全都有定义；因此 $\mathop{}'f\circ{}'g$ 和 $\mathop{}'g\circ{}'f$ 都与恒等映射
同伦，从而得到 $\mathrm C^\bullet(A,B;\Gamma)$ 与
$\mathrm P^\bullet(A,B;\Gamma)$ 的上同调之间所求的同构。
\end{env}

\begin{env}[11.8.7]
\label{0.11.8.7-fr}
\emph{注记}。把（\hyperref[0.11.8.3-fr]{11.8.3}）中的论证应用于
$\mathrm C_\bullet(A)$ 和 $\mathrm L_\bullet(A)$，可知：若 $j$ 和 $p$ 按
（\hyperref[0.11.8.1-fr]{11.8.1}）定义，则 $f=j\circ p$ 仍满足关系式
（\hyperref[0.11.8.3.1-fr]{11.8.3.1}），并且
$|h(\sigma)|\subset|\sigma|$。由此像在
（\hyperref[0.11.8.5-fr]{11.8.5}）中一样，推出
$\mathrm L^\bullet(A;\Gamma)$ 的上同调到
$\mathrm C^\bullet(A;\Gamma)$ 的上同调的同构；这两个复形按显然的方式定义。
这正是（G, I, 3.8.1）中略述其证明的结果。
\end{env}

\begin{env}[11.8.8]
\label{0.11.8.8-fr}
现在重新采用（\hyperref[0.11.8.2-fr]{11.8.2}）的记号和假设，并考虑
$\Sigma(A)\times\Sigma(B)$ 上的系数系所成的复形
$\Gamma^\bullet=(\Gamma^k)$：因此，对每个 $(\sigma,\tau)$，诸
$\Gamma^k_{\sigma,\tau}$ 构成阿贝尔群的复形 $(k\in\mathbb Z)$，并且有交换图
\[
\xymatrix{
\Gamma^k_{\sigma,\tau}\ar[r]\ar[d] &
\Gamma^{k+1}_{\sigma,\tau}\ar[d]\\
\Gamma^k_{\sigma',\tau'}\ar[r] &
\Gamma^{k+1}_{\sigma',\tau'}
}
\]
\oldpage[0\textsubscript{III}]{46}
，其中 $(\sigma,\tau)\leq(\sigma',\tau')$。立即验证可知，
\[
\mathrm C^\bullet(A,B;\Gamma^\bullet)
=\bigl(\mathrm C^h(A,B;\Gamma^k)\bigr)_{(h,k)\in\mathbb Z\times\mathbb Z}
\]
是阿贝尔群的一个双复形，而
\[
\mathrm L^\bullet(A,B;\Gamma^\bullet)
=\bigl(\mathrm L^h(A,B;\Gamma^k)\bigr)
\]
是它的一个子双复形。
\end{env}

\begin{proposition}[11.8.9]
\label{0.11.8.9-fr}
\emph{典范单射
\[
\mathrm L^\bullet(A,B;\Gamma^\bullet)\longrightarrow
\mathrm C^\bullet(A,B;\Gamma^\bullet)
\]
在这两个双复形的上同调之间诱导一个同构。}
\end{proposition}

\begin{proof}
为简化记号，令 $\mathrm C^{\bullet\bullet}=
\mathrm C^\bullet(A,B;\Gamma^\bullet)$ 及
$\mathrm L^{\bullet\bullet}=
\mathrm L^\bullet(A,B;\Gamma^\bullet)$。注意，因为当 $h<0$ 时
$\mathrm C^{hk}=\mathrm L^{hk}=0$，这些双复形的第二个谱序列正则
（\hyperref[0.11.3.3-fr]{11.3.3}）；同态
$\mathrm L^{\bullet\bullet}\to\mathrm C^{\bullet\bullet}$ 因而给出谱序列的态射
\[
{}''E(\mathrm L^{\bullet\bullet})\longrightarrow
{}''E(\mathrm C^{\bullet\bullet})
\]
它在 $E_2$ 项上归结为
\[
\label{0.11.8.9.1-fr}
\tag{11.8.9.1}
H^p_{\mathrm{II}}(H^q_{\mathrm I}(\mathrm L^{\bullet\bullet}))
\longrightarrow
H^p_{\mathrm{II}}(H^q_{\mathrm I}(\mathrm C^{\bullet\bullet})).
\]
可是，对每个 $k\in\mathbb Z$，由
（\hyperref[0.11.8.5-fr]{11.8.5}）
\SourceCorrectionNote{法文印本此处引用（11.8.3），但逐列上同调同构正是命题
（11.8.5）的结论。本译依稳定勘误
EGA-III1-11.8.9-P391-CITATION-11.8.3-001 改用（11.8.5）；法文外交式来源保持不变。}
可知同态 $H^q_{\mathrm I}(\mathrm L^{\bullet,k})\to
H^q_{\mathrm I}(\mathrm C^{\bullet,k})$ 是双射；结论遂由
（\hyperref[0.11.1.5-fr]{11.1.5}）得到。
\end{proof}

\begin{env}[11.8.10]
\label{0.11.8.10-fr}
同样，沿用（\hyperref[0.11.8.6-fr]{11.8.6}）的记号，有双复形的典范同态
\[
\mathrm C^\bullet(A,B;\Gamma^\bullet)\longrightarrow
\mathrm P^\bullet(A,B;\Gamma^\bullet)
\]
（记号的意义显然）；与（\hyperref[0.11.8.9-fr]{11.8.9}）中相同的论证，
这次以（\hyperref[0.11.8.6-fr]{11.8.6}）为依据，表明这个同态仍在上同调上
给出同构。
\end{env}

\subsection{有限型复形的一个引理}
\label{subsection:0.11.9-fr}

\begin{proposition}[11.9.1]
\label{0.11.9.1-fr}
设 $\mathcal C$ 是一个阿贝尔范畴，$K'$ 和 $K''$ 是 $\mathcal C$ 的对象集的
两个子集，满足 $K''\subset K'$ 以及下列条件：
\begin{enumerate}
\item[{\rm(i)}] \emph{对每个对象 $A'\in K'$ 和 $\mathcal C$ 中的每个满射
$u:A\to A'$，都存在对象 $B\in K''$ 和态射 $v:B\to A$，使得 $u\circ v$
为满射。}
\item[{\rm(ii)}] \emph{对每一对对象 $A\in K'$、$B\in K'$ 以及每个满射
$u:A\to B$，$\operatorname{Ker}(u)$ 属于 $K'$。}
\item[{\rm(iii)}] \emph{$K''$ 中两个对象的积属于 $K''$。}
\end{enumerate}
\emph{设 $P_\bullet=(P_i)_{i\in\mathbb Z}$ 是 $\mathcal C$ 中的一个复形，
对每个 $i$ 都有 $H_i(P_\bullet)\in K'$，并且存在 $d$，使得当 $i<d$ 时
$H_i(P_\bullet)=0$。那么，存在 $\mathcal C$ 中的复形
$Q_\bullet=(Q_i)_{i\in\mathbb Z}$，使得对每个 $i$ 都有 $Q_i\in K''$，并且
当 $i<d$ 时 $Q_i=0$；还存在复形态射 $u:Q_\bullet\to P_\bullet$，使得相应的
态射 $H_\bullet(Q_\bullet)\to H_\bullet(P_\bullet)$ 是同构。}
\end{proposition}

\begin{proof}
先证明性质 (i) 的如下推论：

\emph{(i 补) 设 $u:C\to B$ 是 $\mathcal C$ 中的满射，$A$ 是 $K'$ 的对象，
$v:A\to B$ 是 $\mathcal C$ 中的态射；那么存在对象 $D\in K''$、满射
$u':D\to A$ 和态射 $v':D\to C$，使得图}
\[
\xymatrix{
D\ar[r]^{u'}\ar[d]_{v'} & A\ar[d]^{v}\\
C\ar[r]_{u} & B
}
\tag{11.9.1.1}\label{0.11.9.1.1-fr}
\]
\emph{交换。}

事实上，考虑 $\mathcal C$ 中的纤维积 $C\times_B A$ 及使下图交换的典范投影
$p:C\times_B A\to C$、$q:C\times_B A\to A$：
\oldpage[0\textsubscript{III}]{47}
\[
\xymatrix{
C\times_B A\ar[r]^{q}\ar[d]_{p} & A\ar[d]^{v}\\
C\ar[r]_{u} & B
}
\tag{11.9.1.2}\label{0.11.9.1.2-fr}
\]
已知（[27], p.~1-12），$q$ 的余核是 $A$ 对 $v^{-1}(u(C))$ 的商；由于 $u$
是满射，$u(C)=B$ 且 $v^{-1}(u(C))=A$，所以 $q$ 是满射。这时只须对满射
$q:C\times_B A\to A$ 应用 (i)：存在对象 $D\in K''$ 和态射
$w:D\to C\times_B A$，使得 $q\circ w$ 为满射；取 $u'=q\circ w$、
$v'=p\circ w$ 即可。

现在用归纳法证明命题。假设对某个 $i\geq d-1$，已经对 $j\leq i$ 构造了
对象 $Q_j$、态射 $d_j:Q_j\to Q_{j-1}$ 和态射 $u_j:Q_j\to P_j$，使得当
$j<d$ 时 $Q_j=0$，并且对 $j\leq i$ 有 $d_{j-1}\circ d_j=0$ 以及
$d_j\circ u_j=u_{j-1}\circ d_j$；此外，假设下列条件成立：
\begin{enumerate}
\item[$(\mathrm I_i)$] 对 $j\leq i$ 有 $Q_j\in K''$，对 $j<i$ 有
$B_j(Q_\bullet)\in K'$。
\item[$(\mathrm{II}_i)$] 对 $j<i$，由族 $(u_k)_{k\leq i}$ 诱导的同态
$H_j(Q_\bullet)\to H_j(P_\bullet)$ 是同构。
\item[$(\mathrm{III}_i)$] 复合态射
\[
v_i:Z_i(Q_\bullet)\longrightarrow Z_i(P_\bullet)
\longrightarrow H_i(P_\bullet)
\]
（其中左边的箭头是 $u_i$ 的限制，右边的箭头是典范态射）为满射。
\end{enumerate}
注意，由 (ii)，$Z_i(Q_\bullet)$ 是满射 $Q_i\to B_{i-1}(Q_\bullet)$ 的核，
所以根据假设 $(\mathrm I_i)$ 属于 $K'$。再由 (ii) 和假设
$(\mathrm{III}_i)$，还得到 $N_i=\operatorname{Ker}(v_i)$ 属于 $K'$。由
(i 补)，存在 $Q'_{i+1}\in K''$、满射 $d'_{i+1}:Q'_{i+1}\to N_i$ 和态射
$u'_{i+1}:Q'_{i+1}\to P_{i+1}$，使得下图交换：
\[
\xymatrix{
Q'_{i+1}\ar[r]^{d'_{i+1}}\ar[d]_{u'_{i+1}} &
N_i\ar[d]^{u_i}\\
P_{i+1}\ar[r]_{d_{i+1}} & B_i(P_\bullet)
}
\tag{11.9.1.3}\label{0.11.9.1.3-fr}
\]

由于典范态射 $Z_{i+1}(P_\bullet)\to H_{i+1}(P_\bullet)$ 是满射，并且按假设
$H_{i+1}(P_\bullet)\in K'$，由 (i) 可知，存在对象
$Q''_{i+1}\in K''$ 和态射 $u''_{i+1}:Q''_{i+1}\to Z_{i+1}(P_\bullet)$，使得
复合态射 $Q''_{i+1}\to Z_{i+1}(P_\bullet)\to H_{i+1}(P_\bullet)$ 是满射。
若取 $d''_{i+1}:Q''_{i+1}\to N_i$ 为 $0$，则图
\[
\xymatrix{
Q''_{i+1}\ar[r]^{d''_{i+1}}\ar[d]_{u''_{i+1}} &
N_i\ar[d]^{u_i}\\
Z_{i+1}(P_\bullet)\ar[r]_{d_{i+1}} & P_i
}
\tag{11.9.1.4}\label{0.11.9.1.4-fr}
\]
交换，其中下面的水平箭头为 $0$。
\oldpage[0\textsubscript{III}]{48}
于是取
\[
Q_{i+1}=Q'_{i+1}\times Q''_{i+1},\qquad
d_{i+1}=d'_{i+1}+d''_{i+1},\qquad
u_{i+1}=u'_{i+1}+u''_{i+1}.
\]
由于 $d_{i+1}(Q_{i+1})=d'_{i+1}(Q'_{i+1})=N_i\subset Z_i(Q_\bullet)$，有
$d_i\circ d_{i+1}=0$；用通常的记号，还有 $B_i(Q_\bullet)=N_i$，这验证了
$(\mathrm I_{i+1})$。图（\hyperref[0.11.9.1.3-fr]{11.9.1.3}）和
（\hyperref[0.11.9.1.4-fr]{11.9.1.4}）的交换性表明
$d_{i+1}\circ u_{i+1}=u_i\circ d_{i+1}$。按 $N_i$ 的定义，由诸 $u_k$
（$k\leq i+1$）诱导的态射
\[
H_i(Q_\bullet)=Z_i(Q_\bullet)/N_i\longrightarrow H_i(P_\bullet)
\]
就是由 $v_i$ 取商而得到的态射；由于 $v_i$ 是满射，它是同构，由此得到
$(\mathrm{II}_{i+1})$。最后，按定义有
$Q''_{i+1}\subset Z_{i+1}(Q_\bullet)$；$u''_{i+1}$ 的选取表明态射
\[
v_{i+1}:Z_{i+1}(Q_\bullet)\longrightarrow Z_{i+1}(P_\bullet)
\longrightarrow H_{i+1}(P_\bullet)
\]
是满射，因为它在 $Q''_{i+1}$ 上的限制已经是满射；这就得到
$(\mathrm{III}_{i+1})$。因此归纳可以继续，命题得证。
\end{proof}

\begin{corollary}[11.9.2]
\label{0.11.9.2-fr}
设 $A$ 是一个 Noether 环（不一定交换），$P_\bullet=(P_i)_{i\in\mathbb Z}$ 是
右 $A$-模的一个复形。假设诸 $H_i(P_\bullet)$ 是有限型右 $A$-模，并且存在
$d$，使得当 $i<d$ 时 $H_i(P_\bullet)=0$。那么，存在一个由有限秩自由右
$A$-模组成的复形 $Q_\bullet=(Q_i)_{i\in\mathbb Z}$，使得当 $i<d$ 时
$Q_i=0$；还存在复形同态 $u:Q_\bullet\to P_\bullet$，使得由 $u$ 诱导的同态
$H_\bullet(Q_\bullet)\to H_\bullet(P_\bullet)$ 是双射。
\end{corollary}

\begin{proof}
应用（\hyperref[0.11.9.1-fr]{11.9.1}）：取 $\mathcal C$ 为右 $A$-模范畴，
取 $K'$（相应地，$K''$）为有限型右 $A$-模之集（相应地，有限秩自由右
$A$-模之集）。由于假设 $A$ 为 Noether 环，
（\hyperref[0.11.9.1-fr]{11.9.1}）的条件 (i)、(ii)、(iii) 显然成立。
\end{proof}

\begin{env}[11.9.3]
\label{0.11.9.3-fr}
\emph{注记}。
\begin{enumerate}
\item[{\rm(i)}] 在（\hyperref[0.11.9.2-fr]{11.9.2}）的条件下，进一步假设
诸 $P_i$ 是平坦右 $A$-模。那么，对每个左 $A$-模 $M$，复形同态
\[
u\otimes 1:Q_\bullet\otimes_A M\longrightarrow P_\bullet\otimes_A M
\]
仍诱导同调同构 $H_\bullet(Q_\bullet\otimes_A M)\simeq
H_\bullet(P_\bullet\otimes_A M)$，正如我们将在第三章中看到的那样。
\item[{\rm(ii)}] 若不假设 $A$ 为 Noether 环，
（\hyperref[0.11.9.2-fr]{11.9.2}）的结论未必成立；事实上，把它应用于一个
除一项外均约化为 $0$ 的复形，就会推出每个有限型右 $A$-模都有由有限型自由模组成的消解
\SourceCorrectionNote{法文印本此处写作“有限型左 $A$-模”，但本论证和
（11.9.2）始终处理右 $A$-模。本译依稳定勘误
EGA-III1-11.9.3-II-P392-LEFT-RIGHT-MODULE-001 采用“右 $A$-模”；法文外交式
来源保持不变。}，而这一般并不成立（参见 Bourbaki，\emph{交换代数}，
第一章，\S~2，习题~6）。

不过，无须假设 $A$ 为 Noether 环，只须假设诸 $H_i(P_\bullet)$ 有有限的
$\infty$-表示（参见第四章）。
\end{enumerate}
\end{env}

\subsection{有限长度模复形的 Euler--Poincaré 特征}
\label{subsection:0.11.10-fr}

\begin{env}[11.10.1]
\label{0.11.10.1-fr}
设 $A$ 是一个环（不一定交换），
\[
\label{0.11.10.1.1-fr}
\tag{11.10.1.1}
M^\bullet:0\longrightarrow M^0\longrightarrow M^1
\longrightarrow\cdots\longrightarrow M^n\longrightarrow0
\]
是一个由有限长度左 $A$-模组成的有限长度复形。称这个复形的
\emph{Euler--Poincaré 特征}为数
\oldpage[0\textsubscript{III}]{49}
\[
\label{0.11.10.1.2-fr}
\tag{11.10.1.2}
\chi(M^\bullet)=\sum_{i=0}^{n}(-1)^i\operatorname{long}(M^i).
\]
\end{env}

\begin{proposition}[11.10.2]
\label{0.11.10.2-fr}
\emph{对任意由有限长度左 $A$-模组成的有限复形 $M^\bullet$，有
\[
\chi(M^\bullet)=\chi(H^\bullet(M^\bullet))
\]
（其中把 $H^\bullet(M^\bullet)$ 视为微分为零的复形）。特别地，若序列
（\hyperref[0.11.10.1.1-fr]{11.10.1.1}）正合，则
$\chi(M^\bullet)=0$。}
\end{proposition}

\begin{proof}
为简便起见，记 $B^i=B^i(M^\bullet)$、$Z^i=Z^i(M^\bullet)$、
$H^i=H^i(M^\bullet)=Z^i/B^i$；诸 $B^i$、$Z^i$、$H^i$ 都具有有限长度。
由正合序列
\[
0\longrightarrow B^i\longrightarrow Z^i\longrightarrow H^i
\longrightarrow0,
\qquad
0\longrightarrow Z^i\longrightarrow M^i\longrightarrow B^{i+1}
\longrightarrow0
\]
得到关系
\[
\operatorname{long}(Z^i)=\operatorname{long}(H^i)+
\operatorname{long}(B^i),
\qquad
\operatorname{long}(M^i)=\operatorname{long}(Z^i)+
\operatorname{long}(B^{i+1}),
\]
从而
\[
\operatorname{long}(M^i)-\operatorname{long}(H^i)=
\operatorname{long}(B^{i+1})+\operatorname{long}(B^i).
\]
将这个关系乘以 $(-1)^i$，并对 $0\leq i\leq n$ 所得的诸关系求和；注意到
$B^0=B^{n+1}=0$，便得到所求等式。
\end{proof}

\begin{corollary}[11.10.3]
\label{0.11.10.3-fr}
设 $E=(E_r^{pq})$ 是环 $A$ 上模范畴中的一个谱序列。假设诸
$E_2^{pq}$ 都是有限长度 $A$-模，并且只有有限多个序偶 $(p,q)$ 满足
$E_2^{pq}\neq0$。那么，所有复形
$E_r^\bullet=(E_r^{(n)})_{n\in\mathbb Z}$
（\hyperref[0.11.1.1-fr]{11.1.1}）的 Euler--Poincaré 特征
$\chi(E_r^\bullet)$ 都相等。若此外 $E$ 弱收敛，并令
\[
E_\infty^{(n)}=\bigoplus_{p+q=n}E_\infty^{pq}\quad(n\in\mathbb Z),
\]
则还有 $\chi(E_\infty^\bullet)=\chi(E_2^\bullet)$；这里把
$E_\infty^\bullet=(E_\infty^{(n)})_{n\in\mathbb Z}$ 视为微分为零的复形。
\end{corollary}

\begin{proof}
先注意到，若 $E_2^{pq}=0$，则对 $2\leq r<+\infty$ 有
$E_r^{pq}=0$；因此，所有复形 $E_r^\bullet$ 都是由有限长度 $A$-模组成的
有限复形。于是，对任意有限的 $r$，由
（\hyperref[0.11.10.2-fr]{11.10.2}）以及 $H^\bullet(E_r^\bullet)$ 与
$E_{r+1}^\bullet$（视为微分为零的复形）之间的同构，得到
$\chi(E_r^\bullet)=\chi(E_{r+1}^\bullet)$。诸 $E_r^{pq}$ 具有有限长度这一
假设蕴含：对每个序偶 $(p,q)$，序列
$(B_k(E_2^{pq}))_{k\geq2}$ 和 $(Z_k(E_2^{pq}))_{k\geq2}$ 都是稳定的。
又因 $E$ 弱收敛，且除有限多个序偶 $(p,q)$ 外均有 $E_2^{pq}=0$，故存在
整数 $r\geq2$，使得对每个序偶 $(p,q)$ 都有
$E_\infty^{pq}=E_r^{pq}$；关于 $\chi(E_\infty^\bullet)$ 的断言随即得出。
\end{proof}

\section{层上同调补充}
\label{section:0.12-fr}

\subsection{赋环空间上模层的上同调}
\label{subsection:0.12.1-fr}

\begin{env}[12.1.1]
\label{0.12.1.1-fr}
设 $(X,\mathcal O_X)$ 是一个赋环空间。回忆，对任意
$\mathcal O_X$-模层 $\mathcal F$，定义了上同调 $H^\bullet(X,\mathcal F)$；
它是从 $\mathcal O_X$-模层范畴 $C(X)$ 到阿贝尔群范畴的一个泛上同调函子
（T, 2.2），也就是左正合函子
$\mathcal F\to\Gamma(X,\mathcal F)$ 的导出函子。函子
$H^\bullet(X,\mathcal F)$ 同构于
\oldpage[0\textsubscript{III}]{50}
在 $X$ 上\emph{阿贝尔群层}范畴中以同样方式定义的上同调函子在范畴
$C(X)$ 上的限制（G, II, 7.2.1）。
\end{env}

\begin{env}[12.1.2]
\label{0.12.1.2-fr}
令 $A=\Gamma(X,\mathcal O_X)$。由于 $A$ 的每个元素都定义阿贝尔群
$\Gamma(X,\mathcal F)$ 的一个自同态，它由函子性定义
$H^\bullet(X,\mathcal F)$ 这个 $\delta$-函子的一个自同态；这些自同态在
每个 $H^p(X,\mathcal F)$ 上定义一个 $A$-模结构，且算子 $\partial$ 是
$A$-线性的。此外，对任意两个非负整数 $p$、$q$ 以及任意两个
$\mathcal O_X$-模层 $\mathcal F$、$\mathcal G$，有一个称为
\emph{杯积}的 $A$-模同态
\[
\label{0.12.1.2.1-fr}
\tag{12.1.2.1}
H^p(X,\mathcal F)\otimes_A H^q(X,\mathcal G)
\longrightarrow H^{p+q}(X,\mathcal F\otimes_{\mathcal O_X}\mathcal G)
\]
（G, II, 6.6）。这些同态使诸 $H^p(X,\mathcal O_X)$（$p\geq0$）的直和
$S$ 成为一个反交换分次 $A$-代数，并使诸 $H^p(X,\mathcal F)$ 的直和
成为一个分次 $S$-模。

对 $X$ 的任意一个\emph{开覆盖} $\mathfrak U$，我们总用
$\check C^\bullet(\mathfrak U,\mathcal F)$ 表示 $\mathfrak U$ 的神经取值于
系数系统 $\Gamma(U_\alpha,\mathcal F)$ 的\emph{交错}上链复形（这与
（G, II, 5.1）的记号不同）。显然，
$\check C^\bullet(\mathfrak U,\mathcal F)$ 是一个分次 $A$-模，因而这个
复形的上同调群 $\check H^p(\mathfrak U,\mathcal F)$ 具有 $A$-模结构；
此外，典范映射
\[
\check H^p(\mathfrak U,\mathcal F)\longrightarrow H^p(X,\mathcal F)
\]
（G, II, 5.4）是 $A$-线性的。
\end{env}

\begin{env}[12.1.3]
\label{0.12.1.3-fr}
设 $(X',\mathcal O_{X'})$ 是另一个赋环空间，
$f=(\psi,\theta)$ 是从 $X'$ 到 $X$ 的一个态射。

令 $A'=\Gamma(X',\mathcal O_{X'})$；$\psi$-态射 $\theta$ 典范地定义一个
环同态 $A\to A'$。设 $\mathcal F$ 是一个 $\mathcal O_X$-模层，
$\mathcal F'$ 是一个 $\mathcal O_{X'}$-模层；对任意 $f$-态射
$u:\mathcal F\to\mathcal F'$
（\hyperref[0.4.4.1-fr]{0\textsubscript{I}, 4.4.1}），我们将看到，对每个
$p\geq0$，可以定义一个\emph{双同态}
\[
\label{0.12.1.3.1-fr}
\tag{12.1.3.1}
u_p:H^p(X,\mathcal F)\longrightarrow H^p(X',\mathcal F').
\]
事实上，由于 $\psi^*$ 在 $X$ 上阿贝尔群层范畴中是正合的，
$\mathcal F\to H^\bullet(X',\psi^*(\mathcal F))$ 是这个范畴中的一个
$\delta$-函子；而且已知有一个典范的 $\delta$-函子同态
\[
\label{0.12.1.3.2-fr}
\tag{12.1.3.2}
H^\bullet(X,\mathcal F)\longrightarrow H^\bullet(X',\psi^*(\mathcal F))
\]
它由下列条件唯一确定：在次数 0 上，它化为典范同态
$\Gamma(X,\mathcal F)\to\Gamma(X',\psi^*(\mathcal F))$（T, 3.2.2）。
此外，$A$ 的每个元素都确定 $\Gamma(X,\mathcal F)$ 的一个自同态 $\mu$
以及 $\Gamma(X',\psi^*(\mathcal F))$ 的一个自同态 $\mu'$，使得图
\[
\label{0.12.1.3.3-fr}
\tag{12.1.3.3}
\xymatrix{
\Gamma(X,\mathcal F)\ar[r]\ar[d]_\mu
&\Gamma(X',\psi^*(\mathcal F))\ar[d]^{\mu'}\\
\Gamma(X,\mathcal F)\ar[r]
&\Gamma(X',\psi^*(\mathcal F))
}
\]
交换；由泛上同调函子之态射的唯一延拓性质（T, 2.2），可知 $\mu$ 与
$\mu'$ 唯一延拓到上同调，并使类似于
（\hyperref[0.12.1.3.3-fr]{12.1.3.3}）的图交换；这意味着
（\hyperref[0.12.1.3.2-fr]{12.1.3.2}）是 $A$-模同态。现在注意到
\[
f^*(\mathcal F)=\psi^*(\mathcal F)
\otimes_{\psi^*(\mathcal O_X)}\mathcal O_{X'},
\]
因而有一个典范双同态
$\psi^*(\mathcal F)\to f^*(\mathcal F)$，其源是
$\psi^*(\mathcal O_X)$-模层 $\psi^*(\mathcal F)$，目标是
$\mathcal O_{X'}$-模层 $f^*(\mathcal F)$。于是由函子性得到一个函子性的
双同态
\[
\label{0.12.1.3.4-fr}
\tag{12.1.3.4}
H^p(X',\psi^*(\mathcal F))\longrightarrow H^p(X',f^*(\mathcal F))
\]
所对应的环分别为 $A$ 和 $A'$；将这个双同态与
（\hyperref[0.12.1.3.2-fr]{12.1.3.2}）复合，得到关于 $\mathcal F$ 函子性的
典范双同态
\[
\label{0.12.1.3.5-fr}
\tag{12.1.3.5}
\theta_p:H^p(X,\mathcal F)\longrightarrow H^p(X',f^*(\mathcal F)).
\]
最后，由函子性，同态 $u^\sharp:f^*(\mathcal F)\to\mathcal F'$ 导出
$A'$-模同态 $H^p(X',f^*(\mathcal F))\to H^p(X',\mathcal F')$；它与
（\hyperref[0.12.1.3.5-fr]{12.1.3.5}）复合，就得到
（\hyperref[0.12.1.3.1-fr]{12.1.3.1}）。

设 $f'=(\psi',\theta'):X''\to X'$ 是赋环空间的另一个态射，$f''=ff'$
是复合态射。考虑到函子 $\psi'^*$ 与张量积的可交换性
（\hyperref[0.4.3.3-fr]{0\textsubscript{I}, 4.3.3}），立即验证可知，双同态
$H^p(X',f^*(\mathcal F))\to H^p(X'',f'^*(f^*(\mathcal F)))$ 与
（\hyperref[0.12.1.3.5-fr]{12.1.3.5}）的复合，就是相应的双同态
$H^p(X,\mathcal F)\to H^p(X'',f''^*(\mathcal F))$。
\end{env}

\begin{env}[12.1.4]
\label{0.12.1.4-fr}
同态（\hyperref[0.12.1.3.2-fr]{12.1.3.2}）可以如下直接定义：取
$\mathcal F$ 的一个由 $X$ 上阿贝尔群层组成的内射消解
$\mathcal L^\bullet=(\mathcal L^i)$；由于函子 $\psi^*$ 正合，
$\psi^*(\mathcal L^\bullet)$ 是由 $X'$ 上的层组成的
$\psi^*(\mathcal F)$ 的一个消解。若
$\mathcal L'^\bullet=(\mathcal L'^i)$ 是 $X'$ 上阿贝尔群层范畴中
$\psi^*(\mathcal F)$ 的一个内射消解，则存在阿贝尔群层复形的一个态射
$\psi^*(\mathcal L^\bullet)\to\mathcal L'^\bullet$，它与增广相容
（M, V, 1.1.a)），并且在同伦意义下唯一确定。由此得到阿贝尔群复形的同态
\[
\Gamma(X,\mathcal L^\bullet)\longrightarrow
\Gamma(X',\psi^*(\mathcal L^\bullet))\longrightarrow
\Gamma(X',\mathcal L'^\bullet)
\]
它们的复合在取上同调后给出 $\delta$-函子的一个态射
$H^\bullet(X,\mathcal F)\to H^\bullet(X',\psi^*(\mathcal F))$；因为它在
次数 0 上与（\hyperref[0.12.1.3.2-fr]{12.1.3.2}）重合，所以与该态射
相同（T, 2.2）。

现在考虑 $X$ 的一个\emph{开覆盖} $\mathfrak U=(U_\alpha)$，并令
$\mathfrak U'=(f^{-1}(U_\alpha))$ 为 $X'$ 的开覆盖，即 $\mathfrak U$ 的逆像开覆盖。对 $X$ 的
每个开集 $V$，典范同态
$\Gamma(V,\mathcal F)\to\Gamma(f^{-1}(V),f^*(\mathcal F))$ 立即定义
复形同态（参见 G, II, 5.1）
$\check C^\bullet(\mathfrak U,\mathcal F)\to
\check C^\bullet(\mathfrak U',f^*(\mathcal F))$，从而定义典范同态
\[
\label{0.12.1.4.1-fr}
\tag{12.1.4.1}
\theta_p:\check H^p(\mathfrak U,\mathcal F)
\longrightarrow\check H^p(\mathfrak U',f^*(\mathcal F)).
\]

\oldpage[0\textsubscript{III}]{52}
此外，有交换图
\[
\label{0.12.1.4.2-fr}
\tag{12.1.4.2}
\xymatrix{
\check H^p(\mathfrak U,\mathcal F)\ar[r]^{\theta_p}\ar[d]
&\check H^p(\mathfrak U',f^*(\mathcal F))\ar[d]\\
H^p(X,\mathcal F)\ar[r]_{\theta_p}
&H^p(X',f^*(\mathcal F))
}
\]
其中竖直箭头是（G, II, 5.2）的典范同态。为证明
（\hyperref[0.12.1.4.2-fr]{12.1.4.2}）的交换性，考虑 $\mathcal F$ 关于
$\mathfrak U$ 的（交错）上链层复形
$\mathcal C^\bullet(\mathfrak U,\mathcal F)$，使得
$\Gamma(X,\mathcal C^\bullet(\mathfrak U,\mathcal F))=
\check C^\bullet(\mathfrak U,\mathcal F)$（G, II, 5.2）。典范同态
$\Gamma(V,\mathcal F)\to
\Gamma(f^{-1}(V),\psi^*(\mathcal F))$ 定义一个 $\psi$-态射
$\mathcal C^\bullet(\mathfrak U,\mathcal F)\to
\mathcal C^\bullet(\mathfrak U',\psi^*(\mathcal F))$；采用上述记号，有交换图
\[
\xymatrix{
\Gamma(X,\mathcal C^\bullet(\mathfrak U,\mathcal F))\ar[r]\ar[d]
&\Gamma(X',\mathcal C^\bullet(\mathfrak U',\psi^*(\mathcal F)))\ar[d]\\
\Gamma(X,\mathcal L^\bullet)\ar[r]
&\Gamma(X',\psi^*(\mathcal L^\bullet))\ar[r]
&\Gamma(X',\mathcal L'^\bullet)
}
\]
取上同调后得到交换图
\[
\xymatrix{
\check H^p(\mathfrak U,\mathcal F)\ar[r]\ar[d]
&\check H^p(\mathfrak U',\psi^*(\mathcal F))\ar[d]\\
H^p(X,\mathcal F)\ar[r]
&H^p(X',\psi^*(\mathcal F))
}
\]
其中竖直箭头是（G, II, 5.2）的典范同态。只须再把这些图与交换图
\[
\xymatrix{
\check H^p(\mathfrak U',\psi^*(\mathcal F))\ar[r]\ar[d]
&\check H^p(\mathfrak U',f^*(\mathcal F))\ar[d]\\
H^p(X',\psi^*(\mathcal F))\ar[r]
&H^p(X',f^*(\mathcal F))
}
\]
结合起来；
\oldpage[0\textsubscript{III}]{53}
后一个图来自同态 $\psi^*(\mathcal F)\to f^*(\mathcal F)$ 以及
（G, II, 5.2）的典范同态的函子性，由此得到
（\hyperref[0.12.1.4.2-fr]{12.1.4.2}）的交换性。

注意，若 $A=\Gamma(X,\mathcal O_X)$、
$A'=\Gamma(X',\mathcal O_{X'})$，则同态（12.1.4.1）是对应于环 $A$ 和
$A'$ 的模双同态
\SourceCorrectionNote{法文印本此处引用（12.1.4.3），但所讨论的同态正是
（12.1.4.1）。本译依稳定勘误
EGA-III1-12.1.4-P397-CITATION-12143-001 改引（12.1.4.1）；法文外交式来源
保持不变。}。对两个态射的复合，
（\hyperref[0.12.1.4.1-fr]{12.1.4.1}）具有\emph{传递性}，类似于
（\hyperref[0.12.1.3.5-fr]{12.1.3.5}）的传递性。最后注意，在前面的定义中，
可以不从 $\mathcal F$ 的内射消解 $\mathcal L^\bullet$ 出发，而同样可以从
满足对每个 $i$ 和每个 $p>0$ 都有 $H^p(X,\mathcal L^i)=0$ 的消解出发
（G, II, 4.7.1）。
\end{env}

\begin{env}[12.1.5]
\label{0.12.1.5-fr}
设 $\mathcal F$、$\mathcal G$、$\mathcal H$ 是三个 $\mathcal O_X$-模，并考虑
一个 $\mathcal O_X$-同态
$u:\mathcal F\otimes_{\mathcal O_X}\mathcal G\to\mathcal H$；它在上同调上给出
同态
\[
\label{0.12.1.5.1-fr}
\tag{12.1.5.1}
H^p(X,\mathcal F)\otimes_A H^q(X,\mathcal G)
\longrightarrow H^{p+q}(X,\mathcal H)
\]
这是由杯积（\hyperref[0.12.1.2.1-fr]{12.1.2.1}）导出的。我们证明，在
（\hyperref[0.12.1.3-fr]{12.1.3}）的假设和记号下，有交换图
\[
\label{0.12.1.5.2-fr}
\tag{12.1.5.2}
\xymatrix@C=10pt{
H^p(X,\mathcal F)\otimes_A H^q(X,\mathcal G)\ar[r]\ar[d]
&H^{p+q}(X,\mathcal H)\ar[d]\\
H^p(X',f^*(\mathcal F))\otimes_{A'}H^q(X',f^*(\mathcal G))\ar[r]
&H^{p+q}(X',f^*(\mathcal H))
}
\]
其中竖直箭头来自典范同态
（\hyperref[0.12.1.3.5-fr]{12.1.3.5}）。为此，回忆
（\hyperref[0.12.1.5.1-fr]{12.1.5.1}）可以从 $\mathcal F$、$\mathcal G$、
$\mathcal H$ 各自的\emph{典范}消解（G, II, 4.3）
$\mathcal L^\bullet$、$\mathcal M^\bullet$、$\mathcal N^\bullet$ 出发得到；
这些消解由 $\mathcal O_X$-模组成。与 $u$ 对应的 $\mathcal O_X$-模复形的
线性映射
$\mathcal L^\bullet\otimes_{\mathcal O_X}\mathcal M^\bullet
\to\mathcal N^\bullet$ 给出 $A$-模复形的同态
$\Gamma(X,\mathcal L^\bullet)\otimes_A\Gamma(X,\mathcal M^\bullet)
\to\Gamma(X,\mathcal N^\bullet)$；取上同调后得到同态
\[
H^p(\Gamma(X,\mathcal L^\bullet))\otimes_A
H^q(\Gamma(X,\mathcal M^\bullet))
\longrightarrow H^{p+q}(\Gamma(X,\mathcal N^\bullet))
\]
（G, II, 6.6）。显然有交换图
\[
\label{0.12.1.5.3-fr}
\tag{12.1.5.3}
\xymatrix@C=8pt{
\Gamma(X,\mathcal L^\bullet)\otimes_A\Gamma(X,\mathcal M^\bullet)
\ar[r]\ar[d]
&\Gamma(X,\mathcal N^\bullet)\ar[d]\\
\Gamma(X',\psi^*(\mathcal L^\bullet))
\otimes_{\Gamma(X',\psi^*(\mathcal O_X))}
\Gamma(X',\psi^*(\mathcal M^\bullet))\ar[r]
&\Gamma(X',\psi^*(\mathcal N^\bullet))
}
\]
\oldpage[0\textsubscript{III}]{54}
取上同调后得到交换图
\[
\label{0.12.1.5.4-fr}
\tag{12.1.5.4}
\xymatrix@C=7pt{
H^p(X,\mathcal F)\otimes_A H^q(X,\mathcal G)\ar[r]\ar[d]
&H^{p+q}(X,\mathcal H)\ar[d]\\
H^p(\Gamma(X',\psi^*(\mathcal L^\bullet)))
\otimes_{\Gamma(X',\psi^*(\mathcal O_X))}
H^q(\Gamma(X',\psi^*(\mathcal M^\bullet)))\ar[r]
&H^{p+q}(\Gamma(X',\psi^*(\mathcal N^\bullet)))
}
\]
但由于 $\psi^*(\mathcal L^\bullet)$、$\psi^*(\mathcal M^\bullet)$、
$\psi^*(\mathcal N^\bullet)$ 分别是 $\psi^*(\mathcal F)$、
$\psi^*(\mathcal G)$、$\psi^*(\mathcal H)$ 的\emph{消解}，故有交换图
（G, II, 6.6.1）
\[
\label{0.12.1.5.5-fr}
\tag{12.1.5.5}
\xymatrix@C=7pt{
H^p(\Gamma(X',\psi^*(\mathcal L^\bullet)))
\otimes_{\Gamma(X',\psi^*(\mathcal O_X))}
H^q(\Gamma(X',\psi^*(\mathcal M^\bullet)))\ar[r]\ar[d]
&H^{p+q}(\Gamma(X',\psi^*(\mathcal N^\bullet)))\ar[d]\\
H^p(X',\psi^*(\mathcal F))
\otimes_{\Gamma(X',\psi^*(\mathcal O_X))}
H^q(X',\psi^*(\mathcal G))\ar[r]
&H^{p+q}(X',\psi^*(\mathcal H))
}
\]
最后，由函子性有交换图
\[
\label{0.12.1.5.6-fr}
\tag{12.1.5.6}
\xymatrix@C=7pt{
H^p(X',\psi^*(\mathcal F))
\otimes_{\Gamma(X',\psi^*(\mathcal O_X))}
H^q(X',\psi^*(\mathcal G))\ar[r]\ar[d]
&H^{p+q}(X',\psi^*(\mathcal H))\ar[d]\\
H^p(X',f^*(\mathcal F))\otimes_{A'}H^q(X',f^*(\mathcal G))\ar[r]
&H^{p+q}(X',f^*(\mathcal H))
}
\]
结合三个图（\hyperref[0.12.1.5.4-fr]{12.1.5.4}）、
（\hyperref[0.12.1.5.5-fr]{12.1.5.5}）和
（\hyperref[0.12.1.5.6-fr]{12.1.5.6}），就得到所求的交换图
（\hyperref[0.12.1.5.2-fr]{12.1.5.2}）。
\end{env}

\begin{remark}[12.1.6]
\label{0.12.1.6-fr}
\oldpage[0\textsubscript{III}]{55}
沿用（\hyperref[0.12.1.3-fr]{12.1.3}）的记号，设有交换图
\[
\label{0.12.1.6.1-fr}
\tag{12.1.6.1}
\xymatrix{
0\ar[r]&\mathcal F\ar[r]^r\ar[d]_u
&\mathcal G\ar[r]^s\ar[d]_v
&\mathcal H\ar[r]\ar[d]_w&0\\
0\ar[r]&\mathcal F'\ar[r]_{r'}
&\mathcal G'\ar[r]_{s'}
&\mathcal H'\ar[r]&0
}
\]
其中 $r$、$s$ 是 $\mathcal O_X$-模同态，$r'$、$s'$ 是
$\mathcal O_{X'}$-模同态，$u$、$v$、$w$ 是 $f$-态射，并且两行均
\emph{正合}。于是导出交换图
\[
\label{0.12.1.6.2-fr}
\tag{12.1.6.2}
\xymatrix@C=6pt{
\cdots\ar[r]&H^p(X,\mathcal F)\ar[r]\ar[d]_{u_p}
&H^p(X,\mathcal G)\ar[r]\ar[d]_{v_p}
&H^p(X,\mathcal H)\ar[r]^\partial\ar[d]_{w_p}
&H^{p+1}(X,\mathcal F)\ar[r]\ar[d]_{u_{p+1}}&\cdots\\
\cdots\ar[r]&H^p(X',\mathcal F')\ar[r]
&H^p(X',\mathcal G')\ar[r]
&H^p(X',\mathcal H')\ar[r]_\partial
&H^{p+1}(X',\mathcal F')\ar[r]&\cdots
}
\]
事实上，（\hyperref[0.12.1.6.1-fr]{12.1.6.1}）分解为
\[
\xymatrix@R=14pt{
0\ar[r]&\mathcal F\ar[r]\ar[d]
&\mathcal G\ar[r]\ar[d]
&\mathcal H\ar[r]\ar[d]&0\\
0\ar[r]&\psi^*(\mathcal F)\ar[r]\ar[d]
&\psi^*(\mathcal G)\ar[r]\ar[d]
&\psi^*(\mathcal H)\ar[r]\ar[d]&0\\
0\ar[r]&\mathcal F'\ar[r]
&\mathcal G'\ar[r]
&\mathcal H'\ar[r]&0
}
\]
其中中间一行\emph{正合}（\hyperref[0.3.7.2-fr]{0\textsubscript{I}, 3.7.2}）；
只须利用（\hyperref[0.12.1.3.2-fr]{12.1.3.2}）是 $\delta$-函子的一个同态，
以及诸 $H^p(X',\mathcal F')$ 关于 $\mathcal F'$ 构成一个 $\delta$-函子这一
事实即可。
\end{remark}

\begin{env}[12.1.7]
\label{0.12.1.7-fr}
仍在（\hyperref[0.12.1.3-fr]{12.1.3}）的假设和记号下，现在考虑
$\mathcal F=f_*(\mathcal F')=\psi_*(\mathcal F')$ 的情形；我们将看到，在相差
$H^p(X',\mathcal F')$ 的一个自同构的意义下，
（\hyperref[0.12.1.3-fr]{12.1.3}）中定义的双同态
\[
\label{0.12.1.7.1-fr}
\tag{12.1.7.1}
H^p(X,f_*(\mathcal F'))\longrightarrow H^p(X',\mathcal F')
\]
可以作为复合函子 $\mathcal F'\leadsto\Gamma(X,\psi_*(\mathcal F'))$ 的一个
谱序列的\emph{边缘同态}得到（T, 2.4）。
同态（\hyperref[0.12.1.7.1-fr]{12.1.7.1}）在
（\hyperref[0.12.1.4-fr]{12.1.4}）中给出的描述表明，在这里可以如下得到这个
同态：分别取 $\psi_*(\mathcal F')$ 和 $\mathcal F'$ 的内射消解
$\mathcal L^\bullet$ 和 $\mathcal L'^\bullet$，再取一个复形同态
$v:\psi^*(\mathcal L^\bullet)\to\mathcal L'^\bullet$，它位于典范同态
$\psi^*(\psi_*(\mathcal F'))\to\mathcal F'$ “之上”；
\oldpage[0\textsubscript{III}]{56}
然后注意到
$\Gamma(X',\mathcal L'^\bullet)=\Gamma(X,\psi_*(\mathcal L'^\bullet))$，并且
复合同态
\[
\Gamma(X,\mathcal L^\bullet)\longrightarrow
\Gamma(X',\psi^*(\mathcal L^\bullet))
\xrightarrow{\Gamma(v)}\Gamma(X',\mathcal L'^\bullet)
\]
正是
\[
\label{0.12.1.7.2-fr}
\tag{12.1.7.2}
\Gamma(v^\flat):\Gamma(X,\mathcal L^\bullet)
\longrightarrow\Gamma(X,\psi_*(\mathcal L'^\bullet))
\]
（\hyperref[0.3.7.1-fr]{0\textsubscript{I}, 3.7.1}）；而
（\hyperref[0.12.1.7.1-fr]{12.1.7.1}）由
（\hyperref[0.12.1.7.2-fr]{12.1.7.2}）取上同调得到。

另一方面，复合函子 $\mathcal F'\leadsto\Gamma(X,\psi_*(\mathcal F'))$ 的
诸谱序列，是通过在 $X$ 上阿贝尔群层范畴中取复形
$\psi_*(\mathcal L'^\bullet)$ 的一个内射 Cartan--Eilenberg 消解
$\mathcal M^{\bullet\bullet}=(\mathcal M^{ij})$ 得到的；所讨论的谱序列就是
双复形 $\Gamma(X,\mathcal M^{\bullet\bullet})$ 的谱序列（由于当 $i<0$ 或
$j<0$ 时有 $\mathcal M^{ij}=0$，它们是双正则的）。这个双复形的第一个
谱序列\emph{退化}，因为诸层 $\psi_*(\mathcal L'^i)$ 是松弛的
（G, II, 3.1.1），所以当 $q>0$ 时
$H_{\mathrm{II}}^q(\Gamma(X,\mathcal M^{i,\bullet}))
=H^q(\psi_*(\mathcal L'^i))=0$（G, II, 4.4.3）；因此有\emph{双射的}边缘同态
（\hyperref[0.11.1.6-fr]{11.1.6}）
\[
\label{0.12.1.7.3-fr}
\tag{12.1.7.3}
'E_2^{i0}=H^i(H_{\mathrm{II}}^0(\Gamma(X,\mathcal M^{\bullet\bullet})))
\longrightarrow H^i(\Gamma(X,\mathcal M^{\bullet\bullet}))
\]
并且已知（\hyperref[0.11.3.4-fr]{11.3.4}），这个同态由增广
\[
\label{0.12.1.7.4-fr}
\tag{12.1.7.4}
\Gamma(X,\psi_*(\mathcal L'^\bullet))
\longrightarrow\Gamma(X,\mathcal M^{\bullet\bullet})
\]
取上同调而来；而这个增广本身来自增广
$\eta:\psi_*(\mathcal L'^\bullet)\to\mathcal M^{\bullet\bullet}$。另一方面，
第二个谱序列有边缘同态
\[
\label{0.12.1.7.5-fr}
\tag{12.1.7.5}
''E_2^{i0}=H^i(H_{\mathrm I}^0(\Gamma(X,\mathcal M^{\bullet\bullet})))
\longrightarrow H^i(\Gamma(X,\mathcal M^{\bullet\bullet}))
\]
它们由复形同态
$Z_{\mathrm I}^0(\Gamma(X,\mathcal M^{\bullet\bullet}))
\to\Gamma(X,\mathcal M^{\bullet\bullet})$ 取上同调而来
（\hyperref[0.11.3.4-fr]{11.3.4}）。由于 $\psi_*$ 左正合，序列
\[
0\longrightarrow\psi_*(\mathcal F')\longrightarrow
\psi_*(\mathcal L'^0)\longrightarrow\psi_*(\mathcal L'^1)
\]
正合；由 Cartan--Eilenberg 消解的定义
（\hyperref[0.11.4.2-fr]{11.4.2}），因而可以取
$B_{\mathrm I}^0(\mathcal M^{\bullet\bullet})=0$、
$Z_{\mathrm I}^0(\mathcal M^{\bullet\bullet})=\mathcal L^\bullet$。由于图
\[
\xymatrix@C=50pt@R=24pt{
\mathcal L^0\ar[r]^{i^0}&\mathcal M^{00}\\
\psi_*(\mathcal F')\ar[u]^\varepsilon
\ar[r]_{\varepsilon'}\ar[ur]^{\varepsilon''}
&\psi_*(\mathcal L'^0)\ar[u]_{\eta^0}
}
\]
交换，复形的单射
$i:\mathcal L^\bullet\to\mathcal M^{\bullet\bullet}$ 与增广
$\varepsilon$ 和 $\varepsilon''$ 相容。这样便有从 $\mathcal L^\bullet$ 到
$\mathcal M^{\bullet\bullet}$ 的两个复形同态
\[
\xymatrix@C=52pt@R=22pt{
\mathcal L^\bullet\ar[r]^i\ar[dr]_{v^\flat}
&\mathcal M^{\bullet\bullet}\\
&\psi_*(\mathcal L'^\bullet)\ar[u]_\eta
}
\]
\oldpage[0\textsubscript{III}]{57}
它们都与增广 $\varepsilon$ 和 $\varepsilon''$ 相容；由于
$\mathcal L^\bullet$ 是一个\emph{内射消解}，而
$\mathcal M^{\bullet\bullet}$ 由内射层组成，故由（M, V, 1.1a)）可知这两个
同态\emph{同伦}。因此，对应的两个同态
$\Gamma(X,\mathcal L^\bullet)\to\Gamma(X,\mathcal M^{\bullet\bullet})$
也是如此；取上同调后，它们给出\emph{同一个}同态。换言之，我们已经证明，
边缘同态（\hyperref[0.12.1.7.5-fr]{12.1.7.5}），写作
$H^p(X,\psi_*(\mathcal F'))\to H^p(\Gamma(X,\mathcal M^{\bullet\bullet}))$，
是（\hyperref[0.12.1.7.1-fr]{12.1.7.1}）与
（\hyperref[0.12.1.7.3-fr]{12.1.7.3}）的复合；后一个同态写作
$H^p(X',\mathcal F')\to H^p(\Gamma(X,\mathcal M^{\bullet\bullet}))$，并且我们
已经看到它是一个\emph{同构}。断言得证。
\end{env}
\subsection{高阶直接像}
\label{subsection:0.12.2-fr}
\begin{env}[12.2.1]
\label{0.12.2.1-fr}
设 $(X,\mathcal O_X)$、$(Y,\mathcal O_Y)$ 是两个赋环空间，
$f=(\psi,\theta)$ 是从 $X$ 到 $Y$ 的态射，它定义直接像函子
\[
f_*:C(X)\longrightarrow C(Y),
\]
此外该函子正是阿贝尔群层在 $X$ 上的范畴中所定义的函子
$\psi_*$ 限制到 $C(X)$ 所得到的函子。后一个函子是加性的且左正合；
又因为 $X$ 上每个阿贝尔群层都同构于某个阿贝尔群内射层的子层，
于是定义函子 $\psi_*$ 的右导出函子
$\mathcal F\leadsto R^p\psi_*(\mathcal F)$；这些
$R^p\psi_*(\mathcal F)$ 是 $Y$ 上的阿贝尔群层，而 $R^p\psi_*$ 构成
一个泛上同调函子（T, 2.3）。
此外，层 $R^p\psi_*(\mathcal F)$ 是下述预层的层化：
$V\leadsto H^p(f^{-1}(V),\mathcal F)$（T, 3.7.2）。现在设
$\mathcal F$ 是一个 $\mathcal O_X$-模层，则
$H^p(f^{-1}(V),\mathcal F)$ 自然带有
$\Gamma(f^{-1}(V),\mathcal O_X)$-模结构，因而带有
$\Gamma(V,\psi_*(\mathcal O_X))$-模结构；同态
$\theta:\mathcal O_Y\to\psi_*(\mathcal O_X)$ 使其进而带有
$\Gamma(V,\mathcal O_Y)$-模结构。对于这样定义的结构，从开集 $V$ 限制到
开集 $V'\subset V$ 给出一个双同态，因此在每个
$R^p\psi_*(\mathcal F)$ 上定义了一个 $\mathcal O_Y$-模结构；我们将此
$\mathcal O_Y$-模层记作 $R^p f_*(\mathcal F)$，从而 $R^p f_*$ 定义为
从 $C(X)$ 到 $C(Y)$ 的加性函子。
此外，诸 $R^p f_*$ 构成一个 $\partial$-函子。事实上，若
\[
0\longrightarrow\mathcal F'\longrightarrow\mathcal F
\longrightarrow\mathcal F''\longrightarrow0
\]
是 $\mathcal O_X$-模层的正合序列，则上面对 $R^p\psi_*$ 以及
$R^p\psi_*(\mathcal F)$ 上 $\mathcal O_Y$-模结构的描述立即表明，同态
\[
\partial:R^p\psi_*(\mathcal F'')\longrightarrow
R^{p+1}\psi_*(\mathcal F')
\]
在此情形是 $\mathcal O_Y$-模同态。最后，诸 $R^p f_*$ 等同于 $f_*$ 的右导出
函子：事实上，每个 $\mathcal O_X$-模层都有一个由 $\mathcal O_X$-模层组成的
内射消解；由于这样的消解由阿贝尔群的松弛层组成（G, II, 7.1），它可以用来
计算 $R^p\psi_*(\mathcal F)$，因为对任意松弛层 $\mathcal G$ 及 $n\geq1$，
$R^n\psi_*(\mathcal G)=0$（T, 2.4.1，注 3 以及命题
3.3.2 的推论）。因此，诸 $R^p f_*$ 构成从
$C(X)$ 到 $C(Y)$ 的泛上同调函子（T, 2.3）。
\end{env}
\begin{env}[12.2.2]
\label{0.12.2.2-fr}
设 $\mathcal F$ 与 $\mathcal G$ 是两个 $\mathcal O_X$-模层。沿用
（\hyperref[0.12.2.1-fr]{12.2.1}）的记号，对于 $Y$ 的每个开集 $V$，有杯积
同态（\hyperref[0.12.1.2.1-fr]{12.1.2.1}）
\[
H^p(f^{-1}(V),\mathcal F)
\otimes_{\Gamma(f^{-1}(V),\mathcal O_X)}
H^q(f^{-1}(V),\mathcal G)
\longrightarrow
H^{p+q}(f^{-1}(V),\mathcal F\otimes_{\mathcal O_X}\mathcal G)
\]
\oldpage[0\textsubscript{III}]{58}
并且由杯积的定义（G, II, 6.6）立即可知，这些同态与从 $V$ 限制到其开子空间
$V'$（其中 $V$ 为原开集）的过程相交换。另一方面，由 $\theta$ 得到环同态
\[
\Gamma(V,\mathcal O_Y)\longrightarrow\Gamma(V,\psi_*(\mathcal O_X))
=\Gamma(f^{-1}(V),\mathcal O_X)
\]
从而得到张量积的一个典范同态
\[
\begin{aligned}
&H^p(f^{-1}(V),\mathcal F)\otimes_{\Gamma(V,\mathcal O_Y)}
H^q(f^{-1}(V),\mathcal G)\\
&\qquad\longrightarrow
H^p(f^{-1}(V),\mathcal F)
\otimes_{\Gamma(f^{-1}(V),\mathcal O_X)}
H^q(f^{-1}(V),\mathcal G)
\end{aligned}
\]
它同样与从 $V$ 到 $V'$ 的限制相容。于是通过复合得到一个
$\Gamma(V,\mathcal O_Y)$-模同态；对所考虑预层的层化，它定义了一个关于
$\mathcal F$ 和 $\mathcal G$ 函子的典范同态
\[
\label{0.12.2.2.1-fr}
\tag{12.2.2.1}
R^p f_*(\mathcal F)\otimes_{\mathcal O_Y}R^q f_*(\mathcal G)
\longrightarrow
R^{p+q}f_*(\mathcal F\otimes_{\mathcal O_X}\mathcal G).
\]
当 $p=q=0$ 时，此同态退化为（\hyperref[0.4.2.2.1-fr]{0\textsubscript{I}, 4.2.2.1}）。
\end{env}
\begin{proposition}[12.2.3]
\label{0.12.2.3-fr}
\emph{对于任意 $\mathcal O_X$-模层 $\mathcal F$ 以及任意有限秩局部自由的
$\mathcal O_Y$-模层 $\mathcal L$，存在函子性的典范同构
\[
\label{0.12.2.3.1-fr}
\tag{12.2.3.1}
R^p f_*(\mathcal F)\otimes_{\mathcal O_Y}\mathcal L
\xrightarrow{\sim}
R^p f_*(\mathcal F\otimes_{\mathcal O_X}f^*(\mathcal L)).
\]}
\end{proposition}
\begin{proof}
同态（\hyperref[0.12.2.3.1-fr]{12.2.3.1}）由复合下述同态得到；后者是
（\hyperref[0.12.2.2.1-fr]{12.2.2.1}）的特殊情形：
\[
\label{0.12.2.3.2-fr}
\tag{12.2.3.2}
R^p f_*(\mathcal F)\otimes_{\mathcal O_Y}f_*(f^*(\mathcal L))
\longrightarrow
R^p f_*(\mathcal F\otimes_{\mathcal O_X}f^*(\mathcal L))
\]
再与将（\hyperref[0.12.2.3.1-fr]{12.2.3.1}）的首项映到
（\hyperref[0.12.2.3.2-fr]{12.2.3.2}）首项的同态复合；该同态来自典范同态
（\hyperref[0.4.4.3.2-fr]{0\textsubscript{I}, 4.4.3.2}）。为验证当 $\mathcal L$
局部自由时（\hyperref[0.12.2.3.1-fr]{12.2.3.1}）是同构，只须立即归结到
$\mathcal L=\mathcal O_Y$ 的情形，因为问题在 $Y$ 上是局部的，而且所考察的函子
关于 $\mathcal L$ 是加性的。但在此情形，由（\hyperref[0.12.2.2.1-fr]{12.2.2.1}）
的定义，命题归结为验证相应预层同态是双射；这由关系
$f^*(\mathcal O_Y)=\mathcal O_X$ 立即得到。
\end{proof}
\begin{env}[12.2.4]
\label{0.12.2.4-fr}
设 $(Z,\mathcal O_Z)$ 是第三个赋环空间，$g:Y\to Z$ 是赋环空间的态射。由
（G, II, 7.1 和 3.1.1）知，对于每个内射的 $\mathcal O_X$-模层 $\mathcal G$，
$f_*(\mathcal G)$ 是阿贝尔群的松弛层，因而由（\hyperref[0.12.2.1-fr]{12.2.1}）有
$R^p g_*(f_*(\mathcal G))=0$（其中 $p>0$）。
于是（T, 2.4.1）中关于复合函子的 Leray 谱序列适用于复合函子 $g_*f_*$；
存在一个双正则谱序列，其收敛项是函子 $R^\bullet h_*$，其中 $h=g\circ f$，
而其 $E_2$ 项为
\[
\label{0.12.2.4.1-fr}
\tag{12.2.4.1}
E_2^{pq}=R^p g_*(R^q f_*(\mathcal F)).
\]
\end{env}
\begin{env}[12.2.5]
\label{0.12.2.5-fr}
在（\hyperref[0.12.2.4-fr]{12.2.4}）的条件下，我们将直接定义 $\mathcal O_Z$-模层
的典范同态
\[
\label{0.12.2.5.1-fr}
\tag{12.2.5.1}
R^n g_*(f_*(\mathcal F))\longrightarrow R^n h_*(\mathcal F)
\]
\[
\label{0.12.2.5.2-fr}
\tag{12.2.5.2}
R^n h_*(\mathcal F)\longrightarrow g_*(R^n f_*(\mathcal F))
\]
\oldpage[0\textsubscript{III}]{59}
它们可以视为 Leray 谱序列的“边缘同态”（参见
（\hyperref[0.12.1.7-fr]{12.1.7}））。只须在与高阶直接像层相伴的预层上操作即可
（\hyperref[0.12.2.1-fr]{12.2.1}）。为此，取 $Z$ 的任意开集 $W$，以及它在
$Y$ 中的逆像 $g^{-1}(W)$；有一个典范双同态
\[
\label{0.12.2.5.3-fr}
\tag{12.2.5.3}
H^n(g^{-1}(W),f_*(\mathcal F))\longrightarrow
H^n(f^{-1}(g^{-1}(W)),f^*(f_*(\mathcal F)))
\]
相应的环为 $\Gamma(g^{-1}(W),\mathcal O_Y)$ 和
$\Gamma(h^{-1}(W),\mathcal O_X)$；另一方面，典范同态
（\hyperref[0.4.4.3.3-fr]{0\textsubscript{I}, 4.4.3.3}）由函子性给出典范同态
\[
\label{0.12.2.5.4-fr}
\tag{12.2.5.4}
H^n(h^{-1}(W),f^*(f_*(\mathcal F)))\longrightarrow
H^n(h^{-1}(W),\mathcal F)
\]
它们是 $\Gamma(h^{-1}(W),\mathcal O_X)$-模同态。考虑环同态
$\Gamma(W,\mathcal O_Z)\to\Gamma(g^{-1}(W),\mathcal O_Y)$；
可见复合（\hyperref[0.12.2.5.4-fr]{12.2.5.4}）与
（\hyperref[0.12.2.5.3-fr]{12.2.5.3}）得到一个预层同态，进而给出层同态
（\hyperref[0.12.2.5.1-fr]{12.2.5.1}）。
（\hyperref[0.12.2.5.2-fr]{12.2.5.2}）的定义更为简单：按定义，
$R^n h_*(\mathcal F)$ 与预层 $W\leadsto H^n(f^{-1}(g^{-1}(W)),\mathcal F)$
相伴，而 $R^n f_*(\mathcal F)$ 与预层
$V\leadsto H^n(f^{-1}(V),\mathcal F)$ 相伴。因此有一个典范同态
\[
H^n(f^{-1}(g^{-1}(W)),\mathcal F)\longrightarrow
\Gamma(g^{-1}(W),R^n f_*(\mathcal F)),
\]
显然这些同态定义了一个预层同态，而它又定义了
（\hyperref[0.12.2.5.2-fr]{12.2.5.2}）。
\end{env}
\begin{env}[12.2.6]
\label{0.12.2.6-fr}
在（\hyperref[0.12.2.4-fr]{12.2.4}）的假设下，设 $\mathcal F$、$\mathcal G$、
$\mathcal K$ 是三个 $\mathcal O_X$-模层，且
$u:\mathcal F\otimes\mathcal G\to\mathcal K$ 是一个
$\mathcal O_X$-同态。于是有交换图
\[
\xymatrix@C=10pt@R=24pt{
R^p g_*(f_*(\mathcal F))\otimes_{\mathcal O_Z}R^q g_*(f_*(\mathcal G))
\ar[r]\ar[d]&R^{p+q}g_*(f_*(\mathcal K))\ar[d]\\
R^p h_*(\mathcal F)\otimes_{\mathcal O_Z}R^q h_*(\mathcal G)
\ar[r]&R^{p+q}h_*(\mathcal K)
}
\tag{12.2.6.1}\label{0.12.2.6.1-fr}
\]
以及
\[
\xymatrix@C=10pt@R=24pt{
R^p h_*(\mathcal F)\otimes_{\mathcal O_Z}R^q h_*(\mathcal G)
\ar[r]\ar[d]&R^{p+q}h_*(\mathcal K)\ar[d]\\
g_*(R^p f_*(\mathcal F))\otimes_{\mathcal O_Z}g_*(R^q f_*(\mathcal G))
\ar[r]&g_*(R^{p+q}f_*(\mathcal K))
}
\tag{12.2.6.2}\label{0.12.2.6.2-fr}
\]
\oldpage[0\textsubscript{III}]{60}
其中水平箭头来自（\hyperref[0.12.2.2.1-fr]{12.2.2.1}）（最后一个还与
（\hyperref[0.4.2.2.1-fr]{0\textsubscript{I}, 4.2.2.1}）复合），而竖直箭头分别来自
同态（\hyperref[0.12.2.5.1-fr]{12.2.5.1}）和
（\hyperref[0.12.2.5.2-fr]{12.2.5.2}）。
事实上，只须对相应的预层同态验证这一点；回到
（\hyperref[0.12.2.2-fr]{12.2.2}）和
（\hyperref[0.12.2.5-fr]{12.2.5}）中给出的这些同态定义，对于
（\hyperref[0.12.2.6.1-fr]{12.2.6.1}），立即归结为交换图
（\hyperref[0.12.1.5.2-fr]{12.1.5.2}）；对于
（\hyperref[0.12.2.6.2-fr]{12.2.6.2}），验证更加简单。
\end{env}
\subsection{模层的 Ext 函子补充}
\label{subsection:0.12.3-fr}

\begin{env}[12.3.1]
\label{0.12.3.1-fr}
考虑一个赋环空间 $(X,\mathcal O_X)$；我们不再回顾双函子
$\operatorname{Ext}_{\mathcal O_X}^p(X;\mathcal F,\mathcal G)$ 的定义及其主要性质：它从
$\mathcal O_X$-模层范畴取值于 $\Gamma(X,\mathcal O_X)$-模范畴；也不再回顾双函子
$\mathcal{E}xt_{\mathcal O_X}^p(\mathcal F,\mathcal G)$，它从
$\mathcal O_X$-模层范畴取值于其自身；同样不再回顾把二者联系起来的双正则谱序列
$E(\mathcal F,\mathcal G)$（T, 4.2 以及 G, II, 7.3）。
\end{env}

\begin{env}[12.3.2]
\label{0.12.3.2-fr}
按照 (M, XIV, 1) 中的同样方式，定义一个 $\mathcal O_X$-模层 $\mathcal F$ 被一个
$\mathcal O_X$-模层 $\mathcal G$ 扩张的概念，以及等价扩张类之间的复合律：
对模层所作的论证显然适用于任意阿贝尔范畴。(M, XIV, 1.1) 的第二个证明只使用
嵌入内射对象的存在性，因此对 $\mathcal O_X$-模层范畴仍然有效；从而
$\operatorname{Ext}_{\mathcal O_X}^1(X;\mathcal F,\mathcal G)$ 典范地同构于
$\mathcal F$ 被 $\mathcal G$ 扩张的类所成的阿贝尔群。
\end{env}

\begin{proposition}[12.3.3]
\label{0.12.3.3-fr}
\emph{设 $(X,\mathcal O_X)$ 是一个赋环空间，其环层
$\mathcal O_X$ 相干。于是，对任意一对相干的
$\mathcal O_X$-模层 $\mathcal F$、$\mathcal G$ 以及任意 $p\geq0$，
$\mathcal{E}xt_{\mathcal O_X}^p(\mathcal F,\mathcal G)$ 是相干的
$\mathcal O_X$-模层。}
\end{proposition}

\begin{proof}
注意 $\mathcal{E}xt_{\mathcal O_X}^p(\mathcal F,\mathcal G)$ 关于 $\mathcal F$ 构成一个
逆变的上同调函子。由于 $\mathcal F$ 相干，对每个 $p$ 及每个点 $x\in X$，存在
$x$ 的一个开邻域 $U$ 以及一个由 $(\mathcal O_X|U)$-模层组成的正合序列
\SourceCorrectionNote{法文印本在此作“$X$ 的开邻域”；现依稳定校勘
EGA-III1-12.3.3-P404-NEIGHBOURHOOD-X-FOR-x-001，改作点 $x$ 的开邻域。}
\[
0\longrightarrow\mathcal R\longrightarrow\mathcal L_{p-1}
\longrightarrow\cdots\longrightarrow\mathcal L_0
\longrightarrow\mathcal F|U\longrightarrow0
\]
其中每个 $\mathcal L_i$（$0\leq i\leq p-1$）都同构于某个
$\mathcal O_X^{n_i}|U$，而 $\mathcal R$ 相干；这是由
(\hyperref[0.5.3.2-fr]{0\textsubscript{I}, 5.3.2}) 和
(\hyperref[0.5.3.4-fr]{0\textsubscript{I}, 5.3.4}) 对 $p$ 作归纳，并利用
$\mathcal O_X$ 相干这一假设得到的。

现在注意，当 $p\geq1$ 时，对于任意满足 $\mathcal L|U$ 同构于
$\mathcal O_X^n|U$ 的 $\mathcal O_X$-模层 $\mathcal L$，有
\[
\mathcal{E}xt_{\mathcal O_X|U}^p(\mathcal L|U,\mathcal G|U)=0
\]
（T, 4.2.3）\SourceCorrectionNote{法文印本此处使用直立 $\operatorname{Ext}$；
上下文及后续正合序列所论均为层 Ext，现依稳定校勘
EGA-III1-12.3.3-P404-ROMAN-EXT-FOR-SHEAF-EXT-001 使用
$\mathcal{E}xt$。}；因此 (M, V, 7.2) 的论证适用于逆变的上同调函子
\[
\mathcal F\leadsto
\mathcal{H}om_{\mathcal O_X|U}(\mathcal F|U,\mathcal G|U),
\]
并给出正合序列
\[
\mathcal{H}om_{\mathcal O_X|U}(\mathcal L_{p-1},\mathcal G|U)
\longrightarrow
\mathcal{H}om_{\mathcal O_X|U}(\mathcal R,\mathcal G|U)
\longrightarrow
\mathcal{E}xt_{\mathcal O_X|U}^p(\mathcal F|U,\mathcal G|U)
\longrightarrow0
\]
由于该序列的前两项是相干的 $(\mathcal O_X|U)$-模层
（\hyperref[0.5.3.5-fr]{0\textsubscript{I}, 5.3.5}），故第三项也相干
（\hyperref[0.5.3.4-fr]{0\textsubscript{I}, 5.3.4}）。
\end{proof}

\oldpage[0\textsubscript{III}]{61}
\begin{proposition}[12.3.4]
\label{0.12.3.4-fr}
\emph{设 $f:X\to Y$ 是赋环空间之间的平坦态射，并设
$\mathcal F$、$\mathcal G$ 是两个 $\mathcal O_Y$-模层。}
\begin{enumerate}
\item[{\rm(i)}] \emph{存在一个上同调双函子同态
\SourceCorrectionNote{法文印本在（12.3.4.1）的箭头上带同构符号；但
12.3.4 只构造同态，而 12.3.5 在附加相干性条件下才证明其为双射。现依稳定校勘
EGA-III1-12.3.4.1-P405-SPURIOUS-ISOMORPHISM-MARK-001 改为普通箭头。}
\[
\label{0.12.3.4.1-fr}
\tag{12.3.4.1}
f^*(\mathcal{E}xt_{\mathcal O_Y}^{\bullet}(\mathcal F,\mathcal G))
\longrightarrow
\mathcal{E}xt_{\mathcal O_X}^{\bullet}(f^*(\mathcal F),f^*(\mathcal G))
\]
它在次数 $0$ 上退化为典范同态
（\hyperref[0.4.4.6-fr]{0\textsubscript{I}, 4.4.6}）。}

\item[{\rm(ii)}] \emph{存在谱序列之间的典范态射
\[
\label{0.12.3.4.2-fr}
\tag{12.3.4.2}
E(\mathcal F,\mathcal G)\longrightarrow
E(f^*(\mathcal F),f^*(\mathcal G))
\]
它在 $E_2$ 项上退化为下列同态
\[
\label{0.12.3.4.3-fr}
\tag{12.3.4.3}
H^p(Y,\mathcal{E}xt_{\mathcal O_Y}^q(\mathcal F,\mathcal G))
\longrightarrow
H^p(X,\mathcal{E}xt_{\mathcal O_X}^q(f^*(\mathcal F),f^*(\mathcal G)))
\]
这些同态由（\hyperref[0.12.3.4.1-fr]{12.3.4.1}）和
（\hyperref[0.12.1.3.1-fr]{12.1.3.1}）导出。}
\end{enumerate}
\end{proposition}

\begin{proof}
\begin{enumerate}
\item[{\rm(i)}]
由于 $f^*$ 在 $\mathcal O_Y$-模层范畴中是正合函子
（\hyperref[0.6.7.2-fr]{0\textsubscript{I}, 6.7.2}），函子
\[
\mathcal G\leadsto
f^*(\mathcal{H}om_{\mathcal O_Y}(\mathcal F,\mathcal G))
\quad\hbox{和}\quad
\mathcal G\leadsto
\mathcal{H}om_{\mathcal O_X}(f^*(\mathcal F),f^*(\mathcal G))
\]
都是左正合的；由（\hyperref[0.4.4.6-fr]{0\textsubscript{I}, 4.4.6}）典范地得到
它们的导函子之间的同态。

为计算这些导函子，取 $\mathcal G$ 的一个内射消解
$\mathcal L^\bullet=(\mathcal L^i)$，于是有态射
\[
\mathcal H^p(f^*(\mathcal{H}om_{\mathcal O_Y}
(\mathcal F,\mathcal L^\bullet)))
\longrightarrow
\mathcal H^p(\mathcal{H}om_{\mathcal O_X}
(f^*(\mathcal F),f^*(\mathcal L^\bullet)))
\]
它们是层复形的上同调。事实上，由于 $f^*$ 的正合性，有
\[
\mathcal H^p(f^*(\mathcal{H}om_{\mathcal O_Y}
(\mathcal F,\mathcal L^\bullet)))
=f^*(\mathcal H^p(\mathcal{H}om_{\mathcal O_Y}
(\mathcal F,\mathcal L^\bullet)))
=f^*(\mathcal{E}xt_{\mathcal O_Y}^p(\mathcal F,\mathcal G))
\]
按定义。另一方面，$f^*$ 的正合性意味着
$f^*(\mathcal L^\bullet)$ 是 $f^*(\mathcal G)$ 的消解；若
$\mathcal L'^\bullet=(\mathcal L'^{i})$ 是 $f^*(\mathcal G)$ 在
$\mathcal O_X$-模层范畴中的一个内射消解，则存在复形态射
$f^*(\mathcal L^\bullet)\to\mathcal L'^\bullet$，它在同伦意义下唯一，
因而在上同调中确定一个良定义的同态；
将这个同态与前面定义的同态复合，便得到
（\hyperref[0.12.3.4.1-fr]{12.3.4.1}）。

\item[{\rm(ii)}]
沿用上述记号，有一个 $\mathcal O_X$-模层复形的同态
\[
f^*(\mathcal{H}om_{\mathcal O_Y}(\mathcal F,\mathcal L^\bullet))
\longrightarrow
\mathcal{H}om_{\mathcal O_X}(f^*(\mathcal F),\mathcal L'^\bullet).
\]
设 $\mathcal M^{\bullet\bullet}$ 是复形
$\mathcal{H}om_{\mathcal O_Y}(\mathcal F,\mathcal L^\bullet)$ 在
$\mathcal O_Y$-模层范畴中的 Cartan--Eilenberg 内射消解；则由
$f^*$ 的正合性，$f^*(\mathcal M^{\bullet\bullet})$ 是复形
$f^*(\mathcal{H}om_{\mathcal O_Y}(\mathcal F,\mathcal L^\bullet))$ 的
Cartan--Eilenberg 消解；若 $\mathcal M'^{\bullet\bullet}$ 是复形
$\mathcal{H}om_{\mathcal O_X}(f^*(\mathcal F),\mathcal L'^\bullet)$ 的一个
Cartan--Eilenberg 内射消解，则存在
（\hyperref[0.11.4.2-fr]{11.4.2}）一个在同伦意义下唯一的同态
\[
f^*(\mathcal M^{\bullet\bullet})\longrightarrow
\mathcal M'^{\bullet\bullet}
\]
它与上面考察的同态相容；换言之，有一个层双复形的
$f$-态射 $\mathcal M^{\bullet\bullet}\to
\mathcal M'^{\bullet\bullet}$。由此得到双复形模的一个双同态
$\Gamma(Y,\mathcal M^{\bullet\bullet})\to
\Gamma(X,\mathcal M'^{\bullet\bullet})$，它在同伦意义下唯一，并得到一个良定义的谱序列态射
（\hyperref[0.11.3.2-fr]{11.3.2}），而这正是所求的态射
（\hyperref[0.12.3.4.2-fr]{12.3.4.2}）；（\hyperref[0.12.3.4.3-fr]{12.3.4.3}）
的刻画立即由定义得到。
\end{enumerate}
\end{proof}

\begin{proposition}[12.3.5]
\label{0.12.3.5-fr}
\emph{在（\hyperref[0.12.3.4-fr]{12.3.4}）的假设下，再假设环层
$\mathcal O_Y$ 相干；那么，对任意相干的 $\mathcal O_Y$-模层
$\mathcal F$，典范同态（\hyperref[0.12.3.4.1-fr]{12.3.4.1}）都是双射。}
\end{proposition}

\begin{proof}
\oldpage[0\textsubscript{III}]{62}
由于问题在 $Y$ 上是局部的，可以假设存在一个正合序列
\[
0\longrightarrow\mathcal R\longrightarrow\mathcal O_Y^n
\longrightarrow\mathcal F\longrightarrow0,
\]
此时 $\mathcal R$ 也是相干的 $\mathcal O_Y$-模层
（\hyperref[0.5.3.4-fr]{0\textsubscript{I}, 5.3.4}）。为证明同态
\[
f^*(\mathcal{E}xt_{\mathcal O_Y}^p(\mathcal F,\mathcal G))
\longrightarrow
\mathcal{E}xt_{\mathcal O_X}^p(f^*(\mathcal F),f^*(\mathcal G))
\]
是双射，对 $p$ 作归纳；当 $p=0$ 时，命题由
（\hyperref[0.6.7.6.1-fr]{0\textsubscript{I}, 6.7.6.1}）得到。现有交换图
\[
{\tiny
\xymatrix@C=2pt@R=24pt{
f^*(\mathcal{E}xt_{\mathcal O_Y}^{p-1}(\mathcal O_Y^n,\mathcal G))
\ar[r]\ar[d]&
f^*(\mathcal{E}xt_{\mathcal O_Y}^{p-1}(\mathcal R,\mathcal G))
\ar[r]^\partial\ar[d]&
f^*(\mathcal{E}xt_{\mathcal O_Y}^{p}(\mathcal F,\mathcal G))
\ar[r]\ar[d]&
f^*(\mathcal{E}xt_{\mathcal O_Y}^{p}(\mathcal O_Y^n,\mathcal G))
\ar[d]\\
\mathcal{E}xt_{\mathcal O_X}^{p-1}(\mathcal O_X^n,f^*(\mathcal G))
\ar[r]&
\mathcal{E}xt_{\mathcal O_X}^{p-1}(f^*(\mathcal R),f^*(\mathcal G))
\ar[r]^\partial&
\mathcal{E}xt_{\mathcal O_X}^{p}(f^*(\mathcal F),f^*(\mathcal G))
\ar[r]&
\mathcal{E}xt_{\mathcal O_X}^{p}(\mathcal O_X^n,f^*(\mathcal G)).
}}
\]
这是因为 $f^*(\mathcal O_Y)=\mathcal O_X$；又由于 $f^*$ 正合，两行都正合。
此外，有
\[
\mathcal{E}xt_{\mathcal O_Y}^p(\mathcal O_Y^n,\mathcal G)=0
\quad\hbox{对所有 }p>0
\quad\hbox{并且同样}\quad
\mathcal{E}xt_{\mathcal O_X}^p(\mathcal O_X^n,f^*(\mathcal G))=0
\quad\hbox{对所有 }p>0
\]
（T, 4.2.3）。由归纳假设，上图的前两个竖直箭头是同构，而右端各项为 0，故
\[
f^*(\mathcal{E}xt_{\mathcal O_Y}^p(\mathcal F,\mathcal G))
\longrightarrow
\mathcal{E}xt_{\mathcal O_X}^p(f^*(\mathcal F),f^*(\mathcal G))
\]
是一个同构。
\end{proof}

\subsection{直接像函子的超上同调}
\label{subsection:0.12.4-fr}

\begin{env}[12.4.1]
\label{0.12.4.1-fr}
设 $(X,\mathcal O_X)$、$(Y,\mathcal O_Y)$ 是两个赋环空间，
$f:X\to Y$ 是赋环空间的一个态射。对于任意
$\mathcal O_X$-模层复形
$\mathcal K^\bullet=(\mathcal K^i)_{i\in\mathbb Z}$，可以取 $f_*$ 关于
该复形的\emph{超上同调}（\hyperref[0.11.4.4-fr]{11.4.4}），因为
$\mathcal O_Y$-模层的阿贝尔范畴中存在滤过归纳极限且它们正合
（T, 3.1.1）。超上同调的 $\mathcal O_Y$-模层
$R^p f_*(\mathcal K^\bullet)$ 也记作
$\mathcal H^p(f,\mathcal K^\bullet)$ 或
$\mathcal H_f^p(\mathcal K^\bullet)$。回顾
$\mathcal H^\bullet(f,\mathcal K^\bullet)$ 是 $\mathcal O_Y$-模层双复形
$f_*(\mathcal L^{\bullet\bullet})$ 的上同调，其中
$\mathcal L^{\bullet\bullet}$ 是 $\mathcal O_X$-模层范畴中
$\mathcal K^\bullet$ 的一个 Cartan--Eilenberg 内射消解；
$\mathcal H^\bullet(f,\mathcal K^\bullet)$ 是两个谱序列
${}'E(f,\mathcal K^\bullet)$ 和 ${}''E(f,\mathcal K^\bullet)$
的收敛项，它们的 $E_2$ 项由下式给出：
\begin{equation}
\label{0.12.4.1.1-fr}
\tag{12.4.1.1}
{}'E_2^{pq}=\mathcal H^p(\mathcal H^q(f,\mathcal K^\bullet))
\end{equation}
\begin{equation}
\label{0.12.4.1.2-fr}
\tag{12.4.1.2}
{}''E_2^{pq}=\mathcal H^p(f,\mathcal H^q(\mathcal K^\bullet))
\quad(=R^p f_*(\mathcal H^q(\mathcal K^\bullet))).
\end{equation}
在这些公式中，我们采用一般记号 $T(A^\bullet)$ 表示函子作用于复形所得的像
（\hyperref[0.11.2.1-fr]{11.2.1}）；对 $\mathcal O_X$-模层
$\mathcal F$，用 $\mathcal H^p(f,\mathcal F)$ 代替
$R^p f_*(\mathcal F)$。还要回顾，谱序列
${}'E(f,\mathcal K^\bullet)$ 总是\emph{正则的}；当
$\mathcal K^\bullet$ 下有界时，两个谱序列
${}'E(f,\mathcal K^\bullet)$ 和 ${}''E(f,\mathcal K^\bullet)$ 都是
\emph{双正则的}，
\oldpage[0\textsubscript{III}]{63}
或者，当存在整数 $m$，使每个 $\mathcal O_X$-模层都具有长度
$\leq m$ 的松弛消解时，二者也都是双正则的
（\hyperref[0.11.4.4-fr]{11.4.4}）。
\end{env}

\begin{env}[12.4.2]
\label{0.12.4.2-fr}
同样，用 $\mathbf H^\bullet(X,\mathcal K^\bullet)$ 表示函子
$\Gamma$ 关于 $\mathcal O_X$-模层复形 $\mathcal K^\bullet$ 的超上同调；
因此 $\mathbf H^p(X,\mathcal K^\bullet)$ 是
$\Gamma(X,\mathcal O_X)$-模。还可以把
$\mathbf H^\bullet(X,\mathcal K^\bullet)$ 看作
$\mathcal H^\bullet(f,\mathcal K^\bullet)$ 的一个特殊情形，其中 $f$ 是从
$(X,\mathcal O_X)$ 到只含一个点、并赋予环
$\Gamma(X,\mathcal O_X)$ 的赋环空间的态射。

对 $X$ 的任意开集 $V$，用
$\mathbf H^\bullet(V,\mathcal K^\bullet)$ 代替
$\mathbf H^\bullet(V,\mathcal K^\bullet|V)$。
\end{env}

\begin{proposition}[12.4.3]
\label{0.12.4.3-fr}
\emph{对任意整数 $p\in\mathbb Z$，$\mathcal O_Y$-模层
$\mathcal H^p(f,\mathcal K^\bullet)$ 典范同构于 $Y$ 上预层
$U\mapsto\mathbf H^p(f^{-1}(U),\mathcal K^\bullet)$ 的层化。}
\end{proposition}

\begin{proof}
事实上，沿用（\hyperref[0.12.4.1-fr]{12.4.1}）的记号，上同调层
$\mathcal H^p(f_*(\mathcal L^{\bullet\bullet}))$ 是预层
\[
U\longmapsto H^p(\Gamma(U,f_*(\mathcal L^{\bullet\bullet})))
=H^p(\Gamma(f^{-1}(U),\mathcal L^{\bullet\bullet})).
\]
的层化。另一方面，显然
$\mathcal L^{\bullet\bullet}|f^{-1}(U)$ 是
$\mathcal K^\bullet|f^{-1}(U)$ 的 Cartan--Eilenberg 内射消解
（T, 3.1.3），故按定义
\[
H^p(\Gamma(f^{-1}(U),\mathcal L^{\bullet\bullet}))
=\mathbf H^p(f^{-1}(U),\mathcal K^\bullet)
\]
\end{proof}

\begin{proposition}[12.4.4]
\label{0.12.4.4-fr}
\emph{在以下每种情形中，超上同调
$\mathcal H^\bullet(f,\mathcal K^\bullet)$ 关于 $\mathcal K^\bullet$
都是一个上同调函子：}
\begin{enumerate}
\item[\emph{a)}]
\emph{$\mathcal K^\bullet$ 在下有界复形范畴中变动。}
\item[\emph{b)}]
\emph{存在整数 $m$，使每个 $\mathcal O_X$-模层都具有长度
$\leq m$ 的松弛消解。}
\item[\emph{c)}]
\emph{$X$ 是 Noether 空间。}
\end{enumerate}
\end{proposition}

\begin{proof}
情形 a) 和 b) 是（\hyperref[0.11.5.4-fr]{11.5.4}）的特殊情形。
另一方面，情形 c) 由（\hyperref[0.11.5.2-fr]{11.5.2}）得到，因为已知在此情形下
函子 $f_*$ 与归纳极限可交换（G, II, 3.10.1）。
\end{proof}

\begin{env}[12.4.5]
\label{0.12.4.5-fr}
现在考虑 $X$ 的一个开覆盖 $\mathfrak U=(U_\alpha)$。对 $X$ 上任意
\emph{预层}复形 $\mathcal K^\bullet=(\mathcal K^j)$，考虑双复形
$\check C^\bullet(\mathfrak U,\mathcal K^\bullet)$；其指标为 $(i,j)$ 的分量是
$\check C^i(\mathfrak U,\mathcal K^j)$，即覆盖 $\mathfrak U$ 的神经上
取值于 $\mathcal K^j$ 的交错 $i$-上链群（G, II, 5.1）。称这个双复形的上同调为
\emph{覆盖 $\mathfrak U$ 以 $\mathcal K^\bullet$ 为系数的超上同调}，并记作
\[
\mathbf H^\bullet(\mathfrak U,\mathcal K^\bullet)
=H^\bullet(\check C^\bullet(\mathfrak U,\mathcal K^\bullet)).
\]
覆盖的 Leray 谱序列（T, 3.8.1 和 G, II, 5.9.1）可如下推广到超上同调：
\end{env}

\newpage
\begin{proposition}[12.4.6]
\label{0.12.4.6-fr}
\emph{设 $\mathcal K^\bullet=(\mathcal K^i)$ 是一个
$\mathcal O_X$-模层复形。存在一个关于 $\mathcal K^\bullet$ 的正则谱函子，
其收敛项为超上同调 $\mathbf H^\bullet(X,\mathcal K^\bullet)$，且其 $E_2$ 项由}
\begin{equation}
\label{0.12.4.6.1-fr}
\tag{12.4.6.1}
E_2^{pq}=\check H^p(\mathfrak U,h^q(\mathcal K^\bullet))
\end{equation}
\emph{给出，其中 $h^q(\mathcal K^\bullet)$ 表示 $X$ 上的预层复形
$V\mapsto H^q(V,\mathcal K^\bullet)$。若 $\mathcal K^\bullet$ 下有界，
则上述谱序列是双正则的。}
\end{proposition}

\begin{proof}
取 $\mathcal K^\bullet$ 的一个 Cartan--Eilenberg 内射消解
$\mathcal L^{\bullet\bullet}$，并考虑三复形
$\check C^\bullet(\mathfrak U,\mathcal L^{\bullet\bullet})
=(\check C^i(\mathfrak U,\mathcal L^{jk}))$。先把这个三复形看作次数为
$i$ 与 $j+k$ 的双复形。由于 $i$ 只取 $\geq0$ 的值，这个双复形的第二谱序列
是正则的（\hyperref[0.11.3.3-fr]{11.3.3}）并且退化，因为对所有 $q>0$ 都有
$\check H^q(\mathfrak U,\mathcal L^{jk})=0$；这是由于
$\mathcal O_X$-模层 $\mathcal L^{jk}$ 是松弛层
\oldpage[0\textsubscript{III}]{64}
（G, II, 5.2.3）。因此（\hyperref[0.11.1.6-fr]{11.1.6}），由
（G, II, 5.2.2）有典范同构
\[
H^n(\check C^\bullet(\mathfrak U,\mathcal L^{\bullet\bullet}))
\xrightarrow{\sim}
H^n(\Gamma(X,\mathcal L^{\bullet\bullet}))
\]
故按定义（\hyperref[0.12.4.2-fr]{12.4.2}）有同构
\[
H^n(\check C^\bullet(\mathfrak U,\mathcal L^{\bullet\bullet}))
\xrightarrow{\sim}\mathbf H^n(X,\mathcal K^\bullet).
\]
另一方面，把三复形
$\check C^\bullet(\mathfrak U,\mathcal L^{\bullet\bullet})$ 看作次数为
$i+j$ 与 $k$ 的双复形。由于 $k$ 只取 $\geq0$ 的值，这个双复形的第一谱序列
总是正则的；若当 $j<j_0$ 时 $\mathcal L^{jk}=0$，即当
$\mathcal K^\bullet$ 下有界时，它是双正则的
（\hyperref[0.11.3.3-fr]{11.3.3}）。这就是所求的谱序列；事实上，对每个 $j$，
$\mathcal L^{j,\bullet}$ 是 $\mathcal K^j$ 的一个内射消解；因此
$H^q(\check C^i(\mathfrak U,\mathcal L^{j,\bullet}))$ 正是上链复形
$\check C^i(\mathfrak U,h^q(\mathcal K^j))$，命题得证。
\end{proof}

\begin{corollary}[12.4.7]
\label{0.12.4.7-fr}
\emph{若对 $\mathfrak U$ 的神经的每个单形 $\sigma$ 以及每个整数 $i$，当
$q>0$ 时都有 $H^q(U_\sigma,\mathcal K^i)=0$，则有典范同构}
\begin{equation}
\label{0.12.4.7.1-fr}
\tag{12.4.7.1}
\mathbf H^\bullet(\mathfrak U,\mathcal K^\bullet)
\xrightarrow{\sim}\mathbf H^\bullet(X,\mathcal K^\bullet).
\end{equation}
\end{corollary}

\begin{proof}
事实上，该假设蕴含当 $q>0$ 时
$\check C^\bullet(\mathfrak U,h^q(\mathcal K^i))=0$，从而对 $q>0$ 有
$E_2^{pq}=0$。由于谱序列（\hyperref[0.12.4.6.1-fr]{12.4.6.1}）
退化且正则，结论由
$\mathbf H^\bullet(\mathfrak U,\mathcal K^\bullet)$ 的定义
（\hyperref[0.12.4.5-fr]{12.4.5}）以及
（\hyperref[0.11.1.6-fr]{11.1.6}）得到。
\end{proof}

\begin{env}[12.4.8]
\label{0.12.4.8-fr}
设 $(X',\mathcal O_{X'})$ 是另一个赋环空间，并设
$f=(\psi,\theta)$ 是从 $X'$ 到 $X$ 的一个态射。用与
（\hyperref[0.12.1.3-fr]{12.1.3}）和
（\hyperref[0.12.1.4-fr]{12.1.4}）相同的方法，对
$\mathcal O_X$-模层复形 $\mathcal K^\bullet$ 的超上同调定义一个双同态
\begin{equation}
\label{0.12.4.8.1-fr}
\tag{12.4.8.1}
\mathbf H^p(X,\mathcal K^\bullet)
\longrightarrow\mathbf H^p(X',f^*(\mathcal K^\bullet)).
\end{equation}
从 $\mathcal K^\bullet$ 的一个 Cartan--Eilenberg 内射消解
$\mathcal L^{\bullet\bullet}$ 出发；由于 $\psi^*$ 正合，
$\psi^*(\mathcal L^{\bullet\bullet})$ 是
$\psi^*(\mathcal O_X)$-模层范畴中
$\psi^*(\mathcal K^\bullet)$ 的一个 Cartan--Eilenberg 消解。于是存在态射
$\psi^*(\mathcal L^{\bullet\bullet})\to
\mathcal L'^{\bullet\bullet}$，其中 $\mathcal L'^{\bullet\bullet}$ 是
$\psi^*(\mathcal K^\bullet)$ 的一个 Cartan--Eilenberg 内射消解；
\newpage
由此得到上同调的一个态射
\[
\mathbf H^\bullet(X,\mathcal K^\bullet)
\longrightarrow\mathbf H^\bullet(X',\psi^*(\mathcal K^\bullet))~;
\]
再与由函子性从
$\psi^*(\mathcal K^\bullet)\to f^*(\mathcal K^\bullet)$ 所导出的态射复合，
便得到所求的态射（\hyperref[0.12.4.8.1-fr]{12.4.8.1}）。

从（\hyperref[0.12.4.8.1-fr]{12.4.8.1}）和
（\hyperref[0.12.4.3-fr]{12.4.3}）出发，并如
（\hyperref[0.12.2.5-fr]{12.2.5}）那样论证，对于赋环空间的两个态射
$f:X\to Y$、$g:Y\to Z$，可以定义 $\mathcal O_X$-模层复形
$\mathcal K^\bullet$ 的超上同调同态
\begin{equation}
\label{0.12.4.8.2-fr}
\tag{12.4.8.2}
\mathcal H^n(g,f_*(\mathcal K^\bullet))
\longrightarrow\mathcal H^n(h,\mathcal K^\bullet)
\end{equation}
\begin{equation}
\label{0.12.4.8.3-fr}
\tag{12.4.8.3}
\mathcal H^n(h,\mathcal K^\bullet)
\longrightarrow g_*(\mathcal H^n(f,\mathcal K^\bullet)).
\end{equation}
定义的细节留给读者。
\end{env}

\section{同调代数中的逆极限}
\label{section:0.13-fr}

\subsection{Mittag-Leffler 条件}
\label{subsection:0.13.1-fr}

\begin{env}[13.1.1]
\label{0.13.1.1-fr}
设 $\mathcal C$ 是一个存在无限积的阿贝尔范畴（T, 1.5 的公理
AB~3*）；则 $\mathcal C$ 的一个对象的任意一族子对象都存在下确界，
而且 $\mathcal C$ 的对象的每个逆系统都存在逆极限；取逆极限是所考虑的
逆系统的左正合函子（T, 1.8）。设
$(A_\alpha,f_{\alpha\beta})$
\oldpage[0\textsubscript{III}]{65}
是 $\mathcal C$ 的对象的一个逆系统，其指标集 $I$ 向右滤过；令
$A=\varprojlim A_\alpha$，并且对每个 $\alpha\in I$，令
$f_\alpha:A\to A_\alpha$ 为典范态射。对每个 $\alpha\in I$，当
$\alpha\leq\beta$ 时，诸 $f_{\alpha\beta}(A_\beta)$ 构成
$A_\alpha$ 的子对象的一个递减有向族；子对象
\[
A'_\alpha=\inf_{\beta\geq\alpha}f_{\alpha\beta}(A_\beta)
\]
称为 $A_\alpha$ 中的“稳定像”子对象；显然
$f_\alpha(A)\subset A'_\alpha$，并且当 $\alpha\leq\beta$ 时有
$f_{\alpha\beta}(A'_\beta)\subset A'_\alpha$；因此
$(A'_\alpha,f_{\alpha\beta}|A'_\beta)$ 是一个逆系统，且
$A=\varprojlim A'_\alpha$。
\end{env}

\begin{env}[13.1.2]
\label{0.13.1.2-fr}
给定 $\mathcal C$ 中的一个逆系统 $(A_\alpha,f_{\alpha\beta})$，称下述条件为
\emph{Mittag-Leffler 条件}：
\begin{itemize}
\item[\textbf{(ML)}]
\emph{对每个指标 $\alpha$，存在 $\beta\geq\alpha$，使得对每个
$\gamma\geq\beta$，都有
$f_{\alpha\gamma}(A_\gamma)=f_{\alpha\beta}(A_\beta)$。}
\end{itemize}
若诸 $f_{\alpha\beta}$ 都是\emph{满射}，则条件 (ML) 显然成立。
反之，若 (ML) 成立，并且对每个 $\alpha\in I$，$A'_\alpha$ 是
$A_\alpha$ 中的“稳定像”子对象，则当 $\alpha\leq\beta$ 时，
$f_{\alpha\beta}$ 在 $A'_\beta$ 上的限制是一个\emph{满射}
$A'_\beta\to A'_\alpha$：事实上，若 $\gamma\geq\beta$ 使得对
$\delta\geq\gamma$ 有
$f_{\beta\delta}(A_\delta)=f_{\beta\gamma}(A_\gamma)$，则
$A'_\beta=f_{\beta\gamma}(A_\gamma)$；另一方面，这蕴含对
$\delta\geq\gamma$ 有
$f_{\alpha\delta}(A_\delta)=f_{\alpha\gamma}(A_\gamma)$，因而
$A'_\alpha=f_{\alpha\gamma}(A_\gamma)=f_{\alpha\beta}(A'_\beta)$。

还要指出，当诸对象 $A_\alpha$ 在 $\mathcal C$ 中是\emph{Artin 的}时，
条件 (ML) 也成立；这就是说，$A_\alpha$ 的每一族子对象都有一个极小元：
事实上，此时 $A_\alpha$ 的子对象的递减有向族
$(f_{\alpha\beta}(A_\beta))$ 的一个极小元必然就是这些子对象中最小的一个。
\end{env}

\begin{remark}[13.1.3]
\label{0.13.1.3-fr}
当 $\mathcal C$ 例如是集合范畴时，也可以表述条件 (ML)；此时仍可定义
$A_\alpha$ 的“稳定像”子集，并且在
（\hyperref[0.13.1.1-fr]{13.1.1}）和
（\hyperref[0.13.1.2-fr]{13.1.2}）中对此所作的说明仍然成立。
\end{remark}

\subsection{阿贝尔群的 Mittag-Leffler 条件}
\label{subsection:0.13.2-fr}

\begin{proposition}[13.2.1]
\label{0.13.2.1-fr}
\emph{设}
\[
0\longrightarrow A_\alpha\xrightarrow{u_\alpha}B_\alpha
\xrightarrow{v_\alpha}C_\alpha\longrightarrow0
\]
\emph{是阿贝尔群逆系统的一个正合列（这些逆系统具有同一个滤过指标集 $I$）。}
\begin{enumerate}
\item[\emph{(i)}]
\emph{若 $(B_\alpha)$ 满足 (ML)，则 $(C_\alpha)$ 也满足 (ML)。}
\item[\emph{(ii)}]
\emph{若 $(A_\alpha)$ 和 $(C_\alpha)$ 满足 (ML)，则 $(B_\alpha)$ 也满足
(ML)。}
\end{enumerate}
\end{proposition}

\begin{proof}
分别以 $(f_{\alpha\beta})$、$(g_{\alpha\beta})$、
$(h_{\alpha\beta})$ 表示定义逆系统 $(A_\alpha)$、$(B_\alpha)$、
$(C_\alpha)$ 的同态系统。

(i) 假设当 $\lambda\geq\beta$ 时有
$g_{\alpha\beta}(B_\beta)=g_{\alpha\lambda}(B_\lambda)$；由于
$v_\beta$ 和 $v_\lambda$ 都是满射，当 $\lambda\geq\beta$ 时有
\[
h_{\alpha\beta}(C_\beta)
=v_\alpha(g_{\alpha\beta}(B_\beta))
=v_\alpha(g_{\alpha\lambda}(B_\lambda))
=h_{\alpha\lambda}(C_\lambda)
\]

(ii) 设 $\alpha\in I$，并取指标 $\beta\geq\alpha$，使得当
$\lambda\geq\beta$ 时有
$f_{\alpha\beta}(A_\beta)=f_{\alpha\lambda}(A_\lambda)$；再取指标
$\gamma\geq\beta$，使得当 $\lambda\geq\gamma$ 时有
$h_{\beta\gamma}(C_\gamma)=h_{\beta\lambda}(C_\lambda)$。现设
$y_\alpha$ 是 $g_{\alpha\gamma}(B_\gamma)$ 的一个元素；于是
$y_\alpha=g_{\alpha\gamma}(y_\gamma)$，其中 $y_\gamma\in B_\gamma$；令
$y_\beta=g_{\beta\gamma}(y_\gamma)$，从而
$v_\beta(y_\beta)=h_{\beta\gamma}(v_\gamma(y_\gamma))$。对每个
$\lambda\geq\gamma$，由假设存在 $y_\lambda\in B_\lambda$，使得
$h_{\beta\gamma}(v_\gamma(y_\gamma))
=h_{\beta\lambda}(v_\lambda(y_\lambda))
=v_\beta(g_{\beta\lambda}(y_\lambda))$；因此
$v_\beta(y_\beta-g_{\beta\lambda}(y_\lambda))=0$，从而存在
$x_\beta\in A_\beta$
\oldpage[0\textsubscript{III}]{66}
使得
$y_\beta=g_{\beta\lambda}(y_\lambda)+u_\beta(x_\beta)$。由此得到
\[
y_\alpha=g_{\alpha\lambda}(y_\lambda)
+u_\alpha(f_{\alpha\beta}(x_\beta))~;
\]
但由于 $\lambda\geq\beta$，存在 $x_\lambda\in A_\lambda$，使得
$f_{\alpha\beta}(x_\beta)=f_{\alpha\lambda}(x_\lambda)$，最终有
\[
y_\alpha
=g_{\alpha\lambda}(y_\lambda)+u_\alpha(f_{\alpha\lambda}(x_\lambda))
=g_{\alpha\lambda}(y_\lambda+u_\lambda(x_\lambda))
\in g_{\alpha\lambda}(B_\lambda),
\]
这就完成了证明。
\end{proof}

\begin{proposition}[13.2.2]
\label{0.13.2.2-fr}
\emph{设 $I$ 是一个有可数共尾子集的滤过有序集，并设}
\[
0\longrightarrow A_\alpha\xrightarrow{u_\alpha}B_\alpha
\xrightarrow{v_\alpha}C_\alpha\longrightarrow0
\]
\emph{是以 $I$ 为指标集的阿贝尔群逆系统的一个正合列。若 $(A_\alpha)$
满足条件 (ML)，则序列}
\[
0\longrightarrow\varprojlim A_\alpha
\longrightarrow\varprojlim B_\alpha
\longrightarrow\varprojlim C_\alpha\longrightarrow0
\]
\emph{正合。}
\end{proposition}

\begin{proof}
只须证明同态
\[
v=\varprojlim v_\alpha:\varprojlim B_\alpha\longrightarrow
\varprojlim C_\alpha
\]
是满射。设 $z=(z_\alpha)$ 是 $\varprojlim C_\alpha$ 的一个元素，并令
$E_\alpha=v_\alpha^{-1}(z_\alpha)$；显然，在同态
$g_{\alpha\beta}:B_\beta\to B_\alpha$ 的限制之下，诸 $E_\alpha$
构成一个非空集合的逆系统。我们来证明这个逆系统满足条件 (ML)；通过
$u_\alpha$ 把 $A_\alpha$ 视为 $B_\alpha$ 的子集，则对每个
$\alpha\in I$，存在 $\beta\geq\alpha$，使得当 $\lambda\geq\beta$ 时有
$g_{\alpha\beta}(A_\beta)=g_{\alpha\lambda}(A_\lambda)$；我们证明，当
$\lambda\geq\beta$ 时也有
$g_{\alpha\beta}(E_\beta)=g_{\alpha\lambda}(E_\lambda)$。事实上，取
$y_\lambda\in E_\lambda$，并令
$y_\beta=g_{\beta\lambda}(y_\lambda)$、
$y_\alpha=g_{\alpha\lambda}(y_\lambda)$；取
$y'_\alpha\in g_{\alpha\beta}(E_\beta)$，于是对某个
$y'_\beta\in E_\beta$ 有 $y'_\alpha=g_{\alpha\beta}(y'_\beta)$；
有 $y'_\beta-y_\beta=x_\beta\in A_\beta$，而由假设，存在
$x_\lambda\in A_\lambda$ 使得
$g_{\alpha\beta}(x_\beta)=g_{\alpha\lambda}(x_\lambda)$；因而
\[
y'_\alpha=g_{\alpha\beta}(y_\beta)+g_{\alpha\beta}(x_\beta)
=g_{\alpha\lambda}(y_\lambda)+g_{\alpha\lambda}(x_\lambda)
=g_{\alpha\lambda}(y_\lambda+x_\lambda)
\in g_{\alpha\lambda}(E_\lambda),
\]
这就证明了我们的断言。已知（Bourbaki，\emph{Top. gén.}，第 II 章，
第 3 版，§~3，定理 1），在对 $I$ 所作的假设下，满足 (ML) 的非空集合
逆系统具有非空的逆极限；因此存在一点
$y=(y_\alpha)\in\varprojlim E_\alpha$，而按定义，对每个 $\alpha$ 都有
$v_\alpha(y_\alpha)=z_\alpha$，故 $z=v(y)$。证毕。
\end{proof}

\begin{proposition}[13.2.3]
\label{0.13.2.3-fr}
\emph{设 $I$ 满足（\hyperref[0.13.2.2-fr]{13.2.2}）中的假设，并设
$(K_\alpha^\bullet)_{\alpha\in I}$ 是阿贝尔群复形
$K_\alpha^\bullet=(K_\alpha^n)_{n\in\mathbb Z}$ 的一个逆系统，其中微分算子
的次数为 $+1$。对每个 $n$，存在一个具有函子性的典范同态}
\begin{equation}
\label{0.13.2.3.1-fr}
\tag{13.2.3.1}
h_n:H^n(\varprojlim K_\alpha^\bullet)
\longrightarrow\varprojlim H^n(K_\alpha^\bullet).
\end{equation}
\emph{若对每个次数 $n$，阿贝尔群的逆系统
$(K_\alpha^n)_{\alpha\in I}$ 满足 (ML)，则所有同态 $h_n$ 都是满射。
若此外对某个次数 $n$，逆系统
$(H^{n-1}(K_\alpha^\bullet))_{\alpha\in I}$ 满足 (ML)，则同态 $h_n$
是双射。}
\end{proposition}

\begin{proof}
对每个 $n$，令 $K^n=\varprojlim K_\alpha^n$；同态 $h_n$ 的定义来自诸图
\[
\xymatrix{
\cdots\ar[r]&K^{n-1}\ar[r]\ar[d]&K^n\ar[r]\ar[d]
&K^{n+1}\ar[r]\ar[d]&\cdots\\
\cdots\ar[r]&K_\alpha^{n-1}\ar[r]&K_\alpha^n\ar[r]
&K_\alpha^{n+1}\ar[r]&\cdots
}
\]
\oldpage[0\textsubscript{III}]{67}
的交换性，其中 $K^\bullet$ 中的微分算子是
$K_\alpha^\bullet$ 中相应微分算子的逆极限。

考虑正合列
\[
(*_n)\qquad
0\longrightarrow \mathrm B^n(K_\alpha^\bullet)
\longrightarrow \mathrm Z^n(K_\alpha^\bullet)
\longrightarrow H^n(K_\alpha^\bullet)\longrightarrow0
\]
\[
(**_n)\qquad
0\longrightarrow \mathrm Z^{n-1}(K_\alpha^\bullet)
\longrightarrow K_\alpha^{n-1}
\longrightarrow \mathrm B^n(K_\alpha^\bullet)\longrightarrow0.
\]
由假设和命题（\hyperref[0.13.2.1-fr]{13.2.1}, (i)）可知，对每个
$n$，逆系统
$(\mathrm B^n(K_\alpha^\bullet))_{\alpha\in I}$ 满足 (ML)；因此由
（\hyperref[0.13.2.2-fr]{13.2.2}）可知序列
\[
(***_n)\qquad
0\longrightarrow\varprojlim_\alpha\mathrm B^n(K_\alpha^\bullet)
\longrightarrow\varprojlim_\alpha\mathrm Z^n(K_\alpha^\bullet)
\longrightarrow\varprojlim_\alpha H^n(K_\alpha^\bullet)
\longrightarrow0
\]
正合。显然，$\varprojlim_\alpha\mathrm B^n(K_\alpha^\bullet)$ 可与
$K^n$ 的一个包含 $\mathrm B^n(K^\bullet)$ 的子群等同
\SourceCorrectionNote{法文印本及外交式转录此处作 $K^{n+1}$；但
$\mathrm B^n(K^\bullet)\subset K^n$，而同一证明随后也把该逆极限嵌入
$K^n$。当前校勘本已改作 $K^n$。本译采用次数相合的读法，并保留可逆
来源说明；该候选已送交 Canon keeper 去重裁定。}，而
$\varprojlim_\alpha\mathrm Z^n(K_\alpha^\bullet)$ 可与
$\mathrm Z^n(K^\bullet)$ 的一个子群等同；因此 $h_n$ 是满射。
现在再假设逆系统
$(H^{n-1}(K_\alpha^\bullet))_{\alpha\in I}$ 满足 (ML)，则正合列
$(*_{n-1})$ 和命题（\hyperref[0.13.2.1-fr]{13.2.1}, (ii)）说明逆系统
$(\mathrm Z^{n-1}(K_\alpha^\bullet))_{\alpha\in I}$ 满足 (ML)；
此时，把（\hyperref[0.13.2.2-fr]{13.2.2}）应用于正合列
$(**_n)$，可知序列
\[
0\longrightarrow
\varprojlim_\alpha\mathrm Z^{n-1}(K_\alpha^\bullet)
\longrightarrow K^{n-1}\xrightarrow{u}
\varprojlim_\alpha\mathrm B^n(K_\alpha^\bullet)
\longrightarrow0
\]
正合；由于
$\varprojlim_\alpha\mathrm B^n(K_\alpha^\bullet)
\supset\mathrm B^n(K^\bullet)$，而
$\varprojlim_\alpha\mathrm B^n(K_\alpha^\bullet)\to K^n$ 的单射与
$u$ 的复合正是微分算子 $K^{n-1}\to K^n$，所以 $u$ 的满射性蕴含
$\varprojlim_\alpha\mathrm B^n(K_\alpha^\bullet)
=\mathrm B^n(K^\bullet)$，故 $h_n$ 是单射。证毕。
\end{proof}

\begin{env}[13.2.4]
\label{0.13.2.4-fr}
\emph{说明} (13.2.4)。
\begin{enumerate}
\item[\emph{(i)}]
（\hyperref[0.13.2.2-fr]{13.2.2}）的论证（参见 Bourbaki，
\emph{loc. cit.}）表明，只须假设诸 $A_\alpha$ 可赋予\emph{完备可度量空间}的结构，
在这些结构中平移是同胚；定义逆系统 $(A_\alpha)$ 的映射
$f_{\alpha\beta}:A_\beta\to A_\alpha$ 关于所考虑的距离\emph{一致连续}；
并且系统 $(A_\alpha)$ 满足条件
\begin{itemize}
\item[\textbf{(ML')}]
\emph{对每个指标 $\alpha$，存在 $\beta\geq\alpha$，使得对每个
$\gamma\geq\beta$，$f_{\alpha\gamma}(A_\gamma)$ 在
$f_{\alpha\beta}(A_\beta)$ 中稠密，}
\end{itemize}
该命题的结论仍然成立。

这使我们可以给（\hyperref[0.13.2.3-fr]{13.2.3}）补充一个类似结果：
假设对每个 $\alpha$，当 $n<0$ 时有 $K_\alpha^n=0$；再假设当
$n\geq0$ 时 $(K_\alpha^n)_{\alpha\in I}$ 满足 (ML)，并且诸
$A_\alpha=H^0(K_\alpha^\bullet)$ 可赋予满足上述性质的度量空间结构。
那么，当 $n\geq2$ 时，（\hyperref[0.13.2.3-fr]{13.2.3}）的结论不变，
而且 $h_1$ 还是\emph{双射}；这是因为
（\hyperref[0.13.2.2-fr]{13.2.2}）的论证还表明逆系统
$(\mathrm B^1(K_\alpha^\bullet))_{\alpha\in I}$ 满足 (ML)，序列
$(***_1)$ 正合，并且最后由前述结果有
$\varprojlim\mathrm B^1(K_\alpha^\bullet)=\mathrm B^1(K^\bullet)$。
由此特别得到（T, 3.10.2）的断言。

\item[\emph{(ii)}]
可以引入函子 $\varprojlim$ 的右导出函子 $\varprojlim^{(n)}$，并得到比上述结果
更完备的陈述 [28]\SourceCorrectionNote{法文印本及其外交式转录在本句两处均写作归纳极限；
EGA III-2 印本第 218 页的已发表勘误明确要求把这两个符号改为逆极限符号。
本译依稳定勘误 EGA-III1-13.2.4-P411-LIMIT-DIRECTION-ERRATUM-001 使用校正读法；
法文外交式来源保持不变。}。
\end{enumerate}
\end{env}

\subsection{应用：层的逆极限的上同调}
\label{subsection:0.13.3-fr}

\begin{proposition}[13.3.1]
\label{0.13.3.1-fr}
\oldpage[0\textsubscript{III}]{68}
\emph{设 $X$ 是一个拓扑空间，
$(\mathcal F_k)_{k\in\mathbb N}$ 是 $X$ 上阿贝尔群层的一个逆系统，并令
$\mathcal F=\varprojlim_k\mathcal F_k$。假设下列条件成立：}
\begin{enumerate}
\item[\emph{(i)}]
\emph{$X$ 的拓扑存在一个基 $\mathfrak B$，使得对每个
$U\in\mathfrak B$ 和每个 $i\geq0$，逆系统
$(H^i(U,\mathcal F_k))_{k\in\mathbb N}$ 满足 (ML)。}
\item[\emph{(ii)}]
\emph{对每个 $x\in X$ 和每个 $i>0$，都有}
\[
\varinjlim_U\bigl(\varprojlim_k H^i(U,\mathcal F_k)\bigr)=0
\]
\emph{其中 $U$ 遍历 $\mathfrak B$ 中属于 $x$ 的邻域的那些元素。}
\item[\emph{(iii)}]
\emph{定义逆系统 $(\mathcal F_k)$ 的同态
$u_{hk}:\mathcal F_k\to\mathcal F_h$ $(h\leq k)$ 都是满射。}
\end{enumerate}
\emph{在这些条件下，对每个 $i>0$，典范同态}
\[
h_i:H^i(X,\mathcal F)\longrightarrow
\varprojlim_k H^i(X,\mathcal F_k)
\]
\emph{是满射；若此外对某个 $i$，逆系统
$(H^{i-1}(X,\mathcal F_k))_{k\in\mathbb N}$ 满足 (ML)，则 $h_i$ 是双射。}
\end{proposition}

\begin{proof}
\emph{a)} 首先假设诸 $\mathcal F_k$ 以及诸 $u_{hk}$ 的核
$\mathcal N_{hk}$ 都是松弛层；我们来证明，此时命题的条件 (iii) 蕴含当
$i>0$ 时 $H^i(X,\mathcal F)=0$。只须证明：对 $X$ 的每个开集 $U$ 以及
$U$ 的每个由 $U$ 的开集组成的覆盖 $\mathfrak U$，当 $i>0$ 时都有
$\check H^i(\mathfrak U,\mathcal F)=0$。事实上，由此首先得到，对 Čech
上同调，当 $i>0$ 时有 $\check H^i(U,\mathcal F)=0$；再由（G, II, 5.9.2）
应用于 $X$ 的全体开集，得到对每个 $i>0$ 有
$H^i(X,\mathcal F)=0$。由于诸 $\mathcal F_k$ 都是松弛层，对 $i>0$ 有
$\check H^i(\mathfrak U,\mathcal F_k)=0$（G, II, 5.2.3）；对每个 $k$，
考虑覆盖 $\mathfrak U$ 的神经的交错上链复形
$\check C^\bullet(\mathfrak U,\mathcal F_k)$（G, II, 5.1），它们显然构成
阿贝尔群复形的一个逆系统。我们证明，所有映射
$\check C^\bullet(\mathfrak U,\mathcal F_k)\to
\check C^\bullet(\mathfrak U,\mathcal F_h)$ $(h\leq k)$ 都是满射。
根据定义，显然只须证明：
对 $X$ 的每个开集 $V$，映射
$\Gamma(V,\mathcal F_k)\to\Gamma(V,\mathcal F_h)$ 是满射；但由假设，序列
\[
0\longrightarrow\mathcal N_{hk}\longrightarrow\mathcal F_k
\longrightarrow\mathcal F_h\longrightarrow0
\]
正合，因而给出上同调正合列
\[
\Gamma(V,\mathcal F_k)\longrightarrow\Gamma(V,\mathcal F_h)
\longrightarrow H^1(V,\mathcal N_{hk})=0
\]
因为 $\mathcal N_{hk}$ 是松弛层。于是逆系统
$(\check C^\bullet(\mathfrak U,\mathcal F_k))_{k\in\mathbb N}$ 满足
(ML)；对每个 $i\geq0$，
$(\check H^i(\mathfrak U,\mathcal F_k))_{k\in\mathbb N}$ 也满足
(ML)，因为当 $i>0$ 时这是平凡的，而由于
$\check H^0(\mathfrak U,\mathcal F_k)=\Gamma(U,\mathcal F_k)$
（G, II, 5.2.2），由前述结果，当 $i=0$ 时也满足条件 (ML)。因此可应用
（\hyperref[0.13.2.3-fr]{13.2.3}），它表明
\[
\check H^i(\mathfrak U,\mathcal F)
=\varprojlim_k\check H^i(\mathfrak U,\mathcal F_k)=0
\quad\hbox{对所有 }i>0.
\]

\emph{b)} 现在转到一般情形。对每个 $k\in\mathbb N$，考虑
$\mathcal F_k$ 的由松弛层组成的典范消解
$\mathcal C^\bullet(X,\mathcal F_k)
=(\mathcal C^i(X,\mathcal F_k))_{i\geq0}$（G, II, 4.3）。对每个
$i\geq0$，显然
$(\mathcal C^i(X,\mathcal F_k))_{k\in\mathbb N}$ 是松弛层的一个逆系统；
我们证明它满足 \emph{a)} 的条件。
\oldpage[0\textsubscript{III}]{69}
事实上，若当 $h\leq k$ 时 $\mathcal N_{hk}$ 是 $u_{hk}$ 的核，则由 (iii)，
序列
\[
0\longrightarrow\mathcal N_{hk}\longrightarrow\mathcal F_k
\longrightarrow\mathcal F_h\longrightarrow0
\]
正合，而我们的断言来自函子
$\mathcal A\mapsto\mathcal C^i(X,\mathcal A)$ 的正合性
（G, II, 4.3）。令
$\mathcal G^i=\varprojlim_k\mathcal C^i(X,\mathcal F_k)$；于是由
\emph{a)}，当 $j>0$ 且 $i\geq0$ 时有
$H^j(X,\mathcal G^i)=0$。我们来证明
$\mathcal G^\bullet=(\mathcal G^i)_{i\geq0}$ 是层 $\mathcal F$ 的一个消解；
由于当 $j>0$ 时 $H^j(X,\mathcal G^i)=0$，上同调
$H^\bullet(X,\mathcal F)$ 将等于
$H^\bullet(\Gamma(X,\mathcal G^\bullet))$（G, II, 4.7.1）。

显然，对正合列
\[
0\longrightarrow\mathcal F_k
\longrightarrow\mathcal C^0(X,\mathcal F_k)
\longrightarrow\mathcal C^1(X,\mathcal F_k)\longrightarrow\cdots
\]
取逆极限，得到阿贝尔群层的复形
\[
0\longrightarrow\mathcal F\longrightarrow\mathcal G^0
\longrightarrow\mathcal G^1\longrightarrow\cdots.
\]
为了证明我们的断言，只须证明当 $i>0$ 时
$\mathcal H^i(\mathcal G^\bullet)=0$。这个层由预层
$U\mapsto H^i(\Gamma(U,\mathcal G^\bullet))$ 生成（G, II, 4.1）；而复形
$\Gamma(U,\mathcal G^\bullet)$ 是阿贝尔群复形逆系统
\[
(\Gamma(U,\mathcal C^\bullet(X,\mathcal F_k)))_{k\in\mathbb N}
\]
的逆极限（0\textsubscript{I}, 3.2.6）。我们已在 \emph{a)} 中看到，对每个
$i\geq0$，映射
$\Gamma(U,\mathcal C^i(X,\mathcal F_k))\to
\Gamma(U,\mathcal C^i(X,\mathcal F_h))$ $(h\leq k)$ 都是满射；另一方面，
$H^i(U,\mathcal F_k)=
H^i(\Gamma(U,\mathcal C^\bullet(X,\mathcal F_k)))$
因为 $\mathcal C^\bullet(U,\mathcal F_k|U)$ 是由
$\mathcal C^\bullet(X,\mathcal F_k)$ 在 $U$ 上所诱导的典范消解；由假设
(i)，对每个 $U\in\mathfrak B$，可以把
（\hyperref[0.13.2.3-fr]{13.2.3}）应用于复形逆系统
$(\Gamma(U,\mathcal C^\bullet(X,\mathcal F_k)))_{k\in\mathbb N}$，于是
\[
H^i(\Gamma(U,\mathcal G^\bullet))
=\varprojlim_k H^i(U,\mathcal F_k)
\quad\hbox{对所有 }i\geq0.
\]
此时，假设 (ii) 根据定义恰好证明：当 $i>0$ 时，诸层
$\mathcal H^i(\mathcal G^\bullet)$ 都为零。

因此，对每个 $i\geq0$ 有
\[
H^i(X,\mathcal F)=H^i(\Gamma(X,\mathcal G^\bullet))
\quad\hbox{且}\quad
\Gamma(X,\mathcal G^\bullet)
=\varprojlim_k\Gamma(X,\mathcal C^\bullet(X,\mathcal F_k)).
\]
我们刚刚指出，所有映射
$\Gamma(X,\mathcal C^i(X,\mathcal F_k))\to
\Gamma(X,\mathcal C^i(X,\mathcal F_h))$ $(h\leq k)$ 都是满射；故结论再次来自
（\hyperref[0.13.2.3-fr]{13.2.3}）。
\end{proof}

\begin{env}[13.3.2]
\label{0.13.3.2-fr}
\emph{说明} (13.3.2)。
\begin{enumerate}
\item[\emph{(i)}]
陈述（\hyperref[0.13.3.1-fr]{13.3.1}）只有在 $i>0$ 时才有实质内容，
因为当 $i=0$ 时，无须任何假设，$h_i$ 总是同构
（0\textsubscript{I}, 3.2.6）。
\item[\emph{(ii)}]
特别地，若对每个 $k$、每个 $i>0$ 以及每个 $U\in\mathfrak B$ 都有
$H^i(U,\mathcal F_k)=0$，并且对 $U\in\mathfrak B$，诸映射
$\Gamma(U,\mathcal F_k)\to\Gamma(U,\mathcal F_h)$ 都是满射，则
（\hyperref[0.13.3.1-fr]{13.3.1}）的条件 (i) 和 (ii) 都成立。这是应用
（\hyperref[0.13.3.1-fr]{13.3.1}）最常见的情形。
\end{enumerate}
\end{env}
\subsection{Mittag-Leffler 条件与逆系统的伴随分次对象}
\label{subsection:0.13.4-fr}
\label{0.13.4-fr}

\begin{env}[13.4.1]
\label{0.13.4.1-fr}
设 $\mathbf A=(A_k,u_{kh})_{k\in\mathbb Z}$ 是阿贝尔范畴 $C$ 中的一个逆系统；
若存在 $k_0$，使得当 $k<k_0$ 时有 $A_k=0$，则称它\emph{下有界}。
在每个 $A_k$ 上，以如下公式定义一个\emph{滤过}
$(\mathrm F^p(A_k))_{p\in\mathbb Z}$：
\begin{equation}
\label{0.13.4.1.1-fr}
\tag{13.4.1.1}
\mathrm F^p(A_k)=
\begin{cases}
\operatorname{Ker}(A_k\longrightarrow A_{p-1})&\text{当 }p\leq k+1,\\
0&\text{当 }p\geq k+1.
\end{cases}
\end{equation}
\oldpage[0\textsubscript{III}]{70}
因此，由假设有
$\mathrm F^{k_0}(A_k)=A_k$ 和 $\mathrm F^{k+1}(A_k)=0$；换句话说，
所考虑的滤过是\emph{有限的}
（\hyperref[0.11.1.3-fr]{11.1.3}）。这个滤过的伴随分次对象因而为
\[
\operatorname{gr}^p(A_k)=
\operatorname{Ker}(A_k\longrightarrow A_{p-1})
/\operatorname{Ker}(A_k\longrightarrow A_p)
\]
所以，$\operatorname{gr}^p(A_k)$ 同构于
$\operatorname{Ker}(A_k\to A_{p-1})$ 在 $A_k\to A_p$ 下的像；由定义
逆系统的诸态射的传递性，有
\begin{equation}
\label{0.13.4.1.2-fr}
\tag{13.4.1.2}
\operatorname{gr}^p(A_k)=
\operatorname{Ker}(A_p\longrightarrow A_{p-1})
\cap\operatorname{Im}(A_k\longrightarrow A_p)
\end{equation}
又由于按（\hyperref[0.13.4.1.1-fr]{13.4.1.1}）有
$\operatorname{Ker}(A_p\to A_{p-1})=\operatorname{gr}^p(A_p)$，故也有
\begin{equation}
\label{0.13.4.1.3-fr}
\tag{13.4.1.3}
\operatorname{gr}^p(A_k)=
\operatorname{gr}^p(A_p)\cap\operatorname{Im}(A_k\longrightarrow A_p).
\end{equation}
此外，前面的定义还表明，当 $k\leq h$ 时有
\[
u_{kh}(\mathrm F^p(A_h))\subset\mathrm F^p(A_k)
\]
从而对每个 $p\in\mathbb Z$，诸 $\operatorname{gr}^p(u_{kh})$ 定义一个
\emph{逆系统} $(\operatorname{gr}^p(A_k))_{k\in\mathbb Z}$。
\end{env}

\begin{env}[13.4.2]
\label{0.13.4.2-fr}
若当 $k$ 足够大时诸态射 $A_{k+1}\to A_k$ 都是\emph{同构}，则称逆系统
$\mathbf A$ \emph{本质常值}。若诸态射
$A_i\to A_j$ $(j\leq i)$ 都是\emph{满射}，则称逆系统 $\mathbf A$
是\emph{严格的}。当 $\mathbf A$ 严格时，由
（\hyperref[0.13.4.1.3-fr]{13.4.1.3}）可知，当 $p\leq k\leq h$ 时，
典范态射 $\operatorname{gr}^p(A_h)\to\operatorname{gr}^p(A_k)$ 是一个
\emph{同构}；换句话说，逆系统
$(\operatorname{gr}^p(A_k))_{k\in\mathbb Z}$ 本质常值。此时，把诸
$\operatorname{gr}^p(A_p)$（对每个 $p$，都与
$\varprojlim_k\operatorname{gr}^p(A_k)$ 等同）所成的序列记作
$\operatorname{gr}^\bullet(\mathbf A)$，并称为严格逆系统
$\mathbf A=(A_k)$ 的\emph{伴随分次对象}。

现在假设下有界逆系统 $\mathbf A$ 满足 (ML)。由
（\hyperref[0.13.1.2-fr]{13.1.2}）已知，“稳定像”对象所成的逆系统
$\mathbf A'=(A'_k)$ 是\emph{严格的}，而且也是下有界的；此时，与
$\mathbf A'$ 伴随的分次对象 $\operatorname{gr}^\bullet(\mathbf A')$
仍称为与 $\mathbf A$ \emph{伴随的}分次对象，并记作
$\operatorname{gr}^\bullet(\mathbf A)$。
\end{env}

\begin{proposition}[13.4.3]
\label{0.13.4.3-fr}
\emph{设 $\mathbf A=(A_k)_{k\in\mathbb Z}$ 是阿贝尔范畴 $C$ 中一个下有界
逆系统。下列两个条件等价：}
\begin{enumerate}
\item[\emph{a)}] \emph{$\mathbf A$ 满足条件 (ML)。}
\item[\emph{b)}] \emph{对每个 $p\in\mathbb Z$，逆系统
$(\operatorname{gr}^p(A_k))_{k\in\mathbb Z}$ 本质常值。}
\end{enumerate}
\emph{此外，当这些条件成立时，对每个 $p\in\mathbb Z$ 都有典范同构}
\begin{equation}
\label{0.13.4.3.1-fr}
\tag{13.4.3.1}
\operatorname{gr}^p(\mathbf A)\xrightarrow{\sim}
\varprojlim_k\operatorname{gr}^p(A_k).
\end{equation}
\end{proposition}

\begin{proof}
由（\hyperref[0.13.4.1.2-fr]{13.4.1.2}）立刻可知 \emph{a)} 蕴含
\emph{b)}；把同一个公式应用于逆系统 $\mathbf A'$（采用
（\hyperref[0.13.4.2-fr]{13.4.2}）的记号），根据定义就得到同构
（\hyperref[0.13.4.3.1-fr]{13.4.3.1}）。当 $k\leq h$ 时，令
$A_{kh}=\operatorname{Im}(A_h\to A_k)$；若 $k\leq h\leq j$，则有
$A_{kj}\subset A_{kh}\subset A_k$。在 $A_{kh}$ 上取由
$(\mathrm F^p(A_k))$ 诱导的滤过；由定义 $\mathbf A$ 的诸态射的传递性，
立即可验证这个滤过也是 $(\mathrm F^p(A_h))$ 的商滤过；因而
\begin{equation}
\label{0.13.4.3.2-fr}
\tag{13.4.3.2}
\operatorname{gr}^p(A_{kh})=
\operatorname{Im}\bigl(\operatorname{gr}^p(A_h)\longrightarrow
\operatorname{gr}^p(A_k)\bigr).
\end{equation}
\oldpage[0\textsubscript{III}]{71}
现假设 \emph{b)} 成立；对每个 $p\in\mathbb Z$ 和每个 $k\geq p$，存在
整数 $\mathrm L(p,k)$，使得当 $h\geq\mathrm L(p,k)$ 时，
（\hyperref[0.13.4.3.2-fr]{13.4.3.2}）的右端保持不变；由于当
$p<k_0$ 时 $\operatorname{gr}^p(A_k)=0$，所以当 $k$ 固定而 $p$ 遍历
整数 $\leq k$ 时，只有有限多个 $\mathrm L(p,k)$ 非零。令
$\mathrm L(k)=m$ 是这些整数中最大的一个；则对每个 $h\geq m$，有
$A_{kh}\subset A_{km}$，并且按 $m$ 的定义，典范单射
$A_{kh}\to A_{km}$ 定义一个\emph{同构}
$\operatorname{gr}^\bullet(A_{kh})\xrightarrow{\sim}
\operatorname{gr}^\bullet(A_{km})$；由于这些滤过都是\emph{有限的}，
可知上述单射本身也是双射（Bourbaki，\emph{Alg. comm.}，第 III 章，
§~2，n\textsuperscript{o}~8，定理 1），这就证明了 $\mathbf A$ 满足
(ML)。
\end{proof}

\begin{env}[13.4.4]
\label{0.13.4.4-fr}
假设逆极限 $A=\varprojlim_k A_k$ 在 $C$ 中存在。在
（\hyperref[0.13.4.1-fr]{13.4.1}）的定义中，此时可以用 $A$ 替换
$A_k$，而在 $A$ 上这样定义的滤过仍满足
\begin{equation}
\label{0.13.4.4.1-fr}
\tag{13.4.4.1}
\operatorname{gr}^p(A)=\operatorname{gr}^p(A_p)
\cap\operatorname{Im}(A\longrightarrow A_p).
\end{equation}
\end{env}

\begin{corollary}[13.4.5]
\label{0.13.4.5-fr}
\emph{假设 $C$ 是阿贝尔群范畴。若逆系统 $\mathbf A$ 满足 (ML)，并且
$A=\varprojlim_k A_k$，则对每个 $p\in\mathbb Z$ 都有典范同构}
\begin{equation}
\label{0.13.4.5.1-fr}
\tag{13.4.5.1}
\operatorname{gr}^p(A)\xrightarrow{\sim}
\varprojlim_k\operatorname{gr}^p(A_k).
\end{equation}
\end{corollary}
\begin{proof}
事实上，当 $k$ 足够大时有
$\operatorname{Im}(A_k\to A_p)=\operatorname{Im}(A\to A_p)$
（Bourbaki，\emph{Top. gén.}，第 II 章，第 3 版，§~3，
n\textsuperscript{o}~5，定理 1），结论即由
（\hyperref[0.13.4.1.3-fr]{13.4.1.3}）和
（\hyperref[0.13.4.4.1-fr]{13.4.4.1}）得到。
\end{proof}

\subsection{滤过复形的谱序列的逆极限}
\label{subsection:0.13.5-fr}
\begin{env}[13.5.1]
\label{0.13.5.1-fr}
设 $C$ 是一个阿贝尔范畴，$X^\bullet$ 是 $C$ 中对象组成的一个复形，
配备了一个\emph{滤过}
$(\mathrm F^p(X^\bullet))_{p\in\mathbb Z}$，并且存在某个指标 $p_0$ 使得
$\mathrm F^{p_0}(X^\bullet)=X^\bullet$。对每个 $k\in\mathbb Z$，考虑复形
$X_k^\bullet=X^\bullet/\mathrm F^{k+1}(X^\bullet)$；它自然配备由下列对象组成的
滤过：当 $p\leq k$ 时
$\mathrm F^p(X_k^\bullet)=
\mathrm F^p(X^\bullet)/\mathrm F^{k+1}(X^\bullet)$，而当
$p\geq k+1$ 时 $\mathrm F^p(X_k^\bullet)=0$。此外，还有典范态射
$X_{k+1}^\bullet\to X_k^\bullet$，它们使
$\mathbf X^\bullet=(X_k^\bullet)_{k\in\mathbb Z}$ 成为 $C$ 的对象所成的
\emph{滤过复形逆系统}。注意，这个逆系统是\emph{严格的}，并且当
$k<p_0$ 时 $X_k^\bullet=0$。
\end{env}
\begin{env}[13.5.2]
\label{0.13.5.2-fr}
更一般地，考虑 $C$ 的对象所成复形的一个\emph{严格}逆系统
$\mathbf X^\bullet=(X_k^\bullet)_{k\in\mathbb Z}$，并假设它\emph{下有界}；
在每个 $X_k^\bullet$ 上考虑（\hyperref[0.13.4.1-fr]{13.4.1}）中定义的滤过
（这里把问题置于 $C$ 的下有界复形所成的阿贝尔范畴中）。当 $(p\leq k)$ 时，
$X_k^\bullet\to X_p^\bullet$ 成为具有有限滤过的\emph{滤过复形}态射。
滤过复形的谱序列的函子性
（\hyperref[0.11.2.3-fr]{11.2.3}）表明，定义逆系统 $\mathbf X^\bullet$ 的态射
给出态射，从而使
$\mathrm E(\mathbf X^\bullet)=(\mathrm E(X_k^\bullet))$ 成为谱序列的
\emph{逆系统}。
\end{env}
\begin{lemma}[13.5.3]
\label{0.13.5.3-fr}
\emph{假设滤过复形逆系统
$\mathbf X^\bullet=(X_k^\bullet)_{k\in\mathbb Z}$ 如
（\hyperref[0.13.5.2-fr]{13.5.2}）中所得到。则：}
\begin{enumerate}
\item[\emph{a)}]\emph{当 $r\geq p-p_0$ 时，对每个 $k\in\mathbb Z$，有
$\mathrm B_r^{pq}(X_k^\bullet)=\mathrm B_\infty^{pq}(X_k^\bullet)$。}
\item[\emph{b)}]\emph{当 $k+1\leq p+r$ 时，有
$\mathrm Z_r^{pq}(X_k^\bullet)=\mathrm Z_\infty^{pq}(X_k^\bullet)$。}
\item[\emph{c)}]\emph{当 $k+1\geq p+r$ 时，对每个 $h\geq k$，态射
$\mathrm Z_r^{pq}(X_h^\bullet)\to\mathrm Z_r^{pq}(X_k^\bullet)$ 和
$\mathrm B_r^{pq}(X_h^\bullet)\to\mathrm B_r^{pq}(X_k^\bullet)$ 都是同构。}
\end{enumerate}
\end{lemma}
\oldpage[0\textsubscript{III}]{72}
这三个性质直接由（\hyperref[0.11.2.2-fr]{11.2.2}）的定义得到，同时注意到当
$p-r<p_0$ 时有 $\mathrm F^{p-r+1}(X_k^\bullet)=X_k^\bullet$。

\begin{env}[13.5.4]
\label{0.13.5.4-fr}
假设（\hyperref[0.13.5.3-fr]{13.5.3}）的假设成立。那么，对固定的
$p,q,r$（$r$ 有限），逆系统
\[
(\mathrm Z_r^{pq}(X_k^\bullet))_{k\in\mathbb Z},\quad
(\mathrm B_r^{pq}(X_k^\bullet))_{k\in\mathbb Z},\quad
(\mathrm E_r^{pq}(X_k^\bullet))_{k\in\mathbb Z}
\]
都是\emph{本质常值的}；分别以
$\mathrm Z_r^{pq}(\mathbf X^\bullet)$、
$\mathrm B_r^{pq}(\mathbf X^\bullet)$ 和
$\mathrm E_r^{pq}(\mathbf X^\bullet)=
\mathrm Z_r^{pq}(\mathbf X^\bullet)/ \mathrm B_r^{pq}(\mathbf X^\bullet)$
表示它们各自的逆极限。$\mathrm Z_r^{pq}(\mathbf X^\bullet)$ 和
$\mathrm B_r^{pq}(\mathbf X^\bullet)$ 可典范地识别为
$\mathrm E_1^{pq}(\mathbf X^\bullet)$ 的子对象。
由 $d_r^{pq}$ 的定义（M, XV, 1）可知，这些态射（相对于各 $X_k^\bullet$）
也本质常值，因而定义态射
\begin{equation}
\label{0.13.5.4.1-fr}
\tag{13.5.4.1}
d_r^{pq}:\mathrm E_r^{pq}(\mathbf X^\bullet)
\longrightarrow\mathrm E_r^{p+r,q-r+1}(\mathbf X^\bullet)
\end{equation}
满足 $d_r^{p+r,q-r+1}\circ d_r^{pq}=0$；此外，还有从
$\operatorname{Ker}(d_r^{pq})$ 到
$\mathrm Z_{r+1}^{pq}(\mathbf X^\bullet)/ \mathrm B_r^{pq}(\mathbf X^\bullet)$
的典范同构，以及从 $\operatorname{Im}(d_r^{pq})$ 到
$\mathrm B_{r+1}^{p+r,q-r+1}(\mathbf X^\bullet)/
\mathrm B_r^{p+r,q-r+1}(\mathbf X^\bullet)$ 的典范同构。
\end{env}

\begin{lemma}[13.5.5]
\label{0.13.5.5-fr}
\emph{在（\hyperref[0.13.5.3-fr]{13.5.3}）的假设下，当 $s\geq r>p-p_0$ 时，
存在典范单射}
\begin{equation}
\label{0.13.5.5.1-fr}
\tag{13.5.5.1}
i:\mathrm E_s^{pq}(\mathbf X^\bullet)
\longrightarrow\mathrm E_r^{pq}(\mathbf X^\bullet)
\end{equation}
\emph{以及典范同构}
\begin{equation}
\label{0.13.5.5.2-fr}
\tag{13.5.5.2}
j_r:\mathrm E_r^{pq}(\mathbf X^\bullet)
\xrightarrow{\sim}\mathrm E_\infty^{pq}(X_{p+r-1}^\bullet)
\end{equation}
\emph{使得图}
\begin{equation}
\label{0.13.5.5.3-fr}
\tag{13.5.5.3}
\xymatrix{
\mathrm E_s^{pq}(\mathbf X^\bullet)\ar[r]^{j_s}\ar[d]_i
&\mathrm E_\infty^{pq}(X_{p+s-1}^\bullet)\ar[d]\\
\mathrm E_r^{pq}(\mathbf X^\bullet)\ar[r]^{j_r}
&\mathrm E_\infty^{pq}(X_{p+r-1}^\bullet)
}
\end{equation}
\emph{交换}（右侧的竖直箭头来自态射
$X_{p+s-1}^\bullet\to X_{p+r-1}^\bullet$）。
\end{lemma}
\begin{proof}
$i$ 的存在性来自当 $r>p-p_0$ 时
$\mathrm B_r^{pq}(X_k^\bullet)=\mathrm B_\infty^{pq}(X_k^\bullet)$
（（\hyperref[0.13.5.3-fr]{13.5.3}），a）；当 $k+1\leq p+r$ 时有
$\mathrm Z_r^{pq}(X_k^\bullet)=\mathrm Z_\infty^{pq}(X_k^\bullet)$
（（\hyperref[0.13.5.3-fr]{13.5.3}），b），特别地
$\mathrm Z_r^{pq}(X_{p+r-1}^\bullet)=
\mathrm Z_\infty^{pq}(X_{p+r-1}^\bullet)$；另一方面，根据
（\hyperref[0.13.5.3-fr]{13.5.3}，c），
$\mathrm Z_r^{pq}(X_{p+r-1}^\bullet)$ 和
$\mathrm B_r^{pq}(X_{p+r-1}^\bullet)$ 可典范地识别为
$\mathrm Z_r^{pq}(\mathbf X^\bullet)$ 和
$\mathrm B_r^{pq}(\mathbf X^\bullet)$，于是得到 $j_r$ 的存在性及
（\hyperref[0.13.5.5.1-fr]{13.5.5.1}）的交换性。
\end{proof}

\begin{corollary}[13.5.6]
\label{0.13.5.6-fr}
\emph{在（\hyperref[0.13.5.3-fr]{13.5.3}）的假设下，若逆极限
$\varprojlim_r\mathrm E_r^{pq}(\mathbf X^\bullet)$ 或
$\varprojlim_k\mathrm E_\infty^{pq}(X_k^\bullet)$ 之一存在，则另一个也存在，
并且有典范同构}
\begin{equation}
\label{0.13.5.6.1-fr}
\tag{13.5.6.1}
j_\infty:\varprojlim_r\mathrm E_r^{pq}(\mathbf X^\bullet)
\xrightarrow{\sim}
\varprojlim_k\mathrm E_\infty^{pq}(X_k^\bullet).
\end{equation}
\emph{此外，逆系统
$(\mathrm E_r^{pq}(\mathbf X^\bullet))_{r\in\mathbb Z}$ 本质常值
（\hyperref[0.13.4.2-fr]{13.4.2}）的充要条件是逆系统
$(\mathrm E_\infty^{pq}(X_k^\bullet))_{k\in\mathbb Z}$ 本质常值。}
\end{corollary}

\begin{env}[13.5.7]
\label{0.13.5.7-fr}
记 $\mathrm B_\infty^{pq}(\mathbf X^\bullet)$ 和
$\mathrm Z_\infty^{pq}(\mathbf X^\bullet)$ 为
$\mathrm E_1^{pq}(\mathbf X^\bullet)$ 的子对象：前者等于当 $r>p-p_0$ 时的
$\mathrm B_r^{pq}(\mathbf X^\bullet)$，后者等于
$\inf_r\mathrm Z_r^{pq}(\mathbf X^\bullet)$（当后者存在时），从而
$\varprojlim_r\mathrm E_r^{pq}(\mathbf X^\bullet)$
\oldpage[0\textsubscript{III}]{73}
可典范地识别为
$\mathrm E_\infty^{pq}(\mathbf X^\bullet)=
\mathrm Z_\infty^{pq}(\mathbf X^\bullet)/
\mathrm B_\infty^{pq}(\mathbf X^\bullet)$。注意，
$\mathrm Z_r^{pq}(\mathbf X^\bullet)$、
$\mathrm B_r^{pq}(\mathbf X^\bullet)$、
$\mathrm E_r^{pq}(\mathbf X^\bullet)$ $(1\leq r\leq+\infty)$ 以及
$d_r^{pq}$ 都对受（\hyperref[0.13.5.5-fr]{13.5.5}）限制的逆系统
$\mathbf X^\bullet$ \emph{函子性地}依赖，并且
（\hyperref[0.13.5.5-fr]{13.5.5}）和
（\hyperref[0.13.5.6-fr]{13.5.6}）中定义的态射也是函子的。
\end{env}

\subsection{相对于带有限滤过对象的函子的谱序列}
\label{subsection:0.13.6-fr}

\begin{env}[13.6.1]
\label{0.13.6.1-fr}
设 $C$、$C'$ 是两个阿贝尔范畴，$\mathrm T:C\to C'$ 是一个加性协变
函子。假设 $C$ 的每个对象都同构于某个内射对象的子对象，从而右导函子
$\mathrm R^p\mathrm T$ $(p\geq0)$ 存在。
\end{env}

\begin{lemma}[13.6.2]
\label{0.13.6.2-fr}
\emph{设 $A$ 是 $C$ 的一个对象，带有有限滤过
$(\mathrm F^i(A))_{i\in\mathbb Z}$。存在 $A$ 的一个内射消解
$X^\bullet=(X^j)_{j\geq0}$，带有有限滤过
$(\mathrm F^i(X^\bullet))_{i\in\mathbb Z}$，使得关系
$\mathrm F^i(A)=A$（分别地，$\mathrm F^i(A)=0$）蕴含
$\mathrm F^i(X^\bullet)=X^\bullet$（分别地，$\mathrm F^i(X^\bullet)=0$），并且对每个
$i\in\mathbb Z$，$\mathrm F^i(X^\bullet)$ 都是
$\mathrm F^i(A)$ 的内射消解。}
\end{lemma}
\begin{proof}
设 $p$（分别地，$q>p$）为满足 $\mathrm F^p(A)=A$ 的最大指标
（分别地，使 $\mathrm F^q(A)=0$ 的最小指标）。对 $q-p$ 作归纳；当
$q-p=1$ 时引理显然成立。假设已经构造出 $A/\mathrm F^{q-1}(A)$ 的一个
具有所需性质的内射消解 ${X''}^\bullet$。考虑正合列
\[
0\to\mathrm F^{q-1}(A)\to A\to A/\mathrm F^{q-1}(A)\to0,
\]
取 $\mathrm F^{q-1}(A)$ 的一个内射消解 ${X'}^\bullet$，然后选取 $A$ 的一个
内射消解 $X^\bullet$，使得有正合列
\[
0\to {X'}^\bullet\to X^\bullet\to {X''}^\bullet\to0
\]
并且它与前面的正合列相容（M, V, 2.2）。显然，$X^\bullet$ 具有所需的性质。
\end{proof}

\begin{corollary}[13.6.3]
\label{0.13.6.3-fr}
\emph{设 $B$ 是 $C$ 的另一个对象，带有有限滤过
$(\mathrm F^i(B))_{i\in\mathbb Z}$；设 $s$ 为一个整数，并设 $u:A\to B$ 是满足
$u(\mathrm F^i(A))\subset\mathrm F^{i+s}(B)$（对每个 $i\in\mathbb Z$）的态射。
若 $Y^\bullet=(Y^j)_{j\geq0}$ 是 $B$ 的一个内射消解，带有滤过
$(\mathrm F^i(Y^\bullet))_{i\in\mathbb Z}$，且具有
（\hyperref[0.13.6.2-fr]{13.6.2}）中所述的性质，则存在一个与 $u$ 相容的态射
$v:X^\bullet\to Y^\bullet$，使得
$v(\mathrm F^i(X^\bullet))\subset\mathrm F^{i+s}(Y^\bullet)$（对每个
$i\in\mathbb Z$）。此外，任意两个这样的态射 $v$、$v'$ 都同伦。}
\end{corollary}

这立即由前面的构造及（M, V, 2.3）对 $q-p$ 作归纳得到。

\begin{env}[13.6.4]
\label{0.13.6.4-fr}
在（\hyperref[0.13.6.2-fr]{13.6.2}）的假设下，考虑 $C'$ 中的复形
$\mathrm T(X^\bullet)$。由于 $\mathrm F^i(X^\bullet)$ 是 $X^\bullet$ 的直和因子，
该复形显然由复形 $\mathrm T(\mathrm F^i(X^\bullet))$ 滤过。
由（\hyperref[0.13.6.3-fr]{13.6.3}），这个滤过复形的\emph{谱序列}在同构意义下
只依赖于带滤过对象 $A$。其收敛项是上同调 $\mathrm R^\bullet\mathrm T(A)$，
带有如下滤过：
\begin{equation}
\label{0.13.6.4.1-fr}
\tag{13.6.4.1}
\begin{split}
\mathrm F^p(\mathrm R^n\mathrm T(A))
&=\operatorname{Im}\bigl(\mathrm R^n\mathrm T(\mathrm F^p(A))
\longrightarrow\mathrm R^n\mathrm T(A)\bigr)\\
&=\operatorname{Ker}\bigl(\mathrm R^n\mathrm T(A)
\longrightarrow\mathrm R^n\mathrm T(A/\mathrm F^p(A))\bigr)
\end{split}
\end{equation}
（\hyperref[0.11.2.2-fr]{11.2.2}），其 $\mathrm E_1$ 项为
\begin{equation}
\label{0.13.6.4.2-fr}
\tag{13.6.4.2}
\mathrm E_1^{pq}=\mathrm R^{p+q}\mathrm T(\operatorname{gr}^p(A))
\end{equation}
其中 $\operatorname{gr}^p(A)$ 按通常定义为 $\mathrm F^p(A)/\mathrm F^{p+1}(A)$。
由（\hyperref[0.11.2.2-fr]{11.2.2}）显然，收敛项的滤过是有限的，并且对给定的
$p$、$q$，序列
\oldpage[0\textsubscript{III}]{74}
$\mathrm B_r^{pq}(A)=\mathrm B_r^{pq}(\mathrm T(X^\bullet))$
和
$\mathrm Z_r^{pq}(A)=\mathrm Z_r^{pq}(\mathrm T(X^\bullet))$
是稳定的。因此，前述谱序列是\emph{双正则的}
（\hyperref[0.11.1.3-fr]{11.1.3}）。记此谱序列为
$\mathrm E(A)=(\mathrm E_r^{pq}(A))$，并称其为\emph{函子 $\mathrm T$ 相对于带滤过对象
$A$ 的谱序列}。
\end{env}

\begin{env}[13.6.5]
\label{0.13.6.5-fr}
现在假设（\hyperref[0.13.6.3-fr]{13.6.3}）的假设成立，并沿用其中的记号。
由于 $\mathrm F^i(X^\bullet)$（分别地，$\mathrm F^i(Y^\bullet)$）是
$X^\bullet$（分别地，$Y^\bullet$）的直和因子，对每个 $i\in\mathbb Z$ 有
\[
(\mathrm T v)(\mathrm T(\mathrm F^i(X^\bullet)))
\subset\mathrm T(\mathrm F^{i+s}(Y^\bullet))
\]
由（\hyperref[0.11.2.2-fr]{11.2.2}）的定义可知，对 $1\leq r\leq+\infty$，
$\mathrm T v$ 给出态射
$\mathrm B_r^{pq}(\mathrm T(X^\bullet))\to
\mathrm B_r^{p+s,q-s}(\mathrm T(Y^\bullet))$ 
以及态射
$\mathrm Z_r^{pq}(\mathrm T(X^\bullet))\to
\mathrm Z_r^{p+s,q-s}(\mathrm T(Y^\bullet))$，
从而给出态射
\[
w_r:\mathrm E_r^{pq}(A)\to\mathrm E_r^{p+s,q-s}(B)~;
\]
同样，在收敛项上有态射
$u_n:\mathrm R^n\mathrm T(A)\to\mathrm R^n\mathrm T(B)$ 使得
$u_n(\mathrm F^p(\mathrm R^n\mathrm T(A)))\subset
\mathrm F^{p+s}(\mathrm R^n\mathrm T(B))$。

由 $d_r^{pq}$ 的定义（M, XV, 1）还可知下列图表交换：
\[
\xymatrix{
\mathrm E_r^{pq}(A)\ar[r]^{d_r^{pq}}\ar[d]_{w_r}
&\mathrm E_r^{p+r,q-r+1}(A)\ar[d]^{w_r}\\
\mathrm E_r^{p+s,q-s}(B)\ar[r]_{d_r^{p+s,q-s}}
&\mathrm E_r^{p+r+s,q-r-s+1}(B)
}
\]
因此，对于同构 $\alpha_r^{pq}$ 也得到类似的交换图；其具体形式留给读者写出。
最后（\emph{loc. cit.}），对于收敛项也有交换图
\[
\xymatrix{
\mathrm E_\infty^{pq}(A)\ar[r]^{\beta^{pq}}\ar[d]_{w_\infty}
&\operatorname{gr}^p(\mathrm R^{p+q}\mathrm T(A)) \ar[d]^{u_{p+q}}\\
\mathrm E_\infty^{p+s,q-s}(B)\ar[r]_{\beta^{p+s,q-s}}
&\operatorname{gr}^{p+s}(\mathrm R^{p+q}\mathrm T(B)).
}
\]
\end{env}

\begin{env}[13.6.6]
\label{0.13.6.6-fr}
特别地，设存在一个带滤过
$(\mathrm F^i(S))_{i\in\mathbb Z}$ 的环 $S$，以及环同态
\begin{equation}
\label{0.13.6.6.1-fr}
\tag{13.6.6.1}
h:S\to\operatorname{Hom}_C(A,A)
\end{equation}
\oldpage[0\textsubscript{III}]{75}
使得对每个 $t\in\mathrm F^i(S)$ 及每一对 $i$、$j$，都有
$h_t(\mathrm F^j(A))\subset\mathrm F^{i+j}(A)$。
为简便起见，称此时的 $A$ 为滤过环 $S$ 上的\emph{滤过 $S$-$C$-模}。
通过取伴随分次对象，对于 $t\in\mathrm F^j(S)$，每个 $h_t$ 都在
$\operatorname{gr}^\bullet(A)$ 上定义一个齐次次数为 $j$ 的分次自同态
$\overline h_t$；此外，该态射只依赖于 $t$ 在 $\operatorname{gr}^j(S)$ 中的类，
因而得到分次环同态
\[
\overline h:\operatorname{gr}^\bullet(S)\to
\operatorname{Hom}_C^g(\operatorname{gr}^\bullet(A),
\operatorname{gr}^\bullet(A))
\]
右端是 $\operatorname{gr}^\bullet(A)$ 的分次自同态环。称
$\operatorname{gr}^\bullet(A)$ 带有一个
\emph{$\operatorname{gr}^\bullet(S)$-$C$-分次模结构}。

由（\hyperref[0.13.6.5-fr]{13.6.5}），对 $1\leq r\leq+\infty$，每个
$\overline t\in\operatorname{gr}^j(S)$ 在下列双分次对象
$(\mathrm B_r^{pq}(A))_{(p,q)\in\mathbb Z\times\mathbb Z}$、
$(\mathrm Z_r^{pq}(A))_{(p,q)\in\mathbb Z\times\mathbb Z}$ 和
$\mathrm E_r(A)=(\mathrm E_r^{pq}(A))_{(p,q)\in\mathbb Z\times\mathbb Z}$ 上
规范地定义次数为 $(s,-s)$ 的双分次自同态：
在 $\mathrm E_r(A)$ 中（$r$ 有限时），此自同态与 $d_r^{pq}$ 所定义的双分次
自同态交换。这些自同态相对于 $\operatorname{gr}^\bullet(S)$ 以及所考虑的
双分次对象中的加法，满足通常的结合律和分配律。因此，为简便起见，称这些对象为
\emph{$\operatorname{gr}^\bullet(S)$-$C'$-双分次模}；显然，
$\alpha_r^{pq}$ 给出保持这种结构的同构。

对每个整数 $n$，记 $\mathrm B_r^{(n)}(A)$（分别地，
$\mathrm Z_r^{(n)}(A)$、$\mathrm E_r^{(n)}(A)$）为
$\mathrm B_r^\bullet(A)$（分别地，$\mathrm Z_r^\bullet(A)$、
$\mathrm E_r^\bullet(A)$）的分次子对象，由满足 $p+q=n$ 的
$\mathrm B_r^{pq}(A)$（分别地，$\mathrm Z_r^{pq}(A)$、$\mathrm E_r^{pq}(A)$）组成
（$1\leq r\leq+\infty$）。显然，它们是
\emph{$\operatorname{gr}^\bullet(S)$-$C'$-分次模}。
最后，每个 $\overline t\in\operatorname{gr}^j(S)$ 对每个 $n$ 在分次对象
$\operatorname{gr}^\bullet(\mathrm R^n\mathrm T(A))$ 上定义次数为 $j$ 的分次自同态，
于是该对象带有 $\operatorname{gr}^\bullet(S)$-$C'$-分次模结构；当 $p+q=n$ 时，
$\beta^{pq}$ 给出
$\mathrm E^{(n)}(A)$ 到
$\operatorname{gr}^\bullet(\mathrm R^n\mathrm T(A))$ 的同构，并保持上述结构。

注意，当 $C'$ 是阿贝尔群范畴时，$S$-$C'$-模结构（分别地，
分次或双分次 $\operatorname{gr}^\bullet(S)$-$C'$-模结构）不过就是通常的
$S$-模结构（分别地，分次或双分次 $\operatorname{gr}^\bullet(S)$-模结构）。
\end{env}

\subsection{逆极限自变量的导函子}
\label{subsection:0.13.7-fr}

\begin{env}[13.7.1]
\label{0.13.7.1-fr}
设 $C$、$C'$ 为两个阿贝尔范畴，并假设 $C$ 的每个对象都同构于 $C$ 中某个内射对象的子对象；设
$\mathrm T:C\to C'$ 为加性协变函子。考虑 $C$ 中的严格逆系统
$\mathbf A=(A_k)_{k\in\mathbb Z}$，它是\emph{下有界的}；更精确地，假设
$A_k=0$ 对 $k<k_0$ 成立。我们按典范方式在每个 $A_k$ 上赋予滤过
$(\mathrm F^p(A_k))_{p\in\mathbb Z}$，其由公式（\hyperref[0.13.4.1.1-fr]{13.4.1.1}）给出；由于这是一个严格逆系统，典范态射
\begin{equation}
\label{0.13.7.1.1-fr}
\tag{13.7.1.1}
\mathrm F^i(A_h)/\mathrm F^j(A_h)\to
\mathrm F^i(A_k)/\mathrm F^j(A_k)\qquad(h\geq k)
\end{equation}
在 $i\leq j\leq k+1$ 时是同构。还要回忆，对每个 $k$ 都有
$\mathrm F^{k_0}(A_k)=A_k$ 以及 $\mathrm F^{k+1}(A_k)=0$。
\end{env}

\begin{env}[13.7.2]
\label{0.13.7.2-fr}
现在，对每个 $k$ 构造 $A_k$ 的内射消解
$X_k^\bullet=(X_k^j)_{j\geq0}$，使其具有（\hyperref[0.13.6.2-fr]{13.6.2}）的性质。典范态射
$A_{k+1}\to A_k$ 通过（\hyperref[0.13.6.3-fr]{13.6.3}）使我们能够对每个 $k$ 定义一个与滤过相容的复形态射
$X_{k+1}^\bullet\to X_k^\bullet$，从而使
$\mathbf X^\bullet=(X_k^\bullet)_{k\in\mathbb Z}$ 成为复形的逆系统。
\oldpage[0\textsubscript{III}]{76}
此外，我们可以假设这个逆系统是严格的。事实上，根据
（\hyperref[0.13.7.1.1-fr]{13.7.1.1}），$A_k$ 同构于
$A_{k+1}/\mathrm F^{k+1}(A_{k+1})$；因此在构造 $X_{k+1}^\bullet$ 时，可以取
$A_{k+1}/\mathrm F^{k+1}(A_{k+1})$ 的内射消解等于 $X_k^\bullet$。于是由（M, V, 2.3）可知，复形态射
$X_{k+1}^\bullet\to X_k^\bullet$ 可以构造成使其在取商后诱导同构
\[
X_{k+1}^\bullet/\mathrm F^{k+1}(X_{k+1}^\bullet)
\xrightarrow{\sim}X_k^\bullet
\]
，并且保持滤过；这正是（\hyperref[0.13.5.1-fr]{13.5.1}）的条件。
\end{env}

\begin{env}[13.7.3]
\label{0.13.7.3-fr}
按构造，复形的逆系统 $\mathrm T(\mathbf X^\bullet)$，即 $\mathrm T(X_k^\bullet)$ 的逆系统，满足（\hyperref[0.13.5.3-fr]{13.5.3}）的假设。因此（\hyperref[0.13.5.4-fr]{13.5.4}）、（\hyperref[0.13.5.5-fr]{13.5.5}）和（\hyperref[0.13.5.6-fr]{13.5.6}）的结果适用于谱序列 $\mathrm E(\mathrm T(X_k^\bullet))=\mathrm E(A_k)$；当 $1\leq r\leq+\infty$ 时，我们用 $\mathrm E_r^{pq}(\mathbf A)$ 代替
$\mathrm E_r^{pq}(\mathrm T(\mathbf X^\bullet))$（当 $r=+\infty$ 时参见（\hyperref[0.13.5.7-fr]{13.5.7}），类似记号也采用同一约定。特别要注意，由（\hyperref[0.13.6.4.2-fr]{13.6.4.2}）以及系统
$(\operatorname{gr}^p(A_k))$ 本质常值这一事实，有
\begin{equation}
\label{0.13.7.3.1-fr}
\tag{13.7.3.1}
\mathrm E_1^{pq}(\mathbf A)=
\mathrm R^{p+q}\mathrm T(\operatorname{gr}^p(\mathbf A))
\end{equation}
这些结果与（\hyperref[0.13.4.3-fr]{13.4.3}）一起给出下面的命题；该命题最初由 Shih Weishu 用另一种（未发表的）方法证明：
\end{env}

\begin{proposition}[13.7.4]
\label{0.13.7.4-fr}
\emph{（Shih）。--- 设 $n$ 为整数。以下两个条件等价：}
\begin{enumerate}
\item[\emph{a)}] \emph{对于满足 $p+q=n$ 的每一对 $(p,q)$，逆系统
$(\mathrm E_r^{pq}(\mathbf A))_{r\geq2}$ 是本质常值的。}
\item[\emph{b)}] \emph{逆系统
$\mathrm R^n\mathrm T(\mathbf A)=
(\mathrm R^n\mathrm T(A_k))_{k\in\mathbb Z}$ 满足（ML）。}
\end{enumerate}
\emph{此外，当这些条件满足时，有一个典范同构}
\begin{equation}
\label{0.13.7.4.1-fr}
\tag{13.7.4.1}
\operatorname{gr}^p(\mathrm R^n\mathrm T(\mathbf A))
\xrightarrow{\sim}\mathrm E_\infty^{p,n-p}(\mathbf A)
\qquad\emph{对每个 $p\in\mathbb Z$。}
\end{equation}
\end{proposition}
\begin{proof}
事实上，根据（\hyperref[0.13.5.6-fr]{13.5.6}），条件\emph{a)}等价于说，当 $p+q=n$ 时，逆系统
$(\mathrm E_\infty^{pq}(A_k))_{k\in\mathbb Z}$ 是本质常值的。另一方面，
$\operatorname{gr}^p(\mathrm R^n\mathrm T(A_k))$ 典范同构于
$\mathrm E_\infty^{p,n-p}(A_k)$，因此由（\hyperref[0.13.4.3-fr]{13.4.3}）可知，\emph{a)} 与\emph{b)}等价；同构
（\hyperref[0.13.7.4.1-fr]{13.7.4.1}）不过就是将（\hyperref[0.13.5.6.1-fr]{13.5.6.1}）应用于这里的情形。
\end{proof}

\begin{corollary}[13.7.5]
\label{0.13.7.5-fr}
\emph{设 $\mathcal F=(\mathcal F_k)_{k\in\mathbb N}$ 是阿贝尔群层的一个逆系统，满足（\hyperref[0.13.3.1-fr]{13.3.1}）的\emph{（i）}、\emph{（ii）}和\emph{（iii）}条件，并且
并设
$\mathcal F=\varprojlim_k\mathcal F_k$。假设对于函子
$\mathcal G\mapsto\Gamma(X,\mathcal G)$，逆系统
$(\mathrm E_r^{pq}(\mathcal F))_{r\in\mathbb Z}$ 在满足 $p+q=n$ 或 $p+q=n+1$ 的每一对 $(p,q)$ 上都是本质常值的。考虑在 $H^{n+1}(X,\mathcal F)$ 上由
$\mathrm F^p(H^{n+1}(X,\mathcal F))=
\operatorname{Ker}(H^{n+1}(X,\mathcal F)\to
H^{n+1}(X,\mathcal F_{p-1}))$ 定义的滤过。于是有典范同构}
\begin{equation}
\label{0.13.7.5.1-fr}
\tag{13.7.5.1}
\operatorname{gr}^p(H^{n+1}(X,\mathcal F))
\xrightarrow{\sim}\mathrm E_\infty^{p,n-p+1}(\mathcal F)
\qquad\emph{对每个 $p\in\mathbb Z$。}
\end{equation}
\end{corollary}
\begin{proof}
\oldpage[0\textsubscript{III}]{77}
由将（\hyperref[0.13.7.4-fr]{13.7.4}）应用于 $C$ 为 $X$ 上阿贝尔群层范畴、$C'$ 为阿贝尔群范畴且 $\mathrm T=\Gamma$ 的情形，得到一个典范同构
$\operatorname{gr}^p(\mathrm R^{n+1}\Gamma(\mathcal F))
\xrightarrow{\sim}\mathrm E_\infty^{p,n-p+1}(\mathcal F)$ 对每个
$p\in\mathbb Z$。另一方面，根据
（\hyperref[0.13.7.4-fr]{13.7.4}），逆系统
$(H^n(X,\mathcal F_k))_{k\in\mathbb Z}$ 满足（ML），由
（\hyperref[0.13.3.1-fr]{13.3.1}）可推出典范同构
\begin{equation}
\label{0.13.7.5.1b-fr}
\tag{13.7.5.1}
H^{n+1}(X,\mathcal F)\xrightarrow{\sim}
\varprojlim_k H^{n+1}(X,\mathcal F_k).
\end{equation}

由于逆系统 $\mathrm R^{n+1}\Gamma(\mathcal F)$ 根据
（\hyperref[0.13.7.4-fr]{13.7.4}）满足（ML），有典范同构
\[
\operatorname{gr}^p(\mathrm R^{n+1}\Gamma(\mathcal F))
\xrightarrow{\sim}
\varprojlim_k\operatorname{gr}^p(H^{n+1}(X,\mathcal F_k))
\]
（\hyperref[0.13.4.3-fr]{13.4.3}），以及典范同构
\[
\varprojlim_k\operatorname{gr}^p(H^{n+1}(X,\mathcal F_k))
\xrightarrow{\sim}
\operatorname{gr}^p(\varprojlim_k H^{n+1}(X,\mathcal F_k))
\]
（\hyperref[0.13.4.5-fr]{13.4.5}）。因此只需考察同构
（\hyperref[0.13.7.5.1-fr]{13.7.5.1}）与两边滤过的相容性；而这立即由定义以及下图的交换性得到：
\[
\xymatrix{
H^{n+1}(X,\mathcal F)\ar[r]^{\sim}\ar[dr]
&\varprojlim_k H^{n+1}(X,\mathcal F_k)\ar[d]\\
&H^{n+1}(X,\mathcal F_{p-1})
}
\]
对每个 $p$ 均成立。
\end{proof}

\begin{env}[13.7.6]
\label{0.13.7.6-fr}
设 $S$ 是带有滤过
$(\mathrm F^i(S))_{i\in\mathbb Z}$ 的环，满足 $\mathrm F^0(S)=S$（因而
$\operatorname{gr}^i(S)=0$ 对 $i<0$）。假设每个 $A_k$（带有
（\hyperref[0.13.7.1-fr]{13.7.1}）中定义的滤过）都是一个滤过的
$S$-$C$-模（\hyperref[0.13.6.6-fr]{13.6.6}），且当 $k\leq h$ 时，态射 $A_h\to A_k$ 都是滤过的
$S$-$C$-模结构的态射；为简便起见，我们称 $\mathbf A$ 为一个\emph{滤过的
$S$-$C$-模逆系统}。于是立即可见，对于 $k\leq h$，$1\leq r\leq+\infty$，态射
$\mathrm B_r^{pq}(A_h)\to\mathrm B_r^{pq}(A_k)$ 以及
$\mathrm Z_r^{pq}(A_h)\to\mathrm Z_r^{pq}(A_k)$ 是
$\operatorname{gr}^\bullet(S)$-$C'$-双分次模结构的态射
（\hyperref[0.13.6.5-fr]{13.6.5}）；并且对有限的 $r$，族
$(\mathrm Z_r^{pq}(\mathbf A))$、$(\mathrm B_r^{pq}(\mathbf A))$ 和
$(\mathrm E_r^{pq}(\mathbf A))$ 是
$\operatorname{gr}^\bullet(S)$-$C'$-双分次模，前两个是
$(\mathrm E_1^{pq}(\mathbf A))$ 的子模。我们仍记
$\mathrm Z_r^{(n)}(\mathbf A)$、$\mathrm B_r^{(n)}(\mathbf A)$、
$\mathrm E_r^{(n)}(\mathbf A)$ 为前述对象中只取满足 $p+q=n$ 的项所得的相应分次子对象；它们是
$\operatorname{gr}^\bullet(S)$-$C'$-分次模。

当逆系统 $(\mathrm E_r^{pq}(\mathbf A))_{r\in\mathbb Z}$ 本质常值时，$(\mathrm E_\infty^{pq}(\mathbf A))$ 也是
$\operatorname{gr}^\bullet(S)$-$C'$-双分次模，而每个
$\mathrm E_\infty^{(n)}(\mathbf A)$ 都是
$\operatorname{gr}^\bullet(S)$-$C'$-分次模。此外，同构
$\beta^{p,n-p}:\mathrm E_\infty^{p,n-p}(A_k)\xrightarrow{\sim}
\operatorname{gr}^p(\mathrm R^n\mathrm T(A_k))$ 对每个
$k$ 构成一个 $\operatorname{gr}^\bullet(S)$-$C'$-分次模结构下的同构，将
$\mathrm E_\infty^{(n)}(A_k)$ 映到
$\operatorname{gr}^\bullet(\mathrm R^n\mathrm T(A_k))$；在上述条件下，
$\beta^{p,n-p}:\mathrm E_\infty^{(n)}(\mathbf A)\xrightarrow{\sim}
\varprojlim_k\operatorname{gr}^\bullet(\mathrm R^n\mathrm T(A_k))$
也保持这些结构；显然，典范同构
$\operatorname{gr}^\bullet(\mathrm R^n\mathrm T(\mathbf A))
\xrightarrow{\sim}
\varprojlim_k\operatorname{gr}^\bullet(\mathrm R^n\mathrm T(A_k))$，
也保持这些结构。因此，（\hyperref[0.13.7.4.1-fr]{13.7.4.1}）中的同构构成一个
$\operatorname{gr}^\bullet(S)$-$C'$-分次模结构下的同构。
\end{env}

\begin{proposition}[13.7.7]
\label{0.13.7.7-fr}
设 $S$ 为 Noether $\mathfrak J$-进制环。假设 $C$ 是一个阿贝尔范畴，其中每个对象都同构于某个内射对象的子对象，并设
$\mathrm T$ 是从 $C$ 到阿贝尔群范畴的共变加性函子。设
$\mathbf A=(A_k)_{k\in\mathbb Z}$ 为一个
\oldpage[0\textsubscript{III}]{78}
下有界的严格滤过 $S$-$C$-模逆系统（$S$ 上取 $\mathfrak J$-进滤过）。假设对于给定整数
$n$，下述条件成立：
\begin{itemize}
\item[$(\mathrm F_n)$] （\hyperref[0.13.7.3.1-fr]{13.7.3.1}）中的
$\operatorname{gr}^\bullet(S)$-分次模
\[
\mathrm E_1^{(m)}(\mathbf A)=
\bigl(\mathrm R^m\mathrm T(\operatorname{gr}^p(\mathbf A))\bigr)_{p\in\mathbb Z}
\]
在 $m=n$ 和 $m=n+1$ 时都是有限型的。
\end{itemize}
在这些条件下：
\begin{enumerate}
\item[\rm(i)] 逆系统
$(\mathrm R^n\mathrm T(A_k))_{k\in\mathbb Z}$ 和
$(\mathrm R^{n+1}\mathrm T(A_k))_{k\in\mathbb Z}$ 满足（ML）。
\item[\rm(ii)] 若令
\[
\mathrm R'^{\,n}\mathrm T(\mathbf A)=
\varprojlim_k\mathrm R^n\mathrm T(A_k),
\]
则 $\mathrm R'^{\,n}\mathrm T(\mathbf A)$ 是有限型 $S$-模。
\item[\rm(iii)] 在 $\mathrm R'^{\,n}\mathrm T(\mathbf A)$ 上由
\[
\mathrm F^p(\mathrm R'^{\,n}\mathrm T(\mathbf A))=
\operatorname{Ker}\bigl(\mathrm R'^{\,n}\mathrm T(\mathbf A)
\longrightarrow\mathrm R^n\mathrm T(A_{p-1})\bigr)
\quad(p\in\mathbb Z)
\]
定义的滤过是 $\mathfrak J$-良滤过（也就是说，对每个 $p$ 都有
$\mathfrak J\mathrm F^p(\mathrm R'^{\,n}\mathrm T(\mathbf A))
\subset\mathrm F^{p+1}(\mathrm R'^{\,n}\mathrm T(\mathbf A))$，并且当 $p$ 充分大时两端相等）。特别地，这一滤过在
$\mathrm R'^{\,n}\mathrm T(\mathbf A)$ 上定义的拓扑与 $\mathfrak J$-进拓扑相同。
\item[\rm(iv)] 当 $p+q=n$ 或 $p+q=n+1$ 时，逆系统
$(\mathrm E_r^{pq}(\mathbf A))_{r\in\mathbb Z}$ 本质常值；因而
$\mathrm E_\infty^{pq}(\mathbf A)$ 已由（\hyperref[0.13.5.7-fr]{13.5.7}）定义，并且有
$\operatorname{gr}^\bullet(S)$-分次模的典范同构
\begin{equation}
\label{0.13.7.7.1-fr}
\tag{13.7.7.1}
\operatorname{gr}^p(\mathrm R'^{\,n}\mathrm T(\mathbf A))
\xrightarrow{\sim}\mathrm E_\infty^{p,n-p}(\mathbf A)
\quad(p\in\mathbb Z).
\end{equation}
\end{enumerate}
\end{proposition}

\begin{proof}
应当指出，同构（\hyperref[0.13.7.7.1-fr]{13.7.7.1}）使我们可以略用记号，把逆系统
$\mathrm R^n\mathrm T(A_k)$ 的逆极限
$\mathrm R'^{\,n}\mathrm T(\mathbf A)$ 也记作 $\mathrm R^n\mathrm T(\mathbf A)$；这里要考虑同构
（\hyperref[0.13.7.4.1-fr]{13.7.4.1}）。

由于分次环 $\operatorname{gr}^\bullet(S)$ 是 Noether 环
（Bourbaki，\emph{Alg. comm.}，第 III 章，\S~2，第
n\textsuperscript{o}~9 节，定理 2 的推论 5），
$\mathrm E_1^{(m)}(\mathbf A)$ 的
$\operatorname{gr}^\bullet(S)$-分次子模
$\mathrm B_r^{(m)}(\mathbf A)$ 所成的递增序列在 $m=n$ 和 $m=n+1$ 时都是稳定的
（\hyperref[0.13.6.6-fr]{13.6.6}）；因此（\hyperref[0.11.1.10-fr]{11.1.10}）的条件 b）成立。
由此，（\hyperref[0.13.7.4-fr]{13.7.4}）的条件 a）对 $n$ 和 $n+1$ 都成立，这已经证明了（i）。
此外，同构（\hyperref[0.13.7.4.1-fr]{13.7.4.1}）（并考虑
（\hyperref[0.13.7.6-fr]{13.7.6}）中的说明）表明，
$\operatorname{gr}^\bullet(\mathrm R^n\mathrm T(\mathbf A))$ 是一个
$\operatorname{gr}^\bullet(S)$-分次模，并同构于
\[
\mathrm E_\infty^{(n)}(\mathbf A)=
\mathrm Z_\infty^{(n)}(\mathbf A)/\mathrm B_\infty^{(n)}(\mathbf A)~;
\]
由于 $\mathrm Z_\infty^{(n)}(\mathbf A)$ 是
$\mathrm E_1^{(n)}(\mathbf A)$ 的子模，它是有限型的，因而
$\mathrm E_\infty^{(n)}(\mathbf A)$ 也是有限型的。此外，对于滤过
$(\mathrm F^p(\mathrm R'^{\,n}\mathrm T(\mathbf A)))$，由
（\hyperref[0.13.4.5-fr]{13.4.5}）可知，
$\operatorname{gr}^\bullet(\mathrm R^n\mathrm T(\mathbf A))$ 与
$\operatorname{gr}^\bullet(\mathrm R'^{\,n}\mathrm T(\mathbf A))$ 是同构的
$\operatorname{gr}^\bullet(S)$-分次模，这证明了（iv）。最后，断言（ii）和（iii）将由上述结果及下述引理推出：
\end{proof}

\begin{lemma}[13.7.7.2]
\label{0.13.7.7.2-fr}
设 $S$ 为 Noether $\mathfrak J$-进制环，$M$ 为带有余离散滤过
$(\mathrm F^p(M))_{p\in\mathbb Z}$ 的 $S$-模，并且
$\mathfrak J\mathrm F^p(M)\subset\mathrm F^{p+1}(M)$
（这表示 $M$ 是以 $\mathfrak J$-进滤过滤过的环 $S$ 上的滤过模）。再假设
$M$ 对由滤过 $(\mathrm F^p(M))$ 定义的拓扑是分离的。则以下条件等价：
\begin{enumerate}
\item[\rm a)] $M$ 是有限型 $S$-模，且 $(\mathrm F^p(M))$ 是
$\mathfrak J$-良滤过。
\item[\rm b)] $\operatorname{gr}^\bullet(M)$ 是有限型
$\operatorname{gr}^\bullet(S)$-模。
\item[\rm c)] 各 $\operatorname{gr}^p(M)$ 都是有限型 $S$-模，并且当
$p$ 充分大时，典范同态
\begin{equation}
\label{0.13.7.7.3-fr}
\tag{13.7.7.3}
\mathfrak J\otimes_S\operatorname{gr}^p(M)\longrightarrow
\operatorname{gr}^{p+1}(M)
\end{equation}
\oldpage[0\textsubscript{III}]{79}
（它们由 $\mathfrak J\otimes_S\mathrm F^p(M)\to\mathrm F^{p+1}(M)$ 诱导；这里还要考虑到复合同态
\[
\mathfrak J\otimes_S\mathrm F^{p+1}(M)\longrightarrow
\mathfrak J\otimes_S\mathrm F^p(M)\longrightarrow\mathrm F^{p+1}(M)
\]
的像为
$\mathfrak J\mathrm F^{p+1}(M)\subset\mathrm F^{p+2}(M)$）都是满射。
\end{enumerate}
证明参见 Bourbaki，\emph{Alg. comm.}，第 III 章，\S~3，第
n\textsuperscript{o}~1 节，命题 3。
\end{lemma}

\begin{env}[13.7.7.4]
\label{0.13.7.7.4-fr}
为了应用引理（\hyperref[0.13.7.7.2-fr]{13.7.7.2}），只须指出，所考察滤过在
$\mathrm R'^{\,n}\mathrm T(\mathbf A)$ 上定义的拓扑使
$\mathrm R'^{\,n}\mathrm T(\mathbf A)$ 成为分离且完备的 $S$-模；这个拓扑就是离散群
$\mathrm R^n\mathrm T(A_k)$ 的逆极限拓扑。另一方面，若 $A_k=0$ 对
$k<k_0$ 成立，则也有 $\mathrm R^n\mathrm T(A_k)=0$ 对 $k<k_0$ 成立，因而
$\mathrm F^{k_0}(\mathrm R'^{\,n}\mathrm T(\mathbf A))=
\mathrm R'^{\,n}\mathrm T(\mathbf A)$；所以这里确实满足该引理的应用条件。
\end{env}

\begin{corollary}[13.7.8]
\label{0.13.7.8-fr}
若假设 $(\mathrm F_n)$ 成立，则对每个元素 $f\in S$ 都有典范同构
\begin{equation}
\label{0.13.7.8.1-fr}
\tag{13.7.8.1}
\varprojlim_k\bigl((\mathrm R^n\mathrm T(A_k))_f\bigr)
\xrightarrow{\sim}\mathrm R^n\mathrm T(\mathbf A)\otimes_S S_{\{f\}}.
\end{equation}
\end{corollary}

\begin{proof}
事实上，$\mathrm R^n\mathrm T(\mathbf A)$ 是有限型 $S$-模，而
$S_{\{f\}}$ 是 Noether 进制 $S$-代数
（\hyperref[0.7.6.11-fr]{0\textsubscript{I}, 7.6.11}），即
$S_f$ 关于 $\mathfrak J$-预进制拓扑的分离完备化
（\hyperref[0.7.6.2-fr]{0\textsubscript{I}, 7.6.2}）。由
（\hyperref[0.7.7.8-fr]{0\textsubscript{I}, 7.7.8}）和
（\hyperref[0.7.7.1-fr]{0\textsubscript{I}, 7.7.1}）可知，
$\mathrm R^n\mathrm T(\mathbf A)\otimes_S S_{\{f\}}$ 同构于
$\mathrm R^n\mathrm T(\mathbf A)\otimes_S S_f$ 对
$\mathfrak J$-预进制拓扑的分离完备化；此拓扑的一组 $0$ 的基本邻域系为
$(\mathfrak J^p\mathrm R^n\mathrm T(\mathbf A))\otimes_S S_f$，因此
$\mathrm F^p(\mathrm R^n\mathrm T(\mathbf A))\otimes_S S_f$ 也构成这样一组基本邻域系；后者是典范映射
$(\mathrm R^n\mathrm T(\mathbf A))_f\to
(\mathrm R^n\mathrm T(A_{p-1}))_f$ 的核，因而与
$\mathrm R^n\mathrm T(\mathbf A)\otimes_S S_f$ 相关联的分离群可视为
$\varprojlim_k((\mathrm R^n\mathrm T(A_k))_f)$ 的一个子群 $G$。但是，逆系统
$((\mathrm R^n\mathrm T(A_k))_f)$ 显然满足条件（ML），而
$(\mathrm R^n\mathrm T(\mathbf A))_f$ 在每个
$(\mathrm R^n\mathrm T(A_k))_f$ 中的像，等于当 $h\geq k$ 充分大时各
$(\mathrm R^n\mathrm T(A_h))_f$ 的共同像。由此立即得到，$G$ 在
$\varprojlim_k((\mathrm R^n\mathrm T(A_k))_f)$ 中\emph{处处稠密}；由于后一个群是分离且完备的，推论得证。
\end{proof}

\begin{flushright}
\textit{（待续。）}
\end{flushright}

\clearpage
\thispagestyle{empty}
\null
\clearpage
\thispagestyle{plain}
\markboth{相干层的上同调研究}{第三章}
\oldpage[III]{81}
\phantomsection
\label{chapter:III-fr}

\begin{center}
{\large 第三章\par}
\vspace{1.5em}
{\Large\bfseries 相干层的上同调研究\par}
\vspace{1.8em}
{\bfseries 概要\par}
\end{center}

\medskip
\begin{tabular}{@{}r@{\ }p{0.82\textwidth}@{}}
\S~1. & 仿射概形的上同调。\\
\S~2. & 射影态射的上同调研究。\\
\S~3. & 固有态射的有限性定理。\\
\S~4. & 固有态射的基本定理；应用。\\
\S~5. & 相干代数层的一个存在性定理。\\
\S~6. & 局部与整体 Tor 和 Ext 函子；Künneth 公式。\\
\S~7. & 模的协变上同调函子的基变换研究。\\
\S~8. & 射影丛上的对偶定理。\\
\S~9. & 相对上同调与局部上同调；局部对偶。\\
\S~10. & 射影上同调与局部上同调之间的关系。沿除子的形式完备化方法。\\
\S~11. & 整体与局部 Picard 群\footnotemark。
\end{tabular}

\medskip
本章给出相干代数层的上同调的各项基本定理，但从剩余理论导出的那些定理
（对偶定理）除外；后者将成为以后某章的主题。在前一类定理中，实质上有
六个基本定理，分别构成本章前六节的主题。这些结果在下文中将成为不可或缺
的工具，即使所讨论的问题本身并非上同调性质的问题也是如此；读者从
\S~4 起便会看到最初的例子。\S~7 给出一些更具技术性的结果，但它们在应用中
经常使用。最后，在 \S\S~8 至 11 中，我们将发展若干与相干层对偶有关、
对应用特别重要，并且可以在一般剩余理论之前叙述的结果。

\footnotetext{第四章不依赖于 \S\S~8 至 11，并且很可能会先于这些章节出版。}

\oldpage[III]{82}
\S\S~1 和 2 的内容归功于 J.-P. Serre；读者可以看出，我们只须遵循
（FAC）的论述。\S\S~8 和 9 同样受到（FAC）的启发（不过，不同语境所要求
的转换不那么显然）。最后，正如我们在引言中所说，\S~4 应当看作以现代语言
表述 Zariski“全纯函数理论”中的基本“不变性定理”。

最后指出，第 3.4 节的结果（以及（0，13.4 至 13.7）中的预备命题）在
第三章后文中不会使用，因而初次阅读时可以略去。

\setcounter{section}{0}
\section{仿射概形的上同调}
\label{section:III.1-fr}

\subsection{外代数复形的回顾}
\label{subsection:III.1.1-fr}

\begin{env}[1.1.1]
\label{III.1.1.1-fr}
设 $A$ 是一个环，$\mathbf f=(f_i)_{1\leq i\leq r}$ 是 $A$ 中 $r$ 个
元素组成的一个系。与 $\mathbf f$ 对应的\emph{外代数复形
$K_\bullet(\mathbf f)$}是如下定义的链复形（G，I，2.2）：分次
$A$-模 $K_\bullet(\mathbf f)$ 等于按通常方式分次的\emph{外代数
$\bigwedge(A^r)$}，其边缘算子是把 $\mathbf f$ 视为对偶模
$(A^r)^\vee$ 的元素而作的\emph{内乘 $i_{\mathbf f}$}。回忆，
$i_{\mathbf f}$ 是 $\bigwedge(A^r)$ 上次数为 $-1$ 的
\emph{反导子}；若 $(\mathbf e_i)_{1\leq i\leq r}$ 是 $A^r$ 的典范基，
则 $i_{\mathbf f}(\mathbf e_i)=f_i$。条件
$i_{\mathbf f}\circ i_{\mathbf f}=0$ 立即可验。

一个等价定义如下：对每个 $i$，考虑链复形 $K_\bullet(f_i)$，定义为
$K_0(f_i)=K_1(f_i)=A$，而当 $n\neq0,1$ 时 $K_n(f_i)=0$；其边缘算子
由 $d_1:A\to A$ 是\emph{乘以 $f_i$}这一条件定义。于是取
$K_\bullet(\mathbf f)$ 为带有全次数的\emph{张量积
$K_\bullet(f_1)\otimes K_\bullet(f_2)\otimes\cdots\otimes
K_\bullet(f_r)$}（G，I，2.7）；它与上面定义的复形同构，立即可验。
\end{env}

\begin{env}[1.1.2]
\label{III.1.1.2-fr}
对任意 $A$-模 $M$，定义\emph{链复形}
\begin{equation}
\label{III.1.1.2.1-fr}
\tag{1.1.2.1}
K_\bullet(\mathbf f,M)=K_\bullet(\mathbf f)\otimes_A M
\end{equation}
以及\emph{上链复形}（G，I，2.2）
\begin{equation}
\label{III.1.1.2.2-fr}
\tag{1.1.2.2}
K^\bullet(\mathbf f,M)=\operatorname{Hom}_A(K_\bullet(\mathbf f),M).
\end{equation}
若 $g$ 是后一复形的一个 $k$-上链，并令
\[
g(i_1,\ldots,i_k)=g(\mathbf e_{i_1}\wedge\cdots\wedge\mathbf e_{i_k}),
\]
则 $g$ 可等同于从 $[1,r]^k$ 到 $M$ 的一个\emph{交错}映射；由上述定义有
\begin{equation}
\label{III.1.1.2.3-fr}
\tag{1.1.2.3}
d^kg(i_1,i_2,\ldots,i_{k+1})=
\sum_{h=1}^{k+1}(-1)^{h-1}f_{i_h}
g(i_1,\ldots,\widehat{i_h},\ldots,i_{k+1}).
\end{equation}
\end{env}

\oldpage[III]{83}
\begin{env}[1.1.3]
\label{III.1.1.3-fr}
由上述复形，照常得到\emph{同调与上同调 $A$-模}（G，I，2.2）
\begin{align}
\label{III.1.1.3.1-fr}
\tag{1.1.3.1}
H_\bullet(\mathbf f,M)&=H_\bullet(K_\bullet(\mathbf f,M)),\\
\label{III.1.1.3.2-fr}
\tag{1.1.3.2}
H^\bullet(\mathbf f,M)&=H^\bullet(K^\bullet(\mathbf f,M)).
\end{align}
还可定义一个 $A$-同构
$K_i(\mathbf f,M)\xrightarrow{\sim}K^{r-i}(\mathbf f,M)$：把每个链
$z=\sum(\mathbf e_{i_1}\wedge\cdots\wedge\mathbf e_{i_k})\otimes
z_{i_1\ldots i_k}$ 对应到上链 $g_z$，其中
$g_z(j_1,\ldots,j_{r-k})=\varepsilon z_{i_1\ldots i_k}$；
这里，$(j_h)_{1\leq h\leq r-k}$ 是 $[1,r]$ 中严格递增序列
$(i_h)_{1\leq h\leq k}$ 的严格递增补序列，并且
$\varepsilon=(-1)^\nu$，其中 $\nu$ 是 $[1,r]$ 的排列
$i_1,\ldots,i_k,j_1,\ldots,j_{r-k}$ 的逆序数。可以验证
$g_{dz}=d(g_z)$，从而得到同构
\begin{equation}
\label{III.1.1.3.3-fr}
\tag{1.1.3.3}
H^i(\mathbf f,M)\xrightarrow{\sim}H_{r-i}(\mathbf f,M)
\qquad\text{当 }0\leq i\leq r.
\end{equation}
在本章中，我们主要考察上同调模 $H^\bullet(\mathbf f,M)$。

对给定的 $\mathbf f$，立即可见（G，I，2.1），
$M\mapsto H^\bullet(\mathbf f,M)$ 是从 $A$-模范畴到分次 $A$-模范畴
的一个\emph{上同调函子}（T，II，2.1），在次数 $<0$ 和 $>r$ 时为零。
此外，
\begin{equation}
\label{III.1.1.3.4-fr}
\tag{1.1.3.4}
H^0(\mathbf f,M)=\operatorname{Hom}_A(A/(\mathbf f),M),
\end{equation}
其中 $(\mathbf f)$ 表示由 $f_1,\ldots,f_r$ 生成的 $A$ 的理想；
这立即由（\hyperref[III.1.1.2.3-fr]{1.1.2.3}）得到，并且显然
$H^0(\mathbf f,M)$ 等同于 $M$ 中\emph{被 $(\mathbf f)$ 零化}的子模。
同样，由（\hyperref[III.1.1.2.3-fr]{1.1.2.3}）有
\begin{equation}
\label{III.1.1.3.5-fr}
\tag{1.1.3.5}
H^r(\mathbf f,M)=M/\Bigl(\sum_{i=1}^r f_iM\Bigr)
=(A/(\mathbf f))\otimes_A M.
\end{equation}
我们将使用下面这个已知结果；为求完备，回忆其证明：
\end{env}

\begin{proposition}[1.1.4]
\label{III.1.1.4-fr}
设 $A$ 是一个环，$\mathbf f=(f_i)_{1\leq i\leq r}$ 是 $A$ 的有限元素族，
$M$ 是一个 $A$-模。若对 $1\leq i\leq r$，模
$M_{i-1}=M/(f_1M+\cdots+f_{i-1}M)$ 上的数乘映射
$z\mapsto f_i z$ 是单射，则当 $i\neq r$ 时
$H^i(\mathbf f,M)=0$。
\end{proposition}

\begin{proof}
由（\hyperref[III.1.1.3.3-fr]{1.1.3.3}），只须证明对所有 $i>0$，
$H_i(\mathbf f,M)=0$。对 $r$ 作归纳，$r=0$ 的情形显然。令
$\mathbf f'=(f_i)_{1\leq i\leq r-1}$；这一族满足命题中的条件，所以若令
$L_\bullet=K_\bullet(\mathbf f',M)$，则由归纳假设，对 $i>0$ 有
$H_i(L_\bullet)=0$，并由（\hyperref[III.1.1.3.3-fr]{1.1.3.3}）和
（\hyperref[III.1.1.3.5-fr]{1.1.3.5}）有
$H_0(L_\bullet)=M_{r-1}$。为简便起见，令
$K_\bullet=K_\bullet(f_r)=K_0\oplus K_1$，其中 $K_0=K_1=A$，
$d_1:K_1\to K_0$ 是乘以 $f_r$；由定义
（\hyperref[III.1.1.1-fr]{1.1.1}），
$K_\bullet(\mathbf f,M)=K_\bullet\otimes_A L_\bullet$。现有如下引理：

\begin{lemma}[1.1.4.1]
\label{III.1.1.4.1-fr}
设 $K_\bullet$ 是由自由 $A$-模组成、除次数 $0$ 和 $1$ 外均为零的
链复形。对任意由 $A$-模组成的链复形 $L_\bullet$，以及任意指标 $p$，
有正合列
\[
0\longrightarrow H_0(K_\bullet\otimes H_p(L_\bullet))
\longrightarrow H_p(K_\bullet\otimes L_\bullet)
\longrightarrow H_1(K_\bullet\otimes H_{p-1}(L_\bullet))
\longrightarrow0
\]
\end{lemma}

\oldpage[III]{84}
这是 Künneth 谱序列的低阶项正合列的一个特例
（M，XVII，5.2 a）和 G，I，5.5.2）；也可如下直接证明。把 $K_0$
和 $K_1$ 看成链复形（分别在次数 $\neq0$ 和 $\neq1$ 时为零），
便有复形的正合列
\[
0\longrightarrow K_0\otimes L_\bullet
\longrightarrow K_\bullet\otimes L_\bullet
\longrightarrow K_1\otimes L_\bullet\longrightarrow0
\]
对其应用同调正合列
\[
\cdots\longrightarrow H_{p+1}(K_1\otimes L_\bullet)
\xrightarrow{\partial}H_p(K_0\otimes L_\bullet)
\longrightarrow H_p(K_\bullet\otimes L_\bullet)
\longrightarrow H_p(K_1\otimes L_\bullet)
\xrightarrow{\partial}H_{p-1}(K_0\otimes L_\bullet)
\longrightarrow\cdots.
\]
但显然，对所有 $p$，
$H_p(K_0\otimes L_\bullet)=K_0\otimes H_p(L_\bullet)$ 且
$H_p(K_1\otimes L_\bullet)=K_1\otimes H_{p-1}(L_\bullet)$；
此外，立即可验，算子
$\partial:K_1\otimes H_p(L_\bullet)\to K_0\otimes H_p(L_\bullet)$
正是 $d_1\otimes1$。因此，引理由上述正合列以及
$H_0(K_\bullet\otimes H_p(L_\bullet))$ 和
$H_1(K_\bullet\otimes H_{p-1}(L_\bullet))$ 的定义得到。

引理既已证明，（\hyperref[III.1.1.4-fr]{1.1.4}）证明的余下部分立即得到：
归纳假设与引理（\hyperref[III.1.1.4.1-fr]{1.1.4.1}）给出，当
$p\geq2$ 时 $H_p(K_\bullet\otimes L_\bullet)=0$；此外，只要证明
$H_1(K_\bullet,H_0(L_\bullet))=0$，再由引理
（\hyperref[III.1.1.4.1-fr]{1.1.4.1}）即可得到
$H_1(K_\bullet\otimes L_\bullet)=0$。但按定义，
$H_1(K_\bullet,H_0(L_\bullet))$ 正是 $M_{r-1}$ 上数乘映射
$z\mapsto f_rz$ 的核；按假设此核为零，证明完成。
\end{proof}

\begin{env}[1.1.5]
\label{III.1.1.5-fr}
设 $\mathbf g=(g_i)_{1\leq i\leq r}$ 是 $A$ 中 $r$ 个元素组成的另一序列，并置
$\mathbf f\mathbf g=(f_ig_i)_{1\leq i\leq r}$。可以定义复形的典范同态
\begin{equation}
\label{III.1.1.5.1-fr}
\tag{1.1.5.1}
\varphi_{\mathbf g}:K_\bullet(\mathbf f\mathbf g)\longrightarrow
K_\bullet(\mathbf f)
\end{equation}
如下：它是 $A$-线性映射
$(x_1,\ldots,x_r)\mapsto(g_1x_1,\ldots,g_rx_r)$ 从 $A^r$ 到自身的映射在
外代数 $\bigwedge(A^r)$ 上的典范延拓。要验证所得映射确为复形同态，只须作如下
一般观察：若 $u:E\to F$ 是 $A$-线性映射，
$\mathbf x\in F^\vee$ 且 $\mathbf y={}^tu(\mathbf x)\in E^\vee$，则有公式
\begin{equation}
\label{III.1.1.5.2-fr}
\tag{1.1.5.2}
(\bigwedge u)\circ i_{\mathbf y}=i_{\mathbf x}\circ(\bigwedge u)~;
\end{equation}
事实上，两端都是 $\bigwedge F$ 的反导子，只须验证它们在 $F$ 上相同，而这立即由定义得到。

把 $K_\bullet(\mathbf f)$ 与诸 $K_\bullet(f_i)$ 的张量积等同起来时
（\hyperref[III.1.1.1-fr]{1.1.1}），$\varphi_{\mathbf g}$ 是诸
$\varphi_{g_i}$ 的张量积，其中 $\varphi_{g_i}$ 在次数 $0$ 上是恒等映射，在次数 $1$
上是乘以 $g_i$ 的映射。
\end{env}

\begin{env}[1.1.6]
\label{III.1.1.6-fr}
特别地，对满足 $0\leq n\leq m$ 的任意一对整数 $m,n$，有复形同态
\begin{equation}
\label{III.1.1.6.1-fr}
\tag{1.1.6.1}
\varphi_{\mathbf f^{m-n}}:K_\bullet(\mathbf f^m)\longrightarrow
K_\bullet(\mathbf f^n)
\end{equation}
从而有同态
\begin{align}
\label{III.1.1.6.2-fr}
\tag{1.1.6.2}
\varphi_{\mathbf f^{m-n}} &:K^\bullet(\mathbf f^n,M)\longrightarrow
K^\bullet(\mathbf f^m,M),\\
\label{III.1.1.6.3-fr}
\tag{1.1.6.3}
\varphi_{\mathbf f^{m-n}} &:H^\bullet(\mathbf f^n,M)\longrightarrow
H^\bullet(\mathbf f^m,M).
\end{align}

\oldpage[III]{85}
后面这些同态显然满足传递性条件
$\varphi_{\mathbf f^{m-p}}=
\varphi_{\mathbf f^{m-n}}\circ\varphi_{\mathbf f^{n-p}}$（$p\leq n\leq m$）；
因此它们定义了两个 $A$-模的\emph{归纳系}；置
\begin{align}
\label{III.1.1.6.4-fr}
\tag{1.1.6.4}
C^\bullet((\mathbf f),M)&=\varinjlim_n K^\bullet(\mathbf f^n,M),\\
\label{III.1.1.6.5-fr}
\tag{1.1.6.5}
H^\bullet((\mathbf f),M)&=H^\bullet(C^\bullet((\mathbf f),M))
=\varinjlim_n H^\bullet(\mathbf f^n,M),
\end{align}
最后一个等式源于取归纳极限与函子 $H^\bullet$ 可交换（G，I，2.1）。以后将看到
（\hyperref[III.1.4.3-fr]{1.4.3}），$H^\bullet((\mathbf f),M)$ 实际上只依赖于
$A$ 中的\emph{理想 $(\mathbf f)$}（甚至只依赖于 $A$ 上的
$(\mathbf f)$-预进制拓扑），这就说明了这些记号的合理性。

显然，$M\mapsto C^\bullet((\mathbf f),M)$ 是正合的 $A$-线性函子，而
$M\mapsto H^\bullet((\mathbf f),M)$ 是上同调函子。
\end{env}

\begin{env}[1.1.7]
\label{III.1.1.7-fr}
设 $\mathbf f=(f_i)\in A^r$，$\mathbf g=(g_i)\in A^r$；用
$e_{\mathbf g}$ 表示外代数 $\bigwedge(A^r)$ 中左乘向量
$\mathbf g\in A^r$ 的映射；已知有\emph{同伦公式}
\begin{equation}
\label{III.1.1.7.1-fr}
\tag{1.1.7.1}
i_{\mathbf f}e_{\mathbf g}+e_{\mathbf g}i_{\mathbf f}
=\langle\mathbf g,\mathbf f\rangle 1
\end{equation}
成立于 $A$-模 $\bigwedge(A^r)$ 中（$1$ 表示 $\bigwedge(A^r)$ 的恒等自同构）
\SourceCorrectionNote{法文印本及其外交式转录在此两处均写作 $A^r$；
EGA III 已发表勘误（III, 1.1.7）明确要求两处均改为 $\bigwedge(A^r)$。
本译采用该已发表校正，法文外交式来源保持不变。}；这一关系也意味着，
\emph{在复形 $K_\bullet(\mathbf f)$ 中}有
\begin{equation}
\label{III.1.1.7.2-fr}
\tag{1.1.7.2}
de_{\mathbf g}+e_{\mathbf g}d=\langle\mathbf g,\mathbf f\rangle 1.
\end{equation}
若理想 $(\mathbf f)$ 等于 $A$，则存在 $\mathbf g\in A^r$，使得
$\langle\mathbf g,\mathbf f\rangle=\sum_{i=1}^r g_if_i=1$。因而
（G，I，2.4）：
\end{env}

\begin{proposition}[1.1.8]
\label{III.1.1.8-fr}
设由诸 $f_i$ 生成的理想 $(\mathbf f)$ 等于 $A$。则复形
$K_\bullet(\mathbf f)$ 是同伦平凡的，因而对任意 $A$-模 $M$，复形
$K_\bullet(\mathbf f,M)$ 和 $K^\bullet(\mathbf f,M)$ 也都是同伦平凡的。
\end{proposition}

\begin{corollary}[1.1.9]
\label{III.1.1.9-fr}
若 $(\mathbf f)=A$，则对任意 $A$-模 $M$，有
$H^\bullet(\mathbf f,M)=0$ 和 $H^\bullet((\mathbf f),M)=0$。
\end{corollary}

\begin{proof}
事实上，此时对所有 $n$ 都有 $(\mathbf f^n)=A$。
\end{proof}

\begin{remark}[1.1.10]
\label{III.1.1.10-fr}
沿用上述记号，设 $X=\operatorname{Spec}(A)$，$Y$ 是由理想 $(\mathbf f)$ 定义的
$X$ 的闭子概形。我们将在第 \S~9 节证明，$H^\bullet((\mathbf f),M)$ 同构于
与 $Y$ 的闭子集反滤子 $\Phi$ 相对应的上同调
$H_Y^\bullet(X,\widetilde M)$（T，3.2）。我们还将证明，把命题
（\hyperref[III.1.2.3-fr]{1.2.3}）应用于 $X$ 和
$\mathcal F=\widetilde M$，便得到上同调正合列的一个特例
\[
\cdots\longrightarrow H_Y^p(X,\mathcal F)\longrightarrow H^p(X,\mathcal F)
\longrightarrow H^p(X-Y,\mathcal F)\longrightarrow H_Y^{p+1}(X,\mathcal F)
\longrightarrow\cdots.
\]
\end{remark}

\subsection{开覆盖的 Čech 上同调}
\label{subsection:III.1.2-fr}

\begin{notation}[1.2.1]
\label{III.1.2.1-fr}
本节采用如下记号：
\begin{itemize}
\item $X$ 是一个预概形；
\item $\mathcal F$ 是一个拟相干 $\mathcal O_X$-模；
\oldpage[III]{86}
\item $A=\Gamma(X,\mathcal O_X)$，$M=\Gamma(X,\mathcal F)$；
\item $\mathbf f=(f_i)_{1\leq i\leq r}$ 是 $A$ 中元素的一个有限系；
\item $U_i=X_{f_i}$ 是由满足 $f_i(x)\neq0$ 的 $x\in X$ 所成的开集
（\hyperref[0.5.5.2-fr]{0\textsubscript{I}，5.5.2}）；
\item $U=\bigcup_{i=1}^r U_i$；
\item $\mathfrak U$ 是 $U$ 的覆盖 $(U_i)_{1\leq i\leq r}$。
\end{itemize}
\end{notation}

\begin{env}[1.2.2]
\label{III.1.2.2-fr}
设 $X$ 或者是底空间为\emph{Noether} 空间的预概形，或者是底空间\emph{拟紧}的
\emph{概形}。于是已知（\hyperref[I.9.3.3-fr]{I, 9.3.3}）有
$\Gamma(U_i,\mathcal F)=M_{f_i}$。置
\[
U_{i_0i_1\ldots i_p}=\bigcap_{k=0}^p U_{i_k}
=X_{f_{i_0}f_{i_1}\cdots f_{i_p}}
\]
（\hyperref[0.5.5.3-fr]{0\textsubscript{I}，5.5.3}）；因而还有
\begin{equation}
\label{III.1.2.2.1-fr}
\tag{1.2.2.1}
\Gamma(U_{i_0i_1\ldots i_p},\mathcal F)=M_{f_{i_0}f_{i_1}\cdots f_{i_p}}.
\end{equation}
另一方面，（\hyperref[0.1.6.1-fr]{0\textsubscript{I}，1.6.1}）给出
$M_{f_{i_0}f_{i_1}\cdots f_{i_p}}$ 与归纳极限
$\varinjlim_n M_{i_0i_1\ldots i_p}^{(n)}$ 的等同，其中归纳系由
$M_{i_0i_1\ldots i_p}^{(n)}=M$ 组成；当 $m\leq n$ 时，同态
$\varphi_{nm}:M_{i_0i_1\ldots i_p}^{(m)}\to
M_{i_0i_1\ldots i_p}^{(n)}$ 是乘以
$(f_{i_0}f_{i_1}\cdots f_{i_p})^{n-m}$ 的映射。对所有 $n$，用
$C_n^p(M)$ 表示从 $[1,r]^{p+1}$ 到 $M$ 的\emph{交错}映射的集合；这些
$A$-模关于诸 $\varphi_{nm}$ 也组成一个归纳系。若
$C^p(\mathfrak U,\mathcal F)$ 是覆盖 $\mathfrak U$ 关于系数 $\mathcal F$
的\emph{交错} Čech $p$-上链群（G，II，5.1），则由上可写成
\begin{equation}
\label{III.1.2.2.2-fr}
\tag{1.2.2.2}
C^p(\mathfrak U,\mathcal F)=\varinjlim_n C_n^p(M).
\end{equation}
沿用（\hyperref[III.1.1.2-fr]{1.1.2}）的记号，$C_n^p(M)$ 与
$K^{p+1}(\mathbf f^n,M)$ 等同，而映射 $\varphi_{nm}$ 与
（\hyperref[III.1.1.6-fr]{1.1.6}）中定义的映射
$\varphi_{\mathbf f^{n-m}}$ 等同。因此，对每个 $p\geq0$，有关于
$\mathcal F$ 的典范函子同构
\begin{equation}
\label{III.1.2.2.3-fr}
\tag{1.2.2.3}
C^p(\mathfrak U,\mathcal F)\xrightarrow{\sim}C^{p+1}((\mathbf f),M).
\end{equation}
此外，公式（\hyperref[III.1.1.2.3-fr]{1.1.2.3}）和覆盖上同调的定义
（G，II，5.1）表明，诸同构
（\hyperref[III.1.2.2.3-fr]{1.2.2.3}）与余边缘算子相容。
\end{env}

\begin{proposition}[1.2.3]
\label{III.1.2.3-fr}
若 $X$ 是底空间为 Noether 空间的预概形，或是底空间拟紧的概形，则存在关于
$\mathcal F$ 的典范函子同构
\begin{equation}
\label{III.1.2.3.1-fr}
\tag{1.2.3.1}
H^p(\mathfrak U,\mathcal F)\xrightarrow{\sim}H^{p+1}((\mathbf f),M)
\qquad\text{对所有 }p\geq1.
\end{equation}
此外，有关于 $\mathcal F$ 的函子正合列
\begin{equation}
\label{III.1.2.3.2-fr}
\tag{1.2.3.2}
0\longrightarrow H^0((\mathbf f),M)\longrightarrow M
\longrightarrow H^0(\mathfrak U,\mathcal F)
\longrightarrow H^1((\mathbf f),M)\longrightarrow0.
\end{equation}
\end{proposition}

\oldpage[III]{87}
\begin{proof}
诸关系（\hyperref[III.1.2.3.1-fr]{1.2.3.1}）确实立即由
（\hyperref[III.1.2.2-fr]{1.2.2}）中的结论得到。另一方面，
$C^0(\mathfrak U,\mathcal F)=C^1((\mathbf f),M)$；因而
$H^0(\mathfrak U,\mathcal F)$ 与 $C^1((\mathbf f),M)$ 的 $1$-上闭链子群
等同；又因 $M=C^0((\mathbf f),M)$，正合列
（\hyperref[III.1.2.3.2-fr]{1.2.3.2}）无非是由上同调群
$H^0((\mathbf f),M)$ 和 $H^1((\mathbf f),M)$ 的定义所得的正合列。
\end{proof}

\begin{corollary}[1.2.4]
\label{III.1.2.4-fr}
设诸 $X_{f_i}$ 都拟紧，并且存在 $g_i\in\Gamma(U,\mathcal O_X)$，使得
$\sum_i g_i(f_i\vert U)=1\vert U$。则对任意拟相干
$(\mathcal O_X\vert U)$-模 $\mathcal F$，当 $p>0$ 时有
$H^p(\mathfrak U,\mathcal F)=0$
\SourceCorrectionNote{法文印本及其外交式转录在（III, 1.2.4）中依次写作
$\Gamma(U,\mathcal F)$、$\mathcal G$ 和 $H^p(\mathfrak U,\mathcal G)$；
EGA III 已发表勘误明确要求分别改为 $\Gamma(U,\mathcal O_X)$、$\mathcal F$
和 $H^p(\mathfrak U,\mathcal F)$。本译采用该已发表校正，法文外交式来源保持不变。}；若此外 $U=X$，则典范同态
（\hyperref[III.1.2.3.2-fr]{1.2.3.2}）
$M\to H^0(\mathfrak U,\mathcal F)$ 是双射。
\end{corollary}

\begin{proof}
按假设，诸 $U_i=X_{f_i}$ 都拟紧，故 $U$ 也拟紧，于是可归结到 $U=X$
的情形；此时假设蕴含对所有 $p\geq0$ 有
$H^p((\mathbf f),M)=0$（\hyperref[III.1.1.9-fr]{1.1.9}）。于是推论立即由
（\hyperref[III.1.2.3.1-fr]{1.2.3.1}）和
（\hyperref[III.1.2.3.2-fr]{1.2.3.2}）得到。
\end{proof}

应注意，由于 $H^0(\mathfrak U,\mathcal F)=H^0(U,\mathcal F)$
（G，II，5.2.2），这也把（I，1.3.7）作为特例重新证明了一遍。

\begin{remark}[1.2.5]
\label{III.1.2.5-fr}
设 $X$ 是\emph{仿射概形}；则诸 $U_i=X_{f_i}=D(f_i)$ 都是仿射开集，
诸 $U_{i_0i_1\ldots i_p}$ 也如此（但 $U$ 未必仿射）。在这种情形，函子
$\Gamma(X,\mathcal F)$ 和
$\Gamma(U_{i_0i_1\ldots i_p},\mathcal F)$ 关于 $\mathcal F$ 是正合的
（I，1.3.11）。若有拟相干 $\mathcal O_X$-模的正合列
$0\to\mathcal F'\to\mathcal F\to\mathcal F''\to0$，则复形列
\[
0\longrightarrow C^\bullet(\mathfrak U,\mathcal F')
\longrightarrow C^\bullet(\mathfrak U,\mathcal F)
\longrightarrow C^\bullet(\mathfrak U,\mathcal F'')
\longrightarrow0
\]
正合，因而给出上同调正合列
\[
\cdots\longrightarrow H^p(\mathfrak U,\mathcal F')
\longrightarrow H^p(\mathfrak U,\mathcal F)
\longrightarrow H^p(\mathfrak U,\mathcal F'')
\xrightarrow{\partial}H^{p+1}(\mathfrak U,\mathcal F')
\longrightarrow\cdots.
\]
另一方面，若置 $M'=\Gamma(X,\mathcal F')$、
$M''=\Gamma(X,\mathcal F'')$，则列
$0\to M'\to M\to M''\to0$ 正合；因为
$C^\bullet((\mathbf f),M)$ 关于 $M$ 是正合函子，所以还有上同调正合列
\[
\cdots\longrightarrow H^p((\mathbf f),M')
\longrightarrow H^p((\mathbf f),M)
\longrightarrow H^p((\mathbf f),M'')
\xrightarrow{\partial}H^{p+1}((\mathbf f),M')
\longrightarrow\cdots.
\]
既然图表
\[
\xymatrix{
0\ar[r]&C^\bullet(\mathfrak U,\mathcal F')\ar[r]\ar[d]
&C^\bullet(\mathfrak U,\mathcal F)\ar[r]\ar[d]
&C^\bullet(\mathfrak U,\mathcal F'')\ar[r]\ar[d]&0\\
0\ar[r]&C^\bullet((\mathbf f),M')\ar[r]
&C^\bullet((\mathbf f),M)\ar[r]
&C^\bullet((\mathbf f),M'')\ar[r]&0
}
\]
\oldpage[III]{88}
交换，便可推出对每个 $p$，图表
\begin{equation}
\label{III.1.2.5.1-fr}
\tag{1.2.5.1}
\xymatrix{
H^p(\mathfrak U,\mathcal F'')\ar[r]^{\partial}\ar[d]
&H^{p+1}(\mathfrak U,\mathcal F')\ar[d]\\
H^{p+1}((\mathbf f),M'')\ar[r]^{\partial}
&H^{p+2}((\mathbf f),M')
}
\end{equation}
都交换（G，I，2.1.1）。
\end{remark}

\subsection{仿射概形的上同调}
\label{subsection:III.1.3-fr}

\begin{theorem}[1.3.1]
\label{III.1.3.1-fr}
设 $X$ 是仿射概形。对任意拟相干 $\mathcal O_X$-模 $\mathcal F$，
当 $p>0$ 时有 $H^p(X,\mathcal F)=0$。
\end{theorem}

\begin{proof}
设 $\mathfrak U$ 是由仿射开集 $X_{f_i}=D(f_i)$（$1\leq i\leq r$）
组成的 $X$ 的有限覆盖；已知诸 $f_i$ 在
$A=\Gamma(X,\mathcal O_X)$ 中生成的理想等于 $A$。因而由
（\hyperref[III.1.2.4-fr]{1.2.4}）得到
$H^p(\mathfrak U,\mathcal F)=0$（$p>0$）。又因为 $X$ 有由仿射开集组成的
任意精细的有限覆盖（I，1.1.10），Čech 上同调的定义（G，II，5.8）表明
当 $p>0$ 时还有 $\check H^p(X,\mathcal F)=0$。但这同样适用于每个
$f\in A$ 所对应的预概形 $X_f$（I，1.3.6），故
$\check H^p(X_f,\mathcal F)=0$（$p>0$）。由于
$X_f\cap X_g=X_{fg}$，再由（G，II，5.9.2）推出对所有 $p>0$ 均有
$H^p(X,\mathcal F)=0$。
\end{proof}

\begin{corollary}[1.3.2]
\label{III.1.3.2-fr}
设 $Y$ 是预概形，$f:X\to Y$ 是仿射态射（II，1.6.1）。对任意拟相干
$\mathcal O_X$-模 $\mathcal F$，当 $q>0$ 时有
$R^qf_*(\mathcal F)=0$。
\end{corollary}

\begin{proof}
事实上，按定义，$R^qf_*(\mathcal F)$ 是由预层
$U\mapsto H^q(f^{-1}(U),\mathcal F)$ 所关联的 $\mathcal O_Y$-模，
其中 $U$ 遍历 $Y$ 的开集。但仿射开集构成 $Y$ 的一组基；对这样的开集
$U$，$f^{-1}(U)$ 是仿射的（II，1.3.2），故由
（\hyperref[III.1.3.1-fr]{1.3.1}）有
$H^q(f^{-1}(U),\mathcal F)=0$，这就证明了推论。
\end{proof}

\begin{corollary}[1.3.3]
\label{III.1.3.3-fr}
设 $Y$ 是预概形，$f:X\to Y$ 是仿射态射。对任意拟相干
$\mathcal O_X$-模 $\mathcal F$，典范同态
$H^p(Y,f_*(\mathcal F))\to H^p(X,\mathcal F)$（0，12.1.3.1）
对每个 $p$ 都是双射。
\end{corollary}

\begin{proof}
事实上，依（0，12.1.7），只需证明复合函子 $\Gamma f_*$ 的第二个谱序列的
边缘同态
$''E_2^{p0}=H^p(Y,f_*(\mathcal F))\to H^p(X,\mathcal F)$ 是双射。
但该谱序列的 $E_2$ 项由
$''E_2^{pq}=H^p(Y,R^qf_*(\mathcal F))$（G，II，4.17.1）给出；因而由
（\hyperref[III.1.3.2-fr]{1.3.2}）知，当 $q>0$ 时
$''E_2^{pq}=0$，于是谱序列退化；我们的断言随即由（0，11.1.6）得到。
\end{proof}

\begin{corollary}[1.3.4]
\label{III.1.3.4-fr}
设 $f:X\to Y$ 是仿射态射，$g:Y\to Z$ 是态射。
\oldpage[III]{89}
对任意拟相干 $\mathcal O_X$-模 $\mathcal F$，典范同态
$R^pg_*(f_*(\mathcal F))\to R^p(g\circ f)_*(\mathcal F)$（0，12.2.5.1）
对每个 $p$ 都是双射。
\end{corollary}

\begin{proof}
只需注意，由（\hyperref[III.1.3.3-fr]{1.3.3}），对 $Z$ 的每个仿射开集
$W$，典范同态
$H^p(g^{-1}(W),f_*(\mathcal F))\to
H^p(f^{-1}(g^{-1}(W)),\mathcal F)$ 都是双射；这表明，定义典范同态
$R^pg_*(f_*(\mathcal F))\to R^p(g\circ f)_*(\mathcal F)$ 的预层同态
是双射（0，12.2.5）。
\end{proof}

\subsection{对任意预概形上同调的应用}
\label{subsection:III.1.4-fr}

\begin{proposition}[1.4.1]
\label{III.1.4.1-fr}
设 $X$ 是概形，$\mathfrak U=(U_\alpha)$ 是由仿射开集组成的 $X$ 的覆盖。
对任意拟相干 $\mathcal O_X$-模 $\mathcal F$，上同调模
$H^\bullet(X,\mathcal F)$ 和 $H^\bullet(\mathfrak U,\mathcal F)$
（作为 $\Gamma(X,\mathcal O_X)$-模）典范同构。
\end{proposition}

\begin{proof}
事实上，由于 $X$ 是概形，覆盖 $\mathfrak U$ 中开集的每个有限交 $V$
都是仿射的（I，5.5.6），故由（\hyperref[III.1.3.1-fr]{1.3.1}），
当 $q\geq1$ 时有 $H^q(V,\mathcal F)=0$。命题于是由 Leray 定理
（G，II，5.4.1）得到。
\end{proof}

\begin{remark}[1.4.2]
\label{III.1.4.2-fr}
应注意，只要诸 $U_\alpha$ 的有限交都是仿射的，
（\hyperref[III.1.4.1-fr]{1.4.1}）的结论仍然成立，即使不必假设
$X$ 是概形。
\end{remark}

\begin{corollary}[1.4.3]
\label{III.1.4.3-fr}
设 $X$ 是底空间拟紧的概形，$A=\Gamma(X,\mathcal O_X)$，
$\mathbf f=(f_i)_{1\leq i\leq r}$ 是 $A$ 中元素的有限序列，使得诸
$X_{f_i}$（\hyperref[III.1.2.1-fr]{1.2.1} 的记号）都是仿射的。
于是，采用（\hyperref[III.1.2.1-fr]{1.2.1}）的记号，对任意拟相干
$\mathcal O_X$-模 $\mathcal F$，有关于 $\mathcal F$ 函子性的典范同构
\begin{equation}
\label{III.1.4.3.1-fr}
\tag{1.4.3.1}
H^q(U,\mathcal F)\xrightarrow{\sim}H^{q+1}((\mathbf f),M)
\qquad\text{对 }q\geq1
\end{equation}
以及关于 $\mathcal F$ 函子性的正合列
\begin{equation}
\label{III.1.4.3.2-fr}
\tag{1.4.3.2}
0\longrightarrow H^0((\mathbf f),M)\longrightarrow M
\longrightarrow H^0(U,\mathcal F)
\longrightarrow H^1((\mathbf f),M)\longrightarrow0.
\end{equation}
\end{corollary}

\begin{proof}
这立即由（\hyperref[III.1.4.1-fr]{1.4.1}）和
（\hyperref[III.1.2.3-fr]{1.2.3}）得到。
\end{proof}

\begin{env}[1.4.4]
\label{III.1.4.4-fr}
若 $X$ 是\emph{仿射概形}，则由（\hyperref[III.1.2.5-fr]{1.2.5}）和
（\hyperref[III.1.4.1-fr]{1.4.1}），对每个 $q\geq0$，与拟相干
$\mathcal O_X$-模的正合列
$0\to\mathcal F'\to\mathcal F\to\mathcal F''\to0$ 相对应的图表
\begin{equation}
\label{III.1.4.4.1-fr}
\tag{1.4.4.1}
\xymatrix{
H^q(U,\mathcal F'')\ar[r]^{\partial}\ar[d]
&H^{q+1}(U,\mathcal F')\ar[d]\\
H^{q+1}((\mathbf f),M'')\ar[r]^{\partial}
&H^{q+2}((\mathbf f),M')
}
\end{equation}
交换；这里采用（\hyperref[III.1.2.5-fr]{1.2.5}）的记号。
\end{env}

\oldpage[III]{90}
\begin{proposition}[1.4.5]
\label{III.1.4.5-fr}
设 $X$ 是拟紧概形，$\mathcal L$ 是可逆 $\mathcal O_X$-模，并考虑分次环
$A_\bullet=\Gamma_\bullet(\mathcal L)$（0\textsubscript{I}，5.4.6）；则
\[
H^\bullet(\mathcal F,\mathcal L)
=\bigoplus_{n\in\mathbb Z}H^\bullet(X,\mathcal F\otimes\mathcal L^{\otimes n})
\]
是分次 $A_\bullet$-模，并且对每个 $f\in A_n$，有
$(A_\bullet)_{(f)}$-模的典范同构
\begin{equation}
\label{III.1.4.5.1-fr}
\tag{1.4.5.1}
H^\bullet(X_f,\mathcal F)
\xrightarrow{\sim}(H^\bullet(\mathcal F,\mathcal L))_{(f)}.
\end{equation}
\end{proposition}

\begin{proof}
由于 $X$ 是拟紧概形，可以用同一个由仿射开集组成的有限覆盖
$\mathfrak U=(U_i)$ 来计算所有 $\mathcal O_X$-模
$\mathcal F\otimes\mathcal L^{\otimes n}$ 的上同调，并且对每个 $i$，限制
$\mathcal L\vert U_i$ 都同构于 $\mathcal O_X\vert U_i$
（\hyperref[III.1.4.1-fr]{1.4.1}）。于是立即可见，诸
$U_i\cap X_f$ 都是仿射开集（I，1.3.6），所以也可以用覆盖
$\mathfrak U\vert X_f=(U_i\cap X_f)$ 来计算
$H^\bullet(X_f,\mathcal F\otimes\mathcal L^{\otimes n})$
（\hyperref[III.1.4.1-fr]{1.4.1}）。显然，对每个 $f\in A_n$，乘以
$f$ 定义同态
\[
C^\bullet(\mathfrak U,\mathcal F\otimes\mathcal L^{\otimes m})
\longrightarrow
C^\bullet(\mathfrak U,\mathcal F\otimes\mathcal L^{\otimes(m+n)}),
\]
从而定义同态
\[
H^\bullet(\mathfrak U,\mathcal F\otimes\mathcal L^{\otimes m})
\longrightarrow
H^\bullet(\mathfrak U,\mathcal F\otimes\mathcal L^{\otimes(m+n)}),
\]
这证明了第一个断言。另一方面，给定 $f\in A_n$，由（I，9.3.2），
并计及（I，1.3.9，(ii)），有 $(A_\bullet)_{(f)}$-模复形的同构
\[
C^\bullet(\mathfrak U\vert X_f,\mathcal F)
\xrightarrow{\sim}
\left(C^\bullet\left(\mathfrak U,
\bigoplus_{n\in\mathbb Z}\mathcal F\otimes\mathcal L^{\otimes n}
\right)\right)_{(f)}
\]
对这两个复形取上同调，并记得函子 $M\mapsto M_{(f)}$ 在分次
$A_\bullet$-模范畴中正合，便得到同构
（\hyperref[III.1.4.5.1-fr]{1.4.5.1}）。
\end{proof}

\begin{corollary}[1.4.6]
\label{III.1.4.6-fr}
假设（\hyperref[III.1.4.5-fr]{1.4.5}）的条件成立，并且另有
$\mathcal L=\mathcal O_X$。若置 $A=\Gamma(X,\mathcal O_X)$，则对每个
$f\in A$，有 $A_f$-模的典范同构
\[
H^\bullet(X_f,\mathcal F)
\xrightarrow{\sim}(H^\bullet(X,\mathcal F))_f
\]
\end{corollary}

\begin{corollary}[1.4.7]
\label{III.1.4.7-fr}
设 $X$ 是拟紧概形，$f$ 是 $\Gamma(X,\mathcal O_X)$ 的一个元素。
\begin{enumerate}
\item[(i)] 假设开集 $X_f$ 是仿射的。则对任意拟相干 $\mathcal O_X$-模
$\mathcal F$、任意 $i>0$ 及任意 $\xi\in H^i(X,\mathcal F)$，存在整数
$n>0$，使得 $f^n\xi=0$。
\item[(ii)] 反之，假设 $X_f$ 拟紧，并且对 $\mathcal O_X$ 的每个拟相干
理想层 $\mathcal J$ 及每个 $\zeta\in H^1(X,\mathcal J)$，都存在
$n>0$ 使得 $f^n\zeta=0$。则 $X_f$ 是仿射的。
\end{enumerate}
\end{corollary}

\begin{proof}
\begin{enumerate}
\item[(i)] 若 $X_f$ 仿射，则对每个 $i>0$ 有
$H^i(X_f,\mathcal F)=0$（\hyperref[III.1.3.1-fr]{1.3.1}），故断言直接由
（\hyperref[III.1.4.6-fr]{1.4.6}）得到。
\item[(ii)] 由 Serre 判据（II，5.2.1），只需证明对
$\mathcal O_X\vert X_f$ 的每个拟相干理想 $\mathcal K$，有
$H^1(X_f,\mathcal K)=0$。由于 $X_f$ 是拟紧概形 $X$ 中的拟紧开集，
存在 $\mathcal O_X$ 的拟相干理想 $\mathcal J$，使得
$\mathcal K=\mathcal J\vert X_f$（I，9.4.2）。由
（\hyperref[III.1.4.6-fr]{1.4.6}），有
$H^1(X_f,\mathcal K)=(H^1(X,\mathcal J))_f$，而假设蕴含右端为零，
结论随之得到。
\end{enumerate}
\end{proof}

\begin{remark}[1.4.8]
\label{III.1.4.8-fr}
应注意，（\hyperref[III.1.4.7-fr]{1.4.7}，(i)）重新给出了关系
（II，4.5.13.2）的一个更简单的证明。
\end{remark}

\oldpage[III]{91}

\begin{lemma}[1.4.9]
\label{III.1.4.9-fr}
设 $X$ 是拟紧概形，
$\mathfrak U=(U_i)_{1\leq i\leq n}$ 是由仿射开集组成的 $X$ 的有限覆盖，
$\mathcal F$ 是拟相干 $\mathcal O_X$-模。由覆盖 $\mathfrak U$ 定义的层复形
$\mathcal C^\bullet(\mathfrak U,\mathcal F)$（G，II，5.2）于是是拟相干
$\mathcal O_X$-模复形。
\end{lemma}

\begin{proof}
由定义（G，II，5.2），$\mathcal C^p(\mathfrak U,\mathcal F)$ 是诸
$\mathcal F\vert U_{i_0\ldots i_p}$ 在典范嵌入
$U_{i_0\ldots i_p}\to X$ 下的直接像层之直和。由于 $X$ 是概形，
这些嵌入都是仿射态射（I，5.5.6），故诸
$\mathcal C^p(\mathfrak U,\mathcal F)$ 都拟相干（II，1.2.6）。
\end{proof}

\begin{proposition}[1.4.10]
\label{III.1.4.10-fr}
设 $u:X\to Y$ 是分离且拟紧的态射。对任意拟相干 $\mathcal O_X$-模
$\mathcal F$，诸 $R^qu_*(\mathcal F)$ 都是拟相干 $\mathcal O_Y$-模。
\end{proposition}

\begin{proof}
问题关于 $Y$ 是局部的，故可假设 $Y$ 仿射。此时 $X$ 是有限多个仿射开集
$U_i$（$1\leq i\leq n$）的并；令 $\mathfrak U$ 为覆盖 $(U_i)$。
此外，由于 $Y$ 是概形，由（I，5.5.10），对每个仿射开集
$V\subset Y$，典范嵌入 $u^{-1}(V)\to X$ 是仿射态射；因而由
（\hyperref[III.1.4.1-fr]{1.4.1}）和（G，II，5.2）得到典范同构
\begin{equation}
\label{III.1.4.10.1-fr}
\tag{1.4.10.1}
H^\bullet(u^{-1}(V),\mathcal F)
\xrightarrow{\sim}H^\bullet(\Gamma(V,\mathcal K^\bullet))
\end{equation}
其中置 $\mathcal K^\bullet=u_*(\mathcal C^\bullet(\mathfrak U,\mathcal F))$。
由（\hyperref[III.1.4.9-fr]{1.4.9}）和（I，9.2.2），
$\mathcal K^\bullet$ 是拟相干 $\mathcal O_Y$-模复形；另一方面，
它确为一个\emph{层复形}，因为 $\mathcal C^\bullet(\mathfrak U,\mathcal F)$
也是如此。于是由上同调 $\mathcal H^\bullet(\mathcal K^\bullet)$ 的定义
（G，II，4.1），后者由拟相干 $\mathcal O_Y$-模组成（I，4.1.1）。
又因对 $Y$ 中的仿射开集 $V$，函子 $\Gamma(V,\mathcal G)$ 在拟相干
$\mathcal O_Y$-模范畴中关于 $\mathcal G$ 是正合的，故有
（G，II，4.1）
\begin{equation}
\label{III.1.4.10.2-fr}
\tag{1.4.10.2}
H^\bullet(\Gamma(V,\mathcal K^\bullet))
:=\Gamma(V,\mathcal H^\bullet(\mathcal K^\bullet)).
\end{equation}
最后应注意，由（G，II，5.2）给出的典范同态
\[
H^\bullet(\mathfrak U,\mathcal F)\longrightarrow H^\bullet(X,\mathcal F)
\]
之定义可知，若 $V'\subset V$ 是 $Y$ 的另一个仿射开集，则图表
\[
\xymatrix{
H^\bullet(u^{-1}(V),\mathcal F)\ar[r]^{\sim}\ar[d]
&H^\bullet(\Gamma(V,\mathcal K^\bullet))\ar[d]\\
H^\bullet(u^{-1}(V'),\mathcal F)\ar[r]^{\sim}
&H^\bullet(\Gamma(V',\mathcal K^\bullet))
}
\]
交换。因此由上述结果可知，同构
（\hyperref[III.1.4.10.1-fr]{1.4.10.1}）定义了
$\mathcal O_Y$-模的同构
\begin{equation}
\label{III.1.4.10.3-fr}
\tag{1.4.10.3}
R^\bullet u_*(\mathcal F)
\xrightarrow{\sim}\mathcal H^\bullet(\mathcal K^\bullet)
\end{equation}
\oldpage[III]{92}
从而 $R^\bullet u_*(\mathcal F)$ 拟相干。
\end{proof}

此外，由（\hyperref[III.1.4.10.3-fr]{1.4.10.3}）、
（\hyperref[III.1.4.10.2-fr]{1.4.10.2}）和
（\hyperref[III.1.4.10.1-fr]{1.4.10.1}）还得到：

\begin{corollary}[1.4.11]
\label{III.1.4.11-fr}
在（\hyperref[III.1.4.10-fr]{1.4.10}）的假设下，对 $Y$ 的每个仿射开集
$V$，典范同态
\begin{equation}
\label{III.1.4.11.1-fr}
\tag{1.4.11.1}
H^q(u^{-1}(V),\mathcal F)\longrightarrow
\Gamma(V,R^qu_*(\mathcal F))
\end{equation}
对每个 $q\geq0$ 都是同构。
\end{corollary}

\begin{corollary}[1.4.12]
\label{III.1.4.12-fr}
假设（\hyperref[III.1.4.10-fr]{1.4.10}）的条件成立，并且 $Y$ 拟紧。
则存在整数 $r>0$，使得对任意拟相干 $\mathcal O_X$-模 $\mathcal F$
和任意整数 $q>r$，都有 $R^qu_*(\mathcal F)=0$。若 $Y$ 仿射，
则可取这样的整数 $r$：$X$ 存在由 $r$ 个仿射开集组成的覆盖。
\end{corollary}

\begin{proof}
由于可用有限多个仿射开集覆盖 $Y$，由
（\hyperref[III.1.4.11-fr]{1.4.11}），问题归结为证明第二个断言。
若 $\mathfrak U$ 是由 $r$ 个仿射开集组成的 $X$ 的覆盖，则当
$q>r$ 时有 $H^q(\mathfrak U,\mathcal F)=0$，因为
$C^q(\mathfrak U,\mathcal F)$ 的上链是交错的；故结论由
（\hyperref[III.1.4.1-fr]{1.4.1}）得到。
\end{proof}

\begin{corollary}[1.4.13]
\label{III.1.4.13-fr}
假设（\hyperref[III.1.4.10-fr]{1.4.10}）的条件成立，并且
$Y=\operatorname{Spec}(A)$ 仿射。则对任意拟相干 $\mathcal O_X$-模
$\mathcal F$ 和任意 $f\in A$，在典范同构的意义下有
\[
\Gamma(Y_f,R^qu_*(\mathcal F))
=(\Gamma(Y,R^qu_*(\mathcal F)))_f
\]
\end{corollary}

\begin{proof}
事实上，这来自 $R^qu_*(\mathcal F)$ 是拟相干 $\mathcal O_Y$-模
（I，1.3.7）。
\end{proof}

\begin{proposition}[1.4.14]
\label{III.1.4.14-fr}
设 $f:X\to Y$ 是分离拟紧态射，$g:Y\to Z$ 是仿射态射。对任意拟相干
$\mathcal O_X$-模 $\mathcal F$，典范同态
$R^p(g\circ f)_*(\mathcal F)\to g_*(R^pf_*(\mathcal F))$（0，12.2.5.2）
对每个 $p$ 都是双射。
\end{proposition}

\begin{proof}
事实上，对 $Z$ 的每个仿射开集 $W$，$g^{-1}(W)$ 是 $Y$ 的仿射开集。
因此，由（\hyperref[III.1.4.11-fr]{1.4.11}），定义典范同态
\[
R^p(g\circ f)_*(\mathcal F)\longrightarrow g_*(R^pf_*(\mathcal F))
\]
的预层同态（0，12.2.5）是双射。
\end{proof}

\begin{proposition}[1.4.15]
\label{III.1.4.15-fr}
设 $u:X\to Y$ 是分离拟紧态射
\SourceCorrectionNote{法文印本及其外交式转录在（III, 1.4.15）的假设中写作
“分离有限型态射”；EGA III 已发表勘误 $\mathbf{Err}_{\mathrm{III}},25$
明确要求把“有限型”改为“拟紧”。本译采用该已发表校正，法文外交式来源保持不变。}，
$v:Y'\to Y$ 是预概形的平坦态射（0\textsubscript{I}，6.7.1）；令
$u'=u_{(Y')}$，于是有交换图表
\begin{equation}
\label{III.1.4.15.1-fr}
\tag{1.4.15.1}
\xymatrix{
X\ar[d]_u & X'=X_{(Y')}\ar[l]_{v'}\ar[d]^{u'}\\
Y & Y'\ar[l]_v
}
\end{equation}
则对任意拟相干 $\mathcal O_X$-模 $\mathcal F$，
$R^qu_*'(\mathcal F')$（其中
$\mathcal F'=v'^*(\mathcal F)=\mathcal F\otimes_{\mathcal O_Y}\mathcal O_{Y'}$）
对每个 $q\geq0$ 都典范同构于
$R^qu_*(\mathcal F)\otimes_{\mathcal O_Y}\mathcal O_{Y'}
=v^*(R^qu_*(\mathcal F))$。
\end{proposition}

\oldpage[III]{93}
\begin{proof}
典范同态
$\rho:\mathcal F\to v_*'(v'^*(\mathcal F))$（0\textsubscript{I}，4.4.3.2）
由函子性定义同态
\begin{equation}
\label{III.1.4.15.2-fr}
\tag{1.4.15.2}
R^qu_*(\mathcal F)\longrightarrow R^qu_*(v_*'(\mathcal F')).
\end{equation}
另一方面，置 $w=u\circ v'=v\circ u'$，有典范同态
（0，12.2.5.1 和 12.2.5.2）
\begin{equation}
\label{III.1.4.15.3-fr}
\tag{1.4.15.3}
R^qu_*(v_*'(\mathcal F'))\longrightarrow R^qw_*(\mathcal F')
\longrightarrow v_*(R^qu_*'(\mathcal F')).
\end{equation}
复合（\hyperref[III.1.4.15.3-fr]{1.4.15.3}）与
（\hyperref[III.1.4.15.2-fr]{1.4.15.2}），得到同态
\[
\psi:R^qu_*(\mathcal F)\longrightarrow v_*(R^qu_*'(\mathcal F'))
\]
并由此最终得到典范同态（其定义对 $v$ \emph{不使用任何假设}）
\begin{equation}
\label{III.1.4.15.4-fr}
\tag{1.4.15.4}
\psi^\sharp:v^*(R^qu_*(\mathcal F))\longrightarrow R^qu_*'(\mathcal F')
\end{equation}
我们须证明，当 $v$ \emph{平坦}时，它是同构。显然问题关于 $Y$ 和
$Y'$ 是局部的，故可假设 $Y=\operatorname{Spec}(A)$、
$Y'=\operatorname{Spec}(B)$；此外还要用到下述引理。

\begin{lemma}[1.4.15.5]
\label{III.1.4.15.5-fr}
设 $\varphi:A\to B$ 是环同态，$Y=\operatorname{Spec}(A)$、
$X=\operatorname{Spec}(B)$，$f:X\to Y$ 是与 $\varphi$ 对应的态射，
$M$ 是 $B$-模。$\mathcal O_X$-模 $\widetilde M$ 是 $f$-平坦的
（0\textsubscript{I}，6.7.1），当且仅当 $M$ 是平坦 $A$-模。特别地，
态射 $f$ 平坦，当且仅当 $B$ 是平坦 $A$-模。
\end{lemma}

事实上，这由定义（0\textsubscript{I}，6.7.1）和
（0\textsubscript{I}，6.3.3）并计及（I，1.3.4）得到。

回到原证，由（\hyperref[III.1.4.11.1-fr]{1.4.11.1}）及同态
（\hyperref[III.1.4.15.3-fr]{1.4.15.3}）的定义
（参见 0，12.2.5），$\psi$ 此时对应于复合态射
\[
H^q(X,\mathcal F)\xrightarrow{\rho_q}
H^q(X,v_*'(v'^*(\mathcal F)))\xrightarrow{\theta_q}
H^q(X',v'^*(v_*'(v'^*(\mathcal F))))\xrightarrow{\sigma_q}
H^q(X',v'^*(\mathcal F))
\]
其中 $\rho_q$ 和 $\sigma_q$ 是上同调中分别对应于典范态射 $\rho$ 与
$\sigma:v'^*(v_*'(\mathcal G'))\to\mathcal G'$ 的同态，而 $\theta_q$
是关于 $\mathcal O_X$-模 $v_*'(v'^*(\mathcal F))$ 的
$\varphi$-同态（0，12.1.3.1）。但由 $\theta_q$ 的函子性，有交换图表
\[
\xymatrix{
H^q(X,\mathcal F)\ar[rr]^{\rho_q}\ar[dd]_{\theta_q}
&&H^q(X,v_*'(v'^*(\mathcal F)))\ar[dd]^{\theta_q}\\\\
H^q(X',v'^*(\mathcal F))\ar[rr]^{v'^*(\rho_q)}
&&H^q(X',v'^*(v_*'(v'^*(\mathcal F))))
}
\]
\oldpage[III]{94}
又因按定义（0\textsubscript{I}，4.4.3），$v'^*(\rho)$ 是
$\sigma$ 的逆，所以前述复合态射最终正是 $\theta_q$；从而
$\psi^\sharp$ 是所关联的 $B$-同态
\[
H^q(X,\mathcal F)\otimes_A B\longrightarrow H^q(X',\mathcal F').
\]
由于 $u$ 是拟紧的
\SourceCorrectionNote{法文印本及其外交式转录在证明此处仍写作“$u$ 是有限型的”，
而同一命题的已发表勘误已把假设改成“$u$ 拟紧”；有限仿射覆盖只需拟紧性。
Canon 校正现行法文 R2 已把此证明句一并改为“拟紧”，本译据此采用该可逆校正。}，
$X$ 是有限多个仿射开集 $U_i$（$1\leq i\leq r$）的并；令
$\mathfrak U$ 为覆盖 $(U_i)$。此外，$v$ 是仿射态射，故 $v'$ 亦然
（II，1.6.2，(iii)），于是诸 $U_i'=v'^{-1}(U_i)$ 构成 $X'$ 的仿射开覆盖
$\mathfrak U'$。已知（0，12.1.4.2）图表
\[
\xymatrix{
H^q(\mathfrak U,\mathcal F)\ar[r]^{\theta_q}\ar[d]
&H^q(\mathfrak U',\mathcal F')\ar[d]\\
H^q(X,\mathcal F)\ar[r]^{\theta_q}&H^q(X',\mathcal F')
}
\]
交换；又因 $X$ 和 $X'$ 都是概形，竖直箭头均为同构
（\hyperref[III.1.4.1-fr]{1.4.1}）。故只需证明典范
$\varphi$-同态
$\theta_q:H^q(\mathfrak U,\mathcal F)\to H^q(\mathfrak U',\mathcal F')$
所关联的 $B$-同态
\[
H^q(\mathfrak U,\mathcal F)\otimes_A B
\longrightarrow H^q(\mathfrak U',\mathcal F')
\]
是同构。对 $[1,r]$ 中任意 $p+1$ 个指标的序列
$\mathbf s=(i_k)_{0\leq k\leq p}$，置
\[
U_{\mathbf s}=\bigcap_{k=0}^pU_{i_k},\qquad
U_{\mathbf s}'=\bigcap_{k=0}^pU_{i_k}'=v'^{-1}(U_{\mathbf s}),\qquad
M_{\mathbf s}=\Gamma(U_{\mathbf s},\mathcal F),\qquad
M_{\mathbf s}'=\Gamma(U_{\mathbf s}',\mathcal F').
\]
典范映射 $M_{\mathbf s}\otimes_A B\to M_{\mathbf s}'$ 是同构
（I，1.6.5），故典范映射
$C^p(\mathfrak U,\mathcal F)\otimes_A B\to
C^p(\mathfrak U',\mathcal F')$ 是同构；在此同构下，$d\otimes1$
与上边缘算子
$C^p(\mathfrak U',\mathcal F')\to C^{p+1}(\mathfrak U',\mathcal F')$
相等。由于 $B$ 是\emph{平坦} $A$-模，从上同调模的定义立即得到典范映射
$H^q(\mathfrak U,\mathcal F)\otimes_A B\to
H^q(\mathfrak U',\mathcal F')$ 是同构（0\textsubscript{I}，6.1.1）。
这一结果以后将在（\S~6）中推广。
\end{proof}

\begin{corollary}[1.4.16]
\label{III.1.4.16-fr}
设 $A$ 是环，$X$ 是有限型 $A$-概形，$B$ 是在 $A$ 上忠实平坦的
$A$-代数。$X$ 是仿射的，当且仅当 $X\otimes_A B$ 是仿射的。
\end{corollary}

\begin{proof}
条件显然是必要的（I，3.2.2）；下面证明充分性。由于 $X$ 在 $A$ 上分离，
而态射 $\operatorname{Spec}(B)\to\operatorname{Spec}(A)$ 平坦，由
（1.4.15）
\SourceCorrectionNote{法文印本及其外交式转录此处误引（1.4.1）；Canon 稳定决定
EGA-III1-1.4.16-P438-CROSSREFERENCE-001 核定应引用平坦基变换命题
（1.4.15）。本译采用已核定读法，并保留可逆来源说明。}
得到
\begin{equation}
\label{III.1.4.16.1-fr}
\tag{1.4.16.1}
H^i(X\otimes_A B,\mathcal F\otimes_A B)
:=H^i(X,\mathcal F)\otimes_A B
\end{equation}
对每个 $i\geq0$ 和每个拟相干 $\mathcal O_X$-模 $\mathcal F$ 都成立。
若 $X\otimes_A B$ 仿射，则（\hyperref[III.1.4.16.1-fr]{1.4.16.1}）
左端在 $i=1$ 时为零；由于 $B$ 是忠实平坦 $A$-模，
$H^1(X,\mathcal F)$ 也为零。又因 $X$ 是拟紧概形，结论由 Serre 判据
（II，5.2.1）得到。
\end{proof}

\begin{proposition}[1.4.17]
\label{III.1.4.17-fr}
设 $X$ 是预概形，
$0\to\mathcal F\xrightarrow{u}\mathcal G\xrightarrow{v}\mathcal H\to0$
是 $\mathcal O_X$-模的正合列。若 $\mathcal F$ 和 $\mathcal H$ 拟相干，
则 $\mathcal G$ 也拟相干。
\end{proposition}

\oldpage[III]{95}
\begin{proof}
问题关于 $X$ 是局部的，故可假设 $X=\operatorname{Spec}(A)$ 仿射；
此时只需证明 $\mathcal G$ 满足（I，1.4.1）的条件 d~1) 和 d~2)
（取 $V=X$）。条件 d~2) 的验证是立即的：若
$t\in\Gamma(X,\mathcal G)$ 在 $D(f)$ 上的限制为零，则其像
$v(t)\in\Gamma(X,\mathcal H)$ 的限制也为零；因此存在 $m>0$，使得
$f^mv(t)=v(f^mt)=0$（I，1.4.1）。由于 $\Gamma$ 左正合，
$f^mt=u(s)$，其中 $s\in\Gamma(X,\mathcal F)$；又因 $u$ 单射，
$s$ 在 $D(f)$ 上的限制为零，故由（I，1.4.1）存在整数 $n>0$，使得
$f^ns=0$；最终得到 $f^{m+n}t=u(f^ns)=0$。

现在验证 d~1)
\SourceCorrectionNote{法文印本及其外交式转录在两段均写 d~2)；Canon 稳定决定
EGA-III1-1.4.17-P439-D1-D2-001 核定第二段应为延拓条件 d~1)。
本译采用已核定读法，并保留可逆来源说明。}。设
$t'\in\Gamma(D(f),\mathcal G)$；由于 $\mathcal H$
拟相干，存在整数 $m$，使得 $f^mv(t')=v(f^mt')$ 可延拓为截面
$z\in\Gamma(X,\mathcal H)$（I，1.4.1）。但把
（\hyperref[III.1.3.1-fr]{1.3.1}）（或（I，5.1.9.2））应用于拟相干
$\mathcal O_X$-模 $\mathcal F$，可知列
$\Gamma(X,\mathcal G)\to\Gamma(X,\mathcal H)\to0$ 正合，
故存在 $t\in\Gamma(X,\mathcal G)$ 使 $z=v(t)$。若以 $t''$ 表示 $t$
在 $D(f)$ 上的限制，便有 $v(f^mt'-t'')=0$；因此
$f^mt'-t''=u(s')$，其中 $s'\in\Gamma(D(f),\mathcal F)$。
又因 $\mathcal F$ 拟相干，存在整数 $n>0$，使 $f^ns'$ 可延拓为截面
$s\in\Gamma(X,\mathcal F)$。由
$f^{m+n}t'-f^nt''=u(f^ns')$ 可见，$f^{m+n}t'$ 是截面
$f^nt+u(s)\in\Gamma(X,\mathcal G)$ 在 $D(f)$ 上的限制
\SourceCorrectionNote{法文印本及其外交式转录此处误作 $f^nt+u(f^ns)$；Canon 稳定决定
EGA-III1-1.4.17-P439-SECTION-LIFT-001 核定为 $f^nt+u(s)$，因为 $s$
本身延拓 $f^ns'$。本译采用已核定读法，并保留可逆来源说明。}，
证明完毕。
\end{proof}

\section{射影态射的上同调研究}
\label{section:III.2-fr}

\subsection{若干上同调群的显式计算}
\label{subsection:III.2.1-fr}

\begin{env}[2.1.1]
\label{III.2.1.1-fr}
设 $X$ 是预概形，$\mathcal L$ 是可逆 $\mathcal O_X$-模；考虑分次环
（0\textsubscript{I}，5.4.6）
\begin{equation}
\label{III.2.1.1.1-fr}
\tag{2.1.1.1}
S=\Gamma_\bullet(X,\mathcal L)
:=\bigoplus_{n\in\mathbb Z}\Gamma(X,\mathcal L^{\otimes n}).
\end{equation}
设 $(f_i)_{1\leq i\leq r}$ 是 $S$ 中的有限族\emph{齐次}元，且
$f_i\in S_{d_i}$；置 $U_i=X_{f_i}$、$U=\bigcup_iU_i$，并以
$\mathfrak U$ 表示 $U$ 的覆盖 $(U_i)$。对任意拟相干 $\mathcal O_X$-模
$\mathcal F$，置
\begin{align}
\label{III.2.1.1.2-fr}
\tag{2.1.1.2}
H^\bullet(U,\mathcal F(*))
&=\bigoplus_{n\in\mathbb Z}
H^\bullet(U,\mathcal F\otimes\mathcal L^{\otimes n}),\\
\label{III.2.1.1.3-fr}
\tag{2.1.1.3}
H^\bullet(\mathfrak U,\mathcal F(*))
&=\bigoplus_{n\in\mathbb Z}
H^\bullet(\mathfrak U,\mathcal F\otimes\mathcal L^{\otimes n}).
\end{align}
阿贝尔群（\hyperref[III.2.1.1.2-fr]{2.1.1.2}）和
（\hyperref[III.2.1.1.3-fr]{2.1.1.3}）是\emph{双分次}的；这里取
\[
(H^\bullet(U,\mathcal F(*)))_{mn}
=H^m(U,\mathcal F\otimes\mathcal L^{\otimes n})
\]
并对（\hyperref[III.2.1.1.3-fr]{2.1.1.3}）作类似定义。就第二个次数而言，
这些群显然是分次 $S$-模；例如，这可由
$\mathcal F\mapsto H^m(U,\mathcal F)$ 和
$\mathcal F\mapsto H^m(\mathfrak U,\mathcal F)$ 都是函子看出。
\end{env}

\begin{env}[2.1.2]
\label{III.2.1.2-fr}
现在考虑分次 $S$-模（0\textsubscript{I}，5.4.6）
\begin{equation}
\label{III.2.1.2.1-fr}
\tag{2.1.2.1}
M=\Gamma_\bullet(\mathcal L,\mathcal F)
:=H^0(X,\mathcal F(*))
:=\bigoplus_{n\in\mathbb Z}
\Gamma(X,\mathcal F\otimes\mathcal L^{\otimes n}).
\end{equation}
\oldpage[III]{96}
若 $X$ 是底空间为 Noether 空间的预概形，或是拟紧概形，则由
（I，9.3.1），照常置
$U_{i_0i_1\ldots i_p}=\bigcap_{k=0}^pU_{i_k}$ 后，有典范同构
\[
\Gamma(U_{i_0i_1\ldots i_p},\mathcal F(*))
=H^0(U_{i_0i_1\ldots i_p},\mathcal F(*))
=M_{f_{i_0}f_{i_1}\cdots f_{i_p}}.
\]
还可沿用（\hyperref[III.1.2.2-fr]{1.2.2}）的记号，把
$M_{f_{i_0}f_{i_1}\cdots f_{i_p}}$ 认同于
$\varinjlim_n M_{i_0i_1\ldots i_p}^{(n)}$。如下定义次数后，这一认同是分次
$S$-模的同构：若齐次元
$z\in\varinjlim_n M_{i_0i_1\ldots i_p}^{(n)}$，则 $z$ 是齐次元
$x\in M_{i_0i_1\ldots i_p}^{(n)}=M$ 的典范像，且后者的次数为 $m$，
则把 $z$ 的次数定义为
$m-n(d_{i_0}+d_{i_1}+\cdots+d_{i_p})$。考虑到同态
$\varphi_{kn}:M_{i_0i_1\ldots i_p}^{(n)}\to
M_{i_0i_1\ldots i_p}^{(k)}$（\hyperref[III.1.2.2-fr]{1.2.2}）的定义，
立即可知这一定义与所取的 $z$ 的“代表”$x$ 无关。与
（\hyperref[III.1.2.2-fr]{1.2.2}）一样，以 $C_n^p(M)$ 表示从
$[1,r]^{p+1}$ 到 $M$ 的交错映射全体（对每个 $n$）；按上述相同方式在
$\varinjlim_n C_n^p(M)$ 上定义分次 $S$-模结构。仍如
（\hyperref[III.1.2.2-fr]{1.2.2}），有
\begin{equation}
\label{III.2.1.2.2-fr}
\tag{2.1.2.2}
C^p(\mathfrak U,\mathcal F(*))=\varinjlim_n C_n^p(M),
\end{equation}
且两端的同构\emph{保持次数}。于是仍如
（\hyperref[III.1.2.2-fr]{1.2.2}），有
\begin{equation}
\label{III.2.1.2.3-fr}
\tag{2.1.2.3}
C^p(\mathfrak U,\mathcal F(*))
:=C^{p+1}((\mathbf f),M)
:=\varinjlim_n K^{p+1}(\mathbf f^n,M),
\end{equation}
且该同构保持次数：$\varinjlim_n K^{p+1}(\mathbf f^n,M)$ 中一个元素若是
上链 $g\in K^{p+1}(\mathbf f^n,M)$ 的典范像，而诸值
$g(i_0,\ldots,i_p)$ 都属于 $M$ 的同一个齐次分量 $M_m$，则它的次数为
$m-n(d_{i_0}+\cdots+d_{i_p})$；这一数值与所考虑元素的上链代表选择无关。
\end{env}

由于上述诸同构与上边缘算子相容，仍如
（\hyperref[III.1.2.2-fr]{1.2.2}）可得：

\begin{proposition}[2.1.3]
\label{III.2.1.3-fr}
设 $X$ 是底空间为 Noether 空间的预概形，或是拟紧概形。存在关于
$\mathcal F$ 函子性的典范同构
\begin{equation}
\label{III.2.1.3.1-fr}
\tag{2.1.3.1}
H^p(\mathfrak U,\mathcal F(*))
\xrightarrow{\sim}H^{p+1}((\mathbf f),M)
\qquad\text{对所有 }p\geq1.
\end{equation}
此外，还有关于 $\mathcal F$ 函子性的正合列
\begin{equation}
\label{III.2.1.3.2-fr}
\tag{2.1.3.2}
0\longrightarrow H^0((\mathbf f),M)\longrightarrow M
\longrightarrow H^0(\mathfrak U,\mathcal F(*))
\longrightarrow H^1((\mathbf f),M)\longrightarrow0.
\end{equation}
而且，对分次 $S$-模结构（其中 $S$ 是环
（\hyperref[III.2.1.1.1-fr]{2.1.1.1}）），所引入的所有同态次数均为 $0$。
\end{proposition}
\providecommand{\bbGamma}{\boldsymbol{\Gamma}}
\providecommand{\shProj}{\mathcal{P}\mathrm{roj}}

\oldpage[III]{97}
\begin{corollary}[2.1.4]
\label{III.2.1.4-fr}
若 $X$ 是拟紧概形且 $U_i=X_{f_i}$ 均为仿射的，则存在关于 $\mathcal F$
函子性、次数为 $0$ 的典范同构
\begin{equation}
\label{III.2.1.4.1-fr}
\tag{2.1.4.1}
H^p(U,\mathcal F(*))\xrightarrow{\sim}H^{p+1}((\mathbf f),M)
\qquad\text{当 }p\geq1,
\end{equation}
以及关于 $\mathcal F$ 函子性的正合列
\begin{equation}
\label{III.2.1.4.2-fr}
\tag{2.1.4.2}
0\longrightarrow H^0((\mathbf f),M)\longrightarrow M
\longrightarrow H^0(U,\mathcal F(*))
\longrightarrow H^1((\mathbf f),M)\longrightarrow0,
\end{equation}
其中所有同态的次数均为 $0$。
\end{corollary}

\begin{proof}
只需把（1.4.1）应用于（\hyperref[III.2.1.3-fr]{2.1.3}）的结果。
\end{proof}

与（\hyperref[III.2.1.3-fr]{2.1.3}）类似的“局部”命题如下：

\begin{proposition}[2.1.5]
\label{III.2.1.5-fr}
设 $S$ 是正分次环，$f_i$（$(1\leq i\leq r)$）是 $S_+$ 中次数为
$d_i$ 的齐次元，$M$ 是分次 $S$-模。设
$X=\operatorname{Proj}(S)$ 是 $S$ 的齐次素谱，并置
$U_i=D_+(f_i)$、$U=\bigcup_iU_i$ 以及
\[
H^\bullet(U,\widetilde M(*))
=\bigoplus_{n\in\mathbb Z}H^\bullet(U,(M(n))^\sim).
\]
于是存在关于 $M$ 函子性的典范同构；对分次 $S$-模结构，这些同构的次数为
$0$：
\begin{equation}
\label{III.2.1.5.1-fr}
\tag{2.1.5.1}
H^p(U,\widetilde M(*))\xrightarrow{\sim}H^{p+1}((\mathbf f),M)
\qquad\text{当 }p\geq1,
\end{equation}
并且有关于 $M$ 函子性的正合列
\begin{equation}
\label{III.2.1.5.2-fr}
\tag{2.1.5.2}
0\longrightarrow H^0((\mathbf f),M)\longrightarrow M
\longrightarrow H^0(U,\widetilde M(*))
\longrightarrow H^1((\mathbf f),M)\longrightarrow0,
\end{equation}
其中所有同态的次数均为 $0$。
\end{proposition}

\begin{proof}
事实上，按定义（II，2.5.2），有
\[
\Gamma(U_{i_0i_1\ldots i_p},(M(n))^\sim)
=(M_{f_{i_0}\cdots f_{i_p}})_n
\]
故
$\Gamma(U_{i_0i_1\ldots i_p},\widetilde M(*))
=M_{f_{i_0}\cdots f_{i_p}}$。考虑到 $X$ 是概形，其余论证便与证明
（\hyperref[III.2.1.4-fr]{2.1.4}）时相同。
\end{proof}

\begin{env}[2.1.6]
\label{III.2.1.6-fr}
\emph{注记}。---
(i) 在（\hyperref[III.2.1.5-fr]{2.1.5}）的条件下，由
（0\textsubscript{I}，1.3.2），函子
$\Gamma(U_{i_0i_1\ldots i_p},\widetilde M(*))$ 关于 $M$ 是正合的；
与（1.2.5）中相同的论证表明，若
$0\to M'\to M\to M''\to0$ 是分次 $S$-模的正合列
（其中同态次数为 $0$），则对每个 $p\geq0$ 都有交换图表
\begin{equation}
\label{III.2.1.6.1-fr}
\tag{2.1.6.1}
\xymatrix{
H^p(U,\widetilde {M''}(*))\ar[r]^-\partial\ar[d] &
H^{p+1}(U,\widetilde {M'}(*))\ar[d]\\
H^{p+1}((\mathbf f),M'')\ar[r]^-\partial &
H^{p+2}((\mathbf f),M').
}
\end{equation}

(ii) 命题（\hyperref[III.2.1.5-fr]{2.1.5}）最有意义的情形，是 $S$
为由有限多个次数 $1$ 的元素生成的 $A$-代数，且假设 $A$ 为 Noether 环；
\oldpage[III]{98}
事实上，在这种情形下，每个拟相干 $\mathcal O_X$-模都具有
$\widetilde M$ 的形式（II，2.7.7）。
\end{env}

\begin{env}[2.1.7]
\label{III.2.1.7-fr}
我们将把（\hyperref[III.2.1.5-fr]{2.1.5}）应用于如下情形：
$S=A[T_0,\ldots,T_r]$，其中 $A$ 是任意环，$T_i$ 是不定元，并取
$M=S$、$f_i=T_i$。因此，问题实质上归结为计算
$H^\bullet((\mathbf T),S)$，其中
$\mathbf T=(T_i)_{0\leq i\leq r}$。
\end{env}

\begin{lemma}[2.1.8]
\label{III.2.1.8-fr}
若 $S=A[T_0,\ldots,T_r]$，并置
$\mathbf T=(T_i)_{0\leq i\leq r}$，则
\begin{align}
\label{III.2.1.8.1-fr}
\tag{2.1.8.1}
H^i(\mathbf T^n,S)&=0 &&\text{若 }i\neq r+1,\\
\label{III.2.1.8.2-fr}
\tag{2.1.8.2}
H^{r+1}(\mathbf T^n,S)&=S/(\mathbf T^n).
\end{align}
因此，$A$-模 $H^{r+1}(\mathbf T^n,S)$ 在 $A$ 上有一组基，由单项式
\[
T_{\mathbf p}=T_0^{p_0}\cdots T_r^{p_r},
\qquad \mathbf p=(p_0,\ldots,p_r),\quad 0\leq p_i<n
\quad\text{对所有 }i.
\]
模 $(\mathbf T^n)$ 的类组成。
\end{lemma}

\begin{proof}
这是（1.1.3.5）和命题（1.1.4）的直接推论；后者的假设显然成立。
\end{proof}

\begin{env}[2.1.9]
\label{III.2.1.9-fr}
对 $n$ 取归纳极限。关系（\hyperref[III.2.1.8.1-fr]{2.1.8.1}）给出
$H^i((\mathbf T),S)=0$（$i\neq r+1$）。当 $i=r+1$ 时，归纳系由
$S/(\mathbf T^n)$ 组成，而同态
\[
\varphi_{nm}:S/(\mathbf T^n)\longrightarrow S/(\mathbf T^m)
\qquad(0\leq n\leq m)
\]
是乘以 $(T_0\cdots T_r)^{m-n}$。当
$n\geq\sup(p_i)_{0\leq i\leq r}$ 时，以
$\xi_{\mathbf p}^{(n)}=\xi_{p_0,\ldots,p_r}^{(n)}$ 表示
$T_0^{n-p_0}\cdots T_r^{n-p_r}\pmod{(\mathbf T^n)}$ 的类；于是
$\varphi_{nm}(\xi_{\mathbf p}^{(n)})=\xi_{\mathbf p}^{(m)}$，因而这些元素在归纳极限
$H^{r+1}((\mathbf T),S)$ 中具有同一个典范像
$\xi_{\mathbf p}=\xi_{p_0,\ldots,p_r}$。根据
（\hyperref[III.2.1.2-fr]{2.1.2}）中次数的定义，$\xi_{\mathbf p}$ 的次数等于
$-|\mathbf p|=-(p_0+p_1+\cdots+p_r)$。显然，当 $0<p_i\leq n$
且 $0\leq i\leq r$ 时，诸 $\xi_{\mathbf p}^{(n)}$ 构成
$S/(\mathbf T^n)$ 的一组基。因此，由
（\hyperref[III.2.1.8-fr]{2.1.8}）立即得到：
\end{env}

\begin{corollary}[2.1.10]
\label{III.2.1.10-fr}
沿用（\hyperref[III.2.1.8-fr]{2.1.8}）的记号，有
\begin{equation}
\label{III.2.1.10.1-fr}
\tag{2.1.10.1}
H^i((\mathbf T),S)=0\qquad\text{当 }i\neq r+1,
\end{equation}
并且 $H^{r+1}((\mathbf T),S)$ 是自由 $A$-模；其中一组基由满足
$p_i>0$（$0\leq i\leq r$）的元素 $\xi_{p_0,\ldots,p_r}$ 组成。
\end{corollary}

\begin{env}[2.1.11]
\label{III.2.1.11-fr}
\emph{注记}。--- 设 $N$ 是任意 $A$-模并置 $M=S\otimes_A N$；
（\hyperref[III.2.1.8-fr]{2.1.8}）中的论证更一般地表明
\begin{align}
\label{III.2.1.11.1-fr}
\tag{2.1.11.1}
H^i(\mathbf T^n,M)&=0 &&\text{若 }i\neq r+1,\\
\label{III.2.1.11.2-fr}
\tag{2.1.11.2}
H^{r+1}(\mathbf T^n,M)&=(S/(\mathbf T^n))\otimes_A N.
\end{align}
事实上，最后一个公式直接由（1.1.3.5）得到；另一方面，显然有
\[
M/(T_0^nM+\cdots+T_{i-1}^nM)
=\bigl(S/(T_0^nS+\cdots+T_{i-1}^nS)\bigr)\otimes_A N,
\]
而理想 $T_0^nS+\cdots+T_{i-1}^nS$ 是 $A$-模 $S$ 的直和因子；
由此可以把（1.1.4）应用于 $M$，从而得到
（\hyperref[III.2.1.11.1-fr]{2.1.11.1}）。
\end{env}

结合（\hyperref[III.2.1.10-fr]{2.1.10}）和
（\hyperref[III.2.1.5-fr]{2.1.5}），得到：

\begin{proposition}[2.1.12]
\label{III.2.1.12-fr}
设 $A$ 是环，$r$ 是整数 $>0$，且 $X=\mathbf P_A^r$（II，4.1.1）。则：
\begin{enumerate}
\item[{\rm(i)}] 当 $i\neq0,r$ 时，有 $H^i(X,\mathcal O_X(*))=0$。
\item[{\rm(ii)}] 典范同态
$\alpha:S\to H^0(X,\mathcal O_X(*))$（II，2.6.2）是双射。
\oldpage[III]{99}
\item[{\rm(iii)}] $H^r(X,\mathcal O_X(*))$ 是自由 $A$-模，其一组基由元素
$\xi_{p_0,\ldots,p_r}$ 组成，其中 $p_i>0$（$0\leq i\leq r$）；
$\xi_{p_0,\ldots,p_r}$ 的次数为
$-|\mathbf p|=-(p_0+\cdots+p_r)$，且乘积
$T_i\xi_{p_0,\ldots,p_r}$ 等于
$\xi_{p_0,\ldots,p_i-1,\ldots,p_r}$。
\end{enumerate}
\end{proposition}

\begin{proof}
事实上，把正合列（\hyperref[III.2.1.5.2-fr]{2.1.5.2}）应用于
$M=S=A[T_0,\ldots,T_r]$ 时，由
（\hyperref[III.2.1.10.1-fr]{2.1.10.1}）有
$H^0((\mathbf T),S)=0$ 和 $H^1((\mathbf T),S)=0$。又因 $X$ 是
诸 $D_+(T_i)$ 的并（II，2.3.14），命题
（\hyperref[III.2.1.5-fr]{2.1.5}）可应用于 $U=X$。只需再把正合列
（\hyperref[III.2.1.5.2-fr]{2.1.5.2}）中的映射
$S\to H^0(X,\mathcal O_X(*))$ 与典范映射 $\alpha$ 认同；而这由
$H^0(U,\mathcal O_X(*))$ 和 $H^0(\mathfrak U,\mathcal O_X(*))$
的典范认同得到。
\end{proof}

\begin{corollary}[2.1.13]
\label{III.2.1.13-fr}
可能使 $H^i(X,\mathcal O_X(n))\neq0$ 的 $(i,n)$ 只有如下取值：
$i=0$ 且 $n\geq0$，或 $i=r$ 且 $n\leq-(r+1)$。
\end{corollary}

注意，若 $A\neq0$，则对（\hyperref[III.2.1.13-fr]{2.1.13}）中列出的
每一对取值，确有 $H^i(X,\mathcal O_X(n))\neq0$；这由
（\hyperref[III.2.1.12-fr]{2.1.12}）得到，因为此时对所有次数
$n\geq0$，$S_n$ 都是 $\neq0$ 的。

在本章的应用中，我们主要使用如下较弱的结果：

\begin{corollary}[2.1.14]
\label{III.2.1.14-fr}
$A$-模 $H^i(X,\mathcal O_X(n))$ 是有限型自由模；若 $i>0$，则当
$n>0$ 时它们为零。
\end{corollary}

\begin{proposition}[2.1.15]
\label{III.2.1.15-fr}
设 $Y$ 是预概形，$\mathcal E$ 是秩 $r+1$ 的局部自由 $\mathcal O_Y$-模，
$X=\mathbf P(\mathcal E)$ 是由 $\mathcal E$ 定义的射影丛，
$f:X\to Y$ 是结构态射。可能使 $R^if_*(\mathcal O_X(n))\neq0$ 的
$i$ 和 $n$ 只有 $i=0$ 且 $n\geq0$，或 $i=r$ 且 $n\leq-(r+1)$；
此外，典范同态（II，3.3.2）
\[
\alpha:S_{\mathcal O_Y}(\mathcal E)
\longrightarrow\Gamma_*(\mathcal O_X)
=R^0f_*(\mathcal O_X(*))
=\bigoplus_{n\in\mathbb Z}f_*(\mathcal O_X(n))
\]
是同构。
\end{proposition}

\begin{proof}
问题关于 $Y$ 是局部的，故可假设 $Y$ 是环为 $A$ 的仿射概形，且
$\mathcal E=\widetilde E$，其中 $E=A^{r+1}$。考虑到（1.4.11），
问题立即归结为（\hyperref[III.2.1.12-fr]{2.1.12}）。
\end{proof}

\begin{env}[2.1.16]
\label{III.2.1.16-fr}
\emph{注记}。--- 我们以后将补充
（\hyperref[III.2.1.15-fr]{2.1.15}）的结果，证明如下命题。置
\[
\boldsymbol\omega=f^*\left(\bigwedge^{r+1}\mathcal E\right)(-r-1),
\]
这是一个可逆 $\mathcal O_X$-模。则：
\begin{enumerate}
\item[{\rm(i)}] 有典范同构
\begin{equation}
\label{III.2.1.16.1-fr}
\tag{2.1.16.1}
\rho:R^rf_*(\boldsymbol\omega)\xleftarrow{\sim}\mathcal O_Y.
\end{equation}
\item[{\rm(ii)}] 杯积给出的配对（0，12.2.2）
\begin{equation}
\label{III.2.1.16.2-fr}
\tag{2.1.16.2}
R^rf_*(\mathcal O_X(n))\times R^0f_*(\boldsymbol\omega(-n))
\longrightarrow R^rf_*(\boldsymbol\omega)
\end{equation}
与同构 $\rho^{-1}$ 复合后，给出从 $R^rf_*(\mathcal O_X(n))$ 到局部自由
$\mathcal O_Y$-模
\[
R^0f_*(\boldsymbol\omega(-n))
=\left(\bigwedge^{r+1}\mathcal E\right)
\otimes_{\mathcal O_Y}\bigl(S_{\mathcal O_Y}(\mathcal E)\bigr)_{-n}.
\]
的\emph{对偶}的一个\emph{同构}。
\end{enumerate}
\end{env}

\subsection{射影态射的基本定理}
\label{subsection:III.2.2-fr}
\oldpage[III]{100}

\begin{theorem}[2.2.1]
\label{III.2.2.1-fr}
（Serre）设 $Y$ 是 Noether 预概形，$f:X\to Y$ 是紧合态射，
$\mathcal L$ 是一个对 $f$ 丰沛的可逆 $\mathcal O_X$-模。对每个
$\mathcal O_X$-模 $\mathcal F$，置
\[
\mathcal F(n)=\mathcal F\otimes_{\mathcal O_X}\mathcal L^{\otimes n}
\qquad\text{对所有 }n\in\mathbb Z.
\]
于是，对每个凝聚 $\mathcal O_X$-模 $\mathcal F$：
\begin{enumerate}
\item[{\rm(i)}] $R^qf_*(\mathcal F)$ 是凝聚 $\mathcal O_Y$-模。
\item[{\rm(ii)}] 存在整数 $N$，使得当 $n\geq N$ 时，
\[
R^qf_*(\mathcal F(n))=0\qquad\text{对所有 }q>0.
\]
\item[{\rm(iii)}] 存在整数 $N$，使得当 $n\geq N$ 时，典范同态
\[
f^*(f_*(\mathcal F(n)))\longrightarrow\mathcal F(n)
\]
是满射。
\end{enumerate}
\end{theorem}

\begin{proof}
首先注意，若把 $\mathcal L$ 换成 $\mathcal L^{\otimes d}$（$(d>0)$）后
定理成立，则定理的原形式也成立。事实上，此时可以写成
\[
\mathcal F(n)
=(\mathcal F\otimes\mathcal L^{\otimes r})
\otimes\mathcal L^{\otimes hd}
\]
其中 $h>0$ 且 $0\leq r<d$；根据假设，对每个 $r$ 都存在整数 $N_r$，
使得当 $h\geq N_r$ 时，性质 (ii) 和 (iii) 对
$\mathcal O_X$-模 $\mathcal F\otimes\mathcal L^{\otimes r}$ 成立。
取 $N$ 为诸 $dN_r$ 中的最大者，则当 $n\geq N$ 时，(ii) 和 (iii)
成立。因此，可假设 $\mathcal L$ 对 $f$ 极丰沛（II，4.6.11）；
于是存在支配的 $Y$-开浸入 $i:X\to P$，其中
$P=\operatorname{Proj}(\mathcal S)$，$\mathcal S$ 是正分次拟相干
$\mathcal O_Y$-代数，其 $\mathcal S_1$ 为有限型并生成 $\mathcal S$；
而且 $\mathcal L$ 同构于 $i^*(\mathcal O_P(1))$（II，4.4.7）。
但因 $f$ 紧合，$i$ 也紧合（II，5.4.4），故 $i$ 是同构
$X\xrightarrow{\sim}P$。因此，只需考虑
$X=\operatorname{Proj}(\mathcal S)$ 且 $\mathcal L=\mathcal O_X(1)$
的情形。此时，定理（2.2.1）是下列命题的推论。
\end{proof}

\begin{proposition}[2.2.2]
\label{III.2.2.2-fr}
设 $A$ 是 Noether 环，$S$ 是正分次 $A$-代数，其中 $S_1$ 作为
$A$-模有一个含 $r+1$ 个生成元的生成系，并生成代数 $S$。置
$X=\operatorname{Proj}(S)$。对每个凝聚 $\mathcal O_X$-模 $\mathcal F$：
\begin{enumerate}
\item[{\rm(i)}] $A$-模 $H^q(X,\mathcal F)$ 是有限型的。
\item[{\rm(ii)}] 当 $q>r$ 时，有 $H^q(X,\mathcal F)=0$。
\item[{\rm(iii)}] 存在整数 $N$，使得当 $n\geq N$ 时，对所有 $q>0$ 有
$H^q(X,\mathcal F(n))=0$。
\item[{\rm(iv)}] 存在整数 $N$，使得当 $n\geq N$ 时，$\mathcal F(n)$
由其在 $X$ 上的截面生成。
\end{enumerate}
\end{proposition}

\begin{proof}
先说明（2.2.2）如何推出（2.2.1）。在（2.2.1）中，问题已归结为上述
$X=\operatorname{Proj}(\mathcal S)$ 的特殊情形；$Y$ 拟紧，因而可由有限多个
环为 Noether 环的仿射开集覆盖，并且 $\mathcal S_1$ 在每个这样的开集
$U_\alpha$ 上的限制，都由 $\mathcal S_1$ 在 $U_\alpha$ 上的有限多个截面生成。
若已证明（2.2.2），则在（2.2.1）的 (ii) 和 (iii) 中，只需把 $N$
取为对应于诸 $U_\alpha$ 的同类整数中的最大者（并考虑（1.4.11）和
（II，3.4.7））。

为证明（2.2.2），注意 $X$ 可认同于 $P=\mathbf P_A^r$ 的一个闭子概形
（II，3.6.2）；而且，若 $j:X\to P$ 是典范浸入，则
$j_*(\mathcal F)$ 是凝聚 $\mathcal O_P$-模，并有
$j_*(\mathcal F(n))=(j_*(\mathcal F))(n)$（II，3.4.5 和 3.5.2）。
考虑（G，II，定理 4.9.1 的推论），问题因而归结为证明如下特殊情形的
（2.2.2）：
\[
X=\mathbf P_A^r\quad\text{且}\quad S=A[T_0,\ldots,T_r].
\]

\oldpage[III]{101}
由于 $X$ 由 $r+1$ 个仿射开集 $D_+(T_i)$ 覆盖，(ii) 由（1.4.12）得到。
另一方面，(iv) 已在（II，2.7.9）中证明。

我们将同时证明 (i) 和 (iii)。注意，对
$\mathcal F=\mathcal O_X(m)$，这些断言成立
（\hyperref[III.2.1.13-fr]{2.1.13}）；故当 $\mathcal F$ 是有限多个形如
$\mathcal O_X(m_j)$ 的 $\mathcal O_X$-模的直和时，它们也成立。另一方面，根据 (ii)，
当 $q>r$ 时，(i) 和 (iii) 显然成立。我们对 $q$ 作\emph{降归纳}。
已知 $\mathcal F$ 同构于有限多个层 $\mathcal O_X(m_j)$ 的直和
$\mathcal E$ 的一个商（II，2.7.10）；换言之，有正合列
\[
0\longrightarrow\mathcal R\longrightarrow\mathcal E
\longrightarrow\mathcal F\longrightarrow0,
\]
其中 $\mathcal R$ 是凝聚的（0\textsubscript{I}，5.3.3），且
$\mathcal E$ 满足 (i) 和 (iii)。由于 $\mathcal F(n)$ 关于 $\mathcal F$
是正合函子，对每个 $n\in\mathbb Z$ 还有正合列
\[
0\longrightarrow\mathcal R(n)\longrightarrow\mathcal E(n)
\longrightarrow\mathcal F(n)\longrightarrow0
\]
由此得到上同调正合列
\[
H^{q-1}(X,\mathcal E(n))\longrightarrow
H^{q-1}(X,\mathcal F(n))\longrightarrow
H^q(X,\mathcal R(n)).
\]
由于 $\mathcal E(n)$ 是诸 $\mathcal O_X(n+m_j)$ 的直和
（II，2.5.14），$H^{q-1}(X,\mathcal E(n))$ 是有限型的；按归纳假设，
$H^q(X,\mathcal R(n))$ 也是有限型的。因 $A$ 是 Noether 环，得出对每个
$n\in\mathbb Z$，$H^{q-1}(X,\mathcal F(n))$ 都是有限型的，特别地在
$n=0$ 时如此。另一方面，按归纳假设，存在整数 $N$，使得当
$n\geq N$ 时 $H^q(X,\mathcal R(n))=0$；还可假设 $N$ 选得使得当
$n\geq N$ 时 $H^{q-1}(X,\mathcal E(n))=0$，因为 $\mathcal E$ 满足
(iii)。于是，当 $n\geq N$ 时
$H^{q-1}(X,\mathcal F(n))=0$，证明完毕。
\end{proof}

\begin{corollary}[2.2.3]
\label{III.2.2.3-fr}
在（\hyperref[III.2.2.1-fr]{2.2.1}）的假设下，设
$\mathcal F\to\mathcal G\to\mathcal H$ 是凝聚 $\mathcal O_X$-模的正合列。
则存在整数 $N$，使得当 $n\geq N$ 时，序列
\[
f_*(\mathcal F(n))\longrightarrow f_*(\mathcal G(n))
\longrightarrow f_*(\mathcal H(n))
\]
正合。
\end{corollary}

\begin{proof}
以 $\mathcal F'$、$\mathcal G'$、$\mathcal G''$ 分别表示
$\mathcal F\to\mathcal G$ 的核、像和余核；$\mathcal G'$ 是
$\mathcal G\to\mathcal H$ 的核，$\mathcal G''$ 是其像；再以
$\mathcal H''$ 表示这一同态的余核。所有这些 $\mathcal O_X$-模都是凝聚的
（0\textsubscript{I}，5.3.4）。由于 $\mathcal F(n)$ 关于 $\mathcal F$
是正合函子，只需证明当 $n$ 充分大时，下列各序列均正合：
\begin{align*}
0&\longrightarrow f_*(\mathcal F'(n))
\longrightarrow f_*(\mathcal F(n))
\longrightarrow f_*(\mathcal G'(n))\longrightarrow0,\\
0&\longrightarrow f_*(\mathcal G'(n))
\longrightarrow f_*(\mathcal G(n))
\longrightarrow f_*(\mathcal G''(n))\longrightarrow0,\\
0&\longrightarrow f_*(\mathcal G''(n))
\longrightarrow f_*(\mathcal H(n))
\longrightarrow f_*(\mathcal H''(n))\longrightarrow0.
\end{align*}
因而可假设 $0\to\mathcal F\to\mathcal G\to\mathcal H\to0$ 正合。
此时有上同调正合列
\[
0\longrightarrow f_*(\mathcal F(n))
\longrightarrow f_*(\mathcal G(n))
\longrightarrow f_*(\mathcal H(n))
\longrightarrow R^1f_*(\mathcal F(n))\longrightarrow\cdots
\]
结论由（\hyperref[III.2.2.1-fr]{2.2.1}，(ii)）得到。
\end{proof}

\begin{corollary}[2.2.4]
\label{III.2.2.4-fr}
设 $Y$ 是 Noether 预概形，$f:X\to Y$ 是有限型态射，$\mathcal L$ 是
一个对 $f$ 丰沛的可逆 $\mathcal O_X$-模；对每个 $\mathcal O_X$-模
$\mathcal F$，置
$\mathcal F(n)=\mathcal F\otimes_{\mathcal O_X}\mathcal L^{\otimes n}$
（其中 $n\in\mathbb Z$）。设
$\mathcal F\to\mathcal G\to\mathcal H$ 是凝聚 $\mathcal O_X$-模的正合列，
并且 $\operatorname{Im}(\mathcal F\to\mathcal G)$ 的支撑在 $Y$ 上紧合。
则存在整数 $N$，使得当 $n\geq N$ 时，序列
\[
f_*(\mathcal F(n))\longrightarrow f_*(\mathcal G(n))
\longrightarrow f_*(\mathcal H(n))
\]
正合。
\SourceCorrectionNote{已发表勘误 $\mathbf{Err}_{\mathrm{III}},26$ 把印本中
“$\mathcal F$ 与 $\mathcal H$ 的支撑在 $Y$ 上紧合”改为
“$\operatorname{Im}(\mathcal F\to\mathcal G)$ 的支撑在 $Y$ 上紧合”，
并替换证明的相应部分。本译采用该已发表勘误；外交式法文保持不变。}
\end{corollary}

\begin{proof}
与（\hyperref[III.2.2.1-fr]{2.2.1}）证明开头相同的论证表明，若此推论对
$\mathcal L^{\otimes d}$（$(d>0)$）成立，则它对 $\mathcal L$ 也成立；
故只需考虑 $\mathcal L$ 对 $f$ 极丰沛的情形（II，4.6.11）。于是可把
$X$ 认同于一个 $Y$-概形 $Z=\operatorname{Proj}(\mathcal S)$ 中的开集，
其中 $\mathcal S$ 是正分次拟相干 $\mathcal O_Y$-代数，
$\mathcal S_1$ 是有限型的并生成 $\mathcal S$，且
$\mathcal L=i^*(\mathcal O_Z(1))$，这里 $i$ 是典范浸入 $X\to Z$
（II，4.4.7）。

置
$\mathcal G_1=\operatorname{Im}(\mathcal F\to\mathcal G)
:=\operatorname{Ker}(\mathcal G\to\mathcal H)$；它是凝聚的。由于有正合列
$0\to\mathcal G_1(n)\to\mathcal G(n)\to\mathcal H(n)$，且函子
$f_*$ 左正合，对每个 $n$，序列
$0\to f_*(\mathcal G_1(n))\to f_*(\mathcal G(n))\to
f_*(\mathcal H(n))$ 都正合；因此，只需证明当 $n$ 充分大时，序列
$f_*(\mathcal F(n))\to f_*(\mathcal G_1(n))\to0$ 正合。换言之，
只需考虑 $\mathcal H=0$ 且 $\mathcal G$ 的支撑在 $Y$ 上紧合的情形；
于是该支撑在 $Z$ 中是\emph{闭的}（II，5.4.10）。因此，
$\mathcal G'=i_*(\mathcal G)$ 是凝聚 $\mathcal O_Z$-模，且
$\mathcal G'|X=\mathcal G$。已知（I，9.4.3）存在凝聚
$\mathcal O_Z$-模 $\mathcal F'$，使得 $\mathcal F'|X=\mathcal F$。
此外，由于 $X$ 在 $Z$ 中是开集，而
$\operatorname{Supp}(\mathcal G)\subset X$ 在 $Z$ 中闭，立即可按如下方式定义
层的满同态 $u':\mathcal F'\to\mathcal G'$，使 $u'|X$ 等于给定的满同态
$u:\mathcal F\to\mathcal G$：对每个不与
$\operatorname{Supp}(\mathcal G)$ 相交的 $Z$ 中开集 $U$，取 $u'|U=0$；
对每个开集 $U\subset X$，取 $u'|U=u|U$。

与（2.2.2）证明开头相同，可把问题归结为 $Y$ 仿射的情形；因此只需证明，
当 $n$ 充分大时，同态
$\Gamma(u):\Gamma(X,\mathcal F(n))\to\Gamma(X,\mathcal G(n))$ 是满射。
有交换图表
\[
\xymatrix@C=35mm@R=12mm{
\Gamma(Z,\mathcal F'(n)) \ar[r]^{\Gamma(u')} \ar[d]_v &
\Gamma(Z,\mathcal G'(n)) \ar[d]^w \\
\Gamma(X,\mathcal F(n)) \ar[r]_{\Gamma(u)} & \Gamma(X,\mathcal G(n))
}
\]
其中 $v,w$ 是限制同态。$\mathcal G'$ 的定义还表明 $w$ 是\emph{双射}；
另一方面，由（\hyperref[III.2.2.3-fr]{2.2.3}），当 $n$ 充分大时，$\Gamma(u')$ 是\emph{满射}，
故 $\Gamma(u)$ 也为满射。
\end{proof}

\begin{env}[2.2.5]
\label{III.2.2.5-fr}
\emph{注记}。---
(i) 若只假设 $Y$ \emph{局部 Noether}，
（\hyperref[III.2.2.1-fr]{2.2.1}）的断言 (i) 仍成立。事实上，该性质显然
关于 $Y$ 是局部的；另一方面，（\hyperref[III.2.2.1-fr]{2.2.1}）的假设
蕴含：对每个开集 $U\subset Y$，$f$ 在 $f^{-1}(U)$ 上的限制是射影态射
$f^{-1}(U)\to U$（II，5.5.5，(iii)），且
$\mathcal L|f^{-1}(U)$ 对这一态射丰沛（II，4.6.4）。

(ii) 如前所见，若只假设 $X$ 是拟紧概形，或是底空间为 Noether 空间的
预概形，而 $f:X\to Y$ 是拟紧态射，则
（\hyperref[III.2.2.1-fr]{2.2.1}）的断言 (iii) 仍成立
（II，4.6.8）。但应注意，即使假设 $Y$ 是一个\emph{域} $K$ 的谱，
且 $f$ \emph{拟射影}，断言（\hyperref[III.2.2.1-fr]{2.2.1}，(ii)）
也未必成立。例如，设 $X'=\operatorname{Spec}(K[T_0,\ldots,T_r])$，
并令 $X$ 为 $X'$ 的诸仿射开集 $D(T_i)$ $(0\leq i\leq r)$ 之并。
由于浸入 $X\to X'$ 是拟紧的，结构态射 $f:X\to Y$ 是拟仿射的
（II，5.1.10），故 $\mathcal O_X$ 对 $f$ 极丰沛（II，5.1.6）。
但环 $\Gamma(X,\mathcal O_X)$ 可认同于诸分式环
$(K[T_0,\ldots,T_r])_{T_i}$（$0\leq i\leq r$）的交
（I，8.2.1.1），即 $K[T_0,\ldots,T_r]$。因此，由公式（1.4.3.1）和
（1.1.3.5），对每个 $n$ 都有
$H^r(X,\mathcal O_X^{\otimes n})=H^r(X,\mathcal O_X)\neq0$。
\SourceCorrectionNote{Canon 后继处置已确认 EGA3-ZH-DEF-0016：印本与外交式
法文作“$=A\neq0$”，但所引公式给出的是非零的最高次数局部上同调直极限模，
并非多项式环 $A$。本译按最小改动仅删去“$=A$”；法文权威保持不变。}

(iii) 在（2.2.4）关于 $X,Y,f,\mathcal L$ 的假设下，显然，若
$\mathcal H$ 是支撑在 $Y$ 上\emph{紧合}的凝聚 $\mathcal O_X$-模，
则与（2.2.4）类似的论证表明，存在整数 $N$，使得当 $n\geq N$ 时，
对所有 $i>0$ 均有 $R^if_*(\mathcal H(n))=0$。由上同调正合列可知，若
$u:\mathcal F\to\mathcal G$ 是凝聚 $\mathcal O_X$-模的同态，而
$\operatorname{Ker}(u)$ 和 $\operatorname{Coker}(u)$ 的支撑在 $Y$ 上
\emph{紧合}，则存在 $N$，使得当 $n\geq N$ 时，对所有 $i>0$，相应同态
$R^if_*(\mathcal F(n))\to R^if_*(\mathcal G(n))$ 都是\emph{双射}。
\SourceCorrectionNote{本段是已发表勘误 $\mathbf{Err}_{\mathrm{III}},27$
向（III，2.2.5）增补的第 (iii) 项；外交式法文正文保持不变。}
\end{env}

\subsection{分次代数层与模层的应用}
\label{subsection:III.2.3-fr}

\begin{theorem}[2.3.1]
\label{III.2.3.1-fr}
设 $Y$ 是 Noether 预概形，$\mathcal S$ 是正次数分次、拟相干且有限型的
$\mathcal O_Y$-代数，$X=\operatorname{Proj}(\mathcal S)$，
$q:X\to Y$ 是结构态射，$\mathcal M$ 是满足条件 \textbf{(TF)} 的拟相干
分次 $\mathcal S$-模。则存在整数 $N$，使得当 $n\geq N$ 时，典范同态
（II，8.14.5.1）
\[
\alpha_n:\mathcal M_n\longrightarrow
q_*(\shProj_0(\mathcal M(n)))
=q_*((\shProj(\mathcal M))_n)
\]
\oldpage[III]{103}
是双射。换言之，典范同态
\[
\alpha:\mathcal M\longrightarrow\bbGamma_*(\shProj(\mathcal M))
\]
是一个 \textbf{(TN)}-同构。
\end{theorem}

\begin{proof}
只需考虑 $\mathcal M$ 是有限型 $\mathcal S$-模的情形（II，3.4.2）。
由于 $Y$ 拟紧，存在整数 $d>0$，使 $\mathcal S^{(d)}$ 由拟相干
$\mathcal O_Y$-模 $\mathcal S_d$ 生成；后者是有限型的
（II，3.1.10），故因 $Y$ Noether 而凝聚。现在注意，$\mathcal M$ 是
诸 $\mathcal M^{(d,k)}$（$0\leq k<d$）的直和，而且每个
$\mathcal M^{(d,k)}$ 都是有限型拟相干 $\mathcal S^{(d)}$-模；由于问题
在 $Y$ 上是局部的，这由（II，2.1.6，(iii)）得出。显然，只须证明每个典范同态
\[
\alpha:\mathcal M^{(d,k)}
\longrightarrow\bbGamma_*((\shProj(\mathcal M))^{(d,k)})
\]
都是 \textbf{(TN)}-同构。考虑（II，8.14.13），特别是图表（8.14.13.4），
可见问题归结为：$\mathcal S$ 由 $\mathcal S_1$ 生成，且 $\mathcal S_1$
是凝聚 $\mathcal O_Y$-模的情形。

由于 $Y$ Noether，与（\hyperref[III.2.2.2-fr]{2.2.2}）开头相同的论证表明，
只需考虑 $Y=\operatorname{Spec}(A)$、$\mathcal S=\widetilde S$、
$\mathcal M=\widetilde M$ 的情形，其中 $A$ 是 Noether 环，$S_1$ 是有限型
$A$-模，而 $M$ 是有限型分次 $S$-模。我们说明，只须对 $M=S$ 证明定理。
事实上，一般情形下有正合列
\[
L'\longrightarrow L\longrightarrow M\longrightarrow0,
\]
其中 $L$ 和 $L'$ 是形如 $S(m)$ 的分次模的有限直和。若结论对 $M=S$
成立，则它对 $M=S(m)$、从而对 $L$ 与 $L'$ 也成立。于是考虑交换图表
\[
\xymatrix{
\widetilde {L'_n}\ar[r]\ar[d]_{\alpha_n} &
\widetilde {L_n}\ar[r]\ar[d]^{\alpha_n} &
\widetilde {M_n}\ar[r]\ar[d]^{\alpha_n} & 0\\
q_*(\widetilde {L'}(n))\ar[r] &
q_*(\widetilde L(n))\ar[r] &
q_*(\widetilde M(n))\ar[r] & 0 .
}
\]
当 $n$ 充分大时，第二行由（\hyperref[III.2.2.3-fr]{2.2.3}）正合；第一行
也正合，而左边两条竖箭头是同构，故第三条竖箭头也是同构。

现在，为对 $M=S$ 证明定理，先设 $S=A[T_0,\ldots,T_r]$（$T_i$ 为
不定元）；此时断言正是（\hyperref[III.2.1.11-fr]{2.1.11}，(ii)）。
一般情形下，$S$ 可认同于环 $S'=A[T_0,\ldots,T_r]$ 对某个分次理想的
商，因而 $X$ 可认同于 $X'=\mathbf P_A^r$ 的闭子概形（II，2.9.2）。
若 $j$ 是典范单射 $X\to X'$，则 $j_*(\widetilde S(n))$ 正是
$\mathcal O_{X'}$-模 $(\shProj(\widetilde S))(n)$，这里把 $S$ 看作分次
$S'$-模；这立即由（II，2.8.7）得出。由于 $j_*(\widetilde S(n))$ 是满足
\textbf{(TF)} 的 $\mathcal O_{X'}$-模，由前述结果，当 $n$ 充分大时，
典范同态 $\alpha_n:S_n\to\Gamma(X',j_*(\widetilde S(n)))$ 是双射。
而
$\Gamma(X',j_*(\widetilde S(n)))=\Gamma(X,\widetilde S(n))$，故证明完成。
\end{proof}

\oldpage[III]{104}
\begin{corollary}[2.3.2]
\label{III.2.3.2-fr}
在（\hyperref[III.2.3.1-fr]{2.3.1}）的假设下，令
\[
\mathcal S_X=\bigoplus_{n\in\mathbb Z}\mathcal O_X(n),
\]
并设 $\mathcal F$ 是有限型拟相干分次 $\mathcal S_X$-模。则
$\bbGamma_*(\mathcal F)$ 满足条件 \textbf{(TF)}。
\end{corollary}

\begin{proof}
在（\hyperref[III.2.3.1-fr]{2.3.1}）的证明中已见，$X$ 同构于
$\operatorname{Proj}(\mathcal S^{(d)})$（II，3.1.8），因而在 $Y$ 上有限型
（II，3.4.1）。由（II，8.14.9），$\mathcal F$ 同构于一个形如
$\shProj(\mathcal M)$ 的分次 $\mathcal S_X$-模，其中 $\mathcal M$ 是
有限型拟相干分次 $\mathcal S$-模。由
（\hyperref[III.2.3.1-fr]{2.3.1}），$\bbGamma_*(\mathcal F)$ 与
$\mathcal M$ 是 \textbf{(TN)}-同构的，故满足 \textbf{(TF)}。
\end{proof}

\begin{env}[2.3.3]
\label{III.2.3.3-fr}
\emph{附论}。--- 设 $Y$ 是 Noether 预概形，$\mathcal S$ 是满足
（\hyperref[III.2.3.1-fr]{2.3.1}）条件的分次 $\mathcal O_Y$-代数，且
$X=\operatorname{Proj}(\mathcal S)$。令 $K_{\mathcal S}$ 为满足
\textbf{(TF)} 的拟相干分次 $\mathcal S$-模所成的 Abel 范畴，令
$K'_{\mathcal S}$ 为 $K_{\mathcal S}$ 中由满足 \textbf{(TN)} 的
$\mathcal S$-模组成的
子范畴；最后，令 $K_X$ 为有限型拟相干分次 $\mathcal S_X$-模
$\mathcal F$ 所成的范畴（由于 $\mathcal S_X$ 是周期的
（II，8.14.4 和 8.14.12），这等价于说各 $\mathcal F_i$ 都是凝聚
$\mathcal O_X$-模）。于是函子
\[
\mathcal M\longmapsto\shProj(\mathcal M)
\quad\text{在 }K_{\mathcal S},
\qquad
\mathcal F\longmapsto\bbGamma_*(\mathcal F)
\quad\text{在 }K_X
\]
由（II，8.14.8 和 8.14.10）以及
（\hyperref[III.2.3.2-fr]{2.3.2}），定义商范畴
$K_{\mathcal S}/K'_{\mathcal S}$（T，I，1.11）与范畴 $K_X$ 的一个等价
（T，I，1.2）。当 $\mathcal S$ 由 $\mathcal S_1$ 生成时，可把 $K_X$
换成凝聚 $\mathcal O_X$-模的范畴（II，8.14.12）。
\end{env}

\begin{proposition}[2.3.4]
\label{III.2.3.4-fr}
设 $Y$ 是 Noether 预概形。
\begin{enumerate}
\item[{\rm(i)}] 设 $\mathcal S$ 是正次数分次、拟相干且有限型的
$\mathcal O_Y$-代数。令 $X=\operatorname{Proj}(\mathcal S)$，并令
\[
\mathcal S_X=\shProj(\mathcal S)
=\bigoplus_{n\in\mathbb Z}\mathcal O_X(n).
\]
则 $\mathcal S_X$ 是周期分次 $\mathcal O_X$-代数（II，8.14.12），其齐次分量
$(\mathcal S_X)_n=\mathcal O_X(n)$ 是凝聚 $\mathcal O_X$-模；若 $d>0$
是 $\mathcal S_X$ 的一个周期，则 $(\mathcal S_X)_d=\mathcal O_X(d)$
是对 $Y$ 丰沛的可逆 $\mathcal O_X$-模。此外，典范同态
\[
\alpha:\mathcal S\longrightarrow\bbGamma_*(\mathcal S_X)
\]
是一个 \textbf{(TN)}-同构。
\item[{\rm(ii)}] 反过来，设 $q:X\to Y$ 是射影态射，而 $\mathcal S'$ 是
分次 $\mathcal O_X$-代数；其齐次分量 $\mathcal S'_n$ $(n\in\mathbb Z)$
为凝聚 $\mathcal O_X$-模，并且它有一个周期 $d>0$，使 $\mathcal S'_d$
是对 $q$ 丰沛的可逆 $\mathcal O_X$-模。则
\[
\mathcal S=\bigoplus_{n\geq0}q_*(\mathcal S'_n)
\]
是正次数分次、拟相干且有限型的 $\mathcal O_Y$-代数，并且存在一个
$Y$-同构
\[
r:X\xrightarrow{\sim}\operatorname{Proj}(\mathcal S)
\]
使 $r^*(\shProj(\mathcal S))$ 与 $\mathcal S'$ 同构（作为分次
$\mathcal O_X$-代数）。
\end{enumerate}
\end{proposition}

\begin{proof}
(i) 各项断言实际上都已证明；最后一项正是
（\hyperref[III.2.3.2-fr]{2.3.2}）的特例。$\mathcal S_X$ 的周期性已在
（II，8.14.14）中看到，而存在周期 $d>0$ 使 $\mathcal O_X(d)$ 可逆且对
$Y$ 丰沛，正是（II，4.6.18）。最后，对 $0\leq k<d$，
$(\mathcal S_X)^{(d,k)}$ 是有限型 $(\mathcal S_X)^{(d)}$-模
（II，8.14.14）；故由（II，2.1.6，(ii)），每个 $(\mathcal S_X)_n$
都是有限型拟相干 $\mathcal O_X$-模。问题是局部的，而 $\mathcal O_X$
凝聚，所以各 $(\mathcal S_X)_n$ 也凝聚。

(ii) 把周期 $d$ 换成它的一个倍数，便可假设
$\mathcal L=\mathcal S'_d$ 是相对于 $q$ \emph{极丰沛}的可逆
$\mathcal O_X$-模（II，4.6.11）。此外，按假设有
\[
\mathcal S'^{(d)}=\bigoplus_{n\in\mathbb Z}\mathcal L^{\otimes n}
\quad\text{按假设，因而}\quad
\mathcal S^{(d)}=\bigoplus_{n\geq0}q_*(\mathcal L^{\otimes n})~;
\]
已知（II，3.1.8 和 3.2.9），存在 $X'=\operatorname{Proj}(\mathcal S)$
到 $X''=\operatorname{Proj}(\mathcal S^{(d)})$ 的 $Y$-同构 $s$，使得
\oldpage[III]{105}
\[
s^*(\mathcal O_{X''}(n))=\mathcal O_{X'}(nd).
\]
因此，只要证明下述结果，就能建立 $Y$-同构 $X\xrightarrow{\sim}X'$：

\begin{env}[2.3.4.1]
\label{III.2.3.4.1-fr}
设 $Y$ 是 Noether 预概形，$q:X\to Y$ 是射影态射，$\mathcal L$ 是相对于
$q$ 极丰沛的可逆 $\mathcal O_X$-模。则
$\mathcal S=\bigoplus_{n\geq0}q_*(\mathcal L^{\otimes n})$ 是拟相干且
有限型的分次 $\mathcal O_Y$-代数，并且当 $n$ 充分大时
$\mathcal S_n=\mathcal S_1^n$；还存在 $Y$-同构
$r:X\xrightarrow{\sim}P=\operatorname{Proj}(\mathcal S)$，使
$\mathcal L=r^*(\mathcal O_P(1))$。
\end{env}

由于 $q$ 是射影态射，由（II，5.4.4 和 4.4.7）存在 $Y$-同构
$r':X\xrightarrow{\sim}P'=\operatorname{Proj}(\mathcal T)$，其中
$\mathcal T$ 是拟相干 $\mathcal O_Y$-代数，$\mathcal T_1$ 有限型并生成
$\mathcal T$，且 $\mathcal L=r'^*(\mathcal O_{P'}(1))$。若
$q':P'\to Y$ 是结构态射，则
\[
\mathcal S=\bigoplus_{n\geq0}q'_*(\mathcal O_{P'}(n)).
\]
由（\hyperref[III.2.3.1-fr]{2.3.1}），当 $n$ 充分大时，典范同态
$\alpha_n:\mathcal T_n\to\mathcal S_n=q'_*(\mathcal O_{P'}(n))$ 是双射；
由于 $\mathcal T_n=\mathcal T_1^n$，当 $n$ 充分大时，
\emph{更是}有 $\mathcal S_n=\mathcal S_1^n$。此外，分次
$\mathcal O_Y$-代数的典范同态 $\alpha:\mathcal T\to\mathcal S$ 是一个
\textbf{(TN)}-同构；因而
\phantomsection\label{III.2.3.5.1-fr}
\[
\Phi=\operatorname{Proj}(\alpha):
\operatorname{Proj}(\mathcal S)\longrightarrow
\operatorname{Proj}(\mathcal T)
\]
是同构（II，3.6.1）。又由（II，3.5.2），并注意由诸
$\mathcal S(n)$ 限制标量得到的分次 $\mathcal T$-模与 $\mathcal T(n)$
是 \textbf{(TN)}-同构的，故
\[
\Phi_*(\mathcal O_P(n))=\mathcal O_{P'}(n)
\]
对每个 $n$ 都成立（II，3.4.2）。为完成
（\hyperref[III.2.3.4.1-fr]{2.3.4.1}）的证明，只须说明 $\mathcal S$
是有限型 $\mathcal O_Y$-代数。由（\hyperref[III.2.2.1-fr]{2.2.1}），
$\mathcal S_n=q'_*(\mathcal O_{P'}(n))$ 是凝聚 $\mathcal O_Y$-模；若
对 $n\geq n_0$ 有 $\mathcal S_n=\mathcal S_1^n$，则 $\mathcal S$ 由
凝聚模 $\bigoplus_{i\leq n_0}\mathcal S_i$ 生成；断言即由此得出
（I，9.6.2）。

回到（\hyperref[III.2.3.4-fr]{2.3.4}）的证明，并沿用其记号。我们已证明
存在 $Y$-同构 $r'':X\xrightarrow{\sim}X''$，使得对每个
$n\in\mathbb Z$，有 $r''_*(\mathcal L^{\otimes n})=\mathcal O_{X''}(n)$；
记 $q''$ 为结构态射 $X''\to Y$。现在注意，$\mathcal S'$ 是诸分次
$\mathcal S'^{(d)}$-模 $\mathcal S'^{(d,k)}$ 的直和；由于 $\mathcal S'$
的周期性以及各 $\mathcal S'_n$ 是有限型 $\mathcal O_X$-模的假设，
每个后者都是有限型拟相干 $\mathcal S'^{(d)}$-模
（II，8.14.12）。置
$\mathcal F^{(k)}=r''_*(\mathcal S'^{(d,k)})$，于是各
$\mathcal F^{(k)}$ 都是有限型拟相干分次 $\mathcal S_{X''}$-模；故由
（II，8.14.8），典范同态
\[
\beta:\shProj(\bbGamma_*(\mathcal F^{(k)}))
\longrightarrow\mathcal F^{(k)}
\]
是 $\mathcal S_{X''}$-模的同构。但
$q''_*((\mathcal F^{(k)})_n)=q_*((\mathcal S'^{(d,k)})_n)$，而当
$n\geq0$ 时，后一个 $\mathcal O_Y$-模按定义等于
$(\mathcal S^{(d,k)})_n$。换言之，典范单射
$\mathcal S^{(d,k)}\to\bbGamma_*(\mathcal F^{(k)})$ 是
\textbf{(TN)}-同构；因此
\[
\shProj(\mathcal S^{(d,k)})
=\shProj(\bbGamma_*(\mathcal F^{(k)})),
\]
从而
$r''{}^*(\shProj(\mathcal S^{(d,k)}))\simeq\mathcal S'^{(d,k)}$。
最后只须注意，在典范同构意义下，
\[
\shProj(\mathcal S^{(d,k)})
=s_*((\shProj(\mathcal S))^{(d,k)})
\]
（II，8.14.13.1），便证明了 $r^*(\shProj(\mathcal S))$ 与
$\mathcal S'$ 同构。

最后，由（\hyperref[III.2.3.2-fr]{2.3.2}），每个
$\bbGamma_*(\mathcal F^{(k)})$ 都满足条件 \textbf{(TF)}，所以每个
$\mathcal S^{(d,k)}$ 也如此；又由于各 $\mathcal S'_n$ 凝聚，由
（\hyperref[III.2.2.1-fr]{2.2.1}），$\mathcal S_n=q_*(\mathcal S'_n)$
也凝聚。立即可知各 $\mathcal S^{(d,k)}$ 是有限型
$\mathcal S^{(d)}$-模；而（\hyperref[III.2.3.4.1-fr]{2.3.4.1}）表明
$\mathcal S^{(d)}$ 是有限型 $\mathcal O_Y$-代数，故 $\mathcal S$
也如此。
\end{proof}

\oldpage[III]{106}
\begin{proposition}[2.3.5]
\label{III.2.3.5-fr}
设 $Y$ 是整 Noether 预概形，$X$ 是整预概形，$f:X\to Y$ 是射影双有理
态射。则存在凝聚分式理想层 $\mathcal J\subset\mathcal R(Y)$
（II，8.1.2），使 $X$ 与沿 $\mathcal J$ 暴涨得到的预概形
（II，8.1.3）在 $Y$ 上同构。此外，存在 $Y$ 的开集 $U$，使 $f$ 在
$f^{-1}(U)$ 上的限制是 $f^{-1}(U)$ 到 $U$ 的同构（参见 I，6.5.5），
且 $\mathcal J|U$ 可逆。
\end{proposition}

\begin{proof}
由于存在一个相对于 $f$ 极丰沛的可逆 $\mathcal O_X$-模 $\mathcal L$
（II，4.4.2 和 5.3.2），可以应用
（\hyperref[III.2.3.4.1-fr]{2.3.4.1}）；于是 $X$ 可认同于
$\operatorname{Proj}(\mathcal S)$，其中
$\mathcal S=\bigoplus_{n\geq0}f_*(\mathcal L^{\otimes n})$。又已知
$f_*(\mathcal L^{\otimes n})$ 是无挠 $\mathcal O_Y$-模（I，7.4.5），
所以 $\mathcal O_Y$-模 $\mathcal S$ 也是无挠的，从而 $\mathcal S$
可典范地认同于
$\mathcal S\otimes_{\mathcal O_Y}\mathcal R(Y)$ 的一个
$\mathcal O_Y$-子模（I，7.4.1）。后者是一个\emph{简单}层（I，7.3.6），
只要知道它在某个非空开集上的限制，就能确定它；例如可取非空开集
$U'\subset U$，使 $\mathcal L|f^{-1}(U')$ 同构于
$\mathcal O_X|f^{-1}(U')$。按假设，此时
$f_*(\mathcal L^{\otimes n})|U'$ 同构于 $\mathcal O_Y|U'$，因而
$\mathcal S\otimes_{\mathcal O_Y}\mathcal R(Y)$ 是同构于
$\mathcal R(Y)[T]$ 的 $\mathcal R(Y)$-模，其中 $T$ 是不定元；并且
$\mathcal S$ 与由 $f_*(\mathcal L)$ 在
$\mathcal S\otimes_{\mathcal O_Y}\mathcal R(Y)$ 中的典范像所生成的
$\mathcal O_Y$-子代数是 \textbf{(TN)}-同构的
（\hyperref[III.2.3.4.1-fr]{2.3.4.1}）。若把后一个层认同于
$\mathcal R(Y)[T]$，则 $f_*(\mathcal L)$ 的像可认同于
$\mathcal J T$，其中 $\mathcal J$ 是 $\mathcal R(Y)$ 的凝聚
$\mathcal O_Y$-子模（\hyperref[III.2.2.1-fr]{2.2.1}）；它在 $U'$
上的限制同构于 $\mathcal O_Y|U'$，因而 $\mathcal J|U$ 可逆。于是
$\mathcal S$ 与 $\bigoplus_{n\geq0}\mathcal J^n$ 是
\textbf{(TN)}-同构的，证明完成。
\end{proof}

\begin{corollary}[2.3.6]
\label{III.2.3.6-fr}
在（\hyperref[III.2.3.5-fr]{2.3.5}）的假设下，再假设对
$\mathcal R(Y)$ 的每个非零凝聚 $\mathcal O_Y$-子模
$\mathcal J\neq0$，都存在可逆 $\mathcal O_Y$-模 $\mathcal L$，使
\[
\Gamma(Y,\mathcal L\otimes_{\mathcal O_Y}
\mathcal{H}om(\mathcal J,\mathcal O_Y))\neq0~;
\]
则在（\hyperref[III.2.3.5-fr]{2.3.5}）的陈述中，可假设
$\mathcal J$ 是 $\mathcal O_Y$ 的理想层。若存在一个丰沛
$\mathcal O_Y$-模，这一附加条件总是成立。
\end{corollary}

\begin{proof}
事实上，有（0\textsubscript{I}，5.4.2）
\[
\mathcal L\otimes\mathcal{H}om(\mathcal J,\mathcal O_Y)
=\mathcal{H}om(\mathcal L^{-1},\mathcal{H}om(\mathcal J,\mathcal O_Y))
=\mathcal{H}om(\mathcal J\otimes\mathcal L^{-1},\mathcal O_Y)~;
\]
因而假设意味着存在非零同态 $u$，把
$\mathcal J\otimes\mathcal L^{-1}$ 映入 $\mathcal O_Y$。对每个
$(y\in Y)$，$(\mathcal J\otimes\mathcal L^{-1})_y$ 可认同于
$\mathcal O_y$ 的分式域 $\mathcal R(Y)_y$ 的一个
$\mathcal O_y$-子模（I，7.1.5），故 $u_y$ 必为单射；于是 $u$ 是
$\mathcal J\otimes\mathcal L^{-1}$ 到 $\mathcal O_Y$ 的某个理想层
$\mathcal J'$ 上的同构。但
\[
\operatorname{Proj}\left(\bigoplus_{n\geq0}\mathcal J^n\right)
\quad\text{和}\quad
\operatorname{Proj}\left(
\bigoplus_{n\geq0}(\mathcal J\otimes\mathcal L^{-1})^n\right)
\]
在 $Y$ 上同构（II，3.1.8），这就证明了推论的第一项断言。为证明第二项，
注意 $\mathcal F=\mathcal{H}om_{\mathcal O_Y}(\mathcal J,\mathcal O_Y)$
是凝聚且 $\neq0$ 的，因为存在 $Y$ 的开集 $U$，使 $\mathcal J|U$ 可逆。
若 $\mathcal L$ 是丰沛 $\mathcal O_Y$-模，则存在整数 $n$，使
$\mathcal F(n)=\mathcal F\otimes\mathcal L^{\otimes n}$ 由它在 $Y$
上的截面生成（II，4.5.5）；\emph{更是}有
$\Gamma(Y,\mathcal F(n))\neq0$，结论随即得出。
\end{proof}

\begin{corollary}[2.3.7]
\label{III.2.3.7-fr}
设 $X$ 和 $Y$ 是域 $k$ 上的两个整射影概形，$f:X\to Y$ 是双有理
$k$-态射。则 $X$ 在 $k$ 上同构于一个 $Y$-概形，后者由沿 $Y$
的某个闭子概形 $Y'$（不必约化）暴涨而得。
\end{corollary}

\oldpage[III]{107}
\begin{proof}
事实上，$f$ 是射影的（II，5.5.5，(v)）；又因 $Y$ 在 $k$ 上射影，
（\hyperref[III.2.3.6-fr]{2.3.6}）的附加条件成立。于是只须考虑由推论
（\hyperref[III.2.3.6-fr]{2.3.6}）中的凝聚理想层 $\mathcal J$ 所定义的
$Y$ 的闭子概形 $Y'$。
\end{proof}

\begin{env}[2.3.8]
\label{III.2.3.8-fr}
\emph{注记}。第四章研究除子概念时将会看到：若在
（\hyperref[III.2.3.5-fr]{2.3.5}）的陈述中假设诸环
$\mathcal O_y$（$y\in Y$）是唯一分解整环（例如 $Y$ 非奇异时即如此），
则 $X$ 可从 $Y$ 得到：沿 $Y$ 的某个闭子预概形 $Y'$ 作暴涨；其底空间
包含在 $Y-U$ 中。
\end{env}

\subsection{基本定理的一个推广}
\label{subsection:III.2.4-fr}

\begin{theorem}[2.4.1]
\label{III.2.4.1-fr}
设 $Y$ 是 Noether 预概形，$\mathcal S$ 是有限型拟相干
$\mathcal O_Y$-代数。设 $f:X\to Y$ 是射影态射，
$\mathcal S'=f^*(\mathcal S)$，而 $\mathcal M$ 是有限型拟相干
$\mathcal S'$-模。则：
\begin{enumerate}
\item[(i)] 对每个 $p\in\mathbb Z$，$R^pf_*(\mathcal M)$ 都是有限型
$\mathcal S$-模。
\item[(ii)] 再设 $\mathcal L$ 是对 $f$ 丰沛的可逆
$\mathcal O_X$-模，并对每个 $n\in\mathbb Z$ 置
$\mathcal M(n)=\mathcal M\otimes\mathcal L^{\otimes n}$。则存在整数
$N$，使得当 $n\geq N$ 时，对每个 $p>0$ 都有
\[
R^pf_*(\mathcal M(n))=0
\tag{2.4.1.1}\label{III.2.4.1.1-fr}
\]
并且典范同态
\[
f^*(f_*(\mathcal M(n)))\longrightarrow\mathcal M(n)
\]
（0\textsubscript{I}，4.4.3）是满射。
\end{enumerate}
\end{theorem}

\begin{proof}
置 $Y'=\operatorname{Spec}(\mathcal S)$、
$X'=\operatorname{Spec}(\mathcal S')$，于是
$X'=X\times_Y Y'$（II，1.5.5）。设 $g:Y'\to Y$、
$g':X'\to X$ 为结构态射；按定义它们是仿射的。再设
$f'=f_{(Y')}:X'\to Y'$。于是有交换图
\[
\xymatrix{
X\ar[d]_f & X'\ar[l]_{g'}\ar[d]^{f'}\ar[dl]^{h}\\
Y & Y'\ar[l]^{g}
}
\]
且态射 $f'$ 是射影的（II，5.5.5，(iii)）；置
$h=f\circ g'=g\circ f'$。

(i) 把 $X'$ 看成 $X$ 上的仿射概形，并令
$\widetilde{\mathcal M}$ 为与拟相干 $\mathcal S'$-模
$\mathcal M$ 相伴的 $\mathcal O_{X'}$-模（II，1.4.3），于是
$\mathcal M=g'_*(\widetilde{\mathcal M})$。由于 $\mathcal M$ 是有限型
$\mathcal S'$-模，$\widetilde{\mathcal M}$ 是有限型
$\mathcal O_{X'}$-模（II，1.4.5）。又因 $g$ 与 $f'$ 都是有限型的
（II，1.3.7 和 I，6.3.4，(ii)），故 $h$ 是有限型的；因而 $X'$ 是
Noether 的（I，6.3.7），从而 $\widetilde{\mathcal M}$ 是凝聚的。
另一方面，由于 $g'$ 是仿射的，典范同态
$R^pf_*(\mathcal M)\to R^ph_*(\widetilde{\mathcal M})$ 是双射
（\hyperref[III.1.3.4-fr]{1.3.4}）。此外，它还是
$\mathcal S$-模同态。事实上，由定义 $\mathcal M$ 的
$\mathcal S'$-模结构的典范同态
\[
g'_*(\mathcal O_{X'})\otimes_{\mathcal O_X}
g'_*(\widetilde{\mathcal M})\longrightarrow
g'_*(\widetilde{\mathcal M})
\tag{2.4.1.2}\label{III.2.4.1.2-fr}
\]
（这里要记住 $\mathcal S'=g'_*(\mathcal O_{X'})$），可典范地导出同态
\[
f_*(g'_*(\mathcal O_{X'}))\otimes
R^pf_*(g'_*(\widetilde{\mathcal M}))\longrightarrow
R^pf_*(g'_*(\widetilde{\mathcal M})).
\]

\oldpage[III]{108}
（0，12.2.2）。而（\hyperref[III.2.4.1.2-fr]{2.4.1.2}）本身由定义
$\widetilde{\mathcal M}$ 的 $\mathcal O_{X'}$-模结构的同态
\[
\mathcal O_{X'}\otimes\widetilde{\mathcal M}
\longrightarrow\widetilde{\mathcal M}
\]
经应用（0\textsubscript{I}，4.2.2.1）而来，故图
\[
\xymatrix{
f_*(g'_*(\mathcal O_{X'}))\otimes
R^pf_*(g'_*(\widetilde{\mathcal M}))
\ar[r]\ar[d] &
R^pf_*(g'_*(\widetilde{\mathcal M}))\ar[d]\\
h_*(\mathcal O_{X'})\otimes R^ph_*(\widetilde{\mathcal M})
\ar[r] & R^ph_*(\widetilde{\mathcal M})
}
\]
是交换的（0，12.2.6）。把两条水平箭头分别与由典范同态
\[
\mathcal S\longrightarrow f_*(f^*(\mathcal S))=f_*(\mathcal S')
=f_*(g'_*(\mathcal O_{X'}))=h_*(\mathcal O_{X'}),
\]
所导出的同态复合，便得到所断言的结论。另一方面，由于 $g$ 是仿射的，
而 $f'$ 是分离且拟紧的，典范同态
$R^ph_*(\widetilde{\mathcal M})\to
g_*(R^pf'_*(\widetilde{\mathcal M}))$ 是双射
（\hyperref[III.1.4.14-fr]{1.4.14}）；与上面相同的论证表明，它是
$\mathcal S$-模同构（这次使用（0，12.2.6.2）的交换性）。现在，
$f'$ 是射影的，且 $\widetilde{\mathcal M}$ 是凝聚的，所以由
（\hyperref[III.2.2.1-fr]{2.2.1}），
$R^pf'_*(\widetilde{\mathcal M})$ 是凝聚
$\mathcal O_{Y'}$-模；从而 $g_*(R^pf'_*(\widetilde{\mathcal M}))$
是有限型 $\mathcal S$-模（II，1.4.5）。

(ii) 置 $\mathcal L'=g'^*(\mathcal L)$，这是可逆
$\mathcal O_{X'}$-模。对每个 $n\in\mathbb Z$，在相差一个同构的意义下有
\[
g'_*(\widetilde{\mathcal M}\otimes\mathcal L'^{\otimes n})
=g'_*(\widetilde{\mathcal M})\otimes\mathcal L^{\otimes n}
=\mathcal M(n)
\]
（0\textsubscript{I}，5.4.10）。可以对
$\widetilde{\mathcal M}\otimes\mathcal L'^{\otimes n}$ 应用 (i) 中对
$\widetilde{\mathcal M}$ 所作的论证；它证明
\[
R^pf_*(g'_*(\widetilde{\mathcal M}\otimes\mathcal L'^{\otimes n}))
\quad\hbox{同构于}\quad
g_*(R^pf'_*(\widetilde{\mathcal M}\otimes\mathcal L'^{\otimes n})).
\]
又 $\mathcal L'$ 对 $f'$ 丰沛（II，4.6.13，(iii)），所以由
（\hyperref[III.2.2.1-fr]{2.2.1}）可知，存在整数 $N$，使得
\[
R^pf'_*(\widetilde{\mathcal M}\otimes\mathcal L'^{\otimes n})=0
\]
对所有 $p>0$ 以及所有 $n\geq N$ 都成立
\SourceCorrectionNote{法文印本在这里漏去 $p>0$ 的正次数限制；
Canon 稳定决定 EGA-III1-2.4.1-P452-P-POSITIVE-001 核定消失只对
$p>0$ 断言。本译采用已核定读法，并保留可逆来源说明。}，这就证明了
（\hyperref[III.2.4.1.1-fr]{2.4.1.1}）。最后，再由
（\hyperref[III.2.2.1-fr]{2.2.1}）可知，可假定 $N$ 还满足：当
$n\geq N$ 时，典范同态
\[
f'^*(f'_*(\widetilde{\mathcal M}\otimes\mathcal L'^{\otimes n}))
\longrightarrow\widetilde{\mathcal M}\otimes\mathcal L'^{\otimes n}
\]
是满射。由于 $g'_*$ 是正合函子（II，1.4.4），相应同态
\[
g'_*(f'^*(f'_*(\widetilde{\mathcal M}\otimes
\mathcal L'^{\otimes n})))
\longrightarrow
g'_*(\widetilde{\mathcal M}\otimes\mathcal L'^{\otimes n})
=\mathcal M(n)
\]
也是满射。又有
\[
g'_*(f'^*(f'_*(\widetilde{\mathcal M}\otimes
\mathcal L'^{\otimes n})))
=f^*(g_*(f'_*(\widetilde{\mathcal M}\otimes
\mathcal L'^{\otimes n})))
\]
（II，1.5.2）；并且由于 $g_*\circ f'_* = f_*\circ g'_*$，最终得到
\[
g'_*(f'^*(f'_*(\widetilde{\mathcal M}\otimes\mathcal L'^{\otimes n})))
=f^*(f_*(g'_*(\widetilde{\mathcal M}\otimes
\mathcal L'^{\otimes n})))=f^*(f_*(\mathcal M(n))),
\]
证明完成。
\end{proof}

\begin{env}[2.4.2]
\label{III.2.4.2-fr}
特别地，我们将把（\hyperref[III.2.4.1-fr]{2.4.1}）应用于如下情形：
$\mathcal S$ 是正次数分次 $\mathcal O_Y$-代数，且
$\mathcal M=\bigoplus_{k\in\mathbb Z}\mathcal M_k$ 是分次
$\mathcal S'$-模。此时，在 $\mathcal S$ 与 $\mathcal M$ 满足同样有限性

\oldpage[III]{109}
假设的条件下，由（\hyperref[III.2.4.1-fr]{2.4.1}）可知，对每个 $p$，
\[
\bigoplus_{k\in\mathbb Z}R^pf_*(\mathcal M_k)
\]
都是有限型 $\mathcal S$-模；而在
（\hyperref[III.2.4.1-fr]{2.4.1}，(ii)）的附加假设下，存在 $N$，
使得当 $n\geq N$ 时，对每个 $p>0$ 和每个 $k\in\mathbb Z$，都有
$R^pf_*(\mathcal M_k(n))=0$，并且对每个 $k\in\mathbb Z$，典范同态
$f^*(f_*(\mathcal M_k(n)))\to\mathcal M_k(n)$ 是满射。
\end{env}

\subsection{Euler--Poincaré 特征与 Hilbert 多项式}
\label{subsection:III.2.5-fr}

\begin{env}[2.5.1]
\label{III.2.5.1-fr}
设 $A$ 是 Artin 环，$X$ 是在 $Y=\operatorname{Spec}(A)$ 上射影的
$A$-概形。对每个凝聚 $\mathcal O_X$-模 $\mathcal F$，
$H^i(X,\mathcal F)$ $(i\geq0)$ 都是有限型 $A$-模
（\hyperref[III.2.2.1-fr]{2.2.1}）；由于 $A$ 是 Artin 环，它们在这里都具有
有限长度。此外，已知（\hyperref[III.2.2.1-fr]{2.2.1}），除有限多个
$i\geq0$ 外，均有 $H^i(X,\mathcal F)=0$。因此，对每个凝聚
$\mathcal O_X$-模 $\mathcal F$，整数
\[
\chi_A(\mathcal F)=\sum_{i=0}^{\infty}(-1)^i
\operatorname{long}(H^i(X,\mathcal F))
\tag{2.5.1.1}\label{III.2.5.1.1-fr}
\]
都有定义。当 $A$ 是 Artin 局部环时，称 $\chi_A(\mathcal F)$ 为
$\mathcal F$（相对于环 $A$）的 Euler--Poincaré 特征。对
$\mathcal F=\mathcal O_X$，称 $\chi_A(\mathcal O_X)$ 为 $X$
（相对于 $A$）的算术亏格。
\end{env}

\begin{proposition}[2.5.2]
\label{III.2.5.2-fr}
设
\[
0\longrightarrow\mathcal F'\longrightarrow\mathcal F
\longrightarrow\mathcal F''\longrightarrow0
\]
是凝聚 $\mathcal O_X$-模的正合列，则有
\[
\chi_A(\mathcal F)=\chi_A(\mathcal F')+\chi_A(\mathcal F'').
\tag{2.5.2.1}\label{III.2.5.2.1-fr}
\]
\end{proposition}

\begin{proof}
由于 $\mathcal F$、$\mathcal F'$、$\mathcal F''$ 的上同调模除有限多个外
均为零，存在整数 $r>0$，使上同调正合列可写为
\[
0\to H^0(X,\mathcal F')\to H^0(X,\mathcal F)
\to H^0(X,\mathcal F'')\to H^1(X,\mathcal F')\to\cdots
\]
\[
\cdots\to H^r(X,\mathcal F')\to H^r(X,\mathcal F)
\to H^r(X,\mathcal F'')\to0.
\]
已知，在两端均为 0 的有限长度 $A$-模正合列中，各项长度的交错和为零
（0，11.10.1）。应用这一结果，立即得到公式
（\hyperref[III.2.5.2.1-fr]{2.5.2.1}）。
\end{proof}

注意，只要已知 $X$ 是拟紧 $A$-概形，且对每个凝聚
$\mathcal O_X$-模 $\mathcal F$，诸 $A$-模 $H^i(X,\mathcal F)$ 均为有限型，
便可应用（\hyperref[III.2.5.2-fr]{2.5.2}）的结果
（\hyperref[III.1.4.12-fr]{1.4.12}）。

\begin{theorem}[2.5.3]
\label{III.2.5.3-fr}
设 $A$ 是 Artin 局部环，$X$ 是 $Y=\operatorname{Spec}(A)$ 上的射影
概形，$\mathcal L$ 是相对于 $Y$ 极丰沛的可逆 $\mathcal O_X$-模，
$\mathcal F$ 是凝聚 $\mathcal O_X$-模。对每个 $n\in\mathbb Z$，置
$\mathcal F(n)=\mathcal F\otimes_{\mathcal O_X}\mathcal L^{\otimes n}$。
\begin{enumerate}
\item[(i)] 存在唯一多项式 $P\in\mathbb Q[T]$，使得对每个
$n\in\mathbb Z$，都有 $\chi_A(\mathcal F(n))=P(n)$
（称 $P$ 为 $\mathcal F$ 相对于 $A$ 的 Hilbert 多项式）。
\item[(ii)] 当 $n$ 充分大时，有
$\chi_A(\mathcal F(n))=\operatorname{long}_A
\Gamma(X,\mathcal F(n))$。
\item[(iii)] $\chi_A(\mathcal F(n))$ 的最高次项系数 $\geq0$。
\end{enumerate}
\end{theorem}

补充指出，在第四章讨论维数概念的那一节中，我们还将证明：
$\chi_A(\mathcal F(n))$ 的次数等于 $\mathcal F$ 的支撑的维数。

\begin{proof}
由（\hyperref[III.2.2.1-fr]{2.2.1}），只要 $n$ 充分大，对每个 $i>0$
都有 $H^i(X,\mathcal F(n))=0$，

\oldpage[III]{110}
因而当 $n$ 充分大时，有
\[
\chi_A(\mathcal F(n))=\operatorname{long}H^0(X,\mathcal F(n))
=\operatorname{long}\Gamma(X,\mathcal F(n))
\]
这就得到 (ii)。它还表明，当 $n$ 充分大时，
$\chi_A(\mathcal F(n))\geq0$，所以 (iii) 由 (i) 得出。另一方面，(i)
中的唯一性断言是显然的，故只须证明多项式 $P$ 的存在性。

先证明可以假设 $\mathfrak m\mathcal F=0$，其中 $\mathfrak m$ 是 $A$
的极大理想。事实上，存在整数 $s>0$ 使 $\mathfrak m^s=0$，所以
$\mathcal F(n)$ 有有限滤过
\[
\mathcal F(n)\supset\mathfrak m\mathcal F(n)\supset\cdots
\supset\mathfrak m^{s-1}\mathcal F(n)\supset0.
\]
由（\hyperref[III.2.5.2.1-fr]{2.5.2.1}）归纳可得
\SourceCorrectionNote{法文印本此处误引（2.5.1）；Canon 稳定决定
EGA-III1-2.5.3-P453-P454-CITATION-001 核定此处应引用
（2.5.2.1）。本译采用已核定读法，并保留可逆来源说明。}
\[
\chi_A(\mathcal F(n))=\sum_{k=1}^{s}
\chi_A(\mathfrak m^{k-1}\mathcal F(n)/\mathfrak m^k\mathcal F(n))~;
\]
而
$\mathfrak m^{k-1}\mathcal F(n)/\mathfrak m^k\mathcal F(n)
=\mathcal F'_k(n)$，其中
$\mathcal F'_k=\mathfrak m^{k-1}\mathcal F/\mathfrak m^k\mathcal F$，
故断言得证。

于是设 $\mathfrak m\mathcal F=0$。若 $X'$ 是 $X$ 的闭子概形，它是
$\operatorname{Spec}(A)$ 的唯一闭点在结构态射
$X\to\operatorname{Spec}(A)$ 下的逆像，且 $j:X'\to X$ 是典范嵌入，
则有 $\mathcal F=j_*(\mathcal F')$，其中 $\mathcal F'$ 是凝聚
$\mathcal O_{X'}$-模。$X'$ 是 $\operatorname{Spec}(K)$ 上的射影概形，
其中 $K=A/\mathfrak m$。若置 $\mathcal L'=j^*(\mathcal L)$，则
$\mathcal L'$ 相对于 $\operatorname{Spec}(K)$ 极丰沛（II，4.4.10），并且
$\mathcal F(n)=j_*(\mathcal F'(n))$，其中
$\mathcal F'(n)=\mathcal F'\otimes_{\mathcal O_{X'}}
\mathcal L'^{\otimes n}$（0\textsubscript{I}，5.4.10）。由此得出
$\chi_A(\mathcal F(n))=\chi_K(\mathcal F'(n))$（G，II，4.9.1），
所以归结为 $A$ 是域的情形。

现在注意，选取适当的 $r$，可把 $X$ 看成 $P=\mathbf P_A^r$ 的闭子概形
（II，5.5.4，(ii)）。若 $i:X\to P$ 是典范嵌入，则与上面相同可得
$\chi_A(\mathcal F(n))=\chi_A(i_*(\mathcal F)(n))$；因此只须考虑
$X=\mathbf P_A^r=\operatorname{Proj}(S)$、$S=A[T_0,\ldots,T_r]$，
其中 $A$ 是域的情形。

此时有 $\mathcal F=\widetilde M$，其中 $M$ 是有限型分次 $S$-模
（II，2.7.8）。由 Hilbert 合冲定理（M，VIII，6.5），存在 $M$ 的有限型
分次自由 $S$-模有限分解
\[
0\to L_q\to L_{q-1}\to\cdots\to L_1\to M\to0
\]
由于 $M\mapsto\widetilde M$ 是关于 $M$ 的正合函子（II，2.5.4），还有正合列
\[
0\to\widetilde L_q\to\widetilde L_{q-1}\to\cdots
\to\widetilde L_1\to\widetilde M\to0
\]
从而对每个 $n\in\mathbb Z$，列
\[
0\to\widetilde L_q(n)\to\widetilde L_{q-1}(n)\to\cdots
\to\widetilde L_1(n)\to\widetilde M(n)\to0
\]
都是正合的。对 $q$ 归纳应用命题（2.5.2）
\SourceCorrectionNote{法文印本此处误引命题（2.5.1）；Canon 稳定决定
EGA-III1-2.5.3-P453-P454-CITATION-001 核定应引用命题（2.5.2）。
本译采用已核定读法，并保留可逆来源说明。}，得到
\[
\chi_A(\widetilde M(n))=\sum_{j=1}^{q}(-1)^{j+1}
\chi_A(\widetilde L_j(n))
\]
因此，为证明 (i)，只须考虑 $M$ 是有限型分次自由模的情形，进而只须考虑
$M=S(h)$、$h\in\mathbb Z$ 的情形。此时有
$\widetilde M(n)=(M(n))^\sim=(S(n+h))^\sim$（II，2.5.15），
所以最终可知，定理将由下面的引理得出。
\end{proof}

\oldpage[III]{111}
\begin{lemma}[2.5.3.1]
\label{III.2.5.3.1-fr}
设 $A$ 是域，$r$ 是 $>0$ 的整数，且 $X=\mathbf P_A^r$，则对每个
$n\in\mathbb Z$，都有
\[
\chi_A(\mathcal O_X(n))=\binom{n+r}{r}
\]
\end{lemma}

\begin{proof}
事实上，当 $n\geq0$ 时，有
$\chi_A(\mathcal O_X(n))=\operatorname{long}H^0(X,\mathcal O_X(n))$，
它等于诸 $T_i$ 的总次数为 $n$ 的单项式的数目，即
$\binom{n+r}{r}$（\hyperref[III.2.1.12-fr]{2.1.12}）。当
$n\leq-r-1$ 时，同样有
\[
\chi_A(\mathcal O_X(n))=(-1)^r
\operatorname{long}H^r(X,\mathcal O_X(n))~;
\]
若 $n=-r-h$，则 $H^r(X,\mathcal O_X(n))$ 在 $A$ 上的维数等于满足
$\sum_{i=0}^{r}p_i=r+h$ 的正整数 $p_i>0$ 的序列
$(p_i)_{0\leq i\leq r}$ 的数目（\hyperref[III.2.1.12-fr]{2.1.12}），
也就是满足 $\sum_{i=0}^{r}q_i=h-1$ 的非负整数 $q_i\geq0$
$(0\leq i\leq r)$ 的序列数目。因此，它等于
\[
\binom{h+r-1}{r}=(-1)^r\binom{n+r}{r}.
\]
最后，当 $-r\leq n<0$ 时，有 $\binom{n+r}{r}=0$；另一方面，对所有
$i\geq0$，都有 $H^i(X,\mathcal O_X(n))=0$
（\hyperref[III.2.1.12-fr]{2.1.12}）。引理得证。
\end{proof}

\begin{corollary}[2.5.4]
\label{III.2.5.4-fr}
设 $A$ 是 Artin 局部环，$S$ 是由 $S_1$ 生成的有限型分次 $A$-代数，
$M$ 是有限型分次 $S$-模，且 $X=\operatorname{Proj}(S)$。则当 $n$
充分大时，有
\[
\chi_A(\widetilde M(n))=\operatorname{long}M_n
\]
\end{corollary}

\begin{proof}
这是因为，当 $n$ 充分大时，$M_n$ 与 $\Gamma(X,\widetilde M(n))$
同构（\hyperref[III.2.3.1-fr]{2.3.1}）。
\end{proof}

\subsection{应用：丰沛性判别准则}
\label{subsection:III.2.6-fr}

\begin{proposition}[2.6.1]
\label{III.2.6.1-fr}
设 $Y$ 是 Noether 预概形，$f:X\to Y$ 是固有态射，
$\mathcal L$ 是可逆 $\mathcal O_X$-模。下列条件等价：
\begin{enumerate}
\item[a)] $\mathcal L$ 对 $f$ 丰沛。
\item[b)] 对每个凝聚 $\mathcal O_X$-模 $\mathcal F$，存在整数 $N$，
使得当 $n\geq N$ 时，对每个 $q>0$ 都有
\[
R^qf_*(\mathcal F\otimes\mathcal L^{\otimes n})=0
\]
\item[c)] 对 $\mathcal O_X$ 的每个凝聚理想 $\mathcal J$，存在整数 $N$，
使得当 $n\geq N$ 时有
\[
R^1f_*(\mathcal J\otimes\mathcal L^{\otimes n})=0.
\]
\end{enumerate}
\end{proposition}

\begin{proof}
已经证明 a) 蕴含 b)（\hyperref[III.2.2.1-fr]{2.2.1}，(ii)）
\SourceCorrectionNote{法文印本此处误引（2.2.1）(iii)；Canon 稳定决定
EGA-III1-2.6.1-P455-CITATION-PART-001 核定应为 (ii)。本译采用已核定
读法，并保留可逆来源说明。}。
b) 蕴含 c) 是显然的，故只须证明 c) 蕴含 a)。可以限于 $Y$ 为仿射
的情形（II，4.6.4），并在这一情形下证明 $\mathcal L$ 丰沛；只须证明：
当 $h$ 遍历诸 $\mathcal L^{\otimes n}$ $(n>0)$ 在 $X$ 上的截面的集合时，
诸 $X_h$ 中的仿射开集构成 $X$ 的一个覆盖（II，4.5.2）。为此，
证明对 $X$ 的每个闭点 $x$ 及 $x$ 的每个仿射开邻域 $U$，
都存在 $n$ 及 $h\in\Gamma(X,\mathcal L^{\otimes n})$，使
$x\in X_h\subset U$；此时 $X_h$ 必为仿射（I，1.3.6），而这些 $X_h$
的并是 $X$ 的一个开集，包含 $X$ 的所有闭点，因而就是 $X$ 本身，
因为 $X$ 是 Noether 的（I，6.3.7 及 0\textsubscript{I}，2.1.3）。
设 $\mathcal J$（分别地 $\mathcal J'$）是 $\mathcal O_X$ 的拟相干理想层，
它定义 $X$ 的一个约化闭子预概形，其底空间为 $X-U$
（分别地 $(X-U)\cup\{x\}$）（I，5.2.1）。显然，$\mathcal J$ 和
$\mathcal J'$ 都是凝聚的（I，6.1.1），$\mathcal J'\subset\mathcal J$，
且 $\mathcal J''=\mathcal J/\mathcal J'$ 是凝聚 $\mathcal O_X$-模
（0\textsubscript{I}，5.3.3），其支撑为 $\{x\}$，并满足
$\mathcal J''_x=k(x)$。由于 $\mathcal L$ 局部自由，列
\[
0\to\mathcal J'\otimes\mathcal L^{\otimes n}
\to\mathcal J\otimes\mathcal L^{\otimes n}
\to\mathcal J''\otimes\mathcal L^{\otimes n}\to0
\]
对每个 $n$ 都正合；而由假设，存在充分大的 $n$，使
$H^1(X,\mathcal J'\otimes\mathcal L^{\otimes n})=0$。上同调正合列

\oldpage[III]{112}
因此表明，同态
\[
\Gamma(X,\mathcal J\otimes\mathcal L^{\otimes n})
\longrightarrow
\Gamma(X,\mathcal J''\otimes\mathcal L^{\otimes n})
\]
是满射。因此，$\mathcal J''\otimes\mathcal L^{\otimes n}$ 在 $X$ 上满足
$g(x)\neq0$ 的一个截面 $g$，是某个截面
$h\in\Gamma(X,\mathcal J\otimes\mathcal L^{\otimes n})
\subset\Gamma(X,\mathcal L^{\otimes n})$ 的像
（因为由（0\textsubscript{I}，5.4.1），
$\mathcal J\otimes\mathcal L^{\otimes n}$ 是 $\mathcal L^{\otimes n}$
的一个子 $\mathcal O_X$-模）。由定义，有 $h(x)\neq0$，且对
$z\notin U$ 有 $h(z)=0$，证明完成。
\end{proof}

\begin{proposition}[2.6.2]
\label{III.2.6.2-fr}
设 $Y$ 是 Noether 预概形，$f:X\to Y$ 是有限型态射，
$g:X'\to X$ 是满的有限态射，$\mathcal L$ 是可逆 $\mathcal O_X$-模，
且 $\mathcal L'=g^*(\mathcal L)$。假设下列条件成立：存在 $X$ 的一个
在 $Y$ 上固有的子集 $Z$（II，5.4.10），使得对每个 $x\in X-Z$，
或者 $X$ 在点 $x$ 正规，或者 $(g_*(\mathcal O_{X'}))_x$ 是自由
$\mathcal O_x$-模。则 $\mathcal L$ 对 $f$ 丰沛的充要条件是
$\mathcal L'$ 对 $f\circ g$ 丰沛。
\end{proposition}

\begin{proof}
\begin{env}[2.6.2.1]
\label{III.2.6.2.1-fr}
由于 $g$ 是仿射的，这一条件是必要的（II，5.1.12）。为证明它的充分性，
可以假设 $Y$ 为仿射（II，4.6.4）。进一步证明，可以限于 $X$ 约化的情形。
事实上，设 $j:X_{\mathrm{red}}\to X$ 为典范嵌入，并置
$X_1=X_{\mathrm{red}}$、$X'_1=X'\times_X X_1$，从而有交换图
\[
\xymatrix{
X'\ar[d]_g & X'_1\ar[l]_{j'}\ar[d]^{g_1\, .}\\
X & X_1\ar[l]^{j}
}
\tag{2.6.2.2}\label{III.2.6.2.2-fr}
\]
态射 $f\circ j$ 是有限型的（I，6.3.4），$g_1$ 是有限态射
（II，6.1.5，(iii)）；若 $\mathcal L'$ 对 $f\circ g$ 丰沛，则
$j'^*(\mathcal L')$ 对 $f\circ g\circ j'$ 丰沛，因为 $j'$ 是闭浸入
（II，5.1.12 及 I，4.3.2）。置 $Z_1=j^{-1}(Z)$，则 $Z_1$ 在 $Y$
上固有（II，5.4.10）。另一方面，若 $X$ 在点 $x$ 正规，则
$X_{\mathrm{red}}$ 显然也如此；最后，若 $(g_*(\mathcal O_{X'}))_x$
是自由 $\mathcal O_x$-模，则由（II，1.5.2）立即得到
$((g_1)_*(\mathcal O_{X'_1}))_x$ 是自由 $\mathcal O_{X_1,x}$-模。
又因为 $X$ 是 Noether 的（I，6.3.7），若 $j^*(\mathcal L)$ 丰沛，
则 $\mathcal L$ 丰沛（II，4.5.14）；结合
$j'^*(\mathcal L')=g_1^*(j^*(\mathcal L))$，便完成了上述归约。
以下因而假设 $Y$ 仿射且 $X$ 约化。
\end{env}

此时（II，6.6.11）的假设成立，故存在约化 $Y$-预概形 $X_2$ 以及有限
双有理 $Y$-态射 $h:X_2\to X$，使得 $h$ 在 $h^{-1}(X-Z)$ 上的限制
是到 $X-Z$ 的同构，且 $h^*(\mathcal L)$ 丰沛。以 $X_2$ 代替 $X'$，
可见只须在进一步假设 $g$ 具有刚刚列出的 $h$ 的各项性质时证明命题。
我们仍以 $Z$ 表示 $X$ 的一个子预概形，它的底空间为 $Z$，且在 $Y$
上固有（II，5.4.10）。

\begin{env}[2.6.2.3]
\label{III.2.6.2.3-fr}
现在设 $X_1$ 是 $X$ 的闭子预概形，$j:X_1\to X$ 为典范嵌入，
$X'_1=g^{-1}(X_1)=X'\times_X X_1$ 为其逆像，$j':X'_1\to X'$
为典范嵌入，从而有交换图（\hyperref[III.2.6.2.2-fr]{2.6.2.2}）。
置 $\mathcal L_1=j^*(\mathcal L)$、
$\mathcal L'_1=j'^*(\mathcal L')=g_1^*(\mathcal L_1)$，则
$\mathcal L'_1$ 对 $f\circ g\circ j'$ 丰沛（II，5.1.12）。若置
$Z_1=j^{-1}(Z)$，则 $X_1$ 的闭子预概形 $Z_1$ 在 $Y$ 上固有
（II，5.4.2，(ii)）。换言之，（\hyperref[III.2.6.2-fr]{2.6.2}）的
假设对 $X_1$、$\mathcal L_1$、$g_1$ 和 $Z_1$ 成立。
\end{env}

\oldpage[III]{113}
这使我们能够在 $g$ 限制于 $g^{-1}(X-Z)$ 是到 $X-Z$ 的同构的情形下，
用 Noether 归纳法（0\textsubscript{I}，2.2.2）证明
（\hyperref[III.2.6.2-fr]{2.6.2}）
（如（\hyperref[III.2.6.2.1-fr]{2.6.2.1}）所见，这已足够）。
只须证明：若对 $X$ 的每个底空间 $\neq X$ 的闭子预概形 $X_1$，
（\hyperref[III.2.6.2-fr]{2.6.2}）的结论对层 $\mathcal L_1$ 成立，
则该结论对层 $\mathcal L$ 也成立。

\begin{env}[2.6.2.4]
\label{III.2.6.2.4-fr}
置 $\mathcal A=\mathcal O_X$、$\mathcal B=g_*(\mathcal O_{X'})$，
则 $\mathcal B$ 是 $\mathcal R(X)$ 的子 $\mathcal A$-代数，
且为凝聚 $\mathcal A$-模；此外，限制 $\mathcal B|(X-Z)$ 等于
$\mathcal A|(X-Z)$。设 $\mathcal K$ 为 $\mathcal B$ 在 $\mathcal A$
上的导子，即 $\mathcal A$ 的满足 $\mathcal B\mathcal K\subset\mathcal A$
的最大拟相干子 $\mathcal A$-模（也就是 $\mathcal A$-模
$\mathcal B/\mathcal A$ 的零化子（0\textsubscript{I}，5.3.7）），
从而有 $\mathcal B\mathcal K=\mathcal K$。显然，在每个具有邻域 $W_x$，
使得 $g$ 是从 $g^{-1}(W_x)$ 到 $W_x$ 的同构的点上，都有
$\mathcal K_x=\mathcal A_x$；特别地，这对 $X-Z$ 的所有点，以及
$X$ 的每个不可约分支的泛点的某个邻域内的所有点都成立。
考虑 $X$ 中由 $\mathcal K$ 定义的闭子预概形
$Z_1=\operatorname{Spec}(\mathcal A/\mathcal K)$；
它仍在 $Y$ 上固有，因为子空间 $Z_1$ 在 $Z$ 中闭（II，5.4.10）。
此外，$\mathcal K$ 的定义表明
$\mathcal B|(X-Z_1)=\mathcal A|(X-Z_1)$，故总可归约到
$Z=\operatorname{Spec}(\mathcal A/\mathcal K)$ 的情形。
又已知 $X-Z_1$ 是 $X$ 的非空开集，所以总可假设空间 $Z$ 不等于 $X$。
\end{env}

\begin{env}[2.6.2.5]
\label{III.2.6.2.5-fr}
把 $X'$ 看成 $\operatorname{Spec}(\mathcal B)$（因为 $g$ 是仿射的），
并设 $\mathcal K'=\widetilde{\mathcal K}$ 是 $\mathcal O_{X'}$ 的凝聚理想，
满足 $g_*(\mathcal K')=\mathcal K$（II，1.4.1）。
$X'$ 的闭子预概形 $Z'=g^{-1}(Z)=Z\times_X X'$ 由 $\mathcal K'$
定义，且等于 $\operatorname{Spec}(\mathcal B/\mathcal K)$（II，1.4.10）。
由于 $h:Z'\to Z$ 是有限态射（II，6.1.5，(iii)），$Z'$ 在 $Y$
上固有（II，6.1.11 及 5.4.2，(ii)）。

于是，我们须证明：对每个 $x\in X$ 及 $x$ 的每个开邻域 $U$，
都存在某个 $\mathcal L^{\otimes n}$ $(n>0)$ 在 $X$ 上的截面 $s$，
使 $x\in X_s\subset U$（II，4.5.2）。分两种情形：

$1^\circ$ 设 $x\in X-Z$。此时显然可以同时假设 $U\subset X-Z$，
所以开集 $U'=g^{-1}(U)$ 与 $Z'$ 不交。由于假设 $\mathcal L'$ 丰沛，
存在 $n>0$ 以及 $\mathcal L'^{\otimes n}$ 在 $X'$ 上的截面 $s'$，
使 $x'=g^{-1}(x)\in X'_{s'}\subset g^{-1}(U)$（II，4.5.2）。
此外，可假设 $\mathcal K'\otimes\mathcal L'^{\otimes n}$ 由其在 $X'$
上的截面生成（II，4.5.5）；因此，由 $\mathcal K'_{x'}=\mathcal O_{x'}$，
知其中有一个截面 $s''$ 满足 $s''(x')\neq0$。将它乘以 $s'$
（这相当于将 $n$ 换成 $2n$），可见还可假设
$x'\in X'_{s''}\subset g^{-1}(U)$。由（0\textsubscript{I}，5.4.10），
有典范同构
\[
\Gamma(X,\mathcal K\otimes\mathcal L^{\otimes n})
\xrightarrow{\sim}
\Gamma(X',\mathcal K'\otimes\mathcal L'^{\otimes n}).
\]
在此同构下与 $s''$ 对应的 $\mathcal K\otimes\mathcal L^{\otimes n}$
的截面 $s$ 显然具有所要求的性质。

$2^\circ$ 设 $x\in Z$。令 $\mathcal J$ 为 $\mathcal O_X$ 的凝聚理想，
它定义 $X$ 的底空间为 $X-U$ 的约化闭子预概形，并考虑
$\mathcal B$ 中的两个

\oldpage[III]{114}
凝聚理想 $\mathcal J\mathcal B$ 和
$\mathcal J\mathcal K=\mathcal J(\mathcal K\mathcal B)
=\mathcal K(\mathcal J\mathcal B)$，从而有包含关系图
\[
\xymatrix@R=1.5em{
\mathcal J\mathcal B\ar@{^{(}->}[r] & \mathcal B\\
\mathcal J\ar@{^{(}->}[r]\ar@{^{(}->}[u] &
\mathcal A\ar@{^{(}->}[u]\\
\mathcal J\mathcal K\mathcal B=\mathcal J\mathcal K
\ar@{^{(}->}[r]\ar@{^{(}->}[u] &
\mathcal K\ar@{^{(}->}[u]
}
\tag{2.6.2.6}\label{III.2.6.2.6-fr}
\]
设 $\mathcal J'$ 是 $\mathcal O_{X'}$ 的凝聚理想
$(\mathcal J\mathcal B)^\sim$，则 $\mathcal J\mathcal B=g_*(\mathcal J')$，
$\mathcal J'\mathcal K'=(\mathcal J\mathcal K\mathcal B)^\sim$，从而
\[
\mathcal J'/\mathcal J'\mathcal K'
=(\mathcal J\mathcal B/\mathcal J\mathcal K\mathcal B)^\sim
\]
（II，1.4.4）。由于对每个不与 $Z$ 相交的开集 $V$ 都有
$\mathcal J|V=\mathcal J\mathcal K|V$，可见
$\mathcal J'/\mathcal J'\mathcal K'$ 的支撑包含在 $Z'$ 中。
由于 $Z'$ 在 $Y$ 上固有，可以应用（\hyperref[III.2.2.4-fr]{2.2.4}），
得出当 $n$ 充分大时，典范映射
\[
\Gamma(X',\mathcal J'\otimes\mathcal L'^{\otimes n})
\longrightarrow
\Gamma(X',(\mathcal J'/\mathcal J'\mathcal K')
\otimes\mathcal L'^{\otimes n})
\]
是满射。由（0\textsubscript{I}，5.4.10）进而得到，典范映射
\[
\Gamma(X,\mathcal J\mathcal B\otimes\mathcal L^{\otimes n})
\longrightarrow
\Gamma(X,(\mathcal J\mathcal B/\mathcal J\mathcal K\mathcal B)
\otimes\mathcal L^{\otimes n})
\]
是满射。

现在设 $i:Z\to X$、$i':Z'\to X'$ 为典范嵌入，从而有交换图
\[
\xymatrix{
X'\ar[d]_g & Z'\ar[l]_{i'}\ar[d]^{h}\\
X & Z\ar[l]^{i}
}
\]
设 $\mathcal M=i^*(\mathcal L)$、$\mathcal M'=i'^*(\mathcal L')$。
由于 $\mathcal L'$ 丰沛，$\mathcal M'$ 也丰沛（II，5.1.12）；
另一方面，$\mathcal M'=h^*(\mathcal M)$，故由 Noether 归纳假设
（因为 $Z\neq X$），得 $\mathcal M$ 丰沛。因此，当 $n$ 充分大时，
$i^*(\mathcal J/\mathcal J\mathcal K)\otimes\mathcal M^{\otimes n}$
由其在 $Z$ 上的截面生成（II，4.5.5）。由于
$\mathcal J/\mathcal J\mathcal K=i_*(i^*(\mathcal J/\mathcal J\mathcal K))$，
再由（0\textsubscript{I}，5.4.10）可得：对于充分大的 $n$，
存在 $(\mathcal J/\mathcal J\mathcal K)\otimes
\mathcal L^{\otimes n}$ 在 $X$ 上的截面 $s$，使 $s(x)\neq0$；
这是因为，由 $\mathcal J$ 的定义有 $\mathcal J_x=\mathcal A_x$，
而由假设有 $\mathcal K_x\neq\mathcal A_x$。
图（\hyperref[III.2.6.2.6-fr]{2.6.2.6}）表明，$s$ 也是
$(\mathcal J\mathcal B/\mathcal J\mathcal K\mathcal B)
\otimes\mathcal L^{\otimes n}$ 在 $X$ 上的截面；
因此 $s$ 是 $(\mathcal J\mathcal B)\otimes\mathcal L^{\otimes n}$
在 $X$ 上的某个截面 $t$ 的典范像。而由定义，$t$ 模
$(\mathcal J\mathcal K\mathcal B)\otimes\mathcal L^{\otimes n}$
的典范像 $s$ 位于
\[
(\mathcal J\otimes\mathcal L^{\otimes n})/
(\mathcal J\mathcal K\mathcal B\otimes\mathcal L^{\otimes n})
\]
中，所以由（\hyperref[III.2.6.2.6-fr]{2.6.2.6}），$t$ 是
$\mathcal J\otimes\mathcal L^{\otimes n}$ 在 $X$ 上的截面，
\emph{更是} $\mathcal L^{\otimes n}$ 的截面。上面已证明 $t(x)\neq0$，
故 $x\in X_t$；而由 $\mathcal J$ 的定义，$t(y)=0$ 在
$X-U$ 中成立，后者是 $\mathcal O_X/\mathcal J$ 的支撑。
因此 $X_t\subset U$，证明完成。
\end{env}
\end{proof}

\oldpage[III]{115}
\begin{env}[2.6.3]
\label{III.2.6.3-fr}
\emph{注。} 当 $X$ 在 $Y$ 上固有时，可以更简单地证明
（\hyperref[III.2.6.2-fr]{2.6.2}）：仿照 Chevalley 定理（II，6.7.1）
的论证，借助（\hyperref[III.2.6.1-fr]{2.6.1}）和引理（II，6.7.1.1）即可。
\end{env}

\section{固有态射的有限性定理}
\label{section:III.3-fr}

\subsection{逐层归约（dévissage）引理}
\label{subsection:III.3.1-fr}

\begin{definition}[3.1.1]
\label{III.3.1.1-fr}
设 $K$ 是 Abel 范畴。若 $K$ 的对象集的一个子集 $K'$ 满足
$0\in K'$，且对 $K$ 中的每个正合列
\[
0\longrightarrow A'\longrightarrow A\longrightarrow A''\longrightarrow0
\]
只要对象 $A$、$A'$、$A''$ 中有两个属于 $K'$，第三个也属于 $K'$，
则称这个子集是\emph{正合的}。
\end{definition}

\begin{theorem}[3.1.2]
\label{III.3.1.2-fr}
设 $X$ 是 Noether 预概形，以 $K$ 表示凝聚 $\mathcal O_X$-模的
Abel 范畴。设 $K'$ 是 $K$ 的一个正合子集，$X'$ 是 $X$ 的底空间
的一个闭子集。假设对 $X'$ 的每个不可约闭子集 $Y$，记其泛点为 $y$，
% Canon current rebase 2026-08-31: retain the current branch's support hypothesis.
都存在支撑为 $Y$ 的 $\mathcal O_X$-模 $\mathcal G\in K'$，使 $\mathcal G_y$
是维数为 $1$ 的 $k(y)$-向量空间。则每个支撑包含于 $X'$ 的凝聚
$\mathcal O_X$-模都属于 $K'$（特别地，若 $X'=X$，则 $K'=K$）。
\end{theorem}

\begin{proof}
考虑 $X'$ 的闭子集 $Y$ 的下列性质 $P(Y)$：每个支撑包含于 $Y$
的凝聚 $\mathcal O_X$-模都属于 $K'$。由 Noether 归纳原理
（0\textsubscript{I}，2.2.2），只须证明：\emph{若 $Y$ 是 $X'$ 的闭子集，
且对 $Y$ 的每个不同于 $Y$ 的闭子集 $Y'$，性质 $P(Y')$ 都成立，
则 $P(Y)$ 成立}。

因此，设 $\mathcal F\in K$ 的支撑包含于 $Y$，并证明
$\mathcal F\in K'$。仍以 $Y$ 表示 $X$ 的底空间为 $Y$ 的约化
闭子预概形（I，5.2.1）；它由 $\mathcal O_X$ 的一个凝聚理想
$\mathcal J$ 定义。已知（I，9.3.4），存在整数 $n>0$，使
$\mathcal J^n\mathcal F=0$；故对 $1\leq k\leq n$，有凝聚
$\mathcal O_X$-模的正合列
\[
0\longrightarrow
\mathcal J^{k-1}\mathcal F/\mathcal J^k\mathcal F
\longrightarrow \mathcal F/\mathcal J^k\mathcal F
\longrightarrow \mathcal F/\mathcal J^{k-1}\mathcal F
\longrightarrow0
\]
（0\textsubscript{I}，5.3.6 及 5.3.3）。由于 $K'$ 正合，对 $k$
归纳可知，只须证明每个
$\mathcal F_k=\mathcal J^{k-1}\mathcal F/\mathcal J^k\mathcal F$
都属于 $K'$。因而只须在附加假设 $\mathcal J\mathcal F=0$ 下证明
$\mathcal F\in K'$；这一附加假设等价于
$\mathcal F=j_*(j^*(\mathcal F))$，其中 $j$ 是典范嵌入 $Y\to X$。
现在分两种情形：
\begin{enumerate}
\item[a)] $Y$ \emph{可约}。设 $Y=Y'\cup Y''$，其中 $Y'$ 和 $Y''$
是 $Y$ 的不同于 $Y$ 的闭子集。仍以 $Y'$、$Y''$ 表示 $X$ 的底空间
分别为 $Y'$、$Y''$ 的约化闭子预概形，它们分别由 $\mathcal O_X$ 的
凝聚理想 $\mathcal J'$、$\mathcal J''$ 定义。置
\[
\mathcal F'=\mathcal F\otimes_{\mathcal O_X}
(\mathcal O_X/\mathcal J'),\qquad
\mathcal F''=\mathcal F\otimes_{\mathcal O_X}
(\mathcal O_X/\mathcal J'').
\]
典范同态 $\mathcal F\to\mathcal F'$、$\mathcal F\to\mathcal F''$
定义同态 $u:\mathcal F\to\mathcal F'\oplus\mathcal F''$。
证明对每个 $z\notin Y'\cap Y''$，同态
$u_z:\mathcal F_z\to\mathcal F'_z\oplus\mathcal F''_z$ 是\emph{双射}。
事实上，有 $\mathcal J'\cap\mathcal J''=\mathcal J$，因为问题是局部的，

\oldpage[III]{116}
且上述等式由（I，5.2.1 及 1.1.5）得出。若 $z\notin Y''$，则
$\mathcal J'_z=\mathcal J_z$，从而 $\mathcal F'_z=\mathcal F_z$ 且
$\mathcal F''_z=0$，这就证明了这一情形下的断言；若 $z\notin Y'$，
可作同样的论证。因此，$u$ 的核与余核属于 $K$
（0\textsubscript{I}，5.3.4），其支撑包含于 $Y'\cap Y''$，
故由假设属于 $K'$。同理，$\mathcal F'$ 和 $\mathcal F''$ 属于 $K'$，
从而 $\mathcal F'\oplus\mathcal F''$ 也属于，因为 $K'$ 正合。
考察两个正合列
\[
0\longrightarrow\operatorname{Im}u\longrightarrow
\mathcal F'\oplus\mathcal F''\longrightarrow
\operatorname{Coker}u\longrightarrow0,
\]
\[
0\longrightarrow\operatorname{Ker}u\longrightarrow\mathcal F
\longrightarrow\operatorname{Im}u\longrightarrow0
\]
并利用 $K'$ 正合的假设，便得到结论。

\item[b)] $Y$ 不可约，所以 $X$ 的子预概形 $Y$ 是\emph{整的}。
若 $y$ 为其泛点，则 $(\mathcal O_Y)_y=k(y)$；又因 $j^*(\mathcal F)$
是凝聚 $\mathcal O_Y$-模，$\mathcal F_y=(j^*(\mathcal F))_y$
是有限维 $k(y)$-向量空间，记其维数为 $m$。
由假设，存在凝聚 $\mathcal O_X$-模 $\mathcal G$（其支撑必为 $Y$），
使 $\mathcal G_y$ 是维数为 $1$ 的 $k(y)$-向量空间。因此，存在
$k(y)$-同构 $(\mathcal G_y)^m\xrightarrow{\sim}\mathcal F_y$，
它也是 $\mathcal O_Y$-同构。由于 $\mathcal G^m$ 与 $\mathcal F$
凝聚，存在 $y$ 在 $X$ 中的开邻域 $W$ 以及同构
$\mathcal G^m|W\xrightarrow{\sim}\mathcal F|W$
（0\textsubscript{I}，5.2.7）。设 $\mathcal H$ 为此同构的图像；
它是 $(\mathcal G^m\oplus\mathcal F)|W$ 的凝聚子
$(\mathcal O_X|W)$-模，典范同构于 $\mathcal G^m|W$ 及 $\mathcal F|W$。
因此，存在 $\mathcal G^m\oplus\mathcal F$ 的凝聚子 $\mathcal O_X$-模
$\mathcal H_0$，在 $W$ 上诱导 $\mathcal H$，在 $X-Y$ 上诱导 $0$，
因为 $\mathcal G^m$ 与 $\mathcal F$ 的支撑为 $Y$（I，9.4.7）。
典范投影从 $\mathcal G^m\oplus\mathcal F$ 限制而得的
$v:\mathcal H_0\to\mathcal G^m$ 和 $w:\mathcal H_0\to\mathcal F$
于是为凝聚 $\mathcal O_X$-模的同态，且在 $W$ 及 $X-Y$ 上均为同构。
换言之，$v$ 和 $w$ 的核与余核的支撑包含于不同于 $Y$ 的闭集
$Y-(Y\cap W)$。因此它们属于 $K'$；另一方面，
$\mathcal G^m\in K'$，因为 $\mathcal G\in K'$ 且 $K'$ 正合。
再由 $K'$ 的正合性，依次得到 $\mathcal H_0\in K'$，
进而 $\mathcal F\in K'$。证毕。
\end{enumerate}
\end{proof}

\begin{corollary}[3.1.3]
\label{III.3.1.3-fr}
进一步假设 $K$ 的正合子集 $K'$ 还具有如下性质：凝聚
$\mathcal O_X$-模 $\mathcal M\in K'$ 的每个凝聚直因子仍属于 $K'$。
则在（\hyperref[III.3.1.2-fr]{3.1.2}）中，将条件
“$\mathcal G_y$ 是维数为 $1$ 的 $k(y)$-向量空间”换成
$\mathcal G_y\neq0$（这等价于 $\operatorname{Supp}(\mathcal G)=Y$），
结论仍成立。
\end{corollary}

\begin{proof}
事实上，（\hyperref[III.3.1.2-fr]{3.1.2}）的论证只须在情形 b) 中修改。
这时 $\mathcal G_y$ 是维数为 $q>0$ 的 $k(y)$-向量空间，因而有
$\mathcal O_Y$-同构
$(\mathcal G_y)^m\xrightarrow{\sim}(\mathcal F_y)^q$。
（\hyperref[III.3.1.2-fr]{3.1.2}）论证的最后部分于是表明
$\mathcal F^q\in K'$，而关于 $K'$ 的附加假设蕴含 $\mathcal F\in K'$。
\end{proof}

\subsection{有限性定理：通常概形的情形}
\label{subsection:III.3.2-fr}

\begin{theorem}[3.2.1]
\label{III.3.2.1-fr}
设 $Y$ 是局部 Noether 预概形，$f:X\to Y$ 是固有态射。
对每个凝聚 $\mathcal O_X$-模 $\mathcal F$，诸 $\mathcal O_Y$-模
$R^qf_*(\mathcal F)$ 在 $q\geq0$ 时都是凝聚的。
\end{theorem}

\begin{proof}
问题在 $Y$ 上是局部的，故可假设 $Y$ 是 Noether 的，因而 $X$
也是 Noether 的（I，6.3.7）。使定理（\hyperref[III.3.2.1-fr]{3.2.1}）
的结论成立的凝聚 $\mathcal O_X$-模 $\mathcal F$，构成凝聚
$\mathcal O_X$-模范畴 $K$ 的一个正合子集 $K'$。

\oldpage[III]{117}
事实上，设 $0\to\mathcal F'\to\mathcal F\to\mathcal F''\to0$
是凝聚 $\mathcal O_X$-模的正合列。例如，假设 $\mathcal F'$ 和
$\mathcal F''$ 属于 $K'$，则有上同调正合列
\[
R^{q-1}f_*(\mathcal F'')\xrightarrow{\partial}R^qf_*(\mathcal F')
\longrightarrow R^qf_*(\mathcal F)\longrightarrow R^qf_*(\mathcal F'')
\xrightarrow{\partial}R^{q+1}f_*(\mathcal F')
\]
其中两端的四项由假设都是凝聚的，所以由
（0\textsubscript{I}，5.3.4 及 5.3.3），中间项 $R^qf_*(\mathcal F)$
也凝聚。同样可以证明，当 $\mathcal F$ 与 $\mathcal F'$
（分别地 $\mathcal F$ 与 $\mathcal F''$）属于 $K'$ 时，
$\mathcal F''$（分别地 $\mathcal F'$）也属于。
此外，$\mathcal O_X$-模 $\mathcal F\in K'$ 的每个凝聚\emph{直因子}
$\mathcal F'$ 也属于 $K'$：事实上，$R^qf_*(\mathcal F')$ 此时为
$R^qf_*(\mathcal F)$ 的直因子（G，II，4.4.4），故为有限型；
又它拟相干（\hyperref[III.1.4.10-fr]{1.4.10}），而 $Y$ 是 Noether 的，
故它凝聚。由（\hyperref[III.3.1.3-fr]{3.1.3}），只须证明：
当 $X$ \emph{不可约}且泛点为 $x$ 时，存在属于 $K'$ 的凝聚
$\mathcal O_X$-模 $\mathcal F$，使 $\mathcal F_x\neq0$。
事实上，若已证明这一点，就可将其应用于 $X$ 的每个不可约闭子预概形
$Y$；因为若 $j:Y\to X$ 为典范嵌入，则 $f\circ j$ 固有
（II，5.4.2），且若 $\mathcal G$ 是支撑为 $Y$ 的凝聚 $\mathcal O_Y$-模，
则 $j_*(\mathcal G)$ 是凝聚 $\mathcal O_X$-模，并有
\[
R^q(f\circ j)_*(\mathcal G)=R^qf_*(j_*(\mathcal G))
\]
（G，II，4.9.1），故确实满足（\hyperref[III.3.1.3-fr]{3.1.3}）
的应用条件。

由 Chow 引理（II，5.6.2），存在不可约预概形 $X'$ 以及满的
\emph{射影}态射 $g:X'\to X$，使 $f\circ g:X'\to Y$ 是\emph{射影}的。
存在对 $g$ 丰沛的 $\mathcal O_{X'}$-模 $\mathcal L$（II，5.3.1）。
将射影态射基本定理（\hyperref[III.2.2.1-fr]{2.2.1}）应用于
$g:X'\to X$ 和 $\mathcal L$，得到一个整数 $n$，使
$\mathcal F=g_*(\mathcal O_{X'}(n))$ 是凝聚 $\mathcal O_X$-模，
且对每个 $q>0$ 有 $R^qg_*(\mathcal O_{X'}(n))=0$。
此外，由于
\[
g^*(g_*(\mathcal O_{X'}(n)))\longrightarrow\mathcal O_{X'}(n)
\]
在 $n$ 充分大时为满射（\hyperref[III.2.2.1-fr]{2.2.1}），
可假设在 $X$ 的泛点 $x$ 有 $\mathcal F_x\neq0$（II，3.4.7）。
另一方面，因为 $f\circ g$ 射影且 $Y$ 为 Noether 的，
诸 $R^p(f\circ g)_*(\mathcal O_{X'}(n))$ 是\emph{凝聚的}
（\hyperref[III.2.2.1-fr]{2.2.1}）。
又 $R^*(f\circ g)_*(\mathcal O_{X'}(n))$ 是一个 Leray 谱序列的终项，
其 $E_2$ 项由下式给出：
\[
E_2^{pq}=R^pf_*(R^qg_*(\mathcal O_{X'}(n)))~;
\]
上述结果表明，这一谱序列退化；此时已知（0，11.1.6），
$E_2^{p0}=R^pf_*(\mathcal F)$ 同构于
$R^p(f\circ g)_*(\mathcal O_{X'}(n))$，证明完成。
\end{proof}

\begin{corollary}[3.2.2]
\label{III.3.2.2-fr}
设 $Y$ 是局部 Noether 预概形。对每个固有态射 $f:X\to Y$，
每个凝聚 $\mathcal O_X$-模在 $f$ 下的直像都是凝聚 $\mathcal O_Y$-模。
\end{corollary}

\begin{corollary}[3.2.3]
\label{III.3.2.3-fr}
设 $A$ 是 Noether 环，$X$ 是 $A$ 上的固有概形。对每个凝聚
$\mathcal O_X$-模 $\mathcal F$，诸 $H^p(X,\mathcal F)$ 都是有限型
$A$-模；且存在整数 $r>0$，使得对每个凝聚 $\mathcal O_X$-模
$\mathcal F$ 及每个 $p>r$，都有 $H^p(X,\mathcal F)=0$。
\end{corollary}

\begin{proof}
第二个断言已经证明（\hyperref[III.1.4.12-fr]{1.4.12}）；
第一个断言由有限性定理（\hyperref[III.3.2.1-fr]{3.2.1}）
结合（\hyperref[III.1.4.11-fr]{1.4.11}）得出。
\end{proof}

特别地，若 $X$ 是域 $k$ 上的\emph{固有代数概形}，
则对每个凝聚 $\mathcal O_X$-模 $\mathcal F$，
诸 $H^p(X,\mathcal F)$ 都是\emph{有限维} $k$-向量空间。

\begin{corollary}[3.2.4]
\label{III.3.2.4-fr}
设 $Y$ 是局部 Noether 预概形，$f:X\to Y$ 是有限型态射。
对每个支撑在 $Y$ 上固有的凝聚 $\mathcal O_X$-模 $\mathcal F$
（II，5.4.10），诸 $\mathcal O_Y$-模 $R^qf_*(\mathcal F)$ 都凝聚。
\end{corollary}

\oldpage[III]{118}
\begin{proof}
问题在 $Y$ 上是局部的，故可假设 $Y$ 为 Noether 的，因而 $X$
也如此（I，6.3.7）。由假设，$X$ 的每个底空间为
$\operatorname{Supp}(\mathcal F)$ 的闭子预概形 $Z$ 都在 $Y$ 上固有；
换言之，若 $j:Z\to X$ 为典范嵌入，则 $f\circ j:Z\to Y$ 固有。
又可选取 $Z$ 使得 $\mathcal F=j_*(\mathcal G)$，其中
$\mathcal G=j^*(\mathcal F)$ 是凝聚 $\mathcal O_Z$-模（I，9.3.5）。
由（\hyperref[III.1.3.4-fr]{1.3.4}），有
$R^qf_*(\mathcal F)=R^q(f\circ j)_*(\mathcal G)$，
故结论立即由（\hyperref[III.3.2.1-fr]{3.2.1}）得到。
\end{proof}

\subsection{有限性定理的推广（通常概形）}
\label{subsection:III.3.3-fr}

\begin{proposition}[3.3.1]
\label{III.3.3.1-fr}
设 $Y$ 是 Noether 预概形，$\mathcal S$ 是正次数分次、拟相干且有限型的
$\mathcal O_Y$-代数，$Y'=\operatorname{Proj}(\mathcal S)$，而
$g:Y'\to Y$ 为结构态射。设 $f:X\to Y$ 是固有态射，
$\mathcal S'=f^*(\mathcal S)$，
$\mathcal M=\bigoplus_{k\in\mathbb Z}\mathcal M_k$ 是拟相干有限型分次
$\mathcal S'$-模。则
\[
R^pf_*(\mathcal M)=\bigoplus_{k\in\mathbb Z}R^pf_*(\mathcal M_k)
\]
均为有限型分次 $\mathcal S$-模（对每个 $p$）。再假设 $\mathcal S$
由 $\mathcal S_1$ 生成；则对每个 $p\in\mathbb Z$，存在整数 $k_p$，
使得对任意 $k\geq k_p$ 及任意 $r\geq0$，都有
\[
R^pf_*(\mathcal M_{k+r})=\mathcal S_rR^pf_*(\mathcal M_k).
\tag{3.3.1.1}\label{III.3.3.1.1-fr}
\]
\end{proposition}

\begin{proof}
第一个断言与（\hyperref[III.2.4.1-fr]{2.4.1}，(i)）的陈述相同，
只是将“射影态射”换为“固有态射”。而在
（\hyperref[III.2.4.1-fr]{2.4.1}，(i)）的证明中，关于 $f$ 的假设
仅用于证明（沿用该证明的记号）$R^pf'_*(\widetilde{\mathcal M})$
是凝聚 $\mathcal O_{Y'}$-模。在（\hyperref[III.3.3.1-fr]{3.3.1}）
的假设下，$f'$ 固有（II，5.4.2，(iii)），故可原封不动地沿用
（\hyperref[III.2.4.1-fr]{2.4.1}，(i)）的全部证明，只需援用有限性定理
（\hyperref[III.3.2.1-fr]{3.2.1}）。

至于第二个断言，只须注意，存在有限仿射开覆盖 $(U_i)$ 覆盖 $Y$，
使得在 $U_i$ 上，（\hyperref[III.3.3.1.1-fr]{3.3.1.1}）两端的限制
对任意 $k\geq k_{p,i}$ 都相等（II，2.1.6，(ii)）；取 $k_p$
为诸 $k_{p,i}$ 中的最大者即可。
\end{proof}

\begin{corollary}[3.3.2]
\label{III.3.3.2-fr}
设 $A$ 是 Noether 环，$\mathfrak m$ 是 $A$ 的一个理想，$X$ 是固有
$A$-概形，$\mathcal F$ 是凝聚 $\mathcal O_X$-模。则对每个
$p\geq0$，直和 $\bigoplus_{k\geq0}H^p(X,\mathfrak m^k\mathcal F)$
是环 $S=\bigoplus_{k\geq0}\mathfrak m^k$ 上的有限型模；特别地，
存在整数 $k_p\geq0$，使得对任意 $k\geq k_p$ 及任意 $r\geq0$，有
\[
H^p(X,\mathfrak m^{k+r}\mathcal F)
=\mathfrak m^rH^p(X,\mathfrak m^k\mathcal F).
\tag{3.3.2.1}\label{III.3.3.2.1-fr}
\]
\end{corollary}

\begin{proof}
只须在（\hyperref[III.3.3.1-fr]{3.3.1}）中取
$Y=\operatorname{Spec}(A)$、$\mathcal S=\widetilde S$、
$\mathcal M_k=\mathfrak m^k\mathcal F$，并计及
（\hyperref[III.1.4.11-fr]{1.4.11}）
\SourceCorrectionNote{法文印本把上式编号误作（3.3.3.1），并在证明中
循环引用推论（3.3.2）本身；Canon 稳定决定
EGA-III1-3.3.2-P463-NUMBER-AND-SELF-CITATION-001 核定编号为
（3.3.2.1），证明应用命题（3.3.1）。本译采用已核定读法，并保留可逆
来源说明。}。
\end{proof}

这里应回顾，$S$-模 $\bigoplus_{k\geq0}H^p(X,\mathfrak m^k\mathcal F)$
的结构是这样得到的：对每个 $a\in\mathfrak m^r$，考虑映射
\[
H^p(X,\mathfrak m^k\mathcal F)
\longrightarrow H^p(X,\mathfrak m^{k+r}\mathcal F)
\]
它由乘法
$\mathfrak m^k\mathcal F\to\mathfrak m^{k+r}\mathcal F$（由 $a$ 定义）
取上同调而来（\hyperref[III.2.4.1-fr]{2.4.1}）。

\oldpage[III]{119}
\subsection{有限性定理：形式概形的情形}
\label{subsection:III.3.4-fr}

本节的结果（定义（\hyperref[III.3.4.1-fr]{3.4.1}）除外）
在本章后文中不会用到。

\begin{env}[3.4.1]
\label{III.3.4.1-fr}
设 $\mathfrak X$ 与 $\mathfrak S$ 是两个局部 Noether 形式预概形
（I，10.4.2），$f:\mathfrak X\to\mathfrak S$ 是形式预概形的态射。
若 $f$ 满足下列条件，就称它为\emph{固有}态射：
\begin{enumerate}
\item[$1^\circ$] $f$ 是有限型态射（I，10.13.3）。
\item[$2^\circ$] 若 $\mathcal K$ 是 $\mathfrak S$ 的定义理想，并置
$\mathcal J=f^*(\mathcal K)\mathcal O_{\mathfrak X}$、
$X_0=(\mathfrak X,\mathcal O_{\mathfrak X}/\mathcal J)$、
$S_0=(\mathfrak S,\mathcal O_{\mathfrak S}/\mathcal K)$，则态射
$f_0:X_0\to S_0$（由 $f$ 导出，I，10.5.6）固有。
\end{enumerate}

显然，该定义与所考虑的定义理想 $\mathcal K$（属于 $\mathfrak S$）
的选择无关；事实上，若 $\mathcal K'$ 是另一个定义理想，满足
$\mathcal K'\subset\mathcal K$，并置
$\mathcal J'=f^*(\mathcal K')\mathcal O_{\mathfrak X}$、
$X'_0=(\mathfrak X,\mathcal O_{\mathfrak X}/\mathcal J')$、
$S'_0=(\mathfrak S,\mathcal O_{\mathfrak S}/\mathcal K')$，则态射
$f'_0:X'_0\to S'_0$（由 $f$ 导出）使下图交换：
\[
\xymatrix{
X_0\ar[r]^{f_0}\ar[d]_i & S_0\ar[d]^j\\
X'_0\ar[r]_{f'_0} & S'_0
}
\]
其中 $i$ 和 $j$ 均为满的浸入；由（II，5.4.5），$f_0$ 固有与
$f'_0$ 固有等价。

注意，对任意 $n\geq0$，若置
$X_n=(\mathfrak X,\mathcal O_{\mathfrak X}/\mathcal J^{n+1})$、
$S_n=(\mathfrak S,\mathcal O_{\mathfrak S}/\mathcal K^{n+1})$，则态射
$f_n:X_n\to S_n$（由 $f$ 导出，I，10.5.6）对每个 $n$ 都固有，
只要它在 $n=0$ 时固有（II，5.4.6）。
\SourceCorrectionNote{法文印本在上面两处分别写 $f_0:X_0\to Y_0$、
$f_n:X_n\to Y_n$，尽管底概形刚定义为 $S_0$、$S_n$；Canon 稳定决定
EGA-III1-3.4.1-P463-P464-TARGETS-001 核定目标分别为 $S_0$、$S_n$。
本译采用已核定读法，并保留可逆来源说明。}

若 $g:Y\to Z$ 是局部 Noether 通常预概形之间的固有态射，
$Z'$ 是 $Z$ 的闭子集，$Y'$ 是 $Y$ 的闭子集且 $g(Y')\subset Z'$，
则态射 $\widehat g:Y_{/Y'}\to Z_{/Z'}$（即 $g$ 在完备化上的延拓，
I，10.9.1）是形式预概形的固有态射；这由定义及（II，5.4.5）得出。

设 $\mathfrak X$ 与 $\mathfrak S$ 是两个局部 Noether 形式预概形，
$f:\mathfrak X\to\mathfrak S$ 是\emph{有限型}态射（I，10.13.3）。
沿用上述记号，称子集 $Z$（位于 $\mathfrak X$ 的底空间中）
\emph{固有于} $\mathfrak S$（或对 $f$ 固有），如果将它视为
$X_0$ 的子集时，$Z$ \emph{在 $S_0$ 上固有}（II，5.4.10）。
（II，5.4.10）中陈述的通常预概形固有子集的所有性质，对形式预概形的
固有子集仍然成立；这立即由定义得出。
\end{env}

\begin{theorem}[3.4.2]
\label{III.3.4.2-fr}
设 $\mathfrak X$、$\mathfrak Y$ 是两个局部 Noether 形式预概形，
$f:\mathfrak X\to\mathfrak Y$ 是固有态射。对每个凝聚
$\mathcal O_{\mathfrak X}$-模 $\mathcal F$，诸
$\mathcal O_{\mathfrak Y}$-模 $R^qf_*(\mathcal F)$ 在任意 $q\geq0$
时均凝聚。
\end{theorem}

设 $\mathcal J$ 是 $\mathfrak Y$ 的定义理想，
$\mathcal K=f^*(\mathcal J)\mathcal O_{\mathfrak X}$，并考虑
$\mathcal O_{\mathfrak X}$-模
\[
\mathcal F_k=\mathcal F\otimes_{\mathcal O_{\mathfrak Y}}
(\mathcal O_{\mathfrak Y}/\mathcal J^{k+1})
=\mathcal F/\mathcal K^{k+1}\mathcal F\qquad(k\geq0)
\tag{3.4.2.1}\label{III.3.4.2.1-fr}
\]
它们显然构成拓扑 $\mathcal O_{\mathfrak X}$-模的一个\emph{逆系统}，并满足

\oldpage[III]{120}
\[
\mathcal F=\varprojlim_k\mathcal F_k
\]
（由 I，10.11.3）。另一方面，（\hyperref[III.3.4.2-fr]{3.4.2}）将表明，
每个 $R^qf_*(\mathcal F)$ 都凝聚，因而自然带有拓扑
$\mathcal O_{\mathfrak Y}$-模结构（I，10.11.6）；诸 $R^qf_*(\mathcal F_k)$
亦然。典范同态
$\mathcal F\to\mathcal F_k=\mathcal F/\mathcal K^{k+1}\mathcal F$
典范地对应于同态
\[
R^qf_*(\mathcal F)\longrightarrow R^qf_*(\mathcal F_k)
\]
它们对于上述拓扑 $\mathcal O_{\mathfrak Y}$-模结构必定连续
（I，10.11.6），并构成逆系统；取极限便给出函子性的典范同态
\[
R^qf_*(\mathcal F)\longrightarrow\varprojlim_kR^qf_*(\mathcal F_k)
\tag{3.4.2.2}\label{III.3.4.2.2-fr}
\]
这是拓扑 $\mathcal O_{\mathfrak Y}$-模的连续同态。
我们将在证明（\hyperref[III.3.4.2-fr]{3.4.2}）的同时证明下述推论。

\begin{corollary}[3.4.3]
\label{III.3.4.3-fr}
（\hyperref[III.3.4.2.2-fr]{3.4.2.2}）中的每个同态都是拓扑同构。
此外，若 $\mathfrak Y$ 为 Noether 的，则逆系统
$(R^qf_*(\mathcal F/\mathcal K^{k+1}\mathcal F))_{k\geq0}$
满足条件 (ML)（0，13.1.1）。
\end{corollary}

我们先在以下情形建立（\hyperref[III.3.4.2-fr]{3.4.2}）及
（\hyperref[III.3.4.3-fr]{3.4.3}）：$\mathfrak Y$ 是 Noether 仿射形式概形
\SourceCorrectionNote{法文印本及外交式转录此处误作普通概形符号 $Y$；
上下文的形式概形是 $\mathfrak Y$，当前校勘本亦改作 $\mathfrak Y$。
本译采用对象一致的读法，并保留可逆来源说明；该候选已送交 Canon keeper
去重裁定。}（I，10.4.1）。

\begin{corollary}[3.4.4]
\label{III.3.4.4-fr}
在（\hyperref[III.3.4.2-fr]{3.4.2}）的假设下，再假设
$\mathfrak Y=\operatorname{Spf}(A)$，其中 $A$ 是 Noether 进制环。
设 $\mathfrak J$ 是 $A$ 的定义理想，并置
$\mathcal F_k=\mathcal F/\mathfrak J^{k+1}\mathcal F$（$k\geq0$）。
则诸 $H^n(\mathfrak X,\mathcal F)$ 是有限型 $A$-模；逆系统
$(H^n(\mathfrak X,\mathcal F_k))_{k\geq0}$ 对每个 $n$ 满足条件 (ML)；若置
\[
N_{n,k}=\operatorname{Ker}\bigl(H^n(\mathfrak X,\mathcal F)
\longrightarrow H^n(\mathfrak X,\mathcal F_k)\bigr)
\tag{3.4.4.1}\label{III.3.4.4.1-fr}
\]
（由上同调正合列，它也等于
$\operatorname{Im}(H^n(\mathfrak X,\mathfrak J^{k+1}\mathcal F)
\to H^n(\mathfrak X,\mathcal F))$），则诸 $N_{n,k}$ 在
$H^n(\mathfrak X,\mathcal F)$ 上定义一个 $\mathfrak J$-良滤过
（0，13.7.7）
\SourceCorrectionNote{法文印本先定义 $N_{n,k}$，下一句却误作 $N_{k,n}$；
Canon 稳定决定 EGA-III1-3.4.4-P464-N-INDEX-ORDER-001 核定全段统一为
$N_{n,k}$。本译采用已核定读法，并保留可逆来源说明。}；最后，典范同态
\[
H^n(\mathfrak X,\mathcal F)
\longrightarrow\varprojlim_kH^n(\mathfrak X,\mathcal F_k)
\tag{3.4.4.2}\label{III.3.4.4.2-fr}
\]
对每个 $n$ 都是拓扑同构（左端赋予 $\mathfrak J$-进制拓扑，
诸 $H^n(\mathfrak X,\mathcal F_k)$ 赋予离散拓扑）。
\end{corollary}

置
\[
S=\operatorname{gr}(A)=\bigoplus_{k\geq0}
\mathfrak J^k/\mathfrak J^{k+1},\qquad
\mathcal M=\operatorname{gr}(\mathcal F)=\bigoplus_{k\geq0}
\mathfrak J^k\mathcal F/\mathfrak J^{k+1}\mathcal F.
\tag{3.4.4.3}\label{III.3.4.4.3-fr}
\]
已知 $\mathfrak J^\Delta$ 是 $\mathfrak Y$ 的定义理想（I，10.3.1）；设
$\mathcal K=f^*(\mathfrak J^\Delta)\mathcal O_{\mathfrak X}$、
$X_0=(\mathfrak X,\mathcal O_{\mathfrak X}/\mathcal K)$、
$Y_0=(\mathfrak Y,\mathcal O_{\mathfrak Y}/\mathfrak J^\Delta)
=\operatorname{Spec}(A_0)$，其中 $A_0=A/\mathfrak J$。显然，诸
$\mathcal M_k=\mathfrak J^k\mathcal F/\mathfrak J^{k+1}\mathcal F$
是凝聚 $\mathcal O_{X_0}$-模（I，10.11.3）。另一方面，考虑拟相干分次
$\mathcal O_{X_0}$-代数
\[
\mathcal S=\mathcal O_{X_0}\otimes_{A_0}S.
\tag{3.4.4.4}\label{III.3.4.4.4-fr}
\]
乘法定义分次代数的满同态
\[
\mathcal S\longrightarrow
\operatorname{gr}(\mathcal O_{\mathfrak X})
=\bigoplus_{k\geq0}\mathcal K^k/\mathcal K^{k+1},
\]
借此 $\mathcal M$ 成为一个分次 $\mathcal S$-模。

由 $\mathcal F$ 是有限型 $\mathcal O_{\mathfrak X}$-模的假设，首先得到
$\mathcal M$ 是

\oldpage[III]{121}
\emph{有限型}分次 $\mathcal S$-模。事实上，问题在 $\mathfrak X$ 上是局部的，
故处理它时可假设 $\mathfrak X=\operatorname{Spf}(B)$，其中 $B$ 是 Noether
进制环，且 $\mathcal F=N^\Delta$，其中 $N$ 是有限型 $B$-模
（I，10.10.5）。置 $B_0=B/\mathfrak J B$、$N_0=N/\mathfrak JN$。
于是 $X_0=\operatorname{Spec}(B_0)$，而拟相干 $\mathcal O_{X_0}$-模
$\mathcal S$ 和 $\mathcal M$ 分别与下列模相伴：
\[
S'=\bigoplus_{k\geq0}
((\mathfrak J^k/\mathfrak J^{k+1})\otimes_{A_0}B_0),\qquad
M'=\bigoplus_{k\geq0}\mathfrak J^kN/\mathfrak J^{k+1}N.
\]
对每个 $k\geq0$，乘法诱导满射
\[
(\mathfrak J^k/\mathfrak J^{k+1})\otimes_{A_0}N_0
\longrightarrow\mathfrak J^kN/\mathfrak J^{k+1}N,
\qquad \overline a\otimes\overline n\longmapsto\overline{an}.
\]
这些映射定义分次 $S'$-模的满同态
$S'\otimes_{B_0}N_0\longrightarrow M'$。由于 $N_0$ 是有限型 $B_0$-模，
$M'$ 因而是有限型 $S'$-模，从而得到所述断言（I，1.3.13）。
\SourceCorrectionNote{法文印本及其外交式转录在此处把两个自然满同态误写成等式，
并以这些等式直接推出有限生成。Canon D70 的既有勘误
EGA3R2-III1-H02-BATCH-D-GRADED-QUOTIENT 在不增加假设的前提下恢复了上述分次乘法满射论证；
本译采用该核定论证，法文外交式来源保持不变。}

由假设，态射 $f_0:X_0\to Y_0$ \emph{固有}，故可将
（\hyperref[III.3.3.1-fr]{3.3.1}）应用于 $\mathcal S$、$\mathcal M$ 及态射
$f_0$；计及（\hyperref[III.1.4.11-fr]{1.4.11}），可知对\emph{任意 $n\geq0$}，
$\bigoplus_{k\geq0}H^n(X_0,\mathcal M_k)$ 是\emph{有限型}分次 $S$-模。
\SourceCorrectionNote{法文印本及其外交式转录在此处引用命题（3.3.2）；Canon D70 的既有勘误
EGA3R2-III1-H03-BATCH-E-CITATION 核定所应用的是命题（3.3.1）。本译采用核定引文，
法文外交式来源保持不变。}
这证明了（0，13.7.7）中的条件 $(\mathrm F_n)$ 对\emph{任意 $n\geq0$}
都成立，这里考虑的严格逆系统是
$(\mathcal F/\mathfrak J^k\mathcal F)_{k\geq0}$，由 $X_0$ 上的阿贝尔群层组成，
每项带有其自然的
“滤过 $A$-模”结构。因此可应用（0，13.7.7），得知：
\begin{enumerate}
\item[$1^\circ$] 逆系统 $(H^n(\mathfrak X,\mathcal F_k))_{k\geq0}$
满足条件 (ML)。
\item[$2^\circ$] 若 $H'^{,n}=\varprojlim_kH^n(\mathfrak X,\mathcal F_k)$，
则 $H'^{,n}$ 是有限型 $A$-模。
\item[$3^\circ$] 在 $H'^{,n}$ 上，由诸典范同态
$H'^{,n}\to H^n(\mathfrak X,\mathcal F_k)$ 的核所定义的滤过是
$\mathfrak J$-良滤过。
\end{enumerate}

另一方面，若置 $X_k=(\mathfrak X,\mathcal O_{\mathfrak X}/\mathcal K^{k+1})$，
则 $\mathcal F_k$ 是凝聚 $\mathcal O_{X_k}$-模（I，10.11.3）。若 $U$
是 $X_0$ 中的仿射开集，则 $U$ 在每个 $X_k$ 中也都是仿射开集
（I，5.1.9），故 $H^n(U,\mathcal F_k)=0$ 对任意 $n>0$ 及任意 $k$
成立（\hyperref[III.1.3.1-fr]{1.3.1}），且
$H^0(U,\mathcal F_k)\to H^0(U,\mathcal F_h)$ 在 $h\leq k$ 时为满射
（I，1.3.9）。因此，（0，13.3.2）的条件成立，应用（0，13.3.1）可知
$H'^{,n}$ 典范地等同于
$H^n(\mathfrak X,\varprojlim_k\mathcal F_k)=H^n(\mathfrak X,\mathcal F)$；
这就完成了（\hyperref[III.3.4.4-fr]{3.4.4}）的证明。

\begin{env}[3.4.5]
\label{III.3.4.5-fr}
现在回到（\hyperref[III.3.4.2-fr]{3.4.2}）及
（\hyperref[III.3.4.3-fr]{3.4.3}）的证明。先在
$\mathfrak Y=\operatorname{Spf}(A)$ 的情形证明这两个命题，即
（\hyperref[III.3.4.4-fr]{3.4.4}）中考虑的情形。为此，对任意 $g\in A$，
将（\hyperref[III.3.4.4-fr]{3.4.4}）应用于开集
$\mathfrak Y_g=\mathfrak D(g)$（位于 $\mathfrak Y$ 中）上的诱导 Noether
仿射形式概形，它等于 $\operatorname{Spf}(A_{\{g\}})$，以及由
$\mathfrak X$ 在 $f^{-1}(\mathfrak Y_g)$ 上诱导的形式预概形。
注意，$\mathfrak Y_g$ 在预概形
$Y_k=(\mathfrak Y,\mathcal O_{\mathfrak Y}/(\mathfrak J^\Delta)^{k+1})$
中也是仿射开集；而 $\mathcal F_k$ 是凝聚 $\mathcal O_{X_k}$-模，因此有
\[
H^n(f^{-1}(\mathfrak Y_g),\mathcal F_k)
=\Gamma(\mathfrak Y_g,R^nf_*(\mathcal F_k))
\]
对任意 $k\geq0$ 成立，这是由（\hyperref[III.1.4.11-fr]{1.4.11}）得出的。
典范同态
\[
H^n(f^{-1}(\mathfrak Y_g),\mathcal F)
\longrightarrow\varprojlim_k\Gamma(\mathfrak Y_g,R^nf_*(\mathcal F_k))
\]
是同构；但由（0\textsubscript{I}，3.2.6），有
\[
\varprojlim_k\Gamma(\mathfrak Y_g,R^nf_*(\mathcal F_k))
=\Gamma(\mathfrak Y_g,\varprojlim_kR^nf_*(\mathcal F_k))
\]

\oldpage[III]{122}
而层 $R^nf_*(\mathcal F)$ 是由预层
$\mathfrak Y_g\mapsto H^n(f^{-1}(\mathfrak Y_g),\mathcal F)$
（定义在诸 $\mathfrak Y_g$ 上）所关联的
（0\textsubscript{I}，3.2.1），故已证明同态
（\hyperref[III.3.4.2.2-fr]{3.4.2.2}）是双射。再证明
$R^nf_*(\mathcal F)$ 是凝聚 $\mathcal O_{\mathfrak Y}$-模；更精确地说，证明
\[
R^nf_*(\mathcal F)=(H^n(\mathfrak X,\mathcal F))^\Delta.
\tag{3.4.5.1}\label{III.3.4.5.1-fr}
\]
沿用上述记号，由于 $\mathcal F_k$ 是凝聚 $\mathcal O_{X_k}$-模，
有（\hyperref[III.1.4.13-fr]{1.4.13}）
\[
\Gamma(\mathfrak Y_g,R^nf_*(\mathcal F_k))
=(\Gamma(\mathfrak Y,R^nf_*(\mathcal F_k)))_g
=(H^n(\mathfrak X,\mathcal F_k))_g.
\]
诸 $H^n(\mathfrak X,\mathcal F_k)$ 构成满足 (ML) 的逆系统，且其逆极限
$H^n(\mathfrak X,\mathcal F)$ 是有限型 $A$-模。由此得到（0，13.7.8）
\[
\varprojlim_k(H^n(\mathfrak X,\mathcal F_k))_g
=H^n(\mathfrak X,\mathcal F)\otimes_AA_{\{g\}}
=\Gamma(\mathfrak Y_g,(H^n(\mathfrak X,\mathcal F))^\Delta)
\]
其中还用到了将（I，10.10.8）应用于 $A$ 及 $A_{\{g\}}$ 所得的结果。
这就证明了（\hyperref[III.3.4.5.1-fr]{3.4.5.1}），因为
\[
\Gamma(\mathfrak Y_g,R^nf_*(\mathcal F))
=\varprojlim_k\Gamma(\mathfrak Y_g,R^nf_*(\mathcal F_k)).
\]
此时，（\hyperref[III.3.4.2.2-fr]{3.4.2.2}）是凝聚
$\mathcal O_{\mathfrak Y}$-模的同构，因而必为\emph{拓扑}同构
（I，10.11.6）。最后，由关系
$R^nf_*(\mathcal F_k)=(H^n(\mathfrak X,\mathcal F_k))^\Delta$，可知逆系统
$(R^nf_*(\mathcal F_k))_{k\geq0}$ 满足 (ML)（I，10.10.2）。

在形式预概形 $\mathfrak Y$ 为 Noether 仿射的情形证明了
（\hyperref[III.3.4.2-fr]{3.4.2}）及（\hyperref[III.3.4.3-fr]{3.4.3}）后，
立即就能将（\hyperref[III.3.4.2-fr]{3.4.2}）及
（\hyperref[III.3.4.3-fr]{3.4.3}）的第一个断言推广到一般情形，因为它们
在 $\mathfrak Y$ 上都是局部的。至于（\hyperref[III.3.4.3-fr]{3.4.3}）
的第二个断言，由于 $\mathfrak Y$ 为 Noether 的，只须用有限个 Noether
仿射开集 $U_i$ 覆盖它，并注意逆系统 $(R^qf_*(\mathcal F_k))$
在每个 $U_i$ 上的限制都满足 (ML)。
\end{env}

此外，在证明过程中我们还得到了：

\begin{corollary}[3.4.6]
\label{III.3.4.6-fr}
在（\hyperref[III.3.4.4-fr]{3.4.4}）的假设下，典范同态
\[
H^q(\mathfrak X,\mathcal F)
\longrightarrow\Gamma(\mathfrak Y,R^qf_*(\mathcal F))
\tag{3.4.6.1}\label{III.3.4.6.1-fr}
\]
是双射。
\end{corollary}

\section{固有态射的基本定理及其应用}
\label{section:III.4-fr}

\subsection{基本定理}
\label{subsection:III.4.1-fr}

\begin{env}[4.1.1]
\label{III.4.1.1-fr}
设 $X$、$Y$ 是两个通常的 Noether 预概形，$f:X\to Y$ 是固有态射，
$Y'$ 是 $Y$ 的闭子集，$X'$ 是其原像 $f^{-1}(Y')$。我们用
$\widehat X$ 和 $\widehat Y$ 表示形式预概形 $X_{/X'}$ 和 $Y_{/Y'}$，
即 $X$ 和 $Y$ 分别沿 $X'$ 和 $Y'$ 的完备化（I，10.8.5），用
$\widehat f$ 表示 $f$ 到这些完备化上的延拓（I，10.9.1），它是形式
预概形的态射 $\widehat X\to\widehat Y$。对任意凝聚
$\mathcal O_X$-模 $\mathcal F$，我们

\oldpage[III]{123}

用 $\widehat{\mathcal F}$ 表示它的完备化 $\mathcal F_{/X'}$，即沿
$X'$ 的完备化（I，10.8.4），它是凝聚
$\mathcal O_{\widehat X}$-模（I，10.8.8）。
\end{env}

\begin{env}[4.1.2]
\label{III.4.1.2-fr}
设 $\mathcal J$ 是 $\mathcal O_Y$ 的凝聚理想，满足
$\operatorname{Supp}(\mathcal O_Y/\mathcal J)=Y'$（I，5.2.1）；已知
（I，4.4.5）
\[
\mathcal K=f^*(\mathcal J)\mathcal O_X
\]
是 $\mathcal O_X$ 的凝聚理想，并且
\[
\operatorname{Supp}(\mathcal O_X/\mathcal K)=X'.
\]
对每个 $k\geq0$，我们考虑凝聚 $\mathcal O_X$-模
\[
\mathcal F_k
=\mathcal F\otimes_{\mathcal O_Y}(\mathcal O_Y/\mathcal J^{k+1})
=\mathcal F/\mathcal K^{k+1}\mathcal F.
\]
$\mathcal O_Y$-模 $R^nf_*(\mathcal F)$ 和 $R^nf_*(\mathcal F_k)$
对所有 $n$ 都是凝聚的（\hyperref[III.3.2.1-fr]{3.2.1}）。对任意
$k\geq0$ 和任意 $n$，典范同态 $\mathcal F\to\mathcal F_k$
由函子性定义同态
\[
R^nf_*(\mathcal F)\longrightarrow R^nf_*(\mathcal F_k).
\tag{4.1.2.1}\label{III.4.1.2.1-fr}
\]
此外，由于 $\mathcal F_k$ 是 $\mathcal O_X/\mathcal K^{k+1}$-模，
$R^nf_*(\mathcal F_k)$ 是 $\mathcal O_Y/\mathcal J^{k+1}$-模
（0，12.2.1），所以由（\hyperref[III.4.1.2.1-fr]{4.1.2.1}）得到同态
\[
R^nf_*(\mathcal F)\otimes_{\mathcal O_Y}
(\mathcal O_Y/\mathcal J^{k+1})
\longrightarrow R^nf_*(\mathcal F_k).
\tag{4.1.2.2}\label{III.4.1.2.2-fr}
\]
（\hyperref[III.4.1.2.2-fr]{4.1.2.2}）的两端组成两个逆系统，左端的
逆极限正是完备化 $(R^nf_*(\mathcal F))_{/Y'}$，将它记为
$\widehat{R^nf_*(\mathcal F)}$。此外，同态
（\hyperref[III.4.1.2.2-fr]{4.1.2.2}）显然组成一个逆系统，因此取极限
得到典范同态
\[
\varphi_n:\widehat{R^nf_*(\mathcal F)}
\longrightarrow\varprojlim_kR^nf_*(\mathcal F_k).
\tag{4.1.2.3}\label{III.4.1.2.3-fr}
\]
又因为（\hyperref[III.4.1.2.2-fr]{4.1.2.2}）是
$(\mathcal O_Y/\mathcal J^{k+1})$-模同态，所以（I，10.8.3）可以把它
看作伪离散拓扑 $\mathcal O_{\widehat Y}$-模的连续同态
（0\textsubscript{I}，3.8.1）。因此，同态 $\varphi_n$ 是拓扑
$\mathcal O_{\widehat Y}$-模的连续同态。
\end{env}

\begin{env}[4.1.3]
\label{III.4.1.3-fr}
设 $i:\widehat X\to X$ 是（I，10.8.7）中定义的赋环空间的典范态射，
于是有交换图
\[
\xymatrix{
X_k\ar[r]^{h_k}\ar[dr]_{i_k}&\widehat X\ar[d]^i\\
&X
}
\tag{4.1.3.1}\label{III.4.1.3.1-fr}
\]
其中 $X_k$ 是 $X$ 的闭子预概形，由理想 $\mathcal K^{k+1}$ 定义；
$i_k$ 是典范嵌入，$h_k$ 是由底空间上的恒等映射与典范同态
$\mathcal O_{\widehat X}\to\mathcal O_X/\mathcal K^{k+1}$
（I，10.5.2）对应的赋环空间态射。此外，在相差一个典范同构的意义下，
有 $\widehat{\mathcal F}=i^*(\mathcal F)$（I，10.8.8）。已知，在相差
一个典范同构的意义下，有
\[
H^n(X_k,i_k^*(\mathcal F_k))=H^n(X,\mathcal F_k)
\tag{4.1.3.2}\label{III.4.1.3.2-fr}
\]

\oldpage[III]{124}

因为 $\mathcal F_k=(i_k)_*(i_k^*(\mathcal F_k))$（G，II，4.9.1）；
因此，典范同态
$H^n(\widehat X,\widehat{\mathcal F})\to
H^n(X_k,h_k^*(\widehat{\mathcal F}))$（0，12.1.3.5）也可写成
\[
H^n(\widehat X,\widehat{\mathcal F})\longrightarrow H^n(X,\mathcal F_k)
\tag{4.1.3.3}\label{III.4.1.3.3-fr}
\]
这些同态显然组成一个逆系统，从而取极限得到典范同态
\[
\psi_{n,X}:H^n(\widehat X,\widehat{\mathcal F})
\longrightarrow\varprojlim_kH^n(X,\mathcal F_k).
\tag{4.1.3.4}\label{III.4.1.3.4-fr}
\]
把 $X$ 换成形如 $f^{-1}(V)$ 的开集，其中 $V$ 是 $Y$ 的仿射开集，
考虑到（\hyperref[III.1.4.11-fr]{1.4.11}），便得到同态
\[
\psi_{n,V}:H^n(\widehat X\cap f^{-1}(V),\widehat{\mathcal F})
\longrightarrow\varprojlim_k\Gamma(V,R^nf_*(\mathcal F_k))~;
\tag{4.1.3.5}\label{III.4.1.3.5-fr}
\]
这些同态显然与把 $V$ 限制到更小的仿射开集相容，所以最终定义一个
典范层同态
\[
\psi_n:R^n\widehat f_*(\widehat{\mathcal F})
\longrightarrow\varprojlim_kR^nf_*(\mathcal F_k).
\tag{4.1.3.6}\label{III.4.1.3.6-fr}
\]
\end{env}

\begin{env}[4.1.4]
\label{III.4.1.4-fr}
最后，设 $j:\widehat Y\to Y$ 是赋环空间的典范态射（I，10.8.7）；
由于 $R^nf_*(\mathcal F)$ 是凝聚 $\mathcal O_Y$-模
（\hyperref[III.3.2.1-fr]{3.2.1}），在相差一个典范同构的意义下，有
$j^*(R^nf_*(\mathcal F))=\widehat{R^nf_*(\mathcal F)}$
（I，10.8.8），因而有典范同态
\[
\rho_n:\widehat{R^nf_*(\mathcal F)}=j^*(R^nf_*(\mathcal F))
\longrightarrow R^n\widehat f_*(i^*(\mathcal F))
=R^n\widehat f_*(\widehat{\mathcal F})
\tag{4.1.4.1}\label{III.4.1.4.1-fr}
\]
它对一般赋环空间已有定义（见（\hyperref[III.1.4.15-fr]{1.4.15}）
的证明）。我们证明图
\[
\xymatrix{
\widehat{R^nf_*(\mathcal F)}\ar[rr]^{\rho_n}\ar[dr]_{\varphi_n}
&&R^n\widehat f_*(\widehat{\mathcal F})\ar[dl]^{\psi_n}\\
&\displaystyle\varprojlim_kR^nf_*(\mathcal F_k)&
}
\tag{4.1.4.2}\label{III.4.1.4.2-fr}
\]
交换。显然，只须证明相应的预层同态图交换，因而可以限于 $Y$ 为仿射
的情形；问题归结为证明图
\[
\xymatrix{
\widehat{H^n(X,\mathcal F)}\ar[rr]^{\rho_n}\ar[dr]_{\varphi_n}
&&H^n(\widehat X,\widehat{\mathcal F})\ar[dl]^{\psi_{n,X}}\\
&\displaystyle\varprojlim_kH^n(X,\mathcal F_k)&
}
\tag{4.1.4.3}\label{III.4.1.4.3-fr}
\]
交换。而（\hyperref[III.4.1.3.1-fr]{4.1.3.1}）的交换性以及
（\hyperref[III.4.1.3-fr]{4.1.3}）中上同调群之间的关系，立即给出交换图

\oldpage[III]{125}

\[
\xymatrix{
H^n(X,\mathcal F)\ar[r]\ar[dr]&
H^n(\widehat X,\widehat{\mathcal F})
=H^n(\widehat X,i^*(\mathcal F))\ar[d]\\
&H^n(X,\mathcal F_k)=H^n(X_k,i_k^*(\mathcal F))
}
\]
由此立即得到（\hyperref[III.4.1.4.3-fr]{4.1.4.3}）的交换性。
\end{env}

\begin{theorem}[4.1.5]
\label{III.4.1.5-fr}
设 $f:X\to Y$ 是 Noether 预概形的固有态射，$Y'$ 是 $Y$ 的闭子集，
$X'=f^{-1}(Y')$。那么，对任意凝聚 $\mathcal O_X$-模 $\mathcal F$，
$R^n\widehat f_*(\widehat{\mathcal F})$ 是凝聚
$\mathcal O_{\widehat Y}$-模
\SourceCorrectionNote{法文印本及外交式转录此处作
$\mathcal O_{\widehat X}$-模；但 $R^n\widehat f_*$ 是 $\widehat Y$ 上的层，
当前校勘本亦改作 $\mathcal O_{\widehat Y}$。本译采用对象相合的读法，并
保留可逆来源说明；该候选已送交 Canon keeper 去重裁定。}，并且图
（\hyperref[III.4.1.4.2-fr]{4.1.4.2}）中的同态 $\varphi_n$、$\psi_n$
和 $\rho_n$ 都是拓扑同构。
\end{theorem}

\begin{proof}
显然，只须证明 $\varphi_n$ 和 $\psi_n$ 是同构；由于
$R^nf_*(\mathcal F)$ 是凝聚的（\hyperref[III.3.2.1-fr]{3.2.1}），这就
推出 $\widehat{R^nf_*(\mathcal F)}$ 是凝聚的（I，10.8.8），而
$\varphi_n$、$\psi_n$ 和 $\rho_n$ 的双连续性随即自动成立（I，10.11.6）。
\end{proof}

\begin{env}[4.1.6]
\label{III.4.1.6-fr}
\emph{注。}
\begin{enumerate}
\item[(i)] 若置
$\widehat{\mathcal F}_k=
\widehat{\mathcal F}/\widehat{\mathcal K}^{k+1}
\widehat{\mathcal F}$，则立即有 $\mathcal F_k=i_*(\widehat{\mathcal F}_k)$，
并且典范同态（\hyperref[III.4.1.3.6-fr]{4.1.3.6}）正是
（\hyperref[III.3.4.2.2-fr]{3.4.2.2}）中已经定义的同态
\[
R^n\widehat f_*(\widehat{\mathcal F})
\longrightarrow\varprojlim_kR^n\widehat f_*(\widehat{\mathcal F}_k)~;
\tag{4.1.6.1}\label{III.4.1.6.1-fr}
\]
因此，$\psi_n$ 为同构这一事实是（\hyperref[III.3.4.3-fr]{3.4.3}）的
特例。不过，下文将给出直接证明，以避免一般定理
（\hyperref[III.3.4.3-fr]{3.4.3}）所依赖的关于谱序列逆极限的精细论证
（0，13.7）。

\item[(ii)] 由于 $\psi_n$ 是同构，断言 $\varphi_n$ 是同构，等价于
断言 $\rho_n=\psi_n^{-1}\circ\varphi_n$ 是同构。定理
（\hyperref[III.4.1.5-fr]{4.1.5}）特别表明，$R^nf_*$ 的形成与完备化
相交换；因此可以称之为从“代数”理论到“形式”理论的第一比较定理。
\end{enumerate}
我们先建立（\hyperref[III.4.1.5-fr]{4.1.5}）的仿射形式：
\end{env}

\begin{corollary}[4.1.7]
\label{III.4.1.7-fr}
在（\hyperref[III.4.1.5-fr]{4.1.5}）的假设下，再设
$Y=\operatorname{Spec}(A)$，其中 $A$ 是 Noether 环，并且
$\mathcal J=\widetilde{\mathfrak J}$，其中 $\mathfrak J$ 是 $A$ 的理想，
于是 $\mathcal F_k=\mathcal F/\mathfrak J^{k+1}\mathcal F$。典范同态
\[
\varphi_n:\widehat{H^n(X,\mathcal F)}
\longrightarrow\varprojlim_kH^n(X,\mathcal F_k)
\tag{4.1.7.1}\label{III.4.1.7.1-fr}
\]
是同构（其中左端是 $H^n(X,\mathcal F)$ 关于 $\mathfrak J$-预进制拓扑
的分离完备化）。逆系统 $(H^n(X,\mathcal F_k))_{k\geq0}$ 对所有 $n$
都满足条件 (ML)，并且典范同态
\[
\psi_n:H^n(\widehat X,\widehat{\mathcal F})
\longrightarrow\varprojlim_kH^n(X,\mathcal F_k)
\tag{4.1.7.2}\label{III.4.1.7.2-fr}
\]

\oldpage[III]{126}

是同构。最后，在 $H^n(X,\mathcal F)$ 上由典范同态
$H^n(X,\mathcal F)\to H^n(X,\mathcal F_k)$ 的核定义的滤过是
$\mathfrak J$-良滤过（0，13.7.7），并且 $\varphi_n$ 是拓扑同构。
\footnote{下面的证明比原来的证明更简单；此证明以及关于
$H^n(X,\mathcal F)$ 的滤过的补充，均由 J.-P. Serre 告知我们。}
\end{corollary}

\begin{proof}
在本证明中固定整数 $n\geq0$，为简便起见，置
\[
\mathrm H=H^n(X,\mathcal F),\qquad
\mathrm H_k=H^n(X,\mathcal F_k)
\tag{4.1.7.3}\label{III.4.1.7.3-fr}
\]
以及
\[
\mathrm R_k=\operatorname{Ker}(\mathrm H\longrightarrow\mathrm H_k),
\qquad\text{作为子 $A$-模包含于 $\mathrm H$。}
\tag{4.1.7.4}\label{III.4.1.7.4-fr}
\]
上同调正合列
\[
H^n(X,\mathfrak J^{k+1}\mathcal F)\longrightarrow H^n(X,\mathcal F)
\longrightarrow H^n(X,\mathcal F_k)\xrightarrow{\partial}
H^{n+1}(X,\mathfrak J^{k+1}\mathcal F)
\longrightarrow H^{n+1}(X,\mathcal F)
\]
表明还存在关系
\[
\mathrm R_k=\operatorname{Im}\bigl(
H^n(X,\mathfrak J^{k+1}\mathcal F)\longrightarrow H^n(X,\mathcal F)
\bigr)~;
\]
我们置
\[
\begin{aligned}
\mathrm Q_k
&=\operatorname{Ker}\bigl(H^{n+1}(X,\mathfrak J^{k+1}\mathcal F)
\longrightarrow H^{n+1}(X,\mathcal F)\bigr)\\
&=\operatorname{Im}\bigl(H^n(X,\mathcal F_k)
\longrightarrow H^{n+1}(X,\mathfrak J^{k+1}\mathcal F)\bigr).
\end{aligned}
\tag{4.1.7.5}\label{III.4.1.7.5-fr}
\]
于是有正合列
\[
0\longrightarrow\mathrm R_k\longrightarrow\mathrm H
\longrightarrow\mathrm H_k\longrightarrow\mathrm Q_k\longrightarrow0.
\tag{4.1.7.6}\label{III.4.1.7.6-fr}
\]

\begin{env}[4.1.7.7]
\label{III.4.1.7.7-fr}
设 $x$ 是 $\mathfrak J^m$ 的元素（$m\geq0$）；乘以 $x$ 的运算在
$\mathfrak J^k\mathcal F$ 上给出同态
$\mathfrak J^k\mathcal F\to\mathfrak J^{k+m}\mathcal F$，从而诱导同态
\[
\mu_{x,m}:H^n(X,\mathfrak J^k\mathcal F)
\longrightarrow H^n(X,\mathfrak J^{k+m}\mathcal F).
\tag{4.1.7.8}\label{III.4.1.7.8-fr}
\]
若用 $S$ 表示分次 $A$-代数 $\bigoplus_{k\geq0}\mathfrak J^k$，则已知
这些乘法 $\mu_{x,m}$ 在
$\mathrm E=\bigoplus_{k\geq0}H^n(X,\mathfrak J^k\mathcal F)$ 上定义了
分次环 $S$ 上的有限型分次模结构（\hyperref[III.3.3.2-fr]{3.3.2}），
而该分次环是 Noether 的（II，2.1.5）。
\end{env}

\begin{lemma}[4.1.7.9]
\label{III.4.1.7.9-fr}
子模 $\mathrm R_k$ 包含于 $\mathrm H$，并在 $\mathrm H$ 上定义一个
$\mathfrak J$-良滤过。
\end{lemma}

\begin{proof}
首先证明
\[
\mathfrak J^m\mathrm R_k\subset\mathrm R_{k+m},
\tag{4.1.7.10}\label{III.4.1.7.10-fr}
\]
其中在 $\mathrm H=H^n(X,\mathcal F)$ 中乘以元素 $x\in\mathfrak J^m$
的映射就是 $\mu_{x,0}$。对任意 $x\in\mathfrak J^m$，图
\[
\xymatrix{
\mathfrak J^{k+1}\mathcal F\ar[r]^x\ar[d]&
\mathfrak J^{k+m+1}\mathcal F\ar[d]\\
\mathcal F\ar[r]_x&\mathcal F
}
\tag{4.1.7.11}\label{III.4.1.7.11-fr}
\]

\oldpage[III]{127}

是交换的（其中水平箭头是乘以 $x$，竖直箭头是典范嵌入）；所以相应的图
\[
\xymatrix{
H^n(X,\mathfrak J^{k+1}\mathcal F)\ar[r]^{\mu_{x,m}}\ar[d]&
H^n(X,\mathfrak J^{k+m+1}\mathcal F)\ar[d]\\
H^n(X,\mathcal F)\ar[r]_{\mu_{x,0}}&H^n(X,\mathcal F)
}
\]
也是交换的。结合 $\mathrm R_k$ 是
$H^n(X,\mathfrak J^{k+1}\mathcal F)\to H^n(X,\mathcal F)$ 的像这一解释，
这就证明了（\hyperref[III.4.1.7.10-fr]{4.1.7.10}），并且还表明分次
$S$-模 $\mathrm R=\bigoplus_{k\geq0}\mathrm R_k$ 是子 $S$-模
$\mathrm M=\bigoplus_{k\geq0}H^n(X,\mathfrak J^{k+1}\mathcal F)$ 的商，
其中该子模包含于 $\mathrm E$；
由上面的注可知，$\mathrm R$ 是有限型 $S$-模，这等价于
（\hyperref[III.4.1.7.9-fr]{4.1.7.9}）的断言（Bourbaki，
\emph{交换代数}，第 III 章，§~3，第 1 节，定理 1）。
\end{proof}

\begin{env}[4.1.7.12]
\label{III.4.1.7.12-fr}
现在考虑分次 $S$-模
\[
\mathrm N=\bigoplus_{k\geq0}H^{n+1}(X,\mathfrak J^{k+1}\mathcal F)
\]
它的定义与 (4.7.1.8) 中相同；根据（\hyperref[III.3.3.2-fr]{3.3.2}），
它仍然是有限型 $S$-模。关系 $\mathrm Q_k\subset\mathrm N_k$ 对每个
$k$ 都成立，这是由（\hyperref[III.4.1.7.5-fr]{4.1.7.5}）得出的；而将图
（\hyperref[III.4.1.7.11-fr]{4.1.7.11}）中的 $n$ 换成 $n+1$，便可看出
$S_m\mathrm Q_k=\mathfrak J^m\mathrm Q_k\subset\mathrm Q_{k+m}$。
换言之，$\mathrm Q$ 作为分次子 $S$-模包含于 $\mathrm N$，因而是有限型的。
\end{env}

\begin{env}[4.1.7.13]
\label{III.4.1.7.13-fr}
用 $\alpha_m$ 表示典范嵌入 $\mathfrak J^m\to A$，它也可写作
$S_m\to S_0$。因为 $\mathfrak J^{k+1}\mathcal F_k=0$，$A$-模
$H^n(X,\mathcal F_k)$ 被 $\mathfrak J^{k+1}$ 零化；由于 $\mathrm Q_k$
是 $A$-同态
$H^n(X,\mathcal F_k)\to H^{n+1}(X,\mathfrak J^{k+1}\mathcal F)$ 的像，
$\mathrm Q_k$ 作为 $A$-模也被 $\mathfrak J^{k+1}$ 零化；这还意味着，
在 $S$-模 $\mathrm Q$ 中有
\[
\alpha_{k+1}(S_{k+1})\mathrm Q_k=0.
\tag{4.1.7.14}\label{III.4.1.7.14-fr}
\]
由于 $\mathrm Q$ 是有限型 $S$-模，存在整数 $k_0$ 和整数 $h$，使得
$\mathrm Q_{k+h}=S_h\mathrm Q_k$ 对 $k\geq k_0$ 成立（II，2.1.6，(ii)）；
由此关系和（\hyperref[III.4.1.7.14-fr]{4.1.7.14}），可知存在整数
$r>0$，使得
\[
\alpha_r(S_r)\mathrm Q=0.
\tag{4.1.7.15}\label{III.4.1.7.15-fr}
\]
\end{env}

\begin{env}[4.1.7.16]
\label{III.4.1.7.16-fr}
现在注意，典范嵌入
$\mathfrak J^{k+m}\mathcal F\to\mathfrak J^k\mathcal F$ 取上同调后
给出 $A$-同态
\[
\nu_m:H^{n+1}(X,\mathfrak J^{k+m}\mathcal F)
\longrightarrow H^{n+1}(X,\mathfrak J^k\mathcal F)
\tag{4.1.7.17}\label{III.4.1.7.17-fr}
\]
而且，对任意 $x\in\mathfrak J^m$，显然有分解
\[
\xymatrix@C=4.5em{
H^{n+1}(X,\mathfrak J^k\mathcal F)\ar[r]^{\mu_{x,m}}&
H^{n+1}(X,\mathfrak J^{k+m}\mathcal F)\ar[r]^{\nu_m}&
H^{n+1}(X,\mathfrak J^k\mathcal F)
}
\tag{4.1.7.18}\label{III.4.1.7.18-fr}
\]
其复合为 $\mu_{x,0}$。
\end{env}

\oldpage[III]{128}

由此可知，对任意子 $A$-模 $P$（包含于
$H^{n+1}(X,\mathfrak J^k\mathcal F)$），在 $S$-模 $\mathrm N$ 中有
\[
\nu_m(S_mP)=\alpha_m(S_m)P.
\tag{4.1.7.19}\label{III.4.1.7.19-fr}
\]

\begin{lemma}[4.1.7.20]
\label{III.4.1.7.20-fr}
存在整数 $m>0$，使得 $\nu_m(\mathrm Q_{k+m})=0$ 对每个 $k\geq k_0$
都成立。
\end{lemma}

\begin{proof}
事实上，取 $m$ 为 $h$ 的一个倍数，并满足 $\geq r$；由于
$\mathrm Q_{k+m}=S_m\mathrm Q_k$ 对 $k\geq k_0$ 成立，由
（\hyperref[III.4.1.7.19-fr]{4.1.7.19}）和
（\hyperref[III.4.1.7.15-fr]{4.1.7.15}）得到
\[
\nu_m(\mathrm Q_{k+m})=\alpha_m(S_m)\mathrm Q_k
\subset\alpha_r(S_r)\mathrm Q_k=0.
\]
\end{proof}

\begin{env}[4.1.7.21]
\label{III.4.1.7.21-fr}
注意交换图
\[
\xymatrix@C=2.5em{
H^n(X,\mathcal F)\ar[r]&H^n(X,\mathcal F_k)\ar[r]^-\partial&
H^{n+1}(X,\mathfrak J^{k+1}\mathcal F)\ar[r]&H^{n+1}(X,\mathcal F)\\
H^n(X,\mathcal F)\ar[r]\ar[u]&H^n(X,\mathcal F_{k+m})\ar[r]^-\partial\ar[u]&
H^{n+1}(X,\mathfrak J^{k+m+1}\mathcal F)\ar[r]\ar[u]&
H^{n+1}(X,\mathcal F)\ar[u]
}
\]
本身来自交换图
\[
\xymatrix{
0\ar[r]&\mathfrak J^{k+1}\mathcal F\ar[r]&\mathcal F\ar[r]&
\mathcal F_k\ar[r]&0\\
0\ar[r]&\mathfrak J^{k+m+1}\mathcal F\ar[r]\ar[u]&
\mathcal F\ar[r]\ar[u]&\mathcal F_{k+m}\ar[r]\ar[u]&0
}
\]
其中竖直箭头为典范映射。由此得到交换图
\[
\xymatrix@C=2.2em{
0\ar[r]&\mathrm R_k\ar[r]&\mathrm H\ar[r]&\mathrm H_k\ar[r]&
\mathrm Q_k\ar[r]&0\\
0\ar[r]&\mathrm R_{k+m}\ar[r]\ar[u]&\mathrm H\ar[r]\ar@{=}[u]&
\mathrm H_{k+m}\ar[r]\ar[u]&\mathrm Q_{k+m}\ar[r]\ar[u]^{\nu_m}&0
}
\]
其中各行正合。由于当 $k\geq k_0$ 时，最后一个竖直箭头为零
（\hyperref[III.4.1.7.20-fr]{4.1.7.20}），$\mathrm H_{k+m}$ 在
$\mathrm H_k$ 中的像包含于
$\operatorname{Ker}(\mathrm H_k\to\mathrm Q_k)
=\operatorname{Im}(\mathrm H\to\mathrm H_k)$；另一方面，由图的交换性，
该像包含 $\operatorname{Im}(\mathrm H\to\mathrm H_k)$，所以二者相等。
因而，在 $\mathrm H_k$ 中，$\mathrm H_{k'}$ 的像对所有 $k'\geq k+m$
也都相同，这就证明逆系统 $(\mathrm H_k)_{k\geq0}$ 满足条件 (ML)。
此外，任取仿射开集 $U$ 包含于 $X$，有 $H^i(U,\mathcal F_k)=0$
（$i>0$）（\hyperref[III.1.3.1-fr]{1.3.1}），而当 $m>0$
时，映射 $H^0(U,\mathcal F_{k+m})\to H^0(U,\mathcal F_k)$ 是满射
（I，1.3.9）。因此可以

\oldpage[III]{129}

应用（0，13.3.1），于是典范同态
$H^n(\widehat X,\widehat{\mathcal F})\to
\varprojlim_kH^n(X,\mathcal F_k)$ 对每个 $n\geq0$ 都是双射。

因为逆系统 $(\mathrm H/\mathrm R_k)_{k\geq0}$ 是严格的，可以对正合列
\[
0\longrightarrow\mathrm H/\mathrm R_k\longrightarrow\mathrm H_k
\longrightarrow\mathrm Q_k\longrightarrow0
\tag{4.1.7.22}\label{III.4.1.7.22-fr}
\]
取逆极限（0，13.2.2）；由于 $\nu_m(\mathrm Q_{k+m})=0$，有
$\varprojlim_k\mathrm Q_k=0$，从而得到拓扑同构
\[
\varprojlim_k(\mathrm H/\mathrm R_k)\xrightarrow{\sim}
\varprojlim_k\mathrm H_k.
\]
又因为滤过 $(\mathrm R_k)$ 在 $\mathrm H$ 上是 $\mathfrak J$-良滤过，
它在 $\mathrm H$ 上定义了 $\mathfrak J$-预进制拓扑；所以
$\varprojlim_k(\mathrm H/\mathrm R_k)$ 正是 $\mathrm H$ 关于
$\mathfrak J$-预进制拓扑的分离完备化。这就完成了
（\hyperref[III.4.1.7-fr]{4.1.7}）的证明。
\end{env}
\end{proof}

\begin{env}[4.1.8]
\label{III.4.1.8-fr}
最后证明（\hyperref[III.4.1.5-fr]{4.1.5}）：任取仿射开集 $V$ 包含于
$Y$，$\Gamma(V,\widehat{R^nf_*(\mathcal F)})$ 是
$\Gamma(V,R^nf_*(\mathcal F))$ 关于 $\mathfrak J$-预进制拓扑的分离
完备化（若 $\mathcal J|V=\widetilde{\mathfrak J}$），因为
$R^nf_*(\mathcal F)$ 是凝聚 $\mathcal O_Y$-模（I，10.8.4），并且
\[
\Gamma\bigl(V,\varprojlim_kR^nf_*(\mathcal F_k)\bigr)
=\varprojlim_k\Gamma(V,R^nf_*(\mathcal F_k))
\]
（0\textsubscript{I}，3.2.6）；于是 $\varphi_n$ 是拓扑同构这一事实，
由（\hyperref[III.4.1.7-fr]{4.1.7}）和
（\hyperref[III.1.4.11-fr]{1.4.11}）得出。另一方面，仍由
（\hyperref[III.1.4.11-fr]{1.4.11}），结合
（\hyperref[III.4.1.7-fr]{4.1.7}），可知
（\hyperref[III.4.1.3.3-fr]{4.1.3.3}）中的同态 $\psi_{n,V}$ 是同构；
因此，$\psi_n$ 也是同构，这是由
$R^n\widehat f_*(\widehat{\mathcal F})$ 的定义得出的。
\end{env}

\begin{corollary}[4.1.9]
\label{III.4.1.9-fr}
在（\hyperref[III.4.1.4-fr]{4.1.4}）的假设下，任取仿射开集 $V$ 包含于
$Y$，典范同态
\[
H^n(\widehat X\cap f^{-1}(V),\widehat{\mathcal F})
\longrightarrow\Gamma(\widehat Y\cap V,
R^n\widehat f_*(\widehat{\mathcal F}))
\]
是双射。
\end{corollary}

\begin{env}[4.1.10]
\label{III.4.1.10-fr}
\emph{注。}
设 $f:X\to Y$ 是（通常的）Noether 预概形的有限型态射，
$\mathcal F$ 是凝聚 $\mathcal O_X$-模，其支撑在 $Y$ 上是固有的
（II，5.4.10）。这时已知（\hyperref[III.3.2.4-fr]{3.2.4}），
$R^nf_*(\mathcal F)$ 是凝聚 $\mathcal O_Y$-模，这对每个 $n\geq0$ 都成立。
此外，总可以设 $\mathcal F=u_*(\mathcal G)$，其中
$\mathcal G=u^*(\mathcal F)$ 是凝聚 $\mathcal O_Z$-模，$Z$ 表示
$X$ 的一个适当闭子预概形，其底空间为 $\operatorname{Supp}(\mathcal F)$，
而 $u:Z\to X$ 为典范嵌入（I，9.3.5）。若置
\[
\mathcal G_k=\mathcal G\otimes_{\mathcal O_Y}
(\mathcal O_Y/\mathcal J^{k+1}),
\]
则有 $\mathcal G_k=u^*(\mathcal F_k)$、
$R^nf_*(\mathcal F_k)=R^n(f\circ u)_*(\mathcal G_k)$ 和
$R^nf_*(\mathcal F)=R^n(f\circ u)_*(\mathcal G)$
（\hyperref[III.1.3.4-fr]{1.3.4}），最后，由（I，10.9.5）有
\[
R^n\widehat f_*(\widehat{\mathcal F})
=R^n\widehat{(f\circ u)}_*(\widehat{\mathcal G}).
\]
于是可将（\hyperref[III.4.1.5-fr]{4.1.5}）应用于 $\mathcal G$ 和固有
态射 $f\circ u$，从而得出：在这些假设下，
（\hyperref[III.4.1.5-fr]{4.1.5}）的结论对 $\mathcal F$ 和 $f$ 仍然成立。
\end{env}

\subsection{特殊情形与变式}
\label{subsection:III.4.2-fr}

比较定理（\hyperref[III.4.1.5-fr]{4.1.5}）最常用的形式如下：

\begin{proposition}[4.2.1]
\label{III.4.2.1-fr}
设 $Y$ 是局部 Noether 预概形，$f:X\to Y$ 是固有态射，
$\mathcal F$ 是凝聚 $\mathcal O_X$-模。则对每个 $y\in Y$ 和每个 $p$，
$(R^pf_*(\mathcal F))_y$ 是

\oldpage[III]{130}

有限型 $\mathcal O_y$-模，因而对于 $\mathfrak m_y$-预进制拓扑是分离的；
并且有典范拓扑同构
\[
\widehat{(R^pf_*(\mathcal F))_y}
\xrightarrow{\sim}
\varprojlim_kH^p\bigl(f^{-1}(y),
\mathcal F\otimes_{\mathcal O_Y}(\mathcal O_y/\mathfrak m_y^k)\bigr)
\tag{4.2.1.1}\label{III.4.2.1.1-fr}
\]
其中左端是 $(R^pf_*(\mathcal F))_y$ 关于
$\mathfrak m_y$-预进制拓扑的完备化；在右端中，将 $f^{-1}(y)$
对每个 $k\geq0$ 看作预概形
$X\times_Y\operatorname{Spec}(\mathcal O_y/\mathfrak m_y^k)$ 的底空间
（I，3.6.1）。
\end{proposition}

\begin{proof}
由于 $\mathcal O_y$ 是 Noether 局部环，而
$(R^pf_*(\mathcal F))_y$ 是有限型 $\mathcal O_y$-模
（\hyperref[III.3.2.1-fr]{3.2.1}），故关于 $\mathfrak m_y$-预进制拓扑，
$(R^pf_*(\mathcal F))_y$ 是分离的
（0\textsubscript{I}，7.3.5）。当 $Y$ 为 Noether 且点 $y$ 为闭点时，
其余断言由（\hyperref[III.4.1.7-fr]{4.1.7}）得出：把 $Y$ 换成
$y$ 的一个仿射邻域，并取 $Y'=\{y\}$，同时应用（G, II, 4.9.1）。
在一般情形下，置
\[
Y_1=\operatorname{Spec}(\mathcal O_y),\qquad
X_1=X\times_YY_1,\qquad
\mathcal F_1=\mathcal F\otimes_{\mathcal O_Y}\mathcal O_{Y_1}
\]
并令 $f_1=f\times1_{Y_1}:X_1\to Y_1$。于是 $Y_1$ 是 Noether 的，
$f_1$ 是固有的（II，5.4.2，(iii)），而 $\mathcal F_1$ 是凝聚的
（0\textsubscript{I}，5.3.11）。设 $y_1$ 是 $Y_1$ 的唯一闭点；
命题对 $f_1$、$\mathcal F_1$ 和 $y_1$ 成立。又有
$\mathcal O_{y_1}=\mathcal O_y$、$f_1^{-1}(y_1)=f^{-1}(y)$
（I，3.6.5），并且预概形
$X\times_Y\operatorname{Spec}(\mathcal O_y/\mathfrak m_y^k)$ 与
$X_1\times_{Y_1}\operatorname{Spec}(\mathcal O_y/\mathfrak m_y^k)$
典范同一（I，3.3.9）。此外，
$\mathcal F_1\otimes_{\mathcal O_{Y_1}}
(\mathcal O_y/\mathfrak m_y^k)$ 与
$\mathcal F\otimes_{\mathcal O_Y}(\mathcal O_y/\mathfrak m_y^k)$
同一（I，9.1.6）。只需再证明 $R^pf_{1*}(\mathcal F_1)$ 典范同构于
$R^pf_*(\mathcal F)\otimes_{\mathcal O_Y}\mathcal O_{Y_1}$；
这由（\hyperref[III.1.4.15-fr]{1.4.15}）得出，因为局部态射
$\operatorname{Spec}(\mathcal O_y)\to Y$ 是平坦的
（0\textsubscript{I}，6.7.1 及 I，2.4.2）。
\end{proof}

下一推论采用维数理论（第四章）的术语，并且在第四章之前不会使用。

\begin{corollary}[4.2.2]
\label{III.4.2.2-fr}
设 $Y$ 是局部 Noether 预概形，$f:X\to Y$ 是固有态射，$y$ 是
$Y$ 的一个点，$r$ 是 $f^{-1}(y)$ 的维数。则对每个凝聚
$\mathcal O_X$-模 $\mathcal F$，层 $R^pf_*(\mathcal F)$ 在 $y$ 的
某个邻域上当 $p>r$ 时为零。
\end{corollary}

\begin{proof}
事实上，此时
\[
H^p\bigl(f^{-1}(y),
\mathcal F\otimes_{\mathcal O_Y}(\mathcal O_y/\mathfrak m_y^k)\bigr)=0
\]
对每个 $k$ 都成立（G, II, 4.15.2）；所以由
（\hyperref[III.4.2.1-fr]{4.2.1}），
$(R^pf_*(\mathcal F))_y$ 关于 $\mathfrak m_y$-预进制拓扑的
分离完备化为零。由于该拓扑是分离的，亦有
$(R^pf_*(\mathcal F))_y=0$。又因 $R^pf_*(\mathcal F)$ 是凝聚的
（0\textsubscript{I}，5.2.2），故得结论。
\end{proof}

\begin{env}[4.2.3]
\label{III.4.2.3-fr}
结果（\hyperref[III.4.2.1-fr]{4.2.1}）主要用于 $p=0$ 的情形；
于是得到下列推论。
\end{env}

\begin{corollary}[4.2.4]
\label{III.4.2.4-fr}
在（\hyperref[III.4.2.1-fr]{4.2.1}）的假设下，有典范拓扑同构
\[
\widehat{(f_*(\mathcal F))_y}
\xrightarrow{\sim}
\varprojlim_k\Gamma\bigl(f^{-1}(y),
\mathcal F_y/\mathfrak m_y^k\mathcal F_y\bigr).
\]
\end{corollary}

\subsection{Zariski 连通性定理}
\label{subsection:III.4.3-fr}
\label{III.4.3-fr}

本节及下一节的结果推广了 Zariski 的若干著名定理，并且都可由
（\hyperref[III.4.2.4-fr]{4.2.4}）推出。它们是下列定理的推论：

\begin{theorem}[4.3.1]
\label{III.4.3.1-fr}
\textnormal{（连通性定理）。}
设 $Y$ 是局部 Noether 预概形，$f:X\to Y$ 是固有态射。则
$\mathcal A(X)=f_*(\mathcal O_X)$ 是凝聚 $\mathcal O_Y$-代数。

\oldpage[III]{131}

令 $Y'$ 为一个 $Y$-概形，它在 $Y$ 上有限且满足
$\mathcal A(Y')=\mathcal A(X)$；它在相差一个 $Y$-同构的意义下唯一确定
（II，1.3.1 和 6.1.3）。若 $f'=\mathcal A(e)$ 是 $Y$-态射 $X\to Y'$，
它由恒等同构 $e:\mathcal A(Y')\to\mathcal A(X)$ 导出
（II，1.2.7），则 $f'$ 是固有的，$f'_*(\mathcal O_X)$ 同构于
$\mathcal O_{Y'}$，并且纤维 $f'^{-1}(y')$（态射 $f'$ 的纤维）对每个
$y'\in Y'$ 都连通且非空。
\end{theorem}

\begin{proof}
令 $g:Y'\to Y$ 为结构态射。为证明同态
$\theta:\mathcal O_{Y'}\to f'_*(\mathcal O_X)$（它出现在态射 $f'$
的定义中）是双射，只须证明
\[
g_*(\theta):g_*(\mathcal O_{Y'})\longrightarrow
g_*(f'_*(\mathcal O_X))=f_*(\mathcal O_X)
\]
是恒等映射（II，1.4.2），因为 $Y'$ 在 $Y$ 上是仿射的；但这由定义
立即得出，因为 $g_*(\mathcal O_{Y'})=\mathcal A(Y')$ 且
$f_*(\mathcal O_X)=\mathcal A(X)$。$\mathcal A(X)$ 的凝聚性是有限性定理
（\hyperref[III.3.2.1-fr]{3.2.1}）的一个特殊情形。由于 $f$ 是固有的
而 $g$ 是分离的，$f'$ 是固有的（II，5.4.3，(i)）；因此，为完成
（\hyperref[III.4.3.1-fr]{4.3.1}）的证明，只须证明下面的推论
\SourceCorrectionNote{法文印本在这里止于未完短语“只须证明其……”，
Canon 稳定决定 EGA-III1-4.3.1-P475-TRUNCATED-PROOF-001 确认这是印本
遗漏而非转录截断。现依上下文补明“下面的推论”，并保留可逆来源说明。}。
\end{proof}

\begin{corollary}[4.3.2]
\label{III.4.3.2-fr}
在（\hyperref[III.4.3.1-fr]{4.3.1}）的假设下，再假设
$f_*(\mathcal O_X)$ 同构于 $\mathcal O_Y$。则 $f^{-1}(y)$ 是 $f$ 的
连通非空纤维，对每个 $y\in Y$ 都如此。
\end{corollary}

\begin{proof}
$f_*(\mathcal O_X)$ 同构于 $\mathcal O_Y$ 这一假设已经蕴含 $f$ 是
支配的；又因 $f$ 是闭映射，故该态射是满射。与
（\hyperref[III.4.2.1-fr]{4.2.1}）中一样，可归结到 $y$ 在 $Y$ 中
为闭点的情形。$f^{-1}(y)$ 是 Noether 空间，因而只有有限个连通分支；
它是完备化 $\widehat X$（沿 $f^{-1}(y)$）的底空间。若
$Z_i$（$1\leq i\leq n$）是这些连通分支，则显然
$\Gamma(\widehat X,\mathcal O_{\widehat X})$ 是各环
$\Gamma(Z_i,\mathcal O_{\widehat X})$ 的直积；其中每一环都不等于
$0$，因为单位截面不等于 $0$ 于 $\widehat X$ 的每一点。另一方面，将
（\hyperref[III.4.1.5-fr]{4.1.5}）应用于 $\mathcal F=\mathcal O_X$；
它沿 $f^{-1}(y)$ 的完备化是 $\mathcal O_{\widehat X}$，故可见
$\Gamma(\widehat X,\mathcal O_{\widehat X})$ 同构于关于
$\mathfrak m_y$-进拓扑的分离完备环 $\widehat{\mathcal O}_y$，即局部环
$\mathcal O_y$ 的完备化。于是它是局部环，不可能是若干不等于 $0$ 的
环的直积（否则它会有若干互异的极大理想）。所以 $n=1$，推论得证。
\end{proof}

\begin{corollary}[4.3.3]
\label{III.4.3.3-fr}
在（\hyperref[III.4.3.1-fr]{4.3.1}）的假设下，对每个 $y\in Y$，
纤维 $f^{-1}(y)$ 的连通分支集与纤维 $g^{-1}(y)$ 的有限点集一一对应，
其中 $g:Y'\to Y$ 是结构态射（换言之，后者是
$(f_*(\mathcal O_X))_y$ 的极大理想集）。
\end{corollary}

\begin{proof}
事实上，由于 $Y'$ 在 $Y$ 上有限，$g^{-1}(y)$ 是有限离散空间
（II，6.1.7）。又因 $f^{-1}(y)=f'^{-1}(g^{-1}(y))$，该推论由这一
观察和（\hyperref[III.4.3.1-fr]{4.3.1}）得出。
\end{proof}

这样便得到（\hyperref[III.4.3.1-fr]{4.3.1}）中所定义的 $Y$-预概形
$Y'$ 的一个重要解释。分解 $f=g\circ f'$ 对固有态射 $f$ 而言类似于
K.~Stein 对解析空间间全纯映射所得到的分解；以后我们称它为 $f$ 的
\emph{Stein 分解}。

\begin{remark}[4.3.4]
\label{III.4.3.4-fr}
设 $k$ 是域 $k(y)$ 的一个扩张。若预概形
\[
f^{-1}(y)\otimes_{k(y)}k=X\times_Y\operatorname{Spec}(k)
\]
连通，则 $f^{-1}(y)$ 也连通，因为后者是前者在投影态射下的像
（I，3.4.7）。对于预概形态射 $f:X\to Y$ 和点 $y\in Y$，若纤维
$f^{-1}(y)$ 对每个扩张 $k$（它扩张 $k(y)$）都使预概形
$f^{-1}(y)\otimes_{k(y)}k=X\times_Y\operatorname{Spec}(k)$ 连通，
则称该纤维\emph{几何连通}。

\oldpage[III]{132}

在（\hyperref[III.4.3.2-fr]{4.3.2}）的假设下，其结论可加强为：
各纤维 $f^{-1}(y)$ 实际上都\emph{几何连通}。为此，注意对每个扩张
$k$（它扩张 $k(y)$），都存在一个 Noether 局部环 $A$ 和一个局部同态
$\varphi:\mathcal O_y\to A$，使 $A$ 成为\emph{平坦} $\mathcal O_y$-模，
并使 $A$ 的剩余域作为 $k(y)$-扩张同构于 $k$（0，10.3.1）。令
$Y_1=\operatorname{Spec}(A)$，并令 $h:Y_1\to Y$ 为对应于 $\varphi$ 的
局部态射；它把闭点 $y_1$（$Y_1$ 的唯一闭点）映到 $y$（I，2.4.1）。置
$X_1=X\times_YY_1$、$f_1=f\times1_{Y_1}$；则 $f_1$ 是固有的
（II，5.4.2，(iii)），而 $f_1^{-1}(y_1)$ 作为 $k(y_1)$-预概形同构于
$X\times_Y\operatorname{Spec}(k)$。因此，只需证明
$f_{1*}(\mathcal O_{X_1})=\mathcal O_{Y_1}$，便可将
（\hyperref[III.4.3.2-fr]{4.3.2}）应用于 $f_1$。由（I，2.4.2）和
（\hyperref[III.1.4.15-fr]{1.4.15.5}），$h$ 是平坦态射
\SourceCorrectionNote{法文印本此处误作“$g$ 是平坦的”；Canon 稳定决定
EGA-III1-4.3.4-P476-G-H-FLAT-001 核定基变换态射应为 $h$。本译采用
已核定读法，并保留可逆来源说明。}；所以
\[
f_{1*}(\mathcal O_{X_1})
=h^*(f_*(\mathcal O_X))
=h^*(\mathcal O_Y)
=\mathcal O_{Y_1}
\]
这里应用了（\hyperref[III.1.4.15-fr]{1.4.15}）的 $q=0$ 情形。

在一般情形（\hyperref[III.4.3.1-fr]{4.3.1}）下，同样的推理表明
（沿用（\hyperref[III.4.3.1-fr]{4.3.1}）的记号）
\[
f_{1*}(\mathcal O_{X_1})=h^*(g_*(\mathcal O_{Y'})),
\]
并且 Stein 分解 $f_1=g_1\circ f'_1$（$f_1$ 的 Stein 分解）满足
$g_1=g\times1_{Y_1}$（II，1.5.2），相应的有限 $Y_1$-概形为
$Y'_1=Y'\times_YY_1$。考虑到纤维的传递性（I，3.6.4），由
（\hyperref[III.4.3.3-fr]{4.3.3}）可知，$f_1^{-1}(y_1)$ 的连通分支数
等于集合
\[
g_1^{-1}(y_1)=g^{-1}(y)\otimes_{k(y)}k
\]
的元素数。若取 $k$ 为 $k(y)$ 的一个代数闭扩张，该数与所取代数闭
扩张无关，并等于 $g^{-1}(y)$ 的\emph{几何点数}（I，6.4.7），也就是
\emph{可分次数之和}
\[
\sum_i[k(y'_i):k(y)]_s,
\]
其中 $y'_i$ 遍历有限集 $g^{-1}(y)$。这个数也称为 $f^{-1}(y)$ 的
\emph{几何连通分支数}。注意，各 $k(y'_i)$ 正是半局部环
$(f_*(\mathcal O_X))_y$ 的剩余域。
\end{remark}

\begin{proposition}[4.3.5]
\label{III.4.3.5-fr}
设 $X$ 和 $Y$ 是两个局部 Noether 整预概形，$f:X\to Y$ 是固有支配
态射。对每个 $y\in Y$，$f^{-1}(y)$ 的连通分支数不超过
$\mathcal O'_y$ 的极大理想数；这里前者是
$\mathcal O_y$ 在有理函数域 $\mathrm R(X)$ 中的整闭包。
\end{proposition}

\begin{proof}
事实上，对每个开集 $U$（$Y$ 的开集），
\[
\Gamma(U,f_*(\mathcal O_X))=\Gamma(f^{-1}(U),\mathcal O_X)
\]
是各局部环 $\mathcal O_x$ 的交，其中 $x\in f^{-1}(U)$
（I，8.2.1.1）。立即可知，茎 $(f_*(\mathcal O_X))_y$ 是
$\mathrm R(X)$ 的一个包含 $\mathcal O_y$ 的子环。此外，因为
$f_*(\mathcal O_X)$ 是凝聚 $\mathcal O_Y$-模，
$(f_*(\mathcal O_X))_y$ 是有限型 $\mathcal O_y$-模，因而包含于
$\mathcal O'_y$。已知（[13]，卷 I，第 257、259 页），这样一个环
$A$ 的每个极大理想都是 $A$ 与 $\mathcal O'_y$ 的某个极大理想的交，
命题遂得证。
\end{proof}

\begin{definition}[4.3.6]
\label{III.4.3.6-fr}
若一个整局部环的整闭包仍是局部环，则称该环为\emph{单支的}。
若点 $y$ 所在的整预概形为 $Y$，且局部环 $\mathcal O_y$ 是单支的，则称
该点为单支点（特别地，当 $Y$ 在 $y$ 处正规时便是如此）。
\end{definition}

设 $A$ 是整局部环，$K$ 是其分式域。$A$ 为单支环，当且仅当每个子环
$A_1$（它是 $K$ 的子环、包含 $A$ 且为有限 $A$-代数）都是局部环。事实上，
令 $A'$ 为 $A$ 的整闭包；由 Cohen--Seidenberg 第一定理
（Bourbaki，\emph{交换代数}，第五章，\S~2，第 1 号，定理 1），
$A_1$ 的每个极大理想都是 $A'$ 的某个极大理想的迹。因此，若 $A'$
是局部环，$A_1$ 也是局部环。反之，$A'$ 是各有限 $A$-子代数
$A_\alpha$（它们属于 $A'$）所成递增滤过族的归纳极限；若每个
$A_\alpha$ 都是局部环，则由上述同一论证，$A_\alpha$ 的极大理想是在
$A_\alpha$ 上取 $A_\beta$ 的极大理想之迹（当 $A_\alpha\subset A_\beta$ 时），
所以 $A'$ 是
局部环（0，10.3.1.3）。

\oldpage[III]{133}

还应注意，若 Noether 局部环 $A$ 的完备化是整环（此时称 $A$
\emph{解析整}），则 $A$ 是单支的。事实上，令 $\mathfrak m$ 为 $A$
的极大理想，$K$ 为其分式域，$K'$ 为 $\widehat A$ 的分式域；于是
\[
K'=K\otimes_A\widehat A.
\]
令 $A'_F$ 为一个有限 $A$-子代数，它属于 $K$。子环 $B_F$（$K'$ 的
子环，由 $\widehat A$ 和 $A'_F$ 生成）同构于
$A'_F\otimes_A\widehat A$；它是
有限型 $\widehat A$-模，并且是 $A'_F$ 关于 $\mathfrak m$-进拓扑的
完备化（0\textsubscript{I}，7.3.3 和 7.3.6）。由于 $A'_F$ 是半局部环
（Bourbaki，\emph{交换代数}，第四章，\S~2，第 5 号，命题 9 的
推论 3），而它的完备化是整环，$A'_F$ 只能有一个极大理想
$\mathfrak m'_F$，且 $\mathfrak m'_F\cap A=\mathfrak m$；断言得证。

\begin{corollary}[4.3.7]
\label{III.4.3.7-fr}
在（\hyperref[III.4.3.5-fr]{4.3.5}）的假设下，设
$\mathrm R(Y)$ 在 $\mathrm R(X)$ 中的代数闭包的可分次数为 $n$，且
$y\in Y$ 是单支点。则纤维 $f^{-1}(y)$ 至多有 $n$ 个连通分支。
特别地，若 $\mathrm R(Y)$ 在 $\mathrm R(X)$ 中的代数闭包是
$\mathrm R(Y)$ 的纯不可分扩张
\SourceCorrectionNote{法文印本此处误称该代数闭包为 $\mathrm R(X)$ 的
纯不可分扩张；Canon 稳定决定 EGA-III1-4.3.7-P477-RADICIAL-BASE-001
核定底域为 $\mathrm R(Y)$。本译采用已核定读法，并保留可逆来源说明。}，
则 $f^{-1}(y)$ 连通。
\end{corollary}

\begin{proof}
事实上，令 $\mathcal O'_y$ 为 $\mathcal O_y$ 的整闭包；整闭包
$\mathcal O''_y$（即 $\mathcal O_y$ 在 $\mathrm R(X)$ 中的整闭包）也是
$\mathcal O'_y$ 的整闭包。已知若
$\mathcal O'_y$ 是局部环，则 $\mathcal O''_y$ 是一个半局部环，其
极大理想数至多为 $n$（[13]，卷 I，第 289 页，定理 22）。

这个推论基本上正是 Zariski 对代数概形表述其“连通性定理”的形式。
\end{proof}

\begin{remark}[4.3.8]
\label{III.4.3.8-fr}
若在（\hyperref[III.4.3.7-fr]{4.3.7}）的假设之外再假设 $Y$ 在
$y$ 处正规，则纤维 $f^{-1}(y)$ \emph{几何连通}，因为沿用
（\hyperref[III.4.3.4-fr]{4.3.4}）的记号，$g^{-1}(y)$ 只有一个点
$y'$，且 $k(y')$ 是 $k(y)$ 的纯不可分扩张。
\end{remark}

\begin{definition}[4.3.9]
\label{III.4.3.9-fr}
给定局部 Noether 预概形 $Y$，若有限型态射 $f:X\to Y$ 满足：对每个
局部 Noether 不可约预概形 $Y'$ 以及每个支配态射 $g:Y'\to Y$，
$X'=X\times_YY'$ 的每个不可约分支都支配 $Y'$，则称它\emph{普遍开}。
\end{definition}

若 $Y$ 不可约，这等价于说：若 $\eta$、$\eta'$ 分别为 $Y$ 和 $Y'$
的泛点（从而 $g(\eta')=\eta$），并置 $f'=f_{(Y')}$，则 $X'$ 的每个
不可约分支都与 $f'^{-1}(\eta')$ 相交（0\textsubscript{I}，2.1.8）。
因此，这还蕴含：对每个开集 $U$（它属于 $Y$），限制态射
$f^{-1}(U)\to U$（$f$ 的限制）是普遍开的。

\begin{corollary}[4.3.10]
\label{III.4.3.10-fr}
设 $X$、$Y$ 是两个局部 Noether 整预概形，$f:X\to Y$ 是固有、支配
且普遍开的态射。若 $\mathrm R(Y)$ 在 $\mathrm R(X)$ 中的代数闭包是
$\mathrm R(Y)$ 的纯不可分扩张，则每个纤维 $f^{-1}(y)$（$y\in Y$）
都几何连通。
\end{corollary}

\begin{proof}
可只考虑 $Y=\operatorname{Spec}(B)$ 的情形，其中 $B$ 是 Noether 整环。
由（II，7.1.7），存在一个整闭 Noether 局部环 $A$，它支配
$\mathcal O_y$，并以 $\mathrm R(Y)$ 为分式域。令
$Y'=\operatorname{Spec}(A)$，并令 $h:Y'\to Y$ 为典范单射
$B\to A$ 所对应的态射；它是双有理的（因而支配）。此外，若 $y'$
是 $Y'$ 的唯一闭点，则 $h(y')=y$。令 $X'=X\times_YY'$、
$f'=f\times1_{Y'}$；再分别以 $\eta$、$\eta'$、$\xi$ 表示 $Y$、$Y'$、
$X$ 的泛点，于是 $f(\xi)=\eta$ 且 $h^{-1}(\eta)=\{\eta'\}$。
此外，$k(\eta)=k(\eta')=\mathrm R(Y)$，所以 $f'^{-1}(\eta')$ 同构于
$f^{-1}(\eta)$（I，3.6.4）；特别地，由于 $\xi$ 是
$f^{-1}(\eta)$ 的泛点（0\textsubscript{I}，2.1.8），
$f'^{-1}(\eta')$ 只有一个泛点。按假设，$X'$ 的每个不可约分支的
泛点都在 $f'^{-1}(\eta')$ 中，故 $X'$ 必不可约；它的泛点 $\xi'$
就是 $f'^{-1}(\eta')$ 的泛点，并且 $k(\xi')=k(\xi)$。

\oldpage[III]{134}

置 $X''=X'_{\mathrm{red}}$、$f''=f'_{\mathrm{red}}$；于是 $X''$ 是
Noether 整预概形，$f''$ 是固有的（II，5.4.6），而纤维
$f'^{-1}(y')$ 与 $f''^{-1}(y')$ 的底空间相同。此外，
$\mathrm R(X'')=k(\xi')=\mathrm R(X)$，所以 $f''$ 满足
（\hyperref[III.4.3.8-fr]{4.3.8}）的假设，且 $f''^{-1}(y')$ 几何连通。
现在令 $k$ 为 $k(y)$ 的任意扩张；存在一个扩张 $k_1$（它扩张 $k(y)$），
使 $k(y')$ 与 $k$ 都可视为它的子扩张（Bourbaki，\emph{代数}，
第五章，\S~4，命题 2）。按假设，
$f''^{-1}(y')\times_{Y'}\operatorname{Spec}(k_1)$ 连通，并且它与
$f'^{-1}(y')\times_{Y'}\operatorname{Spec}(k_1)$ 有同一既约预概形
（I，5.1.8），所以后者也连通。后者又同构于
$f^{-1}(y)\times_Y\operatorname{Spec}(k_1)$（I，3.6.4），故该预概形
连通；\emph{a fortiori}，根据（\hyperref[III.4.3.4-fr]{4.3.4}）开头的
注释，$f^{-1}(y)\times_Y\operatorname{Spec}(k)$ 也连通，证明完毕。
\end{proof}

\begin{env}[4.3.11]
\label{III.4.3.11-fr}
\emph{注}（4.3.11）。---
\begin{enumerate}
\item[(i)] 上述论证实质上出自 Zariski [20]；不同之处仅在于，他可以取
$A$ 为 $\mathcal O_y$ 的整闭包，因为在经典代数几何的局部环情形，这个
整闭包是 Noether 环。另一方面，Zariski 证明了：若 $Y$ 是
$\mathbf P_k^r$（域 $k$ 上的射影空间）的 Chow 簇，而 $X$ 是
$\mathbf P_k^r\times_kY$ 中定义 $\mathbf P_k^r$ 与 $Y$ 之间 Chow 对应的
闭部分，则投影 $X\to Y$ 是普遍开态射（\emph{同上}，第 82 页的引理）。
看来，这正是某些应用中实际用到的“Chow 坐标”的唯一形式性质。因此在
这类情形中，宜以如下语言取代“射影空间中循环的特殊化”：固有态射的
纤维（可再假设该态射普遍开，或服从其他类似的局部正则性限制）。

\item[(ii)] 在第四章中我们将看到，普遍开态射 $f:X\to Y$ 还可按如下
方式定义（这也说明了这一术语）：对每个态射 $Z\to Y$，态射
$f_{(Z)}:X_{(Z)}\to Z$ 都是开映射。此外还可证明，若 $f$ 满足
（\hyperref[III.4.3.10-fr]{4.3.10}）的假设，那么当 $y,y'$ 是 $Y$ 的
两个点且 $y$ 是 $y'$ 的特殊化时，$f^{-1}(y)$ 的几何连通分支数不超过
$f^{-1}(y')$ 的几何连通分支数。
\end{enumerate}
\end{env}

\begin{corollary}[4.3.12]
\label{III.4.3.12-fr}
在（\hyperref[III.4.3.5-fr]{4.3.5}）的假设下，再假设
$\mathrm R(Y)$ 在 $\mathrm R(X)$ 中代数闭，并令 $y$ 为 $Y$ 的一个
正规点。则 $f^{-1}(y)$ 几何连通，并且存在开邻域 $U$（点 $y$ 在
$Y$ 中的邻域），使 $f_*(\mathcal O_X|f^{-1}(U))$ 同构于 $\mathcal O_Y|U$。
特别地，若 $Y$ 正规（且 $\mathrm R(Y)$ 在 $\mathrm R(X)$ 中代数闭），
则 $f_*(\mathcal O_X)$ 同构于 $\mathcal O_Y$。
\end{corollary}

\begin{proof}
关于 $f^{-1}(y)$ 的第一个断言是
（\hyperref[III.4.3.8-fr]{4.3.8}）的特殊情形。由此可知，若

\oldpage[III]{135}

$X\xrightarrow{f'}Y'\xrightarrow{g}Y$ 是 $f$ 的 Stein 分解
（\hyperref[III.4.3.3-fr]{4.3.3}），则 $g^{-1}(y)$ 只有一个点 $y'$；
此外有
\[
\mathcal O_y\subset\mathcal O_{y'}=(f_*(\mathcal O_X))_y
\subset\mathrm R(X),
\]
而且由于 $\mathcal O_{y'}$ 在 $\mathcal O_y$ 上有限，所以它在
$\mathrm R(Y)$ 上是代数的；因而按假设它包含于 $\mathrm R(Y)$
\SourceCorrectionNote{法文印本把“在 $\mathcal O_y$ 上有限”误推成
“更在域 $\mathrm R(Y)$ 上有限”；Canon 稳定决定
EGA-III1-4.3.12-P479-A-FORTIORI-RY-001 核定这里只能推出代数性。
本译采用已核定读法，并保留可逆来源说明。}。又因 $y$
正规，必有 $\mathcal O_{y'}=\mathcal O_y$；于是 $g$ 在点 $y'$ 处是
局部同构（I，6.5.4），这就证明了推论的第一部分。第二部分由第一部分
得出，因为附加假设蕴含 $g$ 是双射，并在 $Y'$ 每一点的某个邻域上
都是局部同构，因而是同构。
\end{proof}

由于（\hyperref[III.4.3.7-fr]{4.3.7}）是在概形的框架中建立的，因而
可作如下应用：

\begin{proposition}[4.3.13]
\label{III.4.3.13-fr}
设 $A$ 是单支 Noether 局部环，$\mathfrak a$ 是 $A$ 的一个定义理想，
$A_0=A/\mathfrak a$，并令
$S=\operatorname{gr}_{\mathfrak a}(A)$ 为 $A$ 关于 $\mathfrak a$-预进
滤过的相伴分次环。$S$ 是分次 $A_0$-代数，由 $S_1$ 生成，其中 $S_1$
是有限型 $A_0$-模。则 $\operatorname{Proj}(S)$ 是连通 $A_0$-概形。
\end{proposition}

\begin{proof}
令 $\mathfrak m$ 为 $A$ 的极大理想；$Y=\operatorname{Spec}(A)$ 是整
概形，其中点 $y$（对应于 $\mathfrak m$）是唯一闭点。按假设，有
$\mathfrak m^p\subset\mathfrak a\subset\mathfrak m$，其中 $p$ 为某个整数，故
$V(\mathfrak a)=\{\mathfrak m\}$。令
$S'=\bigoplus_{n\geq0}\mathfrak a^n$，并令 $X=\operatorname{Proj}(S')$；
这是一个 $Y$-概形，由爆破理想 $\mathfrak a$ 得到。$X$ 是整的，而结构
态射 $f:X\to Y$ 是双有理的（II，8.1.4），且显然是射影的。因此可应用
（\hyperref[III.4.3.7-fr]{4.3.7}），得到 $f^{-1}(y)$ 连通。另一方面，
$f^{-1}(y)$ 的底空间就是 $\operatorname{Proj}(S'\otimes_AA_0)$ 的底空间
（I，3.6.1 和 II，2.8.10）；由定义 $S'\otimes_AA_0=S$，命题得证。
\end{proof}

\subsection{Zariski“主定理”}
\label{subsection:III.4.4-fr}

\begin{proposition}[4.4.1]
\label{III.4.4.1-fr}
设 $Y$ 是局部 Noether 预概形，$f:X\to Y$ 是固有态射。令 $X'$ 为
满足 $x\in X$ 且在纤维 $f^{-1}(f(x))$ 内孤立的点所成的集合。则
$X'$ 是 $X$ 的开集；若 $f=gf'$ 是 $f$ 的 Stein 分解
（\hyperref[III.4.3.3-fr]{4.3.3}），则 $f'$ 在 $X'$ 上的限制把 $X'$
同构地映到 $Y'$ 的某个开集 $U$ 上所诱导的子预概形，并且
$X'=f'^{-1}(U)$。
\end{proposition}

\begin{proof}
由于 $g^{-1}(f(x))$ 有限且离散
（\hyperref[III.4.3.3-fr]{4.3.3} 和 II，6.1.7），点 $x$ 在
$f^{-1}(f(x))$ 中孤立，当且仅当它在 $f'^{-1}(f'(x))$ 中孤立。因此可
限于 $f'=f$ 的情形，此时 $f_*(\mathcal O_X)=\mathcal O_Y$。若
$x\in X'$，则纤维 $f^{-1}(f(x))$ 由
（\hyperref[III.4.3.2-fr]{4.3.2}）连通，故它必只含点 $x$。由于 $f$
是闭映射，对 $X$ 中点 $x$ 的任意开邻域 $V$，$f(X-V)$ 是 $Y$ 中
不含 $y=f(x)$ 的闭集，因为 $f^{-1}(y)=\{x\}$。若 $U$ 是
$f(X-V)$ 在 $Y$ 中的补集，则 $f^{-1}(U)\subset V$；由此可知，$y$
的一族基本开邻域在 $f$ 下的逆像构成 $x$ 的一族基本开邻域。条件
$f_*(\mathcal O_X)=\mathcal O_Y$ 以及层的直接像之定义
（0\textsubscript{I}，3.4.1 和 4.2.1）于是蕴含：若
$f=(\psi,\theta)$，则同态
$\theta_y^\#:\mathcal O_y\to\mathcal O_x$ 是同构。因而存在 $x$ 的开邻域
$V$ 及 $y$ 的开邻域 $U$，使 $f$ 在 $V$ 上的限制是从 $V$ 到 $U$ 的
同构（I，6.5.4）；并且根据刚才的论证，还可以假设

\oldpage[III]{136}

$f^{-1}(U)=V$。于是由定义立刻得到 $V\subset X'$，证明完毕。
\end{proof}

下一个命题在代数概形的情形由 Chevalley 证明：

\begin{proposition}[4.4.2]
\label{III.4.4.2-fr}
设 $Y$ 是局部 Noether 预概形，$f:X\to Y$ 是态射。下列条件等价：
\begin{enumerate}
\item[a)] $f$ 是有限态射。
\item[b)] $f$ 是仿射且固有的。
\item[c)] $f$ 是固有的，并且对每个 $y\in Y$，$f^{-1}(y)$ 是有限集。
\end{enumerate}
\end{proposition}

\begin{proof}
已知 a) 蕴含 b)（II，6.1.2 和 6.1.11）。若 $f$ 固有且仿射，则态射
$f^{-1}(y)\to\operatorname{Spec}(k(y))$ 也如此
（II，1.6.2，(iii) 和 5.4.2，(iii)）。把有限性定理
（\hyperref[III.3.2.1-fr]{3.2.1}）应用于 $f^{-1}(y)$ 的结构层，可知
$f^{-1}(y)=\operatorname{Spec}(A)$，其中 $A$ 是有限 $k(y)$-代数；故
$f^{-1}(y)$ 是有限集（II，6.1.7），从而 b) 蕴含 c)。最后，
$f^{-1}(y)$ 是 $k(y)$ 上的代数预概形；若 $f^{-1}(y)$ 的点集有限，则
空间 $f^{-1}(y)$ 离散
（I，6.4.4）。沿用（\hyperref[III.4.4.1-fr]{4.4.1}）的记号，此时
$X'=X$，且 $f':X\to Y'$ 是同构；又因 $g$ 是有限态射，c) 蕴含 a)。
\end{proof}

\begin{theorem}[4.4.3]
\label{III.4.4.3-fr}
\textnormal{（Zariski“主定理”。）}
设 $Y$ 是 Noether 预概形，$f:X\to Y$ 是拟射影态射，$X'$ 是所有满足
$x\in X$ 且在纤维 $f^{-1}(f(x))$ 内孤立的点所成的集合。则 $X'$ 是
$X$ 的开子集，并且 $X'$ 上的诱导子预概形同构于某个在 $Y$ 上有限的
$Y$-预概形 $Y'$ 的一个开子集上所诱导的子预概形。
\end{theorem}

\begin{proof}
由假设，存在射影 $Y$-预概形 $Z$，使 $X$ 与 $Z$ 的一个开集上所诱导的
子预概形在 $Y$ 上同构（II，5.3.2 和 5.5.1）。因此只须证明 $f$ 是
射影态射的情形；此时该态射固有（II，5.5.3），结论立即由
（\hyperref[III.4.4.1-fr]{4.4.1}）得到。
\end{proof}

\begin{remark}[4.4.4]
\label{III.4.4.4-fr}
若 $X$ 既约（分别地，不可约且 $X'$ 非空），则在
（\hyperref[III.4.4.3-fr]{4.4.3}）的陈述中，可要求 $Y'$ 既约
（分别地，不可约）。事实上，总可以用 $X'$ 在 $Y'$ 中的概形闭包
$\overline{X'}$ 取代 $Y'$（I，9.5.11 和 II，6.1.5，(i)、(ii)）；已知
若 $X'$ 既约，则 $\overline{X'}$ 也既约（I，9.5.9，(i)）。另一方面，
若 $X'$ 非空且 $X$ 不可约，则前者不可约，因而 $\overline{X'}$
也不可约。
\end{remark}

\begin{corollary}[4.4.5]
\label{III.4.4.5-fr}
设 $Y$ 是局部 Noether 概形，$f:X\to Y$ 是有限型态射，$x$ 是 $X$
中在其纤维 $f^{-1}(f(x))$ 内孤立的点。则 $x$ 在 $X$ 中有一个开邻域，
它同构于某个在 $Y$ 上有限的 $Y$-预概形的一个开子集。
\end{corollary}

\begin{proof}
令 $y=f(x)$；取 $y$ 在 $Y$ 中的仿射开邻域 $U$，以及 $x$ 在 $X$
中的仿射开邻域 $V$，使它包含于 $f^{-1}(U)$。由于 $Y$ 分离，单射
$U\to Y$ 是仿射态射（II，1.6.3）；又因 $V$ 在 $U$ 上仿射
（\emph{同上}），$f$ 在 $V$ 上的限制是仿射态射 $V\to Y$
（II，1.6.2，(ii)）。\emph{a fortiori}，这个限制是拟射影的，因为它是
有限型的（I，6.3.5 和 II，5.3.4，(i)）。于是只须把定理
（\hyperref[III.4.4.3-fr]{4.4.3}）应用于这个限制。
\end{proof}

推论（\hyperref[III.4.4.5-fr]{4.4.5}）可用交换代数的语言表述如下：

\oldpage[III]{137}

\begin{corollary}[4.4.6]
\label{III.4.4.6-fr}
设 $A$ 是 Noether 环，$B$ 是有限型 $A$-代数，$\mathfrak q$ 是 $B$
的素理想，$\mathfrak p$ 是它在 $A$ 中的逆像。假设 $\mathfrak q$ 在
所有逆像为 $\mathfrak p$ 的 $B$ 的素理想所成集合中既是极大的又是
极小的。则存在 $g\in B-\mathfrak q$、有限 $A$-代数 $A'$ 及元素
$f'\in A'$，使 $A$-代数 $B_g$ 与 $A'_{f'}$ 同构。
\end{corollary}

\begin{proof}
只须把（\hyperref[III.4.4.5-fr]{4.4.5}）应用于
$Y=\operatorname{Spec}(A)$ 与 $X=\operatorname{Spec}(B)$；关于
$\mathfrak q$ 的假设正好表示它在其纤维中是孤立点（I，1.1.7）。
\end{proof}

由此得到下述表面上较弱的结果：

\begin{corollary}[4.4.7]
\label{III.4.4.7-fr}
设 $A$ 是 Noether 局部环，$B$ 是有限型 $A$-代数，$\mathfrak n$ 是
$B$ 的素理想，其在 $A$ 中的逆像为极大理想 $\mathfrak m$。假设
$\mathfrak n$ 是 $B$ 的极大理想，并且在所有逆像为 $\mathfrak m$ 的
$B$ 的素理想所成集合中是极小的（这也表示 $B\mathfrak m$ 对
$\mathfrak n$ 是准素的）。则存在有限 $A$-代数 $A'$ 及 $A'$ 的极大
理想 $\mathfrak m'$（它在 $A$ 中的逆像为 $\mathfrak m$），使
$B_{\mathfrak n}$ 与 $A'_{\mathfrak m'}$ 作为 $A$-代数同构。
\end{corollary}

下述（\hyperref[III.4.4.7-fr]{4.4.7}）的特殊情形有时也称为“主定理”：

\begin{corollary}[4.4.8]
\label{III.4.4.8-fr}
在（\hyperref[III.4.4.7-fr]{4.4.7}）的条件下，再假设 $A$ 与 $B$ 都是
整环，并有同一个分式域 $K$。若 $A$ 整闭，则 $B=A$。
\end{corollary}

\begin{proof}
事实上，注（\hyperref[III.4.4.4-fr]{4.4.4}）表明，在应用
（\hyperref[III.4.4.7-fr]{4.4.7}）时，可以假设 $A'$ 是以 $K$ 为
分式域的整环。关于 $A$ 的假设于是给出 $A'=A$，从而
$B_{\mathfrak n}=A$；由于 $A\subset B\subset B_{\mathfrak n}$，便有
$A=B$。
\end{proof}

陈述（\hyperref[III.4.4.8-fr]{4.4.8}）正是 Zariski 赋予其“主定理”的
形式（这里把它推广到任意 Noether 整局部环）。

前面的推论都是全局结果（\hyperref[III.4.4.3-fr]{4.4.3}）的局部变式。
下面给出另一个全局性质的推论：

\begin{corollary}[4.4.9]
\label{III.4.4.9-fr}
设 $X$ 与 $Y$ 是两个局部 Noether 整预概形，$f:X\to Y$ 是分离、局部
有限型且双有理的态射。再假设 $Y$ 正规，并且对所有 $y\in Y$，纤维
$f^{-1}(y)$ 都有限。
\SourceCorrectionNote{法文印本及外交式转录在（III，4.4.9）中只明列 $Y$，
并写“有限型”；Canon 校正现行法文 R2 明列 $X$ 与 $Y$，并改为“局部有限型”。
本译采用该现行可逆读法，外交式法文保持不变。}
则 $f$ 是开浸入；若 $f$ 还是闭映射（特别地，若 $f$ 固有），则 $f$
是同构。
\end{corollary}

\begin{proof}
任取 $x\in X$，置 $y=f(x)$。由于 $f^{-1}(y)$ 是 $k(y)$ 上的局部有限型
概形，它有限这一假设蕴含它离散（I，6.4.4）。
\SourceCorrectionNote{法文印本及外交式转录此处写“$k(y)$ 上的代数概形”；
Canon 校正现行法文 R2 写“$k(y)$ 上的局部有限型概形”。本译采用现行读法。}
此外，$\mathcal O_y$ 整闭，而
$\mathcal O_x$ 与 $\mathcal O_y$ 有同一个分式域（I，7.1.5）。故可应用
（\hyperref[III.4.4.8-fr]{4.4.8}）；若 $f=(\psi,\theta)$，则同态
$\theta_y^\#:\mathcal O_y\to\mathcal O_x$ 是双射。由（I，6.5.4）可知
$f$ 是局部同构。又因 $f$ 分离且 $X$ 是整的，$f$ 是开浸入
（I，8.2.8）。最后一个断言由 $f$ 支配得到。
\end{proof}

\begin{proposition}[4.4.10]
\label{III.4.4.10-fr}
设 $Y$ 是局部 Noether 预概形，$f:X\to Y$ 是局部有限型态射。令 $X'$
为所有在其纤维 $f^{-1}(f(x))$ 内孤立的点 $x\in X$ 所成的集合，则
该集合是 $X$ 的开集。
\end{proposition}

\begin{proof}
问题关于 $X$ 与 $Y$ 都是局部的，故可假设 $X$、$Y$ 是 Noether 仿射
概形且 $f$ 是有限型的。此时 $f$ 是有限型仿射态射，因而拟射影
（II，5.3.4，(i)）；只须应用
（\hyperref[III.4.4.3-fr]{4.4.3}）。
\end{proof}

\oldpage[III]{138}

\begin{corollary}[4.4.11]
\label{III.4.4.11-fr}
设 $Y$ 是局部 Noether 预概形，$f:X\to Y$ 是固有态射。令 $U$ 为所有
满足 $y\in Y$ 且 $f^{-1}(y)$ 离散的点所成的集合，则该集合是 $Y$
的开集，且 $f$ 的限制 $f^{-1}(U)\to U$ 是有限态射。特别地，固有的
拟有限态射 $X\to Y$ 必有限。
\end{corollary}

\begin{proof}
$U$ 在 $Y$ 中的补集是 $X-X'$ 在 $f$ 下的像；由
（\hyperref[III.4.4.10-fr]{4.4.10}），后者在 $X$ 中闭，而 $f$ 是
闭映射，故 $U$ 是开集。又由（II，6.2.2），对每个 $y\in U$，
$f^{-1}(y)$ 都有限；$f$ 的限制 $f^{-1}(U)\to U$ 是固有的
（II，5.4.1），故由（\hyperref[III.4.4.2-fr]{4.4.2}）它是有限态射。
\end{proof}

\begin{env}[4.4.12]
\label{III.4.4.12-fr}
\emph{注}（4.4.12）。---
\begin{enumerate}
\item[(i)] 如（II，6.2.7）中所预告，我们将在第五章证明：若 $Y$ 局部
Noether，则每个拟有限且分离的态射 $f:X\to Y$ 都是拟仿射的，因而是
拟射影的。由此将知，在主定理（\hyperref[III.4.4.3-fr]{4.4.3}）中，
只假设 $f$ 分离且有限型，结论仍成立。事实上，由
（\hyperref[III.4.4.10-fr]{4.4.10}），$X'$ 在 $X$ 中开；又因 $X$
局部 Noether，$f$ 在 $X'$ 上的限制仍是有限型的（I，6.3.5），并且由
$X'$ 的定义是拟有限的，显然也是分离的。因此可把
（\hyperref[III.4.4.3-fr]{4.4.3}）应用于这个限制，结论随即得到。

\item[(ii)] 在第四章中，我们将利用维数理论给出
（\hyperref[III.4.4.10-fr]{4.4.10}）的一个更初等的证明。
\end{enumerate}
\end{env}

\subsection{同态模的完备化}
\label{subsection:III.4.5-fr}

\begin{proposition}[4.5.1]
\label{III.4.5.1-fr}
设 $A$ 是 Noether 环，$\mathfrak J$ 是 $A$ 的一个理想，$X$ 是有限型
$A$-预概形，$\mathcal F$、$\mathcal G$ 是两个凝聚 $\mathcal O_X$-模，
且它们的支集之交在 $Y=\operatorname{Spec}(A)$ 上固有（II，5.4.10）。
则对每个整数 $n\geq0$，
$\operatorname{Ext}^n_{\mathcal O_X}(X;\mathcal F,\mathcal G)$ 是有限型 $A$-模，
并且它关于 $\mathfrak J$-预进制拓扑的分离完备化，典范地等同于（采用
（\hyperref[III.4.1.7-fr]{4.1.7}）的记号）
$\operatorname{Ext}^n_{\mathcal O_{\widehat X}}(\widehat X;\widehat{\mathcal F},
\widehat{\mathcal G})$。
\end{proposition}

\begin{proof}
已知（T，4.2），存在一个双正则谱序列
$\mathrm E(\mathcal F,\mathcal G)$，其收敛项是
$\operatorname{Ext}_{\mathcal O_X}(X;\mathcal F,\mathcal G)$，而其
$\mathrm E_2$ 项由下式给出：
\[
\mathrm E_2^{pq}=\mathrm H^p\bigl(X,
\mathcal{E}\!xt^q_{\mathcal O_X}(\mathcal F,\mathcal G)\bigr).
\]
已知 $\mathcal{E}\!xt^q_{\mathcal O_X}(\mathcal F,\mathcal G)$ 是凝聚
$\mathcal O_X$-模（0，12.3.3），其支集包含于 $\mathcal F$ 与
$\mathcal G$ 的支集之交（T，4.2.2），因而在 $Y$ 上固有（II，5.4.10）。
由（\hyperref[III.3.2.4-fr]{3.2.4}）可知，$\mathrm E_2^{pq}$ 是有限型
$A$-模，从而（0，11.1.8）该谱序列的所有项 $\mathrm E_r^{pq}$ 及其
收敛项也同样如此。另一方面，若
$i:\widehat X\to X$ 是典范态射，则 $\widehat{\mathcal F}$ 与
$\widehat{\mathcal G}$ 分别典范地等同于 $i^*(\mathcal F)$ 与
$i^*(\mathcal G)$，并且 $i$ 平坦（I，10.8.8 与 10.8.9）。此时已知
（0，12.3.4），对每个 $q\geq0$，存在典范的 $i$-态射
\[
u_q:\mathcal{E}\!xt^q_{\mathcal O_X}(\mathcal F,\mathcal G)
\longrightarrow
\mathcal{E}\!xt^q_{\mathcal O_{\widehat X}}
(\widehat{\mathcal F},\widehat{\mathcal G}),
\]
而相应的 $\mathcal O_{\widehat X}$-同态
\[
u_q^\#:i^*\bigl(\mathcal{E}\!xt^q_{\mathcal O_X}(\mathcal F,\mathcal G)\bigr)
\longrightarrow
\mathcal{E}\!xt^q_{\mathcal O_{\widehat X}}
(\widehat{\mathcal F},\widehat{\mathcal G})
\]
是同构（0，12.3.5）；换言之（I，10.8.8），
$\mathcal{E}\!xt^q_{\mathcal O_{\widehat X}}
(\widehat{\mathcal F},\widehat{\mathcal G})$ 典范地等同于完备化
$(\mathcal{E}\!xt^q_{\mathcal O_X}(\mathcal F,\mathcal G))^\wedge$（采用
（\hyperref[III.4.1.7-fr]{4.1.7}）的记号）。于是由比较定理
（\hyperref[III.4.1.10-fr]{4.1.10}），对每个 $p\geq0$，
\[
\mathrm H^p\bigl(\widehat X,
\mathcal{E}\!xt^q_{\mathcal O_{\widehat X}}
(\widehat{\mathcal F},\widehat{\mathcal G})\bigr)
\]
典范地等同于
$\mathrm H^p\bigl(X,\mathcal{E}\!xt^q_{\mathcal O_X}
(\mathcal F,\mathcal G)\bigr)$ 关于 $\mathfrak J$-预进制拓扑的分离完备化。

\oldpage[III]{139}

用 $\mathrm E(\widehat{\mathcal F}, \widehat{\mathcal G})$ 表示（T，4.2）
中对 $\widehat{\mathcal F}$ 与 $\widehat{\mathcal G}$ 定义的双正则谱序列；
若 $\widehat A$ 表示 $A$ 关于 $\mathfrak J$-预进制拓扑的分离完备化，
则在典范同构的意义下有
\[
\mathrm E_2^{pq}(\widehat{\mathcal F},\widehat{\mathcal G})
=\mathrm E_2^{pq}(\mathcal F,\mathcal G)\otimes_A\widehat A
\qquad (0\textsubscript{I}, 7.3.3).
\]

另一方面，已知平坦态射 $i$ 的给定定义了谱序列的典范同态
\[
\varphi:\mathrm E(\mathcal F,\mathcal G)\longrightarrow
\mathrm E(\widehat{\mathcal F},\widehat{\mathcal G})
=\mathrm E(i^*(\mathcal F),i^*(\mathcal G))
\]
它在 $\mathrm E_2$ 项（相应地，在收敛项）上归结为同态
\[
\varphi_2^{pq}:\mathrm H^p\bigl(X,
\mathcal{E}\!xt^q_{\mathcal O_X}(\mathcal F,\mathcal G)\bigr)
\longrightarrow
\mathrm H^p\bigl(\widehat X,
\mathcal{E}\!xt^q_{\mathcal O_{\widehat X}}
(\widehat{\mathcal F},\widehat{\mathcal G})\bigr)
\]
\[
\text{（相应地，}\varphi^n:\operatorname{Ext}^n_{\mathcal O_X}
(X;\mathcal F,\mathcal G)\longrightarrow
\operatorname{Ext}^n_{\mathcal O_{\widehat X}}
(\widehat X;\widehat{\mathcal F},\widehat{\mathcal G})\text{）}
\]
该同态由 $u_q$（相应地，$u_0$）依函子性导出（0，12.3.4）。与
$\widehat A$ 作张量积后，$\varphi_r^{pq}$ 与 $\varphi^n$ 给出
$\widehat A$-模同态
\[
\psi_r^{pq}:\mathrm E_r^{pq}(\mathcal F,\mathcal G)
\otimes_A\widehat A\longrightarrow
\mathrm E_r^{pq}(\widehat{\mathcal F},\widehat{\mathcal G})
\]
\[
\psi^n:\operatorname{Ext}^n_{\mathcal O_X}(X;\mathcal F,\mathcal G)
\otimes_A\widehat A\longrightarrow
\operatorname{Ext}^n_{\mathcal O_{\widehat X}}
(\widehat X;\widehat{\mathcal F},\widehat{\mathcal G}).
\]
由于 $\widehat A$ 是平坦 $A$-模（0\textsubscript{I}，7.3.3），
$\widehat A$-模
$\mathrm E_r^{pq}(\mathcal F,\mathcal G)\otimes_A\widehat A$ 构成一个
双正则谱序列，其收敛项为
$\operatorname{Ext}^n_{\mathcal O_X}(X;\mathcal F,\mathcal G)\otimes_A\widehat A$；
而 $\psi_r^{pq}$ 与 $\psi^n$ 构成谱序列的一个态射。由于
$\psi_2^{pq}$ 都是同构，$\psi^n$ 也都是同构（0，11.1.5）。
\end{proof}

\begin{corollary}[4.5.2]
\label{III.4.5.2-fr}
在（\hyperref[III.4.5.1-fr]{4.5.1}）的假设下，再假设 $A$ 是 Noether
$\mathfrak J$-进制环。则对每个整数 $n\geq0$，
$\operatorname{Ext}^n_{\mathcal O_X}(X;\mathcal F,\mathcal G)$ 典范地等同于
$\operatorname{Ext}^n_{\mathcal O_{\widehat X}}
(\widehat X;\widehat{\mathcal F},\widehat{\mathcal G})$。
\end{corollary}

\begin{proof}
只须注意，$\operatorname{Ext}^n_{\mathcal O_X}
(X;\mathcal F,\mathcal G)$ 作为有限型 $A$-模，关于
$\mathfrak J$-预进制拓扑是分离且完备的（0\textsubscript{I}，7.3.6）。
\end{proof}

（\hyperref[III.4.5.1-fr]{4.5.1}）在 $n=0$ 时的特殊情形可表述如下：

\begin{corollary}[4.5.3]
\label{III.4.5.3-fr}
在（\hyperref[III.4.5.1-fr]{4.5.1}）的假设下，对每个同态
$u:\mathcal F\to\mathcal G$，以
$\widehat u:\widehat{\mathcal F}\to\widehat{\mathcal G}$ 表示完备化后的
同态（I，10.8.4）。则有典范同构
\[
\label{III.4.5.3.1-fr}
\left(\operatorname{Hom}_{\mathcal O_X}(\mathcal F,\mathcal G)\right)^\wedge
\xrightarrow{\sim}
\operatorname{Hom}_{\mathcal O_{\widehat X}}
(\widehat{\mathcal F},\widehat{\mathcal G}),
\tag{4.5.3.1}
\]
其中左端是 $A$-模
$\operatorname{Hom}_{\mathcal O_X}(\mathcal F,\mathcal G)$ 关于
$\mathfrak J$-预进制拓扑的分离完备化；这个同构由同态
$u\mapsto\widehat u$ 取分离完备化而得到。
\end{corollary}

\subsection{形式态射与通常态射之间的关系}

\begin{proposition}[4.6.1]
\label{III.4.6.1-fr}
设 $Y$ 是局部 Noether 预概形，$f:X\to Y$ 是固有态射，$\mathcal F$ 是
凝聚且 $f$-平坦的 $\mathcal O_X$-模，$y$ 是 $Y$ 的一点。假设对某个
整数 $n$ 有
\[
H^n\bigl(f^{-1}(y),\mathcal F\otimes_{\mathcal O_Y}k(y)\bigr)=0.
\]
则在 $Y$ 中存在 $y$ 的一个邻域 $U$，使得
$R^nf_*(\mathcal F)|U=0$；并且对每个整数 $p\geq0$，
（\hyperref[III.4.2.1.1-fr]{4.2.1.1}）的典范同态
\[
\bigl(R^{n-1}f_*(\mathcal F)\bigr)_y
\longrightarrow
H^{n-1}\bigl(f^{-1}(y),
\mathcal F\otimes_{\mathcal O_Y}
(\mathcal O_y/\mathfrak m_y^{p+1})\bigr)
\]
是满射。
\end{proposition}

\begin{proof}
由于 $R^nf_*(\mathcal F)$ 是凝聚 $\mathcal O_Y$-模
（\hyperref[III.3.2.1-fr]{3.2.1}），只要证明
$(R^nf_*(\mathcal F))_y=0$（0\textsubscript{I}，5.2.2），命题的第一个
断言便成立；由（\hyperref[III.4.2.1-fr]{4.2.1}），只须证明
\[
H^n\bigl(f^{-1}(y),\mathcal F\otimes_{\mathcal O_Y}
(\mathcal O_y/\mathfrak m_y^{p+1})\bigr)=0
\]
对每个 $p$ 成立。依假设，这对 $p=0$ 成立；下面对 $p$ 作归纳。置
\[
X_p=X\times_Y\operatorname{Spec}
(\mathcal O_y/\mathfrak m_y^{p+1}),
\]
于是 $X_{p-1}$ 是 $X_p$ 的闭子预概形，且二者有相同的底空间
（I，3.6.1）；归纳假设
\[
H^n\bigl(X_{p-1},\mathcal F\otimes_{\mathcal O_Y}
(\mathcal O_y/\mathfrak m_y^p)\bigr)=0
\]
因而蕴含
\[
H^n\bigl(X_p,\mathcal F\otimes_{\mathcal O_Y}
(\mathcal O_y/\mathfrak m_y^p)\bigr)=0~;
\]
另一方面，从 $\mathcal O_{X_p}$-模的正合序列
\[
0\longrightarrow
\mathfrak m_y^p\mathcal F/\mathfrak m_y^{p+1}\mathcal F
\longrightarrow
\mathcal F/\mathfrak m_y^{p+1}\mathcal F
\longrightarrow
\mathcal F/\mathfrak m_y^p\mathcal F
\longrightarrow0
\]
得到上同调正合序列
\[
H^n\bigl(X_p,
\mathfrak m_y^p\mathcal F/\mathfrak m_y^{p+1}\mathcal F\bigr)
\longrightarrow
H^n\bigl(X_p,\mathcal F/\mathfrak m_y^{p+1}\mathcal F\bigr)
\longrightarrow
H^n\bigl(X_p,\mathcal F/\mathfrak m_y^p\mathcal F\bigr)
\]
故只须证明
\[
\label{III.4.6.1.1-fr}
H^n\bigl(X_p,
\mathfrak m_y^p\mathcal F/\mathfrak m_y^{p+1}\mathcal F\bigr)=0.
\tag{4.6.1.1}
\]
因为一旦如此，$H^n(X_p,\mathcal F/\mathfrak m_y^{p+1}\mathcal F)$ 就是
$H^n(X_p,\mathcal F/\mathfrak m_y^p\mathcal F)$ 的子模，因而由归纳假设
等于 $0$。

\oldpage[III]{140}

现在注意，纤维
$Z=f^{-1}(y)=X\times_Y\operatorname{Spec}(k(y))$ 是 $X_p$ 的闭子预概形，
而 $\mathfrak m_y^p\mathcal F/\mathfrak m_y^{p+1}\mathcal F$ 被
$\mathfrak m_y$ 零化，因而可视为 $\mathcal O_Z$-模；所以
\[
H^n\bigl(Z,\mathfrak m_y^p\mathcal F/
\mathfrak m_y^{p+1}\mathcal F\bigr)
=H^n\bigl(X_p,\mathfrak m_y^p\mathcal F/
\mathfrak m_y^{p+1}\mathcal F\bigr).
\]
下面证明典范的 $\mathcal O_Z$-同态
\[
\label{III.4.6.1.2-fr}
(\mathcal F/\mathfrak m_y\mathcal F)
\otimes_{k(y)}(\mathfrak m_y^p/\mathfrak m_y^{p+1})
\longrightarrow
\mathfrak m_y^p\mathcal F/\mathfrak m_y^{p+1}\mathcal F
\tag{4.6.1.2}
\]
是双射。证明这一点后，由于
$\mathfrak m_y^p/\mathfrak m_y^{p+1}$ 是自由 $k(y)$-模，便有
\[
\begin{split}
H^n\bigl(Z,\mathfrak m_y^p\mathcal F/
\mathfrak m_y^{p+1}\mathcal F\bigr)
&=H^n\bigl(Z,\mathcal F/\mathfrak m_y\mathcal F\bigr)
\otimes_{k(y)}(\mathfrak m_y^p/\mathfrak m_y^{p+1})\\
&=0
\end{split}
\]
（0，12.2.3），因为依假设
$H^n(Z,\mathcal F/\mathfrak m_y\mathcal F)=0$，从而得到
（\hyperref[III.4.6.1.1-fr]{4.6.1.1}）。因此，为证明第一个断言，只剩下
证明（\hyperref[III.4.6.1.2-fr]{4.6.1.2}）是双射。这个问题在 $X$ 上是
逐点的，并且依假设，对每个 $x\in f^{-1}(y)$，$\mathcal F_x$ 都是
平坦 $\mathcal O_y$-模，所以只须应用（0，10.2.1，c）），因为
$\mathcal F_x/\mathfrak m_y\mathcal F_x$ 是域
$k(y)=\mathcal O_y/\mathfrak m_y$ 上的平坦模。

为证明（\hyperref[III.4.6.1-fr]{4.6.1}）的第二个断言，完全如
（\hyperref[III.4.2.1-fr]{4.2.1}）中那样，立即归结到 $Y$ 仿射且 $y$
为闭点的情形。注意，由与（\hyperref[III.4.6.1.1-fr]{4.6.1.1}）相似的
论证，对每个 $k>0$ 有
\[
\label{III.4.6.1.3-fr}
H^n\bigl(X_{k+1},
\mathfrak m_y^k\mathcal F/
\mathfrak m_y^{k+p+1}\mathcal F\bigr)=0.
\tag{4.6.1.3}
\]

\oldpage[III]{141}

由此再用（\hyperref[III.4.2.1-fr]{4.2.1}）推出
\[
\label{III.4.6.1.4-fr}
\bigl(R^nf_*(\mathfrak m_y^k\mathcal F)\bigr)_y=0.
\tag{4.6.1.4}
\]
于是从上同调正合序列得到序列
\[
\bigl(R^{n-1}f_*(\mathcal F)\bigr)_y
\longrightarrow
\bigl(R^{n-1}f_*(\mathcal F/\mathfrak m_y^p\mathcal F)\bigr)_y
\longrightarrow
\bigl(R^nf_*(\mathfrak m_y^p\mathcal F)\bigr)_y=0
\]
正合。又因 $y$ 是闭点且 $Y$ 仿射，由
（\hyperref[III.1.4.11-fr]{1.4.11}）有
\[
R^{n-1}f_*(\mathcal F/\mathfrak m_y^p\mathcal F)
=\bigl(H^{n-1}(X,\mathcal F/\mathfrak m_y^p\mathcal F)\bigr)^\sim
=\bigl(H^{n-1}(f^{-1}(y),
\mathcal F/\mathfrak m_y^p\mathcal F)\bigr)^\sim
\]
（G，II，4.9.1）；而
$H^{n-1}(f^{-1}(y),\mathcal F/\mathfrak m_y^p\mathcal F)$ 是
$(\mathcal O_y/\mathfrak m_y^p)$-模，故
\[
\bigl(R^{n-1}f_*(\mathcal F/\mathfrak m_y^p\mathcal F)\bigr)_y
=H^{n-1}(f^{-1}(y),\mathcal F/\mathfrak m_y^p\mathcal F)
\]
这就完成了（\hyperref[III.4.6.1-fr]{4.6.1}）的证明。
\end{proof}

\begin{corollary}[4.6.2]
\label{III.4.6.2-fr}
设 $Y$ 是局部 Noether 预概形，$f:X\to Y$ 是固有且平坦的态射，
$\mathcal F$、$\mathcal G$ 是两个局部自由 $\mathcal O_X$-模，$y$ 是
$Y$ 的一点。置
\[
X_y=f^{-1}(y)=X\otimes_Yk(y),\qquad
\mathcal F_y=\mathcal F\otimes_{\mathcal O_Y}k(y),\qquad
\mathcal G_y=\mathcal G\otimes_{\mathcal O_Y}k(y),
\]
并假设
\[
\label{III.4.6.2.1-fr}
H^1\bigl(X_y,
\mathcal{H}\!om_{\mathcal O_{X_y}}(\mathcal F_y,\mathcal G_y)\bigr)=0.
\tag{4.6.2.1}
\]
则对每个同态 $u_0:\mathcal F_y\to\mathcal G_y$，存在 $y$ 的一个开邻域
$U$ 以及一个同态
\[
u:\mathcal F|f^{-1}(U)\longrightarrow\mathcal G|f^{-1}(U)
\]
使得 $u_0$ 等于同态 $u\otimes1$。
\end{corollary}

\begin{proof}
事实上，该假设允许我们对凝聚 $\mathcal O_X$-模
$\mathcal H=\mathcal{H}\!om_{\mathcal O_X}(\mathcal F,\mathcal G)$ 应用
（\hyperref[III.4.6.1-fr]{4.6.1}），其中取 $n=1$、$p=0$；这是因为
$\mathcal H$ 局部自由，因而 \emph{a fortiori} 是 $f$-平坦的，并且
$\mathcal O_{X_y}$-模 $\mathcal H\otimes_{\mathcal O_Y}k(y)$ 此时等同于
$\mathcal{H}\!om_{\mathcal O_{X_y}}(\mathcal F_y,\mathcal G_y)$
（0\textsubscript{I}，6.2.2）。可以假设 $Y=\operatorname{Spec}(A)$
是仿射的；这时由（\hyperref[III.1.4.11-fr]{1.4.11}）有
\[
R^0f_*(\mathcal H)=
\bigl(\operatorname{Hom}_{\mathcal O_X}(\mathcal F,\mathcal G)\bigr)^\sim,
\]
因而
\[
\bigl(R^0f_*(\mathcal H)\bigr)_y
=\operatorname{Hom}_{\mathcal O_X}(\mathcal F,\mathcal G)
\otimes_A\mathcal O_y~;
\]
由（\hyperref[III.4.6.1-fr]{4.6.1}），典范同态
\[
\operatorname{Hom}_{\mathcal O_X}(\mathcal F,\mathcal G)
\otimes_A\mathcal O_y
\longrightarrow
\operatorname{Hom}_{\mathcal O_{X_y}}(\mathcal F_y,\mathcal G_y)
\]
是满射，这就证明了该推论，因为
$\operatorname{Hom}_{\mathcal O_X}(\mathcal F,\mathcal G)
\otimes_A\mathcal O_y$ 的每个元素总可写成 $u\otimes(1/s)$ 的形式，其中
$s\notin\mathfrak m_y$ 是 $A$ 的一个元素。
\end{proof}

这个推论可由下一个推论加以补充：

\begin{corollary}[4.6.3]
\label{III.4.6.3-fr}
在（\hyperref[III.4.6.2-fr]{4.6.2}）的假设下，若 $u_0$ 是单射
（相应地，满射、双射），则可假设 $u$ 也具有同一性质。
\end{corollary}

\begin{proof}
可以只考虑 $U=Y$ 的情形。只须证明：若 $u_0$ 是单射（相应地，满射），
则对每个 $x\in f^{-1}(y)$ 有 $\operatorname{Ker}u_x=0$（相应地，
$\operatorname{Coker}u_x=0$）。事实上，$\operatorname{Ker}u$ 与
$\operatorname{Coker}u$ 都是凝聚 $\mathcal O_X$-模
（0\textsubscript{I}，5.3.4），故在 $X$ 中存在 $f^{-1}(y)$ 的一个邻域
$V$，使得 $\operatorname{Ker}u$（相应地，$\operatorname{Coker}u$）在
$V$ 上的限制为 $0$（0\textsubscript{I}，5.2.2）。由于 $f$ 是闭映射，
存在 $y$ 的一个邻域 $U'\subset U$，使 $f^{-1}(U')\subset V$，从而
（\hyperref[III.4.6.3-fr]{4.6.3}）得证。依假设，
\[
u_x\otimes1:
\mathcal F_x\otimes_{\mathcal O_y}k(y)
\longrightarrow
\mathcal G_x\otimes_{\mathcal O_y}k(y)
\]
是单射（相应地，满射）；$\mathcal F_x$ 与 $\mathcal G_x$ 是有限秩自由
$\mathcal O_x$-模，并且 $\mathcal O_x$ 是平坦 $\mathcal O_y$-模。当假设
$u_x\otimes1$ 为单射时，$u_x$ 为单射由（0，10.2.4）得出。当假设
$u_x\otimes1$ 为满射时，\emph{a fortiori}，由它取商所得的同态
\[
\mathcal F_x/\mathfrak m_x\mathcal F_x
\longrightarrow
\mathcal G_x/\mathfrak m_x\mathcal G_x,
\]
是满射。由于 $\mathcal G_x$ 是有限型 $\mathcal O_x$-模，而
$\mathcal O_x$ 是以 $\mathfrak m_x$ 为极大理想的局部环，结论由 Nakayama
引理得出（Bourbaki，\emph{Alg.}，第 VIII 章，\S\,6，第 3 号，命题 6
的推论 4）。
\end{proof}

特别地，由（\hyperref[III.4.6.3-fr]{4.6.3}）得到：

\begin{corollary}[4.6.4]
\label{III.4.6.4-fr}
设 $Y$ 是局部 Noether 预概形，$f:X\to Y$ 是固有且平坦的态射，$y$ 是
$Y$ 的一点，$X_y=X\otimes_Yk(y)$。设 $\mathcal E_0$ 是局部自由
$\mathcal O_{X_y}$-模，并且
\[
\label{III.4.6.4.1-fr}
H^1\bigl(X_y,
\mathcal{H}\!om_{\mathcal O_{X_y}}(\mathcal E_0,\mathcal E_0)\bigr)=0.
\tag{4.6.4.1}
\]
设 $\mathcal F$、$\mathcal G$ 是两个局部自由 $\mathcal O_X$-模，并且
$\mathcal F_y$ 与 $\mathcal G_y$（采用
（\hyperref[III.4.6.2-fr]{4.6.2}）的记号）都同构于 $\mathcal E_0$。
则存在 $y$ 的一个开邻域 $U$，使得 $\mathcal F|f^{-1}(U)$ 与
$\mathcal G|f^{-1}(U)$ 同构。
\end{corollary}

更特别地：

\begin{corollary}[4.6.5]
\label{III.4.6.5-fr}
在（\hyperref[III.4.6.4-fr]{4.6.4}）对 $f$、$X$、$Y$ 的假设下，假设
$H^1(X_y,\mathcal O_{X_y})=0$。若 $\mathcal F$ 与 $\mathcal G$ 是两个
可逆 $\mathcal O_X$-模，且 $\mathcal F_y$ 与 $\mathcal G_y$ 同构，则存在
$y$ 的一个开邻域 $U$，使 $\mathcal F|f^{-1}(U)$ 与
$\mathcal G|f^{-1}(U)$ 同构。
\end{corollary}

\begin{proof}
只须把（\hyperref[III.4.6.4-fr]{4.6.4}）应用于模
$\mathcal F^{-1}\otimes\mathcal G$ 与 $\mathcal O_X$。
\end{proof}

\begin{remarks}[4.6.6]
\label{III.4.6.6-fr}
\begin{enumerate}
\item[(i)] 利用（\hyperref[III.4.6.5-fr]{4.6.5}），我们将在第五章中
建立（II，4.2.7）所预告的射影丛上可逆层的分类。
\item[(ii)] （\hyperref[III.4.6.1-fr]{4.6.1}）的结果将在第 7 节中作为
更一般命题的推论出现。
\end{enumerate}
\end{remarks}

\begin{proposition}[4.6.7]
\label{III.4.6.7-fr}
设 $Z$ 是局部 Noether 预概形，$X$、$Y$ 是两个 $Z$-预概形，并且结构
态射 $g:X\to Z$、$h:Y\to Z$ 都固有。设 $f:X\to Y$ 是 $Z$-态射，
$z$ 是 $Z$ 的一点，并置
\[
f_z=f\times_Z1:
X\otimes_Zk(z)\longrightarrow Y\otimes_Zk(z).
\]
\begin{enumerate}
\item[(i)] 若 $f_z$ 是有限态射（相应地，闭浸入），则存在 $z$ 的一个
开邻域 $U$，使得 $f$ 的限制 $g^{-1}(U)\to h^{-1}(U)$ 是有限态射
（相应地，闭浸入）。
\item[(ii)] 再假设 $g$ 是平坦态射。若 $f_z$ 是同构，则存在 $z$ 的一个
开邻域 $U$，使得 $f$ 的限制 $g^{-1}(U)\to h^{-1}(U)$ 是同构。
\end{enumerate}
\end{proposition}

\begin{proof}
在两种情形中，只须证明：对每个 $y\in h^{-1}(z)$，都存在 $y$ 的一个
邻域 $V_y$，使得 $f$ 的限制 $f^{-1}(V_y)\to V_y$ 是有限态射
（相应地，闭浸入、同构）。事实上，若 $V$ 是诸 $V_y$ 的并，则 $f$
的限制 $f^{-1}(V)\to V$ 是有限态射（相应地，闭浸入、同构）
（II，6.1.1 及 I，4.2.4）。由于 $h$ 是闭映射，存在 $z$ 的一个开邻域
$U$，使得 $h^{-1}(U)\subset V$，命题遂得证。

\noindent\textit{(i)} 首先注意，$f$ 是固有态射（II，5.4.3）。若假设
$f_z$ 有限，则对每个 $y\in h^{-1}(z)$，存在邻域 $V_y$ 使
$f^{-1}(V_y)\to V_y$ 有限，这由
（\hyperref[III.4.4.11-fr]{4.4.11}）得出。

\oldpage[III]{143}

为处理 $f_z$ 是闭浸入的情形，现在可以先假设态射 $f$ 有限，因而
$X=\operatorname{Spec}(\mathcal B)$，其中 $\mathcal B$ 是凝聚
$\mathcal O_Y$-代数；态射 $f$ 对应于典范同态
$u:\mathcal O_Y\to\mathcal B$（II，1.2.7）。若证明对每个
$y\in h^{-1}(z)$，同态 $u_y:\mathcal O_y\to\mathcal B_y$ 都是满射，
则在 $y$ 的某个邻域 $V_y$ 上，$u|V_y$ 是满射，因为层
$\operatorname{Coker}u$ 凝聚（0\textsubscript{I}，5.3.4 及 5.2.2）。
另一方面，有限态射 $f_z$ 对应于同态
\[
v=u\otimes1:
\mathcal O_Y\otimes_{\mathcal O_Z}k(z)
\longrightarrow
\mathcal B\otimes_{\mathcal O_Z}k(z),
\]
而 $f_z$ 是闭浸入这一假设蕴含同态
\[
u_y\otimes1:
\mathcal O_y\otimes_{\mathcal O_z}k(z)
\longrightarrow
\mathcal B_y\otimes_{\mathcal O_z}k(z)
\]
为满射。由于 $\mathcal B_y$ 是有限型 $\mathcal O_y$-模，而
$\mathcal O_y$ 是 Noether 局部环，结论如
（\hyperref[III.4.6.3-fr]{4.6.3}）中一样由 Nakayama 引理得出。

\noindent\textit{(ii)} 与上面相同的论证表明，这次只须由
$u_y\otimes1$ 是双射证明 $u_y$ 是双射。

这由下面的引理得出：
\end{proof}

\begin{lemma}[4.6.7.1]
\label{III.4.6.7.1-fr}
设 $A$、$B$ 是两个 Noether 局部环，$\rho:A\to B$ 是局部同态，
$u:N\to M$ 是 $B$-模同态。假设 $M$ 是平坦 $A$-模，$N$ 是有限型
$B$-模，并且
\[
u\otimes1:N\otimes_Ak\longrightarrow M\otimes_Ak
\]
是单射，其中 $k$ 是 $A$ 的剩余域。则 $N$ 是平坦 $A$-模，且 $u$
是单射。
\end{lemma}

\begin{proof}
为证明第一个断言，必须证明：对任意一对有限型 $A$-模 $P$、$Q$，以及
任意单射 $A$-同态 $v:P\to Q$，同态
$1_N\otimes v:N\otimes_AP\to N\otimes_AQ$ 都是单射。现有交换图
\[
\xymatrix{
N\otimes_AP\ar[r]^{1_N\otimes v}\ar[d]_{u\otimes1_P}
  &N\otimes_AQ\ar[d]^{u\otimes1_Q}\\
M\otimes_AP\ar[r]_{1_M\otimes v}
  &M\otimes_AQ
}
\]
依假设，$1_M\otimes v$ 是单射，故只须证明 $u\otimes1_P$ 也是单射。
设 $\mathfrak m$ 是 $A$ 的极大理想；$A$-模 $N\otimes_AP$ 上的
$\mathfrak m$-进滤过，也是它作为 $B$-模的 $\mathfrak mB$-进滤过。
这个滤过定义的拓扑是分离的，因为 $B$ 是 Noether 环，$\mathfrak mB$
包含于 $B$ 的 Jacobson 根中，并且 $N\otimes_AP$ 是有限型 $B$-模；
这里使用了 $N$ 是有限型 $B$-模且 $P$ 是有限型 $A$-模
（0\textsubscript{I}，7.3.5）。因此，只须证明同态
\[
\operatorname{gr}(u\otimes1_P):
\operatorname{gr}_\bullet(N\otimes_AP)
\longrightarrow
\operatorname{gr}_\bullet(M\otimes_AP)
\]
是单射，其中分次模是相对于 $\mathfrak m$-进滤过而言的（Bourbaki，
\emph{Alg. comm.}，第 III 章，\S\,2，第 8 号，定理 1 的推论 1）。
现在注意，由于 $M$ 是平坦 $A$-模，同态
\[
M\otimes_A(\mathfrak m^nP)
\longrightarrow
\mathfrak m^n(M\otimes_AP)
\]
都是双射；因此典范同态
\[
\varphi_M:
\operatorname{gr}_0(M)\otimes_A\operatorname{gr}_\bullet(P)
\longrightarrow
\operatorname{gr}_\bullet(M\otimes_AP).
\]
也是双射。

\oldpage[III]{144}

另一方面，有交换图
\[
\xymatrix{
\operatorname{gr}_0(N)\otimes_A\operatorname{gr}_\bullet(P)
  \ar[r]^{\operatorname{gr}_0(u)\otimes1}\ar[d]_{\varphi_N}
  &\operatorname{gr}_0(M)\otimes_A\operatorname{gr}_\bullet(P)
  \ar[d]^{\varphi_M}\\
\operatorname{gr}_\bullet(N\otimes_AP)
  \ar[r]_{\operatorname{gr}(u\otimes1)}
  &\operatorname{gr}_\bullet(M\otimes_AP)
}
\]
其中 $\varphi_M$ 是双射，$\varphi_N$ 是满射；此外，依假设
$\operatorname{gr}_0(u)$ 是单射，并且由于
\[
\begin{split}
\operatorname{gr}_0(N)\otimes_A\operatorname{gr}_\bullet(P)
&=\operatorname{gr}_0(N)\otimes_k\operatorname{gr}_\bullet(P),\\
\operatorname{gr}_0(M)\otimes_A\operatorname{gr}_\bullet(P)
&=\operatorname{gr}_0(M)\otimes_k\operatorname{gr}_\bullet(P),
\end{split}
\]
$\operatorname{gr}_0(u)\otimes1$ 也是单射。由此得出
$\operatorname{gr}(u\otimes1)$ 是单射，从而第一个断言得证。第二个断言
由前面的论证取 $P=A$ 得出。
\end{proof}

\begin{proposition}[4.6.8]
\label{III.4.6.8-fr}
设 $Z$ 是局部 Noether 预概形，$X$、$Y$ 是两个 $Z$-预概形，并且结构
态射 $g:X\to Z$、$h:Y\to Z$ 都固有；设 $Z'$ 是 $Z$ 的一个闭子集，
$X'=g^{-1}(Z')$、$Y'=h^{-1}(Z')$ 是它的逆像，
\[
\widehat X=X_{/X'},\qquad
\widehat Y=Y_{/Y'},\qquad
\widehat Z=Z_{/Z'}
\]
是 $X$、$Y$、$Z$ 沿这些闭子集的形式完备化；设 $f:X\to Y$ 是
$Z$-态射，$\widehat f:\widehat X\to\widehat Y$ 是它在完备化上的延拓。
$\widehat f$ 是同构（相应地，闭浸入）的充分必要条件是：存在 $Z'$ 的
一个开邻域 $U$，使得 $f$ 的限制 $g^{-1}(U)\to h^{-1}(U)$ 是同构
（相应地，闭浸入）。
\end{proposition}

\begin{proof}
该条件的充分性是立即的（I，10.14.7）。为证明必要性，仍只须对每个
$y\in Y'$ 证明存在 $y$ 的一个开邻域 $V_y$，使得 $f$ 的限制
$f^{-1}(V_y)\to V_y$ 是同构（相应地，闭浸入）；论证与
（\hyperref[III.4.6.7-fr]{4.6.7}）中相同。由此归结到 $Y=Z$ 且
$Y=\operatorname{Spec}(A)$ 为 Noether 仿射概形的情形。依假设
（I，10.9.1 及 10.14.2），对 $y\in Y'$，纤维 $f^{-1}(y)$ 约化为一点；
又因 $f$ 固有（II，5.4.3），存在 $y$ 的一个开邻域 $U$，使得 $f$ 的
限制 $f^{-1}(U)\to U$ 是有限态射
（\hyperref[III.4.4.11-fr]{4.4.11}）。因此可以先假设 $f$ 是有限态射，
从而 $X=\operatorname{Spec}(B)$，其中 $B$ 是 $A$-代数，且在 $A$ 上有限。若
$Y'=V(\mathfrak J)$，则
\[
\widehat Y=\operatorname{Spf}(\widehat A),\qquad
\widehat X=\operatorname{Spf}(\widehat B),
\]
其中 $\widehat A$ 是 $A$ 关于 $\mathfrak J$-预进制拓扑的分离完备化，
$\widehat B$ 是 $B$ 关于 $\mathfrak JB$-预进制拓扑的分离完备化；等价地，
它是 $A$-模 $B$ 关于 $\mathfrak J$-预进制拓扑的分离完备化。此外，
$\widehat f$ 是与环的典范同态 $\varphi:A\to B$ 的连续延拓
$\widehat\varphi:\widehat A\to\widehat B$ 对应的仿射形式概形态射；
假设即 $\widehat\varphi$ 为满射（相应地，双射）（I，10.14.2）。而
$\widehat\varphi$ 也是把 $\varphi$ 视为 $A$-模同态时的连续延拓；于是
由（I，10.8.14），存在 $Y'$ 的一个开邻域 $U$，使
$\mathcal O_Y$-模同态 $\widetilde\varphi:\widetilde A\to\widetilde B$ 在
$U$ 上的限制为满射（相应地，双射），证明完毕。
\end{proof}

\oldpage[III]{145}

\subsection{一个丰沛性判据}

\begin{theorem}[4.7.1]
\label{III.4.7.1-fr}
设 $Y$ 是局部 Noether 预概形，$f:X\to Y$ 是固有态射，$\mathcal L$ 是
可逆 $\mathcal O_X$-模，$y$ 是 $Y$ 的一点，
$X_y=X\otimes_Yk(y)=f^{-1}(y)$，$g$ 是从 $X_y$ 到 $X$ 的投影。若
$\mathcal L_y=g^*(\mathcal L)=\mathcal L\otimes_{\mathcal O_Y}k(y)$ 在
$X_y$ 上丰沛，则在 $Y$ 中存在 $y$ 的一个开邻域 $U$，使得
$\mathcal L|f^{-1}(U)$ 对 $f$ 在 $f^{-1}(U)$ 上的限制丰沛。
\end{theorem}

\begin{proof}
\noindent\textit{I)} 置
\[
Y'=\operatorname{Spec}(\mathcal O_y),\qquad
X'=X\times_YY',\qquad
\mathcal L'=\mathcal L\otimes_{\mathcal O_Y}\mathcal O_y~;
\]
我们先证明 $\mathcal L'$ 对 $f'=f_{(Y')}$ 丰沛。现有交换图
\[
\xymatrix{
X\ar[d]_f & X'\ar[l]\ar[d]_{f'} & X_y\ar[l]\ar[d]_{f_y}\\
Y & Y'\ar[l] & \operatorname{Spec}(k(y))\ar[l]
}
\]
由于 $f'$ 固有（II，5.4.2，(iii)），且 $\mathcal O_y$ 是 Noether 环，
可只考虑 $Y=Y'=\operatorname{Spec}(\mathcal O_y)$，因而 $X=X'$ 的情形；
假设 $\mathcal L_y$ 对 $f_y$ 丰沛，并证明 $\mathcal L$ 对 $f$ 丰沛
（II，4.6.6）。我们将应用判据
（\hyperref[III.2.6.1-fr]{2.6.1, c)}），并且事实上证明：对每个凝聚
$\mathcal O_X$-模 $\mathcal F$，存在整数 $N$，使得
\[
H^1(X,\mathcal F(n))=0\quad\text{对所有 }n\geq N,
\qquad
\mathcal F(n)=\mathcal F\otimes\mathcal L^{\otimes n}.
\]
注意，$y$ 是 $Y$ 的闭点，对应于 $\mathcal O_y$ 的极大理想
$\mathfrak m$；因此 $X_y$ 是由 $\mathcal O_X$ 的凝聚理想
$\mathcal J=f^*(\widetilde{\mathfrak m})\mathcal O_X=
\mathfrak m\mathcal O_X$ 定义的 $X$ 的闭子预概形（I，4.4.5），而 $g$ 是
典范浸入。考虑分次 $k(y)$-代数
\[
S=\operatorname{gr}(\mathcal O_y)
=\bigoplus_{j\geq0}\mathfrak m^j/\mathfrak m^{j+1},
\]
它是有限型的，因为 $\mathcal O_y$ 是 Noether 环。于是
$\mathcal O_X$-代数 $\mathcal S=f^*(\widetilde S)$ 是拟凝聚且有限型的；
它显然被 $\mathcal J$ 零化，故若置 $\mathcal S_y=g^*(\mathcal S)$，
则 $\mathcal S_y$ 是有限型拟凝聚 $\mathcal O_{X_y}$-代数，并且
$\mathcal S=g_*(\mathcal S_y)$。另一方面，置
\[
\mathcal M_j=\mathfrak m^j\mathcal F/\mathfrak m^{j+1}\mathcal F,
\qquad
\mathcal M=\bigoplus_{j\geq0}\mathcal M_j
=\operatorname{gr}(\mathcal F)~;
\]
因为 $\mathcal F$ 凝聚，$\mathcal M$ 是有限型拟凝聚 $\mathcal S$-模
（0，10.1.1），它也被 $\mathcal J$ 零化。因此若置
\[
\mathcal M'_j=g^*(\mathcal M_j),\qquad
\mathcal M'=g^*(\mathcal M)=\bigoplus_{j\geq0}\mathcal M'_j,
\]
则 $\mathcal M'$ 是有限型拟凝聚分次 $\mathcal S_y$-模，且
$\mathcal M=g_*(\mathcal M')$。此外，若置
$\mathcal M'_j(n)=\mathcal M'_j\otimes\mathcal L_y^{\otimes n}$，则
$\mathcal M'_j(n)=g^*(\mathcal M_j(n))$。由于 $f_y$ 固有
（II，5.4.2，(iii)），且 $\mathcal L_y$ 丰沛，所以 $f_y$ 是射影态射
（II，5.5.4 及 4.6.11）。可以对
$\operatorname{Spec}(k(y))$、$f_y$、$\mathcal S_y$、$\mathcal L_y$ 和
$\mathcal M'$ 应用定理（\hyperref[III.2.4.1-fr]{2.4.1, (ii)}）：存在
整数 $N$，使得当 $n\geq N$ 时，对每个 $q>0$ 和每个 $j$ 都有
$H^q(X_y,\mathcal M'_j(n))=0$；因而对每个 $q>0$ 和每个 $j$ 也有
$H^q(X,\mathcal M_j(n))=0$（G，II，4.9.1）。现在置
\[
\mathcal F(n)_j=\mathcal F(n)/\mathfrak m^{j+1}\mathcal F(n),
\]
于是当 $j\geq1$ 时，
$\mathcal F(n)_{j-1}=\mathcal F(n)_j/\mathcal M_j(n)$，并且
$\mathcal F(n)_0=\mathcal M_0(n)$。现有
$H^1(X,\mathcal F(n)_0)=0$；由上同调正合序列，对每个 $j\geq1$ 有
$H^1(X,\mathcal F(n)_j)=H^1(X,\mathcal F(n)_{j-1})$，故对每个
$j\geq0$ 都有 $H^1(X,\mathcal F(n)_j)=0$。于是由
（\hyperref[III.4.2.1-fr]{4.2.1}）得出 $H^1(X,\mathcal F(n))=0$，从而
完成所需断言的证明。

\oldpage[III]{146}

\noindent\textit{II)} 回到证明开头的记号，并注意总可假设
$Y=\operatorname{Spec}(A)$ 为仿射概形。由于 $f'$ 是有限型的，且
$\mathcal L'$ 对 $f'$ 丰沛，存在整数 $m>0$，使
$\mathcal L'^{\otimes m}$ 对 $f'$ 极丰沛（II，4.6.11）。必要时以
$\mathcal L^{\otimes m}$ 代替 $\mathcal L$，便可只考虑 $\mathcal L'$ 对
$f'$ 极丰沛的情形，并证明 $\mathcal L|f^{-1}(U)$ 对 $f$ 极丰沛。由于
$f'$ 固有，存在一个适当的整数 $r>0$ 和 $Y'$-闭浸入
$j:X'\to P=\mathbf P_{Y'}^r$，使 $\mathcal L'$ 同构于
$j^*(\mathcal O_P(1))$（II，5.5.4，(ii)）。这个浸入典范地对应于满的
$\mathcal O_{X'}$-同态
\[
u:\mathcal O_{X'}^{r+1}\longrightarrow\mathcal L'
\]
（II，4.2.3）。后者对应于 $X'$ 上 $\mathcal L'$ 的 $r+1$ 个截面
$s'_i$（$0\leq i\leq r$）（0\textsubscript{I}，5.1.1），它们生成这个
$\mathcal O_{X'}$-模。依定义，这些截面也是 $Y'$ 上
$f'_*(\mathcal L')$ 的截面；有
\[
f'_*(\mathcal L')=f_*(\mathcal L)\otimes_{\mathcal O_Y}\mathcal O_{Y'}
\]
（0\textsubscript{I}，5.4.10）。又因 $Y$ 仿射，且 $\mathcal O_y$ 是
$A$ 在素理想 $\mathfrak j_y$ 处的局部环，故
$s'_i=s''_i/t_i$，其中 $s''_i$ 是 $Y$ 上 $f_*(\mathcal L)$ 的截面，
$t_i$ 是不属于 $\mathfrak j_y$ 的 $A$ 中元素。由此可知，在 $Y$ 中
存在 $y$ 的一个仿射开邻域 $V$，以及 $f_*(\mathcal L)|V$ 的截面
$s_i$，使 $s'_i=s_i/1$（回忆空间 $Y'$ 包含于 $V$ 中，参见 I，2.4.2）。
于是 $s_i$ 是 $f^{-1}(V)$ 上 $\mathcal L$ 的截面，从而定义同态
\[
v:\bigl(\mathcal O_X|f^{-1}(V)\bigr)^{r+1}
\longrightarrow\mathcal L|f^{-1}(V)
\]
它依假设在 $f^{-1}(y)$ 的所有点处都是满射。由于
$\operatorname{Coker}(v)$ 凝聚（0\textsubscript{I}，5.3.4），其支集闭
（0\textsubscript{I}，5.2.2）；因此存在 $f^{-1}(y)$ 的一个开邻域
$W\subset f^{-1}(V)$，使 $v$ 在 $W$ 上的限制为满射。因为
态射 $f$ 是闭映射，可以假设 $W$ 形如 $f^{-1}(U)$，其中 $U$ 是
$y$ 的一个开邻域；于是结论由（II，4.2.3）得出。
\end{proof}

\subsection{形式预概形的有限态射}

\begin{proposition}[4.8.1]
\label{III.4.8.1-fr}
设 $\mathfrak Y$ 是局部 Noether 形式预概形，$\mathcal K$ 是
$\mathfrak Y$ 的定义理想，$f:\mathfrak X\to\mathfrak Y$ 是形式预概形的
态射。下列条件等价：
\begin{enumerate}
\item[a)] $\mathfrak X$ 局部 Noether，$f$ 是进制态射
  （I，10.12.1），并且若置
  $\mathcal J=f^*(\mathcal K)\mathcal O_{\mathfrak X}$，则由 $f$ 导出的态射
  \[
  f_0:(\mathfrak X,\mathcal O_{\mathfrak X}/\mathcal J)
  \longrightarrow
  (\mathfrak Y,\mathcal O_{\mathfrak Y}/\mathcal K)
  \]
  是有限态射。
\item[b)] $\mathfrak X$ 局部 Noether，并且是一个进制
  $(Y_n)$-归纳系统 $(X_n)$ 的归纳极限，其中态射 $X_0\to Y_0$ 有限。
\item[c)] $\mathfrak Y$ 的每一点都有一个 Noether 仿射形式开邻域 $V$，
  使得 $f^{-1}(V)$ 是仿射形式开集，并且
  $\Gamma(f^{-1}(V),\mathcal O_{\mathfrak X})$ 是有限型
  $\Gamma(V,\mathcal O_{\mathfrak Y})$-模。
\end{enumerate}
\end{proposition}

\begin{proof}
由（I，10.12.3），a) 推出 b) 是立即的。为证明 b) 推出 c)，可以假设
$\mathfrak Y=\operatorname{Spf}(B)$，其中 $B$ 是 Noether 进制环，并且
$\mathcal K=\mathfrak R^\Delta$，其中 $\mathfrak R$ 是 $B$ 的一个定义理想。
依假设，$X_0$ 是仿射概形，其环 $A_0$ 是有限型 $B/\mathfrak R$-模
（II，6.1.3）。由（I，5.1.9），每个 $X_n$ 都是仿射概形；若 $A_n$ 是
它的环，则假设 b) 蕴含：当 $m\leq n$ 时，$A_m$ 同构于
$A_n/\mathfrak R^{m+1}A_n$。由此得出 $\mathfrak X$ 同构于
$\operatorname{Spf}(A)$，其中 $A=\varprojlim A_n$；结论遂由
（0\textsubscript{I}，7.2.9）得出。最后，为证明 c)

\oldpage[III]{147}

推出 a)，仍可只考虑
$\mathfrak Y=\operatorname{Spf}(B)$、$\mathfrak X=\operatorname{Spf}(A)$
的情形，其中 $A$ 是有限 $B$-代数；这时 $A/\mathfrak RA$ 是有限
$B/\mathfrak R$-代数，因而由（I，10.10.9）可知 a) 中的条件成立。
\end{proof}

\begin{definition}[4.8.2]
\label{III.4.8.2-fr}
当（\hyperref[III.4.8.1-fr]{4.8.1}）中的等价性质 a)、b)、c) 成立时，
称态射 $f$ 为\emph{有限的}；也称 $\mathfrak X$ 为\emph{有限形式
$\mathfrak Y$-预概形}，或称其为\emph{在 $\mathfrak Y$ 上有限的}形式预概形。
\end{definition}

\begin{proposition}[4.8.3]
\label{III.4.8.3-fr}
\begin{enumerate}
\item[(i)] 局部 Noether 形式预概形的闭浸入是有限态射。
\item[(ii)] 局部 Noether 形式预概形的两个有限态射之复合仍是有限态射。
\item[(iii)] 设 $\mathfrak X$、$\mathfrak Y$、$\mathfrak S$ 是三个局部
  Noether 形式预概形，$f:\mathfrak X\to\mathfrak S$ 是有限态射，
  $g:\mathfrak Y\to\mathfrak S$ 是态射；则态射
  $\mathfrak X\times_{\mathfrak S}\mathfrak Y\to\mathfrak Y$ 是有限的。
\item[(iv)] 设 $\mathfrak S$ 是局部 Noether 形式预概形，$\mathfrak X'$、
  $\mathfrak Y'$ 是两个局部 Noether 形式预概形，使得
  $\mathfrak X'\times_{\mathfrak S}\mathfrak Y'$ 局部 Noether。若
  $\mathfrak X$、$\mathfrak Y$ 是局部 Noether 形式
  $\mathfrak S$-预概形，$f:\mathfrak X\to\mathfrak X'$、
  $g:\mathfrak Y\to\mathfrak Y'$ 是两个有限 $\mathfrak S$-态射，则
  $f\times_{\mathfrak S}g$ 是有限态射。
\item[(v)] 设 $f:\mathfrak X\to\mathfrak Y$、
  $g:\mathfrak Y\to\mathfrak Z$ 是局部 Noether 形式预概形的两个态射，
  并且 $g$ 是有限型分离态射；若 $g\circ f$ 是有限态射，则 $f$ 是有限态射。
\end{enumerate}
\end{proposition}

\begin{proof}
(i) 是平凡的；借助（\hyperref[III.4.8.1-fr]{4.8.1}）中的判据 a)，
其余断言立即归结为通常预概形态射的相应命题（II，6.1.5）。仿照
（I，10.13.5），细节留给读者。
\end{proof}

\begin{corollary}[4.8.4]
\label{III.4.8.4-fr}
在（I，10.9.9）的假设下，若 $f$ 是有限态射，则其在完备化上的延拓
$\widehat f$ 也是有限态射。
\end{corollary}

\begin{corollary}[4.8.5]
\label{III.4.8.5-fr}
若 $\mathfrak X$ 是在 $\mathfrak Y$ 上有限的形式预概形，
$f:\mathfrak X\to\mathfrak Y$ 是结构态射，则对每个开集
$U\subset\mathfrak Y$，$f^{-1}(U)$ 在 $U$ 上有限。
\end{corollary}

\begin{proposition}[4.8.6]
\label{III.4.8.6-fr}
若 $f:\mathfrak X\to\mathfrak Y$ 是局部 Noether 形式预概形的有限态射，
则 $f_*(\mathcal O_{\mathfrak X})$ 是凝聚
$\mathcal O_{\mathfrak Y}$-代数。
\end{proposition}

\begin{proof}
可将 $f$ 视为态射 $f_n:X_n\to Y_n$ 的一个归纳系统 $(f_n)$ 的归纳极限；
我们证明诸 $f_n$ 是有限态射，且 $f_*(\mathcal O_{\mathfrak X})$ 同构于
诸 $(f_n)_*(\mathcal O_{X_n})$ 的逆极限，这便由（I，10.10.5）证明断言。
只须限于 $\mathfrak Y=\operatorname{Spf}(B)$、
$\mathfrak X=\operatorname{Spf}(A)$ 的情形，并注意：若 $\mathfrak R$ 是
$B$ 的定义理想且 $A$ 是有限型 $B$-模，则
$A/\mathfrak R^{n+1}A$ 是有限型 $B/\mathfrak R^{n+1}B$-模，而 $A$ 是
诸 $A/\mathfrak R^{n+1}A$ 的逆极限。
\end{proof}

反过来：

\begin{proposition}[4.8.7]
\label{III.4.8.7-fr}
设 $\mathfrak Y$ 是局部 Noether 形式预概形，$\mathcal A$ 是凝聚
$\mathcal O_{\mathfrak Y}$-代数。存在一个在 $\mathfrak Y$ 上有限的形式
预概形 $\mathfrak X$，它在唯一 $\mathfrak Y$-同构意义下确定，并且满足
$f_*(\mathcal O_{\mathfrak X})=\mathcal A$，其中
$f:\mathfrak X\to\mathfrak Y$ 是结构态射。
\end{proposition}

\begin{proof}
设 $\mathcal K$ 是 $\mathfrak Y$ 的定义理想，并置
\[
Y_n=(\mathfrak Y,\mathcal O_{\mathfrak Y}/\mathcal K^{n+1}),\qquad
\mathcal A_n=\mathcal A/\mathcal K^{n+1}\mathcal A~;
\]
显然，$\mathcal A_n$ 是有限 $\mathcal O_{Y_n}$-代数，因而定义一个

\oldpage[III]{148}

有限 $Y_n$-预概形 $X_n=\operatorname{Spec}(\mathcal A_n)$（II，6.1.3）；
当 $m\leq n$ 时，典范满同态 $h_{mn}:\mathcal A_n\to\mathcal A_m$
定义态射 $u_{mn}:X_m\to X_n$，使得图
\[
\xymatrix{
X_n\ar[d]_{f_n} & X_m\ar[l]_{u_{mn}}\ar[d]^{f_m}\\
Y_n & Y_m\ar[l]
}
\]
交换（其中 $f_n$ 是结构态射），并且将 $X_m$ 等同于
$X_n\times_{Y_n}Y_m$，这是由（II，1.4.6）立即看出的。归纳系统
$(X_n)$ 的归纳极限形式预概形 $\mathfrak X$ 因而局部 Noether，并且结构态射
$f:\mathfrak X\to\mathfrak Y$，即系统 $(f_n)$ 的归纳极限，是有限态射
（\hyperref[III.4.8.1-fr]{4.8.1} 及 II，10.12.3.1）；此外，我们已在
（\hyperref[III.4.8.6-fr]{4.8.6}）的证明中看到，
$f_*(\mathcal O_{\mathfrak X})$ 是诸 $\mathcal A_n$ 的逆极限，故等于
$\mathcal A$（I，10.10.6）。唯一性断言是下面更一般结果的推论：
\end{proof}

\begin{proposition}[4.8.8]
\label{III.4.8.8-fr}
设 $\mathfrak Y$ 是局部 Noether 形式预概形，$\mathfrak X$、
$\mathfrak X'$ 是两个在 $\mathfrak Y$ 上有限的形式
$\mathfrak Y$-预概形，$f:\mathfrak X\to\mathfrak Y$、
$f':\mathfrak X'\to\mathfrak Y$ 是结构态射。存在从
\[
\operatorname{Hom}_{\mathfrak Y}(\mathfrak X,\mathfrak X')
\quad\hbox{到}\quad
\operatorname{Hom}_{\mathcal O_{\mathfrak Y}}
\bigl(f'_*(\mathcal O_{\mathfrak X'}),f_*(\mathcal O_{\mathfrak X})\bigr)
\footnotemark.
\]
的典范双射。
\footnotetext{最后一个表达式表示 $\mathcal O_{\mathfrak Y}$-代数同态
$f'_*(\mathcal O_{\mathfrak X'})\to f_*(\mathcal O_{\mathfrak X})$ 的集合。}
\end{proposition}

\begin{proof}
该映射 $h\to\mathcal A(h)$ 的定义与（II，1.1.2）中相同；为证明它是
双射，立即归结到 $\mathfrak Y=\operatorname{Spf}(B)$ 是 Noether 仿射
形式概形的情形。此时 $\mathfrak X=\operatorname{Spf}(A)$、
$\mathfrak X'=\operatorname{Spf}(A')$，其中 $A$ 和 $A'$ 是两个有限
$B$-代数，并且 $f_*(\mathcal O_{\mathfrak X})=A^\Delta$、
$f'_*(\mathcal O_{\mathfrak X'})=A'^\Delta$。于是结论来自以下两种一一对应：
一方面，$\mathfrak Y$-态射 $\mathfrak X\to\mathfrak X'$ 与作为代数同态
（必然连续）的 $B$-同态 $A'\to A$ 之间的一一对应（I，10.2.2）；
另一方面，$B$-模同态 $A'\to A$ 与
$\mathcal O_{\mathfrak Y}$-模同态 $A'^\Delta\to A^\Delta$ 之间的一一对应
（I，10.10.2.3）。
\end{proof}

\begin{corollary}[4.8.9]
\label{III.4.8.9-fr}
在（\hyperref[III.4.8.8-fr]{4.8.8}）定义的典范一一对应中，闭浸入
$\mathfrak X\to\mathfrak X'$ 对应于满的 $\mathcal O_{\mathfrak Y}$-代数同态
$f'_*(\mathcal O_{\mathfrak X'})\to f_*(\mathcal O_{\mathfrak X})$。
\end{corollary}

\begin{proof}
这个问题仍是 $\mathfrak Y$ 上局部的，因而归结为局部 Noether 形式预概形的
闭浸入之定义（I，10.14.2）。
\end{proof}

\begin{corollary}[4.8.10]
\label{III.4.8.10-fr}
记号和假设同（\hyperref[III.4.8.1-fr]{4.8.1}）。进制态射 $f$ 是闭浸入
的充分必要条件是 $f_0$ 是（通常预概形的）闭浸入。
\end{corollary}

\begin{proof}
这立即来自（\hyperref[III.4.8.9-fr]{4.8.9}）以及凝聚
$\mathcal O_{\mathfrak Y}$-模同态的满射判据（I，10.11.5）。
\end{proof}

\begin{proposition}[4.8.11]
\label{III.4.8.11-fr}
局部 Noether 形式预概形的态射 $f:\mathfrak X\to\mathfrak Y$

\oldpage[III]{149}

有限的充分必要条件是它固有，且其纤维 $f^{-1}(y)$ 有限（对每个
$y\in\mathfrak Y$）。
\end{proposition}

\begin{proof}
由定义（3.4.1 及 4.8.2），立即归结为 $f_0$ 的同一命题
（记号同（\hyperref[III.4.8.1-fr]{4.8.1}）），这正是（4.4.2）。
\end{proof}

\section{凝聚代数层的一个存在定理}
\label{section:III.5-fr}

\subsection{定理的陈述}
\label{subsection:III.5.1-fr}

\begin{env}[5.1.1]
\label{III.5.1.1-fr}
设 $A$ 是 Noether 进制环，$\mathfrak J$ 是 $A$ 的定义理想，因而 $A$
对 $\mathfrak J$-进制拓扑是分离且完备的。若
$Y=\operatorname{Spec}(A)$，则仿射形式预概形 $\operatorname{Spf}(A)$
可与 $Y$ 沿闭子集 $Y'=\mathrm V(\mathfrak J)$ 的完备化 $\widehat Y$
等同（I，10.10.1）。设 $X$ 是有限型的（通常意义下的）$Y$-预概形，
$f:X\to Y$ 是结构态射；以 $\widehat X$ 表示 $X$ 沿闭子集
$X'=f^{-1}(Y')$ 的完备化，亦即形式 $\widehat Y$-预概形
$X\times_Y\widehat Y$；以 $\widehat f:\widehat X\to\widehat Y$ 表示 $f$
在完备化上的延拓；最后，对任意凝聚 $\mathcal O_X$-模 $\mathcal F$，
以 $\widehat{\mathcal F}$ 表示它的完备化 $\mathcal F_{/X'}$，这是一个
凝聚 $\mathcal O_{\widehat X}$-模。
\end{env}

\begin{proposition}[5.1.2]
\label{III.5.1.2-fr}
在（\hyperref[III.5.1.1-fr]{5.1.1}）的假设与记号下，设 $\mathcal F$
是一个凝聚 $\mathcal O_X$-模，其支集在 $Y$ 上固有（II，5.4.10）。
于是典范同态（4.1.4）
\[
\rho_i:\mathrm H^i(X,\mathcal F)\longrightarrow
\mathrm H^i(\widehat X,\widehat{\mathcal F})
\]
都是同构。
\end{proposition}

\begin{proof}
因为 $\mathrm H^i(X,\mathcal F)$ 是有限型 $A$-模（3.2.4），从而可与它
关于 $\mathfrak J$-预进制拓扑的分离完备化等同
（$0_{\mathrm I}$，7.3.6），故该命题只是（4.1.10）的一个特例。
\end{proof}

回忆：对凝聚 $\mathcal O_X$-模的任意正合列
$0\to\mathcal F'\to\mathcal F\to\mathcal F''\to0$，典范同构 $\rho_i$
均与连接同态交换（（$0$，12.1.6）及（I，10.8.9））。

\begin{corollary}[5.1.3]
\label{III.5.1.3-fr}
设 $\mathcal F$、$\mathcal G$ 是两个凝聚 $\mathcal O_X$-模，并且它们的
支集之交在 $Y$ 上固有。则典范同态
\begin{equation}
\operatorname{Hom}_{\mathcal O_X}(\mathcal F,\mathcal G)
\longrightarrow
\operatorname{Hom}_{\mathcal O_{\widehat X}}
(\widehat{\mathcal F},\widehat{\mathcal G})
\tag{5.1.3.1}
\label{III.5.1.3.1.eq-fr}
\end{equation}
是一个同构；它把每个同态 $u:\mathcal F\to\mathcal G$ 映为其完备化
$\widehat u:\widehat{\mathcal F}\to\widehat{\mathcal G}$。此外，当态射 $f$
是闭的时，$\widehat u$ 为单射（相应地，满射）的充分必要条件是 $u$
为单射（相应地，满射）。
\end{corollary}

\begin{proof}
第一项断言是（4.5.3）的一个特例；这里仍然用到
（\ref{III.5.1.3.1.eq-fr}）左端是有限型 $A$-模，因而可与它的分离完备化
等同。为证明第二项，由（I，10.8.14）可知，$\widehat u$ 为单射
（相应地，满射）的充分必要条件是存在 $X'$ 的一个邻域，使得 $u$
在该邻域中为单射（相应地，满射）。所以结论由下面的引理推出：
\end{proof}

\oldpage[III]{150}

\begin{lemma}[5.1.3.1]
\label{III.5.1.3.1-fr}
在（\hyperref[III.5.1.1-fr]{5.1.1}）的假设下，再假设态射 $f$ 是闭的，
则 $X'$ 在 $X$ 中的每个邻域都等于 $X$。
\end{lemma}

\begin{proof}
首先可归约到 $f(X)=Y$ 的情形。事实上，依假设，$f(X)$ 是 $Y$ 的一个闭子集
$Z$；此外，可用 $f_{\mathrm{red}}$ 替换 $f$（I，6.3.4），从而假设 $X$
与 $Y$ 均既约；于是可用以 $Z$ 为底空间的 $Y$ 的既约闭子预概形替换 $Y$
（I，5.2.2），因为 $A$ 的每个理想都是闭的，所以 $A$ 的每个商环都是
Noether 进制环。此时 $f(X')=Y'$；若 $V$ 是 $X'$ 在 $X$ 中的一个开邻域，
则依假设 $f(X-V)$ 在 $Y$ 中是闭的，并且不与 $Y'$ 相交；但除非 $X-V$
为空，这是不可能的，因为 $\mathfrak J$ 包含于 $A$ 的根中
（I，1.1.15 及 $0_{\mathrm I}$，7.1.10）。结论由此得到。
\end{proof}

如果只考虑支集在 $Y$ 上固有的凝聚 $\mathcal O_X$-模，那么用范畴的语言，
（\hyperref[III.5.1.3-fr]{5.1.3}）就是说函子
$\mathcal F\mapsto\widehat{\mathcal F}$ 从上述 $\mathcal O_X$-模的范畴
到凝聚 $\mathcal O_{\widehat X}$-模范畴是全忠实的，因而给出前一范畴与
后一范畴的一个全子范畴之间的等价（$0$，8.1.6）。存在定理将证明：当 $X$
在 $Y$ 上固有时，这个全子范畴实际上就是所有凝聚
$\mathcal O_{\widehat X}$-模的范畴。确切地说：

\begin{theorem}[5.1.4]
\label{III.5.1.4-fr}
设 $A$ 是 Noether 进制环，$Y=\operatorname{Spec}(A)$，$\mathfrak J$ 是
$A$ 的定义理想，$Y'=\mathrm V(\mathfrak J)$，$f:X\to Y$ 是有限型分离态射，
$X'=f^{-1}(Y')$。设 $\widehat Y=Y_{/Y'}=\operatorname{Spf}(A)$、
$\widehat X=X_{/X'}$ 分别是 $Y$ 与 $X$ 沿 $Y'$ 与 $X'$ 的完备化，
$\widehat f:\widehat X\to\widehat Y$ 是 $f$ 在完备化上的延拓；那么函子
$\mathcal F\mapsto\mathcal F_{/X'}=\widehat{\mathcal F}$ 给出支集在
$\operatorname{Spec}(A)$ 上固有的凝聚 $\mathcal O_X$-模范畴与支集在
$\operatorname{Spf}(A)$ 上固有的凝聚 $\mathcal O_{\widehat X}$-模范畴
之间的等价。
\end{theorem}

换言之，考虑到（\hyperref[III.5.1.3-fr]{5.1.3}）：

\begin{corollary}[5.1.5]
\label{III.5.1.5-fr}
一个 $\mathcal O_{\widehat X}$-模同构于某个支集在
$\operatorname{Spec}(A)$ 上固有的凝聚 $\mathcal O_X$-模的完备化，
其充分必要条件是它凝聚且支集在 $\operatorname{Spf}(A)$ 上固有。
\end{corollary}

最重要的情形是

\begin{corollary}[5.1.6]
\label{III.5.1.6-fr}
假设 $X$ 在 $Y=\operatorname{Spec}(A)$ 上固有。则函子
$\mathcal F\mapsto\widehat{\mathcal F}$ 给出凝聚 $\mathcal O_X$-模范畴与
凝聚 $\mathcal O_{\widehat X}$-模范畴之间的等价。
\end{corollary}

\begin{env}[5.1.7]
\label{III.5.1.7-fr}
\emph{评注。}
考虑到形式预概形上凝聚层的刻画（I，10.11.3），可见在
（\hyperref[III.5.1.1-fr]{5.1.1}）的条件下，给定一个支集在
$\operatorname{Spec}(A)$ 上固有的凝聚 $\mathcal O_X$-模，等价于（置
$Y_n=\operatorname{Spec}(A/\mathfrak J^{n+1})$ 及 $X_n=X\times_Y Y_n$）
给定凝聚 $\mathcal O_{X_n}$-模的一个逆系统 $(\mathcal F_n)$，使得当
$m\leq n$ 时有
\[
\mathcal F_m=\mathcal F_n\otimes_{\mathcal O_{Y_n}}\mathcal O_{Y_m}
\quad\text{（亦即 }\mathcal F_m=\mathcal F_n/\mathfrak J^{m+1}\mathcal F_n\text{）}
\]
并且 $\mathcal F_0$ 的支集是 $X_0$ 的一个在 $Y_0$ 上固有的子集。借助
（I，10.11.4），凝聚 $\mathcal O_X$-模的同态同样可解释为凝聚
$\mathcal O_{X_n}$-模逆系统的同态。

\oldpage[III]{151}

在所有已知的应用中，$A$ 实际上都是 Noether 局部进制环，因而 $Y_n$
都是 Artin 局部环的谱；本段和前面各段的结果在很大程度上把 Noether
局部进制环上的代数几何归约为 Artin 局部环上的代数几何。
\end{env}

\begin{corollary}[5.1.8]
\label{III.5.1.8-fr}
在（\hyperref[III.5.1.4-fr]{5.1.4}）的条件下，映射
$Z\mapsto\widehat Z=Z_{/(Z\cap X')}$ 是从 $X$ 的在 $Y$ 上固有的闭子预概形
$Z$ 之集到 $\widehat X$ 的在 $\widehat Y$ 上固有的闭形式子预概形之集的双射。
\end{corollary}

\begin{proof}
事实上，$\widehat X$ 的一个闭形式子预概形具有形式
$(T,(\mathcal O_{\widehat X}/\mathcal A)|T)$，其中 $\mathcal A$ 是
$\mathcal O_{\widehat X}$ 的一个凝聚理想（I，10.14.2）；若 $T$ 在
$\widehat Y$ 上固有，则由（\hyperref[III.5.1.4-fr]{5.1.4}）可知
$\mathcal O_{\widehat X}/\mathcal A$ 同构于一个形如 $\widehat{\mathcal F}$
的 $\mathcal O_{\widehat X}$-模，其中 $\mathcal F$ 是支集在 $Y$ 上固有的
凝聚 $\mathcal O_X$-模；此外，由（\hyperref[III.5.1.3-fr]{5.1.3}）可知，
典范同态 $\mathcal O_{\widehat X}\to\mathcal O_{\widehat X}/\mathcal A$
具有形式 $\widehat u$，其中 $u:\mathcal O_X\to\mathcal F$ 是
$\mathcal O_X$-模的满同态。因此 $\mathcal F$ 具有形式
$\mathcal O_X/\mathcal N$，其中 $\mathcal N$ 是 $\mathcal O_X$ 的一个凝聚
理想，而且 $\mathcal A=\widehat{\mathcal N}$（I，10.8.8），结论由此得到
（I，10.14.7）。
\end{proof}

\subsection{存在定理的证明：射影与拟射影情形}
\label{subsection:III.5.2-fr}

\begin{env}[5.2.1]
\label{III.5.2.1-fr}
在（\hyperref[III.5.1.4-fr]{5.1.4}）的条件下，如果一个凝聚
$\mathcal O_{\widehat X}$-模同构于某个支集在 $Y$ 上固有的凝聚
$\mathcal O_X$-模 $\mathcal F$ 的完备化 $\widehat{\mathcal F}$，我们暂称
它为可代数化的。
\end{env}

\begin{lemma}[5.2.2]
\label{III.5.2.2-fr}
设 $\mathcal F'$、$\mathcal G'$ 是两个可代数化的
$\mathcal O_{\widehat X}$-模。对任意同态 $u:\mathcal F'\to\mathcal G'$，
$\operatorname{Ker}(u)$、$\operatorname{Im}(u)$ 和
$\operatorname{Coker}(u)$ 都是可代数化的。
\end{lemma}

\begin{proof}
事实上，若 $\mathcal F'=\widehat{\mathcal F}$、
$\mathcal G'=\widehat{\mathcal G}$，其中 $\mathcal F$ 与 $\mathcal G$
是支集在 $Y$ 上固有的凝聚 $\mathcal O_X$-模，则
$u=\widehat v$，其中 $v:\mathcal F\to\mathcal G$ 是一个同态
（\hyperref[III.5.1.3-fr]{5.1.3}）。由函子
$\mathcal F\mapsto\widehat{\mathcal F}$ 的正合性，
$\operatorname{Ker}(\widehat v)$ 同构于
$\widehat{\operatorname{Ker}(v)}$；又因为 $\operatorname{Ker}(v)$ 的支集
包含于 $\mathcal F$ 的支集中，可见 $\operatorname{Ker}(u)$ 可代数化；
对 $\operatorname{Im}(u)$ 与 $\operatorname{Coker}(u)$ 的证明相同。
\end{proof}

\begin{proposition}[5.2.3]
\label{III.5.2.3-fr}
设 $A$ 是 Noether 进制环，$\mathfrak J$ 是 $A$ 的定义理想，
$\mathfrak Y=\operatorname{Spf}(A)$，
$f:\mathfrak X\to\mathfrak Y$ 是形式预概形的固有态射。置
$Y_k=\operatorname{Spec}(A/\mathfrak J^{k+1})$、
$X_k=\mathfrak X\times_{\mathfrak Y}Y_k$，并且对任意
$\mathcal O_{\mathfrak X}$-模 $\mathcal F$，置
\[
\mathcal F_k=\mathcal F\otimes_{\mathcal O_{\mathfrak X}}\mathcal O_{X_k}
=\mathcal F/\mathfrak J^{k+1}\mathcal F.
\]
设 $\mathcal L$ 是可逆 $\mathcal O_{\mathfrak X}$-模，并假设
$\mathcal L_0=\mathcal L/\mathfrak J\mathcal L$ 是丰沛
$\mathcal O_{X_0}$-模；对任意 $\mathcal O_{\mathfrak X}$-模 $\mathcal F$
及任意整数 $n$，置 $\mathcal F(n)=\mathcal F\otimes\mathcal L^{\otimes n}$。
那么，对任意凝聚 $\mathcal O_{\mathfrak X}$-模 $\mathcal F$，存在一个整数
$n_0$，使得对所有 $n\geq n_0$，下列性质成立：
\begin{enumerate}
\item[(i)] 对所有 $k\geq0$，典范同态
$\mathrm H^0(\mathfrak X,\mathcal F(n))\to
\mathrm H^0(\mathfrak X,\mathcal F_k(n))$ 是满射。
\item[(ii)] 对所有 $q>0$，有 $\mathrm H^q(\mathfrak X,\mathcal F(n))=0$。
\end{enumerate}
\end{proposition}

\begin{proof}
已知 $\mathfrak X$ 与 $X_0$ 的底空间相同；诸层
\[
\mathcal M_k=\mathfrak J^k\mathcal F/\mathfrak J^{k+1}\mathcal F
\]
被 $\mathfrak J$ 零化，故可视为凝聚 $\mathcal O_{X_0}$-模
（$0_{\mathrm I}$，5.3.10）；此外，若置
$\mathcal M_k(n)=\mathcal M_k\otimes_{\mathcal O_{X_0}}
\mathcal L_0^{\otimes n}$，立即可见
\[
\mathcal M_k(n)=\mathfrak J^k\mathcal F(n)/
\mathfrak J^{k+1}\mathcal F(n).
\]
注意，因为 $\mathcal L_0$ 对 $f_0:X_0\to Y_0$ 是丰沛的，并且

\oldpage[III]{152}

$f_0$ 是固有的，所以可推出 $f_0$ 是射影的（II，5.5.4）。设
\[
S=\bigoplus_{k\geq0}\mathfrak J^k/\mathfrak J^{k+1}
\]
为 $A$ 的 $\mathfrak J$-进滤过所关联的分次 $A$-代数；它是有限型的，
因为 $A$ 是 Noether 环。若置 $\mathcal S'=f_0^*(\widetilde S)$、
$\mathcal M=\bigoplus_{k\geq0}\mathcal M_k$，则 $\mathcal S'$ 是有限型
拟凝聚 $\mathcal O_{X_0}$-代数，而 $\mathcal M$ 是有限型拟凝聚分次
$\mathcal S'$-模（因为 $\mathcal F_0$ 凝聚且生成 $\mathcal S'$-模
$\mathcal M$）。因此满足定理（2.4.1，(ii)）的应用条件，从而存在
$n_0$，使得当 $n\geq n_0$ 时，对每个 $k$ 都有
\begin{equation}
\mathrm H^q(X_0,\mathcal M_k(n))=0\qquad\text{对所有 }q>0.
\tag{5.2.3.1}
\label{III.5.2.3.1-fr}
\end{equation}
于是，当 $q>0$ 且 $n\geq n_0$ 时，也有
$\mathrm H^q(\mathfrak X,\mathcal M_k(n))=0$，这里把
$\mathcal M_k(n)$ 视为 $\mathcal O_{\mathfrak X}$-模。对正合列
\[
0\to
\frac{\mathfrak J^k\mathcal F(n)}{\mathfrak J^{k+1}\mathcal F(n)}
\to
\frac{\mathfrak J^h\mathcal F(n)}{\mathfrak J^{k+1}\mathcal F(n)}
\to
\frac{\mathfrak J^h\mathcal F(n)}{\mathfrak J^k\mathcal F(n)}
\to0,
\]
应用上同调正合列，首先通过对 $k-h$ 归纳得到：当
$0\leq h<k$、$n\geq n_0$ 且 $q>0$ 时，
\begin{equation}
\mathrm H^q\left(\mathfrak X,
\frac{\mathfrak J^h\mathcal F(n)}{\mathfrak J^k\mathcal F(n)}\right)=0
\tag{5.2.3.2}
\label{III.5.2.3.2-fr}
\end{equation}
特别地，取 $h=0$，得到
\begin{equation}
\mathrm H^q(\mathfrak X,\mathcal F_k(n))=0
\qquad\text{当 }n\geq n_0,\ k\geq0\text{ 且 }q>0.
\tag{5.2.3.3}
\label{III.5.2.3.3-fr}
\end{equation}
取 $h=0$，上同调正合列的另一部分给出正合列
\begin{equation}
\mathrm H^0(\mathfrak X,\mathcal F_{k+1}(n))\to
\mathrm H^0(\mathfrak X,\mathcal F_k(n))\to
\mathrm H^1\left(\mathfrak X,
\frac{\mathfrak J^k\mathcal F(n)}{\mathfrak J^{k+1}\mathcal F(n)}\right)=0
\tag{5.2.3.4}
\label{III.5.2.3.4-fr}
\end{equation}
由此可知，当 $h\leq k$ 时，典范映射
\begin{equation}
\mathrm H^0(\mathfrak X,\mathcal F_k(n))\longrightarrow
\mathrm H^0(\mathfrak X,\mathcal F_h(n))
\tag{5.2.3.5}
\label{III.5.2.3.5-fr}
\end{equation}
是满射。因此对每个 $q$，逆系统
$(\mathrm H^q(\mathfrak X,\mathcal F_k(n)))_{k\geq0}$ 在
$n\geq n_0$ 时满足条件 (ML)。另一方面，$\mathfrak X$ 的每个仿射形式开集
$U$ 在每个 $X_k$ 中也都是仿射开集（I，10.5.2），所以对所有 $q>0$，
有 $\mathrm H^q(U,\mathcal F_k(n))=0$（I，1.3.1），并且当 $h\leq k$
时，$\mathrm H^0(U,\mathcal F_k(n))\to\mathrm H^0(U,\mathcal F_h(n))$
是满射（I，1.3.9）。故满足（$0$，13.3.1）的应用条件，从而当
$n\geq n_0$ 时：
\begin{enumerate}
\item[$1^\circ$] 对所有 $q>0$，映射
\[
\mathrm H^q(\mathfrak X,\mathcal F(n))\longrightarrow
\varprojlim_k\mathrm H^q(\mathfrak X,\mathcal F_k(n))
\]
是双射；因此由（\ref{III.5.2.3.3-fr}），
$\mathrm H^q(\mathfrak X,\mathcal F(n))=0$。
\item[$2^\circ$] 同态
\[
\mathrm H^0(\mathfrak X,\mathcal F(n))\longrightarrow
\varprojlim_k\mathrm H^0(\mathfrak X,\mathcal F_k(n))
\]
是双射；另一方面，因为同态（\ref{III.5.2.3.5-fr}）是满射，故每个同态
\[
\varprojlim_k\mathrm H^0(\mathfrak X,\mathcal F_k(n))
\longrightarrow\mathrm H^0(\mathfrak X,\mathcal F_h(n)),
\]
也都是满射，
\end{enumerate}
这就完成了证明。
\end{proof}

\begin{corollary}[5.2.4]
\label{III.5.2.4-fr}
在（\hyperref[III.5.2.3-fr]{5.2.3}）的假设下，对任意凝聚
$\mathcal O_{\mathfrak X}$-模 $\mathcal F$，存在整数 $N$，使得当
$n\geq N$ 时，$\mathcal F(n)$ 由它在

\oldpage[III]{153}

$\mathfrak X$ 上的截面生成；换言之，$\mathcal F$ 同构于一个形如
$(\mathcal O_{\mathfrak X}(-n))^k$ 的 $\mathcal O_{\mathfrak X}$-模的商。
\end{corollary}

\begin{proof}
因为 $X_0$ 是 Noether 的，由关于 $\mathcal L_0$ 的假设及（II，4.5.5），
存在 $n_0$，使得当 $n\geq n_0$ 时，$\mathcal F_0(n)$ 由它在
$\mathfrak X$ 上的截面生成；另一方面，可假设 $n_0$ 取得足够大，使得当
$n\geq n_0$ 时，同态
\[
\Gamma(\mathfrak X,\mathcal F(n))\longrightarrow
\Gamma(\mathfrak X,\mathcal F_0(n))
\]
是满射（\hyperref[III.5.2.3-fr]{5.2.3}）。因此存在有限多个截面
$s_i\in\Gamma(\mathfrak X,\mathcal F(n))$，其在
$\Gamma(\mathfrak X,\mathcal F_0(n))$ 中的像生成 $\mathcal F_0(n)$
（$0_{\mathrm I}$，5.2.3）。因为 $\mathfrak J$ 包含于 $\mathfrak X$
每一点处局部环的极大理想中，对这些局部环应用 Nakayama 引理可知，
诸 $s_i$ 生成 $\mathcal F(n)$（$0_{\mathrm I}$，5.1.1）。
\end{proof}

\begin{env}[5.2.5]
\label{III.5.2.5-fr}
\emph{存在定理的证明：射影情形。}
记号同（\hyperref[III.5.1.4-fr]{5.1.4}），假设 $f$ 是射影的，因而存在一个
丰沛 $\mathcal O_X$-模 $\mathcal L$（I，5.5.4）。依定义，
$X_n=\widehat X\times_{\widehat Y}Y_n$ 等于 $X$ 的闭子预概形
$X\times_Y Y_n=f^{-1}(Y_n)$；若
$\mathcal L'=\mathcal L\otimes_{\mathcal O_X}\mathcal O_{\widehat X}$
是 $\mathcal L$ 的完备化，则
$\mathcal L'_0=\mathcal L/\mathfrak J\mathcal L$，视为
$\mathcal O_{X_0}$-模；已知 $\mathcal L'_0$ 是丰沛的
（II，4.6.13，(i bis)）。所以可对 $\mathcal L'$ 及任意凝聚
$\mathcal O_{\widehat X}$-模 $\mathcal F$ 应用推论
（\hyperref[III.5.2.4-fr]{5.2.4}）；于是可见 $\mathcal F$ 同构于
\[
\mathcal G=(\mathcal L'^{\otimes(-n)})^k
\]
的一个商，其中 $n>0$ 与 $k>0$ 是适当的整数。显然，$\mathcal G$ 是
$(\mathcal L^{\otimes(-n)})^k$ 的完备化（I，10.8.10），所以是可代数化的。
再考虑典范满同态 $u:\mathcal G\to\mathcal F$，并设
$\mathcal H=\operatorname{Ker}(u)$；它是凝聚 $\mathcal O_{\widehat X}$-模
（$0_{\mathrm I}$，5.3.4）。同样可见，存在可代数化的
$\mathcal O_{\widehat X}$-模 $\mathcal K$ 及同态
$v:\mathcal K\to\mathcal G$，使得 $\mathcal H=\operatorname{Im}(v)$。
于是 $\mathcal F=\operatorname{Coker}(v)$，由
（\hyperref[III.5.2.2-fr]{5.2.2}），$\mathcal F$ 可代数化。
\end{env}

\begin{env}[5.2.6]
\label{III.5.2.6-fr}
\emph{存在定理的证明：拟射影情形。}
记号仍同（\hyperref[III.5.1.4-fr]{5.1.4}），现假设 $f$ 是拟射影的。
于是存在射影态射 $g:Z\to Y$，使得 $X$ 可与 $Z$ 的一个开集上所诱导的
$Y$-预概形等同（II，5.3.2）；若置 $Z'=g^{-1}(Y')$，则
$X'=X\cap Z'$。所以完备化 $\widehat X=X_{/X'}$ 可与完备化
$\widehat Z=Z_{/Z'}$ 在 $\widehat Z$ 的开集 $X\cap Z'$ 上诱导的形式预概形等同
（I，10.8.5）。设 $\mathcal F'$ 是凝聚 $\mathcal O_{\widehat X}$-模，
其支集 $T'$ 在 $\widehat Y$ 上固有；依定义，这意味着存在 $X'$ 的一个闭子预概形，
其底空间为 $T'\subset X'$，使得 $f$ 的限制 $T'\to Y'$ 是固有的；
由此可知 $T'$ 在 $Y$ 上固有，因而在 $Z'$ 中是闭的（II，5.4.10）。
于是 $\mathcal F'$ 是 $\mathcal O_{\widehat Z}$-模 $\mathcal G'$ 在
$\widehat X$ 上诱导的层；这个模由 $\mathcal F'$（定义在
$\widehat Z$ 的开集 $\widehat X$ 上）与 $\widehat Z$ 的开集
$\widehat Z-T'$ 上的层 $0$ 黏合而得，这两个层在交开集
$\widehat X-T'$ 上相同。显然 $\mathcal G'$ 是凝聚的；由
（\hyperref[III.5.2.5-fr]{5.2.5}），存在凝聚 $\mathcal O_Z$-模
$\mathcal G$，使得 $\mathcal G'=\widehat{\mathcal G}$；设 $T$ 为
$\mathcal G$ 的支集，于是 $T'=T\cap Z'$（I，10.8.12）。若 $h$ 是 $g$
在以 $T$ 为底空间的 $Z$ 的既约闭子预概形上的限制，则
$T'=h^{-1}(Y')=T\cap g^{-1}(Y')$，因而 $X\cap T$ 是 $T$ 的一个包含
$T'$ 的开集。

\oldpage[III]{154}

因为 $h$ 是固有的（II，5.4.2），从而是闭的，由
（\hyperref[III.5.1.3.1-fr]{5.1.3.1}）可知 $X\cap T=T$，换言之
$T\subset X$；又因为 $T$ 在 $Z$ 中是闭的，故 $T$ 在 $Y$ 上固有。
若 $\mathcal F$ 是 $\mathcal G$ 在 $X$ 上诱导的层，则其完备化
$\widehat{\mathcal F}$ 是 $\widehat{\mathcal G}$ 在 $\widehat X$ 上诱导的层
（I，10.8.4），因而等于 $\mathcal F'$；证明完成。
\end{env}

\subsection{存在定理的证明：一般情形}
\label{subsection:III.5.3-fr}

\begin{lemma}[5.3.1]
\label{III.5.3.1-fr}
在（\hyperref[III.5.1.4-fr]{5.1.4}）的条件下，若
\[
0\to\mathcal H\to\mathcal F\to\mathcal G\to0
\]
是凝聚 $\mathcal O_{\widehat X}$-模的正合列，并且 $\mathcal G$ 与
$\mathcal H$ 均可代数化，则 $\mathcal F$ 可代数化。
\end{lemma}

\begin{proof}
事实上，假设 $\mathcal G=\widehat{\mathcal B}$、
$\mathcal H=\widehat{\mathcal C}$，其中 $\mathcal B$ 与 $\mathcal C$
是支集在 $Y$ 上固有的凝聚 $\mathcal O_X$-模；$\mathcal F$ 典范地定义了
$A$-模
\[
\operatorname{Ext}^1_{\mathcal O_{\widehat X}}
(\widehat X;\widehat{\mathcal B},\widehat{\mathcal C})
\]
的一个元素（$0$，12.3.2），而这些假设蕴含此 $A$-模典范同构于
\[
\operatorname{Ext}^1_{\mathcal O_X}(X;\mathcal B,\mathcal C)
\]
（4.5.2）；因此存在凝聚 $\mathcal O_X$-模的正合列
$0\to\mathcal C\to\mathcal A\to\mathcal B\to0$，使得与
$\mathcal A$ 对应的 $\operatorname{Ext}^1_{\mathcal O_X}(X;\mathcal B,\mathcal C)$
的元素的典范像，就是与 $\mathcal F$ 对应的
$\operatorname{Ext}^1_{\mathcal O_{\widehat X}}
(\widehat X;\widehat{\mathcal B},\widehat{\mathcal C})$ 的元素。但依定义
（考虑（I，10.8.8，(iii)）），这意味着 $\mathcal F$ 同构于
$\widehat{\mathcal A}$，引理遂得证；因为 $\operatorname{Supp}(\mathcal A)$
包含于 $\operatorname{Supp}(\mathcal B)$ 与
$\operatorname{Supp}(\mathcal C)$ 的并集，故它在 $Y$ 上固有。
\end{proof}

\begin{corollary}[5.3.2]
\label{III.5.3.2-fr}
在（\hyperref[III.5.1.1-fr]{5.1.1}）的条件下，设
$u:\mathcal F\to\mathcal G$ 是凝聚 $\mathcal O_{\widehat X}$-模的同态；
若 $\mathcal G$、$\operatorname{Ker}(u)$ 和 $\operatorname{Coker}(u)$
均可代数化，则 $\mathcal F$ 也可代数化。
\end{corollary}

\begin{proof}
事实上，对同态 $\mathcal G\to\operatorname{Coker}(u)$ 应用引理
（\hyperref[III.5.2.2-fr]{5.2.2}），可知 $\operatorname{Im}(u)$
可代数化；然后只须对正合列
\[
0\to\operatorname{Ker}(u)\to\mathcal F\to\operatorname{Im}(u)\to0.
\]
应用引理（\hyperref[III.5.3.1-fr]{5.3.1}）。
\end{proof}

\begin{lemma}[5.3.3]
\label{III.5.3.3-fr}
在（\hyperref[III.5.1.1-fr]{5.1.1}）的条件下，设 $h:Z\to Y$ 是有限型
态射，$\widehat Z$ 是 $Z$ 沿 $Z'=h^{-1}(Y')$ 的完备化，
$g:Z\to X$ 是固有 $Y$-态射，
$\widehat g:\widehat Z\to\widehat X$ 是它在完备化上的延拓。对任意可代数化的
$\mathcal O_{\widehat Z}$-模 $\mathcal F'$，
$\widehat g_*(\mathcal F')$ 是可代数化的 $\mathcal O_{\widehat X}$-模。
\end{lemma}

\begin{proof}
事实上，若 $\mathcal F$ 是满足 $\mathcal F'=\widehat{\mathcal F}$ 的凝聚
$\mathcal O_Z$-模，则由第一比较定理（4.1.5），
$\widehat g_*(\widehat{\mathcal F})$ 同构于 $g_*(\mathcal F)$ 的完备化。
\end{proof}

\begin{lemma}[5.3.4]
\label{III.5.3.4-fr}
设 $X$ 是（通常意义下的）Noether 概形，$X'$ 是 $X$ 的一个闭子集，
$f:Z\to X$ 是固有态射，$Z'=f^{-1}(X')$，
$\widehat X=X_{/X'}$，$\widehat Z=Z_{/Z'}$，
$\widehat f:\widehat Z\to\widehat X$ 是 $f$ 在完备化上的延拓。设
$\mathcal M$ 是 $\mathcal O_X$ 的凝聚理想，并且若
$U=X-\operatorname{Supp}(\mathcal O_X/\mathcal M)$，则 $f$ 的限制
$f^{-1}(U)\to U$ 是同构。那么，对任意凝聚
$\mathcal O_{\widehat X}$-模 $\mathcal F$，存在整数 $n>0$，使得典范同态
\[
\rho_{\mathcal F}:\mathcal F\longrightarrow
\widehat f_*\bigl(\widehat f^{,*}(\mathcal F)\bigr)
\]
（$0_{\mathrm I}$，4.4.3）的核与余核均被
$\widehat{\mathcal M}^{,n}$ 零化。
\end{lemma}

\begin{proof}
可只考虑 $X=\operatorname{Spec}(B)$ 的情形，其中 $B$ 是 Noether 环，
从而 $X'=\mathrm V(\mathfrak K)$，其中 $\mathfrak K$ 是 $B$ 的理想。
我们将说明，可归约到 $B$ 是 Noether 进制环且 $\mathfrak K$ 是 $B$
的定义理想的情形。事实上，设 $B_1$ 是 $B$ 关于 $\mathfrak K$-预进制
拓扑的分离完备化；若 $\mathfrak K_1=\mathfrak K B_1$，则 $B_1$ 是一个

\oldpage[III]{155}

以 $\mathfrak K_1$ 为定义理想的 Noether 进制环。置
$X_1=\operatorname{Spec}(B_1)$，并设 $h:X_1\to X$ 是典范同态
$B\to B_1$ 所对应的态射；若 $X'_1=h^{-1}(X')$，则
$X'_1=\mathrm V(\mathfrak K_1)$。最后置
$Z_1=Z\times_X X_1=Z_{(X_1)}$、$f_1=f_{(X_1)}:Z_1\to X_1$；这是固有
态射（II，5.4.2）。以 $\widehat X_1$ 表示 $X_1$ 沿 $X'_1$ 的完备化，
以 $\widehat Z_1=Z_1\times_{X_1}\widehat X_1$ 表示 $Z_1$ 沿
$Z'_1=f_1^{-1}(X'_1)$ 的完备化，以 $\widehat f_1$ 表示 $f_1$ 在完备化上的
延拓。立即可见，$h$ 在完备化上的延拓
$\widehat h:\widehat X_1\to\widehat X$ 是一个同构，它对应于 $B_1$
的恒等映射（I，10.9.1）；由此推出相应的同态
$\widehat Z_1\to\widehat Z$ 也是同构，这些同构将 $\widehat f_1$ 与
$\widehat f$ 等同。最后，$\mathcal M_1=h^*(\mathcal M)$ 是
$\mathcal O_{X_1}$ 的凝聚理想，并且
\[
\operatorname{Supp}(\mathcal O_{X_1}/\mathcal M_1)
=h^{-1}(\operatorname{Supp}(\mathcal O_X/\mathcal M))
\]
（I，9.1.13），所以若
$U_1=X_1-\operatorname{Supp}(\mathcal O_{X_1}/\mathcal M_1)$，则
$U_1=h^{-1}(U)$；由此立即得到，$f_1$ 的限制
$f_1^{-1}(U_1)\to U_1$ 是同构（I，3.2.7）；此外，完备化
$\widehat{\mathcal M}$ 与 $\widehat{\mathcal M}_1$ 由 $\widehat h$ 等同
（I，10.9.5）。因此 $X_1$、$X'_1$、$f_1$ 与 $\mathcal M_1$ 满足
（\hyperref[III.5.3.4-fr]{5.3.4}）的全部假设，故此后可假设 $B$ 是
Noether 进制环，$\mathfrak K$ 是 $B$ 的定义理想。这时
$\widehat X=\operatorname{Spf}(B)$，且 $\mathcal F=N^\Delta$，其中 $N$
是有限型 $B$-模，从而 $\mathcal F=\widehat{\mathcal G}$，其中
$\mathcal G$ 是凝聚 $\mathcal O_X$-模 $\widetilde N$（I，10.10.5）；因此
\[
\widehat f^{,*}(\mathcal F)=\widehat{f^*(\mathcal G)}
\]
（I，10.9.5）。此外，由第一比较定理（4.1.5），
\[
\widehat f_*\bigl(\widehat{f^*(\mathcal G)}\bigr)
\simeq\widehat{f_*\bigl(f^*(\mathcal G)\bigr)},
\]
并且由（\hyperref[III.5.1.3-fr]{5.1.3}），典范同态
$\rho_{\mathcal F}$ 正是 $\widehat{\rho}_{\mathcal G}$。而同态
\[
\rho_{\mathcal G}:\mathcal G\longrightarrow f_*\bigl(f^*(\mathcal G)\bigr)
\]
的核 $\mathcal P$ 与余核 $\mathcal R$ 是凝聚 $\mathcal O_X$-模，且依假设，
它们在 $U$ 上的限制显然为零。因此存在整数 $n>0$，使得
$\mathcal M^n\mathcal P=\mathcal M^n\mathcal R=0$（I，9.3.4）；由此得到
$\widehat{\mathcal M}^{,n}\widehat{\mathcal P}=
\widehat{\mathcal M}^{,n}\widehat{\mathcal R}=0$
（I，10.8.8 及 10.8.10）。
\end{proof}

\begin{env}[5.3.5]
\label{III.5.3.5-fr}
\emph{存在定理证明的结束。}
在（\hyperref[III.5.1.4-fr]{5.1.4}）的假设下，我们使用 Noether 归纳原理
（$0_{\mathrm I}$，2.2.2），即假设定理对 $X$ 的每个底空间不同于 $X$
的闭子预概形 $T$ 均成立（当然，$\widehat T$ 表示 $T$ 沿 $T\cap X'$
的完备化）。可假设 $X$ 非空。因为 $f$ 分离且有限型，可以应用 Chow 引理
（II，5.6.1）：于是存在 $Y$-概形 $Z$ 及 $Y$-态射 $g:Z\to X$，使得
结构态射 $h:Z\to Y$ 是拟射影的，态射 $g$ 是射影且满射的；此外，存在
$X$ 的一个非空开集 $U$，使得限制 $g^{-1}(U)\to U$ 是同构。设
$\mathcal M$ 是 $\mathcal O_X$ 的一个凝聚理想，它定义以 $X-U$ 为底空间
的闭子预概形（I，5.2.2）；设 $\mathcal F$ 是凝聚
$\mathcal O_{\widehat X}$-模，其支集 $E$ 在 $Y$ 上固有；以 $\widehat Z$
表示 $Z$ 沿 $h^{-1}(Y')$ 的完备化，以
$\widehat g:\widehat Z\to\widehat X$ 表示 $g$ 在完备化上的延拓。那么
$\widehat g^{,*}(\mathcal F)$ 是凝聚 $\mathcal O_{\widehat Z}$-模，其支集
包含于 $g^{-1}(E)$，从而在 $Y$ 上固有，因为 $g$ 是射影的，故是固有的
（II，5.4.6）。因为 $h$ 是拟射影的，由
（\hyperref[III.5.2.6-fr]{5.2.6}），$\widehat g^{,*}(\mathcal F)$
可代数化。由此推出

\oldpage[III]{156}

$\widehat g_*\bigl(\widehat g^{,*}(\mathcal F)\bigr)$ 是可代数化的
$\mathcal O_{\widehat X}$-模
（\hyperref[III.5.3.3-fr]{5.3.3}），因为 $g$ 是固有的。现在可对
$\mathcal F$ 与 $g$ 应用（\hyperref[III.5.3.4-fr]{5.3.4}）的结果：
同态 $\rho_{\mathcal F}:\mathcal F\to
\widehat g_*\bigl(\widehat g^{,*}(\mathcal F)\bigr)$ 的核 $\mathcal P$
与余核 $\mathcal R$ 被 $\widehat{\mathcal M}^{,n}$ 的某个幂零化；设 $T$
为 $\mathcal M^n$ 所定义的 $X$ 的闭子预概形，其底空间为 $X-U$；
$j:T\to X$ 是典范嵌入，所以其在完备化上的延拓
$\widehat j:\widehat T\to\widehat X$ 是典范嵌入（I，10.14.7）。因此可写成
$\mathcal P=\widehat j_*\bigl(\widehat j^{,*}(\mathcal P)\bigr)$ 及
$\mathcal R=\widehat j_*\bigl(\widehat j^{,*}(\mathcal R)\bigr)$；又因为
$U$ 非空，由归纳假设可知 $\widehat j^{,*}(\mathcal P)$ 与
$\widehat j^{,*}(\mathcal R)$ 是可代数化的 $\mathcal O_{\widehat T}$-模；
由（\hyperref[III.5.3.3-fr]{5.3.3}），$\mathcal P$ 与 $\mathcal R$
可代数化，于是可以应用（\hyperref[III.5.3.2-fr]{5.3.2}），最终证明
$\mathcal F$ 可代数化。

\hfill C.Q.F.D.
\end{env}

\subsection{应用：通常概形态射与形式概形态射的比较。可代数化形式概形}
\label{subsection:III.5.4-fr}

\begin{theorem}[5.4.1]
\label{III.5.4.1-fr}
设 $A$ 是 Noether 进制环，$\mathfrak J$ 是 $A$ 的定义理想，
$S=\operatorname{Spec}(A)$，$S'=\mathrm V(\mathfrak J)$。设
$u:X\to S$ 是固有态射，$v:Y\to S$ 是分离有限型态射，并分别以
$\widehat S$、$\widehat X$、$\widehat Y$ 表示 $S$、$X$、$Y$ 沿
$S'$、$u^{-1}(S')$、$v^{-1}(S')$ 的完备化。对任意 $S$-态射
$f:X\to Y$，设 $\widehat f:\widehat X\to\widehat Y$ 是 $f$ 在完备化上的
延拓，则映射 $f\mapsto\widehat f$ 是双射
\[
\operatorname{Hom}_S(X,Y)\xrightarrow{\ \sim\ }
\operatorname{Hom}_{\widehat S}(\widehat X,\widehat Y).
\]
\end{theorem}

\begin{proof}
先证 $f\mapsto\widehat f$ 是单射。事实上，设 $f$、$g$ 是从 $X$ 到 $Y$
的两个 $S$-态射，且 $\widehat f=\widehat g$。由（I，10.9.4）可知，
存在 $X'=u^{-1}(S')$ 的一个开邻域 $V$，使得 $f$ 与 $g$ 在其中相同。
但因 $u$ 是闭映射，有 $V=X$（\hyperref[III.5.1.3.1-fr]{5.1.3.1}），
故 $f=g$。

现证 $f\mapsto\widehat f$ 是满射。为此，设 $h$ 是一个
$\widehat S$-态射 $\widehat X\to\widehat Y$。置 $Z=X\times_S Y$，
并以 $p:Z\to X$、$q:Z\to Y$ 表示典范投影；$Z$ 在 $S$ 上有限型
（I，6.3.4），因而是 Noether 的。以 $\widehat Z$ 表示其沿
$Z'=p^{-1}(u^{-1}(S'))$ 的完备化；已知 $\widehat Z$ 可典范地等同于
$\widehat X\times_{\widehat S}\widehat Y$，且投影
$\widehat Z\to\widehat X$、$\widehat Z\to\widehat Y$ 分别等同于延拓
$\widehat p$ 与 $\widehat q$（I，10.9.7）。由于 $Y$ 在 $S$ 上分离，
$\widehat Y$ 在 $\widehat S$ 上分离（I，10.15.7），故图态射
\[
\Gamma_h=(1_{\widehat X},h):\widehat X\to\widehat Z
\]
是闭浸入（I，10.15.4）。设 $\mathfrak T$ 是与此浸入相伴的
$\widehat Z$ 的闭形式子预概形，$j:\mathfrak T\to\widehat Z$ 是典范嵌入，
从而 $\Gamma_h=j\circ w$，其中 $w:\widehat X\to\mathfrak T$ 是同构
（I，10.14.3），其逆同构为 $\widehat p\circ j$；此外，
$\mathfrak T$ 显然在 $\widehat S$ 上固有，因为 $\widehat X$ 如此。
由（\hyperref[III.5.1.8-fr]{5.1.8}）可知，存在 $Z$ 的闭子预概形 $T$，
使得 $\mathfrak T=\widehat T=T_{/(T\cap Z')}$，并且 $j=\widehat i$，
其中 $i$ 是典范嵌入 $T\to Z$（I，10.14.7）。此时
$p\circ i:T\to X$ 是同构，因为依假设，它在完备化上的延拓
$\widehat{p\circ i}=\widehat p\circ\widehat i$ 是同构；只须应用

\oldpage[III]{157}

（4.6.8），并注意到如上所述，$S$ 是 $S$ 中 $S'$ 的唯一邻域。
设 $g:X\to T$ 是 $p\circ i$ 的逆同构，并置 $f=q\circ i\circ g$；
这是态射 $X\to Y$，依定义其图为 $\Gamma_f=i\circ g$。由于
$\widehat g$ 是 $\widehat{p\circ i}=w$ 的逆同构，故
\[
\widehat{\Gamma_f}=\widehat i\circ\widehat g=j\circ w=\Gamma_h.
\]
但已知 $\widehat{\Gamma_f}=\Gamma_{\widehat f}$（I，10.9.8），
因而最终有 $h=\widehat f$，证明完成。
\end{proof}

因此，用范畴语言来说，对任意 Noether 进制环 $A$，函子
$X\mapsto\widehat X$ 从 $\operatorname{Spec}(A)$ 上固有概形的范畴到
$\operatorname{Spf}(A)$ 上固有形式概形的范畴是全忠实的
（$0$，8.1.6）；从而它建立第一个范畴与第二个范畴的一个全子范畴之间的
等价。后一子范畴的对象称为\emph{可代数化形式概形}。对于这样的形式概形
$\mathfrak X$，存在 $\operatorname{Spec}(A)$ 上固有的通常概形 $X$，
它在唯一同构意义下唯一确定，并且 $\mathfrak X$ 同构于 $\widehat X$。

\begin{env}[5.4.2]
\label{III.5.4.2-fr}
\emph{评注。}
沿用（\hyperref[III.5.4.1-fr]{5.4.1}）的记号，置
$S_n=\operatorname{Spec}(A/\mathfrak J^{n+1})$、
$X_n=X\times_S S_n$、$Y_n=Y\times_S S_n$。由
（\hyperref[III.5.4.1-fr]{5.4.1}）和（I，10.12.3）可知，给定一个
$S$-态射 $f:X\to Y$，等价于给定一个由 $S_n$-态射
$f_n:X_n\to Y_n$ 构成的 $(S_n)$-进制归纳系统（I，10.12.2）。
\end{env}

\begin{remark}[5.4.3]
\label{III.5.4.3-fr}
与存在定理（\hyperref[III.5.1.6-fr]{5.1.6}）可能给人的印象相反，
存在 $\operatorname{Spf}(A)$ 上固有但不可代数化的形式概形（正如存在并非
来自复代数簇的紧致解析空间）。以后我们将在“模空间理论”中遇到这样的
概形；当基域为 $\mathbb C$ 时，该理论所研究的恰是完备代数簇复结构的
无穷小变形，而已知这种变形可能产生非代数的解析簇。
\end{remark}

\begin{proposition}[5.4.4]
\label{III.5.4.4-fr}
设 $A$ 是 Noether 进制环，$\mathfrak S=\operatorname{Spf}(A)$，
$g:\mathfrak X\to\mathfrak S$、$h:\mathfrak Y\to\mathfrak S$ 是形式概形的
两个固有态射，$f:\mathfrak X\to\mathfrak Y$ 是 $\mathfrak S$-态射。
若 $f$ 是有限态射且 $\mathfrak Y$ 可代数化，则 $\mathfrak X$ 可代数化。
\end{proposition}

\begin{proof}
先注意，关于 $g$ 和 $h$ 的假设已经蕴含 $f$ 固有（3.4.1）；要使 $f$
有限，只须对每个 $y\in\mathfrak Y$，纤维 $f^{-1}(y)$ 都是有限的
（\hyperref[III.4.8.11-fr]{4.8.11}）。由假设，
$\mathcal B=f_*(\mathcal O_{\mathfrak X})$ 是凝聚
$\mathcal O_{\mathfrak Y}$-代数（\hyperref[III.4.8.6-fr]{4.8.6}）。
因此，由存在定理，若 $\mathfrak Y=\widehat Y$ 且 $h=\widehat w$，其中
$w:Y\to\operatorname{Spec}(A)$ 是通常概形的固有态射，则存在凝聚
$\mathcal O_Y$-代数 $\mathcal C$，使得 $\mathcal B=\widehat{\mathcal C}$。
置 $X=\operatorname{Spec}(\mathcal C)$，并以 $u:X\to Y$ 表示结构态射；
于是，由 $\mathcal B$ 构造 $\mathfrak X$ 的定义
（\hyperref[III.4.8.7-fr]{4.8.7}）立即可知，$\mathfrak X$ 典范同构于
$\widehat X$，并且 $f$ 等同于 $\widehat u$（只须在 $Y$ 仿射时验证）。

注意，（\hyperref[III.5.1.8-fr]{5.1.8}）是
（\hyperref[III.5.4.4-fr]{5.4.4}）的特例。
\end{proof}

\begin{theorem}[5.4.5]
\label{III.5.4.5-fr}
设 $A$ 是 Noether 进制环，$\mathfrak J$ 是 $A$ 的定义理想，
$S=\operatorname{Spec}(A)$，
$\mathfrak S=\widehat S=\operatorname{Spf}(A)$，
$f:\mathfrak X\to\mathfrak S$ 是形式概形的固有态射。置
$S_k=\operatorname{Spec}(A/\mathfrak J^{k+1})$、
$X_k=\mathfrak X\times_{\mathfrak S}S_k$，并且对任意
$\mathcal O_{\mathfrak X}$-模 $\mathcal F$，置
$\mathcal F_k=\mathcal F\otimes_{\mathcal O_{\mathfrak X}}\mathcal O_{X_k}
=\mathcal F/\mathfrak J^{k+1}\mathcal F$。

\oldpage[III]{158}

设 $\mathcal L$ 是可逆 $\mathcal O_{\mathfrak X}$-模，并假设
$\mathcal L_0=\mathcal L/\mathfrak J\mathcal L$ 是丰沛
$\mathcal O_{X_0}$-模。则 $\mathfrak X$ 可代数化；并且，若 $X$ 是
固有 $S$-概形，使得 $\mathfrak X$ 同构于 $\widehat X$，则存在丰沛
$\mathcal O_X$-模 $\mathcal M$，使得 $\mathcal L$ 同构于
$\widehat{\mathcal M}$（这蕴含 $X$ 在 $S$ 上是射影的）。
\end{theorem}

\begin{proof}
将（\hyperref[III.5.2.3-fr]{5.2.3}）应用于
$\mathcal F=\mathcal O_{\mathfrak X}$：于是存在整数 $n_0$，使得对
$n\geq n_0$，典范同态
$\Gamma(\mathfrak X,\mathcal L^{\otimes n})\to
\Gamma(X_0,\mathcal L_0^{\otimes n})$ 是满射。可假设所选的
$n\geq n_0$ 足够大，使 $\mathcal L_0^{\otimes n}$ 对 $S_0$ 极丰沛
（II，4.5.10）。由于态射 $f_0:X_0\to S_0$ 固有，
$\Gamma(X_0,\mathcal L_0^{\otimes n})$ 是有限型 $A$-模（3.2.1），
故 $\Gamma(\mathfrak X,\mathcal L^{\otimes n})$ 中存在有限型 $A$-子模
$E$，其在 $\Gamma(X_0,\mathcal L_0^{\otimes n})$ 中的像是后一个模的全部。
于是，对每个 $k\geq0$，考虑由典范嵌入
$E\to\Gamma(\mathfrak X,\mathcal L^{\otimes n})$ 导出的同态
\[
u_k:E/\mathfrak J^{k+1}E\longrightarrow
\Gamma(X_k,\mathcal L_k^{\otimes n})
\]
。注意，$(f_k)_*(\mathcal L_k^{\otimes n})$ 是拟凝聚的，并且
$\Gamma(S_k,(f_k)_*(\mathcal L_k^{\otimes n}))
=\Gamma(X_k,\mathcal L_k^{\otimes n})$，所以 $u_k$ 定义同态
\[
\widetilde u_k:(E/\mathfrak J^{k+1}E)^\sim
\longrightarrow(f_k)_*(\mathcal L_k^{\otimes n}),
\]
从而还定义同态
\[
\widetilde u_k^\sharp:
f_k^*((E/\mathfrak J^{k+1}E)^\sim)\longrightarrow
\mathcal L_k^{\otimes n}.
\]
此外，若置 $\mathcal G_k=f_k^*((E/\mathfrak J^{k+1}E)^\sim)$，则有
$\mathcal G_0=\mathcal G_k/\mathfrak J\mathcal G_k$（I，9.1.5），故
$\widetilde u_0^\sharp:\mathcal G_k/\mathfrak J\mathcal G_k
\to\mathcal L_k^{\otimes n}/\mathfrak J\mathcal L_k^{\otimes n}$ 是由
$\widetilde u_k^\sharp$ 取商而得的。依 $E$ 的定义，
$\widetilde u_0^\sharp$ 正是典范同态
\[
\sigma:f_0^*((f_0)_*(\mathcal L_0^{\otimes n}))\longrightarrow
\mathcal L_0^{\otimes n},
\]
而 $\mathcal L_0^{\otimes n}$ 极丰沛的假设蕴含
$\widetilde u_0^\sharp$ 是满射（II，4.4.3）；由 Nakayama 引理可知，
每个 $\widetilde u_k^\sharp$ 也都是满射。

因此，每个 $\widetilde u_k^\sharp$ 都定义一个 $S_k$-态射
（II，4.2.2）
\[
g_k:X_k\longrightarrow P_k=\mathbf P(E/\mathfrak J^{k+1}E),
\]
又因（II，4.1.3），当 $h\leq k$ 时有
$P_h=P_k\times_{S_k}S_h$；于是，由各 $\widetilde u_k^\sharp$ 之间的关系
及（II，4.2.10），$(g_k)$ 是一个 $(S_k)$-进制归纳系统
（I，10.12.2）。所以这些 $g_k$ 定义形式概形的 $\mathfrak S$-态射
$g:\mathfrak X\to\mathfrak P$，其中 $\mathfrak P$ 是系统 $(P_k)$ 的
归纳极限，亦即 $P=\mathbf P(E)$ 的完备化 $\widehat P$。此外，
$\mathcal L_0^{\otimes n}$ 极丰沛的假设蕴含 $g_0$ 是闭浸入
（II，4.4.3）；由此可知 $g$ 是形式概形的闭浸入（4.8.10），所以
$\mathfrak X$ 可代数化（5.1.8）。由存在定理（5.1.6）可知，
$\mathcal L$ 同构于某个可逆 $\mathcal O_X$-模的完备化
$\widehat{\mathcal M}$。又 $\mathcal L^{\otimes n}$ 是
$\mathcal M^{\otimes n}$ 的完备化（I，10.8.10），且各同态
$\widetilde u_k^\sharp$ 定义唯一确定的同态
\[
v:f^*(\widetilde E)\longrightarrow\mathcal M^{\otimes n}
\]
（5.1.7）；由于 $\widetilde u_k^\sharp$ 都是满射，$\widehat v$ 也是满射
（I，10.11.5），因而 $v$ 也是满射（5.1.3）。此外，由 $v$ 定义的态射
$r:X\to P$（II，4.2.2）在完备化上的延拓是 $g$；由于 $g$ 是闭浸入，
由（5.1.8）和（5.4.1）可知 $r$ 也是闭浸入。因此
$\mathcal M^{\otimes n}$ 极丰沛（II，4.4.6），而 $\mathcal M$ 丰沛
（II，4.5.10）。
\end{proof}

\begin{remark}[5.4.6]
\label{III.5.4.6-fr}
设 $A$ 是 Noether 进制环，$\mathfrak S=\operatorname{Spf}(A)$，
$f:\mathfrak X\to\mathfrak S$ 是形式概形的固有态射。设 $\mathcal N$
是 $\mathcal O_{\mathfrak X}$ 的凝聚理想，使得对 $\mathfrak X$ 的每个
仿射形式开集 $U$，$\Gamma(U,\mathcal N)$ 都是
$\Gamma(U,\mathcal O_{\mathfrak X})$ 的幂零根；此理想的存在性容易由
（I，10.10.2）以及任意环同态 $B\to C$ 都把 $B$ 的幂零根映入 $C$ 的
幂零根这一事实得出。设 $\mathfrak X'$ 是由 $\mathcal N$ 定义的
$\mathfrak X$ 的闭形式子概形（I，10.14.2）；一个值得研究的问题是：
当 $\mathfrak X'$ 可代数化时，$\mathfrak X$ 本身是否也可代数化。
若能对一个通常预概形或形式预概形的结构层 $\mathcal O_{\mathfrak X}$
以平方零理想为核的\emph{扩张}进行分类，例如借助
某些上同调性质的

\oldpage[III]{159}

不变量，或许就能解决这个问题；换言之，需要分类的是满足下述条件的
$\mathcal O_{\mathfrak X}$-代数 $\mathcal A$：
$\mathcal O_{\mathfrak X}$ 同构于 $\mathcal A/\mathcal J$，其中
$\mathcal J$ 是 $\mathcal A$ 的平方零理想。
\end{remark}

\subsection{某些概形的一种分解}
\label{subsection:III.5.5-fr}

\begin{proposition}[5.5.1]
\label{III.5.5.1-fr}
设 $A$ 是 Noether 进制环，$\mathfrak J$ 是 $A$ 的定义理想，
$Y=\operatorname{Spec}(A)$。设 $f:X\to Y$ 是有限型分离态射；置
$Y_0=\operatorname{Spec}(A/\mathfrak J)$，
$X_0=X\times_Y Y_0=f^{-1}(Y_0)$。设 $Z_0$ 是 $X_0$ 的一个开子集，
且在 $Y_0$ 上固有；则 $X$ 中存在开闭子集 $Z$，它在 $Y$ 上固有，
并满足 $Z\cap X_0=Z_0$。
\end{proposition}

\begin{proof}
由假设，$X$ 中存在开集 $T$，使得 $T\cap X_0=Z_0$。以 $\widehat T$
表示由 $X$ 在开集 $T$ 上诱导的概形沿 $Z_0$ 的完备化；由于
$\mathcal O_{\widehat T}$ 的支集是 $Z_0$，而后者在 $Y_0$ 上固有，
故 $\widehat T$ 在 $\widehat Y=\operatorname{Spf}(A)$ 上固有（3.4.1）。
由（5.1.8），存在 $T$ 的闭子概形 $Z$，它在 $Y$ 上固有，并且若
$i:Z\to T$ 是典范嵌入，则 $\widehat i:\widehat Z\to\widehat T$ 是同构
（这里 $\widehat Z$ 是 $Z$ 沿 $Z_0$ 的完备化）。由（4.6.8）可知，
$T$ 中存在 $Z_0$ 的开邻域 $V$，使得 $i$ 的限制
$i^{-1}(V)\to V$ 是同构。然而，$i^{-1}(V)$ 是 $Z$ 中 $Z_0$ 的邻域，
故必等于 $Z$（5.1.3.1）。因此 $Z$ 在 $T$ 中、从而在 $X$ 中是开集，
证明完毕。
\end{proof}

\begin{corollary}[5.5.2]
\label{III.5.5.2-fr}
若 $X_0$ 在 $Y_0$ 上固有，则 $X$ 是两个不交开子集 $Z$ 与 $Z'$ 的并，
其中 $Z$ 在 $Y$ 上固有并包含 $X_0$；此外，$X$ 的任一在 $Y$ 上固有的
闭子集 $P$ 都包含于 $Z$。
\end{corollary}

\begin{proof}
最后一个断言由下述事实得出：$P\cap Z'$ 在 $P$ 中是闭的，因而在
$Y$ 上固有；若 $P\cap Z'$ 非空，则 $f(P\cap Z')$ 是 $Y$ 中非空闭集，
故与 $Y_0$ 相交（5.1.3.1），这与 $Z$ 的定义矛盾。
\end{proof}

\begin{flushright}
\emph{（待续。）}
\end{flushright}

% 印刷页 504 为空白页。
\clearpage
\thispagestyle{empty}
\null
\clearpage
\thispagestyle{plain}

\oldpage[III]{161}
\section*{参考文献（续）}

\begin{flushleft}
[27] P. Gabriel, \emph{Des catégories abéliennes}, thèse, Paris (1961).

\medskip
[28] J. E. Roos, \emph{Sur les foncteurs dérivés de $\varprojlim$.
Applications}, C. R. Acad. Sci. Paris, t. CCLII, 1961, p. 3702--3704.

\medskip
[29] H. Cartan, \emph{Séminaire de l'École Normale Supérieure},
13\textsuperscript{e} année (1960--1961), exposé n\textsuperscript{o} 11.
\end{flushleft}

\clearpage
\oldpage[III]{162}
\section*{符号索引}

\begin{flushleft}
$\mathbf{Ens}$：0，8.1.1。

$h_X(Y)$，$h_X(u)$（$X$、$Y$ 为范畴 $C$ 的对象，$u$ 为 $C$ 中的态射）：
0，8.1.1。

$h_w(Y)$（$Y$ 为 $C$ 的对象，$w$ 为 $C$ 中的态射）：0，8.1.2。
\SourceCorrectionNote{法文印本及外交式符号索引将此条误记为以态射 u 为自变量；
当前 Canon 依 0，8.1.2 中的定义改为以对象 Y 为自变量。}

$Z_k(E_r^{pq})$，$B_k(E_r^{pq})$（$k\geq r$）：0，11.1.1。

$Z_\infty(E_2^{pq})$，$B_\infty(E_2^{pq})$：0，11.1.1。

$K^\bullet$，$K_\bullet$（复形）：0，11.2.1。

$K^{\bullet\bullet}$，$K_{\bullet\bullet}$（双复形）：0，11.2.1。

$B^i(K^\bullet)$，$Z^i(K^\bullet)$，$H^i(K^\bullet)$，
$H^\bullet(K^\bullet)$：0，11.2.1。

$B_i(K_\bullet)$，$Z_i(K_\bullet)$，$H_i(K_\bullet)$，
$H_\bullet(K_\bullet)$：0，11.2.1。

$T(K^\bullet)$，$T(K_\bullet)$（$T$ 为函子）：0，11.2.1。

$E(K^\bullet)$（$K^\bullet$ 为滤过复形）：0，11.2.2。

$K^{i,\bullet}$，$K^{\bullet,j}$，$Z_{\mathrm{II}}^p(K^{i,\bullet})$，
$B_{\mathrm{II}}^p(K^{i,\bullet})$，$H_{\mathrm{II}}^p(K^{i,\bullet})$，
$Z_{\mathrm I}^p(K^{\bullet,j})$，$B_{\mathrm I}^p(K^{\bullet,j})$，
$H_{\mathrm I}^p(K^{\bullet,j})$：0，11.3.1。

$H_{\mathrm{II}}^p(K^{\bullet\bullet})$，
$Z_{\mathrm I}^q(H_{\mathrm{II}}^p(K^{\bullet\bullet}))$，
$B_{\mathrm I}^q(H_{\mathrm{II}}^p(K^{\bullet\bullet}))$，
$H_{\mathrm I}^q(H_{\mathrm{II}}^p(K^{\bullet\bullet}))$，
$H^n(K^{\bullet\bullet})$：0，11.3.1。

$F_{\mathrm I}^p(K^{\bullet\bullet})$，
$F_{\mathrm{II}}^p(K^{\bullet\bullet})$，${}'E(K^{\bullet\bullet})$，
${}''E(K^{\bullet\bullet})$：0，11.3.2。

$K_{i,\bullet}$，$K_{\bullet,j}$，$H_p^{\mathrm{II}}(K_{i,\bullet})$，
$H_p^{\mathrm I}(K_{\bullet,j})$，$H_p^{\mathrm{II}}(K_{\bullet\bullet})$，
$H_p^{\mathrm I}(K_{\bullet\bullet})$，
$H_q^{\mathrm I}(H_p^{\mathrm{II}}(K_{\bullet\bullet}))$，
$H_q^{\mathrm{II}}(H_p^{\mathrm I}(K_{\bullet\bullet}))$，
$H_n(K_{\bullet\bullet})$：0，11.3.5。

$\mathrm R^\bullet T(K^\bullet)$（$T$ 为协变加性函子）：0，11.4.3。

$\mathrm R^\bullet T(K^\bullet,K'^{\bullet})$（$T$ 为协变加性双函子）：
0，11.4.6。

$\mathrm L_\bullet T(K_\bullet)$（$T$ 为协变加性函子）：0，11.6.2。

$\mathrm L_\bullet T(K_\bullet,K'_\bullet)$（$T$ 为协变加性双函子）：
0，11.6.6。

$\mathrm L_\bullet T(K_{\bullet\bullet})$（$T$ 为协变加性函子）：
0，11.7.2。

$|\sigma|$，$\Sigma(A)$，$\mathrm C_\bullet(A)$，$\mathrm L_\bullet(A)$
（$A$ 为有限集，$\sigma$ 为 $A$ 的单形）：0，11.8.1。

$\mathrm C^\bullet(A,B;\mathscr S)$，$\mathrm L^\bullet(A,B;\mathscr S)$
（$A$、$B$ 为有限集，$\mathscr S$ 为系数系统）：0，11.8.4。

$\mathrm P^\bullet(A,B,\mathscr S)$（$A$、$B$ 为有限集，$\mathscr S$
为系数系统）：0，11.8.6。

$\mathrm C^\bullet(A,B;\mathscr S^\bullet)$，
$\mathrm L^\bullet(A,B;\mathscr S^\bullet)$（$A$、$B$ 为有限集，
$\mathscr S^\bullet$ 为系数系统的复形）：0，11.8.8。

$\mathrm P^\bullet(A,B;\mathscr S^\bullet)$（$A$、$B$ 为有限集，
$\mathscr S^\bullet$ 为系数系统的复形）：0，11.8.10。

$\chi(M^\bullet)$（$M^\bullet$ 为有限长度模的有限复形）：0，11.10.1。

$\check C^\bullet(\mathfrak U,\mathcal F)$（$\mathcal F$ 为 $X$ 上的
Abel 群层，$\mathfrak U$ 为 $X$ 的开覆盖）：0，12.1.2。

$\mathrm R^pf_*(\mathcal F)$（$f:X\to Y$ 为环化空间的态射，
$\mathcal F$ 为 $\mathcal O_X$-模层）：0，12.2.1。

$\mathcal H^p(f,\mathcal K^\bullet)$，
$\mathcal H_f^p(\mathcal K^\bullet)$，${}'E(f,\mathcal K^\bullet)$，
${}''E(f,\mathcal K^\bullet)$（$f:X\to Y$ 为环化空间的态射，
$\mathcal K^\bullet$ 为 $\mathcal O_X$-模复形）：0，12.4.1。

$\mathbf H^\bullet(X,\mathcal K^\bullet)$，
$\mathbf H^\bullet(V,\mathcal K^\bullet)$（$X$ 为环化空间，
$\mathcal K^\bullet$ 为 $\mathcal O_X$-模复形，$V$ 为 $X$ 的开集）：
0，12.4.2。

$\operatorname{gr}^\bullet(A)$（$A$ 为满足（ML）的逆系统）：0，13.4.2。

$E(A)$（$A$ 为滤过对象）：0，13.6.4。

$K_\bullet(\mathbf f)$，$i_{\mathbf f}$（$\mathbf f$ 为环中元素的有限族）：
III，1.1.1。

$K^\bullet(\mathbf f,M)$，$K_\bullet(\mathbf f,M)$（$\mathbf f$ 为环
$A$ 中元素的有限族，$M$ 为 $A$-模）：III，1.1.2。

$H^\bullet(\mathbf f,M)$，$H_\bullet(\mathbf f,M)$（$\mathbf f$ 为环
$A$ 中元素的有限族，$M$ 为 $A$-模）：III，1.1.3。

$C^\bullet((\mathbf f),M)$，$H^\bullet((\mathbf f),M)$（$\mathbf f$ 为环
$A$ 中元素的有限族，$M$ 为 $A$-模）：III，1.1.6。

$H^\bullet(U,\mathcal F^{(*)})$，
$H^\bullet(\mathfrak U,\mathcal F^{(*)})$（$\mathcal F$ 为
拟凝聚 $\mathcal O_X$-模，$U$ 为 $X$ 的开集，$\mathfrak U$ 为 $X$ 的
开覆盖）：III，2.1.1。
\end{flushleft}

\clearpage
\oldpage[III]{163}
\section*{术语索引}

\begin{flushleft}
谱序列的收敛项：0，11.1.1。

可代数化（$\mathcal O_X$-模）：III，5.2.1。

可代数化（形式概形）：III，5.4.2。

解析整（环）：III，4.3.6。

拟紧映射：0，9.1.1。

消解的增广：0，11.4.1。

复形的 Euler--Poincaré 特征：0，11.10.1。

凝聚 $\mathcal O_X$-模的 Euler--Poincaré 特征（$X$ 在
$\operatorname{Spec}(A)$ 上射影，$A$ 为 Artin 环）：III，2.5.1。

双交错上链：0，11.8.4。

由双复形定义的复形：0，11.3.1。

外代数复形：III，1.1.1。

（ML）条件：0，13.1.2。

可构造（子集、集合）：0，9.1.2。

可构造（函数）：0，9.3.1。

杯积：0，12.1.2。

范畴中代数结构的双同态：0，8.2.1。

Abel 范畴中的正合（子集）：III，3.1.1。

本质常值（逆系统）：0，13.4.2。

Stein 分解：III，4.3.3。

余离散、余分离、离散、穷尽、有限、分离滤过：0，11.1.3。

范畴的终对象：0，8.1.10。

形式预概形的有限态射：III，4.8.2。

在 $\mathfrak Y$ 上有限的形式预概形、有限形式
$\mathfrak Y$-预概形：III，4.8.2。

典范协变函子 $C\to\operatorname{Hom}(C^o,\mathbf{Ens})$：0，8.1.2。

算术亏格：III，2.5.1。

几何连通（纤维）：III，4.3.4。

范畴上代数结构的同态：0，8.2.1。

函子关于复形 $K^\bullet$ 的超上同调：0，11.4.3。

双函子关于两个复形 $K^\bullet$、$K'^{\bullet}$ 的超上同调：
0，11.4.6。

函子关于复形 $K_\bullet$ 的超同调：0，11.6.2。

双函子关于两个复形 $K_\bullet$、$K'_\bullet$ 的超同调：0，11.6.6。

函子关于双复形的超同调：0，11.7.2。

局部可构造（集合）：0，9.1.11。

范畴对象上的外（内）合成律：0，8.2.1。

\oldpage[III]{164}

滤过 $S$-$C$-模、分次 $\operatorname{gr}^\bullet(S)$-$C$-模、
双分次 $\operatorname{gr}^{\bullet\bullet}(S)$-$C$-模：0，13.6.6。

谱序列的态射：0，11.1.2。

纤维的几何连通分支数：III，4.3.4。

$C$-群对象、$C$-群、$C$-环、$C$-模：0，8.2.3。

边界对象、余边界对象、余循环对象、循环对象：0，11.2.1。

普遍像对象：0，13.1.1。

满足（ML）的逆系统的相伴分次对象：0，13.4.2。

下（上）有界分次对象：0，11.2.1。

形式预概形在 $\mathfrak Y$ 上的固有子集：III，3.4.1。

全（子范畴）：0，8.1.5。

全忠实（函子）：0，8.1.5。

Hilbert 多项式：III，2.5.3。

形式预概形的固有态射：III，3.4.1。

可表（函子）：0，8.1.8。

上同调、右、左、同调、内射、自由、平坦、投射消解：0，11.4.1。

右、内射 Cartan--Eilenberg 消解：0，11.4.2。

左、投射 Cartan--Eilenberg 消解：0，11.6.1。

逆紧（集合）：0，9.1.1。

严格（逆系统）：0，13.4.2。

Abel 范畴中的谱序列：0，11.1.1。

弱收敛、正则、余正则、双正则谱序列：0，11.1.3。

退化谱序列：0，11.1.6。

函子相对于滤过对象的谱序列：0，13.6.4。

双复形的谱序列：0，11.3.2 与 11.3.5。

函子关于复形的超上同调谱序列：0，11.4.3。

函子关于复形的超同调谱序列：0，11.6.2。

函子关于双复形的超同调谱序列：0，11.7.2。

系数系统：0，11.8.4。

单支（环、点）：III，4.3.6。

普遍开（态射）：III，4.3.9。
\end{flushleft}

\clearpage
\oldpage[III]{165}
\section*{目录}

\begin{flushleft}
\textsc{第 0 章。} --- \textbf{预备知识（续）} \dotfill 5

\medskip
\S~8. 可表函子 \dotfill 5

8.1. 可表函子 \dotfill 5

8.2. 范畴中的代数结构 \dotfill 9

\medskip
\S~9. 可构造集 \dotfill 12

9.1. 可构造集 \dotfill 12

9.2. Noether 空间中的可构造集 \dotfill 14

9.3. 可构造函数 \dotfill 16

\medskip
\S~10. 关于平坦模的补充 \dotfill 17

10.1. 平坦模与自由模之间的关系 \dotfill 17

10.2. 平坦性的局部判据 \dotfill 18

10.3. 局部环的平坦扩张的存在性 \dotfill 20

\medskip
\S~11. 同调代数补充 \dotfill 23

11.1. 谱序列回顾 \dotfill 23

11.2. 滤过复形的谱序列 \dotfill 27

11.3. 双复形的谱序列 \dotfill 29

11.4. 函子关于复形 $K^\bullet$ 的超上同调 \dotfill 32

11.5. 超上同调中取余极限 \dotfill 35

11.6. 函子关于复形 $K_\bullet$ 的超同调 \dotfill 39

11.7. 函子关于双复形 $K_{\bullet\bullet}$ 的超同调 \dotfill 41

11.8. 单纯复形上同调的补充 \dotfill 43

11.9. 有限型复形的一个引理 \dotfill 46

11.10. 有限长度模复形的 Euler--Poincaré 特征 \dotfill 48

\medskip
\S~12. 层上同调补充 \dotfill 49

12.1. 赋环空间上模层的上同调 \dotfill 49

12.2. 高阶直接像 \dotfill 57

12.3. 模层的 Ext 函子补充 \dotfill 60

12.4. 直接像函子的超上同调 \dotfill 62

\oldpage[III]{166}

\S~13. 同调代数中的逆极限 \dotfill 64

13.1. Mittag--Leffler 条件 \dotfill 64

13.2. Abel 群的 Mittag--Leffler 条件 \dotfill 65

13.3. 应用：层的逆极限的上同调 \dotfill 68

13.4. Mittag--Leffler 条件与逆系统的相伴分次对象 \dotfill 69

13.5. 滤过复形的谱序列的逆极限 \dotfill 71

13.6. 相对于带有限滤过对象的函子的谱序列 \dotfill 73

13.7. 逆极限自变量的导函子 \dotfill 75

\medskip
\textsc{第 III 章。} --- \textbf{凝聚层的上同调研究} \dotfill 81

\medskip
\S~1. 仿射概形的上同调 \dotfill 82

1.1. 外代数复形的回顾 \dotfill 82

1.2. 开覆盖的 Čech 上同调 \dotfill 85

1.3. 仿射概形的上同调 \dotfill 88

1.4. 对任意预概形上同调的应用 \dotfill 89

\medskip
\S~2. 射影态射的上同调研究 \dotfill 95

2.1. 若干上同调群的显式计算 \dotfill 95

2.2. 射影态射的基本定理 \dotfill 100

2.3. 分次代数层与模层的应用 \dotfill 102

2.4. 基本定理的一个推广 \dotfill 107

2.5. Euler--Poincaré 特征与 Hilbert 多项式 \dotfill 109

2.6. 应用：丰沛性判别准则 \dotfill 111

\medskip
\S~3. 固有态射的有限性定理 \dotfill 115

3.1. 逐层归约（dévissage）引理 \dotfill 115

3.2. 有限性定理：通常概形的情形 \dotfill 116

3.3. 有限性定理的推广（通常概形） \dotfill 118

3.4. 有限性定理：形式概形的情形 \dotfill 119

\medskip
\S~4. 固有态射的基本定理及其应用 \dotfill 122

4.1. 基本定理 \dotfill 122

4.2. 特殊情形与变式 \dotfill 129

4.3. Zariski 连通性定理 \dotfill 130

4.4. Zariski“主定理” \dotfill 135

4.5. 同态模的完备化 \dotfill 138

4.6. 形式态射与通常态射之间的关系 \dotfill 139

4.7. 一个丰沛性判据 \dotfill 145

4.8. 形式预概形的有限态射 \dotfill 146

\oldpage[III]{167}

\S~5. 凝聚代数层的一个存在定理 \dotfill 149

5.1. 定理的陈述 \dotfill 149

5.2. 存在定理的证明：射影与拟射影情形 \dotfill 151

5.3. 存在定理的证明：一般情形 \dotfill 154

5.4. 应用：通常概形态射与形式概形态射的比较。可代数化形式概形
\dotfill 156

5.5. 某些概形的一种分解 \dotfill 159

\medskip
\textsc{参考文献（续）} \dotfill 161

\textsc{符号索引} \dotfill 162

\textsc{术语索引} \dotfill 163

\begin{flushright}
\emph{1961 年 6 月 28 日收稿。}
\end{flushright}
\end{flushleft}
\end{document}
